% ============================================================
%  定位与导航技术文档  ——  主文件（由 Lead 维护，撰稿人请勿修改）
%  编译：xelatex main -> biber main -> xelatex main -> xelatex main
% ============================================================
\documentclass[lang=cn,newtx,10pt,scheme=chinese,color=blue]{elegantbook}

\title{定位与导航技术文档}
\subtitle{从里程计到 SLAM，从建图到自主导航}

\author{算法组}
\institute{RoboMaster 算法组}
\date{2026 年 11 月}
\version{1.0}
\bioinfo{适用对象}{算法组新队员 / 进阶队员}
\extrainfo{配套解析：FAST-LIO2 \textperiodcentered{} Point-LIO \textperiodcentered{} Cartographer \textperiodcentered{} Nav2 \textperiodcentered{} Z-Localization}

% 封面与徽标（由 tools/make_cover.py 生成）
\logo{figures/logo.png}
\cover{figures/cover.png}

% 主题配色：深蓝为主色
\definecolor{structurecolor}{RGB}{22,93,155}
\definecolor{main}{RGB}{0,110,120}
\definecolor{second}{RGB}{190,80,20}
\definecolor{third}{RGB}{0,145,170}
\colorlet{coverlinecolor}{structurecolor}

\setcounter{tocdepth}{2}

% 允许章节标题中出现数学宏（\bm、\frac 等）并安全写入 PDF 书签
\pdfstringdefDisableCommands{%
  \def\bm#1{#1}%
  \def\boldsymbol#1{#1}%
  \def\frac#1#2{#1/#2}%
  \def\mathrm#1{#1}%
  \def\text#1{#1}%
}

% 代码块样式：浅色底、细边框、彩色关键字
\lstset{
  basicstyle=\ttfamily\small,
  keywordstyle=\color{structurecolor}\bfseries,
  commentstyle=\color{green!45!black}\itshape,
  stringstyle=\color{red!55!black},
  frame=single,
  rulecolor=\color{black!18},
  backgroundcolor=\color{black!2},
  framexleftmargin=2mm,
  xleftmargin=2mm,
  aboveskip=5pt,
  belowskip=5pt,
  columns=fullflexible,
  keepspaces=true
}

% 图形库
\usetikzlibrary{arrows.meta, shapes.geometric, shapes.misc, positioning,
  fit, calc, backgrounds, patterns, decorations.pathreplacing}
\usepackage{pgfplots}
\pgfplotsset{compat=1.18}

% ============================================================
%  macros.tex —— 全书公共宏定义（由 Lead 维护，撰稿人请勿修改）
%  使用前已加载：amsmath, amssymb, bm, booktabs, tikz, tcolorbox,
%  listings, enumitem, pgfplots, hyperref, adforn
% ============================================================

% ---------- 数学快捷命令 ----------
\newcommand{\R}{\mathbb{R}}
\newcommand{\SO}{\mathrm{SO}}
\newcommand{\SE}{\mathrm{SE}}
\newcommand{\so}{\mathfrak{so}}
\newcommand{\se}{\mathfrak{se}}
\newcommand{\norm}[1]{\left\lVert #1 \right\rVert}
\newcommand{\abs}[1]{\left\lvert #1 \right\rvert}
\newcommand{\diff}{\mathrm{d}}
\newcommand{\E}{\mathbb{E}}
\DeclareMathOperator{\tr}{tr}
\DeclareMathOperator{\diag}{diag}
\DeclareMathOperator{\Cov}{Cov}
\DeclareMathOperator{\Var}{Var}
\DeclareMathOperator*{\argmin}{arg\,min}
\DeclareMathOperator*{\argmax}{arg\,max}
\newcommand{\degree}{^\circ}
\newcommand{\eg}{e.g.}
\newcommand{\ie}{i.e.}

% ---------- 本章导读（仿 ElegantBook introduction 样式）----------
\newtcolorbox{introbox}{
  enhanced,
  breakable,
  colback=structurecolor!7,
  colframe=structurecolor,
  boxrule=0pt,
  toprule=0.8pt,
  bottomrule=0.8pt,
  sharp corners,
  colbacktitle=structurecolor,
  coltitle=white,
  fonttitle=\bfseries,
  title={本章导读},
  attach boxed title to top center={yshift=-3mm,yshifttext=-1mm},
  boxed title style={sharp corners, boxrule=0pt},
  top=9pt, bottom=7pt, left=3mm, right=3mm,
  before skip=10pt, after skip=10pt
}
\newcommand{\chapterintro}[1]{\begin{introbox}#1\end{introbox}}

% ---------- 要点框 ----------
\newtcolorbox{keybox}[1]{
  enhanced,
  breakable,
  colback=structurecolor!4,
  colframe=structurecolor!55,
  fonttitle=\bfseries\small,
  colbacktitle=structurecolor!12,
  coltitle=structurecolor!80!black,
  title={#1},
  boxrule=0.5pt,
  arc=1.5mm,
  left=2mm, right=2mm, top=1.2mm, bottom=1.2mm,
  before skip=6pt, after skip=6pt,
  drop shadow={opacity=0.10}
}

% ---------- 伪代码环境 ----------
\newtcolorbox{pseudobox}{
  enhanced,
  breakable,
  colback=white,
  colframe=black!25,
  boxrule=0.5pt,
  arc=1mm,
  left=2mm, right=2mm, top=1.5mm, bottom=1.5mm,
  drop shadow={opacity=0.08}
}
\newenvironment{pseudo}[1]{%
  \begin{pseudobox}%
  \noindent\textbf{\color{structurecolor}算法：}\textsf{#1}\par\vspace{1mm}%
  \begin{enumerate}[label=\textbf{\arabic*.}, leftmargin=*, topsep=1pt, itemsep=0.5pt, parsep=0pt]%
}{%
  \end{enumerate}%
  \end{pseudobox}%
}

% ---------- 章末练习（花饰标题，不写入目录）----------
\renewenvironment{problemset}[1][\xchaptertitle~\exercisename]{%
  \vspace*{10pt}%
  \begin{center}%
    \markright{#1}%
    \textcolor{structurecolor}{\Large\bfseries
      \adftripleflourishleft~#1~\adftripleflourishright}%
  \end{center}%
  \begin{enumerate}%
}{%
  \end{enumerate}%
}

% ---------- 本章小结（花饰标题）----------
\newenvironment{summary}{%
  \par\vspace*{12pt}%
  \begin{center}%
    \textcolor{structurecolor}{\Large\bfseries
      \adftripleflourishleft~本章小结~\adftripleflourishright}%
  \end{center}%
  \vspace{-2pt}%
  \noindent\ignorespaces
}{%
  \par\vspace*{4pt}%
}

% ---------- 延伸阅读（花饰标题）----------
\newenvironment{reading}{%
  \par\vspace*{10pt}%
  \begin{center}%
    \textcolor{structurecolor}{\large\bfseries
      \adftripleflourishleft~延伸阅读~\adftripleflourishright}%
  \end{center}%
  \vspace{-2pt}%
  \noindent\ignorespaces
}{%
  \par\vspace*{2pt}%
}

% ---------- 常用 TikZ 样式（可直接在图中使用）----------
\tikzset{
  block/.style   = {rectangle, rounded corners=2pt, draw=structurecolor,
                    thick, fill=structurecolor!8, minimum height=8mm,
                    minimum width=18mm, align=center, font=\small},
  block2/.style  = {rectangle, rounded corners=2pt, draw=structurecolor!70,
                    thick, fill=structurecolor!18, minimum height=8mm,
                    minimum width=18mm, align=center, font=\small},
  arrow/.style   = {-{Stealth[length=2mm]}, thick, draw=black!70},
  darrow/.style  = {{Stealth[length=2mm]}-{Stealth[length=2mm]}, thick,
                    draw=black!70},
  label/.style   = {font=\footnotesize, inner sep=1pt},
}


\addbibresource{reference.bib}

\begin{document}

\maketitle

% 表格斑马纹：只在 table 环境内生效。
% 注意：不能在全局启用 \rowcolors——它会给数学 array（cases/bmatrix/matrix/aligned）
% 的每一行插入颜色盒，导致 \left...\right 定界符无法正常伸缩（括号变小或消失）。
\AtBeginEnvironment{table}{\rowcolors{2}{structurecolor!4}{white}}
\AtEndEnvironment{table}{\rowcolors{1}{}{}}

\frontmatter
\chapter{前言}

\section*{为什么写这本书}

算法组新队员入门定位与导航，最常见的困境不是“找不到资料”，而是“资料太多且互相矛盾”：有人推荐先看《视觉 SLAM 十四讲》，有人让你直接读 FAST-LIO 源码；有人说定位就是卡尔曼滤波，有人说是因子图优化；有人把 Cartographer 当标准答案，有人只用 LIO。结果往往是：视频看了几十个小时，代码 clone 了一堆，真到机器人上还是不知道从哪下手。

这本书试图解决这个问题。它不假设你学过机器人学，也不假设你会 ROS，而是从“机器人怎么知道自己在哪”这个最朴素的问题出发，一层一层把概念、公式、算法、工程串起来。我们尽量做到三件事：
\begin{enumerate}
  \item \textbf{讲清“为什么”}：每个算法先说出它要解决什么问题、在前一个方法上改进了什么，再给公式；
  \item \textbf{推导到能算数}：关键公式从定义开始推，不跳步，并且尽量给一个可以手算的数字例子；
  \item \textbf{落到工程}：每个模块都说明输入输出、频率、参数、失效场景，而不是停在论文层面。
\end{enumerate}

\section*{本书面向谁}

主要读者是 RoboMaster 算法组的新队员，以及任何“会 C++、学过高数线代概率、但没系统学过定位导航”的同学。我们假设你：
\begin{itemize}
  \item 能看懂基本的矩阵运算、求导、概率分布；
  \item 写过一些 C++/Python，能跑通 Linux 命令；
  \item 对机器人有热情，愿意动手。
\end{itemize}
如果你已经做过一两个赛季的定位，这本书的后半部分（SLAM 篇、导航篇、工程篇）仍然值得一读——它会把你在调参中积累的“直觉”补上数学解释。

\section*{怎么读}

全书五篇三十章，建议按顺序读完第一篇（第 1--3 章），之后按需选择：

\begin{itemize}
  \item \textbf{一周速成}：第 1、2、4、6、11、16、17 章。建立“里程计—前端—后端—LIO”的主干。
  \item \textbf{建图方向}：第 1--3、7、11--15 章。重点是 Cartographer、地图表示与回环。
  \item \textbf{导航方向}：第 1、10、11、18--22 章。重点是 Nav2 与 CMU 3D 导航栈。
  \item \textbf{落地实战}：第 23、25--29 章。Z-Localization 的架构、ROS 2 工程、标定与调参。
\end{itemize}

每章开头有“本章导读”，结尾有“本章小结”“思考与练习”“延伸阅读”。建议至少把小结和练习过一遍——练习里的题目大多是“把本章公式代入具体数字”，能算出来才算真懂。

\section*{关于“前端/后端”}

这是本书反复出现的两个词，也是新手最容易混淆的地方。第 1 章先给直觉，第 11 章给出正式定义，并特别用一节讲清楚它们在三种语境下的不同含义：
\begin{itemize}
  \item SLAM 语境：前端造约束，后端解约束；
  \item 定位系统语境：前端是里程计，后端是先验地图配准与因子图；
  \item 导航系统语境：前端是感知（定位+建图），后端是决策（规划+控制）。
\end{itemize}
读任何一章时，请先确认该章用的是哪种语境。第 23 章 Z-Localization 用的是第二种，第 18 章 Nav2 用的是第三种。

\section*{数学约定}

为了让全书符号一致，我们约定：
\begin{itemize}
  \item 向量、矩阵用粗体：$\bm{p},\bm{R}$；标量用普通斜体：$x,\theta$；
  \item 估计值加帽子：$\hat{\bm{x}}$；IMU 原始测量加波浪号：$\tilde{\bm{\omega}}$；
  \item 坐标系：$W$ 世界系、$O$ 里程计系、$B$ 机体系、$L$ 雷达系、$C$ 相机系；
  \item 位姿 $T_{AB}$ 表示“从 $B$ 系到 $A$ 系”的变换，满足 ${}^{A}\bm{p}=T_{AB}\,{}^{B}\bm{p}$；
  \item 重力 $\bm{g}=[0,0,-9.81]^{\top}$，世界系 $z$ 轴向上；
  \item 旋转用右扰动更新：$\bm{R}\leftarrow\bm{R}\exp(\delta\bm{\phi}^{\wedge})$。
\end{itemize}
完整的符号表见第 2 章与附录 A。

\section*{关于开源代码}

本书大量参考了以下开源项目，特此致谢：
\begin{itemize}
  \item \textbf{FAST-LIO / FAST-LIO2 / Point-LIO / VoxelMap}（香港大学 MARS 实验室，Fu Zhang 组）：第 16--17 章的主要素材；
  \item \textbf{Cartographer}（Google）：第 15 章；
  \item \textbf{Nav2}（Open Navigation LLC 与社区）：第 18、22 章；
  \item \textbf{CMU Autonomous Exploration Development Environment}（CMU AirLab，Ji Zhang 组）：第 21 章；
  \item \textbf{Z-Localization}（东北大学 T-DT 战队，作者 mozihe）：第 23 章。该项目的设计说明写得非常坦诚，本章的很多工程判断直接来自作者的分享。
\end{itemize}
书中对代码的描述基于公开文档与论文，若与最新版本有出入，请以官方仓库为准。

\section*{致谢与反馈}

感谢算法组历届队员积累的调试笔记与踩坑记录，它们是本书“工程篇”的素材来源。感谢所有把代码与设计思路开源出来的团队——没有这些开源项目，这本书无从写起。

本书难免有错误与过时之处（尤其涉及具体参数与版本的部分）。如果你发现公式推导有误、代码描述过时、或者某一段实在看不懂，欢迎指出。对一本教材来说，读者的困惑和作者的笔误同样重要。

\tableofcontents

\mainmatter

\part{基础篇：先建立正确的地图}

\chapter{定位、建图与导航：全景导论}\label{ch:01-overview}

\chapterintro{这是全书的“地图”。我们不急着推公式，而是先回答三个最朴素的问题：机器人怎么知道自己在哪里？它怎么知道周围长什么样？它怎么决定往哪走？答案分别对应定位、建图、导航。你会看到这三件事是如何互相纠缠、又如何在工程上被拆开的；也会第一次听到“前端/后端”这两个词，并知道它们在三种不同语境下分别指什么。读完本章，你应该能在脑子里画出整本书的结构，知道后面每一章大概在解决哪一块拼图。}

\section{一个场景：机器人第一次进场地}\label{sec:01-scene}

想象一台刚组装好的哨兵机器人，被放在 RoboMaster 场地的一角。它有一台激光雷达、一个 IMU、四个麦轮，还有一块算力有限的板子。比赛开始前，它需要回答三个问题：

\begin{enumerate}
  \item \textbf{我在哪？} 具体到“我在场地的 $(x,y)$ 坐标、朝向多少度”。这叫\textbf{定位}（localization）。
  \item \textbf{周围长什么样？} 哪里有墙、哪里是通道、哪里可以走。这叫\textbf{建图}（mapping）。
  \item \textbf{我该怎么过去？} 从当前位置到目标点，走哪条路、怎么避开动态的对手。这叫\textbf{导航}（navigation）。
\end{enumerate}

三个问题看起来可以按顺序解决，实际却是一个循环：没有地图，定位只能靠猜；没有定位，建出来的图是歪的；没有可靠的定位和地图，导航就是蒙眼狂奔。本章先把它们的定义和关系理清楚。

\section{三个问题的精确定义}\label{sec:01-def}

\subsection{定位}

\begin{definition}[定位]
定位是估计机器人在某个参考坐标系中的\textbf{位姿}（pose）的过程。位姿包含位置和姿态：
\begin{equation}
\bm{x}=\left(x,\ y,\ z,\ \text{roll},\ \text{pitch},\ \text{yaw}\right),
\end{equation}
在平面移动机器人上通常简化为 $(x,y,\theta)$ 三个量。
\end{definition}

“位姿”这个词在后面会反复出现，务必先建立直觉：位置是“你在哪”，姿态是“你朝哪”。两者缺一不可——只知道位置不知道朝向，机器人不知道往哪走；只知道朝向不知道位置，它不知道自己在场地的哪个角落。

按参考坐标系的不同，定位又分为两类：
\begin{itemize}
  \item \textbf{相对定位}（relative localization）：相对于起点估计位姿，误差会累积。里程计就是典型。
  \item \textbf{绝对定位}（absolute localization）：直接在地图坐标系中估计位姿，误差有界。AMCL、先验地图配准属于这一类。
\end{itemize}

\subsection{建图}

\begin{definition}[建图]
建图是把传感器观测转换成环境表示的过程。环境表示可以是二维占据栅格、三维点云、八叉树、TSDF、语义地图等，取决于用途。
\end{definition}

地图不是“越精细越好”。给导航用的地图要能快速查询“这一点能不能走”；给定位用的地图要能快速做点云配准；给人看的地图要清晰。同一个环境，三种用途往往需要三种不同的地图表示。第 14 章会详细讨论。

\subsection{导航}

\begin{definition}[导航]
导航是在已知（或未知）环境中，从当前位置安全、高效地到达目标位置的过程。一个完整的导航系统通常包含：定位、地图、全局规划、局部规划、运动控制。
\end{definition}

导航不是一个算法，而是一套系统。它把定位和建图的输出当作输入，把速度和转向指令当作输出。第 18 章会给出完整的系统架构。

\subsection{三者的关系}

三者的关系可以用图~\ref{fig:01-triangle} 概括：定位和建图互相依赖（定位需要地图，建图需要位姿），导航同时依赖两者。

\begin{figure}[htbp]
  \centering
  \begin{tikzpicture}[node distance=2.6cm]
    \node[block2, minimum width=28mm] (loc) {定位\\我在哪};
    \node[block, right=of loc, minimum width=28mm] (map) {建图\\环境什么样};
    \node[block2, below=1.9cm of loc, xshift=1.3cm, minimum width=28mm] (nav) {导航\\怎么过去};
    \draw[darrow] (loc) -- node[above, label]{互相依赖} (map);
    \draw[arrow] (loc) -- node[left, label]{位姿} (nav);
    \draw[arrow] (map) -- node[right, label]{地图} (nav);
  \end{tikzpicture}
  \caption{定位、建图、导航三者的关系。定位与建图互为前提，导航同时依赖两者。}
  \label{fig:01-triangle}
\end{figure}

\section{定位问题的分类}\label{sec:01-taxonomy}

初学者最容易混淆的，是“定位”这个词在不同上下文里指的东西完全不同。表~\ref{tab:01-loc-types} 给出四种常见的分类维度。

\begin{table}[htbp]
  \centering
  \caption{定位问题的四种分类维度}\label{tab:01-loc-types}
  \small
  \begin{tabular}{p{2.6cm}p{4.4cm}p{5.6cm}}
    \toprule
    维度 & 两类 & 典型方法 \\
    \midrule
    参考系 & 相对定位 / 绝对定位 & 里程计 / AMCL、先验地图配准 \\
    初始信息 & 局部定位 / 全局定位 & 已知初始位姿 / 绑架问题、重定位 \\
    时间特性 & 单帧定位 / 递推滤波 & 单帧配准 / EKF、粒子滤波 \\
    先验地图 & 有地图 / 无地图 & 地图匹配 / SLAM、UWB 辅助 \\
    \bottomrule
  \end{tabular}
\end{table}

这四组概念经常组合出现。比如“哨兵开机后不知道自己在场地哪里，靠着先验地图找回位置”，就是“绝对定位 + 全局定位 + 有地图”；而“比赛过程中用 FAST-LIO 连续估计位姿”，是“相对定位 + 局部定位 + 无地图”。

理解这张表的意义在于：\textbf{没有“最好的定位算法”，只有“匹配场景的定位方案”}。这也是为什么本书会花大量篇幅讲每种方法的适用条件和失效场景。

\section{SLAM：先有鸡还是先有蛋}\label{sec:01-slam}

把定位和建图放在一起，就出现了著名的“鸡生蛋”问题：定位需要地图，建图需要定位，那到底先有谁？

这个问题的正式名字叫\textbf{同时定位与建图}（Simultaneous Localization and Mapping，SLAM）。它的解决思路不是“先有鸡还是先有蛋”，而是\textbf{同时估计两者}：把位姿序列和地图都当作未知量，用所有观测一起做最大后验估计。

用数学语言写出来，就是
\begin{equation}
\left(\bm{X}^{*},\bm{M}^{*}\right)=\argmax_{\bm{X},\bm{M}}\ p\left(\bm{X},\bm{M}\mid\bm{Z}\right),
\label{eq:01-map}
\end{equation}
其中 $\bm{X}$ 是位姿序列，$\bm{M}$ 是地图，$\bm{Z}$ 是所有观测。对式~\eqref{eq:01-map} 取负对数，就变成一个大规模最小二乘问题——这正是第 13 章“后端”要做的事。

直观上，SLAM 的逻辑是：
\begin{enumerate}
  \item 先用短时间内的观测估计相邻两帧之间的相对运动（这一步不需要地图）；
  \item 把相对运动串起来，得到一条“局部一致”的轨迹；
  \item 随着轨迹变长，误差累积，轨迹开始漂；
  \item 当机器人回到曾经去过的地方（回环），用这个信息把整条轨迹“拉回来”；
  \item 轨迹修正后，地图也跟着修正。
\end{enumerate}

第 1--2 步是前端，第 4--5 步是后端。这就是下一节要讲的划分。

\section{“前端/后端”到底指什么}\label{sec:01-febe}

“前端”和“后端”是本书出现频率最高的两个词，也是最容易被误用的两个词。它们在不同语境下含义不同，先在这里建立直觉，细节分别在第 11 章（SLAM）、第 18 章（导航）展开。

\begin{keybox}{一句话版本}
\textbf{前端负责“局部、快速、增量”的估计，后端负责“全局、较慢、一致”的优化。}
\end{keybox}


三种常见语境：

\begin{enumerate}
  \item \textbf{SLAM 语境}（最常见）：前端 = 数据关联 + 帧间位姿估计 + 局部地图；后端 = 回环检测 + 全局优化 + 地图修正。这是第 11--13 章的内容。
  \item \textbf{定位系统语境}：前端 = 里程计（LIO/VIO），后端 = 先验地图配准 + 因子图 + 状态机。这是第 23 章 Z-Localization 的划分：LIO 自己维护 odom 系，后端负责 $\text{map}\to\text{odom}$ 的变换。
  \item \textbf{导航系统语境}：前端 = 定位与建图（感知），后端 = 规划与控制（决策）。这种说法在工程口语里很常见，但并不严谨——它把“导航”当成了一个更大的系统。
\end{enumerate}

第 11 章会用一张表把三种语境并列对比。你现在只需要记住：\textbf{看到“前端/后端”时，先问一句“相对于什么系统而言”}。

\section{全书地图与阅读路线}\label{sec:01-roadmap}

本书共五篇三十章，结构如图~\ref{fig:01-roadmap} 所示。

\begin{figure}[htbp]
  \centering
  \begin{tikzpicture}[node distance=0.45cm, every node/.style={font=\small}]
    \node[block2, minimum width=13cm, minimum height=9mm] (p1) {第 I 篇\quad 基础篇：坐标系、概率、优化、传感器（第 1--3 章）};
    \node[block, below=of p1, minimum width=13cm, minimum height=9mm] (p2) {第 II 篇\quad 里程计篇：轮式、IMU、激光、视觉、融合、全局定位（第 4--10 章）};
    \node[block2, below=of p2, minimum width=13cm, minimum height=9mm] (p3) {第 III 篇\quad SLAM 篇：前端、后端、建图、Cartographer、LIO（第 11--17 章）};
    \node[block, below=of p3, minimum width=13cm, minimum height=9mm] (p4) {第 IV 篇\quad 导航篇：Nav2、全局/局部规划、CMU 3D 导航、Z-Localization（第 18--24 章）};
    \node[block2, below=of p4, minimum width=13cm, minimum height=9mm] (p5) {第 V 篇\quad 工程篇：ROS 2、标定、调参、评测、实战（第 25--30 章）};
    \draw[arrow] (p1) -- (p2);
    \draw[arrow] (p2) -- (p3);
    \draw[arrow] (p3) -- (p4);
    \draw[arrow] (p4) -- (p5);
    \node[right=0.15cm of p1, font=\scriptsize, text=black!55] {必读};
    \node[right=0.15cm of p2, font=\scriptsize, text=black!55] {核心};
    \node[right=0.15cm of p3, font=\scriptsize, text=black!55] {核心};
    \node[right=0.15cm of p4, font=\scriptsize, text=black!55] {进阶};
    \node[right=0.15cm of p5, font=\scriptsize, text=black!55] {落地};
  \end{tikzpicture}
  \caption{全书结构。建议按顺序阅读第一篇，之后可按兴趣跳读。}
  \label{fig:01-roadmap}
\end{figure}

如果时间有限，有三种读法：

\begin{itemize}
  \item \textbf{速成路线}（一周）：第 1、2、4、6、11、16、17 章。理解“里程计怎么算、前端后端各干什么、FAST-LIO 怎么工作”。
  \item \textbf{建图路线}（两周）：第 1--3、7、11--15 章。重点是 Cartographer 和地图表示。
  \item \textbf{导航路线}（两周）：第 1、10、11、18--22 章。重点是 Nav2 与 CMU 3D 导航栈。
\end{itemize}

\section{给新手的三个建议}\label{sec:01-advice}

\begin{enumerate}
  \item \textbf{先在脑子里跑通数据流，再抠公式。} 每读一个算法，先问：输入是什么？输出是什么？中间经过几个模块？每个模块的频率是多少？把这张图画出来，比记公式有用得多。
  \item \textbf{公式要能“算一个数”。} 看到一个公式，尝试代入具体数字算一遍（比如 $\Delta t=10$ ms、速度 $1$ m/s，走 1 秒误差多少）。能算出数量级，才算真懂。
  \item \textbf{不要跳过标定和坐标变换。} 新手 80\% 的“算法不工作”，最后都发现是外参错了、时间戳错了、TF 树断了。第 26 章值得反复看。
\end{enumerate}

\begin{summary}

\begin{enumerate}
  \item 定位估计位姿，建图表示环境，导航规划并执行运动；三者互相依赖，构成一个闭环。
  \item 定位可按参考系、初始信息、时间特性、先验地图四个维度分类；没有万能的算法，只有匹配场景的方案。
  \item SLAM 通过“同时估计位姿与地图”解决鸡生蛋问题，本质是一个大规模最大后验估计。
  \item 前端负责局部、快速、增量的估计；后端负责全局、较慢、一致的优化。看到这两个词时先明确“相对于哪个系统”。
  \item 全书五篇三十章：基础 → 里程计 → SLAM → 导航 → 工程。
\end{enumerate}

\end{summary}

\begin{problemset}

  \item 用一句话分别定义定位、建图、导航，并各举一个 RoboMaster 场景中的例子。
  \item 为什么“相对定位”的误差会累积，而“绝对定位”的误差有界？请从参考系的角度解释。
  \item 如果机器人开机时不知道自己在场地的哪个角落，它属于哪一种定位问题？可以用哪些方法解决？
  \item 用自己的话解释“前端/后端”，并说明在 SLAM 语境和导航语境下有什么不同。
  \item 画出你心目中一台哨兵机器人的数据流图（传感器 → 算法模块 → 执行器），标出每一路数据的频率量级。
\end{problemset}

\begin{reading}

\begin{itemize}
  \item \cite{thrun2005probabilistic}：第 1 章给出了定位、建图、SLAM 的经典定义与分类，是本章最好的补充。
  \item \cite{barfoot2017state}：第 1 章从状态估计角度统一了定位与 SLAM。
  \item \cite{macenski2022ros2}：了解 ROS 2 的工程背景，为第 25 章做准备。
  \item 本书第 11 章：前端/后端的正式定义与数据流。
\end{itemize}


\end{reading}



\chapter{数学基础：从概率到李群}\label{ch:02-math}

\chapterintro{这一章不追求数学上的完备，只准备“够用”的工具箱。我们按定位与导航里真正会用到的东西来组织：坐标系怎么描述、概率怎么描述不确定性、方差和协方差到底在算什么、如何用最小二乘去拟合一个位姿、以及为什么旋转不能直接相加而要用李群。每一个概念都会从定义讲起，给出计算公式和一个可以手算的小例子。读完这一章，后面所有章节里的公式你至少能“看懂每一行在做什么”。}

\section{坐标系与刚体变换}\label{sec:02-frame}

\subsection{点、向量与坐标系}

一个三维空间中的点，只有在指定了坐标系之后才有坐标。本书约定：用 $\{W\}$ 表示世界系（map 系），$\{B\}$ 表示机体系（默认与 IMU 系重合），$\{L\}$、$\{C\}$ 分别表示雷达系和相机系。点 $P$ 在 $\{B\}$ 系下的坐标记作 ${}^{B}\bm{p}\in\R^{3}$，写成列向量
\begin{equation}
{}^{B}\bm{p}=\begin{bmatrix}p_x\\p_y\\p_z\end{bmatrix}.
\end{equation}
同一个点在不同坐标系下的坐标不同，它们之间通过“旋转 + 平移”联系起来：
\begin{equation}
{}^{W}\bm{p}={}^{W}\bm{R}_{B}\,{}^{B}\bm{p}+{}^{W}\bm{p}_{B},
\label{eq:02-rigid}
\end{equation}
其中 ${}^{W}\bm{R}_{B}\in\R^{3\times3}$ 是旋转矩阵，${}^{W}\bm{p}_{B}\in\R^{3}$ 是 $B$ 系原点在 $W$ 系中的位置。式~\eqref{eq:02-rigid} 就是刚体变换（rigid transformation）的全部内容：先转，再平移。

\subsection{旋转矩阵}

\begin{keybox}{旋转矩阵的定义}
一个 $3\times3$ 矩阵 $\bm{R}$ 是旋转矩阵，当且仅当它满足
\begin{equation}
\bm{R}^{\top}\bm{R}=\bm{I},\qquad \det(\bm{R})=1.
\label{eq:02-so3def}
\end{equation}
所有这样的矩阵构成特殊正交群 $\SO(3)$。
\end{keybox}

为什么是这两个条件？$\bm{R}^{\top}\bm{R}=\bm{I}$ 保证长度和夹角不变：对任意向量 $\bm{v}$，有 $\norm{\bm{R}\bm{v}}^{2}=\bm{v}^{\top}\bm{R}^{\top}\bm{R}\bm{v}=\bm{v}^{\top}\bm{v}=\norm{\bm{v}}^{2}$。而 $\det(\bm{R})=\pm1$ 中取 $+1$ 是为了排除“镜像”，因为镜子里的世界不是旋转出来的。

旋转矩阵的复合就是矩阵乘法：先转 ${}^{B}\bm{R}_{C}$ 再转 ${}^{W}\bm{R}_{B}$，总旋转是 ${}^{W}\bm{R}_{B}{}^{B}\bm{R}_{C}$。注意下标“相邻相消”：
\begin{equation}
{}^{W}\bm{p}={}^{W}\bm{R}_{B}\left({}^{B}\bm{R}_{C}{}^{C}\bm{p}\right)
=\left({}^{W}\bm{R}_{B}{}^{B}\bm{R}_{C}\right){}^{C}\bm{p}.
\end{equation}
这个“下标相消”规则在后面的 TF 树、标定、外参里会反复出现，值得记牢。

\subsection{齐次坐标与 $\SE(3)$}

式~\eqref{eq:02-rigid} 是“先乘再加”，不便于连续复合。引入齐次坐标可以把旋转和平移合并成一次矩阵乘法。定义
\begin{equation}
{}^{W}\bar{\bm{p}}=\begin{bmatrix}{}^{W}\bm{p}\\1\end{bmatrix}\in\R^{4},\qquad
{}^{W}\bm{T}_{B}=\begin{bmatrix}{}^{W}\bm{R}_{B}&{}^{W}\bm{p}_{B}\\\bm{0}^{\top}&1\end{bmatrix}\in\R^{4\times4},
\end{equation}
于是
\begin{equation}
{}^{W}\bar{\bm{p}}={}^{W}\bm{T}_{B}\,{}^{B}\bar{\bm{p}}.
\end{equation}
所有这样的 $\bm{T}$ 构成特殊欧氏群 $\SE(3)$。复合规则同样是矩阵乘法，且
\begin{equation}
\bm{T}^{-1}=\begin{bmatrix}\bm{R}^{\top}&-\bm{R}^{\top}\bm{t}\\\bm{0}^{\top}&1\end{bmatrix}.
\label{eq:02-se3inv}
\end{equation}
验证式~\eqref{eq:02-se3inv}：直接乘开，左上角 $\bm{R}^{\top}\bm{R}=\bm{I}$，右上角 $\bm{R}^{\top}\bm{t}-\bm{R}^{\top}\bm{t}=\bm{0}$。这个逆的闭式表达比数值求逆快得多，在代码里几乎总用这个形式。

\subsection{欧拉角与万向锁}

把旋转拆成绕 $z$、$y$、$x$ 三个轴依次转的角度 $(\psi,\theta,\phi)$（yaw、pitch、roll），叫欧拉角。它直观、人类可读（IMU 输出的姿态常用它显示），但有两个致命问题：

\begin{enumerate}
  \item \textbf{万向锁（gimbal lock）}：当 pitch 接近 $\pm90^\circ$ 时，yaw 和 roll 退化成同一个自由度，姿态表示不唯一，运动学方程出现奇异。
  \item \textbf{插值困难}：两组欧拉角之间直接线性插值，得到的中间姿态往往不是“最短路径”旋转。
\end{enumerate}

因此，内部计算一律用旋转矩阵或四元数，欧拉角只在“给人看”的时候使用。

\section{概率论最小工具箱}\label{sec:02-prob}

定位的本质是“在不确定中估计”。传感器有噪声，运动模型有误差，所以我们用概率分布描述“我认为机器人在哪”。这一节把后面会用到的概率概念逐个讲清楚，尤其是\textbf{方差}和\textbf{协方差}——它们是全书出现频率最高的两个词。

\subsection{随机变量与概率}

一个随机变量 $X$ 是“取值不确定的量”。如果 $X$ 只能取有限个离散值 $x_1,x_2,\dots$，用概率质量函数 $p(X=x_i)$ 描述，满足
\begin{equation}
p(X=x_i)\ge 0,\qquad \sum_i p(X=x_i)=1.
\end{equation}
如果 $X$ 可以在实数轴上连续取值，用概率密度函数 $p(x)$ 描述，满足
\begin{equation}
p(x)\ge 0,\qquad \int_{-\infty}^{+\infty}p(x)\diff x=1,
\end{equation}
并且落在区间 $[a,b]$ 的概率是
\begin{equation}
P(a\le X\le b)=\int_a^b p(x)\diff x.
\end{equation}
注意：$p(x)$ 本身不是概率，它可能大于 1；只有积分出来的面积才是概率。

\subsection{期望}

\begin{keybox}{定义：期望}
离散随机变量的期望是
\begin{equation}
\E[X]=\sum_i x_i\,p(X=x_i),
\end{equation}
连续随机变量的期望是
\begin{equation}
\E[X]=\int_{-\infty}^{+\infty}x\,p(x)\diff x.
\label{eq:02-expect}
\end{equation}
期望就是“加权平均”，权重是概率。
\end{keybox}

\textbf{例}：掷一枚均匀硬币，正面记 1、反面记 0，则
\begin{equation}
\E[X]=1\times 0.5+0\times 0.5=0.5.
\end{equation}

期望最重要的性质是\textbf{线性性}，它对任意随机变量（无论是否独立）都成立：
\begin{equation}
\E[aX+b]=a\E[X]+b,\qquad \E[X+Y]=\E[X]+\E[Y].
\label{eq:02-lin}
\end{equation}
证明第一个式子（离散情形）：
\begin{equation}
\E[aX+b]=\sum_i (ax_i+b)p(x_i)=a\sum_i x_ip(x_i)+b\sum_i p(x_i)=a\E[X]+b.
\end{equation}
最后一步用了 $\sum_i p(x_i)=1$。这个性质是后面所有“误差传播”公式的基础。

\subsection{方差与标准差}

\begin{keybox}{定义：方差}
方差衡量随机变量偏离其期望的程度，定义为
\begin{equation}
\Var(X)=\E\left[(X-\E[X])^{2}\right].
\label{eq:02-vardef}
\end{equation}
它的平方根 $\sigma_X=\sqrt{\Var(X)}$ 叫标准差。
\end{keybox}

把定义展开，可以得到更常用的计算公式：
\begin{equation}
\Var(X)=\E\left[X^{2}-2X\E[X]+(\E[X])^{2}\right]
=\E[X^{2}]-2\E[X]\E[X]+(\E[X])^{2}
=\E[X^{2}]-(\E[X])^{2}.
\label{eq:02-var-expand}
\end{equation}
第二步用了期望的线性性，第三步用了 $\E[\E[X]]=\E[X]$（$\E[X]$ 是一个常数）。所以记住这个等价形式：
\begin{equation}
\boxed{\Var(X)=\E[X^{2}]-\left(\E[X]\right)^{2}}
\end{equation}
即“平方的期望减去期望的平方”。

\textbf{例}：掷一枚均匀骰子，$X\in\{1,2,3,4,5,6\}$，每个概率 $1/6$。
\begin{align}
\E[X]&=\frac{1+2+3+4+5+6}{6}=3.5,\\
\E[X^{2}]&=\frac{1+4+9+16+25+36}{6}=\frac{91}{6}\approx 15.17,\\
\Var(X)&=\frac{91}{6}-3.5^{2}=15.17-12.25\approx 2.92,\\
\sigma_X&=\sqrt{2.92}\approx 1.71.
\end{align}
也就是说，掷骰子的结果平均偏离 3.5 约 1.71。方差是“偏离程度的平方的平均”，开根号后才和原数据同量纲。

方差有两个常用性质，后面推导漂移时会直接用到：
\begin{equation}
\Var(aX+b)=a^{2}\Var(X),
\label{eq:02-var-scale}
\end{equation}
\begin{equation}
\text{若 }X,Y\text{ 相互独立，则}\ \Var(X+Y)=\Var(X)+\Var(Y).
\label{eq:02-var-sum}
\end{equation}
先证式~\eqref{eq:02-var-scale}：
\begin{equation}
\Var(aX+b)=\E\left[(aX+b-\E[aX+b])^{2}\right]
=\E\left[(aX-a\E[X])^{2}\right]
=a^{2}\E\left[(X-\E[X])^{2}\right]=a^{2}\Var(X).
\end{equation}
再证式~\eqref{eq:02-var-sum}。先算
\begin{equation}
\Var(X+Y)=\E\left[(X+Y)^{2}\right]-\left(\E[X+Y]\right)^{2}
=\E[X^{2}]+2\E[XY]+\E[Y^{2}]-\left(\E[X]+\E[Y]\right)^{2}.
\end{equation}
展开右边最后一项：
\begin{equation}
\left(\E[X]+\E[Y]\right)^{2}=(\E[X])^{2}+2\E[X]\E[Y]+(\E[Y])^{2}.
\end{equation}
两式相减得
\begin{equation}
\Var(X+Y)=\left[\E[X^{2}]-(\E[X])^{2}\right]+\left[\E[Y^{2}]-(\E[Y])^{2}\right]+2\left(\E[XY]-\E[X]\E[Y]\right).
\end{equation}
前两个方括号分别是 $\Var(X)$、$\Var(Y)$；而独立意味着 $\E[XY]=\E[X]\E[Y]$，最后一项为零。于是 $\Var(X+Y)=\Var(X)+\Var(Y)$。

\subsection{协方差与协方差矩阵}

方差只描述一个变量。两个变量“一起变化”的程度用协方差描述：
\begin{equation}
\Cov(X,Y)=\E\left[(X-\E[X])(Y-\E[Y])\right]=\E[XY]-\E[X]\E[Y].
\label{eq:02-cov}
\end{equation}
展开方式与方差完全一样。$\Cov(X,X)=\Var(X)$。若 $X,Y$ 独立，则 $\Cov(X,Y)=0$；但反过来不成立——协方差为零只说明“没有线性相关”，可能存在非线性关系。

对多维随机向量 $\bm{x}=[x_1,\dots,x_n]^{\top}$，把所有两两协方差排成矩阵：
\begin{equation}
\bm{\Sigma}=\E\left[(\bm{x}-\E[\bm{x}])(\bm{x}-\E[\bm{x}])^{\top}\right]\in\R^{n\times n},
\label{eq:02-covmat}
\end{equation}
叫协方差矩阵。它的对角元是各方差 $\Sigma_{ii}=\Var(x_i)$，非对角元是协方差 $\Sigma_{ij}=\Cov(x_i,x_j)$。协方差矩阵一定是对称半正定的，因为对任意向量 $\bm{a}$，
\begin{equation}
\bm{a}^{\top}\bm{\Sigma}\bm{a}=\E\left[\left(\bm{a}^{\top}(\bm{x}-\E[\bm{x}])\right)^{2}\right]\ge 0.
\end{equation}

\textbf{协方差的几何意义}：二维情况下，协方差矩阵的特征向量给出分布的“主轴方向”，特征值给出沿主轴的方差大小。特征值越大，说明该方向上不确定性越大，分布椭圆越长。第 9 章用“马氏距离”做外点剔除、第 13 章用“信息矩阵”加权残差，本质都是在利用这个椭圆。

\subsection{高斯分布}

一维高斯（正态）分布的概率密度为
\begin{equation}
p(x)=\frac{1}{\sqrt{2\pi}\,\sigma}\exp\left(-\frac{(x-\mu)^{2}}{2\sigma^{2}}\right),
\label{eq:02-gauss1d}
\end{equation}
记作 $x\sim\mathcal{N}(\mu,\sigma^{2})$。它由两个参数完全确定：均值 $\mu$ 和方差 $\sigma^{2}$。

\textbf{验证它确实是概率密度}：首先换元 $u=(x-\mu)/(\sqrt{2}\sigma)$，则 $\diff x=\sqrt{2}\sigma\diff u$，
\begin{equation}
\int_{-\infty}^{+\infty}p(x)\diff x=\frac{1}{\sqrt{2\pi}\sigma}\int_{-\infty}^{+\infty}e^{-u^{2}}\sqrt{2}\sigma\diff u
=\frac{1}{\sqrt{\pi}}\int_{-\infty}^{+\infty}e^{-u^{2}}\diff u.
\end{equation}
利用高斯积分 $\int_{-\infty}^{+\infty}e^{-u^{2}}\diff u=\sqrt{\pi}$，得到积分等于 1。其次算期望：
\begin{equation}
\E[X]=\int x\,p(x)\diff x=\mu,
\end{equation}
因为密度关于 $x=\mu$ 对称。再算方差，令 $u=(x-\mu)/\sigma$：
\begin{equation}
\Var(X)=\frac{1}{\sqrt{2\pi}\sigma}\int_{-\infty}^{+\infty}(x-\mu)^{2}e^{-\frac{(x-\mu)^{2}}{2\sigma^{2}}}\diff x
=\frac{\sigma^{2}}{\sqrt{2\pi}}\int_{-\infty}^{+\infty}u^{2}e^{-u^{2}/2}\diff u=\sigma^{2},
\end{equation}
其中用到分部积分得到的 $\int u^{2}e^{-u^{2}/2}\diff u=\sqrt{2\pi}$。

多维高斯分布写成
\begin{equation}
p(\bm{x})=\frac{1}{(2\pi)^{n/2}\abs{\bm{\Sigma}}^{1/2}}
\exp\left(-\frac{1}{2}(\bm{x}-\bm{\mu})^{\top}\bm{\Sigma}^{-1}(\bm{x}-\bm{\mu})\right),
\label{eq:02-gaussnd}
\end{equation}
记作 $\bm{x}\sim\mathcal{N}(\bm{\mu},\bm{\Sigma})$。指数里的二次型
\begin{equation}
d^{2}=(\bm{x}-\bm{\mu})^{\top}\bm{\Sigma}^{-1}(\bm{x}-\bm{\mu})
\label{eq:02-mahal}
\end{equation}
叫马氏距离的平方。它和欧氏距离的区别是“按不确定性的方向加权”：方差小的方向权重更大。第 9 章会证明，$d^{2}$ 服从卡方分布，这正是外点剔除检验的依据。

\subsection{条件概率与贝叶斯公式}

条件概率的定义是
\begin{equation}
p(A\mid B)=\frac{p(A,B)}{p(B)},\qquad p(B)>0,
\label{eq:02-cond}
\end{equation}
即“在 $B$ 发生的条件下 $A$ 发生的概率”等于联合概率除以 $B$ 的概率。由对称性 $p(A,B)=p(B\mid A)p(A)$，代入式~\eqref{eq:02-cond} 并交换 $A$、$B$ 的位置，得到贝叶斯公式：
\begin{equation}
p(A\mid B)=\frac{p(B\mid A)\,p(A)}{p(B)}.
\label{eq:02-bayes}
\end{equation}
把它放到定位里：$A$ 是“机器人的位姿”，$B$ 是“传感器观测”。于是
\begin{equation}
\underbrace{p(\text{位姿}\mid\text{观测})}_{\text{后验}}
=\frac{\underbrace{p(\text{观测}\mid\text{位姿})}_{\text{似然}}\underbrace{p(\text{位姿})}_{\text{先验}}}
{\underbrace{p(\text{观测})}_{\text{证据}}}.
\end{equation}
后验就是“看到这帧点云后，我认为自己在哪”。分母与位姿无关，在优化里通常被忽略（当作常数）。整本书的定位与 SLAM 算法，本质上都在算这个后验：滤波器用递推近似它，优化器用最小二乘找它的最大值。

\subsection{独立性}

两个随机变量独立，当且仅当
\begin{equation}
p(X,Y)=p(X)p(Y)\quad\Longleftrightarrow\quad p(X\mid Y)=p(X).
\end{equation}
独立是后面大量公式成立的前提：多个观测独立时，似然可以连乘，取对数后变成求和，最小二乘的各项才能分开写。工程上要特别小心——时间上相邻的传感器读数、同一平面上相邻的点，往往并不独立，第 23 章 Z-Localization 讨论的“平面内点残差相关”就是一例。

\subsection{卡方分布}

\begin{definition}[卡方分布]
设 $X_1,X_2,\dots,X_k$ 是相互独立的标准正态随机变量，$X_i\sim\mathcal{N}(0,1)$，则
\begin{equation}
Q=X_1^{2}+X_2^{2}+\cdots+X_k^{2}
\end{equation}
服从自由度为 $k$ 的卡方分布，记作 $Q\sim\chi^{2}(k)$。
\end{definition}

自由度 $k=1$ 时密度函数为
\begin{equation}
p(q)=\frac{1}{\sqrt{2\pi q}}e^{-q/2},\qquad q\ge 0,
\end{equation}
$k=2$ 时是 $p(q)=\frac12 e^{-q/2}$，$k$ 一般时
\begin{equation}
p(q)=\frac{1}{2^{k/2}\Gamma(k/2)}q^{k/2-1}e^{-q/2},
\end{equation}
其中 $\Gamma$ 是伽马函数（$\Gamma(n)=(n-1)!$）。我们不要求记住这个密度，只需要知道\textbf{给定置信水平 $\alpha$ 时，分位数 $\chi^{2}_{k,\alpha}$ 可以从表里查到}。

为什么工程里关心卡方分布？因为如果残差 $\bm{r}$ 满足零均值高斯分布、协方差为 $\bm{S}$，那么二次型
\begin{equation}
d^{2}=\bm{r}^{\top}\bm{S}^{-1}\bm{r}
\end{equation}
就是若干个标准正态变量的平方和，服从 $\chi^{2}(k)$（$k$ 是残差维数）。于是可以设定一个阈值 $\chi^{2}_{k,0.95}$：如果 $d^{2}$ 超过它，说明这个观测“太离谱”，应当作为外点剔除。这就是第 9 章马氏距离检验、第 23 章 NIS 检验的数学基础。

\section{微积分与优化}\label{sec:02-opt}

\subsection{泰勒展开}

一元函数在 $x_0$ 处的泰勒展开是
\begin{equation}
f(x_0+\delta)=f(x_0)+f'(x_0)\delta+\frac{1}{2}f''(x_0)\delta^{2}+\frac{1}{6}f'''(x_0)\delta^{3}+\cdots
\label{eq:02-taylor}
\end{equation}
它的含义是：用多项式在局部逼近函数。当我们只关心很小的 $\delta$ 时，可以只保留前两项（一阶近似）或前三项（二阶近似），把高阶项写成 $O(\delta^{2})$、$O(\delta^{3})$。

多元函数的一阶近似为
\begin{equation}
f(\bm{x}_0+\delta\bm{x})\approx f(\bm{x}_0)+\nabla f(\bm{x}_0)^{\top}\delta\bm{x},
\end{equation}
其中梯度 $\nabla f=\left[\partial f/\partial x_1,\dots,\partial f/\partial x_n\right]^{\top}$。二阶近似为
\begin{equation}
f(\bm{x}_0+\delta\bm{x})\approx f(\bm{x}_0)+\nabla f^{\top}\delta\bm{x}+\frac12\delta\bm{x}^{\top}\bm{H}\delta\bm{x},
\end{equation}
其中 $\bm{H}$ 是海森矩阵（见下节）。\textbf{整本书的非线性优化，本质上就是反复做这个一阶/二阶近似，然后解一个线性问题。}

\subsection{雅可比矩阵}

\begin{definition}[雅可比矩阵]
设 $\bm{f}:\R^{n}\to\R^{m}$ 是一个向量函数，$\bm{f}(\bm{x})=[f_1(\bm{x}),\dots,f_m(\bm{x})]^{\top}$。它的雅可比矩阵 $\bm{J}\in\R^{m\times n}$ 定义为
\begin{equation}
\bm{J}=\frac{\partial\bm{f}}{\partial\bm{x}}=
\begin{bmatrix}
\dfrac{\partial f_1}{\partial x_1}&\cdots&\dfrac{\partial f_1}{\partial x_n}\\
\vdots&\ddots&\vdots\\
\dfrac{\partial f_m}{\partial x_1}&\cdots&\dfrac{\partial f_m}{\partial x_n}
\end{bmatrix}.
\end{equation}
\end{definition}
雅可比就是“多元函数的导数”，它告诉你：输入变化一点点，输出会怎么变。$\bm{f}(\bm{x}+\delta\bm{x})\approx\bm{f}(\bm{x})+\bm{J}\delta\bm{x}$。

\textbf{例}：$\bm{f}(x,y)=\begin{bmatrix}x^{2}y\\ \sin x\end{bmatrix}$，则
\begin{equation}
\bm{J}=\begin{bmatrix}2xy&x^{2}\\ \cos x&0\end{bmatrix}.
\end{equation}

\subsection{海森矩阵与极值}

\begin{definition}[海森矩阵]
标量函数 $f:\R^{n}\to\R$ 的二阶偏导矩阵
\begin{equation}
\bm{H}=\frac{\partial^{2}f}{\partial\bm{x}^{2}},\qquad H_{ij}=\frac{\partial^{2}f}{\partial x_i\partial x_j}
\end{equation}
叫海森矩阵。
\end{definition}
若 $\nabla f(\bm{x}^{*})=\bm{0}$ 且 $\bm{H}(\bm{x}^{*})$ 正定，则 $\bm{x}^{*}$ 是局部极小点；若负定则是极大点；若不定则是鞍点。这个判别法与一元函数“二阶导数为正”完全对应。

\subsection{最小二乘与正规方程}

\begin{definition}[最小二乘]
给定超定方程组 $\bm{A}\bm{x}\approx\bm{b}$（$\bm{A}\in\R^{m\times n}$，$m>n$），最小二乘解是使残差平方和最小的 $\bm{x}$：
\begin{equation}
\bm{x}^{*}=\argmin_{\bm{x}}\ \norm{\bm{A}\bm{x}-\bm{b}}^{2}.
\end{equation}
\end{definition}

推导：令 $F(\bm{x})=\norm{\bm{A}\bm{x}-\bm{b}}^{2}=(\bm{A}\bm{x}-\bm{b})^{\top}(\bm{A}\bm{x}-\bm{b})$。展开：
\begin{equation}
F(\bm{x})=\bm{x}^{\top}\bm{A}^{\top}\bm{A}\bm{x}-2\bm{b}^{\top}\bm{A}\bm{x}+\bm{b}^{\top}\bm{b}.
\end{equation}
对 $\bm{x}$ 求梯度并令其为零：
\begin{equation}
\nabla F=2\bm{A}^{\top}\bm{A}\bm{x}-2\bm{A}^{\top}\bm{b}=\bm{0},
\end{equation}
得到\textbf{正规方程（normal equation）}
\begin{equation}
\boxed{\bm{A}^{\top}\bm{A}\,\bm{x}=\bm{A}^{\top}\bm{b}}
\label{eq:02-normal}
\end{equation}
若 $\bm{A}^{\top}\bm{A}$ 可逆，则 $\bm{x}=(\bm{A}^{\top}\bm{A})^{-1}\bm{A}^{\top}\bm{b}$。

\textbf{几何解释}：$\bm{A}\bm{x}^{*}$ 是 $\bm{b}$ 在 $\bm{A}$ 的列空间上的正交投影，残差 $\bm{A}\bm{x}^{*}-\bm{b}$ 垂直于列空间。这就是“最小二乘”名字的来源。

\subsection{高斯--牛顿法}

实际问题里残差通常是非线性的：$\bm{r}(\bm{x})$，目标是最小化 $\frac12\norm{\bm{r}(\bm{x})}^{2}$。对残差做一阶泰勒展开：
\begin{equation}
\bm{r}(\bm{x}+\delta\bm{x})\approx\bm{r}(\bm{x})+\bm{J}\delta\bm{x},
\end{equation}
代入目标函数：
\begin{equation}
\frac12\norm{\bm{r}+\bm{J}\delta\bm{x}}^{2}
=\frac12\left(\bm{r}^{\top}\bm{r}+2\bm{r}^{\top}\bm{J}\delta\bm{x}+\delta\bm{x}^{\top}\bm{J}^{\top}\bm{J}\delta\bm{x}\right).
\end{equation}
对 $\delta\bm{x}$ 求导并令其为零：
\begin{equation}
\bm{J}^{\top}\bm{J}\,\delta\bm{x}=-\bm{J}^{\top}\bm{r},
\label{eq:02-gn}
\end{equation}
这就是\textbf{高斯--牛顿方程}。解出 $\delta\bm{x}$ 后更新 $\bm{x}\leftarrow\bm{x}+\delta\bm{x}$，重复直到收敛。它是式~\eqref{eq:02-normal} 的非线性版本，也是第 7 章 ICP、第 13 章因子图、第 15 章 Cartographer 扫描匹配共同的核心。

高斯--牛顿在残差较大或 $\bm{J}^{\top}\bm{J}$ 接近奇异时会不稳定。Levenberg--Marquardt（LM）给它加了一个阻尼项：
\begin{equation}
\left(\bm{J}^{\top}\bm{J}+\lambda\bm{I}\right)\delta\bm{x}=-\bm{J}^{\top}\bm{r},
\label{eq:02-lm}
\end{equation}
$\lambda$ 大时退化为梯度下降（稳定但慢），$\lambda$ 小时退化为高斯--牛顿（快但可能发散）。实际实现中 $\lambda$ 会根据每步误差是变大还是变小自适应调整。

\section{线性代数补充}\label{sec:02-linalg}

\subsection{特征值与特征向量}

\begin{definition}[特征值与特征向量]
对方阵 $\bm{A}\in\R^{n\times n}$，若存在非零向量 $\bm{v}$ 和标量 $\lambda$ 满足
\begin{equation}
\bm{A}\bm{v}=\lambda\bm{v},
\end{equation}
则称 $\lambda$ 为特征值，$\bm{v}$ 为对应的特征向量。
\end{definition}
它的几何意义是：矩阵 $\bm{A}$ 作用在 $\bm{v}$ 上不改变方向，只把长度缩放 $\lambda$ 倍。

对称矩阵（如协方差矩阵、$\bm{J}^{\top}\bm{J}$）有非常好的性质：
\begin{enumerate}
  \item 所有特征值都是实数；
  \item 不同特征值对应的特征向量相互正交；
  \item 可以正交对角化：$\bm{A}=\bm{Q}\bm{\Lambda}\bm{Q}^{\top}$，其中 $\bm{Q}$ 的列是单位正交特征向量，$\bm{\Lambda}=\diag(\lambda_1,\dots,\lambda_n)$。
\end{enumerate}

在定位里，特征值常被用来判断“可观测性”。例如扫描匹配的信息矩阵 $\bm{\Lambda}=\bm{J}^{\top}\bm{J}$，若某个特征值接近零，说明对应方向上“怎么动残差都不变”，即该方向退化（长走廊、大平面就是这样）。第 7 章和第 23 章都会用到这个判据。

\subsection{二次型与正定}

形如 $q(\bm{x})=\bm{x}^{\top}\bm{A}\bm{x}$ 的表达式叫二次型。若对任意 $\bm{x}\neq\bm{0}$ 都有 $q(\bm{x})>0$，称 $\bm{A}$ 正定（$\bm{A}\succ0$），等价于所有特征值大于零。正定二次型的等高线是椭球，椭球的半轴长度是 $1/\sqrt{\lambda_i}$——这就是协方差矩阵“椭圆”的严格解释。

\subsection{奇异值分解}

任意矩阵 $\bm{A}\in\R^{m\times n}$ 可以分解为
\begin{equation}
\bm{A}=\bm{U}\bm{\Sigma}\bm{V}^{\top},
\end{equation}
其中 $\bm{U}\in\R^{m\times m}$、$\bm{V}\in\R^{n\times n}$ 是正交矩阵，$\bm{\Sigma}$ 是对角矩阵（对角元为非负奇异值）。SVD 在点云配准里有关键作用：第 7 章 ICP 的闭式旋转解就是用 $\bm{W}=\bm{U}\bm{\Sigma}\bm{V}^{\top}$ 的 SVD 得到的。

\section{李群与李代数}\label{sec:02-lie}

旋转矩阵有 9 个元素，但只有 3 个自由度（正交约束 6 个）。这种“参数多于自由度”的表示没法直接做“加法更新”——$\bm{R}+\delta\bm{R}$ 一般不再是旋转矩阵。李群/李代数就是解决这个问题的工具。

\subsection{$\SO(3)$ 与 $\so(3)$}

旋转矩阵构成群 $\SO(3)$，它既是群又是光滑流形，所以叫李群。它在单位元处的切空间叫李代数 $\so(3)$，由所有 $3\times3$ 反对称矩阵组成：
\begin{equation}
\so(3)=\left\{\bm{\phi}^{\wedge}\ \middle|\ \bm{\phi}\in\R^{3}\right\},\qquad
\bm{\phi}^{\wedge}=\begin{bmatrix}0&-\phi_3&\phi_2\\ \phi_3&0&-\phi_1\\ -\phi_2&\phi_1&0\end{bmatrix}.
\label{eq:02-hat}
\end{equation}
符号 $(\cdot)^{\wedge}$ 表示“向量变反对称矩阵”，逆运算 $(\cdot)^{\vee}$ 表示“反对称矩阵变回向量”。

指数映射把李代数元素变成旋转矩阵：$\bm{R}=\exp(\bm{\phi}^{\wedge})$。推导它的闭式表达（Rodrigues 公式）。把指数按泰勒级数展开：
\begin{equation}
\exp(\bm{\phi}^{\wedge})=\bm{I}+\bm{\phi}^{\wedge}+\frac{1}{2!}(\bm{\phi}^{\wedge})^{2}+\frac{1}{3!}(\bm{\phi}^{\wedge})^{3}+\cdots
\end{equation}
注意到 $\bm{\phi}^{\wedge}$ 有一个重要性质：若记 $\theta=\norm{\bm{\phi}}$、$\bm{a}=\bm{\phi}/\theta$，则
\begin{equation}
(\bm{a}^{\wedge})^{2}=\bm{a}\bm{a}^{\top}-\bm{I},\qquad (\bm{a}^{\wedge})^{3}=-\bm{a}^{\wedge}.
\end{equation}
第一个式子可以直接验证（对任意单位向量 $\bm{a}$）；第二个由第一个两边乘 $\bm{a}^{\wedge}$ 得到。于是奇数项和偶数项可以分别合并：
\begin{align}
\exp(\theta\bm{a}^{\wedge})
&=\bm{I}+\left(\theta-\frac{\theta^{3}}{3!}+\frac{\theta^{5}}{5!}-\cdots\right)\bm{a}^{\wedge}
+\left(\frac{\theta^{2}}{2!}-\frac{\theta^{4}}{4!}+\cdots\right)(\bm{a}^{\wedge})^{2}\nonumber\\
&=\bm{I}+\sin\theta\,\bm{a}^{\wedge}+(1-\cos\theta)(\bm{a}^{\wedge})^{2}.
\label{eq:02-rodrigues}
\end{align}
这就是\textbf{Rodrigues 公式}。它说明：绕单位轴 $\bm{a}$ 旋转 $\theta$ 角的旋转矩阵只由 $\bm{\phi}=\theta\bm{a}$ 决定。反过来，给定 $\bm{R}$，也能求出 $\bm{\phi}$（对数映射）：
\begin{equation}
\theta=\arccos\left(\frac{\tr(\bm{R})-1}{2}\right),\qquad
\bm{\phi}=\frac{\theta}{2\sin\theta}\left(\bm{R}-\bm{R}^{\top}\right)^{\vee},
\end{equation}
当 $\theta\to0$ 或 $\theta\to\pi$ 时需要特殊处理（代码里通常都有分支）。

\subsection{$\SE(3)$ 与 $\se(3)$}

把平移也纳入进来，$\SE(3)$ 的李代数 $\se(3)$ 是 6 维的：
\begin{equation}
\se(3)=\left\{\bm{\xi}^{\wedge}\ \middle|\ \bm{\xi}=\begin{bmatrix}\bm{\rho}\\ \bm{\phi}\end{bmatrix}\in\R^{6}\right\},\qquad
\bm{\xi}^{\wedge}=\begin{bmatrix}\bm{\phi}^{\wedge}&\bm{\rho}\\\bm{0}^{\top}&0\end{bmatrix}.
\end{equation}
指数映射为
\begin{equation}
\exp(\bm{\xi}^{\wedge})=\begin{bmatrix}\exp(\bm{\phi}^{\wedge})&\bm{J}\bm{\rho}\\\bm{0}^{\top}&1\end{bmatrix},\qquad
\bm{J}=\bm{I}+\frac{1-\cos\theta}{\theta^{2}}\bm{\phi}^{\wedge}+\frac{\theta-\sin\theta}{\theta^{3}}(\bm{\phi}^{\wedge})^{2},
\end{equation}
其中 $\bm{J}$ 叫左雅可比。这里不要求会推导 $\bm{J}$，只要求知道“$\SE(3)$ 上的增量也是 6 维向量”。

\subsection{扰动模型}

在优化中，我们关心的是“给位姿加一个很小的增量”。有两种常见定义：
\begin{align}
\text{左扰动：}&\quad \bm{R}\leftarrow\exp(\delta\bm{\phi}^{\wedge})\bm{R},\\
\text{右扰动：}&\quad \bm{R}\leftarrow\bm{R}\exp(\delta\bm{\phi}^{\wedge}).
\end{align}
两种都对，区别在于雅可比的形式不同。本书约定：位姿用右扰动，即
\begin{equation}
\bm{R}\leftarrow\bm{R}\exp(\delta\bm{\phi}^{\wedge}),\qquad \bm{t}\leftarrow\bm{t}+\bm{R}\,\delta\bm{\rho},
\end{equation}
这样雅可比只和当前位姿有关，推导更整齐。注意 $\delta\bm{\rho}$ 是在机体系下定义的平移增量。

两个旋转复合的对数并不等于两个对数相加，而是有 BCH（Baker--Campbell--Hausdorff）公式：
\begin{equation}
\log\left(\exp(\bm{\phi}_1^{\wedge})\exp(\bm{\phi}_2^{\wedge})\right)
\approx
\begin{cases}
\bm{J}_l(\bm{\phi}_2)^{-1}\bm{\phi}_1+\bm{\phi}_2, & \text{左乘},\\
\bm{J}_r(\bm{\phi}_1)^{-1}\bm{\phi}_2+\bm{\phi}_1, & \text{右乘},
\end{cases}
\end{equation}
当 $\bm{\phi}_2$ 很小时取一阶近似即可。这个公式出现在第 13 章位姿图残差的雅可比推导里。

\subsection{为什么必须在流形上优化}

如果直接对旋转矩阵的 9 个元素做梯度下降，更新后的矩阵不再满足 $\bm{R}^{\top}\bm{R}=\bm{I}$，需要每步重新正交化，既慢又不自然。正确做法是：\textbf{在切空间（李代数）里求增量，再通过指数映射回到流形}：
\begin{equation}
\bm{R}_{k+1}=\bm{R}_k\exp(\delta\bm{\phi}^{\wedge}),
\end{equation}
其中 $\delta\bm{\phi}\in\R^{3}$ 是普通的、可以随便加减的向量。这样优化变量永远是 3 维的，约束自动满足。所有现代 SLAM 后端（g2o、GTSAM、Ceres 的 \lstinline{LocalParameterization}）都是这么做的。

\begin{summary}

\begin{enumerate}
  \item 刚体变换是“旋转 + 平移”，写成 $\SE(3)$ 齐次矩阵后复合就是乘法，逆有闭式表达。
  \item 期望是概率加权平均；方差是偏离期望的平方的期望，有 $\Var(X)=\E[X^{2}]-(\E[X])^{2}$、$\Var(aX+b)=a^{2}\Var(X)$、独立时方差可加三个常用性质。
  \item 协方差矩阵的对角元是方差，非对角元是协方差；它的特征向量给出不确定性的主轴。
  \item 高斯分布由均值与协方差完全确定；马氏距离的平方服从卡方分布，可用于外点剔除。
  \item 贝叶斯公式把“先验 + 似然”变成“后验”，这是所有定位算法的出发点。
  \item 非线性最小二乘 = 泰勒展开 + 正规方程；高斯--牛顿迭代求解，LM 加阻尼保证稳定。
  \item 旋转不能直接相加，要在李代数里求增量再用指数映射更新；Rodrigues 公式给出指数映射的闭式解。
\end{enumerate}

\end{summary}

\begin{problemset}

  \item 掷一枚均匀骰子，求 $X$ 的期望、方差和标准差。若把每个结果乘以 2 再加 1，新的期望和方差是多少？
  \item 证明 $\Cov(X,Y)=\E[XY]-\E[X]\E[Y]$，并说明为什么“协方差为零”不能推出“独立”。
  \item 一维高斯分布中，若 $\mu=2$、$\sigma=3$，求 $P(2\le X\le 5)$（提示：用标准正态分布表，$P(-1\le Z\le1)\approx0.6827$）。
  \item 对 $\bm{f}(\bm{x})=\bm{R}\bm{x}+\bm{t}$，求雅可比 $\partial\bm{f}/\partial\bm{x}$；若 $\bm{R}$ 用增量 $\exp(\delta\bm{\phi}^{\wedge})$ 扰动，求 $\partial\bm{f}/\partial\delta\bm{\phi}$。
  \item 用 Rodrigues 公式计算绕 $z$ 轴转 $90^\circ$ 的旋转矩阵，并与直觉核对。
  \item 设 $\bm{A}=\begin{bmatrix}2&1\\1&2\end{bmatrix}$，求特征值、特征向量，并说明它对应的二次型椭球的半轴方向。
\end{problemset}

\begin{reading}

\begin{itemize}
  \item \cite{barfoot2017state}：第 2--4 章系统讲解概率、估计与李群，适合作为本章的补充教材。
  \item \cite{sola2017eskf}：用 50 多页把四元数、SO(3)、ESKF 讲透，是 IMU 状态估计必读。
  \item \cite{thrun2005probabilistic}：第 2 章讲贝叶斯滤波，第 6 章讲栅格地图，经典且清晰。
  \item 3Blue1Brown 的线性代数系列视频：直观理解特征值、SVD 和矩阵变换。
\end{itemize}


\end{reading}



\chapter{传感器与观测模型}\label{ch:03-sensors}

\chapterintro{所有定位算法都建立在传感器之上。这一章把机器人常用的五类传感器——IMU、轮式编码器、激光雷达、相机、UWB/GNSS——逐个拆开：它们测什么、输出什么、噪声长什么样、在什么情况下会骗你。我们会为每类传感器建立“测量模型”，也就是“测量值 = 真值 + 零偏 + 噪声”这样的数学表达。后面所有估计器（卡尔曼滤波、因子图、ICP）都是在这个模型上做概率推断。}

\section{观测与测量模型}\label{sec:03-model}

\begin{definition}[测量模型]
测量模型（measurement model）描述传感器读数与真实物理量之间的统计关系，通常写成
\begin{equation}
\bm{z}=\bm{h}(\bm{x})+\bm{n},
\end{equation}
其中 $\bm{z}$ 是测量值，$\bm{x}$ 是待估计的真实状态，$\bm{h}(\cdot)$ 是已知的观测函数，$\bm{n}$ 是\textbf{观测噪声}。
\end{definition}

观测噪声通常假设为零均值高斯：$\bm{n}\sim\mathcal{N}(\bm{0},\bm{R})$，$\bm{R}$ 叫噪声协方差（在卡尔曼滤波里也记作测量噪声矩阵）。这个假设不是“真理”，而是工程上的折中：高斯分布只需要均值和方差两个数，且高斯经过线性变换后仍是高斯，这让滤波器的推导变得极其简洁。第 2 章已经证明，若 $\bm{n}\sim\mathcal{N}(\bm{0},\bm{R})$，则负对数似然正比于马氏范数 $\norm{\bm{z}-\bm{h}(\bm{x})}_{\bm{R}^{-1}}^{2}$——这正是最小二乘的来源。

\begin{keybox}{为什么噪声协方差这么重要}
残差同样大，协方差小的观测会被“更相信”，协方差大的会被“更怀疑”。调参时如果发现某个传感器“压不住”或“带偏”了融合结果，第一反应应该是检查它的协方差是否设得合理。
\end{keybox}

\section{IMU：惯性测量单元}\label{sec:03-imu}

\subsection{它测什么}

IMU 通常由两部分组成：
\begin{itemize}
  \item \textbf{陀螺仪}（gyroscope）：测量角速度 $\bm{\omega}=[\omega_x,\omega_y,\omega_z]^{\top}$，单位 rad/s 或 $^\circ$/s。
  \item \textbf{加速度计}（accelerometer）：测量比力（specific force）$\bm{f}$，单位 m/s$^{2}$。注意它测的不是“加速度”，而是“加速度减去重力”，静止时读数是 $\approx 9.81$ m/s$^{2}$ 指向天空。
\end{itemize}

现代 IMU 多为 MEMS 工艺，体积小、成本低，但噪声大、零偏随温度漂移。工业级/战术级 IMU 精度高一个到几个数量级，价格也高一个到几个数量级。

\subsection{测量模型}

IMU 的测量模型是全书第一个正式的传感器模型，务必理解每一项：
\begin{align}
\tilde{\bm{\omega}}&=\bm{\omega}+\bm{b}_g+\bm{n}_g,\label{eq:03-gyro}\\
\tilde{\bm{a}}&=\bm{R}_{WB}^{\top}\left(\bm{a}_W-\bm{g}\right)+\bm{b}_a+\bm{n}_a,\label{eq:03-accel}
\end{align}
其中：
\begin{itemize}
  \item $\tilde{\bm{\omega}},\tilde{\bm{a}}$ 是\textbf{测量值}（带弯弯的波浪号在全书统一表示原始测量）；
  \item $\bm{\omega},\bm{a}_W$ 是真实的角速度和世界系加速度；
  \item $\bm{b}_g,\bm{b}_a$ 是\textbf{零偏}（bias），即传感器在零输入时的非零输出，缓慢变化；
  \item $\bm{n}_g,\bm{n}_a$ 是零均值高斯白噪声；
  \item $\bm{g}$ 是重力向量，本书约定 $\bm{g}=[0,0,-9.81]^{\top}$（世界系 $z$ 轴向上）；
  \item $\bm{R}_{WB}^{\top}$ 把世界系向量转到机体系。
\end{itemize}

\begin{keybox}{加速度计的“坑”}
式~\eqref{eq:03-accel} 最容易被新手写错的地方是重力符号和参考系。静止时 $\bm{a}_W=\bm{0}$，测量值应为 $\tilde{\bm{a}}=-\bm{R}_{WB}^{\top}\bm{g}$。如果模拟数据里 IMU 读数模长不是 $9.81$ 左右，或者静止时方向不对，先别怀疑算法，先检查模型实现。
\end{keybox}


零偏不是常数，而是缓慢漂移的随机过程。常用模型是\textbf{随机游走}：
\begin{equation}
\dot{\bm{b}}_g=\bm{n}_{bg},\qquad \dot{\bm{b}}_a=\bm{n}_{ba},
\label{eq:03-bias-rw}
\end{equation}
其中 $\bm{n}_{bg},\bm{n}_{ba}$ 是零均值高斯白噪声。随机游走意味着零偏的方差随时间的增长而线性增长——这也是 IMU 不能长时间单独定姿的根本原因。

\subsection{噪声参数与离散化}

厂商给出的噪声密度通常是连续时间参数（如 $\sigma_g=0.001$ rad/s/$\sqrt{\text{Hz}}$）。把它离散到采样周期 $\Delta t$ 上，噪声标准差变为
\begin{equation}
\sigma_{g,d}=\frac{\sigma_g}{\sqrt{\Delta t}},\qquad \sigma_{bg,d}=\sigma_{bg}\sqrt{\Delta t}.
\label{eq:03-discrete}
\end{equation}
为什么采样越快、单步噪声越大？因为连续白噪声的方差是无穷的，我们实际观测的是它在 $\Delta t$ 上的积分；采样间隔越短，积分区间越小，但单位和归一化方式使得离散标准差与 $\sqrt{\Delta t}$ 成反比。第 6 章会从功率谱密度出发给出这个关系的完整推导。

\subsection{关键指标}

选购或评估 IMU 时关注这些量：
\begin{itemize}
  \item \textbf{零偏稳定性}（bias instability）：零偏在长时间内的波动，单位 $^\circ$/h 或 $\mu$g。越小越好。
  \item \textbf{角度随机游走}（ARW）：陀螺噪声，单位 $^\circ/\sqrt{\text{h}}$。
  \item \textbf{速度随机游走}（VRW）：加速度计噪声。
  \item \textbf{量程}（range）：超过量程会\textbf{饱和}，读数被截断——这是 FAST-LIO 类算法在剧烈运动下的典型失效原因。
  \item \textbf{带宽}与采样率：常见 100--1000 Hz。
\end{itemize}

图~\ref{fig:03-imu} 展示了 IMU 测量模型中各分量如何叠加成读数。

\begin{figure}[htbp]
  \centering
  \begin{tikzpicture}[node distance=1.2cm]
    \node[block2, minimum width=20mm] (true) {真实角速度\\$\bm{\omega}$};
    \node[block, right=0.8cm of true, minimum width=20mm] (bias) {零偏\\$\bm{b}_g(t)$};
    \node[block, right=0.8cm of bias, minimum width=20mm] (noise) {白噪声\\$\bm{n}_g$};
    \node[block2, right=0.8cm of noise, minimum width=20mm] (meas) {测量\\$\tilde{\bm{\omega}}$};
    \draw[arrow] (true) -- (bias);
    \draw[arrow] (bias) -- (noise);
    \draw[arrow] (noise) -- (meas);
    \node[below=0.5cm of bias, font=\footnotesize, text=black!60] {缓慢漂移（随机游走）};
    \node[below=0.5cm of noise, font=\footnotesize, text=black!60] {高频抖动};
  \end{tikzpicture}
  \caption{IMU 角速度的测量模型：真值叠加零偏与白噪声。}
  \label{fig:03-imu}
\end{figure}

\section{轮式编码器}\label{sec:03-wheel}

\subsection{原理与分辨率}

增量式编码器装在电机或轮轴上，轮子每转一圈输出固定数量的脉冲。设编码器线数为 $N$（每圈 $N$ 个脉冲）、减速比为 $i$、轮半径为 $r$，则轮子每转一圈，电机端产生 $N$ 个脉冲，轮端产生 $Ni$ 个脉冲。若采用四倍频计数（A/B 相正交，上升沿和下降沿都计数），单圈计数为 $4Ni$，于是每个计数对应的行进距离为
\begin{equation}
\Delta s=\frac{2\pi r}{4Ni}.
\label{eq:03-encoder}
\end{equation}
\textbf{例}：$N=500$、$i=20$、$r=0.075$ m，则
\begin{equation}
\Delta s=\frac{2\pi\times0.075}{4\times500\times20}\approx5.9\times10^{-6}\ \text{m}\approx5.9\ \mu\text{m}.
\end{equation}
分辨率远高于实际精度——真正的误差来自轮子打滑、地面不平、轮胎磨损，而不是编码器本身。

\subsection{输出}

编码器本身只给出“转了多少”，需要结合时间得到轮速 $v=\Delta s/\Delta t$，再结合底盘运动学（第 5 章）得到车体线速度和角速度。ROS 2 中通常由底盘驱动节点发布 \lstinline{nav_msgs/Odometry}，并广播 odom 到 base\_link 的变换。

\section{激光雷达}\label{sec:03-lidar}

\subsection{分类}

按扫描方式：
\begin{itemize}
  \item \textbf{机械旋转式}：电机带动激光模块旋转，360$^\circ$ 扫描，线数 16/32/64/128。代表：Velodyne、Ouster、Hesai。点云规则、每线一个 \lstinline{ring} 编号。
  \item \textbf{半固态/固态}：MEMS 微振镜、转镜、OPA 等。视场角有限（如 $120^\circ\times90^\circ$），点云非重复。代表：Livox Mid-360、Avia。
  \item \textbf{面阵/Flash}：一次曝光成图，无扫描线。
\end{itemize}

按测距原理：dToF（直接飞行时间）、iToF（间接飞行时间）、FMCW（调频连续波，可测速、抗干扰）。这个细节对算法影响不大，知道即可。

\subsection{关键参数}

\begin{table}[htbp]
  \centering
  \caption{激光雷达关键参数}\label{tab:03-lidar-params}
  \small
  \begin{tabular}{p{3.2cm}p{5.2cm}p{4.6cm}}
    \toprule
    参数 & 含义 & 对算法的影响 \\
    \midrule
    线数 / 视场角 & 垂直方向的分辨能力 & 影响点云密度与建图完整性 \\
    水平角分辨率 & 相邻两束激光的角度间隔 & 影响远处点云稀疏程度 \\
    测距范围与盲区 & 最近/最远有效距离 & 盲区内的障碍物“看不见” \\
    点频 & 每秒输出点数 & 影响单帧点数与匹配耗时 \\
    精度 & 单点测距误差（cm 级） & 影响点面残差噪声模型 \\
    回波模式 & 单回波/双回波 & 影响玻璃、雨雾场景 \\
    \bottomrule
  \end{tabular}
\end{table}

\subsection{点云与运动畸变}

\begin{definition}[点云]
点云是三维点的集合 $\mathcal{P}=\{\bm{p}_i\}_{i=1}^{N}$，每个点通常带有坐标、反射强度、时间戳，机械式雷达还带 \lstinline{ring} 编号。
\end{definition}

激光雷达不是“瞬间拍一张照片”，而是逐点扫描。一帧点云里，第一个点和最后一个点的采样时间可能相差一个扫描周期（10 Hz 雷达约 100 ms）。如果机器人在这一帧内运动了，点云就会“拖影”——同一个平面被扫成斜面。这叫\textbf{运动畸变}，必须用 IMU 或匀速模型做\textbf{去畸变}（第 16 章详述）。

\begin{keybox}{Livox 的非重复扫描}
Livox 系列采用非重复花瓣式扫描，单帧点云不形成规则的行列结构，不能直接套用“按 ring 处理”的算法。它的优势是同样点频下覆盖更均匀、近距离更密，适合 FAST-LIO 这类直接配准方法。代价是点云无序，很多传统特征提取方法需要改造。
\end{keybox}

\section{相机}\label{sec:03-camera}

相机输出二维图像，信息丰富但缺少深度。按获取深度的方式分：
\begin{itemize}
  \item \textbf{单目}：结构最简单，但尺度不可观（需要 IMU 或已知物体提供尺度）。
  \item \textbf{双目}：通过视差恢复深度，$d=\dfrac{f\,b}{Z}$（$f$ 焦距、$b$ 基线、$Z$ 深度），远处精度迅速下降。
  \item \textbf{RGB-D}：结构光或 ToF 直接给深度，室内近距离好用，阳光下易失效。
\end{itemize}

\begin{keybox}{内参与外参}
\textbf{内参}（intrinsic）描述相机自身的成像几何，如焦距 $f_x,f_y$、主点 $c_x,c_y$、畸变系数。\textbf{外参}（extrinsic）描述相机与另一个坐标系（如 IMU、车体）之间的旋转和平移。两者都要标定，混用会直接导致投影错位。第 26 章给出完整标定流程。
\end{keybox}


相机在定位里的角色通常是“提供视觉约束”：特征点重投影误差（第 8 章）、光度误差、回环场景识别。纯视觉定位的优点是信息丰富、成本低；缺点是对光照、纹理、动态物体敏感，且单目有尺度问题。

\section{其他传感器}\label{sec:03-others}

\begin{itemize}
  \item \textbf{UWB}（超宽带）：通过无线信号飞行时间测距，多个基站可解算三维位置。精度 10--30 cm，室内可用，但多径、遮挡、基站几何分布会显著影响精度。第 23 章会看到它如何作为先验地图缺失时的位置辅助。
  \item \textbf{GNSS}（卫星导航）：室外绝对定位，RTK 可达厘米级；室内、隧道、城市峡谷不可用。
  \item \textbf{二维码/ArUco}：视觉标签，精度高但需要布置，适合定点校准或小范围定位。
  \item \textbf{磁力计}：提供航向参考，易受电机、金属干扰，在机器人上通常只做辅助。
  \item \textbf{轮式触觉/IMU 之外的里程计}：如视觉里程计、雷达里程计，本质上都是“用其他传感器推算运动”，见第 4 章。
\end{itemize}

\section{时间同步与坐标系}\label{sec:03-sync}

\subsection{时间同步}

\begin{definition}[时间同步]
时间同步是让不同传感器的数据带有同一时间基准的过程，包括：硬件触发（PPS、PTP）、软件时间戳对齐、以及估计并补偿传感器之间的固定延迟。
\end{definition}

为什么重要？假设机器人以 $2$ m/s 运动，两个传感器时间戳相差 $10$ ms，位置差就是 $2\times0.01=2$ cm。对于厘米级定位，这个误差已经不可忽略；对点云配准，时间错位会让点云“错层”，残差永远降不下去。

工程上的检查清单：
\begin{enumerate}
  \item 所有传感器的时间戳是否来自同一个时钟源（\lstinline{/clock} 或系统时钟）？
  \item 回放 bag 时是否加了 \lstinline{--clock}、所有节点是否设置了 \lstinline{use_sim_time:=true}？
  \item 雷达时间戳是“帧起始”还是“帧结束”？IMU 时间戳是否已做延迟补偿？
  \item 外参标定文件中，平移和旋转的方向是否与公式一致？
\end{enumerate}

\subsection{外参}

\begin{definition}[外参]
外参是传感器坐标系相对于机器人本体坐标系（或另一个传感器坐标系）的位姿，记作 $\bm{T}_{BL}=(\bm{R}_{BL},{}^{B}\bm{p}_{L})$，表示“雷达系到机体系的变换”。
\end{definition}

外参错误的典型表现：点云与图像对不齐、地图出现双层墙、LIO 建图“越走越歪”、静止时点云漂移。所有外参都必须标定并在配置文件中正确填写（第 26 章）。

\section{传感器选型对比}\label{sec:03-compare}

\begin{table}[htbp]
  \centering
  \caption{常用传感器对比（数量级参考）}\label{tab:03-compare}
  \small
  \begin{tabular}{lcccl}
    \toprule
    传感器 & 频率 & 精度 & 成本 & 主要失效场景 \\
    \midrule
    轮式编码器 & 50--100 Hz & 1--10\% 行程 & 极低 & 打滑、悬空、地面沉降 \\
    MEMS IMU & 100--1000 Hz & 姿态短时好 & 低 & 零偏漂移、饱和、振动 \\
    2D 激光 & 10--40 Hz & cm 级 & 中 & 无高度信息、玻璃、长走廊 \\
    3D 激光 & 10--20 Hz & cm 级 & 中高 & 雨雾、强光、稀疏远处 \\
    单目相机 & 20--60 Hz & 尺度不可观 & 低 & 光照、纹理、动态物体 \\
    RGB-D & 10--30 Hz & 近距离 cm 级 & 中 & 阳光、远距离、反光 \\
    UWB & 10--100 Hz & 10--30 cm & 中 & 多径、遮挡、基站几何差 \\
    GNSS-RTK & 1--20 Hz & cm--m 级 & 高 & 室内、遮挡、多路径 \\
    \bottomrule
  \end{tabular}
\end{table}

这张表的用法不是“选最贵的”，而是“看清每种传感器的失效场景，用另一种去补它”。这也是第 9 章多传感器融合的出发点。

\begin{summary}

\begin{enumerate}
  \item 测量模型是“读数 = 真值 + 零偏 + 噪声”，高斯噪声假设让最小二乘与卡尔曼滤波得以成立。
  \item IMU 测角速度与比力，零偏缓慢漂移，量程饱和是剧烈运动下的典型问题；离散噪声与 $\sqrt{\Delta t}$ 相关。
  \item 编码器分辨率很高，但轮式里程计的真正误差来自打滑与地面。
  \item 激光雷达逐点扫描产生运动畸变，机械式与固态的点云结构差异很大，算法选型要匹配。
  \item 相机内参、外参、时间同步三者任何一个出错，都会表现为“算法不收敛”。
  \item 没有完美的传感器，只有互补的组合。
\end{enumerate}

\end{summary}

\begin{problemset}

  \item 静止时 IMU 加速度计读数应该是多少？如果模长是 $9.5$ 或 $10.2$，可能是什么原因？
  \item 一台 16 线雷达、水平角分辨率 $0.2^\circ$，在 10 m 处的相邻两点间距约是多少？
  \item 机器人以 $1.5$ m/s 运动，雷达与 IMU 时间戳相差 $20$ ms，会产生多大的位置误差？如果这个误差一直存在，点云配准会表现为什么现象？
  \item 为什么单目相机有“尺度不可观”问题？双目为什么可以恢复尺度？请写出视差与深度的关系式并解释。
  \item 列出你所用机器人上所有传感器，写出它们的时间戳来源和坐标系名称，画出 TF 树草图。
\end{problemset}

\begin{reading}

\begin{itemize}
  \item \cite{thrun2005probabilistic}：第 2--5 章详细讨论传感器模型与噪声。
  \item \cite{sola2017eskf}：IMU 噪声模型与离散化的标准参考。
  \item \cite{barfoot2017state}：第 6--7 章从估计理论角度讲传感器模型。
  \item 各传感器厂商数据手册：噪声密度、零偏稳定性、量程等参数的实际含义。
\end{itemize}


\end{reading}




\part{里程计篇：机器人如何推算自己的位置}

\chapter{里程计总论：相对定位的艺术}\label{ch:04-odom-intro}

\chapterintro{从本章开始，我们正式进入「机器人如何知道自己在哪」这个核心问题。本章先不碰任何具体传感器，只建立一张地图：什么叫位姿，什么叫位姿递推，为什么里程计天然是「相对」的，它的误差为什么注定越积越大，以及我们该用什么指标评价它。学完本章，你应当能回答三个问题：里程计在整个定位系统里处于什么位置；为什么所有里程计都会漂移、漂移的数学形状是什么；为什么评价里程计要看相对误差而不是绝对误差。后六章分别是本章四种里程计来源的展开。}

\section{为什么需要里程计}\label{sec:04-why}

设想你在一间没有窗户的房间里闭着眼走路。你记得自己先向前迈了三步、再向右转九十度、又走了两步。凭这些「自己相对于自己的运动」，你能大致画出走过的轨迹，但你并不知道房间的绝对位置，也不知道自己的朝向与地理北方的夹角。更麻烦的是，每迈一步你都会有几厘米的估计误差，走得越久，累积的位置偏差越大。

机器人的处境完全相同，只是它可以靠传感器把「迈了几步」测得更准。定位问题通常被拆成两个不同的问题：

\begin{itemize}
  \item \textbf{我在世界的哪里？}——需要地图或外部基准（GPS、已知信标、先验地图）。这类问题叫绝对定位，详见第~\ref{ch:10-global-loc}~章。
  \item \textbf{我相对出发时动了多少？}——只需要传感器对自身运动的测量，连续、高频，不需要任何外部设施。这类问题叫相对定位，也就是本章的主角：里程计。
\end{itemize}

为什么不能只靠绝对定位？三个现实原因。其一是\textbf{频率不够}：GNSS 在室外通常只有 1--10\,Hz，而差速底盘的控制回路常常要 50--100\,Hz 的位姿反馈，用 1\,Hz 的位姿闭环，控制器在两个更新之间只能靠「猜」。其二是\textbf{室内不可用}：RoboMaster 赛场是典型室内环境，GNSS 信号被建筑强烈衰减甚至完全遮挡。其三是\textbf{绝对定位会失败}：基于地图的定位在长廊、对称场地、动态障碍物遮挡时会短时失效，而控制与规划不能因为定位模块「卡了一下」就停车。

所以真实系统几乎总是「里程计 + 绝对定位修正」：里程计提供高频、平滑、局部一致的位姿，绝对定位（或 SLAM 后端、回环检测）以低频把它拉回全局正确的位置——里程计负责「连续」，绝对定位负责「不漂」。

\begin{figure}[htbp]
  \centering
  \begin{tikzpicture}[node distance=7mm and 11mm]
    \node[block] (s) {传感器\\编码器 / IMU\\激光 / 相机};
    \node[block, right=of s] (o) {里程计\\位姿递推};
    \node[block, right=of o] (tf) {TF 树\\odom$\to$base};
    \node[block2, below=of o] (b) {后端优化\\回环检测};
    \node[block2, below=of tf] (m) {map$\to$odom\\校正量};
    \draw[-stealth, thick] (s) -- node[label, above] {增量} (o);
    \draw[-stealth, thick] (o) -- node[label, above] {高频位姿} (tf);
    \draw[-stealth, thick] (o) -- (b);
    \draw[-stealth, thick] (b) -- (m);
    \draw[-stealth, thick] (m) -- (tf);
    \draw[-stealth, thick, dashed] (tf.north) to[out=90, in=90] (o.north);
  \end{tikzpicture}
  \caption{里程计在定位系统中的位置。里程计把传感器增量递推成高频位姿，直接发布 \texttt{odom}$\to$\texttt{base} 变换；后端优化低频地给出 \texttt{map}$\to$\texttt{odom} 的修正量，两者叠加才得到全局位姿。}
  \label{fig:04-system}
\end{figure}

\section{位姿、绝对定位与相对定位}\label{sec:04-pose}

先把最基础的对象说清楚。全书统一使用如下坐标系约定（见第~\ref{ch:02-math}~章）：世界系 $W$ 的 $z$ 轴向上（ENU 式），机体系 $B$ 与 IMU 系重合，重力 $\bm{g}=[0,0,-9.81]^{\top}$。

\begin{definition}[位姿]
机器人机体系 $B$ 在世界系 $W$ 中的\textbf{位姿}（pose）是一个二元组
\begin{equation}
  T_{WB} = \left(\bm{R}_{WB},\ {}^{W}\bm{p}_{B}\right),
  \label{eq:04-pose}
\end{equation}
其中 $\bm{R}_{WB}\in SO(3)$ 是旋转矩阵，${}^{W}\bm{p}_{B}\in\mathbb{R}^{3}$ 是机体原点在世界系中的位置。它作用在任意一点上：
\begin{equation}
  {}^{W}\bm{q} = \bm{R}_{WB}\,{}^{B}\bm{q} + {}^{W}\bm{p}_{B}.
  \label{eq:04-action}
\end{equation}
所有满足式~\eqref{eq:04-action} 的 $(\bm{R},\bm{p})$ 构成特殊欧氏群 $SE(3)$，即 $T_{WB}\in SE(3)$。
\end{definition}

\begin{definition}[绝对定位与相对定位]
\textbf{绝对定位}（absolute localization）指估计机体在世界系 $W$ 中的位姿 $T_{WB}$，其结果是相对某个固定原点而言的；\textbf{相对定位}（relative localization）指只估计两个时刻之间的位姿差 $T_{k}^{-1}T_{k+1}$，其结果是相对自身历史而言的。
\end{definition}

\begin{definition}[里程计与位姿递推]
\textbf{里程计}（odometry）是利用传感器测得的相邻时刻运动增量，从一个已知初始位姿出发，递推地累加出连续位姿序列的算法。把增量复合回历史位姿的过程称为\textbf{位姿递推}（pose propagation）：
\begin{equation}
  T_{k+1} = T_{k}\,\Delta T_{k},
  \label{eq:04-propagation}
\end{equation}
其中 $\Delta T_{k}\in SE(3)$ 是第 $k$ 步的运动增量。
\end{definition}

递推式~\eqref{eq:04-propagation} 的合理性来自位姿复合。设 $T_{AB}$ 表示「从 $B$ 系到 $A$ 系」的变换，则连续复合满足
\begin{equation}
  T_{AC} = T_{AB}\,T_{BC}.
  \label{eq:04-compose}
\end{equation}
把式~\eqref{eq:04-action} 代入验证：对任意点 ${}^{C}\bm{q}$，有
\begin{align}
  {}^{A}\bm{q}
  &= \bm{R}_{AB}\left(\bm{R}_{BC}\,{}^{C}\bm{q} + {}^{B}\bm{p}_{C}\right) + {}^{A}\bm{p}_{B} \notag\\
  &= \left(\bm{R}_{AB}\bm{R}_{BC}\right){}^{C}\bm{q}
     + \left(\bm{R}_{AB}\,{}^{B}\bm{p}_{C} + {}^{A}\bm{p}_{B}\right).
  \label{eq:04-compose-proof}
\end{align}
因此复合后的旋转与平移分别是
\begin{equation}
  \bm{R}_{AC} = \bm{R}_{AB}\bm{R}_{BC},
  \qquad
  {}^{A}\bm{p}_{C} = \bm{R}_{AB}\,{}^{B}\bm{p}_{C} + {}^{A}\bm{p}_{B}.
\end{equation}
逆变换同样可以逐项写出：在式~\eqref{eq:04-action} 两边左乘 $\bm{R}_{WB}^{\top}$，得
\begin{equation}
  {}^{B}\bm{q} = \bm{R}_{WB}^{\top}{}^{W}\bm{q} - \bm{R}_{WB}^{\top}{}^{W}\bm{p}_{B},
\end{equation}
即
\begin{equation}
  T_{BW} = \left(\bm{R}_{WB}^{\top},\ -\bm{R}_{WB}^{\top}\,{}^{W}\bm{p}_{B}\right).
  \label{eq:04-inverse}
\end{equation}
注意式~\eqref{eq:04-inverse} 中平移不是简单取负，而是「转回本体系再取负」——新手最常见的错误之一。

\section{位姿递推与指数映射}\label{sec:04-exp}

里程计输出的原始量通常是速度或加速度。例如轮式底盘给出线速度 $v$ 与角速度 $\omega$，IMU 给出角速度与比力。要把「速度乘时间」变成一个位姿增量 $\Delta T_{k}\in SE(3)$，就需要指数映射。

设第 $k$ 步内机体在本体系下的速度与角速度为常值，记为 $\bm{v}_{k},\bm{\omega}_{k}\in\mathbb{R}^{3}$，时间间隔 $\Delta t$。定义六维运动量
\begin{equation}
  \bm{\xi}_{k} = \begin{bmatrix} \bm{\rho}_{k} \\ \bm{\phi}_{k}\end{bmatrix}
  = \begin{bmatrix} \bm{v}_{k}\Delta t \\ \bm{\omega}_{k}\Delta t\end{bmatrix}
  \in\mathbb{R}^{6}.
  \label{eq:04-xi}
\end{equation}
则位姿增量为指数映射
\begin{equation}
  \Delta T_{k} = \exp\left(\bm{\xi}_{k}^{\wedge}\right),
  \qquad
  \bm{\xi}^{\wedge} = \begin{bmatrix} \bm{\phi}^{\wedge} & \bm{\rho}\\ \bm{0}^{\top} & 0\end{bmatrix}.
  \label{eq:04-exp-se3}
\end{equation}
下面把 $\exp$ 的每一块完整推出来。先定义反对称矩阵算子：对 $\bm{\omega}=[\omega_1,\omega_2,\omega_3]^{\top}$，
\begin{equation}
  \bm{\omega}^{\wedge} = \begin{bmatrix}
    0 & -\omega_3 & \omega_2\\
    \omega_3 & 0 & -\omega_1\\
    -\omega_2 & \omega_1 & 0
  \end{bmatrix},
  \qquad
  \bm{\omega}^{\wedge}\bm{x} = \bm{\omega}\times\bm{x}.
  \label{eq:04-skew}
\end{equation}
它的关键性质是 $\bm{\omega}^{\wedge\top} = -\bm{\omega}^{\wedge}$（反对称）。

\textbf{第一步：旋转矩阵的微分方程。}设机体以常值角速度 $\bm{\omega}$ 旋转，任意一个固连在机体上的向量 ${}^{B}\bm{q}$ 在世界系中的像为 ${}^{W}\bm{q}(t)=\bm{R}(t){}^{B}\bm{q}$。由于该向量在机体系中不变，有
\begin{equation}
  \dot{}^{W}\bm{q}(t) = \bm{R}(t)\left(\bm{\omega}\times{}^{B}\bm{q}\right)
  = \bm{R}(t)\bm{\omega}^{\wedge}{}^{B}\bm{q}.
\end{equation}
另一方面直接对 ${}^{W}\bm{q}=\bm{R}\,{}^{B}\bm{q}$ 求导得 $\dot{}^{W}\bm{q}=\dot{\bm{R}}{}^{B}\bm{q}$。比较两式，并注意到 ${}^{B}\bm{q}$ 任意，得到
\begin{equation}
  \dot{\bm{R}} = \bm{R}\,\bm{\omega}^{\wedge}.
  \label{eq:04-rdot}
\end{equation}

\textbf{第二步：把解写成级数。}常系数线性微分方程 \eqref{eq:04-rdot} 在初值 $\bm{R}(0)=\bm{I}$ 下的解可以用幂级数表示。把 $\bm{R}(t)=\sum_{n\ge 0}\frac{t^{n}}{n!}\bm{A}^{n}$ 代入式~\eqref{eq:04-rdot} 逐项求导（左端 $\sum \frac{t^{n-1}}{(n-1)!}\bm{A}^{n}$，右端 $\sum\frac{t^{n}}{n!}\bm{A}^{n}\bm{\omega}^{\wedge}$），比较 $t^{n}$ 的系数即知取 $\bm{A}=\bm{\omega}^{\wedge}$ 时逐项相等，因此
\begin{equation}
  \bm{R}(t) = \exp\left(t\bm{\omega}^{\wedge}\right)
  = \sum_{n=0}^{\infty}\frac{t^{n}}{n!}\left(\bm{\omega}^{\wedge}\right)^{n}.
  \label{eq:04-series}
\end{equation}

\textbf{第三步：化简级数。}令 $\bm{\phi}=\bm{\omega}t$，$\theta=\norm{\bm{\phi}}$，$\bm{u}=\bm{\phi}/\theta$（当 $\theta\ne0$）。先求单位轴幂次的闭式。取 $\bm{u}=\bm{e}_3=[0,0,1]^{\top}$，由式~\eqref{eq:04-skew} 直接算出
\begin{equation}
  \bm{e}_3^{\wedge} = \begin{bmatrix}0&-1&0\\1&0&0\\0&0&0\end{bmatrix},\qquad
  \left(\bm{e}_3^{\wedge}\right)^{2} = \begin{bmatrix}-1&0&0\\0&-1&0\\0&0&0\end{bmatrix}
  = \bm{e}_3\bm{e}_3^{\top} - \bm{I}.
\end{equation}
再用它算三次幂：
\begin{equation}
  \left(\bm{e}_3^{\wedge}\right)^{3}
  = \left(\bm{e}_3^{\wedge}\right)^{2}\bm{e}_3^{\wedge}
  = \begin{bmatrix}-1&0&0\\0&-1&0\\0&0&0\end{bmatrix}
    \begin{bmatrix}0&-1&0\\1&0&0\\0&0&0\end{bmatrix}
  = \begin{bmatrix}0&1&0\\-1&0&0\\0&0&0\end{bmatrix}
  = -\bm{e}_3^{\wedge}.
\end{equation}
对任意单位轴 $\bm{u}$，存在旋转矩阵 $\bm{Q}$ 使 $\bm{u}=\bm{Q}\bm{e}_3$，而 $\bm{u}^{\wedge}=\bm{Q}\bm{e}_3^{\wedge}\bm{Q}^{\top}$，故上述两条关系在一般轴下同样成立：
\begin{equation}
  \left(\bm{u}^{\wedge}\right)^{2} = \bm{u}\bm{u}^{\top} - \bm{I},
  \qquad
  \left(\bm{u}^{\wedge}\right)^{3} = -\bm{u}^{\wedge}.
  \label{eq:04-axis-powers}
\end{equation}
把式~\eqref{eq:04-axis-powers} 反复使用，任意整数次幂都可降到 $\bm{I}$、$\bm{u}^{\wedge}$、$(\bm{u}^{\wedge})^{2}$ 三档。代入式~\eqref{eq:04-series} 并按这三档重新分组（设 $\bm{\phi}=\theta\bm{u}$）：
\begin{align}
  \exp\left(\bm{\phi}^{\wedge}\right)
  &= \bm{I}
   + \left(\sum_{j=0}^{\infty}\frac{(-1)^{j}\theta^{2j+1}}{(2j+1)!}\right)\bm{u}^{\wedge}
   + \left(\sum_{j=1}^{\infty}\frac{(-1)^{j+1}\theta^{2j}}{(2j)!}\right)\left(\bm{u}^{\wedge}\right)^{2} \notag\\
  &= \bm{I} + \sin\theta\,\bm{u}^{\wedge}
   + \left(1-\cos\theta\right)\left(\bm{u}^{\wedge}\right)^{2}.
  \label{eq:04-rodrigues}
\end{align}
其中第一个括号正是 $\sin\theta$ 的泰勒级数，第二个括号是 $1-\cos\theta$ 的泰勒级数。式~\eqref{eq:04-rodrigues} 就是 \textbf{Rodrigues 公式}（李群指数映射的系统推导见 \cite{barfoot2017state,sola2017eskf}）。把 $\bm{u}^{\wedge}$ 与 $(\bm{u}^{\wedge})^{2}$ 按分量展开，令 $c=\cos\theta$、$s=\sin\theta$、$\bm{u}=[u_x,u_y,u_z]^{\top}$：
\begin{equation}
  \bm{R} = \begin{bmatrix}
    c+u_x^{2}(1-c) & u_xu_y(1-c)-u_zs & u_xu_z(1-c)+u_ys\\
    u_yu_x(1-c)+u_zs & c+u_y^{2}(1-c) & u_yu_z(1-c)-u_xs\\
    u_zu_x(1-c)-u_ys & u_zu_y(1-c)+u_xs & c+u_z^{2}(1-c)
  \end{bmatrix}.
  \label{eq:04-rotmat}
\end{equation}
\textbf{第四步：小角度近似。}当 $\theta\to0$ 时，$\sin\theta=\theta-\theta^{3}/6+O(\theta^{5})$，$1-\cos\theta=\theta^{2}/2+O(\theta^{4})$，因此
\begin{equation}
  \exp\left(\bm{\phi}^{\wedge}\right)
  = \bm{I} + \bm{\phi}^{\wedge} + \frac{1}{2}\left(\bm{\phi}^{\wedge}\right)^{2} + O(\theta^{3}),
  \qquad
  \bm{R}\approx\bm{I}+\bm{\phi}^{\wedge}\quad(\theta\ll1).
  \label{eq:04-rot-approx}
\end{equation}
这就是「一帧内转角很小」时的线性化，它是后面所有滤波与优化推导的核心近似。

\textbf{第五步：平移部分。}对 $SE(3)$ 的指数映射，把式~\eqref{eq:04-exp-se3} 写成块形式并求解微分方程。设 $T(t)=\exp(t\bm{\xi}^{\wedge})$，其块形式为
\begin{equation}
  T(t) = \begin{bmatrix}\bm{R}(t) & \bm{p}(t)\\ \bm{0}^{\top} & 1\end{bmatrix},
\end{equation}
两边对 $t$ 求导并利用 $T(t)$ 与 $\bm{\xi}^{\wedge}$ 交换，得到 $\dot{\bm{R}}=\bm{R}\bm{\phi}^{\wedge}$ 与 $\dot{\bm{p}}=\bm{R}\bm{\rho}$。第一式已由式~\eqref{eq:04-rodrigues} 解出。第二式两边积分：
\begin{equation}
  \bm{p}(t) = \int_{0}^{t}\bm{R}(s)\bm{\rho}\,\diff s
  = \int_{0}^{t}\left[\bm{I} + \frac{\sin(s\theta)}{\theta}\bm{\phi}^{\wedge}
  + \frac{1-\cos(s\theta)}{\theta^{2}}\left(\bm{\phi}^{\wedge}\right)^{2}\right]\bm{\rho}\,\diff s.
\end{equation}
逐项积分，注意 $\int_0^t\sin(s\theta)\,\diff s=(1-\cos(t\theta))/\theta$ 与 $\int_0^t(1-\cos(s\theta))\,\diff s=t-\sin(t\theta)/\theta$，并记 $t=1$（$\bm{\xi}$ 已包含时间），得到
\begin{equation}
  \bm{p} = \bm{J}\bm{\rho},\qquad
  \bm{J} = \bm{I} + \frac{1-\cos\theta}{\theta^{2}}\bm{\phi}^{\wedge}
  + \frac{\theta-\sin\theta}{\theta^{3}}\left(\bm{\phi}^{\wedge}\right)^{2}.
  \label{eq:04-J}
\end{equation}
于是完整的 $SE(3)$ 指数映射为
\begin{equation}
  \exp\left(\bm{\xi}^{\wedge}\right)
  = \begin{bmatrix}\exp\left(\bm{\phi}^{\wedge}\right) & \bm{J}\bm{\rho}\\ \bm{0}^{\top} & 1\end{bmatrix}.
  \label{eq:04-exp-final}
\end{equation}
当 $\norm{\bm{\phi}}$ 很小时，$\bm{J}\approx\bm{I}$，式~\eqref{eq:04-exp-final} 退化为
\begin{equation}
  \Delta T \approx \begin{bmatrix}\bm{I}+\bm{\phi}^{\wedge} & \bm{\rho}\\ \bm{0}^{\top} & 1\end{bmatrix}
  = \bm{I} + \bm{\xi}^{\wedge},
  \label{eq:04-exp-small}
\end{equation}
这正好对上式~\eqref{eq:04-propagation} 的直觉：$\Delta\bm{R}\approx\bm{I}+\bm{\omega}^{\wedge}\Delta t$，$\Delta\bm{p}\approx\bm{v}\Delta t$。

\begin{example}[一次增量到底转了多少]
设 $\bm{\omega}=[0,0,1]^{\top}\,\mathrm{rad/s}$，$\Delta t=0.1\,\mathrm{s}$，则 $\bm{\phi}=[0,0,0.1]^{\top}$，$\theta=0.1\,\mathrm{rad}$，转轴 $\bm{u}=\bm{e}_3$。由式~\eqref{eq:04-rodrigues}，$\sin0.1=0.099833$，$\cos0.1=0.995004$，故
\begin{equation}
  \Delta\bm{R} = \begin{bmatrix}0.995004 & -0.099833 & 0\\ 0.099833 & 0.995004 & 0\\ 0&0&1\end{bmatrix}
  \approx \begin{bmatrix}0.995 & -0.0998 & 0\\ 0.0998 & 0.995 & 0\\ 0&0&1\end{bmatrix},
\end{equation}
与一阶近似 $\bm{I}+0.1\bm{e}_3^{\wedge}$ 相比，对角元差约 $5\times10^{-6}$，在这个尺度上完全可以忽略。若同时有 $\bm{v}=[1,0,0]^{\top}\,\mathrm{m/s}$，则位移增量是 $\bm{J}\bm{\rho}\approx[0.1,0,0]^{\top}\,\mathrm{m}$，误差来自 $\bm{J}$ 的高阶项，量级约 $10^{-5}\,\mathrm{m}$。
\end{example}

\begin{keybox}{指数映射的三条实用结论}
\begin{enumerate}
  \item $\exp(\bm{\xi}^{\wedge})$ 把「本体系下的速度 $\times$ 时间」统一变成 $SE(3)$ 增量，旋转和平移可以一起处理；
  \item 小角度时 $\exp(\bm{\phi}^{\wedge})\approx\bm{I}+\bm{\phi}^{\wedge}$，误差量级 $O(\theta^{2})$，一般 $\theta<0.1\,\mathrm{rad}$（约 $5.7\degree$）时工程上可接受；
  \item 指数映射与矩阵乘法不可交换：$\exp(\bm{\xi}_1^{\wedge})\exp(\bm{\xi}_2^{\wedge})\ne\exp((\bm{\xi}_1+\bm{\xi}_2)^{\wedge})$。因此位姿递推必须做矩阵（或四元数）乘法，不能直接把增量相加。
\end{enumerate}
\end{keybox}

\section{里程计的四种来源}\label{sec:04-sources}

同一件事——估计本体的运动——可以用完全不同的物理量来测量。表~\ref{tab:04-sources} 给出四种主流方案的数量级对比。记住这张表比记住任何公式都有用：它决定了你在什么场景该信谁。

\begin{table}[htbp]
  \centering
  \caption{四类里程计的数量级对比（典型值，随器件与场景变化）}
  \label{tab:04-sources}
  \small
  \begin{tabular}{lcccc}
    \toprule
    来源 & 典型频率 & 直接观测量 & 相对精度（短期） & 主要失效场景 \\
    \midrule
    轮式编码器 & 50--100\,Hz & 轮转角 & 行程的 1--10\,\% & 打滑、悬空、地面沉降 \\
    IMU & 100--1000\,Hz & 角速度、比力 & 姿态 0.1--1$\degree$，位置漂移极快 & 零偏、振动、温漂 \\
    激光雷达 & 5--20\,Hz & 点云 & 行程的 0.1--1\,\% & 长廊、空旷、无结构 \\
    视觉相机 & 20--60\,Hz & 图像 & 行程的 0.5--5\,\% & 弱纹理、强光、动态物体 \\
    \bottomrule
  \end{tabular}
\end{table}

\textbf{轮式}（第~\ref{ch:05-wheel}~章）最便宜、频率最高、短时最稳，但它是「假设轮子不滑」的里程计：轮子一打滑，它测的就是轮子的转数而非车体的运动，误差立刻变成系统性偏差。\textbf{IMU}（第~\ref{ch:06-imu}~章）不依赖外部环境，短时角速度精度极高，是姿态估计的主力；但加速度要积分两次才得到位置，零偏误差随时间平方增长。\textbf{激光雷达}（第~\ref{ch:07-lidar-odom}~章）深度精度在厘米级，点云配准能给出很准的帧间位姿，是 FAST-LIO 系列的基础，代价是数据量大、对几何结构敏感。\textbf{视觉}（第~\ref{ch:08-visual-odom}~章）信息量最丰富，但受光照、纹理、动态物体影响，单目还有尺度不可观的问题。四种来源强弱互补，这正是第~\ref{ch:09-fusion}~章要讲的融合动机。

\section{累积误差与漂移}\label{sec:04-drift}

里程计最本质的缺陷是：它把每一步的误差都攒了下来。本节把这个「攒」字算清楚。

\textbf{第一步：复习期望与方差。}（定义见第~\ref{ch:02-math}~章。）设 $X$ 是连续型随机变量，密度为 $p(x)$，则
\begin{equation}
  \E[X] = \int x\,p(x)\,\diff x,
  \qquad
  \Var(X) = \E\left[(X-\E[X])^{2}\right] = \E[X^{2}] - \left(\E[X]\right)^{2}.
  \label{eq:04-var-def}
\end{equation}
方差衡量 $X$ 围绕均值的散布程度；它的平方根 $\sigma=\sqrt{\Var(X)}$ 叫标准差，与 $X$ 同量纲。

\textbf{第二步：两个独立变量之和的方差。}由式~\eqref{eq:04-var-def}，
\begin{align}
  \Var(X+Y)
  &= \E\left[(X+Y)^{2}\right] - \left(\E[X+Y]\right)^{2} \notag\\
  &= \E[X^{2}]+2\E[XY]+\E[Y^{2}]
     - \left(\E[X]^{2}+2\E[X]\E[Y]+\E[Y]^{2}\right) \notag\\
  &= \left(\E[X^{2}]-\E[X]^{2}\right) + \left(\E[Y^{2}]-\E[Y]^{2}\right)
     + 2\left(\E[XY]-\E[X]\E[Y]\right) \notag\\
  &= \Var(X)+\Var(Y)+2\Cov(X,Y).
  \label{eq:04-var-sum}
\end{align}
若 $X$ 与 $Y$ 独立，则 $\E[XY]=\E[X]\E[Y]$，协方差项为零，于是
\begin{equation}
  \Var(X+Y)=\Var(X)+\Var(Y).
  \label{eq:04-var-independent}
\end{equation}

\textbf{第三步：$n$ 步的累积误差。}设第 $i$ 步的位置误差为 $\epsilon_i$，各步相互独立、同分布，均值为零、方差为 $\sigma^{2}$。$n$ 步后的总误差为 $E_n=\sum_{i=1}^{n}\epsilon_i$。对式~\eqref{eq:04-var-independent} 反复使用 $n-1$ 次（每加一项只用一次两变量结论），得
\begin{equation}
  \Var(E_n) = \sum_{i=1}^{n}\Var(\epsilon_i) = n\sigma^{2},
  \qquad
  \sigma_{n} = \sqrt{n}\,\sigma.
  \label{eq:04-drift-var}
\end{equation}
这就是全书反复出现的\textbf{随机游走}（random walk）结论：误差的方差随步数线性增长，标准差随步数按 $\sqrt{n}$ 增长。

\textbf{第四步：换成时间与距离。}设里程计频率为 $f=1/\Delta t$，运行时间 $t$，则 $n=ft$，于是 $\sigma(t)=\sigma\sqrt{f\,t}$：它不是随时间线性增长，而是随时间开方增长，这一点常常被误解。但如果每步误差有非零均值 $b$（称为\textbf{系统误差}或\textbf{零偏}），则
\begin{equation}
  \E[E_n] = \sum_{i=1}^{n}\E[\epsilon_i] = nb,
  \label{eq:04-bias}
\end{equation}
偏差随时间线性增长，是比随机游走更危险的一类误差。真实里程计是两者叠加：随时间一长，系统误差必然主导。

\begin{example}[一步一厘米，一百秒后差多少]
设 $\Delta t=0.1\,\mathrm{s}$，每步位置误差标准差 $\sigma=0.01\,\mathrm{m}$，各步独立。运行 $t=100\,\mathrm{s}$ 共 $n=1000$ 步，则 $\sigma_n=\sqrt{1000}\times0.01\approx0.316\,\mathrm{m}$。若同样的误差里还藏着 $b=0.5\,\mathrm{mm}$ 的均值偏差，则 $1000$ 步后平均偏差是 $1000\times0.0005=0.5\,\mathrm{m}$，反而比随机部分更大。
\end{example}

\textbf{第五步：漂移率。}为了横向比较不同里程计，定义\textbf{漂移率}（drift rate）
\begin{equation}
  \rho_{\text{drift}} = \frac{\norm{\hat{\bm{p}}(t)-\bm{p}(t)}_{\text{RMS}}}{L(t)},
  \label{eq:04-drift-rate}
\end{equation}
即位置误差的均方根除以累计行程 $L(t)$，常用百分数表示。轮式里程计典型值 1--10\,\%，激光/视觉 0.1--1\,\%。注意漂移率是「误差/距离」而不是「误差/时间」，因为大多数里程计误差主要来自走过的距离与转弯次数，而不是纯粹的时间。

\begin{figure}[htbp]
  \centering
  \begin{tikzpicture}[x=1.15cm,y=0.9cm]
    \draw[-stealth] (0,0) -- (5.2,0) node[right]{\footnotesize 时间 $t$};
    \draw[-stealth] (0,0) -- (0,4.4) node[above]{\footnotesize 误差};
    \draw[thick, red!60!black, domain=0:4.6, samples=40] plot (\x, {0.75*\x});
    \draw[thick, structurecolor, domain=0:4.6, samples=40] plot (\x, {1.05*sqrt(\x)});
    \node[font=\footnotesize, red!60!black] at (3.3,3.4) {系统误差 $\propto t$};
    \node[font=\footnotesize, structurecolor] at (3.3,1.9) {随机误差 $\propto\sqrt{t}$};
    \draw[dashed, black!40] (2,0) -- (2,1.05*1.414) -- (0,1.05*1.414);
    \node[font=\footnotesize, below] at (2,0) {$2t_0$};
    \node[font=\footnotesize, left] at (0,1.48) {$\sqrt{2}\,\sigma_0$};
  \end{tikzpicture}
  \caption{两类误差的增长形状。随机误差按 $\sqrt{t}$ 增长（蓝线），系统误差按 $t$ 线性增长（红线）。长时间运行后，系统的零偏一定是主要矛盾，这也是为什么所有里程计都需要外部修正。}
  \label{fig:04-drift}
\end{figure}

\section{里程计的评价}\label{sec:04-eval}

评价里程计不能用「我觉得挺准」这种描述，工程上使用两个标准指标（工具链见第~\ref{ch:28-evaluation}~章，这里只给定义）。

\textbf{绝对轨迹误差 ATE。}设估计位姿序列为 $\{\hat{T}_i\}_{i=1}^{N}$，真值序列为 $\{T_i\}$。先求一个全局对齐变换 $\bm{S}\in SE(3)$（通常是第一帧对齐或 Umeyama 对齐，消除初始位姿的任意性），定义第 $i$ 个误差
\begin{equation}
  \bm{E}_i = T_i^{-1}\,\bm{S}\,\hat{T}_i \in SE(3).
  \label{eq:04-ate-error}
\end{equation}
绝对轨迹误差（absolute trajectory error, ATE）定义为平移部分的均方根：
\begin{equation}
  \mathrm{ATE}_{\text{RMSE}} = \sqrt{\frac{1}{N}\sum_{i=1}^{N}\norm{\mathrm{trans}\left(\bm{E}_i\right)}^{2}},
  \label{eq:04-ate}
\end{equation}
其中 $\mathrm{trans}(\cdot)$ 表示取 $SE(3)$ 的平移分量。

\textbf{相对位姿误差 RPE。}ATE 会随着轨迹变长而系统性变大，不适合单独评价「里程计本身」。相对位姿误差（relative pose error, RPE）只比较固定间隔 $\Delta$ 内的相对运动：
\begin{equation}
  \bm{E}_i = \left(T_i^{-1}T_{i+\Delta}\right)^{-1}
  \left(\hat{T}_i^{-1}\hat{T}_{i+\Delta}\right),
  \label{eq:04-rpe-error}
\end{equation}
再对平移部分取 RMSE 得到 $\mathrm{RPE}_{\text{RMSE}}$；对旋转部分取 $\mathrm{SO}(3)$ 上的角度误差 RMSE 得到 $\mathrm{RPE}_{\text{rot}}$。

为什么里程计要特别看 RPE？因为里程计的输出本来就是一系列相对增量，它没有全局原点；完美的里程计应当「任意两帧之间的相对变换都准」，式~\eqref{eq:04-rpe-error} 恰好把全局漂移消掉，只留下局部一致性。

\begin{keybox}{ATE 与 RPE 的使用场景}
\begin{itemize}
  \item \textbf{ATE}：评价「整个系统」——把里程计 + 回环 + 后端作为一个黑箱，看最终轨迹与真值的偏离。SLAM 论文一般报 ATE。
  \item \textbf{RPE}：评价「里程计本身」——只看局部相对运动，不受累计漂移干扰。调里程计参数（如 ICP 阈值、噪声协方差）时应多看 RPE。
  \item 两者必须一起看：ATE 小可能是回环在「救场」，掩盖了里程计本身很差；RPE 小但 ATE 大，说明里程计准但存在未修正的长期漂移。
\end{itemize}
\end{keybox}

\section{工程视角：里程计模块的输入输出}\label{sec:04-engineering}

不管底层用哪种传感器，一个里程计节点在工程上的接口几乎是固定的：输入带时间戳的传感器测量与外参 $T_{SB}$，输出每个时刻的位姿 $T_{OB}$（$O$ 为里程计系，连续、无跳变）及协方差，并约定 $T_{OB}(t_0)=\bm{I}$。全局修正由 \texttt{map}$\to$\texttt{odom} 这个额外的变换承担，绝不直接改里程计输出——这样短期控制器拿到的位姿不会因为一次回环而突变。

\begin{pseudo}{通用里程计递推循环}
  \item 初始化：$T_{OB}\leftarrow\bm{I}$，$t\leftarrow t_0$
  \item 当收到时间戳 $t_k$ 的测量 $\bm{z}_k$ 时：
  \item \quad 计算时间间隔 $\Delta t\leftarrow t_k-t_{k-1}$
  \item \quad 由 $\bm{z}_k$ 得到本体系速度估计 $(\hat{\bm{v}}_k,\hat{\bm{\omega}}_k)$
  \item \quad 构造增量 $\bm{\xi}_k\leftarrow[\hat{\bm{v}}_k\Delta t;\ \hat{\bm{\omega}}_k\Delta t]$
  \item \quad $T_{OB}\leftarrow T_{OB}\exp(\bm{\xi}_k^{\wedge})$
  \item \quad 按式~\eqref{eq:04-drift-var} 同款逻辑传播协方差
  \item \quad 发布 $T_{OB}$、协方差与 \texttt{odom}$\to$\texttt{base} 变换
\end{pseudo}

\begin{summary}

\begin{enumerate}
  \item 定位问题分为绝对定位与相对定位；里程计解决相对定位，特点是高频、连续、局部一致，但会累积漂移。
  \item 位姿 $T_{WB}=(\bm{R}_{WB},{}^{W}\bm{p}_B)$ 属于 $SE(3)$；复合 $T_{AC}=T_{AB}T_{BC}$ 对应旋转相乘、平移按「先转再加」；逆变换的平移是 $-\bm{R}^{\top}\bm{p}$ 而不是 $-\bm{p}$。
  \item 位姿递推是 $T_{k+1}=T_{k}\Delta T_k$，增量由指数映射给出：$\Delta T_k=\exp(\bm{\xi}_k^{\wedge})$；$\mathrm{SO}(3)$ 的闭式是 Rodrigues 公式 $\bm{R}=\bm{I}+\sin\theta\,\bm{u}^{\wedge}+(1-\cos\theta)(\bm{u}^{\wedge})^{2}$，小角度时 $\bm{R}\approx\bm{I}+\bm{\phi}^{\wedge}$。
  \item $SE(3)$ 指数映射为 $\exp(\bm{\xi}^{\wedge})=[\exp(\bm{\phi}^{\wedge}),\bm{J}\bm{\rho};\bm{0}^{\top},1]$，其中 $\bm{J}$ 由式~\eqref{eq:04-J} 给出，小角度时 $\bm{J}\approx\bm{I}$。
  \item 四种里程计来源各有短长：轮式便宜但怕打滑，IMU 高频但位置漂移快，激光/视觉几何精度高但怕退化与光照。
  \item 独立同分布误差的方差相加：$\Var(\sum\epsilon_i)=n\sigma^{2}$，标准差按 $\sqrt{n}$ 增长；有零偏时均值按 $n$ 线性增长，长时间后前者必被后者超过。
  \item 评价里程计用 ATE 与 RPE，前者看整体、后者看局部一致性；工程上里程计输出接在 \texttt{odom}$\to$\texttt{base}，全局修正接在 \texttt{map}$\to$\texttt{odom}。
\end{enumerate}


\end{summary}

\begin{problemset}

  \item 设初始位姿 $T_{WB}=(\bm{R}_{WB},{}^{W}\bm{p}_B)$，验证式~\eqref{eq:04-compose} 的逆变换为 $T_{BW}=(\bm{R}_{WB}^{\top},-\bm{R}_{WB}^{\top}{}^{W}\bm{p}_B)$，并用一个具体的二维例子（$\bm{R}$ 为 $90\degree$ 旋转，${}^{W}\bm{p}_B=[1,0,0]^{\top}$）算一遍。
  \item 用 Rodrigues 公式计算 $\bm{\omega}=[0,0,\pi/2]^{\top}\,\mathrm{rad/s}$、$\Delta t=1\,\mathrm{s}$ 时的 $\Delta\bm{R}$，并与 $\bm{I}+\bm{\phi}^{\wedge}$ 比较误差的量级。
  \item 证明：若每步误差独立且标准差为 $\sigma$，走 $n$ 步后的位置标准差为 $\sigma\sqrt{n}$；若每步误差都带同样大小的均值 $b$，则期望误差为 $nb$。当 $n$ 从 $10^{2}$ 增到 $10^{4}$ 时，两者比例如何变化？
  \item 某轮式里程计在 $10\,\mathrm{m}$ 直线行程中位置误差 $0.3\,\mathrm{m}$，在 $50\,\mathrm{m}$ 行程中位置误差 $1.2\,\mathrm{m}$。分别计算漂移率，并判断误差更像随机游走还是系统偏差。
  \item 一个团队报告「ATE 为 $0.05\,\mathrm{m}$，RPE 为 $0.4\,\mathrm{m}$」，请你解释这组数字可能意味着什么，以及应当先排查哪个模块。
\end{problemset}

\begin{reading}

\begin{itemize}
  \item Sola 的 ESKF 教程（\href{https://arxiv.org/abs/1711.02508}{arXiv:1711.02508}）：$SO(3)/SE(3)$ 求导与误差状态。
  \item Barfoot《State Estimation for Robotics》（\href{https://books.google.com/books?id=jfZFDwAAQBAJ}{Cambridge, 2017}）：李群上的估计理论。
\end{itemize}


\end{reading}

\chapter{轮式里程计：从编码器到航迹推算}\label{ch:05-wheel}

\chapterintro{轮式里程计是所有移动机器人最基础、最廉价、也最常被低估的定位手段：两个编码器加初中数学就能给出 100\,Hz 的位姿。但「便宜」不等于「简单」——轮子会打滑，轮距会标不准，地面不会永远平。本章从编码器的一个脉冲讲起，一路推到差速底盘的圆弧航迹推算、麦轮底盘的逆运动学，以及一套可落地的标定流程。学完本章，你应当能独立写出一个差速里程计节点，并知道它的误差从哪来、该怎么标、协方差该怎么填。}

\section{编码器：把一个脉冲变成一段位移}\label{sec:05-encoder}

车轮的运动最终由编码器变成数字。\textbf{增量式编码器}（incremental encoder）输出 A、B 两路正交方波，转子每转过一个固定角度产生一个脉冲，数脉冲得到「转了多少」，比较 A、B 相位得到「往哪转」；\textbf{绝对式编码器}（absolute encoder）每个角度对应唯一编码，上电即知绝对角度。本书主要用增量式，因为便宜、用定时器正交解码即可读取。

设编码器\textbf{线数}（每转脉冲数，pulses per revolution, PPR）为 $N$，电机与轮之间的\textbf{减速比}为 $i$，车轮半径为 $r$。工程上普遍采用\textbf{四倍频}（quadrature）计数：A、B 两相的上升沿与下降沿都计数，因此电机每转一圈得到 $4N$ 个计数。电机转 $i$ 圈车轮转一圈，于是车轮每转一圈对应
\begin{equation}
  C_{\text{rev}} = 4Ni
  \label{eq:05-counts-rev}
\end{equation}
个计数。车轮转一圈走过弧长 $2\pi r$，所以单个计数对应的行进距离是
\begin{equation}
  \Delta s = \frac{2\pi r}{4Ni}.
  \label{eq:05-ds}
\end{equation}
式~\eqref{eq:05-ds} 里有三个容易搞错的量：$N$ 是\textbf{电机轴}上的线数而不是轮轴上的；$i$ 是减速比，$i>1$ 表示减速；四倍频的因子 4 只在 A/B 正交解码时成立，若只数单相上升沿则没有这个 4。

\begin{example}[一个计数有多长]
取 $N=13$、四倍频后电机每转 $52$ 个计数、减速比 $i=30$、轮半径 $r=0.065\,\mathrm{m}$，则
\begin{equation}
  \Delta s=\frac{2\pi\times0.065}{4\times13\times30}
  =\frac{0.4084}{1560}\approx2.62\times10^{-4}\,\mathrm{m},
\end{equation}
即每个计数约 $0.26\,\mathrm{mm}$；若车轮以 $1\,\mathrm{m/s}$ 前进，则每秒产生约 $3817$ 个计数。
\end{example}

单个计数太粗怎么办？用时间换精度：在固定周期 $\Delta t$ 内数计数 $n$，得到平均速度 $\hat{v}=n\Delta s/\Delta t$。但计数本身有 $\pm1$ 的量化误差，因此速度估计的量化噪声为 $\Delta s/\Delta t$；仍用上例，$\Delta t=10\,\mathrm{ms}$ 时噪声约 $2.6\,\mathrm{cm/s}$，对 $1\,\mathrm{m/s}$ 的车是 $2.6\,\%$。这解释了为什么低速时编码器测速很差。

\begin{keybox}{编码器分辨率的三个坑}
\begin{enumerate}
  \item \textbf{线数不等于计数}：$N$ 线四倍频后是 $4N$ 计数；某些厂商标的是「每转计数」。
  \item \textbf{减速比方向}：$i=30$ 表示电机转 30 圈车轮转 1 圈；若把 $i$ 当成增速比，结果会差 $i^{2}$ 倍。
  \item \textbf{分辨率不等于精度}：$0.26\,\mathrm{mm}$ 只是量化步长，真实精度还取决于轮半径误差、打滑与安装偏心。
\end{enumerate}
\end{keybox}

\section{两轮差速运动学}\label{sec:05-diff}

\subsection{从轮速到车体速度}

考虑两轮差速底盘：左右两个驱动轮位于车体两侧，轮心连线过车体中心，轮距为 $L$。记左、右轮的线速度（轮缘相对地面的速度）为 $v_L$、$v_R$，车体中心的速度为 $v$、角速度为 $\omega$。刚体上任意一点的速度等于「中心速度 + 角速度叉乘该点相对中心的位置」。右轮相对中心的位置是 $(0,-L/2)$（机体系 $x$ 向前、$y$ 向左），而 $\bm{\omega}=[0,0,\omega]^{\top}$，于是右轮接地点的速度为
\begin{equation}
  \bm{v}_{R} = \bm{v} + \bm{\omega}\times\bm{r}_{R}
  = \begin{bmatrix}v\\0\\0\end{bmatrix}
  + \begin{bmatrix}0\\0\\\omega\end{bmatrix}\times\begin{bmatrix}0\\-L/2\\0\end{bmatrix}
  = \begin{bmatrix}v+\dfrac{\omega L}{2}\\[2mm] 0\\0\end{bmatrix}.
  \label{eq:05-wheel-point}
\end{equation}
轮子不打滑时，轮缘速度沿前进方向的分量就是 $v_R$，横向分量被地面约束为零。左轮同理，于是
\begin{equation}
  v_R = v + \frac{\omega L}{2},
  \qquad
  v_L = v - \frac{\omega L}{2}.
  \label{eq:05-wheel-v}
\end{equation}
把两式相加、相减，解出车体速度：
\begin{equation}
  v = \frac{v_R+v_L}{2},
  \qquad
  \omega = \frac{v_R-v_L}{L}.
  \label{eq:05-v-omega}
\end{equation}
式~\eqref{eq:05-v-omega} 的物理含义很直白：$v$ 是两轮速度的平均（决定直线快慢），$\omega$ 是两轮速度之差除以轮距（决定转向快慢）。当 $v_R=v_L$ 时 $\omega=0$，车走直线；当 $v_R=-v_L$ 时 $v=0$，车原地转。

\subsection{非完整约束}

车体在世界系中的位姿是三维的 $(x,y,\theta)$，但上式只有两个独立输入 $(v,\omega)$。把车体速度投影到世界系：
\begin{equation}
  \dot{x}=v\cos\theta,\qquad \dot{y}=v\sin\theta,\qquad \dot{\theta}=\omega.
  \label{eq:05-kinematics}
\end{equation}
消去 $v$ 得到
\begin{equation}
  \dot{x}\sin\theta-\dot{y}\cos\theta = 0.
  \label{eq:05-nonholonomic}
\end{equation}
这是一个不能积分成「位置约束」的微分约束，叫\textbf{非完整约束}（nonholonomic constraint）：普通轮式车不能侧向平移，必须先转再走。它决定了轮式车不能像全向底盘那样直接横移，也决定了路径规划必须考虑曲率（见第~\ref{ch:19-global-planning}~章）。

\begin{figure}[htbp]
  \centering
  \begin{tikzpicture}[scale=0.82]
    \draw[thick, fill=structurecolor!8, rounded corners=3pt] (-1.15,-0.7) rectangle (1.35,0.7);
    \draw[thick, fill=black!18] (-0.55,0.7) rectangle (0.55,0.95);
    \draw[thick, fill=black!18] (-0.55,-0.95) rectangle (0.55,-0.7);
    \draw[<->, thin] (1.05,0.825) -- (1.05,-0.825) node[midway, right, font=\footnotesize] {轮距 $L$};
    \draw[-stealth, very thick, structurecolor] (0,0) -- (1.15,0) node[above right, font=\footnotesize] {$v$};
    \draw[-stealth, thin] (0,0) -- (0,1.35) node[above, font=\footnotesize] {$y_B$};
    \draw[-stealth, thin] (0,0) -- (-1.55,0) node[left, font=\footnotesize] {$x_B$};
    \draw[-stealth, thick, red!60!black] (0.35,0.35) arc (45:200:0.5);
    \node[font=\footnotesize, red!60!black] at (0.12,-0.5) {$\omega$};
    \draw[-stealth, thick] (-1.0,0.825) -- (-0.25,0.825) node[pos=0.5, above, font=\footnotesize] {$v_L$};
    \draw[-stealth, thick] (-1.0,-0.825) -- (-0.25,-0.825) node[pos=0.5, below, font=\footnotesize] {$v_R$};
    \draw[dashed, black!40] (-1.15,0) -- (1.35,0);
    \node[font=\footnotesize] at (1.95,0.45) {车体};
  \end{tikzpicture}
  \caption{两轮差速底盘的俯视图（机体系 $x_B$ 向前、$y_B$ 向左）。车体中心速度 $v$ 是两轮速度的平均，角速度 $\omega$ 由两轮速度差除以轮距决定；轮距 $L$ 的标定误差会直接变成角速度的比例误差。}
  \label{fig:05-diff}
\end{figure}

\section{离散航迹推算：圆弧模型与直线近似}\label{sec:05-dead-reckoning}

里程计在每个采样周期执行一次「已知上一时刻位姿与本周期的 $v,\omega$，求新位姿」。这就是\textbf{航迹推算}（dead reckoning，航位推算），它的概率版本（把轮速误差建模成噪声）见 \cite{thrun2005probabilistic}。关键在于：一个周期内 $v$、$\omega$ 怎么变化？最常用的是\textbf{常值假设}：周期内 $v$、$\omega$ 不变。

\subsection{圆弧模型}

在常值假设且 $\omega\ne0$ 时，车体走一段圆弧。由式~\eqref{eq:05-kinematics}，第 $k$ 个周期初始位姿为 $(x_k,y_k,\theta_k)$，则时刻 $t\in[0,\Delta t]$ 的朝向为 $\theta(t)=\theta_k+\omega t$，而 $\dot{x}=v\cos\theta$ 从 $0$ 到 $\Delta t$ 的积分为 $x_{k+1}-x_k=\int_{0}^{\Delta t}v\cos(\theta_k+\omega t)\diff t$。做换元 $u=\theta_k+\omega t$，则 $\diff u=\omega\,\diff t$，$\diff t=\diff u/\omega$，积分限从 $t=0,\Delta t$ 变为 $u=\theta_k,\ \theta_k+\omega\Delta t$：
\begin{equation}
  x_{k+1}-x_k
  = \frac{v}{\omega}\int_{\theta_k}^{\theta_k+\omega\Delta t}\cos u\,\diff u
  = \frac{v}{\omega}\left[\sin\left(\theta_k+\omega\Delta t\right)-\sin\theta_k\right].
  \label{eq:05-arc-x}
\end{equation}
同理对 $\dot{y}=v\sin\theta$ 做同样的换元（$\int\sin u\,\diff u=-\cos u$），得 $y_{k+1}-y_k=-\frac{v}{\omega}[\cos(\theta_k+\omega\Delta t)-\cos\theta_k]$。再加上朝向更新 $\theta_{k+1}=\theta_k+\omega\Delta t$，就得到完整的\textbf{圆弧模型}：
\begin{equation}
  \begin{aligned}
    x_{k+1} &= x_k + \frac{v}{\omega}\left[\sin\left(\theta_k+\omega\Delta t\right)-\sin\theta_k\right],\\
    y_{k+1} &= y_k - \frac{v}{\omega}\left[\cos\left(\theta_k+\omega\Delta t\right)-\cos\theta_k\right],\\
    \theta_{k+1} &= \theta_k+\omega\Delta t.
  \end{aligned}
  \label{eq:05-arc}
\end{equation}

\textbf{圆弧模型的几何解释。}记转角 $\Delta\theta=\omega\Delta t$、转弯半径 $R=v/\omega$。用和差化积公式 $\sin a-\sin b=2\cos\frac{a+b}{2}\sin\frac{a-b}{2}$，取 $a=\theta_k+\Delta\theta$、$b=\theta_k$，式~\eqref{eq:05-arc-x} 变为
\begin{equation}
  x_{k+1}-x_k
  = \frac{v}{\omega}\cdot2\cos\left(\theta_k+\frac{\Delta\theta}{2}\right)\sin\frac{\Delta\theta}{2}
  = \underbrace{2R\sin\frac{\Delta\theta}{2}}_{\text{弦长 }d}
  \cos\left(\theta_k+\frac{\Delta\theta}{2}\right).
\end{equation}
对 $y$ 分量做同样的换元积分，再使用和差化积公式 $\cos a-\cos b=-2\sin\frac{a+b}{2}\sin\frac{a-b}{2}$（取 $a=\theta_k+\Delta\theta$、$b=\theta_k$），得到
\begin{equation}
  y_{k+1}-y_k
  = -\frac{v}{\omega}\left[\cos\left(\theta_k+\Delta\theta\right)-\cos\theta_k\right]
  = d\sin\left(\theta_k+\frac{\Delta\theta}{2}\right),
\end{equation}
其中 $d=2R\sin(\Delta\theta/2)$ 就是同一步的弦长。因此圆弧模型可以统一地理解为：\textbf{一步的位移等于弦长 $d$，方向是中间朝向 $\theta_k+\Delta\theta/2$}。这比硬背公式好记，也直接给出下面的直线极限。

\subsection{直线模型与极限}

当 $\omega\to0$ 时 $v/\omega\to\infty$、$\Delta\theta\to0$，式~\eqref{eq:05-arc} 是 $0/0$ 型极限。对 $x$ 分量做一阶展开：
\begin{equation}
  \sin\left(\theta_k+\Delta\theta\right)-\sin\theta_k
  = \sin\theta_k\cos\Delta\theta+\cos\theta_k\sin\Delta\theta-\sin\theta_k.
\end{equation}
当 $\Delta\theta\to0$ 时 $\cos\Delta\theta\to1$、$\sin\Delta\theta\to\Delta\theta+O(\Delta\theta^{3})$，故上式 $=\Delta\theta\cos\theta_k+O(\Delta\theta^{2})$。于是
\begin{equation}
  x_{k+1}-x_k = \frac{v}{\omega}\cdot\Delta\theta\cos\theta_k + O(\Delta\theta^{2})
  = v\Delta t\cos\theta_k + O(\Delta t^{2}),
\end{equation}
其中用了 $v\Delta\theta/\omega = v\Delta t$。对 $y$ 分量同理可得 $y_{k+1}-y_k\to v\Delta t\sin\theta_k$。于是 $\omega=0$ 时圆弧模型退化为\textbf{直线模型}：
\begin{equation}
  x_{k+1}=x_k+v\Delta t\cos\theta_k,\quad
  y_{k+1}=y_k+v\Delta t\sin\theta_k,\quad
  \theta_{k+1}=\theta_k.
  \label{eq:05-line}
\end{equation}

\subsection{三种离散化误差对比}

同一个物理运动，用不同离散公式近似，误差量级不同。以式~\eqref{eq:05-kinematics} 的 $x$ 分量为例如下对比。

\textbf{欧拉法}用周期初朝向：$x_{k+1}=x_k+v\Delta t\cos\theta_k$。真实位移为 $\int_0^{\Delta t}v\cos(\theta_k+\omega t)\diff t$，把余弦在 $t=0$ 展开得 $\cos(\theta_k+\omega t)=\cos\theta_k-\omega t\sin\theta_k+O(t^{2})$，逐项积分：
\begin{equation}
  \Delta x_{\text{true}} = v\Delta t\cos\theta_k - \frac{1}{2}v\omega\Delta t^{2}\sin\theta_k+O(\Delta t^{3}),
\end{equation}
故欧拉法局部误差是 $O(v\omega\Delta t^{2})$，是\textbf{一阶}方法。\textbf{中点法}用中间朝向 $x_{k+1}=x_k+v\Delta t\cos(\theta_k+\omega\Delta t/2)$，把余弦在区间中点展开后奇次项对称抵消，局部误差为 $O(\Delta t^{3})$，是\textbf{二阶}方法。\textbf{圆弧公式}在常值假设下没有离散化误差。所以急转弯时不要图省事用欧拉法；经验上 $\abs{\omega}\Delta t<0.05\,\mathrm{rad}$ 时欧拉法误差可忽略，否则应当升级。

\begin{keybox}{圆弧模型的三条实践结论}
\begin{enumerate}
  \item 一步走的是圆弧：位移是弦长 $d=2R\sin(\Delta\theta/2)$，方向是中间朝向；$\omega\to0$ 时 $d\to v\Delta t$，退化为直线。
  \item 圆弧公式与直线公式可以统一写成「$d$ 沿中间朝向投影」，代码里用一个 \texttt{if} 判断 $\abs{\omega}$ 是否小于阈值即可。
  \item 当 $v$、$\omega$ 在一个周期内变化剧烈时，任何离散模型都有误差；解决办法是提高采样频率。
\end{enumerate}
\end{keybox}

\begin{pseudo}{差速轮式航迹推算（单周期）}
  \item 读取左、右轮累计计数 $c_L,c_R$，与上一周期作差得 $\Delta c_L,\Delta c_R$
  \item 换算成轮子行进距离：$\Delta s_{L,R}\leftarrow\Delta c_{L,R}\cdot\Delta s$
  \item 计算轮速：$v_{L,R}\leftarrow\Delta s_{L,R}/\Delta t$
  \item 由式~\eqref{eq:05-v-omega} 得 $v\leftarrow(v_R+v_L)/2$，$\omega\leftarrow(v_R-v_L)/L$
  \item 若 $\abs{\omega}<\epsilon$（$\epsilon$ 取 $10^{-3}\,\mathrm{rad/s}$）：
  \item \quad $x\leftarrow x+v\Delta t\cos\theta$， $y\leftarrow y+v\Delta t\sin\theta$
  \item 否则：
  \item \quad $x\leftarrow x+\frac{v}{\omega}[\sin(\theta+\omega\Delta t)-\sin\theta]$
  \item \quad $y\leftarrow y-\frac{v}{\omega}[\cos(\theta+\omega\Delta t)-\cos\theta]$
  \item 更新 $\theta\leftarrow\theta+\omega\Delta t$，发布位姿与协方差
\end{pseudo}

\section{麦轮与全向底盘}\label{sec:05-mecanum}

差速底盘只能前进和转向，横向移动要靠多次转向完成。麦轮（Mecanum wheel）解决了这个问题：每个轮子的轮缘上装有一圈与轮轴成 $45\degree$ 的辊子，轮子既能沿自身滚动方向驱动，又能沿辊子轴向被动滑动；四个轮子配合，底盘就能同时实现前后、左右、自转三个自由度的运动。

\subsection{麦轮逆运动学推导}

采用机体系：$x$ 向前、$y$ 向左、$\omega$ 逆时针为正。四个轮心位置为 $\bm{r}_1=(a,b)$、$\bm{r}_2=(a,-b)$、$\bm{r}_3=(-a,b)$、$\bm{r}_4=(-a,-b)$，其中 $2a$ 为前后轮距、$2b$ 为左右轮距（$a,b>0$）。第 $i$ 个轮子的接地点速度（把该点当作刚体上一点）为
\begin{equation}
  \bm{u}_i = \bm{v}+\omega\hat{\bm{z}}\times\bm{r}_i
  = \begin{bmatrix}v_x-\omega r_{iy}\\ v_y+\omega r_{ix}\end{bmatrix},
  \label{eq:05-contact-vel}
\end{equation}
其中 $\bm{v}=[v_x,v_y]^{\top}$ 是底盘中心速度。麦轮的约束是：接地点速度沿「轮子约束方向」的分量必须等于轮子线速度 $s_i$，而垂直方向的分量由辊子自由滑动吸收。记第 $i$ 个轮子的约束方向单位向量为 $\bm{e}_i$，则
\begin{equation}
  s_i = \bm{u}_i\cdot\bm{e}_i.
  \label{eq:05-mecanum-proj}
\end{equation}
由于辊子与轮轴成 $45\degree$，四个轮子的约束方向与坐标轴也成 $45\degree$。取 $\bm{e}_1=\frac{1}{\sqrt2}[1,-1]^{\top}$、$\bm{e}_2=\frac{1}{\sqrt2}[1,1]^{\top}$、$\bm{e}_3=\frac{1}{\sqrt2}[1,1]^{\top}$、$\bm{e}_4=\frac{1}{\sqrt2}[1,-1]^{\top}$，逐轮代入式~\eqref{eq:05-contact-vel} 与式~\eqref{eq:05-mecanum-proj}。以轮 1 为例，$\bm{r}_1=(a,b)$，由式~\eqref{eq:05-contact-vel} 得 $\bm{u}_1=[v_x-\omega b,\ v_y+\omega a]^{\top}$，于是
\begin{equation}
  s_1=\frac{1}{\sqrt{2}}\left[(v_x-\omega b)-(v_y+\omega a)\right]
  =\frac{1}{\sqrt{2}}\left[v_x-v_y-(a+b)\omega\right].
\end{equation}
其余三轮同样计算，得到麦轮的\textbf{逆运动学}（由底盘速度求四个轮速）：
\begin{equation}
  \begin{aligned}
    s_1 &= \tfrac{1}{\sqrt{2}}\left[v_x-v_y-(a+b)\omega\right],\\
    s_2 &= \tfrac{1}{\sqrt{2}}\left[v_x+v_y+(a+b)\omega\right],\\
    s_3 &= \tfrac{1}{\sqrt{2}}\left[v_x+v_y-(a+b)\omega\right],\\
    s_4 &= \tfrac{1}{\sqrt{2}}\left[v_x-v_y+(a+b)\omega\right].
  \end{aligned}
  \label{eq:05-mecanum-inv}
\end{equation}

\subsection{正运动学：从轮速回到速度}

里程计需要的是反方向：已知四个轮速，求底盘速度。把式~\eqref{eq:05-mecanum-inv} 的四式做三种线性组合：四式相加时含 $v_y$ 与 $\omega$ 的项成对抵消，只剩 $4v_x$；取 $s_2+s_3-s_1-s_4$ 时含 $v_x$ 与 $\omega$ 的项抵消，只剩 $4v_y$；取 $s_2+s_4-s_1-s_3$ 时只剩 $4(a+b)\omega$。整理得
\begin{equation}
  \begin{aligned}
    v_x &= \frac{\sqrt{2}}{4}\left(s_1+s_2+s_3+s_4\right),\\
    v_y &= \frac{\sqrt{2}}{4}\left(s_2+s_3-s_1-s_4\right),\\
    \omega &= \frac{\sqrt{2}}{4(a+b)}\left(s_2+s_4-s_1-s_3\right),
  \end{aligned}
  \qquad\Longleftrightarrow\qquad
  \begin{bmatrix}v_x\\v_y\\\omega\end{bmatrix}
  = \frac{\sqrt{2}}{4}\begin{bmatrix}
    1 & 1 & 1 & 1\\
    -1 & 1 & 1 & -1\\
    -\frac{1}{a+b} & \frac{1}{a+b} & -\frac{1}{a+b} & \frac{1}{a+b}
  \end{bmatrix}
  \begin{bmatrix}s_1\\s_2\\s_3\\s_4\end{bmatrix}.
  \label{eq:05-mecanum-fwd}
\end{equation}
积分 $(v_x,v_y,\omega)$ 即可得到位姿；工程上把轮子角速度乘以轮半径得到 $s_i$ 后直接代入。这张矩阵就是麦轮里程计的全部「物理模型」，剩下的事情都是噪声与打滑。

\subsection{为什么麦轮里程计更容易漂}

式~\eqref{eq:05-mecanum-fwd} 是在「辊子与地面纯滚动、无侧向滑移」假设下推出来的，而现实中这个假设很差：辊子只在与地面接触的瞬间近似纯滚动，其余时刻存在横向滑动，滑移量随负载和速度变化；四个轮子的半径、减速比、编码器分辨率、轮胎磨损很难一致，线性叠加会把差异放大成系统性偏差；麦轮对地面起伏极敏感，一个 1\,mm 的凸起就足以改变接触点速度方向；横移时四个轮子都在侧向滑移，编码器却只测到滚动分量，$v_y$ 的估计必然偏小。因此在 RoboMaster 这类高速对抗场景中，麦轮里程计通常只用来做短时速度反馈，位置估计要靠 IMU、激光或视觉兜底（第~\ref{ch:06-imu}~章与第~\ref{ch:09-fusion}~章）。

\begin{figure}[htbp]
  \centering
  \begin{tikzpicture}[scale=0.72]
    \draw[thick, fill=structurecolor!8, rounded corners=3pt] (-1.6,-1.3) rectangle (1.6,1.3);
    \draw[thick, fill=black!18] (1.2,0.95) rectangle (1.55,1.25);
    \draw[thick, fill=black!18] (1.2,-1.25) rectangle (1.55,-0.95);
    \draw[thick, fill=black!18] (-1.55,0.95) rectangle (-1.2,1.25);
    \draw[thick, fill=black!18] (-1.55,-1.25) rectangle (-1.2,-0.95);
    \draw[-stealth, thick, red!60!black] (1.375,1.1) -- ++(0.5,-0.5);
    \node[font=\scriptsize, red!60!black] at (2.1,0.75) {$\bm{e}_1$};
    \draw[-stealth, thick, red!60!black] (1.375,-1.1) -- ++(0.5,0.5);
    \node[font=\scriptsize, red!60!black] at (2.1,-0.75) {$\bm{e}_2$};
    \draw[-stealth, thick, red!60!black] (-1.375,1.1) -- ++(0.5,0.5);
    \node[font=\scriptsize, red!60!black] at (-2.1,0.75) {$\bm{e}_3$};
    \draw[-stealth, thick, red!60!black] (-1.375,-1.1) -- ++(0.5,-0.5);
    \node[font=\scriptsize, red!60!black] at (-2.1,-0.75) {$\bm{e}_4$};
    \draw[-stealth, very thick, structurecolor] (0,0) -- (0.9,0.45) node[above, font=\footnotesize] {$\bm{v}$};
    \draw[-stealth, thick, structurecolor!70] (2.35,0) arc (0:150:0.45);
    \node[font=\footnotesize, structurecolor] at (2.75,0.35) {$\omega$};
    \draw[dashed, black!40] (-1.6,-0.75) -- (1.6,-0.75);
    \draw[dashed, black!40] (0.75,-1.3) -- (0.75,1.3);
    \node[font=\footnotesize, below] at (0,-1.5) {$2b$（左右轮距）};
    \node[font=\footnotesize, right] at (1.72,0.2) {$2a$（前后轮距）};
    \node[font=\scriptsize, black!60] at (0.1,0.15) {辊子 $45\degree$};
  \end{tikzpicture}
  \caption{麦轮底盘俯视图。四轮约束方向 $\bm{e}_i$ 与坐标轴成 $45\degree$；每个轮子只能约束接地点速度在 $\bm{e}_i$ 方向的分量，另一个方向靠辊子滑动吸收——这既是全向移动的来源，也是打滑误差的来源。}
  \label{fig:05-mecanum}
\end{figure}

\section{误差来源与标定}\label{sec:05-calib}

\subsection{误差从哪里来}

由式~\eqref{eq:05-v-omega} 与式~\eqref{eq:05-ds}，轮式里程计的参数误差可以分成两类。\textbf{轮半径误差}：$r$ 有相对误差 $e_r$ 时，$\Delta s$ 以及 $v$、$\omega$ 都同比放大，整条轨迹被等比例放大而形状不变，这叫\textbf{尺度误差}。\textbf{轮距误差}：$L$ 有相对误差 $e_L$ 时，$\omega$ 同比缩小而 $v$ 不变，转弯半径 $R=v/\omega$ 被放大，轨迹形状被扭曲，且转角越大误差越大。此外还有打滑（轮子转但车没走）、地面沉降（有效半径变小）、轮胎磨损与安装偏心等。前两类是\textbf{系统性}误差，可以通过标定压到很小；打滑是\textbf{随机且时变}的，只能靠外部传感器修正。

\subsection{标定流程：直线定半径、旋转定轮距}

\textbf{第一步：直线标定轮半径。}让机器人沿直线走过一段已知距离 $D$（卷尺测量），记录编码器给出的左右轮累计距离 $\hat{s}_R$、$\hat{s}_L$，并记编码器算出的行驶距离为 $\hat{D}=\frac{1}{2}(\hat{s}_R+\hat{s}_L)$。设标称半径为 $r_{\text{nom}}$，真实半径为 $r$。由 $\Delta s\propto r$ 可知真实距离与编码器距离之比等于半径之比：
\begin{equation}
  \frac{D}{\hat{D}} = \frac{r}{r_{\text{nom}}}
  \ \Longrightarrow\
  r = r_{\text{nom}}\cdot\frac{D}{\hat{D}}.
  \label{eq:05-r-calib}
\end{equation}
这里 $r$ 是式~\eqref{eq:05-ds} 中唯一的线性因子，其余量（$N$、$i$、$4$）在标定期间固定，所以这个比例关系成立。

\textbf{第二步：旋转标定轮距。}让机器人原地旋转 $n$ 整圈（$n$ 取 10 圈以减小单圈误差），实际转角 $\theta_{\text{true}}=2\pi n$。编码器给出的转角为 $\hat{\theta}=(\hat{s}_R-\hat{s}_L)/L_{\text{nom}}$，其中 $\hat{s}_{R,L}$ 是左右轮累计走过的弧长。真实轮距满足同一关系，由 $\hat{\theta}/\theta_{\text{true}}=L_{\text{true}}/L_{\text{nom}}$ 得
\begin{equation}
  L = L_{\text{nom}}\cdot\frac{\theta_{\text{true}}}{\hat{\theta}}.
  \label{eq:05-L-calib}
\end{equation}
注意：原地旋转时左右轮弧长只依赖半径与转角，而两者之差又正比于轮距，所以旋转实验\textbf{只标定 $L$}；半径必须由直线实验单独标定，否则两者会混在一起。

\textbf{第三步：多组实验的最小二乘拟合。}式~\eqref{eq:05-L-calib} 只用了一次实验。更稳的做法是做 $M$ 组不同圈数的旋转实验，最小化
\begin{equation}
  J(L)=\sum_{j=1}^{M}\left(\theta_j-\frac{d_j}{L}\right)^{2},
  \qquad d_j=\hat{s}_{R,j}-\hat{s}_{L,j},
  \label{eq:05-ls-cost}
\end{equation}
其中 $\theta_j$ 是第 $j$ 组实验的真实转角，$d_j$ 是左右轮累计距离之差。对 $L$ 求导并令其为零，由链式法则
\begin{equation}
  \frac{\diff J}{\diff L}
  = \sum_{j=1}^{M}2\left(\theta_j-\frac{d_j}{L}\right)\cdot\frac{d_j}{L^{2}}=0.
\end{equation}
两边乘以 $L^{2}/2$（$L\ne0$），再移项得正规方程 $\sum_j\theta_j d_j=\sum_j d_j^{2}/L$，解出
\begin{equation}
  \hat{L} = \frac{\sum_{j=1}^{M}d_j^{2}}{\sum_{j=1}^{M}\theta_j d_j}.
  \label{eq:05-ls-solution}
\end{equation}
目标函数对 $1/L$ 是二次的，所以有闭式解。工程上可以把 $r$ 与 $L$ 一起做二维最小二乘，但要注意可辨识性：直线实验只约束 $r$，旋转实验只约束 $L$，两组实验缺一不可。

\begin{keybox}{标定的正确顺序}
\begin{enumerate}
  \item 先检查硬件：轮胎磨损、轮子偏心、编码器线数与减速比是否和配置一致；
  \item 直线标定 $r$：走 10\,m 以上，往返各三次取平均，消除打滑不对称；
  \item 旋转标定 $L$：原地转 10 圈，正反各一次，用式~\eqref{eq:05-ls-solution} 拟合；
  \item 回到实际场地验证漂移率：走一个 20\,m 的矩形回路，看位置闭合误差；若仍大，说明打滑是主因，标定救不了。
\end{enumerate}
\end{keybox}

\section{工程实现：ROS 2 中的轮式里程计}\label{sec:05-engineering}

在 ROS 2 中，轮式里程计的标准接口是三件事：话题 \texttt{/odom}（类型 \texttt{nav\_msgs/msg/Odometry}，包含位姿、速度及它们的协方差）、TF 变换 \texttt{odom}$\to$\texttt{base\_link}（由里程计节点维护，不跳变），以及约定 \texttt{odom} 的初始位姿为单位阵、之后只由里程计递推。协方差是 $6\times6$ 矩阵按行优先展开的 36 个数，通常只填 $x,y,\theta$ 对应的对角元（索引 0、7、35），数值随累计行程增大。

协方差的量级可以这样估：设单步距离误差标准差为 $\sigma_s$，走 $n$ 步后按独立误差假设有 $\sigma_x^{2}\approx n\sigma_s^{2}$（这里用了第~\ref{ch:04-odom-intro}~章方差相加的结论）。很多团队的 \texttt{odom} 协方差是随手填的常量，这会让下游融合算法（第~\ref{ch:09-fusion}~章）错误地信任里程计，是排错时的常见根源。

差速里程计的核心计算只有十几行，但几个符号一错就全错，下面给出参考实现（注释与字符串为英文）：

\begin{lstlisting}[language=C++, caption={差速里程计核心计算（节选）}, label={lst:05-core}]
// inputs: wheel distance increments (m) and time step (s)
double v = 0.5 * (ds_l + ds_r) / dt;
double w = (ds_r - ds_l) / dt / wheel_separation;
if (std::fabs(w) < 1e-3) {
    x += v * dt * std::cos(theta);
    y += v * dt * std::sin(theta);
} else {
    x += v / w * (std::sin(theta + w * dt) - std::sin(theta));
    y -= v / w * (std::cos(theta + w * dt) - std::cos(theta));
    theta += w * dt;
}
\end{lstlisting}

参数含义：\texttt{wheel\_separation} 就是式~\eqref{eq:05-v-omega} 中的 $L$，\texttt{wheel\_radius} 决定 \texttt{dist\_per\_count}，\texttt{publish\_rate} 决定 $\Delta t$，\texttt{enable\_odom\_tf} 为真时节点自己发布 \texttt{odom}$\to$\texttt{base\_link}，与外部的 TF 广播器二选一。

\begin{summary}

\begin{enumerate}
  \item 增量式编码器四倍频后电机每转 $4N$ 计数，车轮每转 $4Ni$ 计数，单计数位移 $\Delta s=2\pi r/(4Ni)$。
  \item 差速底盘由刚体速度关系得 $v_R=v+\omega L/2$、$v_L=v-\omega L/2$，反解得 $v=(v_R+v_L)/2$、$\omega=(v_R-v_L)/L$；轮式底盘满足非完整约束 $\dot{x}\sin\theta-\dot{y}\cos\theta=0$。
  \item 常值 $v,\omega$ 假设下航迹推算有闭式圆弧解：位移等于弦长 $d=2R\sin(\Delta\theta/2)$，方向为中间朝向；$\omega\to0$ 时退化为直线模型。
  \item 欧拉法局部误差 $O(\Delta t^{2})$，中点法与圆弧模型更高阶；急转弯时应使用圆弧模型或提高采样频率。
  \item 麦轮逆运动学由接地点速度在约束方向上的投影推出，正运动学是它的线性反解；麦轮漂移主要来自辊子滑移与四轮不一致。
  \item 轮半径误差导致轨迹尺度误差，轮距误差导致曲率误差；直线实验标定 $r$，旋转实验用最小二乘标定 $L$。
  \item 工程上里程计发布 \texttt{/odom} 与 \texttt{odom}$\to$\texttt{base\_link}，协方差应按累计行程传播。
\end{enumerate}


\end{summary}

\begin{problemset}

  \item 编码器 $N=1024$ 线、四倍频、直驱、$r=0.1\,\mathrm{m}$，求单计数位移与 $10\,\mathrm{ms}$ 周期下的速度量化噪声。
  \item 差速车 $L=0.5\,\mathrm{m}$，左右轮 $0.6$、$0.9\,\mathrm{m/s}$，求 $v$、$\omega$、$R$ 及 1\,s 后的位姿（初始朝向 0）。
  \item 用式~\eqref{eq:05-arc} 计算 $v=1\,\mathrm{m/s}$、$\omega=2\,\mathrm{rad/s}$、$\Delta t=0.1\,\mathrm{s}$、$\theta_k=0$ 的一步位移，并与欧拉法、中点法比较误差量级。
  \item 麦轮底盘 $a=0.25\,\mathrm{m}$、$b=0.22\,\mathrm{m}$，四个轮速均为 $1\,\mathrm{m/s}$，求底盘速度；若 $s_4$ 因打滑偏小 $10\,\%$，估算 $\omega$ 的偏差。
  \item 用式~\eqref{eq:05-ls-solution} 处理 $(d_j,\theta_j)=(1.0,6.28)$、$(2.0,12.57)$、$(3.0,18.85)$，求 $\hat{L}$ 并说明它为何只适用于「原地旋转」。
\end{problemset}

\begin{reading}

\begin{itemize}
  \item ROS 2 \texttt{diff\_drive\_controller} 文档（\href{https://control.ros.org/}{control.ros.org}）：参数与协方差语义。
  \item Thrun 等《Probabilistic Robotics》第 5 章：里程计运动模型的概率表述。
  \item 麦轮底盘厂家手册：核对轮距、辊子角度与运动学矩阵的符号约定。
\end{itemize}


\end{reading}

\chapter{IMU 姿态解算与航迹推算}\label{ch:06-imu}

\chapterintro{IMU 是本篇最特殊的传感器：它测的不是位移，而是角速度与比力；数据最「脏」，噪声、零偏、温漂全在里面；但它又是唯一在轮子打滑、激光退化、相机过曝时仍然工作的传感器。本章把 IMU 从「一个报姿态的模块」讲清楚：测量模型怎么来、噪声怎么离散化、姿态怎么用四元数更新、横滚俯仰为什么能从重力读出来而偏航不能、互补滤波与 Mahony 为什么有效、ESKF 的 15 维公式每一块怎么推、以及为什么纯 IMU 位置积分几十秒就废。公式较多，请动笔跟着推。}

\section{IMU 测量模型与噪声}\label{sec:06-model}

IMU 通常包含三轴陀螺仪与三轴加速度计。陀螺仪测的是机体系相对惯性系的角速度；加速度计测的是\textbf{比力}（specific force），即单位质量所受的非引力合力，在地表附近近似等于「加速度减重力」。沿用全书符号：$\bm{g}=[0,0,-9.81]^{\top}$ 为世界系下的重力向量，$\bm{\omega}$ 为机体系下的真实角速度，$\bm{a}_W$ 为机体系原点在世界系下的真实加速度。记 IMU 测量值为 $\tilde{\bm{\omega}}$、$\tilde{\bm{a}}$，则
\begin{equation}
  \tilde{\bm{\omega}} = \bm{\omega} + \bm{b}_g + \bm{n}_g,
  \qquad
  \tilde{\bm{a}} = \bm{R}_{WB}^{\top}\left(\bm{a}_W - \bm{g}\right) + \bm{b}_a + \bm{n}_a.
  \label{eq:06-measure}
\end{equation}
逐项解释：$\bm{b}_g,\bm{b}_a$ 是陀螺仪、加速度计的\textbf{零偏}（bias），是一条缓慢变化的非零常值；$\bm{n}_g,\bm{n}_a$ 是零均值高斯白噪声；$\bm{R}_{WB}^{\top}(\bm{a}_W-\bm{g})$ 说明加速度计测的是本体系下的比力——静止时 $\bm{a}_W=0$，读数不是零，而是 $-\bm{R}^{\top}\bm{g}$，这就是为什么静置的 IMU 加速度读数约为 $9.81\,\mathrm{m/s^2}$ 且指向「上」。零偏本身也在缓慢游走，建模为
\begin{equation}
  \dot{\bm{b}}_g = \bm{n}_{bg},
  \qquad
  \dot{\bm{b}}_a = \bm{n}_{ba},
  \label{eq:06-bias-rw}
\end{equation}
其中 $\bm{n}_{bg},\bm{n}_{ba}$ 也是零均值白噪声。这种「白噪声的积分」叫\textbf{随机游走}（random walk）。

\subsection{从连续噪声到离散噪声}

仿真与滤波都需要离散时间的噪声参数，而传感器手册给的是连续时间的噪声密度。下面把两套参数的关系推出来。连续时间白噪声 $\bm{n}(t)\in\mathbb{R}^{3}$ 定义为
\begin{equation}
  \E\left[\bm{n}(t)\right]=\bm{0},
  \qquad
  \E\left[\bm{n}(t)\bm{n}(t')^{\top}\right]=\sigma_c^{2}\,\delta(t-t')\,\bm{I},
  \label{eq:06-white-cont}
\end{equation}
其中 $\delta(\cdot)$ 是狄拉克函数，$\sigma_c$ 是\textbf{连续噪声密度}，单位是「物理量单位 / $\sqrt{\mathrm{Hz}}$」（陀螺仪常见的 $0.01\degree/\mathrm{s}/\sqrt{\mathrm{Hz}}$ 即属此类）。滤波器在周期 $\Delta t$ 内实际用到的是该区间上的平均噪声
\begin{equation}
  \bm{n}_d = \frac{1}{\Delta t}\int_{t_k}^{t_k+\Delta t}\bm{n}(t)\,\diff t.
\end{equation}
它的方差由式~\eqref{eq:06-white-cont} 的二重积分得到：
\begin{align}
  \E\left[\bm{n}_d\bm{n}_d^{\top}\right]
  &= \frac{1}{\Delta t^{2}}\int_{t_k}^{t_k+\Delta t}\!\!\int_{t_k}^{t_k+\Delta t}
     \E\left[\bm{n}(t)\bm{n}(t')^{\top}\right]\diff t\,\diff t' \notag\\
  &= \frac{\sigma_c^{2}}{\Delta t^{2}}\int_{t_k}^{t_k+\Delta t}\!\!\int_{t_k}^{t_k+\Delta t}
     \delta(t-t')\,\bm{I}\,\diff t\,\diff t' \notag\\
  &= \frac{\sigma_c^{2}}{\Delta t^{2}}\int_{t_k}^{t_k+\Delta t}\bm{I}\,\diff t
  = \frac{\sigma_c^{2}}{\Delta t}\,\bm{I}.
  \label{eq:06-white-disc}
\end{align}
第二个等号用了「$\delta$ 函数只在 $t=t'$ 处非零」的性质：先对 $t'$ 积分得到 $\bm{I}$，再对 $t$ 积分得到 $\Delta t\,\bm{I}$。于是离散噪声的标准差为
\begin{equation}
  \sigma_d = \frac{\sigma_c}{\sqrt{\Delta t}}.
  \label{eq:06-sigma-disc}
\end{equation}
注意 $\Delta t$ 越小，$\sigma_d$ 越大：单次采样的噪声更「尖」，但积分一段时间后的总方差 $\sigma_d^{2}\Delta t=\sigma_c^{2}\Delta t$ 随时间线性增长，与物理直觉一致。

对零偏随机游走式~\eqref{eq:06-bias-rw} 在一个周期上积分，得 $\bm{b}_{k+1}-\bm{b}_k=\int_{t_k}^{t_k+\Delta t}\bm{n}_b(t)\diff t$，其方差为
\begin{equation}
  \E\left[\left(\bm{b}_{k+1}-\bm{b}_k\right)\left(\bm{b}_{k+1}-\bm{b}_k\right)^{\top}\right]
  = \sigma_{bc}^{2}\,\Delta t\,\bm{I},
\end{equation}
即零偏离散化后的标准差为
\begin{equation}
  \sigma_{bd} = \sigma_{bc}\sqrt{\Delta t}.
  \label{eq:06-bias-disc}
\end{equation}
与白噪声相反，随机游走的离散标准差随 $\Delta t$ 的平方根增大。

\begin{example}[把手册参数变成滤波器参数]
某 IMU 陀螺噪声密度 $0.01\degree/\mathrm{s}/\sqrt{\mathrm{Hz}}$，零偏随机游走 $0.001\degree/\mathrm{s}/\sqrt{\mathrm{s}}$，滤波周期 $\Delta t=5\,\mathrm{ms}$。由式~\eqref{eq:06-sigma-disc}，单周期角速度噪声标准差为 $0.01/\sqrt{0.005}\approx0.141\degree/\mathrm{s}$；由式~\eqref{eq:06-bias-disc}，单周期零偏随机游走标准差为 $0.001\times\sqrt{0.005}\approx7.1\times10^{-5}\,\degree/\mathrm{s}$。两者相差三个数量级，这就是为什么陀螺零偏可以近似当常值、而噪声必须每个周期都处理。
\end{example}

\section{姿态表示与四元数更新}\label{sec:06-quaternion}

\textbf{姿态解算}（attitude estimation）指利用陀螺仪积分、加速度计重力观测（以及磁力计）估计机体相对世界系的旋转 $\bm{R}_{WB}$。工程上几乎都用单位四元数存储：$\bm{q}=[q_w,q_x,q_y,q_z]^{\top}$，$\norm{\bm{q}}=1$。四元数乘法定义为
\begin{equation}
  \bm{q}_1\otimes\bm{q}_2
  = \begin{bmatrix}
    w_1w_2-\bm{v}_1^{\top}\bm{v}_2\\
    w_1\bm{v}_2+w_2\bm{v}_1+\bm{v}_1\times\bm{v}_2
  \end{bmatrix},
  \label{eq:06-qmul}
\end{equation}
其中 $\bm{v}=[q_x,q_y,q_z]^{\top}$ 是四元数的向量部分。单位四元数到旋转矩阵的完整换算是
\begin{equation}
  \bm{R}(\bm{q}) = \begin{bmatrix}
    1-2(q_y^{2}+q_z^{2}) & 2(q_xq_y-q_wq_z) & 2(q_xq_z+q_wq_y)\\
    2(q_xq_y+q_wq_z) & 1-2(q_x^{2}+q_z^{2}) & 2(q_yq_z-q_wq_x)\\
    2(q_xq_z-q_wq_y) & 2(q_yq_z+q_wq_x) & 1-2(q_x^{2}+q_y^{2})
  \end{bmatrix}.
  \label{eq:06-q2R}
\end{equation}

\textbf{运动学方程。}设机体角速度在机体系下为 $\bm{\omega}$，考察极短时间 $\diff t$ 内的姿态变化。由于 $\bm{\omega}$ 是机体系下的量，增量旋转应右乘：
\begin{equation}
  \bm{q}(t+\diff t) = \bm{q}(t)\otimes\exp\left(\tfrac{1}{2}\bm{\omega}\diff t\right).
\end{equation}
当 $\diff t$ 很小时，四元数指数映射一阶近似为 $\exp(\tfrac{1}{2}\bm{\omega}\diff t)\approx[1,\ \tfrac{1}{2}\bm{\omega}\diff t]$（完整形式为 $[\cos(\theta/2),\ \bm{u}\sin(\theta/2)]$，$\theta=\norm{\bm{\omega}}\diff t$）。代入并取极限：
\begin{align}
  \frac{\bm{q}(t+\diff t)-\bm{q}(t)}{\diff t}
  &= \frac{1}{\diff t}\left(\bm{q}(t)\otimes\begin{bmatrix}1\\ \tfrac{1}{2}\bm{\omega}\diff t\end{bmatrix}
     - \bm{q}(t)\otimes\begin{bmatrix}1\\ \bm{0}\end{bmatrix}\right) \notag\\
  &= \bm{q}(t)\otimes\begin{bmatrix}0\\ \tfrac{1}{2}\bm{\omega}\end{bmatrix}
  = \frac{1}{2}\,\bm{q}(t)\otimes\begin{bmatrix}0\\ \bm{\omega}\end{bmatrix}.
\end{align}
于是得到四元数运动学方程
\begin{equation}
  \dot{\bm{q}} = \frac{1}{2}\,\bm{q}\otimes\begin{bmatrix}0\\ \bm{\omega}\end{bmatrix}.
  \label{eq:06-qdot}
\end{equation}
离散实现中直接用估计角速度做增量更新：
\begin{equation}
  \bm{q}_{k+1} = \bm{q}_k\otimes\exp\left(\tfrac{1}{2}\left(\tilde{\bm{\omega}}_k-\hat{\bm{b}}_{g,k}\right)\Delta t\right),
  \label{eq:06-q-update}
\end{equation}
指数项按 $\theta=\norm{\bm{\omega}}\Delta t$、$\bm{u}=\bm{\omega}/\theta$ 算成 $[\cos(\theta/2),\ \bm{u}\sin(\theta/2)]$，更新后归一化。

\begin{keybox}{四元数更新的两个易错点}
\begin{enumerate}
  \item 左乘还是右乘：$\bm{q}_{WB}\otimes\delta\bm{q}$ 是「绕机体系轴转」，$\delta\bm{q}\otimes\bm{q}_{WB}$ 是「绕世界系轴转」。用机体系角速度必须右乘。
  \item 归一化：数值积分会让范数缓慢偏离 1，每次更新后应归一化，否则式~\eqref{eq:06-q2R} 得到的矩阵不再正交。
\end{enumerate}
\end{keybox}

\section{由加速度计估计重力方向与横滚俯仰}\label{sec:06-gravity}

静止或匀速运动时 $\bm{a}_W\approx\bm{0}$，式~\eqref{eq:06-measure} 的加速度计读数（忽略噪声与零偏）为
\begin{equation}
  \tilde{\bm{a}} = -\bm{R}_{WB}^{\top}\bm{g}.
  \label{eq:06-static-acc}
\end{equation}
记重力大小 $g_0=9.81$，则 $-\bm{g}=[0,0,g_0]^{\top}$。设 $\bm{R}_{WB}=\bm{R}_z(\psi)\bm{R}_y(\theta)\bm{R}_x(\phi)$（先横滚 $\phi$、再俯仰 $\theta$、最后偏航 $\psi$），下面展开式~\eqref{eq:06-static-acc}。先算两个基本旋转矩阵的乘积：
\begin{equation}
  \bm{R}_y(\theta)\bm{R}_x(\phi)
  = \begin{bmatrix}
    \cos\theta & 0 & \sin\theta\\
    0 & 1 & 0\\
    -\sin\theta & 0 & \cos\theta
  \end{bmatrix}
  \begin{bmatrix}
    1 & 0 & 0\\
    0 & \cos\phi & -\sin\phi\\
    0 & \sin\phi & \cos\phi
  \end{bmatrix}
  = \begin{bmatrix}
    \cos\theta & \sin\theta\sin\phi & \sin\theta\cos\phi\\
    0 & \cos\phi & -\sin\phi\\
    -\sin\theta & \cos\theta\sin\phi & \cos\theta\cos\phi
  \end{bmatrix}.
\end{equation}
记这个矩阵为 $\bm{M}$。由于 $\bm{R}_{WB}=\bm{R}_z(\psi)\bm{M}$，而 $\bm{R}_z$ 只作用于 $x,y$ 分量、不改变 $z$ 分量，所以 $\bm{R}_{WB}^{\top}\bm{e}_z$（即重力方向）的 $z$ 分量等于 $\bm{M}^{\top}\bm{e}_z$ 的 $z$ 分量。$\bm{M}^{\top}$ 的第三行就是 $\bm{M}$ 的第三列 $[\sin\theta\cos\phi,\ -\sin\phi,\ \cos\theta\cos\phi]$，于是
\begin{equation}
  \tilde{\bm{a}} = g_0\begin{bmatrix}
    \sin\theta\cos\phi\\
    -\sin\phi\\
    \cos\theta\cos\phi
  \end{bmatrix}.
  \label{eq:06-acc-expand}
\end{equation}
由第二、三分量解出横滚，由第一、三分量解出俯仰：
\begin{equation}
  \phi = \arctan2\left(-\tilde{a}_y,\ \tilde{a}_z\right),
  \qquad
  \theta = \arctan2\left(\tilde{a}_x,\ \sqrt{\tilde{a}_y^{2}+\tilde{a}_z^{2}}\right).
  \label{eq:06-roll-pitch}
\end{equation}
有些资料写成 $\phi=\arctan2(a_y,a_z)$、$\theta=\arctan2(-a_x,\sqrt{a_y^2+a_z^2})$，那是把重力方向取反的约定；两种写法只差整体符号，使用时必须与自己的加速度计符号约定一致。

\textbf{为什么偏航不可观测？}式~\eqref{eq:06-acc-expand} 中没有 $\psi$：绕世界系 $z$ 轴的旋转矩阵 $\bm{R}_z(\psi)$ 被 $\bm{R}_z^{\top}$ 作用在 $\bm{e}_z$ 上时满足 $\bm{R}_z^{\top}\bm{e}_z=\bm{e}_z$，因此 $\bm{R}_{WB}^{\top}\bm{g}$ 与 $\psi$ 无关。物理上，重力方向只提供「哪边是下」，无法提供「哪边是北」——绕重力轴的旋转不改变任何加速度计读数。这就是所谓\textbf{重力对齐}（gravity alignment）：加速度计只能把姿态对齐到重力方向（两个自由度），第三个自由度（偏航）必须靠磁力计、轮速、激光或视觉提供。

\begin{figure}[htbp]
  \centering
  \begin{tikzpicture}[scale=0.72]
    \draw[-stealth, thick] (0,0) -- (3.6,0) node[right, font=\footnotesize] {$x_W$};
    \draw[-stealth, thick] (0,0) -- (0,3.4) node[above, font=\footnotesize] {$z_W$（上）};
    \draw[-stealth, very thick, red!60!black] (1.2,2.8) -- (1.2,1.2) node[midway, right, font=\footnotesize] {$\bm{g}$};
    \draw[thick, fill=structurecolor!10, rotate around={-25:(1.2,1.7)}] (0.55,1.5) rectangle (1.85,1.9);
    \draw[-stealth, structurecolor, rotate around={-25:(1.2,1.7)}] (1.2,1.7) -- (1.85,1.7) node[right, font=\footnotesize] {$x_B$};
    \draw[-stealth, structurecolor, rotate around={-25:(1.2,1.7)}] (1.2,1.7) -- (1.2,2.35) node[above, font=\footnotesize] {$z_B$};
    \draw[dashed, black!40] (1.2,1.7) -- (2.2,1.7);
    \draw (1.85,1.7) arc (0:-25:0.65);
    \node[font=\footnotesize] at (2.35,1.5) {$\theta$};
    \node[font=\footnotesize, align=center] at (1.0,-0.4) {静止时加速度计测的是比力\\$\tilde{\bm{a}}=-\bm{R}_{WB}^{\top}\bm{g}$};
  \end{tikzpicture}
  \caption{加速度计测量重力方向示意。静止时读数只与「$z_B$ 轴相对重力的倾角」有关，与绕重力轴的偏航角 $\psi$ 无关，因此重力只能观测横滚与俯仰两个自由度。}
  \label{fig:06-gravity}
\end{figure}

\section{互补滤波与 Mahony 滤波}\label{sec:06-complementary}

陀螺积分短期准、长期漂；加速度计长期无漂、短期噪声大。两者恰好互补，于是有了\textbf{互补滤波}（complementary filter）。设真实角度为 $\theta$，陀螺积分得到 $\theta_g$（含低频漂移），加速度计得到 $\theta_a$（含高频噪声）。取一阶低通、高通滤波器
\begin{equation}
  H_{\text{LP}}(s)=\frac{\omega_c}{s+\omega_c},
  \qquad
  H_{\text{HP}}(s)=\frac{s}{s+\omega_c},
  \label{eq:06-filters}
\end{equation}
其中 $\omega_c$ 是截止角频率。两者之和恒为 1：$H_{\text{LP}}+H_{\text{HP}}=(s+\omega_c)/(s+\omega_c)=1$，于是
\begin{equation}
  \hat{\theta}(s) = H_{\text{LP}}(s)\,\theta_a(s) + H_{\text{HP}}(s)\,\theta_g(s).
  \label{eq:06-complementary}
\end{equation}
这就是互补滤波：两个通道的传递函数互补成 1，低频信加速度计（不受漂移影响），高频信陀螺（不受噪声影响）。用后向差分 $s\approx(1-z^{-1})/\Delta t$ 离散化，得
\begin{equation}
  \hat{\theta}_k = \alpha\left(\hat{\theta}_{k-1}+\omega_k\Delta t\right) + (1-\alpha)\,\theta_{a,k},
  \qquad
  \alpha = \frac{\omega_c\Delta t}{1+\omega_c\Delta t}.
  \label{eq:06-comp-discrete}
\end{equation}
工程上常写成 $\alpha=\tau/(\tau+\Delta t)$，$\tau$ 是时间常数。$\alpha$ 越接近 1 越信任陀螺，越接近 0 越信任加速度计；典型值 $0.95$--$0.99$。

\textbf{Mahony 滤波。}互补滤波逐轴处理横滚俯仰，难以直接推广到三维。Mahony 的做法是构造一个三维「姿态误差」向量，再用 PI 控制器把它反馈到角速度上。加速度计测得的方向归一化为 $\hat{\bm{a}}=\tilde{\bm{a}}/\norm{\tilde{\bm{a}}}$，当前姿态估计预测的「上」方向为 $\hat{\bm{g}}_b=\hat{\bm{R}}^{\top}\bm{e}_z$。定义误差
\begin{equation}
  \bm{e} = \hat{\bm{g}}_b\times\hat{\bm{a}}.
  \label{eq:06-mahony-err}
\end{equation}
\textbf{为什么这样定义？}设真实姿态与估计姿态差一个小角度 $\delta\bm{\theta}$，即 $\bm{R}=\hat{\bm{R}}\exp(\delta\bm{\theta}^{\wedge})$。真实「上」方向在机体系下为 $\bm{u}=\bm{R}^{\top}\bm{e}_z$，代入 $\bm{R}^{\top}=\exp(-\delta\bm{\theta}^{\wedge})\hat{\bm{R}}^{\top}\approx(\bm{I}-\delta\bm{\theta}^{\wedge})\hat{\bm{R}}^{\top}$：
\begin{equation}
  \bm{u} \approx \left(\bm{I}-\delta\bm{\theta}^{\wedge}\right)\hat{\bm{g}}_b
  = \hat{\bm{g}}_b - \delta\bm{\theta}\times\hat{\bm{g}}_b
  = \hat{\bm{g}}_b + \hat{\bm{g}}_b\times\delta\bm{\theta},
\end{equation}
其中最后一步用了 $-\delta\bm{\theta}\times\hat{\bm{g}}_b=\hat{\bm{g}}_b\times\delta\bm{\theta}$。又因为静止时 $\hat{\bm{a}}$ 就是 $\bm{u}$，代入式~\eqref{eq:06-mahony-err}：
\begin{equation}
  \bm{e} = \hat{\bm{g}}_b\times\left(\hat{\bm{g}}_b+\hat{\bm{g}}_b\times\delta\bm{\theta}\right)
  = \hat{\bm{g}}_b\times\left(\hat{\bm{g}}_b\times\delta\bm{\theta}\right)
  = \hat{\bm{g}}_b\left(\hat{\bm{g}}_b^{\top}\delta\bm{\theta}\right)-\delta\bm{\theta}
  = -\left(\bm{I}-\bm{P}\right)\delta\bm{\theta},
  \label{eq:06-mahony-linear}
\end{equation}
最后一步用了向量三重积公式 $\bm{a}\times(\bm{b}\times\bm{c})=\bm{b}(\bm{a}\cdot\bm{c})-\bm{c}(\bm{a}\cdot\bm{b})$ 与 $\bm{P}=\hat{\bm{g}}_b\hat{\bm{g}}_b^{\top}$。式~\eqref{eq:06-mahony-linear} 说明两件事：$\bm{e}$ 等于「姿态误差中垂直于重力方向的那一部分」取负；平行于重力方向（即偏航）的误差被投影矩阵 $\bm{I}-\bm{P}$ 滤掉，所以 Mahony 只能修正横滚俯仰。

既然 $\bm{e}$ 指向「把估计转回真值」的方向，把它按 PI 律加进角速度即可收敛：
\begin{equation}
  \bm{\omega}_{\text{cmd}} = \tilde{\bm{\omega}} - \hat{\bm{b}}_g + K_p\,\bm{e} + K_i\int\bm{e}\,\diff t,
  \qquad
  \dot{\hat{\bm{b}}}_g = -K_i\,\bm{e}.
  \label{eq:06-mahony}
\end{equation}
比例项 $K_p\bm{e}$ 负责快速纠正倾斜；积分项同时估计陀螺零偏——当姿态稳定、$\bm{e}\to0$ 时 $\hat{\bm{b}}_g$ 停止变化，残差被积分吸收。注意 $\bm{e}$ 恒为零只说明估计方向与加速度计一致，偏航仍自由漂移，所以 Mahony 只适用于「横滚俯仰必须准、偏航交给别的传感器」的场景（第~\ref{ch:09-fusion}~章）。

\section{ESKF：误差状态卡尔曼滤波}\label{sec:06-eskf}

\subsection{先复习普通卡尔曼滤波}

线性高斯系统 $\bm{x}_{k+1}=\bm{F}_k\bm{x}_k+\bm{w}_k$、$\bm{z}_k=\bm{H}_k\bm{x}_k+\bm{v}_k$，其中 $\bm{w}_k\sim\mathcal{N}(\bm{0},\bm{Q}_k)$、$\bm{v}_k\sim\mathcal{N}(\bm{0},\bm{R}_k)$ 且相互独立。卡尔曼滤波 \cite{kalman1960} 的预测与更新为
\begin{equation}
  \hat{\bm{x}}_k^{-}=\bm{F}_{k-1}\hat{\bm{x}}_{k-1},
  \qquad
  \bm{P}_k^{-}=\bm{F}_{k-1}\bm{P}_{k-1}\bm{F}_{k-1}^{\top}+\bm{Q}_{k-1},
  \label{eq:06-kf-predict}
\end{equation}
\begin{equation}
  \bm{K}_k=\bm{P}_k^{-}\bm{H}_k^{\top}\left(\bm{H}_k\bm{P}_k^{-}\bm{H}_k^{\top}+\bm{R}_k\right)^{-1},
  \qquad
  \hat{\bm{x}}_k=\hat{\bm{x}}_k^{-}+\bm{K}_k\left(\bm{z}_k-\bm{H}_k\hat{\bm{x}}_k^{-}\right),
  \label{eq:06-kf-update}
\end{equation}
\begin{equation}
  \bm{P}_k=\left(\bm{I}-\bm{K}_k\bm{H}_k\right)\bm{P}_k^{-}.
  \label{eq:06-kf-cov}
\end{equation}
预测把不确定性按动力学传播，更新用测量收缩不确定性；$\bm{K}$ 是「测量比先验可信多少」的加权。这些公式成立的前提是系统线性、噪声高斯——IMU 的姿态部分是矩阵乘法，不是向量加法，直接套 EKF 会在线性化处引入误差，于是有了 ESKF。

\subsection{误差状态的定义}

\textbf{ESKF 误差状态}（error-state Kalman filter）的核心思想：把真实状态写成「名义状态 $\oplus$ 小误差」。对 IMU 惯导取 15 维状态
\begin{equation}
  \bm{x} = \left(\bm{p},\ \bm{v},\ \bm{R},\ \bm{b}_g,\ \bm{b}_a\right),
  \qquad
  \delta\bm{x} = \left(\delta\bm{p},\ \delta\bm{v},\ \delta\bm{\theta},\ \delta\bm{b}_g,\ \delta\bm{b}_a\right)\in\mathbb{R}^{15},
  \label{eq:06-state}
\end{equation}
其中 $\bm{p},\bm{v}\in\mathbb{R}^{3}$ 是世界系下的位置与速度，$\bm{R}\in SO(3)$，$\bm{b}_g,\bm{b}_a\in\mathbb{R}^{3}$。真实状态与名义状态的关系为
\begin{equation}
  \bm{p}=\hat{\bm{p}}+\delta\bm{p},\quad
  \bm{v}=\hat{\bm{v}}+\delta\bm{v},\quad
  \bm{R}=\hat{\bm{R}}\exp\left(\delta\bm{\theta}^{\wedge}\right),\quad
  \bm{b}_g=\hat{\bm{b}}_g+\delta\bm{b}_g,\quad
  \bm{b}_a=\hat{\bm{b}}_a+\delta\bm{b}_a.
  \label{eq:06-error-def}
\end{equation}
注意姿态不能写成「$\hat{\bm{R}}+\delta\bm{\theta}$」——旋转矩阵相加后不再正交，必须写成名义旋转右乘一个小旋转。

\subsection{误差状态运动方程}

把式~\eqref{eq:06-error-def} 代入真实状态的动力学，减去名义状态的动力学，保留一阶项。

\textbf{位置误差。}由 $\dot{\bm{p}}=\bm{v}$ 与 $\dot{\hat{\bm{p}}}=\hat{\bm{v}}$ 直接相减得
\begin{equation}
  \delta\dot{\bm{p}} = \delta\bm{v}.
  \label{eq:06-dp}
\end{equation}
\textbf{速度误差。}真实加速度与名义加速度分别为
\begin{equation}
  \bm{a}_W = \bm{R}\left(\tilde{\bm{a}}-\bm{b}_a-\bm{n}_a\right)+\bm{g},
  \qquad
  \hat{\bm{a}}_W = \hat{\bm{R}}\left(\tilde{\bm{a}}-\hat{\bm{b}}_a\right)+\bm{g}.
\end{equation}
两式相减得 $\delta\dot{\bm{v}}=\bm{a}_W-\hat{\bm{a}}_W$。代入 $\bm{R}=\hat{\bm{R}}(\bm{I}+\delta\bm{\theta}^{\wedge})$ 与 $\bm{b}_a=\hat{\bm{b}}_a+\delta\bm{b}_a$：
\begin{align}
  \bm{a}_W-\hat{\bm{a}}_W
  &= \hat{\bm{R}}\left(\bm{I}+\delta\bm{\theta}^{\wedge}\right)
     \left(\tilde{\bm{a}}-\hat{\bm{b}}_a-\delta\bm{b}_a-\bm{n}_a\right)
     - \hat{\bm{R}}\left(\tilde{\bm{a}}-\hat{\bm{b}}_a\right) \notag\\
  &= \hat{\bm{R}}\Big[\left(\bm{I}+\delta\bm{\theta}^{\wedge}\right)\left(\tilde{\bm{a}}-\hat{\bm{b}}_a\right)
     -\left(\bm{I}+\delta\bm{\theta}^{\wedge}\right)\left(\delta\bm{b}_a+\bm{n}_a\right)
     -\left(\tilde{\bm{a}}-\hat{\bm{b}}_a\right)\Big].
\end{align}
第二项里 $\delta\bm{\theta}^{\wedge}(\delta\bm{b}_a+\bm{n}_a)$ 是两个小量相乘，属于二阶项，按一阶近似舍去；第一项与第三项相消后剩下 $\delta\bm{\theta}^{\wedge}(\tilde{\bm{a}}-\hat{\bm{b}}_a)$。于是
\begin{equation}
  \delta\dot{\bm{v}}
  = \hat{\bm{R}}\left[\delta\bm{\theta}^{\wedge}\left(\tilde{\bm{a}}-\hat{\bm{b}}_a\right)
    -\delta\bm{b}_a-\bm{n}_a\right]
  = -\hat{\bm{R}}\left(\tilde{\bm{a}}-\hat{\bm{b}}_a\right)^{\wedge}\delta\bm{\theta}
    -\hat{\bm{R}}\,\delta\bm{b}_a-\hat{\bm{R}}\,\bm{n}_a,
  \label{eq:06-dv}
\end{equation}
最后一步用了 $\delta\bm{\theta}^{\wedge}\bm{x}=-\bm{x}^{\wedge}\delta\bm{\theta}$。式~\eqref{eq:06-dv} 说明三件事：姿态误差会「漏」到速度里（耦合项 $-\hat{\bm{R}}\bm{a}_b^{\wedge}$，其中 $\bm{a}_b=\tilde{\bm{a}}-\hat{\bm{b}}_a$）；加速度计零偏误差直接进速度；噪声经当前姿态旋转到世界系。

\textbf{姿态误差。}对 $\bm{R}=\hat{\bm{R}}\exp(\delta\bm{\theta}^{\wedge})$ 求导，代入真实/名义动力学 $\dot{\bm{R}}=\bm{R}\bm{\omega}^{\wedge}$、$\dot{\hat{\bm{R}}}=\hat{\bm{R}}\hat{\bm{\omega}}^{\wedge}$（$\hat{\bm{\omega}}=\tilde{\bm{\omega}}-\hat{\bm{b}}_g$）：
\begin{align}
  \dot{\bm{R}} &= \dot{\hat{\bm{R}}}\exp(\delta\bm{\theta}^{\wedge})
                 +\hat{\bm{R}}\frac{\diff}{\diff t}\exp(\delta\bm{\theta}^{\wedge})
  \approx \hat{\bm{R}}\hat{\bm{\omega}}^{\wedge}\left(\bm{I}+\delta\bm{\theta}^{\wedge}\right)
  +\hat{\bm{R}}\delta\dot{\bm{\theta}}^{\wedge},
\end{align}
另一方面 $\bm{\omega}=\tilde{\bm{\omega}}-\bm{b}_g-\bm{n}_g=\hat{\bm{\omega}}-\delta\bm{b}_g-\bm{n}_g$，故
\begin{equation}
  \dot{\bm{R}} = \bm{R}\bm{\omega}^{\wedge}
  \approx \hat{\bm{R}}\left(\bm{I}+\delta\bm{\theta}^{\wedge}\right)
  \left(\hat{\bm{\omega}}^{\wedge}-\delta\bm{b}_g^{\wedge}-\bm{n}_g^{\wedge}\right).
\end{equation}
令两式右边相等，两边左乘 $\hat{\bm{R}}^{\top}$，只保留一阶项：
\begin{equation}
  \hat{\bm{\omega}}^{\wedge}+\hat{\bm{\omega}}^{\wedge}\delta\bm{\theta}^{\wedge}+\delta\dot{\bm{\theta}}^{\wedge}
  = \hat{\bm{\omega}}^{\wedge}-\delta\bm{b}_g^{\wedge}-\bm{n}_g^{\wedge}
    +\delta\bm{\theta}^{\wedge}\hat{\bm{\omega}}^{\wedge}.
\end{equation}
消去两边的 $\hat{\bm{\omega}}^{\wedge}$，移项得
\begin{equation}
  \delta\dot{\bm{\theta}}^{\wedge}
  = \delta\bm{\theta}^{\wedge}\hat{\bm{\omega}}^{\wedge}-\hat{\bm{\omega}}^{\wedge}\delta\bm{\theta}^{\wedge}
    -\delta\bm{b}_g^{\wedge}-\bm{n}_g^{\wedge}.
\end{equation}
用反对称矩阵恒等式 $\bm{a}^{\wedge}\bm{b}^{\wedge}-\bm{b}^{\wedge}\bm{a}^{\wedge}=(\bm{a}\times\bm{b})^{\wedge}$，第一项为 $(\delta\bm{\theta}\times\hat{\bm{\omega}})^{\wedge}$；而 $\delta\bm{\theta}\times\hat{\bm{\omega}}=-\hat{\bm{\omega}}\times\delta\bm{\theta}=-\hat{\bm{\omega}}^{\wedge}\delta\bm{\theta}$。脱掉 $\wedge$ 得
\begin{equation}
  \delta\dot{\bm{\theta}} = -\hat{\bm{\omega}}^{\wedge}\,\delta\bm{\theta}-\delta\bm{b}_g-\bm{n}_g.
  \label{eq:06-dtheta}
\end{equation}
\textbf{零偏误差。}零偏的随机游走对真实与名义状态相同，相减后噪声保留：
\begin{equation}
  \delta\dot{\bm{b}}_g = \bm{n}_{bg},
  \qquad
  \delta\dot{\bm{b}}_a = \bm{n}_{ba}.
  \label{eq:06-dbias}
\end{equation}

\subsection{离散化、状态转移矩阵与过程噪声}

把式~\eqref{eq:06-dp}、\eqref{eq:06-dv}、\eqref{eq:06-dtheta}、\eqref{eq:06-dbias} 写成 $\delta\dot{\bm{x}}=\bm{A}\,\delta\bm{x}+\bm{B}\bm{n}$，$\bm{A}$ 是 $15\times15$ 分块矩阵：
\begin{equation}
  \bm{A} = \begin{bmatrix}
    \bm{0} & \bm{I} & \bm{0} & \bm{0} & \bm{0}\\
    \bm{0} & \bm{0} & -\hat{\bm{R}}\bm{a}_b^{\wedge} & \bm{0} & -\hat{\bm{R}}\\
    \bm{0} & \bm{0} & -\hat{\bm{\omega}}^{\wedge} & -\bm{I} & \bm{0}\\
    \bm{0} & \bm{0} & \bm{0} & \bm{0} & \bm{0}\\
    \bm{0} & \bm{0} & \bm{0} & \bm{0} & \bm{0}
  \end{bmatrix},
  \qquad \bm{a}_b=\tilde{\bm{a}}-\hat{\bm{b}}_a.
  \label{eq:06-A}
\end{equation}
一阶离散化 $\bm{F}=\bm{I}+\bm{A}\Delta t$，逐块写出：
\begin{equation}
  \bm{F} = \begin{bmatrix}
    \bm{I} & \bm{I}\Delta t & \bm{0} & \bm{0} & \bm{0}\\
    \bm{0} & \bm{I} & -\hat{\bm{R}}\bm{a}_b^{\wedge}\Delta t & \bm{0} & -\hat{\bm{R}}\Delta t\\
    \bm{0} & \bm{0} & \bm{I}-\hat{\bm{\omega}}^{\wedge}\Delta t & -\bm{I}\Delta t & \bm{0}\\
    \bm{0} & \bm{0} & \bm{0} & \bm{I} & \bm{0}\\
    \bm{0} & \bm{0} & \bm{0} & \bm{0} & \bm{I}
  \end{bmatrix}.
  \label{eq:06-F}
\end{equation}
逐块解释：第一行表示位置误差由速度误差积分而来；第二行表示速度误差由「姿态误差耦合的加速度」与「加速度计零偏误差」驱动，这两块分别是 $\partial\dot{\bm{v}}/\partial\delta\bm{\theta}$ 与 $\partial\dot{\bm{v}}/\partial\delta\bm{b}_a$ 乘以 $\Delta t$；第三行表示姿态误差由「自身随角速度旋转」与「陀螺零偏误差」驱动；第四、五行表示零偏误差本身不传播，只有噪声驱动。过程噪声协方差（一阶离散）为
\begin{equation}
  \bm{Q}_d = \mathrm{diag}\left(\bm{0},\ \sigma_{a}^{2}\Delta t\,\bm{I},\ \sigma_{g}^{2}\Delta t\,\bm{I},\ \sigma_{bg}^{2}\Delta t\,\bm{I},\ \sigma_{ba}^{2}\Delta t\,\bm{I}\right),
  \label{eq:06-Qd}
\end{equation}
其中 $\sigma_a,\sigma_g$ 是加速度计、陀螺仪的连续噪声密度，$\sigma_{bg},\sigma_{ba}$ 为零偏随机游走密度；系数 $\Delta t$ 正是式~\eqref{eq:06-white-disc}、\eqref{eq:06-bias-disc} 的结论。位置块为零，因为位置误差的随机性完全来自速度误差的积分。

\subsection{观测方程、更新与重置}

\textbf{位置观测。}若外部（激光、视觉、轮速积分）给出世界系位置 $\bm{z}_p=\bm{p}+\bm{n}_p$，残差 $\bm{r}=\bm{z}_p-\hat{\bm{p}}$，则
\begin{equation}
  \bm{H}_p = \left[\bm{I}\ \ \bm{0}\ \ \bm{0}\ \ \bm{0}\ \ \bm{0}\right]\in\mathbb{R}^{3\times15}.
\end{equation}
\textbf{本体系速度观测。}轮式里程计给出机体系速度 $\bm{z}_v=\bm{R}_{WB}^{\top}\bm{v}_W+\bm{n}_v$。对姿态误差求雅可比：$\bm{R}^{\top}\approx(\bm{I}-\delta\bm{\theta}^{\wedge})\hat{\bm{R}}^{\top}$，故
\begin{equation}
  \bm{z}_v \approx \hat{\bm{R}}^{\top}\hat{\bm{v}}+\hat{\bm{R}}^{\top}\delta\bm{v}
  -\delta\bm{\theta}^{\wedge}\hat{\bm{R}}^{\top}\hat{\bm{v}}
  = \hat{\bm{R}}^{\top}\hat{\bm{v}}+\hat{\bm{R}}^{\top}\delta\bm{v}
  +\left(\hat{\bm{R}}^{\top}\hat{\bm{v}}\right)^{\wedge}\delta\bm{\theta},
\end{equation}
其中用了 $-\delta\bm{\theta}^{\wedge}\bm{w}=\bm{w}^{\wedge}\delta\bm{\theta}$。于是
\begin{equation}
  \bm{H}_v = \left[\bm{0}\ \ \hat{\bm{R}}^{\top}\ \ \left(\hat{\bm{R}}^{\top}\hat{\bm{v}}\right)^{\wedge}\ \ \bm{0}\ \ \bm{0}\right]\in\mathbb{R}^{3\times15}.
  \label{eq:06-Hv}
\end{equation}
这两项很有代表性：位置观测只「看」位置块；速度观测同时通过 $\hat{\bm{R}}^{\top}$ 与姿态误差耦合，所以它还能顺带修正姿态。残差 $\bm{r}=\bm{z}-\bm{h}(\hat{\bm{x}})$，按式~\eqref{eq:06-kf-update} 算 $\bm{K}$ 与 $\delta\hat{\bm{x}}$，再注入名义状态：
\begin{equation}
  \hat{\bm{p}}\leftarrow\hat{\bm{p}}+\delta\hat{\bm{p}},\quad
  \hat{\bm{v}}\leftarrow\hat{\bm{v}}+\delta\hat{\bm{v}},\quad
  \hat{\bm{R}}\leftarrow\hat{\bm{R}}\exp\left(\delta\hat{\bm{\theta}}^{\wedge}\right),\quad
  \hat{\bm{b}}_g\leftarrow\hat{\bm{b}}_g+\delta\hat{\bm{b}}_g,\quad
  \hat{\bm{b}}_a\leftarrow\hat{\bm{b}}_a+\delta\hat{\bm{b}}_a.
  \label{eq:06-inject}
\end{equation}
注入后误差状态被「用掉」，应重置为零。由于旋转的注入不是加法，协方差需要乘重置雅可比：
\begin{equation}
  \bm{P}\leftarrow\bm{G}\,\bm{P}\,\bm{G}^{\top},
  \qquad
  \bm{G}=\mathrm{diag}\left(\bm{I},\ \bm{I},\ \bm{I}-\tfrac{1}{2}\left[\delta\hat{\bm{\theta}}\right]^{\wedge},\ \bm{I},\ \bm{I}\right).
  \label{eq:06-reset}
\end{equation}
$\bm{G}$ 的姿态块来自「注入后新误差状态与旧误差状态之差」的一阶展开；$\delta\hat{\bm{\theta}}$ 很小时 $\bm{G}\approx\bm{I}$，许多实现直接省略这一步。完整推导见 \cite{sola2017eskf}，李群视角见 \cite{barfoot2017state}。

\begin{figure}[htbp]
  \centering
  \begin{tikzpicture}[node distance=6mm and 13mm]
    \node[block] (imu) {IMU 测量\\$\tilde{\bm{\omega}},\tilde{\bm{a}}$};
    \node[block, right=of imu] (nom) {名义状态\\积分预测};
    \node[block, right=of nom] (f) {误差协方差\\$\bm{F},\bm{Q}_d$ 传播};
    \node[block, below=of f] (upd) {卡尔曼更新\\$\bm{K}$、$\delta\hat{\bm{x}}$};
    \node[block, left=of upd] (inj) {误差注入\\与 $\bm{P}$ 重置};
    \node[block, left=of inj] (out) {输出位姿\\$T_{WB}$};
    \node[block2, above=of upd] (obs) {外部观测\\位置 / 速度};
    \draw[-stealth, thick] (imu) -- (nom);
    \draw[-stealth, thick] (nom) -- (f);
    \draw[-stealth, thick] (f) -- (upd);
    \draw[-stealth, thick] (obs) -- (upd);
    \draw[-stealth, thick] (upd) -- (inj);
    \draw[-stealth, thick] (inj) -- (out);
    \draw[-stealth, thick, dashed] (inj.north) |- (nom.south);
  \end{tikzpicture}
  \caption{ESKF 的单次迭代。IMU 驱动名义状态积分与误差状态协方差传播；有观测时做卡尔曼更新，把 $\delta\hat{\bm{x}}$ 注入名义状态并重置协方差；注入后的名义状态继续作为下一次积分的起点。}
  \label{fig:06-eskf}
\end{figure}

\begin{pseudo}{ESKF 单次迭代（IMU 预测 + 观测更新）}
  \item 读取 $\tilde{\bm{\omega}},\tilde{\bm{a}}$ 与时间间隔 $\Delta t$
  \item 名义状态预测：$\hat{\bm{a}}_W\leftarrow\hat{\bm{R}}(\tilde{\bm{a}}-\hat{\bm{b}}_a)+\bm{g}$
  \item \quad $\hat{\bm{p}}\leftarrow\hat{\bm{p}}+\hat{\bm{v}}\Delta t+\tfrac{1}{2}\hat{\bm{a}}_W\Delta t^{2}$，$\hat{\bm{v}}\leftarrow\hat{\bm{v}}+\hat{\bm{a}}_W\Delta t$
  \item \quad $\hat{\bm{R}}\leftarrow\hat{\bm{R}}\exp((\tilde{\bm{\omega}}-\hat{\bm{b}}_g)^{\wedge}\Delta t)$
  \item 协方差预测：由式~\eqref{eq:06-F}、\eqref{eq:06-Qd} 算 $\bm{P}\leftarrow\bm{F}\bm{P}\bm{F}^{\top}+\bm{Q}_d$
  \item 若本周期有观测 $\bm{z}$：$\bm{r}\leftarrow\bm{z}-\bm{h}(\hat{\bm{x}})$，$\bm{S}\leftarrow\bm{H}\bm{P}\bm{H}^{\top}+\bm{R}$
  \item \quad $\bm{K}\leftarrow\bm{P}\bm{H}^{\top}\bm{S}^{-1}$，$\delta\hat{\bm{x}}\leftarrow\bm{K}\bm{r}$
  \item \quad 按式~\eqref{eq:06-inject} 注入，按式~\eqref{eq:06-reset} 重置 $\bm{P}$
\end{pseudo}

\textbf{ESKF 与 EKF 的对比。}普通 EKF 直接在 15 维状态上线性化，姿态块要处理「旋转矩阵加向量」的非法运算，且线性化点离真值越远误差越大。ESKF 的线性化点始终在名义状态附近，误差状态每次更新后被重置、始终很小，因此线性化误差小；姿态更新用指数映射保持正交性；协方差描述小误差，数值条件更好。代价是多一套注入与重置的簿记。

\section{航迹推算与零速修正}\label{sec:06-strapdown}

\textbf{航迹推算}（strapdown，捷联惯导）就是按测量模型把加速度积分两次。由式~\eqref{eq:06-measure} 去掉噪声并用估计零偏：
\begin{equation}
  \hat{\bm{a}}_W = \hat{\bm{R}}\left(\tilde{\bm{a}}-\hat{\bm{b}}_a\right)+\bm{g},
  \label{eq:06-acc-world}
\end{equation}
然后欧拉积分
\begin{equation}
  \hat{\bm{v}}_{k+1}=\hat{\bm{v}}_k+\hat{\bm{a}}_{W,k}\Delta t,
  \qquad
  \hat{\bm{p}}_{k+1}=\hat{\bm{p}}_k+\hat{\bm{v}}_k\Delta t+\tfrac{1}{2}\hat{\bm{a}}_{W,k}\Delta t^{2}.
  \label{eq:06-integration}
\end{equation}
\textbf{为什么纯 IMU 位置会很快不可用？}设加速度计零偏残差为 $b$，则速度误差按 $bt$ 增长、位置误差按 $\tfrac{1}{2}bt^{2}$ 增长。取 $b=0.01\,\mathrm{m/s^2}$、$t=60\,\mathrm{s}$，得 $\Delta p\approx\tfrac{1}{2}\times0.01\times60^{2}=18\,\mathrm{m}$。更隐蔽的是姿态误差：一个 $0.01\degree$ 的倾角误差 $\epsilon=1.75\times10^{-4}\,\mathrm{rad}$ 会把重力「漏」进水平加速度，产生 $g_0\epsilon\approx1.7\times10^{-3}\,\mathrm{m/s^{2}}$ 的虚假加速度，60\,s 后位置误差约 $\tfrac{1}{2}g_0\epsilon t^{2}\approx3.1\,\mathrm{m}$。这就是为什么纯 IMU 必须配合外部观测使用，也是第~\ref{ch:09-fusion}~章融合动机的定量依据。

\textbf{零速修正}（zero velocity update, ZUPT）是最廉价的修正手段：机器人静止时真实速度为零，把「速度为零」当作一次观测即可约束速度与零偏。检测条件通常是「角速度模长小于阈值、加速度模长接近 $g_0$、且持续若干周期」。触发后构造伪观测 $\bm{z}_v=\bm{0}$：
\begin{equation}
  \bm{r} = \bm{0}-\hat{\bm{v}},\qquad
  \bm{H}=\left[\bm{0}\ \ \bm{I}\ \ \bm{0}\ \ \bm{0}\ \ \bm{0}\right],
\end{equation}
按式~\eqref{eq:06-kf-update} 更新。零速时速度误差被观测直接压小，同时通过 $\bm{F}$ 的耦合把零偏误差也约束住，这就是「ZUPT 能估零偏」的机理。工程上 $\bm{R}$ 不能过大（没效果）也不能过小（误检测时会把状态错误拉回）。

\begin{keybox}{ZUPT 的三条工程经验}
\begin{enumerate}
  \item 检测比修正更重要：先保证「确实静止」再修正，宁可漏检不要误检。
  \item 静止判据用滑动窗口（如连续 10 个周期都满足），单周期判据容易被撞击或振动骗过。
  \item ZUPT 只约束速度与零偏，不约束位置；位置漂移需要外部绝对观测来消除。
\end{enumerate}
\end{keybox}

\section{工程实现：ROS 2 中的 IMU}\label{sec:06-engineering}

ROS 2 的 \texttt{sensor\_msgs/msg/Imu} 包含 \texttt{orientation}（四元数）、\texttt{angular\_velocity}（rad/s）、\texttt{linear\_acceleration}（m/s$^2$，是比力而非世界系加速度）以及三组协方差。最常见的坑：

\begin{itemize}
  \item \textbf{单位与坐标系}：角速度必须转成 rad/s，IMU 轴系必须通过外参与全书约定（$x$ 前、$y$ 左、$z$ 上）统一，否则重力方向符号全错。
  \item \textbf{协方差}：填的是测量噪声而非零偏随机游走，混填会让滤波器过度自信。
  \item \textbf{振动}：高速运动时振动使加速度模长远超重力，静止检测会失效。
\end{itemize}

\textbf{IMU 预积分}（preintegration）是另一种使用方式：在两个关键帧 $i,j$ 之间，把 IMU 测量积分成与初始状态无关的相对运动量
\begin{equation}
  \Delta\bm{R}_{ij}=\prod_{k=i}^{j-1}\exp\left(\left(\tilde{\bm{\omega}}_k-\hat{\bm{b}}_g\right)^{\wedge}\Delta t\right),
  \quad
  \Delta\bm{v}_{ij}=\sum_{k=i}^{j-1}\Delta\bm{R}_{ik}\left(\tilde{\bm{a}}_k-\hat{\bm{b}}_a\right)\Delta t,
  \label{eq:06-preint}
\end{equation}
$\Delta\bm{p}_{ij}$ 则由 $\Delta\bm{v}_{ik}$ 与加速度再做一次积分得到。
从而在优化时用 $\bm{R}_j=\bm{R}_i\Delta\bm{R}_{ij}$、$\bm{v}_j=\bm{v}_i+\bm{g}\Delta t_{ij}+\bm{R}_i\Delta\bm{v}_{ij}$、$\bm{p}_j=\bm{p}_i+\bm{v}_i\Delta t_{ij}+\tfrac{1}{2}\bm{g}\Delta t_{ij}^{2}+\bm{R}_i\Delta\bm{p}_{ij}$ 恢复状态，不必重复积分原始 IMU 数据。零偏更新后只需对预积分量做一阶修正。其噪声传播与图优化中的用法见第~\ref{ch:13-backend}~章与 \cite{forster2017preintegration}。

\begin{summary}

\begin{enumerate}
  \item IMU 测量模型：$\tilde{\bm{\omega}}=\bm{\omega}+\bm{b}_g+\bm{n}_g$，$\tilde{\bm{a}}=\bm{R}^{\top}(\bm{a}_W-\bm{g})+\bm{b}_a+\bm{n}_a$；零偏按随机游走建模。
  \item 白噪声离散化 $\sigma_d=\sigma_c/\sqrt{\Delta t}$，零偏随机游走 $\sigma_{bd}=\sigma_{bc}\sqrt{\Delta t}$，两者都由积分方差推出。
  \item 四元数运动学 $\dot{\bm{q}}=\tfrac12\bm{q}\otimes[0,\bm{\omega}]$，离散更新 $\bm{q}_{k+1}=\bm{q}_k\otimes\exp(\tfrac12\bm{\omega}\Delta t)$，必须右乘并归一化。
  \item 静止时加速度计读数为 $-\bm{R}^{\top}\bm{g}$，可解出横滚与俯仰；偏航不可观测，因为绕重力轴旋转不改变读数。
  \item 互补滤波把低频的加速度计与高频的陀螺按 $H_{\text{LP}}+H_{\text{HP}}=1$ 合成；Mahony 用叉乘误差 $\bm{e}=\hat{\bm{g}}_b\times\hat{\bm{a}}$ 加 PI 反馈，误差只含垂直重力方向的分量。
  \item ESKF 以 $\bm{R}=\hat{\bm{R}}\exp(\delta\bm{\theta}^{\wedge})$ 定义误差状态，15 维 $\bm{F}$ 由式~\eqref{eq:06-F} 给出，每块都对应一条一阶误差传播；更新后要注入并重置。
  \item 纯 IMU 位置误差随 $t^{2}$ 增长：$b_a=0.01\,\mathrm{m/s^2}$ 时 60\,s 约 18\,m，$0.01\degree$ 倾角误差时约 3\,m。
  \item ZUPT 用「速度为零」的伪观测约束速度与零偏，检测的可靠性比修正本身更重要。
\end{enumerate}


\end{summary}

\begin{problemset}

  \item 由式~\eqref{eq:06-white-cont} 完整重推 $\sigma_d=\sigma_c/\sqrt{\Delta t}$ 与 $\sigma_{bd}=\sigma_{bc}\sqrt{\Delta t}$，并说明两者随 $\Delta t$ 的变化方向为何相反。
  \item 取 $\bm{q}=[\cos(\pi/8),\ 0,\ 0,\ \sin(\pi/8)]$，用式~\eqref{eq:06-q2R} 求 $\bm{R}$，并验证它是绕 $z$ 轴转 $\pi/4$。
  \item 用式~\eqref{eq:06-roll-pitch} 计算 $\tilde{\bm{a}}=[0.5,\ 0.3,\ 9.7]^{\top}$ 对应的横滚与俯仰，并解释为什么偏航无法从这组数得到。
  \item 写出 Mahony 误差 $\bm{e}$ 的一阶展开，证明 $\bm{e}=-(\bm{I}-\hat{\bm{g}}_b\hat{\bm{g}}_b^{\top})\delta\bm{\theta}$，并说明为什么偏航分量不进入误差。
  \item 推导式~\eqref{eq:06-F} 中「速度误差对姿态误差」那一块的来源，并说明为什么系数是 $\hat{\bm{R}}\bm{a}_b^{\wedge}$ 而不是 $\bm{a}_b^{\wedge}\hat{\bm{R}}$。
  \item 设加速度计零偏残差 $0.02\,\mathrm{m/s^2}$、陀螺零偏残差 $0.05\degree/\mathrm{s}$，估算 30\,s 后的位置误差量级，并判断哪个误差源主导。
\end{problemset}

\begin{reading}

\begin{itemize}
  \item Sola 的 ESKF 教程（\href{https://arxiv.org/abs/1711.02508}{arXiv:1711.02508}）：本章第 6 节几乎逐式对应。
  \item Barfoot《State Estimation for Robotics》（\href{https://books.google.com/books?id=jfZFDwAAQBAJ}{Cambridge, 2017}）：李群视角下的误差状态与协方差传播。
  \item Mahony 2008 与 Madgwick 2011 的开源实现（\href{https://x-io.co.uk/open-source-imu-and-ahrs-algorithms/}{x-io.co.uk}）：工程细节与参数整定。
\end{itemize}


\end{reading}

\chapter{激光里程计与扫描匹配}\label{ch:07-lidar-odom}

\chapterintro{激光雷达给的是「一圈点的三维坐标」，它不告诉你自己动了多少，却把周围几何形状以毫米级精度记录下来。于是激光里程计的核心问题变成：给定两帧点云，如何求出它们之间的相对位姿？这就是\textbf{扫描匹配}（scan matching）。本章从「最小化对应点距离」这一句话出发，把 ICP 的闭式解一步步推完，再讲点到面残差的 Gauss-Newton 求解、NDT 的高斯化建模、LOAM 的特征匹配框架，最后用信息矩阵的特征值解释「为什么激光里程计在长廊里会失效」。学完本章，你应当能看懂 FAST-LIO、LOAM、Cartographer 前端里那些矩阵是怎么来的。}

\section{为什么用 LiDAR 做里程计}\label{sec:07-why}

激光雷达的优势是几何精度高、不受光照影响、测距范围远：一个 16 线雷达在 10\,m 处的测距误差通常只有几厘米，而两帧点云之间的相对位姿可以由成百上千个点共同约束，理论上精度可以做到行程的 $0.1\,\%$--$1\,\%$。它不依赖纹理，所以在白墙、暗处、逆光等视觉失效的场景仍然可用。

缺点同样明显：点云稀疏且分布不均（近处密、远处疏），对几何结构敏感（长廊、大平面、空旷场地里约束不足），并且存在\textbf{运动畸变}——雷达转一圈需要 $100\,\mathrm{ms}$，这期间机器人已经动了，一帧点云的各个点其实是在不同位姿下测到的（去畸变的方法见第~\ref{ch:16-lio-evolution}~章）。此外还有数据量大、算力消耗高的问题。

激光里程计和 SLAM 的关系是：它只负责「相邻帧之间怎么动」，不建全局地图、不做回环。SLAM 篇（第~\ref{ch:11-slam-framework}~章）会把它当作前端的一部分。

\section{点云配准的问题定义}\label{sec:07-registration}

\begin{definition}[扫描匹配与点云配准]
给定源点云 $\mathcal{P}=\{\bm{p}_i\}_{i=1}^{N}$ 与目标点云 $\mathcal{Q}=\{\bm{q}_j\}_{j=1}^{M}$（都在同一坐标系下表达，例如雷达系），\textbf{点云配准}（point cloud registration）指求一个刚体变换 $T=(\bm{R},\bm{t})\in SE(3)$，使变换后的源点云与目标点云尽量重合；\textbf{扫描匹配}（scan matching）特指以激光扫描为输入的配准，它是激光里程计的核心步骤。
\end{definition}

最自然的度量是「对应点距离平方和」：假设已经知道每个 $\bm{p}_i$ 与哪个 $\bm{q}_i$ 对应，则
\begin{equation}
  \min_{\bm{R}\in SO(3),\ \bm{t}\in\mathbb{R}^{3}}
  \sum_{i=1}^{N}\norm{\bm{R}\bm{p}_i+\bm{t}-\bm{q}_i}^{2}.
  \label{eq:07-icp-problem}
\end{equation}
但「哪个点对应哪个点」本身也不知道，于是问题变成鸡生蛋：知道对应关系就能求变换，知道变换就能找对应关系。ICP（iterative closest point）的解法是交替迭代：用当前变换找最近邻作为对应，再用对应求变换。

\begin{definition}[最近邻与 KD 树]
对点 $\bm{p}$，它在目标点云中的\textbf{最近邻}（nearest neighbor）是 $\argmin_{\bm{q}\in\mathcal{Q}}\norm{\bm{p}-\bm{q}}$。若对每个点都暴力扫描全部 $M$ 个点，单次查找是 $O(M)$；\textbf{KD 树}（k-d tree）把点云按坐标轴递归二分建树，把单次最近邻查询降到平均 $O(\log M)$，是 ICP 能实时运行的关键数据结构。
\end{definition}

\section{ICP：从闭式解到迭代流程}\label{sec:07-icp}

\subsection{点到点 ICP 的闭式解}

先解决「对应关系已知」时的子问题。目标函数
\begin{equation}
  J(\bm{R},\bm{t}) = \sum_{i=1}^{N}\norm{\bm{R}\bm{p}_i+\bm{t}-\bm{q}_i}^{2}
\end{equation}
对 $\bm{t}$ 是二次的，可以直接求极值。对 $\bm{t}$ 求偏导（对向量求导，梯度为零向量）：
\begin{equation}
  \frac{\partial J}{\partial \bm{t}}
  = \sum_{i=1}^{N}2\left(\bm{R}\bm{p}_i+\bm{t}-\bm{q}_i\right)=\bm{0}
  \ \Longrightarrow\
  N\bm{t} = \sum_{i=1}^{N}\bm{q}_i - \bm{R}\sum_{i=1}^{N}\bm{p}_i.
\end{equation}
记质心 $\bar{\bm{p}}=\frac{1}{N}\sum_i\bm{p}_i$、$\bar{\bm{q}}=\frac{1}{N}\sum_i\bm{q}_i$，则
\begin{equation}
  \bm{t} = \bar{\bm{q}}-\bm{R}\bar{\bm{p}}.
  \label{eq:07-icp-t}
\end{equation}
把式~\eqref{eq:07-icp-t} 代回 $J$，并记去质心坐标 $\tilde{\bm{p}}_i=\bm{p}_i-\bar{\bm{p}}$、$\tilde{\bm{q}}_i=\bm{q}_i-\bar{\bm{q}}$：
\begin{align}
  \bm{R}\bm{p}_i+\bm{t}-\bm{q}_i
  &= \bm{R}\left(\tilde{\bm{p}}_i+\bar{\bm{p}}\right)+\bar{\bm{q}}-\bm{R}\bar{\bm{p}}-\left(\tilde{\bm{q}}_i+\bar{\bm{q}}\right)
  = \bm{R}\tilde{\bm{p}}_i-\tilde{\bm{q}}_i,
\end{align}
于是
\begin{equation}
  J = \sum_{i=1}^{N}\norm{\bm{R}\tilde{\bm{p}}_i-\tilde{\bm{q}}_i}^{2}
  = \sum_{i=1}^{N}\left(\tilde{\bm{p}}_i^{\top}\tilde{\bm{p}}_i+\tilde{\bm{q}}_i^{\top}\tilde{\bm{q}}_i
  -2\tilde{\bm{q}}_i^{\top}\bm{R}\tilde{\bm{p}}_i\right).
  \label{eq:07-icp-decoupled}
\end{equation}
前两项与 $\bm{R}$ 无关，因此最小化 $J$ 等价于最大化
\begin{equation}
  \sum_{i=1}^{N}\tilde{\bm{q}}_i^{\top}\bm{R}\tilde{\bm{p}}_i
  = \tr\left(\bm{R}\sum_{i=1}^{N}\tilde{\bm{p}}_i\tilde{\bm{q}}_i^{\top}\right)
  = \tr\left(\bm{R}\bm{W}\right),
  \label{eq:07-icp-trace}
\end{equation}
其中第一步用了 $\bm{a}^{\top}\bm{b}=\tr(\bm{b}\bm{a}^{\top})$，并定义了 $3\times3$ 矩阵
\begin{equation}
  \bm{W} = \sum_{i=1}^{N}\tilde{\bm{p}}_i\tilde{\bm{q}}_i^{\top}.
  \label{eq:07-icp-W}
\end{equation}
现在问题变成「在 $SO(3)$ 上最大化 $\tr(\bm{R}\bm{W})$」。对 $\bm{W}$ 做奇异值分解（SVD）：
\begin{equation}
  \bm{W} = \bm{U}\bm{\Sigma}\bm{V}^{\top},
  \qquad
  \bm{\Sigma}=\diag(\sigma_1,\sigma_2,\sigma_3),\quad \sigma_1\ge\sigma_2\ge\sigma_3\ge0.
\end{equation}
代入目标函数并利用迹的循环置换性质 $\tr(ABC)=\tr(BCA)$：
\begin{equation}
  \tr(\bm{R}\bm{W}) = \tr\left(\bm{R}\bm{U}\bm{\Sigma}\bm{V}^{\top}\right)
  = \tr\left(\bm{\Sigma}\,\bm{V}^{\top}\bm{R}\,\bm{U}\right).
\end{equation}
记 $\bm{A}=\bm{V}^{\top}\bm{R}\bm{U}$。由于 $\bm{V},\bm{R},\bm{U}$ 都是正交阵，$\bm{A}$ 也是正交阵，其列向量是单位向量，故对角元满足 $A_{ii}\le1$。于是
\begin{equation}
  \tr(\bm{R}\bm{W}) = \sum_{k=1}^{3}\sigma_k A_{kk}
  \le \sum_{k=1}^{3}\sigma_k\norm{A_{kk}}
  \le \sum_{k=1}^{3}\sigma_k.
\end{equation}
上界 $\sum_k\sigma_k$ 当且仅当 $A_{kk}=1$（$k=1,2,3$）时取到，此时 $\bm{A}=\bm{I}$，即
\begin{equation}
  \bm{V}^{\top}\bm{R}\bm{U}=\bm{I}
  \ \Longrightarrow\
  \bm{R}=\bm{V}\bm{U}^{\top}.
  \label{eq:07-icp-R}
\end{equation}
\textbf{反射情形。}式~\eqref{eq:07-icp-R} 只有在 $\det(\bm{V}\bm{U}^{\top})=+1$ 时才是合法的旋转；若 $\det(\bm{V}\bm{U}^{\top})=-1$，说明最优解是「旋转 + 镜像」，必须改成旋转。此时取
\begin{equation}
  \bm{R} = \bm{V}\,\diag(1,1,-1)\,\bm{U}^{\top},
  \label{eq:07-icp-Rfix}
\end{equation}
它等价于把 $\bm{V}$ 的最后一列取反，在保持正交的前提下使 $\det=+1$，并且是满足该约束时使 $\sum_k\sigma_k A_{kk}$ 最大的选择（因为取反只影响最小的 $\sigma_3$ 对应的那一项）。求出 $\bm{R}$ 后，由式~\eqref{eq:07-icp-t} 得 $\bm{t}=\bar{\bm{q}}-\bm{R}\bar{\bm{p}}$。

\begin{keybox}{ICP 闭式解的三步记忆法}
\begin{enumerate}
  \item 去质心：$\tilde{\bm{p}}_i=\bm{p}_i-\bar{\bm{p}}$，$\tilde{\bm{q}}_i=\bm{q}_i-\bar{\bm{q}}$，把平移分离出去；
  \item 攒矩阵 $\bm{W}=\sum\tilde{\bm{p}}_i\tilde{\bm{q}}_i^{\top}$，做 SVD $\bm{W}=\bm{U}\bm{\Sigma}\bm{V}^{\top}$；
  \item 取 $\bm{R}=\bm{V}\bm{U}^{\top}$（$\det<0$ 时用 $\bm{V}\mathrm{diag}(1,1,-1)\bm{U}^{\top}$），$\bm{t}=\bar{\bm{q}}-\bm{R}\bar{\bm{p}}$。
\end{enumerate}
\end{keybox}

\subsection{迭代流程、对应关系与收敛}

闭式解要求对应关系已知，而实际只知点云。ICP 用「最近邻 + 闭式解」交替逼近，单次迭代流程如下面的伪代码所示。

\begin{pseudo}{点到点 ICP 的单次迭代}
  \item 输入：源点云 $\mathcal{P}$、目标点云 $\mathcal{Q}$、当前变换 $T=(\bm{R},\bm{t})$
  \item 把源点云变换到目标系：$\bm{p}_i'\leftarrow\bm{R}\bm{p}_i+\bm{t}$
  \item 对每个 $\bm{p}_i'$，在 KD 树中查最近邻 $\bm{q}_{c(i)}$
  \item 剔除坏对应：距离 $\norm{\bm{p}_i'-\bm{q}_{c(i)}}>\tau_d$ 或法向夹角大的点对丢弃
  \item 用剩下的点对，按式~\eqref{eq:07-icp-W}、\eqref{eq:07-icp-R}、\eqref{eq:07-icp-t} 求新变换 $\Delta T$
  \item 更新 $T\leftarrow\Delta T\,T$（或右乘，取决于变换定义约定）
  \item 若残差相对变化小于阈值或达到最大迭代次数则停止，否则回到第 2 步
\end{pseudo}

几点工程经验：\textbf{初值}很重要——ICP 只有局部收敛性，两帧之间转得太快或位移太大时最近邻会连错，导致收敛到错误解，所以实际系统会用 IMU 或匀速模型预测初值；\textbf{外点剔除}必须有——动态物体、遮挡边缘、雷达噪点都会产生距离很大的错误对应，通常按距离阈值（如 $3\sigma$）或法向一致性剔除；\textbf{收敛判据}一般用相邻两次迭代的残差下降量，而不是只看迭代次数。

\textbf{ICP 的精度与局限。}点到点 ICP 的残差直接是欧氏距离，它隐含假设「场景中任意方向都有足够约束」。在走廊里，沿走廊方向平移几乎不改变点到点距离，这个方向就不可观（第~\ref{sec:07-degeneracy}~节详述）。点到面 ICP 通过引入法向信息缓解了这个问题，也是现代 LIO 系统的默认选择。

\begin{figure}[htbp]
  \centering
  \begin{tikzpicture}[scale=1.0]
    \fill[structurecolor!6] (0,0) rectangle (2.6,3.0);
    \fill[structurecolor!6] (4.2,0) rectangle (6.8,3.0);
    \node[font=\footnotesize] at (1.3,3.3) {源点云 $\mathcal{P}$（当前帧）};
    \node[font=\footnotesize] at (5.5,3.3) {目标点云 $\mathcal{Q}$（上一帧）};
    \foreach \x/\y in {0.6/2.4, 1.2/1.8, 0.7/1.0, 1.6/2.6, 1.8/0.6, 2.1/1.5} {
      \fill[structurecolor] (\x,\y) circle (1.6pt);
    }
    \foreach \x/\y in {4.7/2.5, 5.3/1.9, 4.8/1.1, 5.7/2.7, 5.9/0.7, 6.2/1.6} {
      \fill[red!60!black] (\x,\y) circle (1.6pt);
    }
    \draw[-stealth, thick, black!55] (0.6,2.4) -- (4.7,2.5);
    \draw[-stealth, thick, black!55] (1.2,1.8) -- (5.3,1.9);
    \draw[-stealth, thick, black!55] (0.7,1.0) -- (4.8,1.1);
    \draw[-stealth, thick, black!55] (1.6,2.6) -- (5.7,2.7);
    \draw[-stealth, thick, black!55] (1.8,0.6) -- (5.9,0.7);
    \draw[-stealth, thick, black!55] (2.1,1.5) -- (6.2,1.6);
    \node[font=\footnotesize] at (3.4,0.25) {最近邻对应关系};
  \end{tikzpicture}
  \caption{点到点 ICP 的一步。对当前帧每个点找到上一帧的最近点，用这些对应点对求闭式刚体变换；迭代若干次后对应关系与变换同时收敛。}
  \label{fig:07-icp}
\end{figure}

\subsection{点到面 ICP 与 Gauss-Newton 求解}

点到点残差对「沿平面方向的滑动」不敏感，点到面残差则直接把点压到平面上。设目标点云中与 $\bm{p}_i$ 对应的平面参数为 $(\bm{n}_i,\bm{q}_i)$，其中 $\bm{n}_i$ 是单位法向、$\bm{q}_i$ 是平面上一点。\textbf{点到面残差}（point-to-plane residual）定义为点 $\bm{p}_i$ 到该平面的有符号距离：
\begin{equation}
  r_i = \bm{n}_i^{\top}\left(\bm{R}\bm{p}_i+\bm{t}-\bm{q}_i\right).
  \label{eq:07-p2plane}
\end{equation}
它是非线性的（$\bm{R}$ 在旋转）。设当前估计附近的旋转增量为 $\delta\bm{\theta}$，用一阶近似 $\bm{R}\approx(\bm{I}+\delta\bm{\theta}^{\wedge})\bm{R}_0$；为简洁起见把 $\bm{R}_0\bm{p}_i$ 记作 $\bm{p}_i$（即先变换到当前估计的坐标系），并把 $\bm{t}$ 写成增量形式：
\begin{align}
  r_i &\approx \bm{n}_i^{\top}\left[\left(\bm{I}+\delta\bm{\theta}^{\wedge}\right)\bm{p}_i+\bm{t}-\bm{q}_i\right] \notag\\
  &= \bm{n}_i^{\top}\left(\bm{p}_i-\bm{q}_i\right)+\bm{n}_i^{\top}\bm{t}
     +\bm{n}_i^{\top}\delta\bm{\theta}^{\wedge}\bm{p}_i.
\end{align}
最后一项用 $\bm{n}^{\top}\delta\bm{\theta}^{\wedge}\bm{p}=\delta\bm{\theta}^{\top}\left(\bm{p}\times\bm{n}\right)$（由混合积 $\bm{n}\cdot(\delta\bm{\theta}\times\bm{p})=\delta\bm{\theta}\cdot(\bm{p}\times\bm{n})$ 得到），于是
\begin{equation}
  r_i \approx \underbrace{\bm{n}_i^{\top}\left(\bm{p}_i-\bm{q}_i\right)}_{r_i^{0}}
  + \left(\bm{p}_i\times\bm{n}_i\right)^{\top}\delta\bm{\theta}+\bm{n}_i^{\top}\bm{t}.
\end{equation}
把待求量写成六维 $\delta\bm{x}=[\delta\bm{\theta};\ \bm{t}]\in\mathbb{R}^{6}$，则
\begin{equation}
  r_i \approx r_i^{0} + \bm{A}_i\,\delta\bm{x},
  \qquad
  \bm{A}_i = \left[\left(\bm{p}_i\times\bm{n}_i\right)^{\top}\ \ \bm{n}_i^{\top}\right]\in\mathbb{R}^{1\times6}.
  \label{eq:07-p2plane-jac}
\end{equation}
最小化 $\sum_i w_i r_i^{2}$（$w_i$ 为权重，可选鲁棒核），令梯度为零：
\begin{equation}
  \frac{\partial}{\partial\delta\bm{x}}\sum_i w_i\left(r_i^{0}+\bm{A}_i\delta\bm{x}\right)^{2}
  = 2\sum_i w_i\bm{A}_i^{\top}\left(r_i^{0}+\bm{A}_i\delta\bm{x}\right)=\bm{0}.
\end{equation}
整理成\textbf{正规方程}（normal equation）：
\begin{equation}
  \left(\sum_i w_i\bm{A}_i^{\top}\bm{A}_i\right)\delta\bm{x}
  = -\sum_i w_i\bm{A}_i^{\top}r_i^{0}.
  \label{eq:07-gn}
\end{equation}
左边矩阵 $\bm{H}=\sum_i w_i\bm{A}_i^{\top}\bm{A}_i\in\mathbb{R}^{6\times6}$ 就是\textbf{信息矩阵}（information matrix）的近似，右边是梯度。解出 $\delta\bm{x}$ 后更新 $\bm{R}\leftarrow\exp(\delta\bm{\theta}^{\wedge})\bm{R}$、$\bm{t}\leftarrow\bm{t}+\delta\bm{t}$，再重新找对应、重新线性化——这就是 Gauss-Newton 版的 ICP。实践中常加 Levenberg-Marquardt 阻尼 $\bm{H}+\lambda\bm{I}$ 防止病态。

\begin{figure}[htbp]
  \centering
  \begin{tikzpicture}[scale=1.1]
    \draw[thick, structurecolor] (0,0.6) -- (5.2,-0.9);
    \node[font=\footnotesize, structurecolor] at (4.2,0.2) {目标平面};
    \draw[-stealth, thick, red!60!black] (2.6,-0.15) -- (2.6,1.35) node[above, font=\footnotesize] {$\bm{n}_i$};
    \fill[black] (1.2,1.9) circle (1.8pt) node[above right, font=\footnotesize] {$\bm{p}_i$};
    \fill[black] (2.15,-0.06) circle (1.8pt) node[below right, font=\footnotesize] {$\bm{q}_i$};
    \draw[dashed, black!50] (1.2,1.9) -- (2.15,-0.06);
    \node[font=\footnotesize] at (1.35,0.75) {$r_i$};
    \draw[dashed, black!30] (1.2,1.9) -- (2.6,0.24);
  \end{tikzpicture}
  \caption{点到面残差。点 $\bm{p}_i$ 到目标平面的有符号距离 $r_i=\bm{n}_i^{\top}(\bm{p}_i-\bm{q}_i)$；与点到点距离相比，它只惩罚「离开平面」的方向，对沿平面滑动不敏感，因此对走廊、墙面等结构收敛更好。}
  \label{fig:07-p2plane}
\end{figure}

\section{NDT：用高斯分布代替最近邻}\label{sec:07-ndt}

ICP 的每一步都要找最近邻，代价高且对初值敏感。NDT（normal distributions transform）换了个思路：先把目标点云切成体素，每个体素内的点拟合成一个高斯分布；配准时不再找「最近点」，而是看当前帧的点落在目标体素高斯分布里的「概率」有多大。

\begin{definition}[NDT 与体素高斯]
把目标点云按固定边长划分成体素，第 $j$ 个体素内若有 $N_j$ 个点，定义其均值与协方差
\begin{equation}
  \bm{\mu}_j=\frac{1}{N_j}\sum_{i=1}^{N_j}\bm{x}_i,
  \qquad
  \bm{\Sigma}_j=\frac{1}{N_j}\sum_{i=1}^{N_j}\left(\bm{x}_i-\bm{\mu}_j\right)\left(\bm{x}_i-\bm{\mu}_j\right)^{\top}.
  \label{eq:07-ndt-gauss}
\end{equation}
当 $N_j$ 太小时该体素不参与配准（或需要正则化 $\bm{\Sigma}_j$）。
\end{definition}

对当前帧的每个点 $\bm{p}_i$，用变换 $T$ 把它变到目标系得 $\bm{p}_i'=\bm{R}\bm{p}_i+\bm{t}$，找到它所在体素的高斯参数，定义该点的「匹配得分」
\begin{equation}
  s_i(T) = \exp\left(-\frac{1}{2}\left(\bm{p}_i'-\bm{\mu}_j\right)^{\top}\bm{\Sigma}_j^{-1}\left(\bm{p}_i'-\bm{\mu}_j\right)\right).
\end{equation}
整体目标是最小化负得分（等价于让点尽量落在高斯中心）：
\begin{equation}
  \min_{T}\ J(T) = -\sum_{i=1}^{N}\exp\left(-\frac{1}{2}\bm{d}_i^{\top}\bm{\Sigma}_j^{-1}\bm{d}_i\right),
  \qquad \bm{d}_i=\bm{p}_i'-\bm{\mu}_j.
  \label{eq:07-ndt-cost}
\end{equation}
\textbf{Newton 法求解。}记 $\bm{g}_i=\bm{\Sigma}_j^{-1}\bm{d}_i$，把 $s_i$ 对变换增量 $\delta\bm{x}=[\delta\bm{\theta};\ \delta\bm{t}]$ 求导。先算 $\bm{d}_i$ 的雅可比：$\bm{p}_i'=\bm{R}\bm{p}_i+\bm{t}\approx\bm{p}_i+\delta\bm{\theta}^{\wedge}\bm{p}_i+\bm{t}$，故
\begin{equation}
  \frac{\partial\bm{d}_i}{\partial\delta\bm{x}}
  = \left[-\bm{p}_i^{\wedge}\ \ \bm{I}\right]\in\mathbb{R}^{3\times6},
  \label{eq:07-ndt-jac}
\end{equation}
其中 $-\bm{p}_i^{\wedge}$ 来自 $\delta\bm{\theta}^{\wedge}\bm{p}_i=-\bm{p}_i^{\wedge}\delta\bm{\theta}$。对 $s_i=\exp(-\frac12\bm{d}_i^{\top}\bm{g}_i)$ 求导：
\begin{equation}
  \frac{\partial s_i}{\partial\delta\bm{x}}
  = -\frac{1}{2}s_i\,\frac{\partial\left(\bm{d}_i^{\top}\bm{\Sigma}_j^{-1}\bm{d}_i\right)}{\partial\delta\bm{x}}
  = -s_i\left(\frac{\partial\bm{d}_i}{\partial\delta\bm{x}}\right)^{\top}\bm{\Sigma}_j^{-1}\bm{d}_i,
\end{equation}
最后一步用了 $\partial(\bm{d}^{\top}\bm{\Sigma}^{-1}\bm{d})/\partial\bm{x}=2(\partial\bm{d}/\partial\bm{x})^{\top}\bm{\Sigma}^{-1}\bm{d}$，因子 2 与前面的 $1/2$ 抵消。于是梯度
\begin{equation}
  \bm{J}_{\text{NDT}} = \frac{\partial J}{\partial\delta\bm{x}}
  = \sum_{i=1}^{N}s_i\left(\frac{\partial\bm{d}_i}{\partial\delta\bm{x}}\right)^{\top}\bm{\Sigma}_j^{-1}\bm{d}_i.
\end{equation}
Hessian 用 Gauss-Newton 近似 $\bm{H}_{\text{NDT}}\approx\sum_i s_i(\partial\bm{d}_i/\partial\delta\bm{x})^{\top}\bm{\Sigma}_j^{-1}(\partial\bm{d}_i/\partial\delta\bm{x})$，然后迭代
\begin{equation}
  \delta\bm{x} = -\bm{H}_{\text{NDT}}^{-1}\bm{J}_{\text{NDT}}.
\end{equation}
\textbf{NDT 与 ICP 的对比：}NDT 不需要显式的最近邻一一对应，因此收敛域更大、对初值更宽容；代价是体素化与高斯拟合并不能表达精细结构（体素内的平面/边缘会被「糊」成高斯），精度通常略低于点到面 ICP。两者在现代系统中经常混用，例如 Cartographer 的 2D 前端用相关性搜索，而 FAST-LIO 系列用点到面 ICP。原始论文见 \cite{biber2003ndt}。

\section{相关性搜索与分支定界}\label{sec:07-csm}

2D 激光场景下还有一种简单粗暴但极其鲁棒的方法：\textbf{相关扫描匹配}（correlative scan matching, CSM）。它把位姿空间离散成栅格 $(x_k,y_k,\theta_k)$，对每个候选位姿计算「扫描点落在障碍栅格上的数量」作为得分，取最高分。由于是穷举搜索，它不依赖初值、不会陷入局部极小，但计算量是 $O(N_xN_yN_\theta)$。

工程上用\textbf{分支定界}（branch and bound）加速：把位姿空间组织成树，每个节点保存一个上界（该子树内所有位姿得分的上界），优先搜索上界高的节点，一旦某节点的上界低于当前最优解就整枝剪掉。这样绝大多数分支都不用展开。Cartographer 的 2D 前端正是这个思路（详细推导见第~\ref{ch:15-cartographer}~章与 \cite{hess2016cartographer}）。本章只需要记住：在 2D 场景、初值不可靠时，相关性搜索比 ICP 更稳。

\section{LOAM 框架：特征匹配的激光里程计}\label{sec:07-loam}

LOAM（lidar odometry and mapping）是 3D 激光里程计的里程碑式工作 \cite{zhang2014loam}，后续的 LeGO-LOAM \cite{shan2018lego}、LIO-SAM \cite{shan2020liosam}、KISS-ICP \cite{vizzo2023kissicp} 都受它影响。它的核心是「不再逐点找最近邻，而是先提取线、面特征，再让特征之间匹配」。

\textbf{第一步：特征提取。}对每个点 $\bm{p}_i$，取同一扫描线上前后各若干点组成集合 $\mathcal{S}$（如前后各 5 个点），定义曲率
\begin{equation}
  \kappa_i = \frac{1}{\abs{\mathcal{S}}}\sum_{j\in\mathcal{S}}\norm{\bm{p}_j-\bm{p}_i}^{2}.
  \label{eq:07-curvature}
\end{equation}
$\kappa_i$ 大说明该点邻域起伏剧烈，是\textbf{边缘点}（线特征）；$\kappa_i$ 小说明邻域平坦，是\textbf{平面点}（面特征）。按曲率排序后取最大的若干作为边缘点、最小的若干作为平面点，同时剔除可能属于同一物体边缘的邻近点，保证特征在空间中分散。

\textbf{第二步：帧间匹配。}对当前帧的每个边缘点 $\bm{p}_i$，在上一帧的边缘点中找最近的两个点 $\bm{p}_a,\bm{p}_b$（构成一条线），定义\textbf{点到线残差}为点到直线的距离：
\begin{equation}
  r_{\text{edge}} = \frac{\norm{\left(\bm{p}_i-\bm{p}_a\right)\times\left(\bm{p}_i-\bm{p}_b\right)}}{\norm{\bm{p}_a-\bm{p}_b}}.
\end{equation}
这个公式来自「三角形面积 = 底乘高 / 2 = 叉乘模长 / 2」：分子是平行四边形的面积，除以底边长度就是高，即点到直线的距离。对平面点 $\bm{p}_i$，在上一帧平面点中找最近的三个点 $\bm{p}_a,\bm{p}_b,\bm{p}_c$，用法向
\begin{equation}
  \bm{n} = \frac{\left(\bm{p}_a-\bm{p}_b\right)\times\left(\bm{p}_c-\bm{p}_b\right)}{\norm{\left(\bm{p}_a-\bm{p}_b\right)\times\left(\bm{p}_c-\bm{p}_b\right)}}
\end{equation}
定义\textbf{点到面残差} $r_{\text{plane}}=\bm{n}^{\top}(\bm{p}_i-\bm{p}_b)$。把两类残差按式~\eqref{eq:07-gn} 组装正规方程，用 LM 迭代求解帧间位姿。

\textbf{第三步：帧间 → 帧到地图。}LOAM 把里程计与建图分开：里程计以 10\,Hz 估计相邻帧位姿 $T_{k-1,k}$，再以更低频率把当前帧配准到已建立的地图（局部地图点云）上，得到全局一致的位姿 $T_{WB}$。位姿合成关系为
\begin{equation}
  T_{WB}(k) = T_{WB}(k-1)\,T_{k-1,k},
  \label{eq:07-pose-compose}
\end{equation}
而帧到地图的配准会对 $T_{WB}(k)$ 做一次修正。两阶段的好处是：高频部分只做局部配准（快），低频部分做全局配准（准）。下面的伪代码给出一次迭代的骨架。

\begin{pseudo}{LOAM 单帧处理流程}
  \item 输入：当前帧点云 $\mathcal{P}_k$、上一帧位姿 $T_{k-1}$、局部地图 $\mathcal{M}$
  \item 提取曲率 $\kappa_i$，选出边缘点集与平面点集
  \item 用匀速模型给初值 $T_k\leftarrow T_{k-1}T_{k-2}^{-1}T_{k-1}$
  \item 重复若干次：
  \item \quad 把 $\mathcal{P}_k$ 变换到上一帧系，为边缘点找最近线段、为平面点找最近平面
  \item \quad 组装点到线、点到面残差，解正规方程得增量 $\delta\bm{x}$
  \item \quad 更新 $T_k\leftarrow\exp(\delta\bm{x}^{\wedge})T_k$，剔除大残差点
  \item 把 $\mathcal{P}_k$ 按 $T_k$ 加入局部地图 $\mathcal{M}$，输出位姿 $T_k$
\end{pseudo}

\begin{figure}[htbp]
  \centering
  \begin{tikzpicture}[node distance=6mm and 10mm]
    \node[block] (raw) {原始点云\\$10\,\mathrm{Hz}$};
    \node[block, right=of raw] (feat) {特征提取\\曲率 $\kappa_i$};
    \node[block, right=of feat] (scan) {帧间匹配\\点到线 / 点到面};
    \node[block, below=of scan] (map) {帧到地图\\局部地图配准};
    \node[block, left=of map] (pose) {位姿输出\\$T_{WB}$};
    \draw[-stealth, thick] (raw) -- (feat);
    \draw[-stealth, thick] (feat) -- (scan);
    \draw[-stealth, thick] (scan) -- node[label, right] {粗位姿} (map);
    \draw[-stealth, thick] (map) -- (pose);
    \draw[-stealth, thick, dashed] (pose.north) to[out=90, in=270] (scan.south);
  \end{tikzpicture}
  \caption{LOAM 的数据流。高频的帧间匹配给出粗位姿，低频的帧到地图配准把它拉回全局一致；两阶段结构后来被 LeGO-LOAM、LIO-SAM 等继承。}
  \label{fig:07-loam}
\end{figure}

\section{退化分析：为什么长廊里会失效}\label{sec:07-degeneracy}

激光里程计的位姿解由正规方程~\eqref{eq:07-gn} 决定，其中 $\bm{H}=\sum_i w_i\bm{A}_i^{\top}\bm{A}_i$ 是 $6\times6$ 信息矩阵。$\bm{H}$ 的特征值刻画了「各个方向上有多少约束」：特征值大说明该方向被点云牢牢约束，特征值小说明该方向几乎无约束。

\begin{definition}[退化]
若信息矩阵 $\bm{H}$ 的某个特征值接近零（或远小于最大特征值），则对应的特征向量方向称为\textbf{退化方向}（degenerate direction）：在这个方向上，位姿解对残差几乎不敏感，估计结果由噪声决定而不是由数据决定。
\end{definition}

\textbf{长廊的例子。}一条无限长的走廊，两侧是平行墙面。沿走廊方向平移时，每个点到墙面的距离完全不变，点到面残差不变，因此该方向对应的 $\bm{A}_i$ 列全为零，$\bm{H}$ 的最小特征值趋于零。类似地，空旷场地里围绕竖直轴的旋转、大平面场景里垂直于平面的平移，都可能退化。

\textbf{检测方法。}对 $\bm{H}$ 做特征分解 $\bm{H}=\sum_k\lambda_k\bm{v}_k\bm{v}_k^{\top}$（特征值与特征向量定义见第~\ref{ch:02-math}~章），两种常用判据：
\begin{itemize}
  \item 最小特征值阈值：$\lambda_{\min}<\lambda_{\text{th}}$ 判为退化；
  \item 条件数：$\lambda_{\max}/\lambda_{\min}>\kappa_{\text{th}}$ 判为病态。
\end{itemize}
退化方向就是 $\lambda_{\min}$ 对应的特征向量 $\bm{v}_{\min}$；工程上可以把解在这个方向上「冻结」或「用先验填充」：把 $\bm{H}$ 在该方向上的元素替换成一个很大的值（等价于加入极强的先验），使解不再依赖无约束的数据方向。

\begin{keybox}{退化处理的三种手段}
\begin{enumerate}
  \item \textbf{换传感器}：加入 IMU 提供姿态与加速度先验，这是 LIO 紧耦合（第~\ref{ch:16-lio-evolution}~章）的核心动机。
  \item \textbf{换残差}：点到面比点到点约束更强，加入法向、平面、地面约束可以补充缺失方向。
  \item \textbf{换判据}：实时监控 $\lambda_{\min}$ 与条件数，一旦判为退化就降低该方向权重、拒绝输出不可信位姿，而不是硬算。
\end{enumerate}
\end{keybox}

\begin{summary}

\begin{enumerate}
  \item 激光里程计 = 扫描匹配：给定两帧点云求 $T$ 使对应点距离最小。
  \item 点到点 ICP 的闭式解：去质心后最大化 $\tr(\bm{R}\bm{W})$，$\bm{W}=\sum\tilde{\bm{p}}_i\tilde{\bm{q}}_i^{\top}$，SVD 得 $\bm{R}=\bm{V}\bm{U}^{\top}$（$\det<0$ 时修正），$\bm{t}=\bar{\bm{q}}-\bm{R}\bar{\bm{p}}$。
  \item ICP 迭代 = 找最近邻 + 闭式解，初值、外点剔除、收敛判据是三个工程要点；KD 树把最近邻查询降到 $O(\log M)$。
  \item 点到面残差线性化后得到 $\bm{A}_i=[(\bm{p}_i\times\bm{n}_i)^{\top},\bm{n}_i^{\top}]$，正规方程 $\bm{H}\delta\bm{x}=-\bm{b}$ 就是 Gauss-Newton 的一步。
  \item NDT 用体素高斯代替最近邻，得分是负指数函数，Newton 法迭代；收敛域更大但精度略低。
  \item 相关性搜索穷举位姿空间，鲁棒但慢；分支定界把它降到实用，是 Cartographer 的核心。
  \item LOAM 用曲率提取边缘点与平面点，点到线、点到面残差两阶段（帧间 + 帧到地图）求解，位姿按 $T_{WB}(k)=T_{WB}(k-1)T_{k-1,k}$ 合成。
  \item 退化由信息矩阵 $\bm{H}$ 的小特征值检测，长廊、空旷、大平面是典型场景；对策是换传感器、换残差、换判据。
\end{enumerate}


\end{summary}

\begin{problemset}

  \item 给两组对应点：$\mathcal{P}=\{(0,0),(1,0),(0,1)\}$、$\mathcal{Q}=\{(0,1),(1,1),(0,2)\}$，手工算出 $\bar{\bm{p}}$、$\bar{\bm{q}}$，验证这是纯平移 $[0,1]^{\top}$，并用式~\eqref{eq:07-icp-t} 反推 $\bm{t}$。
  \item 推导式~\eqref{eq:07-icp-t}：从 $\partial J/\partial\bm{t}=0$ 出发，写出每一步，并解释为什么平移可以「先分离出去」。
  \item 设 $\bm{W}=\bm{U}\diag(2,1,0.1)\bm{V}^{\top}$ 且 $\det(\bm{V}\bm{U}^{\top})=-1$，写出正确的 $\bm{R}$ 并说明为什么不能直接用 $\bm{V}\bm{U}^{\top}$。
  \item 用式~\eqref{eq:07-p2plane-jac} 写出一个点的 $\bm{A}_i$，并说明为什么沿平面方向的位移在残差中不出现。
  \item 解释 NDT 得分函数式~\eqref{eq:07-ndt-cost} 中协方差逆矩阵的作用：如果某个体素内点云近似一条直线，$\bm{\Sigma}_j^{-1}$ 会把哪个方向的残差放大？
  \item 走廊场景中，若只用点到面残差，$\bm{H}$ 的哪个特征值会接近零？请给出特征向量的物理含义，并设计一个检测方案。
\end{problemset}

\begin{reading}

\begin{itemize}
  \item Besl \& McKay 1992 的 ICP 原始论文（\href{https://ieeexplore.ieee.org/document/121791}{IEEE TPAMI}）：闭式解与收敛性分析。
  \item Biber \& Strasser 2003 的 NDT 论文：体素高斯与 2D 激光配准。
  \item Zhang \& Singh 2014 的 LOAM 论文（\href{https://www.ri.cmu.edu/pub_files/2014/7/Ji_LidarMapping_RSS2014_v8.pdf}{RSS}）：特征提取与两阶段框架，建议对照本章的伪代码与图~\ref{fig:07-loam} 阅读。
\end{itemize}


\end{reading}

\chapter{视觉里程计与 VIO}\label{ch:08-visual-odom}

\chapterintro{相机每秒送来几十张图像，每张图里有成千上万个「特征」，本质上是把三维世界压扁成二维。视觉里程计（VO）要做的，就是从这一串二维投影里反推出相机的三维运动。本章按「相机怎么成像 → 特征怎么匹配 → 运动怎么求解」的顺序展开：先把针孔模型与畸变讲透，再推对极几何与三角化，接着用重投影误差把位姿优化写出来，然后介绍直接法与 VIO 的三条主流路线（MSCKF、VINS-Mono、ORB-SLAM3），最后说清单目尺度不可观这一根本局限。学完本章，你应当能看懂 ORB-SLAM 与 VINS 前端在算什么。}

\section{相机模型与畸变}\label{sec:08-camera}

\subsection{针孔模型}

最简单的成像模型是\textbf{针孔模型}（pinhole model）：光线从物体出发，穿过光心，落在成像平面上。设相机系 $C$ 的原点在光心，$z$ 轴朝前（光轴方向），$x$ 轴向右、$y$ 轴向下（与图像坐标一致）。空间点 ${}^{C}\bm{X}=[X,Y,Z]^{\top}$ 投影到成像平面上的坐标为 $(X/Z,Y/Z)$（用相似三角形：成像平面距光心 $f$，则像高与物高的比等于像距与物距的比）。

像素坐标还要经过缩放与平移。设焦距以像素为单位为 $f_x,f_y$（$f_x=f/s_x$，$s_x$ 是单个像素的物理宽度），主点（光轴与成像平面的交点）像素坐标为 $(c_x,c_y)$，则
\begin{equation}
  u = f_x\frac{X}{Z}+c_x,
  \qquad
  v = f_y\frac{Y}{Z}+c_y.
  \label{eq:08-projection}
\end{equation}
写成齐次坐标：
\begin{equation}
  Z\begin{bmatrix}u\\v\\1\end{bmatrix}
  = \underbrace{\begin{bmatrix}
    f_x & 0 & c_x\\
    0 & f_y & c_y\\
    0 & 0 & 1
  \end{bmatrix}}_{\bm{K}}
  \begin{bmatrix}X\\Y\\Z\end{bmatrix}.
  \label{eq:08-K}
\end{equation}
矩阵 $\bm{K}$ 叫\textbf{内参矩阵}（intrinsic matrix），它只与相机本身有关；$f_x,f_y,c_x,c_y$ 叫\textbf{内参}。记号 $\pi(\cdot)$ 表示把三维点投影到像素平面，$\pi^{-1}$ 表示「给定深度」时的反投影。

\begin{example}[一组典型的 USB 相机内参]
分辨率 $640\times480$，$f_x=f_y=500$，$c_x=320$、$c_y=240$。空间点 $[1,0.5,5]^{\top}\,\mathrm{m}$ 的投影像素为
\begin{equation}
  u=500\times\frac{1}{5}+320=420,
  \qquad
  v=500\times\frac{0.5}{5}+240=290,
\end{equation}
即位于图像右半边、中心偏下。若该点深度从 $5\,\mathrm{m}$ 变到 $10\,\mathrm{m}$，像素 $u$ 会从 420 移到 $500/10+320=370$：深度翻倍，像素向主点靠近一半。
\end{example}

\subsection{径向与切向畸变}

真实镜头不是理想针孔，直线在图像里会变成曲线。工程上采用如下畸变模型：先把像素坐标转成归一化坐标 $x=(u-c_x)/f_x$、$y=(v-c_y)/f_y$，令 $r^{2}=x^{2}+y^{2}$，
\begin{align}
  x_{\text{dist}} &= x\left(1+k_1r^{2}+k_2r^{4}+k_3r^{6}\right)
  + 2p_1xy+p_2\left(r^{2}+2x^{2}\right), \label{eq:08-dist-x}\\
  y_{\text{dist}} &= y\left(1+k_1r^{2}+k_2r^{4}+k_3r^{6}\right)
  + p_1\left(r^{2}+2y^{2}\right)+2p_2xy, \label{eq:08-dist-y}
\end{align}
其中 $k_1,k_2,k_3$ 是\textbf{径向畸变}系数（$k_1<0$ 表现为桶形畸变），$p_1,p_2$ 是\textbf{切向畸变}系数（由镜头与成像面不平行引起）。最后再乘回内参得到畸变后的像素坐标。

\textbf{去畸变。}标定得到的是「畸变后像素 → 归一化坐标」的映射，而后续几何计算需要「无畸变像素」。由于式~\eqref{eq:08-dist-x} 不可解析求逆，工程上用迭代法：以畸变点 $x_d$ 为初值，反复执行
\begin{equation}
  x \leftarrow \frac{x_d-2p_1xy-p_2\left(r^{2}+2x^{2}\right)}{1+k_1r^{2}+k_2r^{4}+k_3r^{6}},
  \label{eq:08-undistort}
\end{equation}
（$y$ 同理）直到 $x$ 的变化小于阈值。多数库（OpenCV \texttt{undistort}）的做法是预先算好整幅图像的映射表，运行时查表，速度更快。

\begin{keybox}{相机模型的三个易错点}
\begin{enumerate}
  \item 图像 $v$ 轴向下：$c_y$ 与 $y$ 方向不要搞反，否则三角化出的三维点会上下颠倒。
  \item 内参来自标定：$f_x$ 与 $f_y$ 可以不相等（非方形像素），不要强行假设相等。
  \item 畸变系数作用于\textbf{归一化坐标}：先除内参再去畸变，顺序不能反。
\end{enumerate}
\end{keybox}

\section{特征点法视觉里程计}\label{sec:08-feature}

VO 的两大流派是特征点法与直接法。本节先讲特征点法：先在图像里找「好认的点」，再在两张图之间匹配它们，最后用匹配点求运动。

\begin{definition}[特征点与描述子匹配]
\textbf{特征点}（feature point / keypoint）是图像中可重复检测、对光照和小视角变化稳定的位置，如角点。\textbf{描述子}（descriptor）是描述该点邻域外观的向量，用于判断两个特征点是否是同一物理点。\textbf{描述子匹配}（descriptor matching）指按描述子距离为两帧特征点建立对应关系，常用欧氏距离（SIFT）或汉明距离（ORB 的二进制描述子）。
\end{definition}

\textbf{检测与描述。}FAST 特征检测角点：以像素 $p$ 为中心、半径 3 的圆上取 16 个像素，若其中连续一段（如 9 个）都明显亮于或暗于 $p$，则 $p$ 是角点。ORB 在 FAST 基础上加了方向（用灰度质心法算主方向）与旋转不变的 BRIEF 二进制描述子（256 bit），匹配时用汉明距离（两串二进制不同位的个数），可以用位运算快速计算。匹配后常做\textbf{比值检验}（ratio test）：对每个特征点，若最近与次近描述子距离之比小于 0.7 才接受匹配，这能剔除大量误匹配。

\textbf{运动求解的两条路。}匹配点若都是二维像素，用对极几何求相对位姿（2D-2D）；若其中一帧有深度（RGB-D 或三角化过），用 PnP 求位姿（3D-2D）；若两帧都有三维点，用 ICP（第~\ref{ch:07-lidar-odom}~章）。下面依次推导。

\subsection{对极几何：2D-2D 的约束}

设同一空间点在两帧相机系下的坐标为 $\bm{X}_1,\bm{X}_2$，两帧间的位姿变换为 $\bm{X}_2=\bm{R}\bm{X}_1+\bm{t}$，归一化像点为 $\bm{x}_1=\bm{X}_1/Z_1$、$\bm{x}_2=\bm{X}_2/Z_2$（即 $Z[\bm{x}^{\top},1]^{\top}=\bm{X}$）。关键观察：$\bm{x}_2$、$\bm{R}\bm{x}_1$、$\bm{t}$ 三个向量共面，因为 $\bm{X}_2=Z_1\bm{R}\bm{x}_1+\bm{t}$，而 $\bm{x}_2$ 与 $\bm{X}_2$ 平行。三个共面向量的混合积为零：
\begin{equation}
  \bm{x}_2^{\top}\left(\bm{t}\times\bm{R}\bm{x}_1\right)=0
  \quad\Longleftrightarrow\quad
  \bm{x}_2^{\top}\bm{t}^{\wedge}\bm{R}\,\bm{x}_1=0,
  \label{eq:08-epipolar}
\end{equation}
其中 $\bm{t}^{\wedge}$ 是反对称矩阵，满足 $\bm{t}^{\wedge}\bm{x}=\bm{t}\times\bm{x}$。式~\eqref{eq:08-epipolar} 就是\textbf{对极约束}（epipolar constraint），它说明「一对匹配点」就提供了关于 $\bm{R},\bm{t}$ 的一个标量方程。定义
\begin{equation}
  \bm{E}=\bm{t}^{\wedge}\bm{R},\qquad
  \bm{F}=\bm{K}_2^{-\top}\bm{E}\,\bm{K}_1^{-1},
  \label{eq:08-E-F}
\end{equation}
其中 $\bm{E}$ 叫\textbf{本质矩阵}（essential matrix），$\bm{F}$ 叫\textbf{基础矩阵}（fundamental matrix，作用在像素坐标上）。$\bm{E}$ 有 5 个自由度（$\bm{R}$ 三个、$\bm{t}$ 方向两个，尺度不可观），$\bm{F}$ 有 7 个自由度；作为 $3\times3$ 矩阵它们有 9 个元素，去掉整体尺度剩 8 个，再减去秩 2 约束正好是 7 个。

\textbf{八点法。}把 $\bm{F}=[f_1,\dots,f_9]^{\top}$ 拉直，每对匹配点给一个线性方程
\begin{equation}
  \left[u_2u_1,\ u_2v_1,\ u_2,\ v_2u_1,\ v_2v_1,\ v_2,\ u_1,\ v_1,\ 1\right]\bm{f}=0.
\end{equation}
$N\ge8$ 对点堆成 $\bm{A}\bm{f}=\bm{0}$，取 $\bm{A}$ 的最小奇异值对应的右奇异向量作为 $\bm{f}$。由于噪声，这样得到的 $\bm{F}$ 通常是满秩的，需要\textbf{强制秩 2}：对 $\bm{F}$ 做 SVD $\bm{F}=\bm{U}\diag(\sigma_1,\sigma_2,\sigma_3)\bm{V}^{\top}$，把最小的 $\sigma_3$ 置零，得 $\hat{\bm{F}}=\bm{U}\diag(\sigma_1,\sigma_2,0)\bm{V}^{\top}$。再从 $\hat{\bm{E}}=\bm{K}_2^{\top}\hat{\bm{F}}\bm{K}_1$ 分解出 $\bm{R},\bm{t}$（SVD 分解有 4 组解，用「点必须在相机前方」筛选）。工程上通常先用 RANSAC 剔除误匹配，再解 $\bm{E}$。

\subsection{三角化}

已知 $\bm{R},\bm{t}$ 与匹配点 $\bm{x}_1,\bm{x}_2$，要恢复该点的三维坐标。由 $\bm{X}_2=\bm{R}\bm{X}_1+\bm{t}$ 与 $\bm{X}_i=Z_i[\bm{x}_i^{\top},1]^{\top}$：
\begin{equation}
  Z_2\begin{bmatrix}\bm{x}_2\\1\end{bmatrix}
  = Z_1\bm{R}\begin{bmatrix}\bm{x}_1\\1\end{bmatrix}+\bm{t}.
  \label{eq:08-triangulate}
\end{equation}
两边左乘 $\bm{x}_2^{\wedge}$ 消掉左边的未知量 $Z_2$：
\begin{equation}
  \bm{0} = Z_1\,\bm{x}_2^{\wedge}\bm{R}\begin{bmatrix}\bm{x}_1\\1\end{bmatrix}
  +\bm{x}_2^{\wedge}\bm{t}
  \ \Longrightarrow\
  Z_1\,\bm{x}_2^{\wedge}\bm{R}\bar{\bm{x}}_1 = -\bm{x}_2^{\wedge}\bm{t},
  \label{eq:08-tri-ls}
\end{equation}
其中 $\bar{\bm{x}}_i=[\bm{x}_i^{\top},1]^{\top}$。这是关于 $Z_1$ 的超定线性方程（三个方程，秩 2）。用最小二乘解：两边左乘 $(\bm{x}_2^{\wedge}\bm{R}\bar{\bm{x}}_1)^{\top}$，得
\begin{equation}
  Z_1 = -\frac{\left(\bm{x}_2^{\wedge}\bm{R}\bar{\bm{x}}_1\right)^{\top}\left(\bm{x}_2^{\wedge}\bm{t}\right)}
  {\norm{\bm{x}_2^{\wedge}\bm{R}\bar{\bm{x}}_1}^{2}}.
  \label{eq:08-Z1}
\end{equation}
代回得 $\bm{X}_1=Z_1\bar{\bm{x}}_1$，再用 $\bm{X}_2=\bm{R}\bm{X}_1+\bm{t}$ 得到另一帧坐标。注意三角化的精度与视差有关：两帧间基线太短（视差小）时，深度对像素噪声极其敏感，这是单目 VO 在慢速直线运动时深度估计很差的根本原因。

\subsection{PnP 与重投影误差}

当一帧有深度信息时，问题变成 3D-2D：已知三维点 $\bm{X}_j$ 与像素观测 $\bm{u}_{ij}$，求相机位姿 $T_i$。最直接的做法是最小化\textbf{重投影误差}（reprojection error）：
\begin{equation}
  \bm{e}_{ij} = \bm{u}_{ij}-\pi\left(\bm{R}_i\bm{X}_j+\bm{t}_i\right),
  \qquad
  \min_{T_i}\sum_j\norm{\bm{e}_{ij}}^{2}.
  \label{eq:08-reproj}
\end{equation}
它是 $\bm{R}$ 的非线性函数，用 Gauss-Newton 迭代。先给出投影对三维点的雅可比。由式~\eqref{eq:08-projection}，
\begin{equation}
  \frac{\partial\pi}{\partial\bm{p}}
  = \begin{bmatrix}
    \dfrac{f_x}{Z} & 0 & -\dfrac{f_xX}{Z^{2}}\\[2mm]
    0 & \dfrac{f_y}{Z} & -\dfrac{f_yY}{Z^{2}}
  \end{bmatrix}\in\mathbb{R}^{2\times3},
  \qquad \bm{p}=\begin{bmatrix}X\\Y\\Z\end{bmatrix}.
  \label{eq:08-dpi-dp}
\end{equation}
再给出三维点对位姿增量 $\delta\bm{\xi}=[\delta\bm{\theta};\delta\bm{t}]$ 的雅可比。$\bm{p}=\bm{R}\bm{X}+\bm{t}$，对 $\delta\bm{\theta}$ 用 $\bm{R}\leftarrow\exp(\delta\bm{\theta}^{\wedge})\bm{R}\approx(\bm{I}+\delta\bm{\theta}^{\wedge})\bm{R}$：
\begin{equation}
  \frac{\partial\bm{p}}{\partial\delta\bm{\theta}}=-\bm{p}^{\wedge}\in\mathbb{R}^{3\times3},
  \qquad
  \frac{\partial\bm{p}}{\partial\delta\bm{t}}=\bm{I}\in\mathbb{R}^{3\times3}.
\end{equation}
链式法则得 $2\times6$ 的雅可比
\begin{equation}
  \bm{J}_{ij}=\frac{\partial\bm{e}_{ij}}{\partial\delta\bm{\xi}}
  = -\frac{\partial\pi}{\partial\bm{p}}\Big|_{\bm{p}=\bm{R}_i\bm{X}_j+\bm{t}_i}
  \left[-\bm{p}^{\wedge}\ \ \bm{I}\right]\in\mathbb{R}^{2\times6}.
  \label{eq:08-reproj-jac}
\end{equation}
正规方程为 $\left(\sum_j\bm{J}_{ij}^{\top}\bm{J}_{ij}\right)\delta\bm{\xi}=-\sum_j\bm{J}_{ij}^{\top}\bm{e}_{ij}$，与第~\ref{ch:07-lidar-odom}~章的点到面完全同构。若相机内参也未知，可以同时对 $\bm{K}$ 求导做自标定，工程上一般先离线标定好。

\textbf{PnP 的初值与外点。}重投影误差是局部方法，初值差会发散。常用的初始化方法有 DLT（直接线性变换）与 EPnP：DLT 把投影方程写成
\begin{equation}
  u\,\bm{p}_3^{\top}\bm{X}= \bm{p}_1^{\top}\bm{X},
  \qquad
  v\,\bm{p}_3^{\top}\bm{X}= \bm{p}_2^{\top}\bm{X},
\end{equation}
其中 $\bm{P}=\bm{K}[\bm{R}\ \bm{t}]\in\mathbb{R}^{3\times4}$ 有 12 个未知数，$N\ge6$ 个点即可线性求解（SVD 取零空间）；EPnP 则把三维点表示成 4 个虚拟控制点的加权和，把未知量降到 12 维以内，更快更稳。外点用 RANSAC：随机取最小点集求 PnP，统计内点数，迭代取内点最多的模型，最后用全部内点做一次重投影误差优化。

\subsection{光束法平差}

把多帧、多个点的重投影误差放在一起联合优化，就叫\textbf{光束法平差}（bundle adjustment, BA）：
\begin{equation}
  \min_{\{T_i\},\{\bm{X}_j\}}\sum_{(i,j)\in\mathcal{O}}\norm{\bm{u}_{ij}-\pi\left(\bm{R}_i\bm{X}_j+\bm{t}_i\right)}^{2},
  \label{eq:08-BA}
\end{equation}
其中 $\mathcal{O}$ 是所有观测。BA 的变量包含位姿与三维点，规模很大，但雅可比矩阵高度稀疏（每个观测只涉及一个位姿与一个点），可以用舒尔补（Schur complement）把点变量消掉，只解位姿（推导见第~\ref{ch:13-backend}~章）。视觉 SLAM 的后端本质上就是「带鲁棒核与边缘化的增量 BA」。

\begin{pseudo}{特征点法 VO 的单帧流程}
  \item 输入：当前帧图像 $I_k$、上一帧状态
  \item 提取 FAST 角点，计算 ORB 描述子
  \item 与上一帧描述子做匹配，比值检验剔除误匹配
  \item RANSAC 估计本质矩阵 $\bm{E}$（或基础矩阵 $\bm{F}$），得到初始 $\bm{R},\bm{t}$
  \item 三角化匹配点得到三维坐标；若上一帧已有三维点则改做 PnP
  \item 以重投影误差做 Gauss-Newton 局部优化，剔除大残差点
  \item 输出位姿 $T_{k}$，把新三角化的点加入地图
\end{pseudo}

\begin{figure}[htbp]
  \centering
  \begin{tikzpicture}[scale=1.0]
    \draw[thick, fill=structurecolor!8] (1.0,0.2) rectangle (1.2,2.4);
    \node[font=\footnotesize, rotate=90] at (1.1,1.95) {成像平面};
    \fill[black] (0,0) circle (1.8pt) node[below left, font=\footnotesize] {光心 $O$};
    \fill[black] (3.6,2.6) circle (1.8pt) node[right, font=\footnotesize] {$\bm{X}$};
    \draw[thick] (0,0) -- (3.6,2.6);
    \fill[red!60!black] (1.1,0.794) circle (1.6pt) node[right, font=\footnotesize] {$\bm{x}$};
    \draw[dashed, black!40] (0,0) -- (0,2.2) node[above, font=\footnotesize] {$f$};
    \draw[-stealth, thick] (0,0) -- (0.9,0) node[right, font=\footnotesize] {$x$};
    \draw[-stealth, thick] (0,0) -- (0,0.9) node[above, font=\footnotesize] {$y$};
    \draw[-stealth, thick] (0,0) -- (-0.7,-0.7) node[below left, font=\footnotesize] {$z$};
  \end{tikzpicture}
  \caption{针孔成像模型。三维点 $\bm{X}$ 与光心连线交成像平面于像点 $\bm{x}$，像素坐标由 $u=f_xX/Z+c_x$ 给出；焦距 $f$、深度 $Z$ 与像高构成相似三角形。}
  \label{fig:08-pinhole}
\end{figure}

\section{直接法与光流}\label{sec:08-direct}

特征点法要先检测、匹配，再丢弃大量信息（一幅图只用几百个点）。直接法（direct method）反过来：直接假设「同一个空间点在两帧中的像素灰度不变」，最小化\textbf{光度误差}（photometric error）：
\begin{equation}
  e_{i} = I_1\left(\bm{u}_i\right)-I_2\left(\pi\left(\bm{R}\,\pi^{-1}\left(\bm{u}_i,d_i\right)+\bm{t}\right)\right),
  \label{eq:08-photometric}
\end{equation}
其中 $d_i$ 是像素 $\bm{u}_i$ 对应的深度。目标是最小化 $\sum_i e_i^{2}$，用 Gauss-Newton 对位姿（以及可能的深度）迭代，雅可比通过图像梯度链式得到：
\begin{equation}
  \frac{\partial e_i}{\partial\delta\bm{\xi}}
  = -\nabla I_2^{\top}\frac{\partial\pi}{\partial\bm{p}}\left[-\bm{p}^{\wedge}\ \ \bm{I}\right].
\end{equation}
直接法不需要特征，弱纹理区域只要有小梯度就能用，还能利用全部像素；代价是光度误差对光照变化、相机自动曝光、卷帘快门非常敏感，且需要较好的深度初值。DSO \cite{engel2018dso} 是直接法的代表作，它联合优化光度参数与几何参数，在无纹理、高动态范围场景表现突出。

\textbf{光流}可以看作直接法的「前身」：为每个像素在下一帧里找到灰度最相似的位置（LK 光流），再把这些像素位移当作特征匹配的替代品。光流匹配的位移通常很小，所以适合视频这种帧间运动小的场景。

\section{VIO：视觉与惯性的融合}\label{sec:08-vio}

\begin{definition}[VIO]
\textbf{视觉惯性里程计}（visual-inertial odometry, VIO）指把相机与 IMU 的数据联合起来估计位姿与速度：相机提供长期无漂移的几何约束（但频率低、有尺度模糊），IMU 提供高频运动与尺度信息（但漂移快）。二者互补，是现代无人机、机器人的主流定位方案。
\end{definition}

\textbf{尺度问题：单目为什么不能定尺度。}由式~\eqref{eq:08-K}，投影只依赖 $X/Z$ 的比值。若把相机系下所有三维点与平移同时放大 $s$ 倍：$\bm{X}\to s\bm{X}$、$\bm{t}\to s\bm{t}$，则
\begin{equation}
  \frac{s\bm{X}}{sZ}=\frac{\bm{X}}{Z},
\end{equation}
像素投影完全不变。因此单目视觉只能恢复「相差一个未知尺度」的轨迹，速度、位置、地图尺度都不可观。IMU 的加速度计提供绝对加速度（真实物理量），能在运动激励足够时把尺度观测出来，这是 VIO 相对纯单目 VO 的核心优势。

\textbf{MSCKF。}MSCKF（multi-state constraint Kalman filter）\cite{mourikis2007msckf} 用 EKF 维护一个滑动窗口内的多个相机位姿，特征点不进入状态。当某个特征点离开视野时，它的多次观测给窗口内的位姿提供了一组约束；为了不把特征点坐标加进状态，用\textbf{零空间投影}把这些约束投影到「与特征点无关」的子空间，得到一个只含位姿的残差，再做 EKF 更新。好处是状态维度小、计算量可控；代价是信息有损（投影后约束变弱）。

\textbf{VINS-Mono。}VINS-Mono \cite{qin2018vinsmono} 走优化路线：IMU 预积分（第~\ref{ch:06-imu}~章）连接相邻关键帧，视觉观测提供重投影残差，二者放进同一个滑动窗口做 BA；窗口滑出时用边缘化（marginalization）把旧状态的信息转成先验，保留约束、丢弃变量。它是「视觉惯性紧耦合 + 滑动窗口优化」的经典实现。

\textbf{ORB-SLAM3。}ORB-SLAM3 \cite{campos2021orbslam3} 把特征点、多地图（ATLAS）、视觉惯性紧耦合整合在一起，支持单目/双目/RGB-D 与 IMU 的组合，是工程完整度最高的视觉 SLAM 系统之一。它的前端用 ORB 特征与对极几何/PnP，后端用局部 BA + 回环 + 位姿图。

\begin{figure}[htbp]
  \centering
  \begin{tikzpicture}[node distance=6mm and 11mm]
    \node[block] (cam) {相机\\30\,Hz};
    \node[block, right=of cam] (vis) {视觉前端\\特征 + 匹配};
    \node[block, below=of cam] (imu) {IMU\\200\,Hz};
    \node[block, right=of imu] (pre) {预积分\\$\Delta\bm{R},\Delta\bm{v},\Delta\bm{p}$};
    \node[block, right=of vis, yshift=-15mm] (win) {滑动窗口\\联合优化};
    \node[block, right=of win] (out) {位姿、速度\\零偏、地图};
    \draw[-stealth, thick] (cam) -- (vis);
    \draw[-stealth, thick] (imu) -- (pre);
    \draw[-stealth, thick] (vis) -- node[label, above] {重投影残差} (win);
    \draw[-stealth, thick] (pre) -- node[label, below] {IMU 残差} (win);
    \draw[-stealth, thick] (win) -- (out);
    \draw[-stealth, thick, dashed] (out.south) |- (pre.south);
  \end{tikzpicture}
  \caption{VIO 的信息流。视觉前端给出重投影残差，IMU 预积分给出帧间运动约束，二者在滑动窗口里联合优化；优化结果反过来用于下一轮预积分与边缘化。}
  \label{fig:08-vio}
\end{figure}

\section{视觉里程计的局限}\label{sec:08-limits}

\begin{itemize}
  \item \textbf{尺度}：单目尺度不可观（上面已证明）；双目/IMU 可以补上，但需要基线或运动激励。
  \item \textbf{光照}：强光、逆光、暗处、曝光变化会让特征/光度误差失效；直接法比特征点法更敏感。
  \item \textbf{纹理}：白墙、地板、大面积单色物体上没有可用的角点或梯度，VO 会退化（与激光的长廊问题本质相同）。
  \item \textbf{动态物体}：行人、其他机器人会引入错误的特征匹配，需要 RANSAC、语义分割或运动一致性检验。
  \item \textbf{滚动快门}：高速旋转时 CMOS 逐行曝光会引入几何畸变，需要专门建模。
  \item \textbf{尺度漂移}：即使有 IMU，长时间慢速运动时尺度仍可能缓慢漂移，需要绝对观测（GPS、地图）来闭合。
\end{itemize}

这也是为什么真实的定位系统几乎都是「视觉 + IMU + 轮速 + 激光」多传感器融合：没有任何单一传感器能在所有场景下可靠工作（第~\ref{ch:09-fusion}~章）。

\begin{keybox}{选择 VO 方案的三条原则}
\begin{enumerate}
  \item 有纹理、要精度：特征点法 + 回环（ORB-SLAM3 一类）；
  \item 弱纹理、要速度：直接法（DSO 一类）或光流；
  \item 要尺度、要鲁棒：一定要上 IMU（VIO），纯单目 VO 不适合作为长期定位的唯一来源。
\end{enumerate}
\end{keybox}

\begin{summary}

\begin{enumerate}
  \item 针孔模型：$Z[u,v,1]^{\top}=\bm{K}[X,Y,Z]^{\top}$，内参含 $f_x,f_y,c_x,c_y$；畸变由 $k_1,k_2,k_3,p_1,p_2$ 描述，作用在归一化坐标上。
  \item 特征点法：FAST 检测、ORB 描述子、汉明距离匹配与比值检验，是 VO 前端的基础。
  \item 对极约束 $\bm{x}_2^{\top}\bm{t}^{\wedge}\bm{R}\bm{x}_1=0$ 来自三个向量共面；八点法 + 秩 2 约束求 $\bm{F}$，再分解出 $\bm{R},\bm{t}$。
  \item 三角化把匹配点变成三维点，深度精度强烈依赖视差；视差小则深度不可靠。
  \item PnP 最小化重投影误差，雅可比由 $\partial\pi/\partial\bm{p}$ 与 $[-\bm{p}^{\wedge},\bm{I}]$ 链式相乘得到；初值可用 DLT/EPnP，外点用 RANSAC。
  \item BA 是联合优化多帧位姿与三维点，雅可比稀疏，可用舒尔补消元。
  \item 直接法最小化光度误差，适应弱纹理但怕光照变化；DSO 是代表。
  \item 单目尺度不可观（投影对整体缩放不变）；VIO 用 IMU 补尺度与高频运动，MSCKF、VINS-Mono、ORB-SLAM3 是三条代表性路线。
\end{enumerate}


\end{summary}

\begin{problemset}

  \item 用 $f_x=f_y=500$、$c_x=c_y=320$ 的内参，求空间点 $[2,-1,4]^{\top}$ 的像素坐标；再求 $[4,-2,8]^{\top}$ 的像素坐标，验证两者相同（说明尺度不可观）。
  \item 从 $\bm{X}_2=\bm{R}\bm{X}_1+\bm{t}$ 与 $\bm{x}_i=\bm{X}_i/Z_i$ 出发，完整推导对极约束~\eqref{eq:08-epipolar}，并说明每一步用了什么性质。
  \item 写出八点法的线性方程组，说明为什么最少需要 8 对点，以及为什么解出的 $\bm{F}$ 需要强制秩 2。
  \item 用式~\eqref{eq:08-dpi-dp} 与 $[-\bm{p}^{\wedge},\bm{I}]$ 推导 $\bm{J}_{ij}$，并写出 $\bm{J}_{ij}$ 的形状与每一块的含义。
  \item 若两帧间基线为 $0.05\,\mathrm{m}$、目标深度 $5\,\mathrm{m}$，像素噪声 $0.5$ 像素、$f=500$，估算三角化深度的相对误差量级。
  \item 比较特征点法与直接法的优缺点，并说明为什么 VIO 能解决单目的尺度问题。
\end{problemset}

\begin{reading}

\begin{itemize}
  \item DSO 论文（\href{https://ieeexplore.ieee.org/document/8220081}{IEEE TPAMI 2018}）：直接法的光度误差与模型参数联合优化。
  \item VINS-Mono 论文与开源代码（\href{https://github.com/HKUST-Aerial-Robotics/VINS-Mono}{GitHub}）：预积分 + 滑动窗口 + 边缘化的完整实现。
  \item ORB-SLAM3 论文与代码（\href{https://github.com/UZ-SLAMLab/ORB_SLAM3}{GitHub}）：多地图与视觉惯性紧耦合的工程范本。
\end{itemize}


\end{reading}

\chapter{多传感器融合定位}\label{ch:09-fusion}

\chapterintro{前面五章分别讲了轮式、IMU、激光、视觉四种里程计，每一种都有自己的死穴：轮式怕打滑、IMU 漂移快、激光怕退化、视觉怕光照。真实系统从来不会只押注一个传感器，而是把它们的信息按各自的可信度加权组合——这就是\textbf{多传感器融合}。本章先讲清融合的层次与松紧耦合的区别，再把卡尔曼融合、序贯更新与批量更新的等价性推完，接着定义信息矩阵与马氏距离并用卡方检验做外点剔除，最后讲时间同步这个「看起来是工程细节、实际上决定成败」的问题。学完本章，你应当能设计一个多传感器融合的架构，并解释为什么 FAST-LIO 和 VINS 要那么写。}

\section{融合的层次与松紧耦合}\label{sec:09-levels}

\subsection{三个融合层次}

按信息被处理的深度，融合分三个层次：

\begin{itemize}
  \item \textbf{数据级融合}（data-level fusion）：直接融合原始测量。例如把激光点云与相机像素做深度对齐，或把 IMU 与轮速放进同一个滤波器。信息损失最小，但对时间同步与标定要求最高。
  \item \textbf{特征级融合}（feature-level fusion）：各传感器先提取特征（点云特征、图像特征、轮速积分位姿），再融合特征。多数 LIO/VIO 系统属于这一层。
  \item \textbf{决策级融合}（decision-level fusion）：各传感器独立算出位姿，再对结果做加权平均或投票。实现最简单，但丢失了传感器间的相关性信息。
\end{itemize}

\subsection{松耦合与紧耦合}

\begin{definition}[松耦合与紧耦合]
\textbf{松耦合}（loosely coupled）：每个传感器子系统先独立估计出一个位姿/速度，再把估计结果作为「观测」送进融合滤波器。\textbf{紧耦合}（tightly coupled）：不使用子系统的中间位姿结果，而是把原始测量（原始点云残差、原始像素残差、IMU 预积分）直接放进同一个滤波器或优化器，联合估计状态。
\end{definition}

\begin{figure}[htbp]
  \centering
  \begin{tikzpicture}[node distance=5mm and 9mm]
    \node[block] (l1) {LiDAR 里程计};
    \node[block, below=of l1] (v1) {视觉里程计};
    \node[block, below=of v1] (i1) {IMU 积分};
    \node[block, right=26mm of v1] (f1) {融合滤波器};
    \node[block, right=of f1] (o1) {位姿};
    \draw[-stealth, thick] (l1) -| node[label, above, pos=0.25] {位姿 + 协方差} (f1.north);
    \draw[-stealth, thick] (v1) -- (f1);
    \draw[-stealth, thick] (i1) -| (f1.south);
    \draw[-stealth, thick] (f1) -- (o1);
    \node[font=\footnotesize, structurecolor] at (2.6,-2.6) {(a) 松耦合};
    \begin{scope}[xshift=6.4cm]
      \node[block] (l2) {LiDAR 点云};
      \node[block, below=of l2] (v2) {图像};
      \node[block, below=of v2] (i2) {IMU};
      \node[block, right=26mm of v2] (f2) {联合优化\\（原始残差）};
      \node[block, right=of f2] (o2) {位姿};
      \draw[-stealth, thick] (l2) -| (f2.north);
      \draw[-stealth, thick] (v2) -- (f2);
      \draw[-stealth, thick] (i2) -| (f2.south);
      \draw[-stealth, thick] (f2) -- (o2);
      \node[font=\footnotesize, structurecolor] at (2.6,-2.6) {(b) 紧耦合};
    \end{scope}
  \end{tikzpicture}
  \caption{松耦合与紧耦合。松耦合先把各传感器压成「位姿 + 协方差」再融合，实现简单但会丢失相关性；紧耦合把原始测量直接送入联合估计，信息保留最完整，代价是系统复杂度与标定要求更高。}
  \label{fig:09-coupling}
\end{figure}

为什么紧耦合通常更好？松耦合把子系统的估计当作独立观测，但子系统的误差其实与原数据相关（例如激光里程计位姿误差与 IMU 预积分误差共享同一段运动），当成独立会低估不确定性、导致融合权重失真；更重要的是，松耦合在子系统失效时（如激光退化）无法利用其他传感器的原始信息补救。紧耦合则可以在同一个优化里用 IMU 先验补上激光缺失的方向，这正是 FAST-LIO 在长廊里仍能工作的原因（第~\ref{ch:16-lio-evolution}~章）。

\section{卡尔曼融合的统一形式}\label{sec:09-kf}

\subsection{状态增广与观测堆叠}

多传感器融合最直接的实现是\textbf{状态增广}：把所有需要估计的量放进同一个状态向量 $\bm{x}$，每个传感器的测量都写成 $\bm{z}_i=\bm{H}_i\bm{x}+\bm{v}_i$。这样问题就退化成「一个状态、多组观测」的卡尔曼滤波。假设先验 $\bm{x}\sim\mathcal{N}(\hat{\bm{x}}_0,\bm{P}_0)$，两组独立观测 $\bm{z}_1=\bm{H}_1\bm{x}+\bm{v}_1$、$\bm{z}_2=\bm{H}_2\bm{x}+\bm{v}_2$，噪声 $\bm{v}_i\sim\mathcal{N}(\bm{0},\bm{R}_i)$。批量融合的结果是
\begin{equation}
  \hat{\bm{x}} = \left(\bm{P}_0^{-1}+\bm{H}_1^{\top}\bm{R}_1^{-1}\bm{H}_1+\bm{H}_2^{\top}\bm{R}_2^{-1}\bm{H}_2\right)^{-1}
  \left(\bm{P}_0^{-1}\hat{\bm{x}}_0+\bm{H}_1^{\top}\bm{R}_1^{-1}\bm{z}_1+\bm{H}_2^{\top}\bm{R}_2^{-1}\bm{z}_2\right).
  \label{eq:09-batch}
\end{equation}

\subsection{序贯更新与批量更新的等价性}

实现上更常见的是序贯更新：先更新 $\bm{z}_1$，把结果作为先验再更新 $\bm{z}_2$。下面证明两者等价。

\textbf{第一步：更新 $\bm{z}_1$。}令 $\bm{\Lambda}_0=\bm{P}_0^{-1}$、$\bm{b}_0=\bm{P}_0^{-1}\hat{\bm{x}}_0$。卡尔曼更新的信息形式为
\begin{equation}
  \bm{\Lambda}_1 = \bm{\Lambda}_0+\bm{H}_1^{\top}\bm{R}_1^{-1}\bm{H}_1,
  \qquad
  \bm{b}_1 = \bm{b}_0+\bm{H}_1^{\top}\bm{R}_1^{-1}\bm{z}_1,
  \qquad
  \hat{\bm{x}}_1=\bm{\Lambda}_1^{-1}\bm{b}_1.
  \label{eq:09-seq1}
\end{equation}
\textbf{第二步：以 $\hat{\bm{x}}_1$ 为先验更新 $\bm{z}_2$。}
\begin{equation}
  \bm{\Lambda}_2 = \bm{\Lambda}_1+\bm{H}_2^{\top}\bm{R}_2^{-1}\bm{H}_2,
  \qquad
  \bm{b}_2 = \bm{b}_1+\bm{H}_2^{\top}\bm{R}_2^{-1}\bm{z}_2,
  \qquad
  \hat{\bm{x}}_2=\bm{\Lambda}_2^{-1}\bm{b}_2.
\end{equation}
把式~\eqref{eq:09-seq1} 代入：
\begin{equation}
  \bm{\Lambda}_2 = \bm{P}_0^{-1}+\bm{H}_1^{\top}\bm{R}_1^{-1}\bm{H}_1+\bm{H}_2^{\top}\bm{R}_2^{-1}\bm{H}_2,
\end{equation}
\begin{equation}
  \bm{b}_2 = \bm{P}_0^{-1}\hat{\bm{x}}_0+\bm{H}_1^{\top}\bm{R}_1^{-1}\bm{z}_1+\bm{H}_2^{\top}\bm{R}_2^{-1}\bm{z}_2.
\end{equation}
这与式~\eqref{eq:09-batch} 完全相同。结论：\textbf{信息形式下，多观测融合就是信息矩阵与信息向量的相加}：
\begin{equation}
  \bm{\Lambda} = \bm{P}_0^{-1}+\sum_i\bm{H}_i^{\top}\bm{R}_i^{-1}\bm{H}_i,
  \qquad
  \bm{b} = \bm{P}_0^{-1}\hat{\bm{x}}_0+\sum_i\bm{H}_i^{\top}\bm{R}_i^{-1}\bm{z}_i.
  \label{eq:09-information-sum}
\end{equation}
这个结论带来两个实用推论：第一，观测的\textbf{顺序不影响结果}（序贯更新可以任意排列）；第二，每个观测对信息矩阵的贡献是 $\bm{H}_i^{\top}\bm{R}_i^{-1}\bm{H}_i$，与状态维度无关——这就是为什么紧耦合里可以把上千个点云残差直接累加，而不必把它们变成「一个位姿观测」。

\begin{keybox}{松耦合/紧耦合的取舍}
\begin{enumerate}
  \item 松耦合适合「传感器模块已经很成熟、接口稳定」的场合，调试简单、故障隔离好；
  \item 紧耦合适合「单传感器会失效、需要互相补位」的场合，精度更高但对标定、时间同步、状态维数要求更高；
  \item 工程折中：关键子系统紧耦合（如 IMU + 激光），外围模块松耦合（如全局定位）。
\end{enumerate}
\end{keybox}

\section{信息矩阵、马氏距离与卡方检验}\label{sec:09-mahalanobis}

\subsection{信息矩阵}

\begin{definition}[信息矩阵]
协方差矩阵 $\bm{P}$ 的逆 $\bm{\Lambda}=\bm{P}^{-1}$ 叫\textbf{信息矩阵}（information matrix）。$\bm{P}$ 描述不确定性的大小，$\bm{\Lambda}$ 描述「信息量」的大小：方差越小、信息越多、$\Lambda$ 越大。在线性高斯系统中，融合后的信息矩阵等于各观测信息矩阵之和（式~\eqref{eq:09-information-sum}）。
\end{definition}

\subsection{马氏距离及其卡方分布}

考虑一个观测残差 $\bm{r}=\bm{z}-\bm{H}\hat{\bm{x}}$。不同残差分量量纲不同，不能直接平方相加，需要按各自的不确定性归一化。设噪声与状态误差独立，真实状态为 $\bm{x}$，先验误差 $\bm{e}=\hat{\bm{x}}-\bm{x}\sim\mathcal{N}(\bm{0},\bm{P})$，则
\begin{equation}
  \bm{r} = \bm{z}-\bm{H}\hat{\bm{x}}
  = \left(\bm{H}\bm{x}+\bm{v}\right)-\bm{H}\left(\bm{x}+\bm{e}\right)
  = \bm{v}-\bm{H}\bm{e}.
  \label{eq:09-residual}
\end{equation}
它的协方差为（用了 $\bm{v}$ 与 $\bm{e}$ 独立、均为零均值）
\begin{equation}
  \E\left[\bm{r}\bm{r}^{\top}\right]
  = \E\left[\bm{v}\bm{v}^{\top}\right]+\bm{H}\E\left[\bm{e}\bm{e}^{\top}\right]\bm{H}^{\top}
  = \bm{R}+\bm{H}\bm{P}\bm{H}^{\top} = \bm{S}.
  \label{eq:09-S}
\end{equation}
\textbf{马氏距离}（Mahalanobis distance）定义为用 $\bm{S}$ 归一化后的平方：
\begin{equation}
  d^{2} = \bm{r}^{\top}\bm{S}^{-1}\bm{r}
  = \bm{r}^{\top}\left(\bm{H}\bm{P}\bm{H}^{\top}+\bm{R}\right)^{-1}\bm{r}.
  \label{eq:09-mahalanobis}
\end{equation}

\textbf{为什么它服从卡方分布？}对 $\bm{S}$ 做特征分解 $\bm{S}=\bm{U}\bm{\Lambda}_S\bm{U}^{\top}$（$\bm{U}$ 正交、$\bm{\Lambda}_S$ 对角正定），定义
\begin{equation}
  \bm{y} = \bm{\Lambda}_S^{-1/2}\bm{U}^{\top}\bm{r}.
\end{equation}
由于 $\bm{r}$ 是零均值高斯，$\bm{y}$ 也是零均值高斯；其协方差为
\begin{equation}
  \E\left[\bm{y}\bm{y}^{\top}\right]
  = \bm{\Lambda}_S^{-1/2}\bm{U}^{\top}\bm{S}\,\bm{U}\bm{\Lambda}_S^{-1/2}
  = \bm{\Lambda}_S^{-1/2}\bm{\Lambda}_S\bm{\Lambda}_S^{-1/2}=\bm{I}.
\end{equation}
即 $\bm{y}$ 的各分量是独立标准正态变量。于是
\begin{equation}
  d^{2} = \bm{r}^{\top}\bm{S}^{-1}\bm{r}
  = \bm{r}^{\top}\bm{U}\bm{\Lambda}_S^{-1}\bm{U}^{\top}\bm{r}
  = \bm{y}^{\top}\bm{y} = \sum_{i=1}^{m}y_i^{2}\ \sim\ \chi^{2}(m),
  \label{eq:09-chi2}
\end{equation}
其中 $m$ 是残差维数。卡方分布的定义见第~\ref{ch:02-math}~章。直观上：$d^{2}$ 是「$m$ 个独立标准正态变量的平方和」，它超过某个阈值的概率可以查表。

\textbf{用于外点剔除。}若观测是内点，$d^{2}$ 应落在卡方分布的高概率区间；若超出 $1-\alpha$ 分位数，就认为该观测异常并拒绝。对 $m=3$、$\alpha=0.05$，阈值 $\chi^{2}_{3,0.95}=7.815$；$\alpha=0.01$ 时阈值 $11.345$。工程上要注意：外点剔除必须基于「先验 + 观测」的$\bm{S}$，而不是只基于观测噪声 $\bm{R}$，否则会把「先验本来就不确定」的情况误判为外点。

\begin{pseudo}{序贯融合与卡方检验}
  \item 初始化先验 $\hat{\bm{x}}_0,\ \bm{P}_0$
  \item 对每个传感器观测 $\bm{z}_i=\bm{H}_i\bm{x}+\bm{v}_i$：
  \item \quad 计算残差 $\bm{r}_i\leftarrow\bm{z}_i-\bm{H}_i\hat{\bm{x}}$
  \item \quad 计算新息协方差 $\bm{S}_i\leftarrow\bm{H}_i\bm{P}\bm{H}_i^{\top}+\bm{R}_i$
  \item \quad 计算 $d_i^{2}\leftarrow\bm{r}_i^{\top}\bm{S}_i^{-1}\bm{r}_i$，若 $d_i^{2}$ 超过卡方阈值则拒绝该观测
  \item \quad $\bm{K}_i\leftarrow\bm{P}\bm{H}_i^{\top}\bm{S}_i^{-1}$
  \item \quad $\hat{\bm{x}}\leftarrow\hat{\bm{x}}+\bm{K}_i\bm{r}_i$，$\bm{P}\leftarrow(\bm{I}-\bm{K}_i\bm{H}_i)\bm{P}$
  \item 输出融合位姿与协方差
\end{pseudo}

\section{因子图视角：融合就是最大后验}\label{sec:09-factorgraph}

卡尔曼滤波是「逐时刻递推」的融合，而因子图（factor graph）是「全局一起解」的融合。把所有测量写成残差 $\bm{r}_i(\bm{x})$，假设残差零均值、协方差 $\bm{\Sigma}_i$，则最大后验估计（MAP）为
\begin{equation}
  \hat{\bm{x}} = \argmin_{\bm{x}}\sum_i\norm{\bm{r}_i(\bm{x})}_{\bm{\Sigma}_i}^{2}
  = \argmin_{\bm{x}}\sum_i\bm{r}_i(\bm{x})^{\top}\bm{\Sigma}_i^{-1}\bm{r}_i(\bm{x}),
  \label{eq:09-map}
\end{equation}
其中 $\norm{\bm{r}}^{2}_{\bm{\Sigma}}=\bm{r}^{\top}\bm{\Sigma}^{-1}\bm{r}$ 是马氏距离平方。对线性高斯情形，每个残差的雅可比为 $\bm{H}_i$，令梯度为零：
\begin{equation}
  \sum_i\bm{H}_i^{\top}\bm{\Sigma}_i^{-1}\left(\bm{r}_i^{0}+\bm{H}_i\delta\bm{x}\right)=\bm{0}
  \ \Longrightarrow\
  \left(\sum_i\bm{H}_i^{\top}\bm{\Sigma}_i^{-1}\bm{H}_i\right)\delta\bm{x}
  = -\sum_i\bm{H}_i^{\top}\bm{\Sigma}_i^{-1}\bm{r}_i^{0}.
  \label{eq:09-gn}
\end{equation}
左边正是式~\eqref{eq:09-information-sum} 的信息矩阵累加。也就是说：\textbf{滤波与优化在「信息形式」下是同一个东西}——滤波是优化在马尔可夫假设下的一种增量、递推解法。因子图的详细推导与求解见第~\ref{ch:13-backend}~章，可参考 \cite{dellaert2017factorgraphs}。

\section{时间同步与延时补偿}\label{sec:09-time}

前面所有公式都隐含一个假设：参与融合的测量是「同一时刻」的。现实中这是最容易出错的地方。

\textbf{为什么 10\,ms 会致命？}设机器人以 $2\,\mathrm{m/s}$ 运动，时间戳错 $10\,\mathrm{ms}$，则位姿的对应位置偏差约
\begin{equation}
  \delta p \approx v\,\delta t = 2\times0.01=0.02\,\mathrm{m}=2\,\mathrm{cm}.
\end{equation}
对激光里程计来说，$2\,\mathrm{cm}$ 的偏差在配准时会变成系统性残差，滤波器会把它当作「真实运动」去修正，于是状态被持续拖偏；如果时间错的是旋转量（角速度 $1\,\mathrm{rad/s}$ 错 $10\,\mathrm{ms}$ 即 $0.57\degree$），姿态也会被带偏。更隐蔽的是，时间戳错误往往表现为「缓慢漂移 + 偶发跳变」，很难归因。

\textbf{补偿方法。}设测量时刻为 $t$，需要把它对齐到统一时刻 $t_0$。若已有两个时刻的位姿 $T_a(T_a\text{ at }t_a)$、$T_b(T_b\text{ at }t_b)$，用 $\mathrm{SE}(3)$ 插值：
\begin{equation}
  T(t) = T_a\exp\left(\alpha\,\log\left(T_a^{-1}T_b\right)\right),
  \qquad
  \alpha=\frac{t-t_a}{t_b-t_a},
  \label{eq:09-interp}
\end{equation}
其中 $\log(\cdot)$ 是 $\exp(\cdot)$ 的逆映射（把 $SE(3)$ 元素变回六维 twist）。$\alpha=0$ 时得到 $T_a$，$\alpha=1$ 时得到 $T_b$，中间按指数映射平滑过渡。若只有速度，也可以做一阶补偿：$T(t)\approx T_a\exp((t-t_a)\bm{\xi}^{\wedge})$。

工程上还有两件事必须做：一是\textbf{统一时钟}——相机、IMU、激光的驱动都把时间戳打到同一时钟源（ROS 2 用 \texttt{node->now()} 或硬件触发）；二是\textbf{检查频率与延迟}——用 \texttt{ros2 topic hz} 与 \texttt{ros2 topic delay} 确认没有丢帧与过大延迟。ROS 2 的 \texttt{message\_filters} 提供 \texttt{ApproximateTime} 策略，可以在时间戳不完全一致时按最近邻配对，但它只是「配对」，并不能消除真正的延时误差。

\begin{figure}[htbp]
  \centering
  \begin{tikzpicture}[scale=1.0]
    \draw[-stealth, thick] (0,0) -- (10.5,0) node[right, font=\footnotesize] {时间};
    \foreach \x/\lab in {1/相机, 3/IMU, 5/激光, 7/相机, 9/IMU} {
      \draw[thick] (\x,0.1) -- (\x,-0.35);
      \node[font=\footnotesize, rotate=45, anchor=north east] at (\x,-0.35) {\lab};
    }
    \draw[thick, structurecolor] (3,0.25) -- (7,0.25);
    \draw[-stealth, structurecolor] (3,0.25) -- (3,0.55);
    \draw[-stealth, structurecolor] (7,0.25) -- (7,0.55);
    \node[font=\footnotesize, structurecolor, above] at (5,0.55) {IMU 预积分区间};
    \draw[<->, red!60!black] (5,1.35) -- (6.6,1.35) node[midway, above, font=\footnotesize] {$\delta t$ 延时};
    \draw[dashed, red!60!black] (5,1.3) -- (5.8,-0.3);
    \draw[dashed, red!60!black] (6.6,1.3) -- (6.6,-0.3);
  \end{tikzpicture}
  \caption{时间同步示意。不同传感器的时间戳必须对齐到同一时钟；若激光帧的位姿被错误地关联到 $\delta t$ 之前的状态，残差里就会混入 $v\,\delta t$ 的系统偏差。}
  \label{fig:09-time}
\end{figure}

\section{典型融合架构}\label{sec:09-arch}

\textbf{FAST-LIO（激光 + IMU 紧耦合）}\cite{xu2022fastlio2}：用迭代扩展卡尔曼滤波（IEKF）把激光点到面残差直接作为观测更新，状态含位置、速度、姿态、零偏与重力；IMU 负责预测，激光负责更新；点云通过增量式 KD 树直接配准到地图，不需要提取特征。它的核心优势是把「激光退化方向靠 IMU 先验补」写进了同一个滤波器。详见第~\ref{ch:17-advanced-lio}~章。

\textbf{VINS-Mono（视觉 + IMU 紧耦合）}\cite{qin2018vinsmono}：视觉重投影残差与 IMU 预积分残差放进滑动窗口联合优化，窗口滑出时边缘化；视觉前端用光流跟踪，后端用 Ceres 求解。详见第~\ref{ch:16-lio-evolution}~章与第~\ref{ch:13-backend}~章。

\textbf{Z-Localization（前后端分离）}\cite{zloc2026}：前端做高频的里程计与局部跟踪，后端做低频的全局匹配与位姿图优化；前后端通过位姿与协方差解耦，兼顾实时性与全局一致性。详见第~\ref{ch:23-zlocalization}~章。

\begin{keybox}{融合系统的排错顺序}
\begin{enumerate}
  \item 先查时间戳：每个话题的 header.stamp 是否单调、是否来自同一时钟；
  \item 再查外参：传感器之间的 $T_{SB}$ 是否标定、符号是否正确；
  \item 再查协方差：$\bm{R}$、$\bm{Q}$ 的量级是否与真实噪声匹配（过大/过小都会让融合偏科）；
  \item 最后才调参数：在数据一致的前提下，融合参数才是有意义的。
\end{enumerate}
\end{keybox}

\begin{summary}

\begin{enumerate}
  \item 融合分数据级、特征级、决策级；按接口分松耦合与紧耦合，紧耦合保留原始测量、精度更高，代价是复杂度与标定要求。
  \item 状态增广后多传感器融合就是「一个状态、多组观测」；信息形式下融合 = 信息矩阵与信息向量相加：$\bm{\Lambda}=\bm{P}_0^{-1}+\sum\bm{H}_i^{\top}\bm{R}_i^{-1}\bm{H}_i$。
  \item 序贯更新与批量更新在信息形式下完全等价，观测顺序不影响结果。
  \item 信息矩阵 $\bm{\Lambda}=\bm{P}^{-1}$；马氏距离 $d^{2}=\bm{r}^{\top}(\bm{H}\bm{P}\bm{H}^{\top}+\bm{R})^{-1}\bm{r}$，在零均值高斯假设下服从 $\chi^{2}(m)$。
  \item 卡方检验用于外点剔除：$d^{2}>\chi^{2}_{m,1-\alpha}$ 时拒绝观测；$\bm{S}$ 必须包含先验不确定性。
  \item 因子图把融合写成 MAP：$\min\sum\norm{\bm{r}_i}^{2}_{\bm{\Sigma}_i}$，线性化后的正规方程与信息滤波一致。
  \item 时间同步是融合成败的关键：$2\,\mathrm{m/s}$ 下 $10\,\mathrm{ms}$ 错位就是 $2\,\mathrm{cm}$；用 $\mathrm{SE}(3)$ 插值补偿，并用统一时钟源避免错误。
  \item 典型架构：FAST-LIO（激光-IMU 紧耦合）、VINS-Mono（视觉-IMU 紧耦合）、Z-Localization（前后端分离）。
\end{enumerate}


\end{summary}

\begin{problemset}

  \item 用一维情形验证式~\eqref{eq:09-batch}：两个观测 $z_1=x+v_1$、$z_2=x+v_2$，方差分别为 $r_1,r_2$，先验方差 $p_0$，分别用批量与序贯公式求后验方差，验证相同。
  \item 证明式~\eqref{eq:09-information-sum}：从卡尔曼更新的信息形式出发，写出 $\bm{P}^{-1}$ 的更新式，并说明为什么观测顺序无关。
  \item 推导式~\eqref{eq:09-residual} 与式~\eqref{eq:09-S}，说明为什么 $\bm{r}$ 的协方差不是 $\bm{R}$。
  \item 对 $m=2$、$\alpha=0.05$ 求卡方阈值（查表或记忆），并解释为什么外点剔除阈值不能直接用 $\bm{R}$ 归一化。
  \item 某系统相机在 $t=1.000\,\mathrm{s}$、IMU 在 $t=1.008\,\mathrm{s}$ 各有一组数据，机器人速度 $1.5\,\mathrm{m/s}$、角速度 $0.8\,\mathrm{rad/s}$，估算不对齐造成的位移与角度偏差。
  \item 解释为什么「FAST-LIO 在长廊里仍能用」而「纯激光里程计会退化」，用本章的紧耦合概念与第~\ref{ch:07-lidar-odom}~章的退化分析回答。
\end{problemset}

\begin{reading}

\begin{itemize}
  \item Thrun 等《Probabilistic Robotics》第 2--3 章：高斯滤波与信息形式的经典推导。
  \item Dellaert 的因子图教程（\href{https://gtsam.org/}{gtsam.org}）：从 MAP 到因子图的完整视角。
  \item FAST-LIO2 论文与代码（\href{https://github.com/hku-mars/FAST_LIO}{GitHub}）：紧耦合 IEKF 的工程实现。
\end{itemize}


\end{reading}

\chapter{全局定位与重定位}\label{ch:10-global-loc}

\chapterintro{里程计回答的是「我相对刚才动了多少」，但它永远回答不了「我现在在地图的哪里」。机器人可能被放到赛场的任意角落开机，可能在比赛中被撞得转了三圈，也可能在长廊里迷失——这些情况都需要\textbf{全局定位}：在一张先验地图中确定自己的绝对位姿。本章讲清楚三件事：为什么里程计无法解决全局定位、粒子滤波为什么是解决它的标准工具（从贝叶斯递归一路推到重采样）、以及 AMCL 这个 ROS 里最常见的实现是怎么做的。学完本章，你应当能解释「绑架问题」为什么难，以及 KLD 采样公式里每个符号的含义。}

\section{局部定位与全局定位}\label{sec:10-local-global}

第~\ref{ch:04-odom-intro}~章已经区分了绝对定位与相对定位。从地图的使用方式看，还有一个更实用的分法：

\begin{itemize}
  \item \textbf{局部定位}（local localization，也叫位姿跟踪）：已知初始位姿（或初始位姿误差很小），在运动过程中持续修正。它依赖局部地图匹配，初值附近只有一个「吸引域」，一旦偏差超过这个范围就会跟丢。
  \item \textbf{全局定位}（global localization）：初始位姿未知，需要在地图全局搜索所有可能的位姿，并从中挑出与当前观测一致的解。
\end{itemize}

里程计只能做局部定位：它提供的位姿是相对初始时刻的，初始时刻在地图里的位置它并不知道；而且里程计的漂移是累积的，时间越长，它离真实位姿越远，最终会超出任何匹配算法的收敛域。所以全局定位必须有「地图先验」和「全局搜索」两样东西。

\textbf{绑架问题}（kidnapped robot problem）是全局定位的极端版本：机器人在正常运行时被突然「瞬移」到地图的另一个位置（现实中可能是被搬起、被撞飞、或开机后被人挪动）。此时它自己的里程计还停留在旧位置，所有基于「上一时刻位姿」的跟踪都会失效，必须能够主动怀疑自己、重新搜索全局。

\section{基于地图的定位方法}\label{sec:10-methods}

按匹配与搜索的方式，基于地图的定位可以分为四类：

\begin{itemize}
  \item \textbf{栅格匹配}：把激光扫描与占据栅格地图做相关匹配，在位姿空间里搜索得分最高的位置。2D 场景下典型代表是相关性扫描匹配（第~\ref{ch:07-lidar-odom}~章）。
  \item \textbf{特征匹配}：从地图与当前扫描中提取几何特征（直线、角点、平面），用特征对应关系解算位姿。结构化场景（走廊、房间）效果好，但特征提取本身可能失败。
  \item \textbf{描述子检索 + 配准}：先用地点的全局描述子在地图数据库中检索候选位置，再做局部配准。Scan Context \cite{kim2018scancontext} 是典型代表：把点云按「环—扇区」划分成二维矩阵，每个元素取该区域最大高度，得到一个对旋转敏感、对平移鲁棒的描述子；检索时按列移位匹配，找到最相似的历史地点。这类方法适合大场景重定位与回环检测。
  \item \textbf{粒子滤波}：用一组带权样本表示位姿的置信度分布，能同时维护多个假设，天然适合全局定位与绑架恢复。AMCL（adaptive Monte Carlo localization）就是它的工程实现 \cite{fox2003amcl}。
\end{itemize}

本章重点讲粒子滤波，因为它是唯一能同时处理「多峰分布」与「全局搜索」的通用框架，也是理解 AMCL、Cartographer 定位模式、以及第~\ref{ch:23-zlocalization}~章重定位模块的基础。

\section{贝叶斯滤波递归}\label{sec:10-bayes}

全局定位的本质是「在给定历史观测与运动的情况下，推断当前位姿的后验分布」。记时刻 $t$ 的位姿为 $\bm{x}_t$，控制/里程计输入为 $\bm{u}_t$，观测为 $\bm{z}_t$，历史为 $\bm{z}_{1:t}=\{\bm{z}_1,\dots,\bm{z}_t\}$、$\bm{u}_{1:t}$。定义\textbf{置信度}（belief）
\begin{equation}
  \mathrm{bel}(\bm{x}_t) = p\left(\bm{x}_t\mid\bm{z}_{1:t},\bm{u}_{1:t}\right).
  \label{eq:10-belief}
\end{equation}
下面从条件概率与贝叶斯公式把它写成递推形式。由条件概率定义 $p(a\mid b)=p(a,b)/p(b)$，有
\begin{equation}
  p\left(\bm{x}_t\mid\bm{z}_{1:t},\bm{u}_{1:t}\right)
  = \frac{p\left(\bm{z}_t\mid\bm{x}_t,\bm{z}_{1:t-1},\bm{u}_{1:t}\right)
    p\left(\bm{x}_t\mid\bm{z}_{1:t-1},\bm{u}_{1:t}\right)}
  {p\left(\bm{z}_t\mid\bm{z}_{1:t-1},\bm{u}_{1:t}\right)}.
  \label{eq:10-bayes-step1}
\end{equation}
分母只与观测有关、与 $\bm{x}_t$ 无关，可以写成一个归一化常数 $\eta$。对第一个因子用\textbf{马尔可夫假设}（观测只依赖当前状态）：$p(\bm{z}_t\mid\bm{x}_t,\bm{z}_{1:t-1},\bm{u}_{1:t})=p(\bm{z}_t\mid\bm{x}_t)$。对第二个因子，把 $\bm{x}_{t-1}$ 作为中间变量引入全概率公式：
\begin{equation}
  p\left(\bm{x}_t\mid\bm{z}_{1:t-1},\bm{u}_{1:t}\right)
  = \int p\left(\bm{x}_t\mid\bm{x}_{t-1},\bm{z}_{1:t-1},\bm{u}_{1:t}\right)
    p\left(\bm{x}_{t-1}\mid\bm{z}_{1:t-1},\bm{u}_{1:t}\right)\diff\bm{x}_{t-1}.
\end{equation}
再用一次马尔可夫假设：$\bm{x}_t$ 只依赖 $\bm{x}_{t-1}$ 与当前输入 $\bm{u}_t$，即 $p(\bm{x}_t\mid\bm{x}_{t-1},\bm{z}_{1:t-1},\bm{u}_{1:t})=p(\bm{x}_t\mid\bm{x}_{t-1},\bm{u}_t)$；且 $\bm{x}_{t-1}$ 不依赖未来的输入 $\bm{u}_t$，故 $p(\bm{x}_{t-1}\mid\bm{z}_{1:t-1},\bm{u}_{1:t})=\mathrm{bel}(\bm{x}_{t-1})$。代回式~\eqref{eq:10-bayes-step1} 得
\begin{equation}
  \mathrm{bel}(\bm{x}_t)
  = \eta\,p\left(\bm{z}_t\mid\bm{x}_t\right)
  \int p\left(\bm{x}_t\mid\bm{x}_{t-1},\bm{u}_t\right)\mathrm{bel}(\bm{x}_{t-1})\diff\bm{x}_{t-1}.
  \label{eq:10-bayes-recursion}
\end{equation}
式~\eqref{eq:10-bayes-recursion} 就是\textbf{贝叶斯滤波递归}：$\int p(\bm{x}_t\mid\bm{x}_{t-1},\bm{u}_t)\mathrm{bel}(\bm{x}_{t-1})\diff\bm{x}_{t-1}$ 是「预测」，用运动模型把上一时刻的置信度推演到当前时刻；乘以 $p(\bm{z}_t\mid\bm{x}_t)$ 是「更新」，用当前观测给预测加权；$\eta$ 是归一化。卡尔曼滤波、粒子滤波、直方图滤波都是它的不同实现：区别只在「怎么表示 bel」。

\section{粒子滤波}\label{sec:10-pf}

贝叶斯递归里的积分一般没有解析解，粒子滤波的解法是：用 $N$ 个带权样本（粒子）近似置信度：
\begin{equation}
  \mathrm{bel}(\bm{x}_t) \approx \sum_{i=1}^{N} w_t^{i}\,\delta\left(\bm{x}_t-\bm{x}_t^{i}\right),
  \label{eq:10-particles}
\end{equation}
其中 $\bm{x}_t^{i}$ 是第 $i$ 个粒子（一个位姿假设），$w_t^{i}$ 是它的归一化权重（$\sum_i w_t^i=1$），$\delta$ 是狄拉克函数。位姿的估计值可以取加权均值，也可以取权重最大的粒子。

\subsection{重要性采样}

直接按 $\mathrm{bel}$ 采样做不到（$\mathrm{bel}$ 本身就是要估计的量）。\textbf{重要性采样}（importance sampling）的思路是：从一个容易采样的\textbf{提议分布}（proposal distribution）$q(\bm{x})$ 采样，再用权重修正它与 $\mathrm{bel}$ 的差异。对任意函数 $f$，其期望为
\begin{align}
  \E_{\mathrm{bel}}\left[f(\bm{x})\right]
  &= \int f(\bm{x})\,\mathrm{bel}(\bm{x})\diff\bm{x} \notag\\
  &= \int f(\bm{x})\,\frac{\mathrm{bel}(\bm{x})}{q(\bm{x})}\,q(\bm{x})\diff\bm{x} \notag\\
  &= \E_{q}\left[f(\bm{x})\,w(\bm{x})\right],
  \qquad
  w(\bm{x}) = \frac{\mathrm{bel}(\bm{x})}{q(\bm{x})}.
  \label{eq:10-is}
\end{align}
第二个等号是在积分里「乘一个 $q$ 再除一个 $q$」，前提是 $q(\bm{x})>0$ 在所有 $\mathrm{bel}(\bm{x})>0$ 处成立。实践中 $\mathrm{bel}$ 只知道到常数因子（归一化常数未知），因此改用未归一化权重 $\tilde{w}(\bm{x})=\mathrm{bel}(\bm{x})/q(\bm{x})$，再归一化：
\begin{equation}
  w^{i} = \frac{\tilde{w}^{i}}{\sum_{j=1}^{N}\tilde{w}^{j}}.
  \label{eq:10-w-normalize}
\end{equation}

\subsection{序贯重要性采样}

把重要性采样用到贝叶斯递归上，得到\textbf{序贯重要性采样}（sequential importance sampling, SIS）。关键选择：用「运动模型」作为提议分布，即从 $q(\bm{x}_t\mid\bm{x}_{t-1}^{i},\bm{u}_t)=p(\bm{x}_t\mid\bm{x}_{t-1}^{i},\bm{u}_t)$ 采样当前粒子。此时权重可以递推。由权重定义
\begin{equation}
  w_t^{i} = \frac{\mathrm{bel}(\bm{x}_t^{i})}{q(\bm{x}_t^{i}\mid\bm{x}_{t-1}^{i},\bm{u}_t)}.
\end{equation}
把式~\eqref{eq:10-bayes-recursion} 的 $\mathrm{bel}(\bm{x}_t^{i})\propto p(\bm{z}_t\mid\bm{x}_t^{i})\int p(\bm{x}_t^{i}\mid\bm{x}_{t-1},\bm{u}_t)\mathrm{bel}(\bm{x}_{t-1})\diff\bm{x}_{t-1}$ 代入，并用粒子近似把积分写成求和：
\begin{equation}
  \mathrm{bel}(\bm{x}_t^{i}) \propto p\left(\bm{z}_t\mid\bm{x}_t^{i}\right)
  \sum_{j=1}^{N}w_{t-1}^{j}\,p\left(\bm{x}_t^{i}\mid\bm{x}_{t-1}^{j},\bm{u}_t\right).
\end{equation}
对我们从 $\bm{x}_{t-1}^{i}$ 采样得到的那个粒子，求和里只有 $j=i$ 的项被采样到，其余项在粒子近似下不贡献，于是
\begin{equation}
  \mathrm{bel}(\bm{x}_t^{i}) \propto p\left(\bm{z}_t\mid\bm{x}_t^{i}\right)w_{t-1}^{i}\,p\left(\bm{x}_t^{i}\mid\bm{x}_{t-1}^{i},\bm{u}_t\right).
\end{equation}
把它代回权重定义，提议分布恰好把这个运动模型因子约掉：
\begin{equation}
  w_t^{i} \propto \frac{p\left(\bm{z}_t\mid\bm{x}_t^{i}\right)w_{t-1}^{i}\,p\left(\bm{x}_t^{i}\mid\bm{x}_{t-1}^{i},\bm{u}_t\right)}{p\left(\bm{x}_t^{i}\mid\bm{x}_{t-1}^{i},\bm{u}_t\right)}
  = w_{t-1}^{i}\,p\left(\bm{z}_t\mid\bm{x}_t^{i}\right).
  \label{eq:10-weight-recursion}
\end{equation}
这就是粒子滤波最常用的权重更新式：\textbf{新权重 = 旧权重 × 观测似然}。它把「运动模型采样」与「观测加权」两件事分得干干净净。归一化后即得式~\eqref{eq:10-w-normalize}。

\begin{keybox}{粒子滤波的三步循环}
\begin{enumerate}
  \item \textbf{预测}：对每个粒子，按运动模型 $p(\bm{x}_t\mid\bm{x}_{t-1}^{i},\bm{u}_t)$ 采样一个新位姿（把里程计增量加上噪声）；
  \item \textbf{加权}：用观测模型算 $p(\bm{z}_t\mid\bm{x}_t^{i})$，乘到旧权重上并归一化；
  \item \textbf{重采样}：按权重重新抽取 $N$ 个粒子，权重重置为 $1/N$。
\end{enumerate}
循环往复，粒子会逐渐聚集到与观测一致的位置。
\end{keybox}

\subsection{重采样与样本贫化}

只加权不重采样，运行一段时间后会出现\textbf{权重退化}：绝大多数粒子的权重趋近于零，有效粒子数急剧减少，计算资源被浪费在无意义的粒子上。重采样的做法是「按权重的比例复制粒子」：权重大的粒子被多复制，权重小的被淘汰。重采样后所有粒子权重相等，再继续下一轮。

重采样也有代价：它会让粒子多样性下降（同一个高权重粒子被复制多份），长期运行后可能出现\textbf{样本贫化}（sample impoverishment）——粒子全都挤在一个假设附近，丧失探索其他可能位置的能力。缓解手段包括：只在必要时重采样、加入少量随机粒子、使用低方差重采样算法。

\textbf{何时重采样？}用\textbf{有效粒子数}（effective sample size）判断：
\begin{equation}
  N_{\text{eff}} = \frac{1}{\sum_{i=1}^{N}\left(w^{i}\right)^{2}}.
  \label{eq:10-neff}
\end{equation}
它的含义是：这组带权粒子「等效于」多少个等权粒子。当所有权重相等（$w^i=1/N$）时，$N_{\text{eff}}=N$；当一个粒子独占权重时，$N_{\text{eff}}=1$。典型做法是 $N_{\text{eff}}<N/2$ 时才重采样。

\textbf{系统重采样}是最常用的低方差重采样算法：把权重累积成区间 $[0,1)$，产生一个在 $[0,1/N)$ 内均匀分布的随机数 $u$，然后依次取 $u,\ u+1/N,\ u+2/N,\dots$，看每个点落在哪个粒子的累积区间里，就复制该粒子。它只用一次随机数，方差比「独立重复抽样」小。

\begin{pseudo}{系统重采样}
  \item 输入：粒子 $\{\bm{x}^{i}\}_{i=1}^{N}$ 与权重 $\{w^{i}\}$（$\sum_i w^i=1$）
  \item 计算累积权重：$c_1\leftarrow w^{1}$，$c_i\leftarrow c_{i-1}+w^{i}$
  \item 采样起点 $u\sim\mathcal{U}\left[0,\ 1/N\right)$，令 $r\leftarrow u$
  \item 对 $k=1,\dots,N$：
  \item \quad 找到最小的 $i$ 使 $c_i\ge r$，复制粒子 $\bm{x}^{i}$ 到新集合
  \item \quad $r\leftarrow r+1/N$
  \item 输出：新粒子集合，权重全部置为 $1/N$
\end{pseudo}

\begin{figure}[htbp]
  \centering
  \begin{tikzpicture}[scale=0.9]
    \foreach \dx/\lab in {0/{(a) 预测}, 4.4/{(b) 加权}, 8.8/{(c) 重采样}} {
      \draw[thick, black!30] (\dx-0.4,-0.3) rectangle (\dx+3.4,3.3);
      \node[font=\footnotesize] at (\dx+1.5,3.6) {\lab};
    }
    \foreach \x/\y/\r in {0.4/2.4/3, 1.2/1.1/3, 2.0/2.6/3, 0.8/0.4/3, 2.6/1.6/3, 2.2/0.6/3} {
      \draw[structurecolor, fill=structurecolor!25] (\x,\y) circle (\r pt);
    }
    \node[font=\footnotesize, structurecolor] at (1.4,-0.7) {粒子按运动模型散开};
    \foreach \x/\y/\r in {4.8/2.4/10, 5.6/1.1/4, 6.4/2.6/7, 5.2/0.4/2, 7.0/1.6/6, 6.6/0.6/3} {
      \draw[structurecolor, fill=structurecolor!25] (\x,\y) circle (\r pt);
    }
    \node[font=\footnotesize, structurecolor] at (5.8,-0.7) {圆的面积正比于权重};
    \foreach \x/\y/\r in {9.6/2.4/6, 9.6/2.2/6, 10.4/2.6/6, 9.8/0.9/5, 10.2/1.0/5, 9.9/2.5/6} {
      \draw[structurecolor, fill=structurecolor!25] (\x,\y) circle (\r pt);
    }
    \node[font=\footnotesize, structurecolor] at (10.2,-0.7) {高权重粒子被复制};
  \end{tikzpicture}
  \caption{粒子滤波的一轮迭代：预测让粒子按运动模型散开并加入噪声；加权让与观测一致的粒子变大；重采样把高权重粒子复制、低权重粒子淘汰。}
  \label{fig:10-pf}
\end{figure}

\section{AMCL：工程化的粒子滤波}\label{sec:10-amcl}

AMCL \cite{fox2003amcl} 是 ROS 中实现 2D 全局定位的标准模块。它的输入是占据栅格地图、激光扫描与里程计，输出是 \texttt{map}$\to$\texttt{odom} 变换（以及位姿协方差）。它把粒子滤波拆成三个具体模型。

\subsection{运动模型：从里程计增量采样}

里程计给出相邻时刻的位姿增量，把它分解成三段：先原地转 $\delta_{\text{rot1}}$，再直行 $\delta_{\text{trans}}$，最后原地转 $\delta_{\text{rot2}}$。真实运动与里程计读数有偏差，AMCL 按如下方式加噪声再执行（$\alpha_1$--$\alpha_4$ 是噪声系数）：
\begin{align}
  \hat{\delta}_{\text{rot1}} &= \delta_{\text{rot1}}
  + \mathrm{noise}\left(\alpha_1\delta_{\text{rot1}}^{2}+\alpha_2\delta_{\text{trans}}^{2}\right), \notag\\
  \hat{\delta}_{\text{trans}} &= \delta_{\text{trans}}
  + \mathrm{noise}\left(\alpha_3\delta_{\text{trans}}^{2}+\alpha_4\left(\delta_{\text{rot1}}^{2}+\delta_{\text{rot2}}^{2}\right)\right), \notag\\
  \hat{\delta}_{\text{rot2}} &= \delta_{\text{rot2}}
  + \mathrm{noise}\left(\alpha_1\delta_{\text{rot2}}^{2}+\alpha_2\delta_{\text{trans}}^{2}\right),
  \label{eq:10-motion-model}
\end{align}
其中 $\mathrm{noise}(\sigma^{2})$ 表示从零均值高斯分布中采样，方差为 $\sigma^{2}$。这个模型的物理含义是：转得越多、走得越远，误差越大；而且「走」与「转」的误差相互耦合（走的时候会带一点转，转的时候会带一点走）。

\subsection{观测模型：似然场}

给定一个粒子位姿，怎么算「在这个位姿下看到当前激光扫描的概率」？AMCL 用\textbf{似然场}（likelihood field）模型：对每条激光束，取它的末端点在全局地图中的位置，查「该点到最近障碍物的距离」$d$，然后
\begin{equation}
  p\left(\bm{z}_t\mid\bm{x}_t^{i}\right) = \prod_{k=1}^{K}\left[
  z_{\text{hit}}\cdot\exp\left(-\frac{d_k^{2}}{2\sigma_{\text{hit}}^{2}}\right)
  + \frac{z_{\text{rand}}}{z_{\max}}\right],
  \label{eq:10-likelihood-field}
\end{equation}
其中 $K$ 是激光束数，$d_k$ 是第 $k$ 条束末端到最近障碍的距离，$z_{\text{hit}}$ 是「命中」权重、$z_{\text{rand}}$ 是随机噪声权重，$z_{\max}$ 是最大量程。这个模型的好处是不需要做射线投射（ray casting），只需查一张预先算好的距离场（ESDF 的思想），速度快；缺点是它对「光束穿过墙」不敏感（只看到端点）。

\subsection{KLD 采样：自适应粒子数}

粒子数 $N$ 怎么定？固定值要么太浪费（简单场景），要么不够用（复杂场景）。AMCL 采用 \textbf{KLD 采样}（Kullback-Leibler divergence sampling）：边采样边观察「落到了多少个不同的状态格（bin）」$k$，当 $k$ 不再明显增长时停止。所需粒子数由下式给出：
\begin{equation}
  n = \frac{1}{2\varepsilon^{2}}\,\chi^{2}_{k-1,\,1-\delta},
  \label{eq:10-kld}
\end{equation}
其中各符号含义如下：
\begin{itemize}
  \item $\varepsilon$：允许的误差尺度。本书按提纲约定，把它对应的 KL 散度上界记作 $\varepsilon^{2}$（有的文献直接把该上界记作 $\varepsilon$，公式就写成 $n=\chi^{2}_{k-1,1-\delta}/(2\varepsilon)$，两者只差一个记号约定）；
  \item $\delta$：允许的失败概率，即真实分布与样本分布偏离超过上界的概率；
  \item $k$：当前已经被采样到的不同 bin 的数量，bin 越多样本越分散，需要的粒子越多；
  \item $\chi^{2}_{k-1,1-\delta}$：自由度为 $k-1$ 的卡方分布的 $1-\delta$ 分位数（卡方分布见第~\ref{ch:02-math}~章）。
\end{itemize}
这个公式来自多项分布的统计推断：$k$ 个 bin 上的样本计数服从多项分布，其 Pearson 统计量 $\sum_j(N_j-np_j)^{2}/(np_j)$ 近似服从 $\chi^{2}(k-1)$；要让经验分布与真实分布的 KL 散度以概率 $1-\delta$ 不超过上界，就需要 $n$ 至少取上式。直观上：bin 越多、要求越准、要求越有把握，粒子就越多。

\begin{example}[KLD 采样的粒子数]
设当前采样到 $k=30$ 个 bin，$\delta=0.01$（要求 $99\,\%$ 的把握），查表得 $\chi^{2}_{29,0.99}\approx49.6$。若取 $\varepsilon=0.2$，则
\begin{equation}
  n=\frac{49.6}{2\times0.2^{2}}=\frac{49.6}{0.08}\approx620.
\end{equation}
也就是说这个场景下 $600$ 多个粒子就够了。若 $\varepsilon$ 缩小到 $0.1$（精度要求翻倍），粒子数会变成约 $2480$——这说明「精度」与「计算量」是平方关系，调参时要格外小心。
\end{example}

\section{重定位与绑架恢复}\label{sec:10-relocal}

有了全局定位的框架，重定位的具体策略还需要考虑工程可行性：

\begin{itemize}
  \item \textbf{初始化}：机器人开机时，把粒子均匀撒在地图的可行区域（或按已知的初始区域撒），第一帧观测后权重会迅速集中在少数假设上。若初始位置已知（如比赛固定起点），可以直接在起点附近撒，收敛更快。
  \item \textbf{丢失恢复}：持续监控粒子权重与观测似然。如果长时间似然很低、$N_{\text{eff}}$ 很小，说明当前假设与观测不符，可能是跟丢了；此时注入随机粒子（在全局或局部重新撒），让滤波器有机会「重新找到自己」。
  \item \textbf{多假设}：人在长廊里也会犹豫，粒子滤波天然维护多个假设。不要过早把粒子聚成一个——在对称场景（两条相似的走廊）里，保留多个峰比强行选一个更安全。
  \item \textbf{开关式重定位}：把「局部跟踪」与「全局重定位」做成两个模式，局部跟踪置信度低时切到全局模式，置信度恢复后切回。这是第~\ref{ch:23-zlocalization}~章 Z-Localization 前后端分离思想的一种体现。
\end{itemize}

\begin{keybox}{全局定位的四个工程要点}
\begin{enumerate}
  \item 粒子数要自适应：用 KLD 采样，别硬编码一个巨大值；
  \item 观测模型要快：似然场查表比射线投射快一个数量级；
  \item 重采样要克制：$N_{\text{eff}}$ 低于阈值再重采样，避免样本贫化；
  \item 要知道自己可能错：永远保留随机粒子与多假设，才能从绑架中恢复。
\end{enumerate}
\end{keybox}

\begin{summary}

\begin{enumerate}
  \item 局部定位是位姿跟踪，全局定位要在整张地图上搜索；里程计只能做局部，因为它不知道初始位姿且会累积漂移。
  \item 基于地图的定位分栅格匹配、特征匹配、描述子检索 + 配准、粒子滤波四类；Scan Context 是描述子检索的代表。
  \item 贝叶斯滤波递归：$\mathrm{bel}(\bm{x}_t)=\eta\,p(\bm{z}_t\mid\bm{x}_t)\int p(\bm{x}_t\mid\bm{x}_{t-1},\bm{u}_t)\mathrm{bel}(\bm{x}_{t-1})\diff\bm{x}_{t-1}$，由条件概率、贝叶斯公式与两处马尔可夫假设推出。
  \item 粒子滤波用带权样本近似置信度；序贯重要性采样的权重递推为 $w_t^{i}\propto w_{t-1}^{i}p(\bm{z}_t\mid\bm{x}_t^{i})$。
  \item 重采样按权重复制粒子，代价是样本贫化；用 $N_{\text{eff}}=1/\sum(w^{i})^{2}$ 判断是否需要重采样，系统重采样方差最小。
  \item AMCL 由运动模型采样、似然场观测模型、KLD 自适应粒子数三部分组成；KLD 公式 $n=\chi^{2}_{k-1,1-\delta}/(2\varepsilon^{2})$ 中 $\varepsilon$ 是误差尺度、$\delta$ 是失败概率、$k$ 是 bin 数。
  \item 绑架问题需要「怀疑自己」的机制：随机粒子注入、多假设、开关式重定位是常用手段。
\end{enumerate}


\end{summary}

\begin{problemset}

  \item 从条件概率定义出发，重推贝叶斯滤波递归式~\eqref{eq:10-bayes-recursion}，并在每一步旁标注用了哪条马尔可夫假设。
  \item 解释为什么提议分布取运动模型时，权重递推式~\eqref{eq:10-weight-recursion} 里没有运动模型项；如果提议分布改成「带噪声的运动模型」，权重会变成什么形式？
  \item 设 $N=5$，权重为 $\{0.4,0.3,0.2,0.05,0.05\}$，计算 $N_{\text{eff}}$，并手动执行一次系统重采样（给定 $u=0.03$），写出新粒子集合。
  \item 用式~\eqref{eq:10-kld} 计算 $k=50$、$\delta=0.05$、$\varepsilon=0.15$ 时的粒子数（$\chi^{2}_{49,0.95}\approx66.3$）。
  \item 说明似然场模型式~\eqref{eq:10-likelihood-field} 中 $z_{\text{hit}}$、$z_{\text{rand}}$、$z_{\max}$ 各自的物理作用，以及 $d_k$ 是怎么算出来的。
  \item 设计一个「绑架检测」方案：用哪些量、什么阈值、检出后如何恢复？请结合本章内容说明。
\end{problemset}

\begin{reading}

\begin{itemize}
  \item Thrun 等《Probabilistic Robotics》第 8 章：粒子滤波与 Monte Carlo 定位的系统推导。
  \item Fox 等 2003 年的 AMCL 论文（\href{https://ieeexplore.ieee.org/document/1249223}{IEEE TPAMI}）：似然场模型与 KLD 采样。
  \item Scan Context 论文与代码（\href{https://github.com/irapkaist/scancontext}{GitHub}）：大场景地点检索的工程实现。
\end{itemize}


\end{reading}


\part{SLAM 篇：前端、后端与建图}

% 本章使用 TikZ 的 arrows.meta / positioning 库（main.tex 已加载，此处保证单章自检也能编译）
\usetikzlibrary{arrows.meta, positioning}

\chapter{SLAM 框架：前端与后端}\label{ch:11-slam-framework}

\chapterintro{从这一章开始进入 SLAM 篇。前面两部分我们先后回答了三个问题：位置、姿态与坐标变换怎么描述（第~\ref{ch:02-math}~章），传感器能测到什么（第~\ref{ch:03-sensors}~章），以及只靠里程计能不能定位（第~\ref{ch:04-odom-intro}~章至第~\ref{ch:10-global-loc}~章）。里程计给出的答案有一个绕不开的宿命：误差随时间累积，而它自己没有任何机制把误差``拉回来''。SLAM 要回答的正是这个问题——在\emph{没有}先验地图的未知环境里，机器人能否一边估计自己的轨迹、一边把地图建出来？本章不写具体算法，而是先把骨架立起来：把 SLAM 写成最大后验（MAP）估计，定义前端与后端，讲清它们各自负责什么、凭什么可以分开，并特别澄清一个新手最容易混淆的点——在 SLAM、定位系统、导航系统三种语境下，``前端/后端''指的并不是同一件事。读完本章，你应该能看懂任何一套 SLAM 系统的框图，并在读源码时迅速判断``这段代码属于哪一层''。}

\section{先有鸡还是先有蛋：SLAM 到底难在哪}\label{sec:11-problem}

先做一个思想实验。假设机器人手里已经有一张精确的地图 $\mathcal{M}$，那么定位就是回答``在当前观测 $\bm{z}$ 下，我在哪''，即求后验 $p(\bm{x}\mid \bm{z},\mathcal{M})$。这正是第~\ref{ch:10-global-loc}~章的粒子滤波定位在做的事：地图是已知的，状态只有一个位姿。

反过来，假设有人告诉机器人它每一时刻的精确位姿 $\bm{x}_0,\bm{x}_1,\dots,\bm{x}_N$，那么建图就是把这些位姿下的观测``贴''到地图上做融合，这正是第~\ref{ch:14-mapping}~章要讲的占据栅格、TSDF 等内容：位姿是已知的，地图是待求量。

SLAM 的处境是两者都没有。机器人既不知道位姿，也没有地图，只能拿着一串观测 $\bm{Z}=\{\bm{z}_1,\dots,\bm{z}_M\}$ 和一串输入 $\bm{U}=\{\bm{u}_1,\dots,\bm{u}_N\}$（轮速、IMU 读数等），指望同时把它们都估出来。这就是著名的``鸡生蛋''困境：定位需要地图，建图需要位姿。

\begin{definition}[SLAM]
同时定位与建图（Simultaneous Localization and Mapping，SLAM）是指：机器人在未知环境中运动时，仅利用自身携带的传感器观测（可包含里程计等 proprioceptive 输入），\textbf{在线地、同时地}估计自身位姿轨迹并构建环境地图的问题。
\end{definition}

把它写成联合后验的形式，就是本章所有内容的出发点：
\begin{equation}
p\big(X,\mathcal{M}\mid \bm{Z},\bm{U}\big),
\qquad X=\{\bm{x}_0,\bm{x}_1,\dots,\bm{x}_N\},\quad \mathcal{M}=\text{地图}.
\label{eq:11-joint}
\end{equation}
式~\eqref{eq:11-joint} 中 $X$ 是待估轨迹，$\mathcal{M}$ 是待估地图，条件是全部观测 $\bm{Z}$ 与输入 $\bm{U}$。这个式子看着简单，却藏着两个根本困难。

\textbf{困难一：维度爆炸。}一张 $100\ \mathrm{m}\times 100\ \mathrm{m}$、分辨率 $5\ \mathrm{cm}$ 的二维栅格地图有 $2000\times 2000=4\times 10^{6}$ 个格子，每个格子一个未知量；再乘上上千个位姿，联合状态的维数是百万量级。想直接套用卡尔曼滤波把这么大的协方差矩阵求逆，内存与计算都无法接受。

\textbf{困难二：数据关联。}机器人并不知道``这一帧点云里的第 37 个点''与``上一帧里的哪个点''观测的是同一个物理表面。连''哪一帧和哪一帧到过同一个地方''都不知道。这个''配对''过程一旦出错，就会把一个错误的约束喂给优化器，让整张地图扭曲。因此我们先把这件事定义清楚。

\begin{definition}[数据关联]
数据关联（data association）是指在两次或多次传感器观测之间，确定``哪些观测量对应同一个物理实体''的过程。激光 SLAM 中它表现为点与点、点与面、帧与子图之间的对应；视觉 SLAM 中它表现为特征点匹配；在整个系统层面，它表现为``当前帧是否回到了历史某处''，即回环检测。
\end{definition}

数据关联可以分为三个层次，后面各章会依次展开：\textbf{帧间关联}（当前帧与上一帧，见第~\ref{ch:12-frontend}~章）、\textbf{帧到地图关联}（当前帧与局部地图，见第~\ref{ch:12-frontend}~章）、\textbf{全局关联}（当前帧与历史所有子图，即回环检测，见第~\ref{ch:12-frontend}~章与第~\ref{ch:15-cartographer}~章）。

\begin{keybox}{一句话记住 SLAM 的两张皮}
\textbf{数据关联}把传感器的原始数据变成``谁和谁有关系''；\textbf{状态估计}把这些关系变成``一致的轨迹和地图''。SLAM 系统的一切复杂度，都来自这两件事互相依赖：关联错了，估计就歪；估计歪了，下一帧的关联更容易错。
\end{keybox}

\section{从最大后验到最小二乘}\label{sec:11-map}

上面说的是``难在哪''，这一节说''怎么写成可计算的式子''。SLAM 的标准做法是把式~\eqref{eq:11-joint} 当成一个最大后验估计问题。

\begin{definition}[最大后验估计]
最大后验估计（Maximum A Posteriori，MAP）是指：在已知观测 $\bm{z}$ 的条件下，取使后验概率密度最大的状态作为估计值，即 $\hat{\bm{x}}=\argmax_{\bm{x}} p(\bm{x}\mid\bm{z})$。它与最大似然估计（ML，$\argmax p(\bm{z}\mid\bm{x})$）的区别是：MAP 额外乘上了先验 $p(\bm{x})$，因此可以融合``观测之外的已有知识''。
\end{definition}

由贝叶斯公式（定义见第~\ref{ch:02-math}~章）有
\begin{align}
p(X,\mathcal{M}\mid \bm{Z},\bm{U})
  &= \frac{p(\bm{Z}\mid X,\mathcal{M})\,p(X,\mathcal{M}\mid\bm{U})}{p(\bm{Z}\mid\bm{U})}.
\label{eq:11-bayes-posterior}
\end{align}
分母 $p(\bm{Z}\mid\bm{U})$ 与待估量 $X,\mathcal{M}$ 无关，在对 $X,\mathcal{M}$ 取 $\argmax$ 时可以丢掉，于是
\begin{equation}
p(X,\mathcal{M}\mid \bm{Z},\bm{U}) \;\propto\; \underbrace{p(\bm{Z}\mid X,\mathcal{M})}_{\text{似然}}\;\underbrace{p(X,\mathcal{M}\mid\bm{U})}_{\text{先验}}.
\label{eq:11-propto}
\end{equation}

接下来对式~\eqref{eq:11-propto} 做两个标准假设，把联合概率拆成一个个小因子——这一步是把``一个巨大的概率''变成``一堆可分别计算的残差''的关键。

\textbf{假设 A（观测条件独立）：}给定位姿与地图后，各次观测互相独立，
\begin{equation}
p(\bm{Z}\mid X,\mathcal{M})=\prod_{j=1}^{M}p\big(\bm{z}_j\mid \bm{x}_{i(j)},\mathcal{M}\big),
\end{equation}
其中 $\bm{x}_{i(j)}$ 表示第 $j$ 个观测所涉及的那个位姿。直观理解：传感器每一帧的噪声不会``串''到下一帧。

\textbf{假设 B（一阶马尔可夫性）：}位姿轨迹的生成只依赖上一时刻，
\begin{equation}
p(X,\mathcal{M}\mid\bm{U})=p(\bm{x}_0)\,p(\mathcal{M})\prod_{i=1}^{N}p\big(\bm{x}_i\mid\bm{x}_{i-1},\bm{u}_i\big).
\end{equation}
这相当于说：知道了上一帧位姿和这一帧的里程计读数，这一帧位姿就与更早的历史无关。

把 A、B 代入式~\eqref{eq:11-propto}，得到
\begin{equation}
p(X,\mathcal{M}\mid \bm{Z},\bm{U}) \;\propto\;
p(\bm{x}_0)\,p(\mathcal{M})\prod_{i=1}^{N}p(\bm{x}_i\mid\bm{x}_{i-1},\bm{u}_i)
\prod_{j=1}^{M}p(\bm{z}_j\mid\bm{x}_{i(j)},\mathcal{M}).
\label{eq:11-factorized}
\end{equation}

式~\eqref{eq:11-factorized} 是一个连乘。连乘不好优化，但取对数就变成连加，而 $\log$ 是单调递增函数，不会改变极值点的位置：
\begin{align}
X^{\star},\mathcal{M}^{\star}
  &= \argmax_{X,\mathcal{M}} p(X,\mathcal{M}\mid\bm{Z},\bm{U}) \notag\\
  &= \argmax_{X,\mathcal{M}} \log p(X,\mathcal{M}\mid\bm{Z},\bm{U}) \notag\\
  &= \argmax_{X,\mathcal{M}} \Big\{ \log p(\bm{x}_0)+\log p(\mathcal{M})
     +\sum_{i=1}^{N}\log p(\bm{x}_i\mid\bm{x}_{i-1},\bm{u}_i)
     +\sum_{j=1}^{M}\log p(\bm{z}_j\mid\bm{x}_{i(j)},\mathcal{M}) \Big\}.
\label{eq:11-log-posterior}
\end{align}
注意 $\argmax f = \argmin(-f)$，把上式整体加负号、再改写成最小化，就得到``负对数后验''。这一步是整个 SLAM 后端（以及所有状态估计）的公共起点。

现在代入高斯假设。若某次测量满足 $\bm{z}=\bm{h}(\bm{x})+\bm{n}$，噪声 $\bm{n}\sim\mathcal{N}(\bm{0},\bm{\Sigma})$（高斯的定义见第~\ref{ch:02-math}~章），则
\begin{equation}
p(\bm{z}\mid\bm{x})=\frac{1}{\sqrt{(2\pi)^k\det\bm{\Sigma}}}\exp\!\Big(-\frac12\big(\bm{z}-\bm{h}(\bm{x})\big)^{\top}\bm{\Sigma}^{-1}\big(\bm{z}-\bm{h}(\bm{x})\big)\Big).
\end{equation}
两边取负对数，记残差 $\bm{r}(\bm{x})=\bm{z}-\bm{h}(\bm{x})$：
\begin{align}
-\log p(\bm{z}\mid\bm{x})
  &= \frac12\,\bm{r}^{\top}\bm{\Sigma}^{-1}\bm{r}
   + \frac12\log\det\bm{\Sigma} + \frac{k}{2}\log 2\pi .
\label{eq:11-neglog}
\end{align}
后两项与 $\bm{x}$ 无关，优化时可以扔掉。于是定义\textbf{马氏范数}（Mahalanobis norm，其几何含义见第~\ref{ch:09-fusion}~章）
\begin{equation}
\norm{\bm{r}}_{\bm{\Sigma}^{-1}}^{2} \;\triangleq\; \bm{r}^{\top}\bm{\Sigma}^{-1}\bm{r},
\end{equation}
它相当于先把残差除以``该方向上的标准差''再做平方和：$\bm{\Sigma}$ 大（这个方向测不准）的方向权重小，$\bm{\Sigma}$ 小（测不准的程度低）的方向权重大。

把式~\eqref{eq:11-neglog} 代回式~\eqref{eq:11-log-posterior}，并记每个因子对应一个残差 $\bm{r}_k$，最终得到 SLAM 的标准形式：
\begin{equation}
X^{\star} \;=\; \argmin_{X}\;\sum_{k=1}^{K}\frac12\,\norm{\bm{r}_k(X)}_{\bm{\Sigma}_k^{-1}}^{2}.
\label{eq:11-ls}
\end{equation}

式~\eqref{eq:11-ls} 中每一类残差的物理含义如下（$K=N+M+1$ 个因子）：
\begin{itemize}
  \item \textbf{先验项} $\norm{\bm{r}_0}_{\bm{\Sigma}_0^{-1}}^2$：$-\log p(\bm{x}_0)$，来自起点先验或边缘化得到的信息；
  \item \textbf{运动项} $\norm{\bm{r}_i}_{\bm{\Sigma}_i^{-1}}^2$：$-\log p(\bm{x}_i\mid\bm{x}_{i-1},\bm{u}_i)$，即里程计/IMU 预积分约束，量化``相邻两帧应该差多少''；
  \item \textbf{观测项} $\norm{\bm{r}_j}_{\bm{\Sigma}_j^{-1}}^2$：$-\log p(\bm{z}_j\mid\bm{x}_{i(j)},\mathcal{M})$，即扫描匹配、重投影、点到面等几何约束，量化''当前位姿下，地图上应该看到什么''。
\end{itemize}

\begin{example}[一维两帧：最小二乘就是加权平均]
设只有两个位姿 $x_0,x_1$（一维）。先验说 $x_0=0$，标准差 $\sigma_0=0.1$；里程计说 $x_1-x_0=1$，$\sigma_u=0.2$；观测说 $x_1=0.9$，$\sigma_z=0.1$。信息矩阵就是 $1/\sigma^2$，于是
\begin{equation}
J(x_0,x_1)=\frac{(x_0-0)^2}{2\times 0.01}+\frac{(x_1-x_0-1)^2}{2\times 0.04}+\frac{(x_1-0.9)^2}{2\times 0.01}.
\end{equation}
分别对 $x_0,x_1$ 求偏导并令其为零：
\begin{align}
\frac{\partial J}{\partial x_0} &= \frac{x_0}{0.01}-\frac{x_1-x_0-1}{0.04}=0
  \;\Longrightarrow\; 5x_0-x_1+1=0,\\
\frac{\partial J}{\partial x_1} &= \frac{x_1-x_0-1}{0.04}+\frac{x_1-0.9}{0.01}=0
  \;\Longrightarrow\; 5x_1-x_0-4.6=0.
\end{align}
联立解得 $x_0=-1/60\approx-0.017$，$x_1=11/12\approx0.917$。可以看到：观测的 $\sigma$ 更小（更可信），所以 $x_1$ 停在 $0.9$ 附近而不是里程计说的 $1.0$；同时 $x_0$ 被微微拉离 $0$，以缓和``里程计说走了 1 米''与``观测说在 0.9 米''的矛盾。这就是最小二乘全部的工作方式——\textbf{在互相矛盾的约束之间按可信度做折中}。
\end{example}

\begin{keybox}{本部分的地基}
式~\eqref{eq:11-ls} 说明 SLAM 在数学上就是\textbf{一个大规模稀疏非线性最小二乘问题}。``大规模''来自百万级地图与上千位姿，``稀疏''来自每个因子只涉及极少数变量。第~\ref{ch:13-backend}~章会专门讲怎么高效地解它。但解它之前，得先把残差造出来——造残差靠数据关联，这就是前端。
\end{keybox}

\section{前端与后端：职责划分}\label{sec:11-fe-be}

既然 SLAM $=$ 数据关联 $+$ 状态估计，一个自然的工程划分就出现了：把``造残差''和``解残差''拆成两个模块。这就是前端与后端。

\begin{definition}[前端（Front-end）]
前端是 SLAM 系统中直接处理传感器原始数据、完成数据关联与局部位姿估计的部分。它的输出不是最终地图或最终轨迹，而是一组\textbf{相对位姿约束}（以及关键帧、回环候选），供后端使用。
\end{definition}

\begin{definition}[后端（Back-end）]
后端是 SLAM 系统中接收前端给出的因子/约束、对整条轨迹（以及地图状态）做全局优化与一致性维护的部分。它的输出是全局一致的位姿轨迹、状态协方差以及更新后的地图，供定位、规划与显示使用。
\end{definition}

两者为什么必须分开？因为它们的\textbf{目标、节奏和失败模式完全不同}，见表~\ref{tab:11-fe-be}。

\begin{table}[htbp]
  \centering
  \caption{SLAM 前端与后端的职责对比}\label{tab:11-fe-be}
  \small
  \begin{tabular}{p{0.14\linewidth}p{0.38\linewidth}p{0.38\linewidth}}
    \toprule
    维度 & 前端（Front-end） & 后端（Back-end） \\
    \midrule
    一句话 & 负责\emph{局部一致} & 负责\emph{全局一致} \\
    输入 & 传感器原始数据（scan、图像、IMU、轮速） & 前端给出的因子/约束（相对位姿、回环边） \\
    输出 & 相对位姿、关键帧、回环候选 & 全局位姿、协方差、地图更新量 \\
    运行频率 & 高（$10\sim100$ Hz，与传感器同量级） & 低（$0.1\sim5$ Hz，或事件触发） \\
    关键算法 & 预处理、扫描匹配、数据关联、外点剔除 & 因子图/位姿图优化、鲁棒核、边缘化、稀疏求解 \\
    失败表现 & 漂移、跟踪丢失、关联错乱 & 全局变形、优化不收敛、回环错闭 \\
    能否替代 & 可整体换掉（LOAM 换 Cartographer 前端） & 相对稳定（g2o/GTSAM/Ceres 通用） \\
    \bottomrule
  \end{tabular}
\end{table}

表~\ref{tab:11-fe-be} 里最关键的一行是``一句话''。把它展开成一句可背诵的话：\textbf{前端把数据变成约束，后端把约束变成一致的轨迹与地图；前端可以错，但不能慢；后端可以慢，但不能错。}

所谓``全局一致''是什么意思？设想机器人绕一个矩形走廊走了一圈回到起点。由于里程计的累积误差，前端给出的轨迹是一条\emph{不闭合}的螺旋线：最后一帧的位姿和第一帧差了两米。后端通过回环检测发现``这两帧其实是同一个地方''，于是把这个``差两米''的信息作为一条强约束加进去，重新优化全部轨迹，让首尾闭合。优化之后，整条轨迹上的每一帧位姿都会被微调——这就是后端的作用范围：它改的是整条轨迹，而前端只改当前帧。

\section{三种语境下的``前端/后端''（重点）}\label{sec:11-three-contexts}

``前端/后端''这个词在中文技术圈用得极多，但在三种语境下指的东西完全不同。新手最容易犯的错是：看完 SLAM 的前后端划分，转头去看定位系统或导航系统的框图，发现``前端''指的不是同一个东西，于是怀疑自己没学懂。其实不是你错，是术语重载了。本节把三种语境并排摆出来（表~\ref{tab:11-contexts}），并给出判断规则。

\begin{table}[htbp]
  \centering
  \caption{三种语境下``前端/后端''的所指}\label{tab:11-contexts}
  \footnotesize
  \begin{tabular}{p{0.15\linewidth}p{0.27\linewidth}p{0.27\linewidth}p{0.19\linewidth}}
    \toprule
    语境 & ``前端''指什么 & ``后端''指什么 & 本书在哪讲 \\
    \midrule
    SLAM（建图） & 里程计 + 数据关联：预处理、帧间/帧到地图匹配、关键帧选取 & 全局优化 + 回环：因子图/位姿图优化、地图更新 & 第~\ref{ch:11-slam-framework}--\ref{ch:15-cartographer}~章 \\
    定位系统 & 实时状态估计：LIO/VIO、ESKF，输出高频位姿与协方差 & 先验地图配准 + 因子图更新，低频修正漂移并判断可信度 & 第~\ref{ch:23-zlocalization}~章 \\
    导航系统 & 感知：定位 + 建图 + 障碍检测，回答``我在哪、周围是什么'' & 决策与控制：行为树、规划器、控制器，回答``我该往哪走'' & 第~\ref{ch:18-nav-overview}~章 \\
    \bottomrule
  \end{tabular}
\end{table}

\textbf{语境一：SLAM 意义下。}这是本章与第~\ref{ch:12-frontend}--\ref{ch:15-cartographer}~章默认使用的划分。前端 $=$ 里程计 $+$ 数据关联，后端 $=$ 优化 $+$ 回环。划分依据是\textbf{数学结构}：前端产出因子，后端求解式~\eqref{eq:11-ls}。

\textbf{语境二：定位系统意义下。}一个定位系统（比如 RoboMaster 赛场上跑的 Z-Localization）并不建图，它面对的是一张\textbf{已有的}先验地图。此时``前端''是保持高频、连续的实时状态估计（LIO 或 ESKF），``后端''是与先验地图做配准、并用因子图把``先验地图约束''和``里程计约束''融合起来的那一层。划分依据是\textbf{信息的来源}：前端只用自己的传感器，后端引入了先验地图这个外部信息源。

\textbf{语境三：导航系统意义下。}在 Nav2 这类导航框架里，``前端/后端''是按\textbf{信息流方向}划分的：定位与建图属于``感知前端''，回答``世界现在是什么样''；规划与控制属于``决策后端''，回答``接下来怎么做''。划分依据是\textbf{抽象层次}：前者输出世界状态，后者输出动作。

\begin{keybox}{判断规则：先问信息链的终点}
看到``前端/后端''，先别急着套定义，先问一句：\textbf{这条信息链的终点是什么？}终点是``地图与轨迹''，就是语境一；终点是``当前位姿''，就是语境二；终点是``速度指令''，就是语境三。三种语境并不冲突：同一个系统里可以同时存在三对前后端——Z-Localization 内部有语境二的一对，它又把位姿喂给 Nav2，成为语境三中``感知前端''的一部分。
\end{keybox}

\begin{note}
本书的默认约定：第~\ref{ch:11-slam-framework}~章到第~\ref{ch:15-cartographer}~章，以及第~\ref{ch:16-lio-evolution}~章、第~\ref{ch:17-advanced-lio}~章，``前端/后端''一律指语境一；第~\ref{ch:23-zlocalization}~章按语境二讲；第~\ref{ch:18-nav-overview}~章按语境三讲。写作与阅读时请以各章开头的一句话声明为准。
\end{note}

\section{数据流：一帧数据在系统里走过的路}\label{sec:11-dataflow}

把上面的职责划分画成图，就是图~\ref{fig:11-dataflow}。

\begin{figure}[htbp]
  \centering
  \begin{tikzpicture}[font=\small]
    \node[block]  (s)  at (0,0)      {传感器};
    \node[block2] (fe) at (3.8,0)    {前端\\{\footnotesize 匹配 + 数据关联}};
    \node[block]  (kf) at (7.6,0)    {关键帧\\{\footnotesize 相对位姿约束}};
    \node[block2] (be) at (11.4,0)   {后端\\{\footnotesize 全局优化}};
    \node[block]  (lc) at (7.6,2.1)  {回环检测};
    \node[block]  (out) at (11.4,-2.1) {全局位姿\\{\footnotesize TF / 地图}};
    \node[block]  (lm) at (3.8,-2.1) {局部地图};
    \draw[arrow] (s)  -- node[label,above]{原始数据} (fe);
    \draw[arrow] (fe) -- node[label,above]{关键帧} (kf);
    \draw[arrow] (kf) -- node[label,above]{约束} (be);
    \draw[arrow] (be) -- (out);
    \draw[arrow] (out) -- node[label,below]{地图更新} (lm);
    \draw[arrow] (lm) -- node[label,left,align=center]{局部\\地图} (fe);
    \draw[arrow, dashed] (kf.north) -- (lc.south);
    \draw[arrow, dashed] (lc.east) -| node[label,above,pos=0.35]{回环约束} (be.north);
  \end{tikzpicture}
  \caption{SLAM 系统的总体数据流。实线为在线主通路：传感器 $\rightarrow$ 前端 $\rightarrow$ 关键帧与约束 $\rightarrow$ 后端 $\rightarrow$ 全局位姿与地图；下方是``地图回喂前端''的闭环。虚线为回环检测这条\emph{全局数据关联}支路：它由前端发起、向后端提供回环约束。正是这条支路把``局部一致''升级为``全局一致''。}\label{fig:11-dataflow}
\end{figure}

沿着图~\ref{fig:11-dataflow} 把各个箭头上流动的数据说清楚，读源码时就能对上号：
\begin{enumerate}
  \item \textbf{传感器 $\rightarrow$ 前端}：原始 scan/图像/IMU，带时间戳，通常先做时间同步（第~\ref{ch:03-sensors}~章）。
  \item \textbf{前端 $\rightarrow$ 关键帧}：每帧的相对位姿估计及其协方差，以及``这一帧够不够格当关键帧''的判断（第~\ref{ch:12-frontend}~章）。
  \item \textbf{关键帧 $\rightarrow$ 后端}：一条条因子。最典型的是位姿图里的二元边（第~\ref{ch:13-backend}~章）。
  \item \textbf{后端 $\rightarrow$ 全局位姿}：优化后的轨迹，发布成 TF 变换与地图坐标系（第~\ref{ch:25-ros2}~章会讲 TF 树）。
  \item \textbf{全局位姿 / 地图 $\rightarrow$ 局部地图 $\rightarrow$ 前端}：前端做帧到地图匹配时需要一块``附近的地图''，这块地图由后端维护的全局地图裁剪而来。
  \item \textbf{回环检测}：注意它在图里是``从前端出发''的。回环检测本质是\emph{全局范围的数据关联}，属于前端职责；它只负责``找出候选 + 校验''，把约束交给后端去优化。这是第~\ref{ch:12-frontend}~章要专门解释的一个常见误解。
\end{enumerate}

\section{SLAM 系统的分类}\label{sec:11-taxonomy}

不同文献里的 SLAM 系统五花八门，其实可以从五个维度切开。每一项的取舍都对应着具体的工程权衡。

\begin{enumerate}
  \item \textbf{滤波 vs 优化。}滤波类（EKF-SLAM、FastSLAM）把状态和协方差一路递推，优点是天然增量、能做在线预测，缺点是协方差矩阵随状态维数平方增长，且线性化点一旦固定就无法回头修正。优化类（因子图、位姿图）把历史状态全部保留，用迭代最小二乘求解，优点是精度高、能全局重线性化，缺点是计算量随规模增长（靠稀疏性与增量求解缓解）。\textbf{现在主流是优化类}，原因见 \S\ref{sec:11-why-opt}。
  \item \textbf{直接法 vs 特征法。}直接法直接用像素灰度/点云强度构造残差（如 DSO\cite{engel2018dso}），省去特征提取，在弱纹理场景更稳，但对光度变化与初值敏感；特征法先提取点/线/面特征再匹配（如 ORB-SLAM3\cite{campos2021orbslam3}），对光照鲁棒、容易做回环，但弱纹理处会丢特征。激光 SLAM 介于两者之间：``点到面''用的是几何结构而非人工特征。
  \item \textbf{2D vs 3D。}2D 假设机器人在地面运动、地图是一层栅格，计算量小、结构简单，适合单层场地（第~\ref{ch:15-cartographer}~章）；3D 要恢复完整六自由度与三维地图，计算与内存大一个量级，但能处理坡道、多层、空中（第~\ref{ch:16-lio-evolution}~章与第~\ref{ch:17-advanced-lio}~章）。
  \item \textbf{单机 vs 多机。}多机 SLAM 要让多个机器人共享地图并保持坐标系一致，核心难点从``估计''变成``通信与一致性''。
  \item \textbf{离线 vs 在线。}离线（如用 rosbag 建图）可以用全局优化反复迭代，追求最好结果；在线必须满足实时性，只能``够用就好、逐步逼近''。
\end{enumerate}

\section{为什么现在的主流是``前端 LIO + 后端因子图''}\label{sec:11-why-opt}

把本节前面的内容串起来，就能理解为什么近十年的主流架构几乎都是``前端高频里程计（LIO/VIO）$+$ 后端低频因子图''：LIO-SAM\cite{shan2020liosam}、FAST-LIO2\cite{xu2022fastlio2}、Cartographer\cite{hess2016cartographer} 都是这个骨架。原因有五条。

\begin{enumerate}
  \item \textbf{符合信息流。}前端能拿到的是``相邻两帧之间的关系''，这天然是相对约束；后端要回答的是``全局怎么摆''，天然是全局优化。把两者硬塞进一个滤波器，等于用递推的方式做全局问题，事倍功半。
  \item \textbf{计算量可扩展。}式~\eqref{eq:11-ls} 的海森矩阵是稀疏的：每个因子只连接少数变量。1000 个位姿的位姿图，海森矩阵的非零元只占百分之几，稀疏 Cholesky 分解可以做到毫秒级；而同样规模的 EKF 要维护 $6000\times6000$ 的稠密协方差并每次求逆\cite{kummerle2011g2o}，根本跑不动。
  \item \textbf{精度更高。}优化在后端每轮都重新线性化，而 EKF 只在一步处线性化一次，误差会沿着轨迹累积。因子图的表述方式在数学上等价于``对所有历史观测做联合最大似然''，这正是式~\eqref{eq:11-ls} 想做的事\cite{dellaert2017factorgraphs}。
  \item \textbf{可恢复性。}前端总会漂移，甚至短时跟丢。只要后端还保存着历史因子，一次成功的回环或重定位就能把整条轨迹``拉正''；滤波器模式下历史状态已经被边缘化掉，想改也改不了。
  \item \textbf{工程解耦。}前端可以整体替换（今天用 LOAM，明天换 Cartographer，甚至换成视觉），后端则可以一直用现成的图优化库（如 g2o、GTSAM、Ceres\cite{kummerle2011g2o}），基本不用动。这种接口稳定性对多人协作的算法组极其重要。
\end{enumerate}

代价也要说清楚：后端是滞后的，前端已经漂了几十米，后端才通过回环把误差摊平；如果前端漂移大到连回环都找不到（描述子已经对不上），后端也无能为力。所以``前端准''和``后端稳''是互补的两件事，缺一不可。

\begin{summary}
\begin{enumerate}
  \item SLAM 是``定位需要地图、建图需要位姿''的鸡生蛋问题，形式化为联合后验 $p(X,\mathcal{M}\mid\bm{Z},\bm{U})$（式~\eqref{eq:11-joint}）。
  \item 做两个假设（观测条件独立、一阶马尔可夫）后，联合后验拆成先验、运动、观测三类因子的连乘；取负对数并把噪声建模为高斯，就得到标准最小二乘形式 $\min_X\sum_k\frac12\norm{\bm{r}_k}_{\bm{\Sigma}_k^{-1}}^2$（式~\eqref{eq:11-ls}）。
  \item 马氏范数 $\norm{\bm{r}}^2_{\bm{\Sigma}^{-1}}=\bm{r}^{\top}\bm{\Sigma}^{-1}\bm{r}$ 的含义是``按各方向可信度加权的平方误差''，信息矩阵 $\bm{\Sigma}^{-1}$ 越大该约束越硬。
  \item 前端 $=$ 数据关联 $+$ 局部位姿估计，追求实时与鲁棒，负责\emph{局部一致}；后端 $=$ 全局优化 $+$ 回环，追求精度与可恢复，负责\emph{全局一致}（表~\ref{tab:11-fe-be}）。
  \item ``前端/后端''有三种语境：SLAM 下按数学结构分（造因子/解因子），定位系统下按信息来源分（自身传感器/先验地图），导航系统下按抽象层次分（感知/决策）。判断规则是``先问信息链的终点''（表~\ref{tab:11-contexts}）。
  \item 回环检测属于前端（全局数据关联），它产出约束，由后端优化；把回环检测归到后端是常见误解。
  \item 主流架构``前端 LIO $+$ 后端因子图''的五条理由：符合信息流、计算可扩展、精度更高、可恢复、工程解耦；代价是后端滞后。
\end{enumerate}


\end{summary}

\begin{problemset}
  \item 在式~\eqref{eq:11-factorized} 的两个假设中，如果去掉``观测条件独立''，式子会变成什么样？举一个现实中不独立的例子（提示：同一帧内相邻点的噪声、动态物体）。
  \item 用本节的加权平均方法重算：若观测标准差改为 $\sigma_z=0.3$（比里程计更不可信），$x_0,x_1$ 会偏向哪边？请给出解析解并与前面的 $0.917$ 对比。
  \item 为什么说``回环检测属于前端''？如果把它硬放到后端，会在工程上带来什么麻烦？请从运行频率和数据依赖两个角度回答。
  \item Z-Localization 里``前端''指 LIO、``后端''指先验地图配准 $+$ 因子图。请说明为什么这里的``后端''不需要维护地图，而 SLAM 的后端需要。
  \item 假设你只有一只 2D 激光雷达和轮式里程计，在一个 $30\ \mathrm{m}\times20\ \mathrm{m}$ 的单层场地里建图，你会选滤波式还是优化式？说明你的取舍依据。
\end{problemset}

\begin{reading}
\begin{enumerate}
  \item Thrun, Burgard, Fox，《Probabilistic Robotics》第 10--11 章\cite{thrun2005probabilistic}：SLAM 的概率表述与 EKF-SLAM 的原点，用来对照``滤波 vs 优化''。
  \item Dellaert \& Kaess, ``Factor Graphs for Robot Perception''\cite{dellaert2017factorgraphs}：因子图视角下的 SLAM，读完能理解为什么式~\eqref{eq:11-ls} 可以画成一张图。
  \item Kümmerle 等，g2o 论文\cite{kummerle2011g2o}：把``稀疏最小二乘''落到代码，适合读完后去翻 g2o 的 \lstinline{OptimizableGraph} 与稀疏求解器部分。
  \item LIO-SAM 论文\cite{shan2020liosam}：``前端里程计 + 后端因子图''的教科书级实现，读完再回看本章图~\ref{fig:11-dataflow} 会非常直观。
\end{enumerate}


\end{reading}

% 本章使用 TikZ 的 arrows.meta / positioning 库（main.tex 与 tests 均已加载，此行幂等无害）
\usetikzlibrary{arrows.meta, positioning}

\chapter{SLAM 前端详解}\label{ch:12-frontend}

\chapterintro{第~\ref{ch:11-slam-framework}~章把 SLAM 拆成了``造约束的前端''和``解约束的后端''。本章专门讲前端：它从传感器原始数据出发，经过预处理、数据关联、位姿估计、外点剔除、关键帧选取、局部地图维护，最终给后端交出一串可信的相对位姿约束。前端是整套系统里离传感器最近、也最容易出错的一环——后端的数学再漂亮，前端喂进去一条错误约束，整张地图照样会歪。本章会逐个环节讲清楚``它解决什么问题、怎么做、什么时候会失效''，重点讲透三件事：数据关联为什么是前端的心脏、为什么``帧到局部地图''比``帧到帧''稳、以及回环检测为什么属于前端而不属于后端。}

\section{前端流水线}\label{sec:12-pipeline}

一个典型的激光 SLAM 前端每收到一帧点云，都会走一遍图~\ref{fig:12-pipeline} 的流程。视觉前端结构类似，只是把``点云''换成``图像 + 特征''。

\begin{figure}[htbp]
  \centering
  \begin{tikzpicture}[font=\footnotesize]
    \node[block] (p1) at (0,0)    {预处理};
    \node[block] (p2) at (2.8,0)  {数据关联};
    \node[block] (p3) at (5.6,0)  {位姿估计};
    \node[block] (p4) at (8.4,0)  {外点剔除};
    \node[block] (p5) at (11.2,0) {关键帧判定};
    \node[block, fill=black!5] (out) at (11.2,-1.9) {关键帧 + 约束};
    \node[block] (lm)  at (2.8,-1.9) {局部地图};
    \draw[arrow] (p1) -- (p2);
    \draw[arrow] (p2) -- (p3);
    \draw[arrow] (p3) -- (p4);
    \draw[arrow] (p4) -- (p5);
    \draw[arrow] (p5) -- (out);
    \draw[arrow] (out) -- node[label,below]{加入窗口} (lm);
    \draw[arrow] (lm) -- node[label,left,align=center]{匹配\\目标} (p2);
  \end{tikzpicture}
  \caption{SLAM 前端流水线。主链是``预处理 $\rightarrow$ 数据关联 $\rightarrow$ 位姿估计 $\rightarrow$ 外点剔除 $\rightarrow$ 关键帧判定''；下方闭环表示关键帧被加入局部地图窗口，成为下一帧匹配的目标。}\label{fig:12-pipeline}
\end{figure}

每一步的输入输出见表~\ref{tab:12-pipeline}。读源码时，只要能把一个函数映射到表里的某一行，就基本知道它在干什么。

\begin{table}[htbp]
  \centering
  \caption{前端各环节的输入与输出}\label{tab:12-pipeline}
  \small
  \begin{tabular}{p{0.18\linewidth}p{0.36\linewidth}p{0.36\linewidth}}
    \toprule
    环节 & 输入 & 输出 \\
    \midrule
    预处理 & 原始 scan/图像、IMU、时间戳 & 去畸变点云、降采样点云、地面/动态点标记 \\
    数据关联 & 当前帧 + 局部地图 & 点—点、点—面、点—体素对应关系 \\
    位姿估计 & 对应关系 + 初值 & 当前帧相对位姿 $\hat{T}$ 及协方差 \\
    外点剔除 & 对应关系 + 残差 & 内点集合、内点比例、是否跟踪失败 \\
    关键帧判定 & 位姿、视角变化、观测新区域比例 & 是否为关键帧的布尔量 \\
    \bottomrule
  \end{tabular}
\end{table}

为什么要拆成这么多步？因为每一步的\textbf{失败模式}不同，拆开才能分别监控、分别兜底。比如``内点比例骤降''是关联出了问题的信号，而``海森矩阵最小特征值过小''是几何退化（\S\ref{sec:12-fail}）的信号。把所有步骤揉成一个黑盒，出了问题就无从下手。

\section{数据关联：前端的心脏}\label{sec:12-association}

\begin{definition}[数据关联（工程视角）]
在 SLAM 前端中，数据关联指：给当前帧中的每个观测（点、特征、整帧），在局部地图或历史帧中找到它所对应的那个物理实体，并给出对应关系集合 $\mathcal{C}=\{(i,j)\}$，其中 $i$ 是当前帧观测的编号，$j$ 是地图元素的编号。
\end{definition}

第~\ref{ch:11-slam-framework}~章已经说明关联错了后果严重，这里把它落到具体做法上。最常用的三种手段如下。

\textbf{（1）最近邻（Nearest Neighbor，NN）。}对当前帧的每个点 $\bm{p}_i$，在局部地图里找欧氏距离最近的元素：
\begin{equation}
j^{\star}(i) = \argmin_{j}\ \norm{\bm{p}_i-\bm{q}_j}_2 ,
\qquad \text{若}\ \norm{\bm{p}_i-\bm{q}_{j^{\star}}}_2 > \tau_d\ \text{则丢弃}.
\label{eq:12-nn}
\end{equation}
式中 $\tau_d$ 是距离阈值，通常取``期待的最大位移加上若干倍点间距''。直接用 NN 的问题是对称性：点 $i$ 认为 $j$ 最近，但 $j$ 可能认为别的点更近。

\textbf{（2）相互最近邻（Mutual NN，MNN）。}要求双向成立，即 $j^{\star}(i)=j$ 且 $i^{\star}(j)=i$ 才保留。这在两帧点密度接近时能剔除大量错误对应，代价是内点变少。

\textbf{（3）RANSAC 外点剔除。}先随机抽最少的点对（2D 刚体变换最少 2 对，3D 最少 3 对）解出一个候选变换，统计有多少点对在该变换下残差小于阈值（内点），重复若干次取内点最多的一次。它能容忍高比例的错关联，代价是需要多次抽样。

图~\ref{fig:12-badassoc} 展示了错误关联的破坏力：当机器人从侧面靠近时，由于初值不准，$A$ 墙上的点被错误地关联到 $B$ 墙，优化器为了``满足''这条约束，会把两堵墙拉近，于是在地图上出现一堵``重影墙''。

\begin{figure}[htbp]
  \centering
  \begin{tikzpicture}[font=\footnotesize]
    % legend
    \fill[structurecolor] (0.35,3.55) circle (1.4pt);
    \node[label, anchor=west] at (0.5,3.55) {当前帧观测};
    \fill[red!60!black] (3.6,3.55) circle (1.4pt);
    \node[label, anchor=west] at (3.75,3.55) {局部地图点};
    % left: correct association
    \draw[line width=3pt, black!25] (1.0,0) -- (1.0,3.0);
    \node[label] at (1.0,-0.35) {同一面墙};
    \foreach \y in {0.5,1.5,2.5}{
      \fill[structurecolor] (0.72,\y) circle (1.4pt);
      \fill[red!60!black] (1.28,\y) circle (1.4pt);
      \draw[arrow] (0.72,\y) -- (1.28,\y);
    }
    \node[align=center] at (1.0,-1.05) {正确关联：\\箭头短且彼此一致};
    % right: wrong association
    \begin{scope}[xshift=6.4cm]
    \draw[line width=3pt, black!25] (0,0) -- (0,3.0);
    \draw[line width=3pt, black!25] (2.4,0) -- (2.4,3.0);
    \node[label] at (0,-0.35) {墙 A};
    \node[label] at (2.4,-0.35) {墙 B};
    \fill[structurecolor] (-0.3,0.4) circle (1.4pt);
    \fill[structurecolor] (-0.3,1.5) circle (1.4pt);
    \fill[structurecolor] (-0.3,2.6) circle (1.4pt);
    \fill[red!60!black] (2.7,0.4) circle (1.4pt);
    \fill[red!60!black] (2.7,1.5) circle (1.4pt);
    \fill[red!60!black] (2.7,2.6) circle (1.4pt);
    \draw[arrow, red!60!black] (-0.3,0.4) -- (2.7,2.6);
    \draw[arrow, red!60!black] (-0.3,1.5) -- (2.7,1.5);
    \draw[arrow, red!60!black] (-0.3,2.6) -- (2.7,0.4);
    \node[align=center, red!60!black] at (1.2,-1.05) {错误关联：\\A 墙的点被配到 B 墙};
    \end{scope}
  \end{tikzpicture}
  \caption{左：正确关联——两次观测落在同一面墙上，对应箭头短且彼此一致；右：错误关联——初值偏差使 A 墙的观测点被配到 B 墙，对应箭头跨越两墙、互相交叉。后者会在优化中产生巨大残差，迫使轨迹或地图扭曲，形成``重影墙''。}\label{fig:12-badassoc}
\end{figure}

\begin{keybox}{前端的第一原则}
前端的第一原则不是``算得准''，而是\textbf{``别把错误关联喂给后端''}。一条错误约束的破坏力，远大于十条正确约束的收益：正确约束只是让结果更精确一点，错误约束足以让整段轨迹塌陷。所以外点剔除、鲁棒核（第~\ref{ch:13-backend}~章）、回环几何校验这些``保险''，在工程上比提升匹配精度更重要。
\end{keybox}

\section{从``帧到帧''升级到``帧到局部地图''}\label{sec:12-ftm}

第~\ref{ch:07-lidar-odom}~章讲的扫描匹配是帧到帧（frame-to-frame）：把当前帧和上一帧配准，得到两帧之间的相对位姿。这在里程计里够用，但在 SLAM 前端里不够，原因是误差会沿着链式匹配累积。

假设每一帧匹配都有独立、零均值、标准差为 $\sigma$ 的误差。那么第 $k$ 帧相对第 0 帧的误差是 $k$ 个误差之和。由方差的可加性（独立随机变量之和的方差等于方差之和，见第~\ref{ch:02-math}~章）：
\begin{equation}
\Var\Big(\sum_{i=1}^{k}\epsilon_i\Big) = \sum_{i=1}^{k}\Var(\epsilon_i) = k\sigma^{2}
\quad\Longrightarrow\quad
\sigma_k = \sqrt{k}\,\sigma .
\label{eq:12-drift}
\end{equation}
$k=100$ 时标准差已经放大 10 倍。这就是漂移的数学本质：误差按``随机游走''的方式增长。

帧到局部地图（frame-to-map）改动了一个关键设定：\textbf{匹配目标不再是单帧，而是多帧融合成的局部地图}。于是问题变成``当前帧在该地图坐标系下处于什么位姿''，误差不沿着链传递。

\begin{definition}[局部地图]
局部地图（local map）是以当前位姿为中心、由最近若干个关键帧的观测融合而成的、用于当前帧匹配的地图子集。它只覆盖机器人附近的一小块区域，因而可以实时查询与更新。
\end{definition}

局部地图的工程组织方式有三种，各有取舍：
\begin{itemize}
  \item \textbf{滑动窗口}：保留最近 $N$ 个关键帧的原始点云，匹配时全部拼起来。实现最简单，内存随 $N$ 线性增长。
  \item \textbf{体素地图}：把点云按体素栅格（\S\ref{sec:12-preprocess}）降采样后长期保存，查询时取当前位姿邻域的体素。内存可控，适合大场景。
  \item \textbf{子图（submap）}：把连续若干帧融合成一张概率栅格（第~\ref{ch:15-cartographer}~章），匹配时直接对栅格做插值查询。Cartographer 用的就是它。
\end{itemize}

为什么帧到地图更稳？有三个原因。第一，\textbf{目标更``干净''}：多帧融合相当于对同一个表面做了平均，若 $n$ 次独立观测的噪声标准差为 $\sigma$，平均后降到 $\sigma/\sqrt{n}$。第二，\textbf{约束更多}：一帧只和上一帧比，只能约束两帧之间的相对关系；而帧到地图同时受地图中所有可见表面的约束。第三，\textbf{误差不进入目标}：帧到帧时，上一帧本身带着误差，等于用``错的东西''当基准；帧到地图时，基准是后端优化过的全局一致地图（或至少是融合后的局部地图），误差被摊平了。

\begin{note}
\textbf{帧到帧并非一无是处。}它不需要维护地图、计算量小，非常适合做``里程计''（第~\ref{ch:07-lidar-odom}~章），也常被用作帧到地图匹配的初值。实际系统里两者往往同时存在：先用帧到帧（或 IMU 预测）给出初值，再用帧到地图精化。
\end{note}

\section{关键帧策略}\label{sec:12-keyframe}

\begin{definition}[关键帧]
关键帧（keyframe）是从连续帧序列中挑选出的、代表局部观测的子集帧。只有关键帧参与地图维护与后端约束构建；非关键帧只用于帧间追踪，不进入优化。
\end{definition}

为什么要有关键帧？因为相邻帧之间信息高度冗余：机器人以 $0.5\ \mathrm{m/s}$ 运动、雷达 $10\ \mathrm{Hz}$ 时，相邻两帧只差 $5\ \mathrm{cm}$，把它们都塞进后端，节点数会爆炸而信息量几乎不增加。

关键帧的触发条件通常有四类，工程上取它们的``或''：
\begin{enumerate}
  \item \textbf{平移足够大}：$\norm{\bm{t}_k-\bm{t}_{\mathrm{last}}}_2 > d_{\max}$，典型 $d_{\max}=0.5\sim1.0\ \mathrm{m}$；
  \item \textbf{旋转足够大}：$\abs{\Delta\theta_k} > \theta_{\max}$，典型 $\theta_{\max}=10^{\circ}\sim20^{\circ}$（3D 时要分别看 roll/pitch/yaw 或转角轴）；
  \item \textbf{时间足够长}：距上一关键帧超过 $t_{\max}$，典型 $0.5\sim2\ \mathrm{s}$，防止静止时长时间没有关键帧、导致地图更新停滞；
  \item \textbf{信息量足够大}：当前帧观测到的``新区域''（从未见过的栅格/体素）占比超过阈值，或匹配协方差的迹（trace，见第~\ref{ch:02-math}~章）超过阈值。这条最接近本质，但实现最复杂。
\end{enumerate}

合起来写成一个判定式：
\begin{equation}
\text{新关键帧} \;\Longleftrightarrow\;
\Big(\norm{\bm{t}_k-\bm{t}_{\mathrm{last}}}_2>d_{\max}\Big)\ \vee\
\Big(\abs{\Delta\theta_k}>\theta_{\max}\Big)\ \vee\
\Big(t_k-t_{\mathrm{last}}>t_{\max}\Big)\ \vee\
\Big(\text{新区域占比}>\eta\Big).
\label{eq:12-kf}
\end{equation}
式~\eqref{eq:12-kf} 中四个条件互相独立，工程上取``或''是为了在任意一种情况下都能及时更新地图：走得快时靠平移条件，原地转身时靠旋转条件，静止时靠时间条件，进入新房间时靠信息量条件。

关键帧太多或太少都有问题，见表~\ref{tab:12-kf}。取值的``甜点''取决于场景：室内小场地 $0.5\ \mathrm{m}$ 一帧比较合适，大场地可以放宽到 $1\sim2\ \mathrm{m}$。

\begin{table}[htbp]
  \centering
  \caption{关键帧数量对系统的影响}\label{tab:12-kf}
  \small
  \begin{tabular}{p{0.16\linewidth}p{0.36\linewidth}p{0.36\linewidth}}
    \toprule
    情形 & 好处 & 代价 \\
    \midrule
    关键帧过多 & 地图细节丰富、约束密集、回环易检测 & 后端节点多、优化变慢、内存与存储膨胀 \\
    关键帧过少 & 后端轻量、优化快 & 约束稀疏、回环描述子区分度低、地图轮廓破碎 \\
    \bottomrule
  \end{tabular}
\end{table}

\section{点云预处理}\label{sec:12-preprocess}

预处理决定了``喂给匹配的数据质量''。这里讲三个最常用、也最影响效果的步骤。

\subsection{体素降采样}

原始雷达一帧可能有十万个点，直接匹配太慢。体素降采样（voxel grid downsampling）把空间切成边长为 $r$ 的小立方体（体素），每个体素只保留一个代表点。体素索引由坐标直接算出：
\begin{equation}
(i,j,k)=\left(\left\lfloor\frac{x-x_{\min}}{r}\right\rfloor,\;
\left\lfloor\frac{y-y_{\min}}{r}\right\rfloor,\;
\left\lfloor\frac{z-z_{\min}}{r}\right\rfloor\right),
\label{eq:12-voxel-index}
\end{equation}
其中 $\lfloor\cdot\rfloor$ 是向下取整，$x_{\min},y_{\min},z_{\min}$ 是点云包围盒的最小角。同一体素内的点被划分为一组。

代表点取体素质心，即
\begin{equation}
\bar{\bm{p}}_{ijk}=\frac{1}{n_{ijk}}\sum_{m=1}^{n_{ijk}}\bm{p}_m ,
\label{eq:12-voxel-centroid}
\end{equation}
其中 $n_{ijk}$ 是该体素内的点数。为什么取质心而不是随便取一个点？因为质心是这个体素内所有点的最小二乘代表：我们要找一个点 $\bm{q}$，使得到体内所有点的距离平方和最小，
\begin{equation}
J(\bm{q})=\sum_{m=1}^{n}\norm{\bm{q}-\bm{p}_m}_2^{2}.
\end{equation}
对 $\bm{q}$ 求导并令其为零：
\begin{equation}
\frac{\partial J}{\partial \bm{q}}=2\sum_{m=1}^{n}\big(\bm{q}-\bm{p}_m\big)=\bm{0}
\;\Longrightarrow\;
n\bm{q}=\sum_{m=1}^{n}\bm{p}_m
\;\Longrightarrow\;
\bm{q}=\frac{1}{n}\sum_{m=1}^{n}\bm{p}_m .
\end{equation}
可见式~\eqref{eq:12-voxel-centroid} 不是``约定俗成''，而是``误差平方和最小''的直接结果。分辨率 $r$ 的取法：太小则降采样无效，太大则丢失结构，经验值是略大于点间距（如 $0.1\sim0.2\ \mathrm{m}$）。

\subsection{地面分割（RANSAC 平面拟合）}

地面点占一帧点云的很大比例，且对匹配帮助不大（平坦地面在任何水平平移下都能配上，反而会``吸住''错误的位姿），所以先把它分出来。最常用的方法是 RANSAC 拟合平面。

平面用单位法向 $\bm{n}$ 和截距 $d$ 表示：$\bm{n}^{\top}\bm{p}=d$，$\norm{\bm{n}}=1$。取三个不共线的点 $\bm{p}_1,\bm{p}_2,\bm{p}_3$，法向由叉积给出：
\begin{equation}
\bm{n}=\frac{(\bm{p}_2-\bm{p}_1)\times(\bm{p}_3-\bm{p}_1)}
{\norm{(\bm{p}_2-\bm{p}_1)\times(\bm{p}_3-\bm{p}_1)}},
\qquad d=\bm{n}^{\top}\bm{p}_1 .
\label{eq:12-plane}
\end{equation}
点到该平面的距离是 $\abs{\bm{n}^{\top}\bm{p}-d}$（因为 $\bm{n}$ 是单位向量，投影长度就是距离）。距离小于阈值 $\tau_p$ 的点称为内点。重复抽样 $k$ 次，取内点最多的平面作为地面。

抽多少次才够？设内点占全体的比例为 $w$，一次抽样三个点都落在内点的概率是 $w^{3}$（近似认为独立），则 $k$ 次中都失败的概率为 $(1-w^{3})^{k}$，至少成功一次的概率为
\begin{equation}
P_{\mathrm{succ}}=1-(1-w^{3})^{k}.
\end{equation}
要求 $P_{\mathrm{succ}}\ge P_{0}$，解出
\begin{equation}
k\;\ge\;\frac{\log(1-P_{0})}{\log(1-w^{3})}.
\label{eq:12-ransac-k}
\end{equation}
代入数字：若地面点占六成（$w=0.6$），要求 $P_0=0.99$，则 $1-w^{3}=1-0.216=0.784$，
\begin{equation}
k\ge\frac{\log 0.01}{\log 0.784}=\frac{-4.605}{-0.243}=18.9\;\Rightarrow\;k=19 .
\end{equation}
所以 20 次抽样就能以 99\% 的概率找到地面，这就是 RANSAC 快的原因。

\subsection{去畸变与动态点剔除}

激光雷达一帧的采集需要时间（$0.1\ \mathrm{s}$ 量级），机器人在这段时间里一直在动，于是``同一帧''里的点其实是在不同位姿下测到的——这就是运动畸变。去畸变要用 IMU 或匀速模型把每个点补偿到同一时刻，做法与第~\ref{ch:16-lio-evolution}~章相同。动态点（行人、另一个机器人、传送带）则会违反``环境静止''的假设，简单做法是检查该点在后几帧是否稳定存在，或直接用鲁棒核把大残差降权（第~\ref{ch:13-backend}~章）。

\section{描述子与回环检测}\label{sec:12-loop}

回环检测要回答一个全局问题：\textbf{当前帧是不是回到了很久以前来过的地方？}这属于``全局范围的数据关联''。它不能靠暴力匹配所有历史帧（太慢），所以要先算一个紧凑的``场景描述子''做粗筛，再做几何校验。

激光 SLAM 里最常用的描述子之一是 Scan Context\cite{kim2018scancontext}。它的构造非常直观：把点云按极坐标切成 $N_r$ 个同心环（ring）、$N_s$ 个扇区（sector），得到一个 $N_r\times N_s$ 的矩阵，第 $(r,s)$ 个元素取落入该 bin 的所有点的最大高度：
\begin{equation}
I_{rs}=\max\Big\{\,z_k \;\Big|\; \frac{r-1}{N_r}R< \rho_k\le\frac{r}{N_r}R,\;
\frac{2\pi(s-1)}{N_s}\le\varphi_k<\frac{2\pi s}{N_s}\Big\},
\label{eq:12-scancontext}
\end{equation}
其中 $(\rho_k,\varphi_k,z_k)$ 是第 $k$ 个点的柱坐标，$R$ 是最大量程。若某 bin 没有点，取 $0$。图~\ref{fig:12-scancontext} 是它的示意。

\begin{figure}[htbp]
  \centering
  \begin{tikzpicture}[font=\footnotesize]
    \draw[thick, structurecolor] (0,0) circle (2.4);
    \foreach \r in {0.6,1.2,1.8}
      \draw[structurecolor!55] (0,0) circle (\r);
    \foreach \a in {0,30,...,330}
      \draw[structurecolor!40] (0,0) -- (\a:2.4);
    \fill[structurecolor!35] (0:0.6) -- (30:0.6) arc (30:0:0.6) -- cycle;
    \fill[structurecolor!35] (0:1.2) -- (30:1.2) arc (30:0:1.2) -- cycle;
    \node[label] at (0,2.62) {环 $r$（径向距离）};
    \node[label, align=center] at (0,-2.85) {扇区 $s$（方位角）};
    \node[label, structurecolor] at (15:1.5) {bin $(r,s)$};
  \end{tikzpicture}
  \caption{Scan Context 的环—扇区划分。极坐标平面被切为 $N_r$ 个同心环与 $N_s$ 个扇区，每个 bin 记录该区域内点的最大高度，于是环境被压缩成一个 $N_r\times N_s$ 的矩阵描述子。}\label{fig:12-scancontext}
\end{figure}

为什么要这么切？因为\textbf{平移会让点云的笛卡尔坐标剧变，但旋转只会让极坐标矩阵的列发生循环移位}。所以比较两帧时，只需在列方向上做循环移位搜索：
\begin{equation}
D(I^{q},I^{c})=\min_{\delta\in\{0,1,\dots,N_s-1\}}
\frac{1}{N_s}\sum_{s=1}^{N_s}\Big(1-\frac{\bm{c}^{q\top}_{s}\bm{c}^{c}_{s+\delta}}
{\norm{\bm{c}^{q}_{s}}\norm{\bm{c}^{c}_{s+\delta}}}\Big),
\label{eq:12-sc-distance}
\end{equation}
其中 $\bm{c}_s\in\R^{N_r}$ 是矩阵的第 $s$ 列，$\delta$ 是候选的列移位量，相当于``假设机器人转过 $\delta\cdot360^{\circ}/N_s$''；下标 $s+\delta$ 在超过 $N_s$ 时按模 $N_s$ 回绕（循环移位）。两列方向一致时余弦相似度为 1、距离为 0；正交时为 1。取所有移位中的最小值 $D$，若 $D$ 小于阈值，就认为``这两帧很可能在同一个地方''。

\begin{keybox}{回环检测为什么属于前端}
回环检测的本质是\textbf{数据关联}——判断``这两帧观测的是不是同一个地方''，其输出是一条相对位姿约束（回环边），交给后端去优化。它不负责改轨迹，只负责``提出并校验假设''。把它放到后端会有两个问题：一是后端要接触原始点云，破坏了``前端造因子、后端解因子''的接口；二是回环检测必须按帧执行（频率高），而后端是低频的。所以正确划分是：\textbf{回环检测归前端，回环优化归后端。}
\end{keybox}

Scan Context 只是粗筛：$D$ 小不代表真是回环（相似的走廊、对称的房间都会骗过它），还必须做几何校验。典型做法是用粗筛给出的位姿作初值，做一次 ICP/NDT（第~\ref{ch:07-lidar-odom}~章）或分支定界搜索（第~\ref{ch:15-cartographer}~章），用匹配得分与内点比例确认；校验不过就丢掉这个候选。

\section{前端的失效模式}\label{sec:12-fail}

前端失效是 SLAM 最常见的故障来源。表~\ref{tab:12-fail} 把四类典型失效整理成``场景—信号—对策''，工程上可以照着做监控。

\begin{table}[htbp]
  \centering
  \caption{前端典型失效模式与对策}\label{tab:12-fail}
  \footnotesize
  \begin{tabular}{p{0.13\linewidth}p{0.20\linewidth}p{0.31\linewidth}p{0.26\linewidth}}
    \toprule
    失效类型 & 典型场景 & 可观测信号 & 对策 \\
    \midrule
    几何退化 & 长走廊、隧道、空旷场地：某一方向没有约束 & 海森矩阵最小特征值接近 0、条件数骤增、匹配得分对该方向不敏感 & 引入 IMU 强约束、对该方向降权或锁定、切换匹配特征 \\
    动态物体 & 行人、对方机器人、传送带 & 内点比例骤降、残差分布重尾、匹配得分忽高忽低 & 动态点剔除、鲁棒核降权、多帧一致性检验 \\
    传感器失效 & 遮挡、强光、雨雾、断流 & 点数骤减、强度分布异常、时间戳不连续 & 健康检查、故障切换、用运动模型预测维持 \\
    时间不同步 & 多传感器融合、急加速 & 残差随速度增大、标定外参漂移 & 硬件触发同步/PTP、在线估计时延、重新标定 \\
    \bottomrule
  \end{tabular}
\end{table}

\textbf{几何退化}值得单独强调，因为它最隐蔽：机器人走在一条笔直的长走廊里，两侧墙壁平行，前后方向几乎没有可区分的结构，于是这个方向的位姿是``不可观''的。此时匹配照样会给出一个数，但这个数完全不可信。检测方法是对海森矩阵做特征值分解，若最小特征值 $\lambda_{\min}$ 太小（相对于最大特征值 $\lambda_{\max}$），就说明该方向退化。对策是``承认自己不知道''：在这个方向上相信 IMU 或保持上一时刻的估计，而不是相信匹配结果。

\begin{pseudo}{前端单帧处理主循环}
  \item 读取一帧点云、IMU 与时间戳，做时间同步
  \item 去畸变，体素降采样，地面分割，标记动态点
  \item 用 IMU 或匀速模型预测当前位姿作为初值
  \item 在局部地图中按最近邻建立对应关系，用距离阈值粗筛
  \item 以初值为起点迭代求解位姿（扫描匹配），期间用鲁棒核降权
  \item 统计内点比例：低于阈值则判定跟踪失败，进入重定位流程
  \item 检查海森矩阵条件数：过大则标记该方向退化，切换融合策略
  \item 按式~\eqref{eq:12-kf} 判定是否生成关键帧
  \item 若是关键帧：加入局部地图窗口，计算描述子并检索回环候选
  \item 把相对位姿约束与回环候选发送给后端
\end{pseudo}

\begin{summary}
\begin{enumerate}
  \item 前端流水线是``预处理 $\rightarrow$ 数据关联 $\rightarrow$ 位姿估计 $\rightarrow$ 外点剔除 $\rightarrow$ 关键帧判定''（图~\ref{fig:12-pipeline}），分步的目的是让每类失败都可监控、可兜底。
  \item 数据关联是前端的心脏，常用 NN（式~\eqref{eq:12-nn}）、MNN、RANSAC 三件套；一条错误关联的破坏力大于十条正确关联的收益。
  \item 帧到帧匹配的误差按 $\sqrt{k}$ 放大（式~\eqref{eq:12-drift}）；帧到局部地图匹配把目标换成多帧融合地图，噪声按 $\sigma/\sqrt{n}$ 降低，且误差不进入匹配目标，所以更稳。
  \item 关键帧由平移、旋转、时间、信息量四类条件触发（式~\eqref{eq:12-kf}），数量过多与过少都有代价（表~\ref{tab:12-kf}）。
  \item 体素降采样用式~\eqref{eq:12-voxel-index} 划分、式~\eqref{eq:12-voxel-centroid} 代表，质心是``误差平方和最小''的解；RANSAC 平面拟合的抽样次数由式~\eqref{eq:12-ransac-k} 决定。
  \item Scan Context 把点云压成``环—扇区''矩阵（式~\eqref{eq:12-scancontext}），用列移位最小距离（式~\eqref{eq:12-sc-distance}）做旋转不变的粗筛，再用几何校验确认。
  \item 回环检测属于前端（全局数据关联），回环优化属于后端；前端典型失效有几何退化、动态物体、传感器失效、时间不同步四类（表~\ref{tab:12-fail}）。
\end{enumerate}


\end{summary}

\begin{problemset}
  \item 设每帧匹配误差标准差 $\sigma=1\ \mathrm{cm}$，$k=400$ 帧后按式~\eqref{eq:12-drift} 的标准差是多少？若改成帧到地图匹配、局部地图由 $n=9$ 帧融合，目标噪声标准差降为多少？两者相差几个数量级？
  \item 在式~\eqref{eq:12-nn} 中，如果 $\tau_d$ 取得过大和过小分别会发生什么？请结合``重影墙''现象说明。
  \item 为什么 RANSAC 的抽样次数只与内点比例 $w$ 有关，而与点的总数无关？这个结论在什么条件下不成立？
  \item Scan Context 的列移位为什么能处理旋转、却处理不了明显的平移？如果两帧之间存在较大平移，$D$ 会怎样变化？
  \item 机器人走在长走廊里，匹配器仍然输出了位姿。你会用哪些信号判断它其实已经退化？请给出至少两个可计算的指标。
  \item 把表~\ref{tab:12-fail} 中``动态物体''一行扩写成一段伪代码：你会怎样判断一个点是不是动态点？
\end{problemset}

\begin{reading}
\begin{enumerate}
  \item Kim \& Kim, ``Scan Context''\cite{kim2018scancontext}：本章回环描述子的原始论文，重点看第三节的列移位距离与第四节的候选校验。
  \item 第~\ref{ch:07-lidar-odom}~章：扫描匹配与 ICP/NDT 的完整推导，本章默认你已经掌握，只做了``升级''讨论。
  \item Cartographer 论文\cite{hess2016cartographer} 的前端部分：帧到子图匹配的工业级实现，可与本章 \S\ref{sec:12-ftm} 对照阅读。
  \item FAST-LIO2\cite{xu2022fastlio2}：直接法前端的代表，看看``不做特征提取''时数据关联如何组织。
\end{enumerate}


\end{reading}

% 本章使用 TikZ 的 arrows.meta / positioning 库（main.tex 与 tests 均已加载，此行幂等无害）
\usetikzlibrary{arrows.meta, positioning}

\chapter{SLAM 后端详解}\label{ch:13-backend}

\chapterintro{第~\ref{ch:12-frontend}~章的前端把原始数据变成一条条相对位姿约束，本章的后端要回答：拿到这一大堆互相矛盾的约束之后，怎么算出一条``全局一致''的轨迹？这是全书数学密度最高的章节之一。我们会从第~\ref{ch:11-slam-framework}~章的 MAP 表述出发，完整推导三件事：高斯噪声下``最大后验''如何变成``最小二乘''并如何迭代求解；位姿图的残差与雅可比如何一步一步求出来；束调整（BA）里让实时成为可能的``舒尔补''到底在做什么。最后两节讲工程上躲不开的两个问题：错误回环如何毁掉整张图（鲁棒核），以及滑窗优化里旧状态怎么``移出去而不丢信息''（边缘化）。}

\section{从最大后验到最小二乘}\label{sec:13-map2ls}

第~\ref{ch:11-slam-framework}~章已经做过完整推导：把联合后验拆成先验、运动、观测三类因子，在高斯噪声、观测独立、系统一阶马尔可夫的假设下取负对数，得到
\begin{equation}
X^{\star}=\argmin_{X}\;\sum_{k=1}^{K}\frac12\norm{\bm{r}_k(X)}_{\bm{\Sigma}_k^{-1}}^{2},
\qquad
\norm{\bm{r}}_{\bm{\Sigma}^{-1}}^{2}=\bm{r}^{\top}\bm{\Sigma}^{-1}\bm{r}.
\label{eq:13-ls}
\end{equation}
其中 $\bm{\Sigma}_k^{-1}$ 是第 $k$ 个因子的信息矩阵（定义见第~\ref{ch:09-fusion}~章），$\bm{r}_k$ 是该因子对应的残差，$X$ 是全部待估状态（位姿，可能还有地图点）。式~\eqref{eq:13-ls} 称为\textbf{非线性最小二乘}问题：``非线性''指 $\bm{r}_k(X)$ 一般不是 $X$ 的线性函数（旋转矩阵、投影、对数映射都是非线性的）。它成立的三个前提决定了本章方法的边界：噪声已知且为高斯（否则二范数会被外点带偏，见 \S\ref{sec:13-robust}）、数据关联已确定（关联错了该式本身发现不了）、线性化只是局部的（初值太差会收敛到错误的局部极小）。

\section{Gauss-Newton 与 Levenberg-Marquardt}\label{sec:13-gn}

式~\eqref{eq:13-ls} 没有解析解，只能迭代。高斯-牛顿（Gauss-Newton，GN）的思路是：在当前估计 $X$ 附近把残差线性化，把非线性问题变成一系列线性子问题。

\subsection{线性化与正规方程}

设当前估计为 $X$，给状态一个小的增量 $\delta$（在流形上用 ``$\oplus$'' 表示，定义见第~\ref{ch:02-math}~章，例如 $\bm{R}\oplus\delta\bm{\theta}=\bm{R}\exp(\delta\bm{\theta}^{\wedge})$）。对残差做一阶泰勒展开：
\begin{equation}
\bm{r}_k(X\oplus\delta)\approx\bm{r}_k(X)+\bm{J}_k\delta,
\qquad
\bm{J}_k=\frac{\partial\bm{r}_k}{\partial\delta}\in\R^{m_k\times n},
\label{eq:13-taylor}
\end{equation}
其中 $m_k$ 是残差维数，$n$ 是全部状态维数。代入式~\eqref{eq:13-ls}，第 $k$ 项变成
\begin{align}
\frac12\norm{\bm{r}_k+\bm{J}_k\delta}_{\bm{\Sigma}_k^{-1}}^{2}
&=\frac12\big(\bm{r}_k+\bm{J}_k\delta\big)^{\top}\bm{\Sigma}_k^{-1}\big(\bm{r}_k+\bm{J}_k\delta\big)\notag\\
&=\frac12\bm{r}_k^{\top}\bm{\Sigma}_k^{-1}\bm{r}_k
 +\delta^{\top}\bm{J}_k^{\top}\bm{\Sigma}_k^{-1}\bm{r}_k
 +\frac12\delta^{\top}\bm{J}_k^{\top}\bm{\Sigma}_k^{-1}\bm{J}_k\delta .
\label{eq:13-quadratic}
\end{align}
第一项与 $\delta$ 无关，可以丢掉。于是整体目标函数变成 $\delta$ 的二次型：
\begin{equation}
F(\delta)=\delta^{\top}\bm{b}+\frac12\delta^{\top}\bm{H}\delta,
\qquad
\bm{H}=\sum_{k}\bm{J}_k^{\top}\bm{\Sigma}_k^{-1}\bm{J}_k,\quad
\bm{b}=\sum_{k}\bm{J}_k^{\top}\bm{\Sigma}_k^{-1}\bm{r}_k .
\label{eq:13-Hb}
\end{equation}
对 $\delta$ 求导并令其为零：
\begin{equation}
\frac{\partial F}{\partial\delta}=\bm{b}+\bm{H}\delta=\bm{0}
\;\Longrightarrow\;
\bm{H}\delta=-\bm{b}.
\label{eq:13-normal}
\end{equation}
式~\eqref{eq:13-normal} 就是\textbf{正规方程}（定义见第~\ref{ch:02-math}~章），$\bm{H}$ 称为\textbf{海森矩阵}。严格的海森矩阵还应包含 $\sum_k\sum_i r_{k,i}\nabla^2 r_{k,i}$ 这样的二阶项，GN 把它忽略了：当残差接近零或残差的二阶导数很小时，这一项可以忽略；代价是初值很差时 GN 可能不收敛甚至发散。求解得到增量后按 $X\leftarrow X\oplus\delta$ 更新状态，再重新计算残差与雅可比，进入下一轮，这就是 GN 迭代。

\subsection{Levenberg-Marquardt：给步长加个``刹车''}

GN 的问题在于它假设线性化在整个步长内都有效。初值差时，算出来的 $\delta$ 可能很大，一步跨到线性化失效的区域，目标函数反而上升。Levenberg-Marquardt（LM）的解决办法很朴素：\textbf{在目标函数里惩罚过大的步长}，即求解
\begin{equation}
\min_{\delta}\;\frac12\norm{\bm{r}+\bm{J}\delta}_{\bm{\Sigma}^{-1}}^{2}+\frac12\lambda\norm{\delta}_2^{2}.
\label{eq:13-lm-objective}
\end{equation}
对 $\delta$ 求导：
\begin{equation}
\frac{\partial}{\partial\delta}\Big[\frac12(\bm{r}+\bm{J}\delta)^{\top}\bm{\Sigma}^{-1}(\bm{r}+\bm{J}\delta)+\frac{\lambda}{2}\delta^{\top}\delta\Big]
=\bm{J}^{\top}\bm{\Sigma}^{-1}\bm{r}+\bm{J}^{\top}\bm{\Sigma}^{-1}\bm{J}\delta+\lambda\delta .
\end{equation}
令其为零，得到 LM 的线性方程：
\begin{equation}
\big(\bm{H}+\lambda\bm{I}\big)\delta=-\bm{b},
\label{eq:13-lm-eq}
\end{equation}
其中 $\bm{I}$ 是单位矩阵。式~\eqref{eq:13-lm-eq} 的物理意义很清楚：$\lambda\to0$ 时它退化为 GN 方程~\eqref{eq:13-normal}，适合初值好、残差小的情形；$\lambda\to\infty$ 时 $\bm{H}$ 相对可忽略，$\delta\approx-\bm{b}/\lambda$，退化为沿负梯度方向的一小步（梯度下降），适合初值差的情形。工程上常取 $\lambda$ 与 $\bm{H}$ 对角线同量级，即 $(\bm{H}+\lambda\,\mathrm{diag}(\bm{H}))\delta=-\bm{b}$，以抵消不同状态量纲不同的影响。

LM 的迭代策略是``试探—接受/拒绝''：算出一个 $\delta$ 后先试探，若目标函数真的下降，就接受并\textbf{减小} $\lambda$（$\lambda\leftarrow\lambda/2$，更信任 GN）；若上升，就拒绝并\textbf{增大} $\lambda$（$\lambda\leftarrow2\lambda$，更保守）。这套策略使 LM 在初值差时也能稳定收敛，是 Ceres、g2o、GTSAM 的默认选择。

\begin{pseudo}{Gauss-Newton / Levenberg-Marquardt 迭代}
  \item 输入：初始状态 $X_0$、最大迭代次数、收敛阈值
  \item 初始化 $\lambda=\lambda_0$（如 $10^{-4}$），计算当前目标函数值 $F_0$
  \item \textbf{repeat}
  \item \quad 对每个因子计算残差 $\bm{r}_k$ 与雅可比 $\bm{J}_k$
  \item \quad 累加 $\bm{H}=\sum_k\bm{J}_k^{\top}\bm{\Sigma}_k^{-1}\bm{J}_k$，$\bm{b}=\sum_k\bm{J}_k^{\top}\bm{\Sigma}_k^{-1}\bm{r}_k$
  \item \quad 求解 $(\bm{H}+\lambda\bm{I})\delta=-\bm{b}$（用稀疏 Cholesky 分解）
  \item \quad 若 $\norm{\delta}$ 小于阈值：收敛，跳出
  \item \quad 试探 $X'=X\oplus\delta$，计算 $F'$
  \item \quad 若 $F'<F$：接受 $X\leftarrow X'$，$\lambda\leftarrow\lambda/2$
  \item \quad 否则：拒绝，保留 $X$，$\lambda\leftarrow2\lambda$
  \item \textbf{until} 达到最大迭代次数
\end{pseudo}

\subsection{海森矩阵为什么是稀疏的}

式~\eqref{eq:13-Hb} 里 $\bm{H}$ 是全部因子的 ${\bm{J}_k^{\top}\bm{\Sigma}_k^{-1}\bm{J}_k}$ 之和。关键事实是：\textbf{每个因子只涉及少数几个状态变量}。一个里程计因子只连第 $i$ 帧和第 $i+1$ 帧；一个点到面因子只连一帧位姿；回环因子只连两帧位姿。于是 $\bm{J}_k$ 在绝大多数列上都是零，$\bm{H}$ 的绝大部分元素也是零。

图~\ref{fig:13-sparse} 画出了 5 个位姿的位姿图对应的 $\bm{H}$ 结构（每个位姿 2 维，共 10 维）。对角块来自``自己与自己''的项，相邻块来自里程计边，左下与右上那个非零块来自回环边。一个 $N$ 帧的位姿图，$\bm{H}$ 是 $2N\times2N$（3D 时 $6N\times6N$）的矩阵，但非零元只有 $O(N)$ 个。正是这种稀疏性，让 ``$N=1000$ 也能毫秒级求解''成为可能。

\begin{figure}[htbp]
  \centering
  \begin{tikzpicture}[scale=0.36]
    \foreach \i in {0,...,9} \foreach \j in {0,...,9}
      \draw[black!12] (\i,\j) rectangle ++(1,1);
    % diagonal blocks (self terms)
    \foreach \b in {0,2,4,6,8}
      \fill[structurecolor] (\b,\b) rectangle ++(2,2);
    % odometry edges: off-diagonal 2x2 blocks
    \foreach \a/\b in {0/2,2/0,2/4,4/2,4/6,6/4,6/8,8/6}
      \fill[structurecolor!45] (\a,\b) rectangle ++(2,2);
    % loop edge blocks (0,8) and (8,0)
    \fill[red!60!black] (0,8) rectangle ++(2,2);
    \fill[red!60!black] (8,0) rectangle ++(2,2);
    \node[label] at (5,-1.0) {$5$ 个位姿、每个 $2$ 维，共 $10\times10$};
  \end{tikzpicture}
  \caption{位姿图海森矩阵的稀疏结构。深色为对角块（单个位姿自身项），中色为相邻位姿之间的里程计块，红色为回环边产生的远距离非零块。}\label{fig:13-sparse}
\end{figure}

\section{位姿图}\label{sec:13-posegraph}

\begin{definition}[因子图]
因子图（factor graph）是一种二分图，用来表示``一堆变量通过一堆约束联系在一起''的概率模型：图中有两类节点——\textbf{变量节点}（待估状态，如位姿、地图点）与\textbf{因子节点}（每个观测/先验构成的约束）。因子节点与它涉及的所有变量节点相连。式~\eqref{eq:13-ls} 的每一项就对应一个因子。
\end{definition}

\begin{definition}[位姿图]
位姿图（pose graph）是因子图的一个特例：变量节点只有位姿 $T_1,\dots,T_N\in SE(3)$，因子节点只有两类——相邻帧之间的``里程计因子''和回环产生的``回环因子''，每个因子都是一条二元边 $(i,j)$，约束 $T_i$ 与 $T_j$ 的相对位姿。位姿图是 SLAM 后端最常用的模型，因为它把地图点全部消掉，只优化位姿，规模小得多。
\end{definition}

\subsection{残差的定义}

对每条边 $(i,j)$，前端给出的测量是一个相对位姿 $Z_{ij}$（带信息矩阵 $\bm{\Omega}_{ij}$）。若两个位姿完全正确，那么用 $T_i,T_j$ 算出来的相对位姿应等于测量值：
\begin{equation}
Z_{ij}\;\overset{?}{=}\;T_i^{-1}T_j .
\end{equation}
二者之差用``误差变换'' $E_{ij}$ 表示：
\begin{equation}
E_{ij}=Z_{ij}^{-1}\,T_i^{-1}T_j\in SE(3).
\label{eq:13-error-transform}
\end{equation}
$E_{ij}$ 是单位变换就说明没有误差。用李代数把它映射成 $\R^{6}$ 中的向量（对数映射，定义见第~\ref{ch:02-math}~章）：
\begin{equation}
\bm{r}_{ij}=\log\big(Z_{ij}^{-1}T_i^{-1}T_j\big)^{\vee}\in\R^{6}.
\label{eq:13-pg-residual}
\end{equation}
式~\eqref{eq:13-pg-residual} 中 $(\cdot)^{\vee}$ 把反对称矩阵变回向量，$\log$ 是李群到李代数的对数映射。为什么要取 $\log$ 而不是直接减？因为 $SE(3)$ 不是向量空间，两个变换``相减''没有定义；$\log$ 把``差多少''映射到切空间上，才能用欧氏范数量化。

\subsection{雅可比的逐步推导（二维情形）}

3D 的完整推导要用到 BCH 公式，篇幅很长。为了把每一步都写清楚，我们在二维 $SE(2)$ 上推导，最后给出 3D 的结论。二维位姿写成 $T_i=(\bm{R}_i,\bm{t}_i)$，其中 $\bm{R}_i$ 是 $2\times2$ 旋转矩阵，$\bm{t}_i\in\R^{2}$。

先把式~\eqref{eq:13-pg-residual} 展开成具体分量。预测的相对位姿是
\begin{equation}
\hat{\bm{R}}=\bm{R}_i^{\top}\bm{R}_j,\qquad
\hat{\bm{t}}=\bm{R}_i^{\top}(\bm{t}_j-\bm{t}_i).
\end{equation}
误差变换 $E=Z^{-1}\hat{T}$ 的旋转与平移分别是
\begin{equation}
\bm{R}_E=\bm{R}_z^{\top}\hat{\bm{R}}=\bm{R}_z^{\top}\bm{R}_i^{\top}\bm{R}_j,
\qquad
\bm{t}_E=\bm{R}_z^{\top}(\hat{\bm{t}}-\bm{t}_z).
\label{eq:13-E-2d}
\end{equation}
残差取 $\bm{r}=(\bm{t}_E,\theta_E)$，其中 $\theta_E$ 是 $\bm{R}_E$ 对应的转角：
\begin{equation}
\theta_E=\mathrm{atan2}\big(R_E(2,1),\,R_E(1,1)\big).
\end{equation}
残差维数 3，正好等于单个二维位姿的自由度（2 平移 $+$ 1 旋转）。

现在施加扰动。采用\textbf{右扰动}（扰动加在自身坐标系，见第~\ref{ch:02-math}~章）：$T_i\leftarrow T_i\exp(\delta_i)$，$T_j\leftarrow T_j\exp(\delta_j)$，其中 $\delta=(\delta\bm{\rho},\delta\theta)\in\R^{3}$。对 $SE(2)$ 的指数映射做一阶展开（$\delta\theta$ 很小时 $\sin\delta\theta\approx\delta\theta$、$1-\cos\delta\theta\approx0$）：
\begin{equation}
\exp(\delta)\approx\Big(\bm{I}+\begin{bmatrix}0&-\delta\theta\\\delta\theta&0\end{bmatrix},\ \delta\bm{\rho}\Big)
=\Big(\bm{I}+[\delta\theta]_\times,\ \delta\bm{\rho}\Big),
\label{eq:13-exp2d}
\end{equation}
其中 $[\delta\theta]_\times$ 是二维反对称矩阵。于是
\begin{equation}
\bm{R}_i\to\bm{R}_i(\bm{I}+[\delta\theta_i]_\times),\qquad
\bm{t}_i\to\bm{t}_i+\bm{R}_i\delta\bm{\rho}_i .
\label{eq:13-perturb-i}
\end{equation}

\textbf{第一步：算 $\bm{t}_E$ 对 $\delta_i$ 的导数。}把式~\eqref{eq:13-perturb-i} 代入 $\hat{\bm{t}}$：
\begin{align}
\hat{\bm{t}}\to\big(\bm{R}_i(\bm{I}+[\delta\theta_i]_\times)\big)^{\top}
\big((\bm{t}_j-\bm{t}_i)-\bm{R}_i\delta\bm{\rho}_i\big)
&=\big(\bm{I}-[\delta\theta_i]_\times\big)\bm{R}_i^{\top}(\bm{t}_j-\bm{t}_i)
 -\big(\bm{I}-[\delta\theta_i]_\times\big)\delta\bm{\rho}_i\notag\\
&=\hat{\bm{t}}-[\delta\theta_i]_\times\hat{\bm{t}}-\delta\bm{\rho}_i+O(\delta^{2}).
\end{align}
这里用了 $(\bm{R}_i(\bm{I}+[\delta\theta]_\times))^{\top}=(\bm{I}+[\delta\theta]_\times)^{\top}\bm{R}_i^{\top}=(\bm{I}-[\delta\theta]_\times)\bm{R}_i^{\top}$，以及 $\bm{R}_i^{\top}\bm{R}_i=\bm{I}$。所以
\begin{equation}
\delta\hat{\bm{t}}=-\delta\bm{\rho}_i-[\delta\theta_i]_\times\hat{\bm{t}} .
\end{equation}
再由 $\bm{t}_E=\bm{R}_z^{\top}(\hat{\bm{t}}-\bm{t}_z)$ 得
\begin{equation}
\delta\bm{t}_E=\bm{R}_z^{\top}\delta\hat{\bm{t}}
=-\bm{R}_z^{\top}\delta\bm{\rho}_i-\delta\theta_i\,\bm{R}_z^{\top}[\bm{e}_z]_\times\hat{\bm{t}}
=-\bm{R}_z^{\top}\delta\bm{\rho}_i-\delta\theta_i\,[\bm{e}_z]_\times\bm{R}_z^{\top}\hat{\bm{t}},
\label{eq:13-dte-i}
\end{equation}
最后一步用了二维情形的交换关系 $\bm{R}_z^{\top}[\bm{e}_z]_\times=[\bm{e}_z]_\times\bm{R}_z^{\top}$（旋转与``逆时针转 90 度''可交换）。式~\eqref{eq:13-dte-i} 说明：$\delta\bm{\rho}_i$ 经 $\bm{R}_z^{\top}$ 旋转后直接进入 $\bm{t}_E$，而 $\delta\theta_i$ 的效果是让 $\hat{\bm{t}}$ 绕原点转了一下，幅度与 $\hat{\bm{t}}$ 的长度成正比。

\textbf{第二步：算 $\bm{t}_E$ 对 $\delta_j$ 的导数。}同理
\begin{align}
\hat{\bm{t}}\to\bm{R}_i^{\top}\big(\bm{t}_j+\bm{R}_j\delta\bm{\rho}_j-\bm{t}_i\big)
&=\hat{\bm{t}}+\bm{R}_i^{\top}\bm{R}_j\delta\bm{\rho}_j
=\hat{\bm{t}}+\hat{\bm{R}}\,\delta\bm{\rho}_j,
\end{align}
于是
\begin{equation}
\delta\bm{t}_E=\bm{R}_z^{\top}\hat{\bm{R}}\,\delta\bm{\rho}_j=\bm{R}_E\,\delta\bm{\rho}_j,
\end{equation}
即 $\partial\bm{t}_E/\partial\delta\bm{\rho}_j=\bm{R}_E$，而 $\partial\bm{t}_E/\partial\delta\theta_j=\bm{0}$（$\hat{\bm{t}}$ 与 $\bm{R}_j$ 无关）。

\textbf{第三步：算 $\theta_E$ 的导数。}由式~\eqref{eq:13-E-2d}，
\begin{equation}
\bm{R}_E\to\bm{R}_z^{\top}\big(\bm{R}_i(\bm{I}+[\delta\theta_i]_\times)\big)^{\top}\bm{R}_j
=\bm{R}_z^{\top}(\bm{I}-[\delta\theta_i]_\times)\bm{R}_i^{\top}\bm{R}_j
=\bm{R}_E-\bm{R}_E[\delta\theta_i]_\times .
\end{equation}
最后一步把 $[\delta\theta_i]_\times$ 移到最右边（同样利用交换性）。而 $\bm{R}_E-\bm{R}_E[\delta\theta]_\times=\bm{R}_E(\bm{I}-[\delta\theta]_\times)=\bm{R}_E\bm{R}(-\delta\theta)$，相当于转角减少 $\delta\theta_i$，所以
\begin{equation}
\frac{\partial\theta_E}{\partial\delta\theta_i}=-1 .
\end{equation}
对 $\delta_j$ 同理有 $\bm{R}_E\to\bm{R}_E(\bm{I}+[\delta\theta_j]_\times)=\bm{R}_E\bm{R}(\delta\theta_j)$，转角增加 $\delta\theta_j$，故 $\partial\theta_E/\partial\delta\theta_j=1$。平移扰动对转角没有影响。

\textbf{第四步：拼成矩阵。}把上面三步的结果按 $\delta=(\delta\bm{\rho},\delta\theta)$ 的列顺序拼起来，得到 $3\times6$ 的分块雅可比：
\begin{equation}
\frac{\partial\bm{r}}{\partial(\delta_i,\delta_j)}=
\left[
\begin{array}{cc|cc}
-\bm{R}_z^{\top} & -[\bm{e}_z]_\times\bm{R}_z^{\top}\hat{\bm{t}} & \bm{R}_E & \bm{0}\\[2pt]
\bm{0}^{\top} & -1 & \bm{0}^{\top} & 1
\end{array}
\right]\in\R^{3\times6}.
\label{eq:13-pg-jac}
\end{equation}
（为紧凑起见，矩阵中每个 2×2 或 2×1 的块各写成一个元素；按行展开即为标准的 $3\times6$ 矩阵。）式~\eqref{eq:13-pg-jac} 中左半 $3\times3$ 是残差对 $T_i$ 的雅可比，右半 $3\times3$ 是对 $T_j$ 的雅可比，所有量都在当前估计处求值，每轮迭代重算一次。读它时不必死记：平移残差对自己平移扰动的响应是 $-\bm{R}_z^{\top}$（把误差转到测量坐标系），旋转残差对自己是 $-1$、对对方是 $+1$（因为相对位姿是 $T_i^{-1}T_j$，一正一负），而两处旋转扰动都会``串''到平移残差里。

\subsection{三维情形}

三维位姿图残差仍是式~\eqref{eq:13-pg-residual}，只是 $\bm{r}_{ij}\in\R^{6}$（3 平移 $+$ 3 旋转）。按同样的``右扰动 + 一阶 BCH''推导（BCH 公式见第~\ref{ch:02-math}~章），可得
\begin{equation}
\frac{\partial\bm{r}_{ij}}{\partial\delta\bm{\xi}_i}
=-\bm{J}_l^{-1}(\bm{r}_{ij})\,\mathrm{Ad}\big(T_j^{-1}T_i\big),
\qquad
\frac{\partial\bm{r}_{ij}}{\partial\delta\bm{\xi}_j}
=\bm{J}_r^{-1}(\bm{r}_{ij}),
\label{eq:13-pg-jac3d}
\end{equation}
其中 $\mathrm{Ad}(\cdot)$ 是 $SE(3)$ 的伴随矩阵，
\begin{equation}
\mathrm{Ad}(T)=\begin{bmatrix}\bm{R}&[\bm{t}]_\times\bm{R}\\\bm{0}&\bm{R}\end{bmatrix}\in\R^{6\times6},
\end{equation}
$\bm{J}_l^{-1},\bm{J}_r^{-1}$ 是 $SE(3)$ 左/右雅可比的逆。当残差较小时 $\bm{J}_l^{-1}\approx\bm{J}_r^{-1}\approx\bm{I}$，很多实现直接采用简化形式
\begin{equation}
\frac{\partial\bm{r}_{ij}}{\partial\delta\bm{\xi}_i}\approx-\mathrm{Ad}\big(T_j^{-1}T_i\big),
\qquad
\frac{\partial\bm{r}_{ij}}{\partial\delta\bm{\xi}_j}\approx\bm{I},
\end{equation}
这也是 g2o\cite{kummerle2011g2o} 中位姿图边的经典形式。式~\eqref{eq:13-pg-jac3d} 的推导与二维完全同构：把式~\eqref{eq:13-perturb-i} 里的 $\bm{R}_i(\bm{I}+[\delta\theta]_\times)$ 换成 $\bm{R}_i\exp([\delta\bm{\theta}]_\times)$，把标量转角换成 $\R^{3}$ 旋转向量，把二维叉乘换成 $3\times3$ 反对称矩阵，再用 BCH 的一阶展开把``$\exp$ 的乘积''并成``对数的和''。

\section{从位姿图到 BA：舒尔补}\label{sec:13-schur}

位姿图把地图点消掉了。但如果我们要同时优化``相机位姿 + 地图点''，模型里就多了一大类变量。这就是光束法平差（Bundle Adjustment，BA，定义见第~\ref{ch:08-visual-odom}~章）：以重投影误差为残差，同时优化相机位姿与三维点。

\subsection{BA 的残差与雅可比}

设第 $i$ 个相机位姿为 $T_i\in SE(3)$（世界到相机的变换），第 $j$ 个地图点在世界系下为 $\bm{X}_j\in\R^{3}$。投影与重投影残差为
\begin{equation}
\bm{u}_{ij}=\pi(\bm{R}_i\bm{X}_j+\bm{t}_i),\qquad
\bm{r}_{ij}=\bm{z}_{ij}-\pi(\bm{R}_i\bm{X}_j+\bm{t}_i)\in\R^{2},
\end{equation}
其中 $\pi(\bm{p})=\big(f_x p_x/p_z+c_x,\ f_y p_y/p_z+c_y\big)^{\top}$ 是针孔投影（$f_x,f_y$ 为焦距，$c_x,c_y$ 为主点），$\bm{z}_{ij}$ 是观测到的像素坐标。残差由两部分雅可比组成。先算投影对相机坐标的导数：
\begin{equation}
\frac{\partial\pi}{\partial\bm{p}}=\begin{bmatrix}
\dfrac{f_x}{p_z}&0&-\dfrac{f_x p_x}{p_z^{2}}\\[4pt]
0&\dfrac{f_y}{p_z}&-\dfrac{f_y p_y}{p_z^{2}}
\end{bmatrix}\in\R^{2\times3}.
\end{equation}
再算相机坐标对位姿扰动 $\delta\bm{\xi}=(\delta\bm{\rho},\delta\bm{\phi})$ 的导数。位姿扰动 $T\leftarrow T\exp(\delta\bm{\xi})$ 作用在点上，一阶展开为 $\bm{p}\to\bm{p}+[\delta\bm{\phi}]_\times\bm{p}+\delta\bm{\rho}$，写成矩阵即
\begin{equation}
\frac{\partial\bm{p}}{\partial\delta\bm{\xi}}=\begin{bmatrix}\bm{I}&-[\bm{p}]_\times\end{bmatrix}\in\R^{3\times6}
\end{equation}
（因为 $[\delta\bm{\phi}]_\times\bm{p}=-[\bm{p}]_\times\delta\bm{\phi}$）。对点 $\bm{X}_j$ 的导数是 $\partial\bm{p}/\partial\bm{X}_j=\bm{R}_i$。链式相乘，最终
\begin{equation}
\bm{J}^{\mathrm{pose}}_{ij}=-\frac{\partial\pi}{\partial\bm{p}}\begin{bmatrix}\bm{I}&-[\bm{p}_{ij}]_\times\end{bmatrix}\in\R^{2\times6},
\qquad
\bm{J}^{\mathrm{point}}_{ij}=-\frac{\partial\pi}{\partial\bm{p}}\bm{R}_i\in\R^{2\times3}.
\label{eq:13-ba-jac}
\end{equation}
看式~\eqref{eq:13-ba-jac} 的维度就知道：\textbf{一条观测只涉及一个位姿（6 维）和一个点（3 维）}，所以 $\bm{H}$ 一定是极端稀疏的。

\subsection{分块与舒尔补}

把所有变量排成``相机在前、点在后''的顺序，海森矩阵自然分成四块：
\begin{equation}
\begin{bmatrix}
\bm{H}_{cc}&\bm{H}_{cp}\\[2pt]
\bm{H}_{pc}&\bm{H}_{pp}
\end{bmatrix}
\begin{bmatrix}\delta\bm{\xi}_c\\[2pt]\delta\bm{X}_p\end{bmatrix}
=
\begin{bmatrix}-\bm{b}_c\\[2pt]-\bm{b}_p\end{bmatrix}.
\label{eq:13-ba-block}
\end{equation}
其中 $\bm{H}_{cc}$ 是``相机—相机''块（尺寸 $6N_c\times6N_c$），$\bm{H}_{pp}$ 是``点—点''块（$3N_p\times3N_p$），$\bm{H}_{cp}=\bm{H}_{pc}^{\top}$ 是交叉块。

关键观察：\textbf{每个三维点只在``观测到它的那些相机''里出现，点与点之间没有任何因子}。所以 $\bm{H}_{pp}$ 是\textbf{块对角}的：
\begin{equation}
\bm{H}_{pp}=\mathrm{diag}\big(\bm{H}_{1},\bm{H}_{2},\dots,\bm{H}_{N_p}\big),
\qquad \bm{H}_{j}\in\R^{3\times3}.
\end{equation}
块对角矩阵的逆就是逐块求逆，每个只有 $3\times3$，极其便宜。这正是舒尔补能用的前提。

\begin{definition}[舒尔补]
对分块矩阵 $\begin{bmatrix}\bm{A}&\bm{B}\\\bm{C}&\bm{D}\end{bmatrix}$（$\bm{D}$ 可逆），称 $\bm{A}-\bm{B}\bm{D}^{-1}\bm{C}$ 为其（关于 $\bm{D}$ 的）舒尔补（Schur complement）。它在消去分块线性方程组中的变量 $\delta_p$ 时自然出现。
\end{definition}

下面一步一步消去 $\delta\bm{X}_p$。把式~\eqref{eq:13-ba-block} 写成两个方程：
\begin{align}
\bm{H}_{cc}\delta\bm{\xi}_c+\bm{H}_{cp}\delta\bm{X}_p&=-\bm{b}_c, \label{eq:13-ba-r1}\\
\bm{H}_{pc}\delta\bm{\xi}_c+\bm{H}_{pp}\delta\bm{X}_p&=-\bm{b}_p. \label{eq:13-ba-r2}
\end{align}
由式~\eqref{eq:13-ba-r2} 解出 $\delta\bm{X}_p$（$\bm{H}_{pp}$ 块对角可逆）：
\begin{equation}
\delta\bm{X}_p=\bm{H}_{pp}^{-1}\big(-\bm{b}_p-\bm{H}_{pc}\delta\bm{\xi}_c\big).
\label{eq:13-ba-dp}
\end{equation}
代入式~\eqref{eq:13-ba-r1}：
\begin{align}
\bm{H}_{cc}\delta\bm{\xi}_c+\bm{H}_{cp}\bm{H}_{pp}^{-1}\big(-\bm{b}_p-\bm{H}_{pc}\delta\bm{\xi}_c\big)&=-\bm{b}_c\notag\\
\bm{H}_{cc}\delta\bm{\xi}_c-\bm{H}_{cp}\bm{H}_{pp}^{-1}\bm{H}_{pc}\delta\bm{\xi}_c&=-\bm{b}_c+\bm{H}_{cp}\bm{H}_{pp}^{-1}\bm{b}_p\notag\\
\Big(\bm{H}_{cc}-\bm{H}_{cp}\bm{H}_{pp}^{-1}\bm{H}_{pc}\Big)\delta\bm{\xi}_c&=-\bm{b}_c+\bm{H}_{cp}\bm{H}_{pp}^{-1}\bm{b}_p .
\label{eq:13-reduced}
\end{align}
式~\eqref{eq:13-reduced} 就是\textbf{缩减相机系统}（reduced camera system），它的系数矩阵
\begin{equation}
\bm{S}=\bm{H}_{cc}-\bm{H}_{cp}\bm{H}_{pp}^{-1}\bm{H}_{pc}
\label{eq:13-schur}
\end{equation}
正是舒尔补。解出 $\delta\bm{\xi}_c$ 后，回代式~\eqref{eq:13-ba-dp} 就得到所有点的增量。

为什么这能``实时''？设 $N_c=200$ 个相机、$N_p=20000$ 个点：直接解 $\bm{H}\delta=-\bm{b}$ 是 $6\times200+3\times20000=61200$ 维，点变量极多、填充（fill-in）严重；舒尔补后只解 $6\times200=1200$ 维的 $\bm{S}$，点变量则通过 $3\times3$ 块逐块求逆、逐块回代，计算量与点数成线性。两者规模比约为 $(61200/1200)^{3}\approx1.3\times10^{5}$，这就是``BA 能实时''的秘密。位姿图可以看作``把所有点都边缘化掉之后的 BA''，这也是它规模小、速度快的原因。

\begin{pseudo}{基于舒尔补的 BA 一轮求解}
  \item 遍历所有观测，累加 $\bm{H}_{cc},\bm{H}_{cp},\bm{H}_{pp},\bm{b}_c,\bm{b}_p$
  \item 对每个点 $j$：求 $\bm{H}_j^{-1}$（$3\times3$），计算 $\bm{H}_j^{-1}\bm{H}_{pc}$ 的对应行块与 $\bm{H}_j^{-1}\bm{b}_p$
  \item 累加 $\bm{S}=\bm{H}_{cc}-\sum_j\bm{H}_{cp,j}\bm{H}_j^{-1}\bm{H}_{pc,j}$ 与右端
  \item 用稀疏 Cholesky 解 $\bm{S}\,\delta\bm{\xi}_c=$ 右端
  \item 对每个点 $j$：回代 $\delta\bm{X}_j=\bm{H}_j^{-1}(-\bm{b}_{p,j}-\bm{H}_{pc,j}\delta\bm{\xi}_c)$
  \item 用 $\oplus$ 更新全部位姿与点，进入下一轮迭代
\end{pseudo}

\section{回环与鲁棒核}\label{sec:13-robust}

回环闭合是后端最有价值、也最危险的功能。图~\ref{fig:13-factorgraph} 上半部分是正常情况下的一条链式因子图：位姿 $x_0\to x_1\to\cdots\to x_5$，相邻帧有里程计因子，$x_5$ 与 $x_0$ 之间有一条回环因子。正是这条回环边把``漂移了几十厘米''的轨迹拉回闭合。下半部分是错误回环：把``看起来很像、其实不是同一个地方''的两帧连起来，优化器为了满足这条假约束，会把整段轨迹拧成一团。

\begin{figure}[htbp]
  \centering
  \begin{tikzpicture}[font=\footnotesize]
    % panel A: correct loop closure
    \foreach \i in {0,...,5}{
      \node[circle, draw=structurecolor, thick, fill=structurecolor!12, minimum size=6.5mm] (x\i) at (\i*1.5,1.2) {$x_{\i}$};
    }
    \foreach \i [evaluate=\i as \j using int(\i+1)] in {0,...,4}{
      \node[rectangle, fill=structurecolor, minimum size=2.4mm, inner sep=0pt] (f\i) at (\i*1.5+0.75,0.25) {};
      \draw[black!60] (f\i) -- (x\i);
      \draw[black!60] (f\i) -- (x\j);
    }
    \node[rectangle, fill=structurecolor, minimum size=2.4mm, inner sep=0pt] (fp) at (-1.0,1.2) {};
    \draw[black!60] (fp) -- (x0);
    \node[label] at (-1.0,1.8) {先验};
    \node[rectangle, fill=red!60!black, minimum size=2.4mm, inner sep=0pt] (fl) at (3.75,2.35) {};
    \draw[red!60!black] (fl) -- (x0);
    \draw[red!60!black] (fl) -- (x5);
    \node[label, red!60!black] at (3.75,2.85) {回环因子};
    % panel B: wrong loop closure
    \foreach \i in {0,...,4}{
      \node[circle, draw=structurecolor, thick, fill=structurecolor!12, minimum size=6.5mm] (y\i) at (\i*1.5,-2.3) {$x_{\i}$};
    }
    \node[rectangle, fill=red!60!black, minimum size=2.4mm, inner sep=0pt] (wl) at (2.9,-0.85) {};
    \draw[red!60!black] (wl) -- (y1);
    \draw[red!60!black] (wl) -- (y3);
    \node[label, red!60!black] at (2.9,-0.45) {错误回环};
    \node[label, anchor=west, align=left] at (7.6,1.2) {上：正确回环，\\轨迹被拉闭合\\[3pt]下：错误回环，\\轨迹被拧成一团};
  \end{tikzpicture}
  \caption{上：正确的回环。圆圈是位姿变量，实心方块是因子节点；相邻位姿由里程计因子相连，$x_5$ 与 $x_0$ 之间的回环因子把累积漂移拉回闭合。下：错误回环。$x_1$ 与 $x_3$ 被错误地连起来，优化器为了满足这条假约束会把轨迹拧成一团，破坏地图一致性。}\label{fig:13-factorgraph}
\end{figure}

问题来了：最小二乘对残差是\textbf{平方}惩罚，一条错边的残差很大时，平方会把它的``话语权''放大到压倒所有正确约束。鲁棒核函数（robust kernel / M-estimator）的思路是：\textbf{让大残差的惩罚增长得比平方慢}，从而自动给可疑约束降权。

\begin{definition}[鲁棒核函数]
鲁棒核函数是一个非负、单调递增、在 $s=0$ 处为 $0$ 的函数 $\rho(s)$，用于把``残差的平方''$\norm{\bm{r}}^2_{\bm{\Omega}}$ 替换成 $\rho\big(\norm{\bm{r}}^2_{\bm{\Omega}}\big)$，使优化目标变为
\begin{equation}
\min_{X}\;\sum_{k}\rho\big(s_k\big),\qquad s_k=\norm{\bm{r}_k}_{\bm{\Omega}_k}^{2}.
\end{equation}
当 $\rho(s)=s$ 时退化为普通最小二乘；$\rho$ 在 $s$ 很大时增长得比线性慢（次线性），就得到鲁棒性。
\end{definition}

最常用的是 Huber 核：
\begin{equation}
\rho(s)=\begin{cases}
s, & s\le\delta^{2},\\[2pt]
2\delta\sqrt{s}-\delta^{2}, & s>\delta^{2},
\end{cases}
\label{eq:13-huber}
\end{equation}
其中 $\delta$ 是``软阈值''，$s$ 是残差平方范数（注意 $\sqrt{s}=\norm{\bm{r}}_{\bm{\Omega}}$）。在 $s=\delta^2$ 处两侧值相等（$\delta^2$），导数也相等，所以 $\rho$ 是连续可导的。

\textbf{推导它对迭代求解的影响。}目标函数对状态的梯度为
\begin{equation}
\frac{\partial}{\partial X}\sum_k\rho(s_k)
=\sum_k\rho'(s_k)\frac{\partial s_k}{\partial X}
=\sum_k\rho'(s_k)\cdot 2\,\bm{J}_k^{\top}\bm{\Omega}_k\bm{r}_k,
\end{equation}
其中用了 $\partial s_k/\partial X=2\bm{J}_k^{\top}\bm{\Omega}_k\bm{r}_k$。令梯度为零、约掉常数 2，得
\begin{equation}
\sum_k w_k\,\bm{J}_k^{\top}\bm{\Omega}_k\bm{r}_k=\bm{0},
\qquad w_k=\rho'(s_k).
\end{equation}
再用 GN 的线性化 $\bm{r}_k+\bm{J}_k\delta$ 代替 $\bm{r}_k$，就得到\textbf{加权正规方程}：
\begin{equation}
\Big(\sum_k w_k\bm{J}_k^{\top}\bm{\Omega}_k\bm{J}_k\Big)\delta
=-\sum_k w_k\bm{J}_k^{\top}\bm{\Omega}_k\bm{r}_k .
\label{eq:13-irls}
\end{equation}
式~\eqref{eq:13-irls} 与普通 GN 的唯一区别是每个因子的信息矩阵被乘上了权重 $w_k$。对 Huber 核，$\rho'(s)=1$（$s\le\delta^2$）或 $\rho'(s)=\delta/\sqrt{s}$（$s>\delta^2$），于是
\begin{equation}
w_k=\min\Big(1,\;\frac{\delta}{\norm{\bm{r}_k}_{\bm{\Omega}_k}}\Big).
\label{eq:13-huber-weight}
\end{equation}
这就是 Huber 核的全部魔力：\textbf{残差小于阈值时权重为 1（完全相信），残差大于阈值时权重与残差成反比}。残差再大，它对法方程的贡献也不会超过阈值对应的量级。这种``用当前残差算权重、再解加权最小二乘''的迭代方式，称为迭代重加权最小二乘（Iteratively Reweighted Least Squares，IRLS）。

\begin{keybox}{鲁棒核不能``修正''错误关联}
鲁棒核只是\textbf{降权}，不是\textbf{修正}。一条错误的回环边即便权重降到很低，它仍然在``往错误方向拉''，只是拉得轻了；当错误边很多、或者错误边与其他边形成``一致但错误''的组合时，降权救不了场。所以工程上要三层保险：前端回环几何校验（第~\ref{ch:12-frontend}~章）、后端鲁棒核降权、以及\textbf{开关变量}（给每条回环边一个 $s_k\in[0,1]$ 的连续变量，与位姿一起优化，收敛到 0 表示``这条边被判定为假''）。开源实现里 GTSAM 的 \lstinline{SwitchableConstraint} 就是这个思想\cite{dellaert2017factorgraphs}。
\end{keybox}

\section{边缘化：把旧状态``移出去''但不丢信息}\label{sec:13-marginalization}

滑窗优化的困境是：窗口里只能保留最近 $M$ 个位姿，更旧的状态必须移出窗口，否则计算量无限增长。但直接把旧状态``删掉''会丢信息——旧状态和保留状态之间是相关的，删掉就等于假装它们独立，会引入不一致。

正确做法是\textbf{边缘化}：把旧状态从概率分布里积掉，留下它对保留状态的影响。在线性化框架下，这件事正好又是舒尔补。

\begin{definition}[边缘化]
边缘化（marginalization）是指从联合概率分布 $p(X_m,X_r)$ 中，通过对被移除变量 $X_m$ 积分得到保留变量 $X_r$ 的边缘分布：
\begin{equation}
p(X_r)=\int p(X_m,X_r)\,\diff X_m .
\end{equation}
在高斯情形下，积分的结果仍是一个高斯分布，其信息矩阵恰好是舒尔补。
\end{definition}

推导与 \S\ref{sec:13-schur} 的块消元完全相同：把联合目标函数在旧线性化点附近写成二次型，按 $(\delta_m,\delta_r)$ 分块，从 $\delta_m$ 那一块解出
\begin{equation}
\bm{H}_{mm}\delta_m+\bm{H}_{mr}\delta_r+\bm{b}_m=\bm{0}
\;\Longrightarrow\;
\delta_m=\bm{H}_{mm}^{-1}\big(-\bm{b}_m-\bm{H}_{mr}\delta_r\big),
\end{equation}
再代回另一块，就得到只含 $\delta_r$ 的二次型：
\begin{equation}
F(\delta_r)=\frac12\delta_r^{\top}\underbrace{\Big(\bm{H}_{rr}-\bm{H}_{rm}\bm{H}_{mm}^{-1}\bm{H}_{mr}\Big)}_{\bm{\Lambda}}\delta_r
+\underbrace{\Big(\bm{b}_r-\bm{H}_{rm}\bm{H}_{mm}^{-1}\bm{b}_m\Big)}_{\bm{\eta}}^{\top}\delta_r+\text{常数}.
\label{eq:13-marg-prior}
\end{equation}
式~\eqref{eq:13-marg-prior} 中的 $\bm{\Lambda}$ 与 $\bm{\eta}$ 定义了一个关于 $\delta_r$ 的先验因子：$\bm{\Lambda}$ 是它的信息矩阵，$\bm{\eta}$ 是信息向量。下一轮优化时，把这个先验因子加到目标函数里，就等价于``没有丢掉旧状态的信息''。

边缘化有一个绕不开的副作用：\textbf{线性化误差}。式~\eqref{eq:13-marg-prior} 是在旧线性化点 $\bar{X}$ 附近做的二次近似，它只在 $\bar{X}$ 附近有效。随着后续迭代把 $X_r$ 更新到离 $\bar{X}$ 越来越远，这个先验会``拉''住新的估计，产生偏差。更麻烦的是，如果每一轮都用新的线性化点重新计算 $\bm{\Lambda},\bm{\eta}$，边缘化的``零点''会漂移，导致同一信息被重复计数或丢失（不一致）。工程上的标准解法是\textbf{固定线性化点}（First Estimate Jacobian，FEJ）：边缘化产生的先验，其雅可比始终在第一次线性化点处计算，不再更新。这是 VINS-Mono\cite{qin2018vinsmono}、OKVIS 等滑窗 VIO 的关键实现细节。

\section{求解器工程}\label{sec:13-solver}

最后把``用什么库''这件事说清楚。表~\ref{tab:13-solver} 对比了三个最常用的后端库。

\begin{table}[htbp]
  \centering
  \caption{三个常用后端求解器的定位}\label{tab:13-solver}
  \footnotesize
  \begin{tabular}{p{0.10\linewidth}p{0.20\linewidth}p{0.30\linewidth}p{0.30\linewidth}}
    \toprule
    库 & 模型组织方式 & 优势 & 典型使用场景 \\
    \midrule
    g2o & 图（顶点/边），C++ & 轻量、可定制、贴合 ROS 生态 & 教学与研究，Cartographer 早期版本 \\
    GTSAM & 因子图，C++ & 增量式求解（iSAM2）、理论完备 & LIO-SAM、需要在线增量的系统 \\
    Ceres & 通用非线性最小二乘 & 自动微分、鲁棒核齐全、性能强 & Cartographer、VINS-Mono \\
    \bottomrule
  \end{tabular}
\end{table}

三者底层是同一件事：构造 $\bm{H}\delta=-\bm{b}$，用稀疏 Cholesky（如 SuiteSparse 的 CHOLMOD）求解；区别只在\textbf{如何组织因子}与\textbf{如何利用增量信息}。g2o 用``顶点—边''的图结构，最贴近 $\bm{H}$ 的块稀疏本质，但不是增量式的，每轮都要重新分解；GTSAM 的 iSAM2 用贝叶斯树（Bayes tree）只更新被新因子影响的部分，适合在线长期运行；Ceres 提供自动微分与齐全的鲁棒核，写残差最省事，但通常整图重解。量级参考（桌面 CPU）：1000 个节点的位姿图一轮优化 $10\sim50\ \mathrm{ms}$；200 个相机 $+$ 2 万个点的 BA 一轮 $50\sim200\ \mathrm{ms}$。所以后端通常\textbf{降频运行}（每 $1\sim5\ \mathrm{s}$ 一次，或检测到回环时触发），而不是每帧都优化。

\begin{summary}
\begin{enumerate}
  \item 后端要解的是第~\ref{ch:11-slam-framework}~章得到的非线性最小二乘式~\eqref{eq:13-ls}（前提：噪声高斯、关联正确、初值不太差）。Gauss-Newton 用一阶泰勒展开把残差线性化，得到正规方程 $\bm{H}\delta=-\bm{b}$，其中 $\bm{H}=\sum\bm{J}^{\top}\bm{\Sigma}^{-1}\bm{J}$、$\bm{b}=\sum\bm{J}^{\top}\bm{\Sigma}^{-1}\bm{r}$（式~\eqref{eq:13-Hb}）；海森矩阵稀疏的根本原因是每个因子只连少数变量（图~\ref{fig:13-sparse}）。
  \item LM 在目标函数上加 $\frac{\lambda}{2}\norm{\delta}^2$ 的惩罚，得到 $(\bm{H}+\lambda\bm{I})\delta=-\bm{b}$（式~\eqref{eq:13-lm-eq}）；$\lambda$ 大时接近梯度下降、小时接近 GN，按``下降就减、上升就增''调节，因此初值差时也稳。
  \item 位姿图是只有位姿变量的因子图，残差取误差变换的李代数 $\bm{r}_{ij}=\log(Z_{ij}^{-1}T_i^{-1}T_j)^{\vee}$（式~\eqref{eq:13-pg-residual}）；二维情形可逐步求出 $3\times6$ 的分块雅可比（式~\eqref{eq:13-pg-jac}），3D 情形用伴随矩阵写成式~\eqref{eq:13-pg-jac3d}。
  \item BA 的 $\bm{H}$ 按相机与点分块后，$\bm{H}_{pp}$ 块对角，消去点变量得到缩减相机系统，其系数正是舒尔补 $\bm{S}=\bm{H}_{cc}-\bm{H}_{cp}\bm{H}_{pp}^{-1}\bm{H}_{pc}$（式~\eqref{eq:13-schur}），求解规模从``点数''降到``相机数''。
  \item 鲁棒核把残差平方换成次线性增长函数；Huber 核给出的 IRLS 权重是 $w_k=\min(1,\delta/\norm{\bm{r}_k})$（式~\eqref{eq:13-huber-weight}），本质是给可疑约束降权，但\textbf{不能修正}错误关联。
  \item 边缘化用舒尔补把旧状态积掉，得到保留状态的先验因子（式~\eqref{eq:13-marg-prior}）；代价是线性化误差，工程上用固定线性化点（FEJ）缓解。g2o / GTSAM / Ceres 底层都是稀疏最小二乘，区别在因子组织方式与是否增量求解，后端通常降频运行。
\end{enumerate}


\end{summary}

\begin{problemset}
  \item 在式~\eqref{eq:13-ls} 中，若某个因子的信息矩阵 $\bm{\Sigma}^{-1}$ 取为单位阵，最小二乘就退化为``残差平方和''。此时 Huber 核与直接对残差做截断（把大残差直接丢弃）有什么区别？
  \item 用 $SU(2)$ 或 $SO(2)$ 的直觉解释：为什么位姿图残差必须取 $\log$，而不能直接把两个变换矩阵相减再取范数？举一个数值反例。
  \item 在式~\eqref{eq:13-pg-jac} 中令 $\bm{R}_z=\bm{I}$、$\hat{\bm{t}}=\bm{0}$，请写出简化后的雅可比矩阵，并解释它的物理含义。
  \item 假设 BA 有 $N_c=100$ 个相机、$N_p=10^4$ 个点。直接解与舒尔补解各要处理多少维的线性系统？若 Cholesky 复杂度为 $O(n^3)$，两者相差多少倍？
  \item 边缘化产生的先验因子为什么不能每轮用新的线性化点重算？如果不固定线性化点，可能出现什么现象（提示：信息被重复计数或丢失）？
\end{problemset}

\begin{reading}
\begin{enumerate}
  \item Dellaert \& Kaess, ``Factor Graphs for Robot Perception''\cite{dellaert2017factorgraphs}：因子图与增量求解的系统教程，第 3--4 章可配合本章 \S\ref{sec:13-posegraph} 阅读。
  \item Kümmerle 等，g2o 论文\cite{kummerle2011g2o}：本章位姿图部分的工程对照，重点看 ``Types of Edges'' 与稀疏求解器；Barfoot《State Estimation for Robotics》\cite{barfoot2017state} 第 7 章可补 $SE(3)$ 雅可比与伴随矩阵的细节。
\end{enumerate}


\end{reading}

% 本章使用 TikZ 的 arrows.meta / positioning 库（main.tex 与 tests 均已加载，此行幂等无害）
\usetikzlibrary{arrows.meta, positioning}

\chapter{地图与建图}\label{ch:14-mapping}

\chapterintro{前两章都在回答``我在哪''，本章回答``世界是什么样''。地图是 SLAM 的输出，也是导航的输入：没有地图，全局规划无从谈起；地图错了，定位与避障都会跟着错。本章沿着``地图给谁用''这条线索，把最常用的五种地图讲清楚：占据栅格（概率化的二维地图）、八叉树（省内存的三维占据地图）、TSDF（能重建光滑表面的截断符号距离场）、ESDF（能给规划提供梯度的距离场）、以及点云地图。重点是两条完整推导：占据栅格的 log-odds 更新如何从贝叶斯公式一步步推出来，以及 TSDF 的加权平均为什么是``最小二乘''的自然结果。}

\section{地图的用途决定地图的形式}\label{sec:14-usage}

新手常问``哪种地图最好''，这个问题没有答案，因为地图是为用途服务的。同一片场地，用来做全局定位的``地图''可能是一张点云，用来做路径规划的``地图''可能是一张二维栅格，用来做局部避障的``地图''可能是占据栅格，用来做轨迹优化的``地图''则是 ESDF。表~\ref{tab:14-map-usage} 给出常见的对应关系。

\begin{table}[htbp]
  \centering
  \caption{地图类型与用途的对应关系（$\checkmark$ 表示适合，$\circ$ 表示勉强可用）}\label{tab:14-map-usage}
  \footnotesize
  \begin{tabular}{lccccc}
    \toprule
    地图类型 & 定位 & 全局规划 & 局部避障 & 表面重建 & 内存 \\
    \midrule
    点云地图 & $\checkmark$ & $\circ$ & $\circ$ & $\checkmark$ & 大 \\
    占据栅格（2D） & $\checkmark$ & $\checkmark$ & $\checkmark$ & $\circ$ & 中 \\
    八叉树（3D 占据） & $\checkmark$ & $\checkmark$ & $\checkmark$ & $\circ$ & 中 \\
    TSDF & $\circ$ & $\circ$ & $\circ$ & $\checkmark$ & 大 \\
    ESDF & $\circ$ & $\checkmark$ & $\checkmark$ & $\circ$ & 大 \\
    \bottomrule
  \end{tabular}
\end{table}

表~\ref{tab:14-map-usage} 的读法是：\textbf{要定位就选能算似然的}（点云、占据栅格），\textbf{要规划就选能算代价的}（栅格、ESDF），\textbf{要重建就选能算距离的}（TSDF）。实际系统里往往同时维护好几张地图，各司其职。

\begin{keybox}{地图的三条工程准则}
\begin{enumerate}
  \item \textbf{分辨率决定一切}：$5\ \mathrm{cm}$ 的栅格与 $20\ \mathrm{cm}$ 的栅格是两种应用，前者能看清桌腿，后者只能看墙。分辨率减半，内存变四倍（2D）或八倍（3D）。
  \item \textbf{只维护``用得到''的区域}：机器人周围 $20\ \mathrm{m}$ 内的地图需要高频更新，远处的地图只需保存；这是八叉树、体素地图存在的理由。
  \item \textbf{地图必须能``变旧''}：环境会变（人走了、箱子搬了），地图若永远只增不减，就会把历史当成现实。所有工程地图都要么带时间戳，要么允许权重衰减。
\end{enumerate}
\end{keybox}

\section{占据栅格地图与 log-odds 更新}\label{sec:14-occupancy}

\begin{definition}[占据栅格]
占据栅格（occupancy grid）把空间切成规则的格子（2D 是正方形，3D 是立方体），每个格子 $m_i$ 只取两个状态：被占据（$m_i=1$，有障碍）或空闲（$m_i=0$，可通行）。由于传感器有噪声，我们不存``是/否''，而存``被占据的概率'' $p(m_i=1\mid \bm{z}_{1:t})$。
\end{definition}

一张 $100\ \mathrm{m}\times100\ \mathrm{m}$、分辨率 $5\ \mathrm{cm}$ 的二维栅格有 $2000\times2000=4\times10^{6}$ 个格子；若用单精度浮点存概率，需要 $16\ \mathrm{MB}$。这还只是一层。所以存储格式的优化从第一天就要考虑——这就是为什么工程上不存概率，而存``概率的对数''（log-odds）。

\subsection{从贝叶斯公式推导 log-odds}

目标：给定从第 $1$ 帧到第 $t$ 帧的全部观测 $\bm{z}_{1:t}$，求每个格子被占据的后验概率 $p(m_i\mid \bm{z}_{1:t})$。为简洁，记 $m=m_i$。由贝叶斯公式（定义见第~\ref{ch:02-math}~章）：
\begin{equation}
p(m\mid \bm{z}_{1:t})=\frac{p(\bm{z}_t\mid m,\bm{z}_{1:t-1})\,p(m\mid \bm{z}_{1:t-1})}{p(\bm{z}_t\mid \bm{z}_{1:t-1})}.
\label{eq:14-bayes1}
\end{equation}
\textbf{假设（观测只依赖当前地图状态）：}$p(\bm{z}_t\mid m,\bm{z}_{1:t-1})=p(\bm{z}_t\mid m)$。也就是说，这一帧的观测只和``这个格子现在是什么''有关，与历史无关。于是
\begin{equation}
p(m\mid \bm{z}_{1:t})=\frac{p(\bm{z}_t\mid m)\,p(m\mid \bm{z}_{1:t-1})}{p(\bm{z}_t\mid \bm{z}_{1:t-1})}.
\label{eq:14-bayes2}
\end{equation}
对``未占据''事件 $\neg m$（即 $m=0$）写同样的式子：
\begin{equation}
p(\neg m\mid \bm{z}_{1:t})=\frac{p(\bm{z}_t\mid \neg m)\,p(\neg m\mid \bm{z}_{1:t-1})}{p(\bm{z}_t\mid \bm{z}_{1:t-1})}.
\label{eq:14-bayes3}
\end{equation}
注意式~\eqref{eq:14-bayes2} 与式~\eqref{eq:14-bayes3} 的分母相同。两式相除，分母消掉：
\begin{equation}
\frac{p(m\mid \bm{z}_{1:t})}{p(\neg m\mid \bm{z}_{1:t})}
=
\frac{p(\bm{z}_t\mid m)}{p(\bm{z}_t\mid\neg m)}
\cdot
\frac{p(m\mid \bm{z}_{1:t-1})}{p(\neg m\mid \bm{z}_{1:t-1})}.
\label{eq:14-odds}
\end{equation}
左边是当前时刻的``占据比''，右边第二项是上一时刻的``占据比''，第一项是这一帧观测带来的``修正因子''。

定义占据比（odds）与对数占据比（log-odds）：
\begin{equation}
\mathrm{odds}(p)=\frac{p}{1-p},
\qquad
l=\log\frac{p}{1-p}.
\label{eq:14-logodds-def}
\end{equation}
$\mathrm{odds}$ 把概率 $(0,1)$ 映射到 $(0,\infty)$，$\log$ 再把它映射到 $(-\infty,\infty)$——这正是我们想要的：概率接近 $0$ 或 $1$ 时不会``撞墙''，而且乘性更新变成加性更新。对式~\eqref{eq:14-odds} 两边取对数，得到本章的核心递推式：
\begin{equation}
l_t=l_{t-1}+\log\frac{p(\bm{z}_t\mid m)}{p(\bm{z}_t\mid\neg m)}.
\label{eq:14-logodds-recursion}
\end{equation}

再处理一下修正因子。对式~\eqref{eq:14-bayes1}（把 $t$ 换成单帧）用一次贝叶斯，可得
\begin{equation}
\frac{p(m\mid\bm{z}_t)}{p(\neg m\mid\bm{z}_t)}
=\frac{p(\bm{z}_t\mid m)}{p(\bm{z}_t\mid\neg m)}\cdot\frac{p(m)}{p(\neg m)},
\end{equation}
取对数即
\begin{equation}
\log\frac{p(\bm{z}_t\mid m)}{p(\bm{z}_t\mid\neg m)}
=\underbrace{\log\frac{p(m\mid\bm{z}_t)}{1-p(m\mid\bm{z}_t)}}_{l(m\mid\bm{z}_t)}
-\underbrace{\log\frac{p(m)}{1-p(m)}}_{l(m)}.
\end{equation}
把它代入式~\eqref{eq:14-logodds-recursion}，就得到教材里最常见的形式：
\begin{equation}
\boxed{\,l_t(m)=l_{t-1}(m)+l(m\mid\bm{z}_t)-l(m)\,}
\label{eq:14-logodds-update}
\end{equation}
其中 $l(m\mid\bm{z}_t)$ 是``只用这一帧的观测算出的占据 log-odds''（称为\textbf{逆观测模型}的值），$l(m)$ 是先验 log-odds。

式~\eqref{eq:14-logodds-update} 的物理含义：\textbf{每来一帧观测，就把``这一帧认为它有多大概率被占据''与``本来认为它有多大概率''之差，加到历史累积值上}。如果先验取 $p(m)=0.5$（即 $l(m)=0$，完全不知道），更新就简化为 $l_t=l_{t-1}+l(m\mid\bm{z}_t)$，即``观测直接累加''。

\subsection{逆观测模型与数值例子}

激光雷达的逆观测模型非常直观：一条射线如果打到了某个格子，该格子就是``命中''（占据）；射线穿过但没打到的格子就是``空闲''；没有被任何射线穿过的格子``未知''，不更新。用两个常数近似：
\begin{equation}
p(m\mid\bm{z}_t)=\begin{cases}
p_{\mathrm{occ}}, & \text{射线终点命中该格},\\
p_{\mathrm{free}}, & \text{射线穿过该格},\\
0.5\ (\text{不更新}), & \text{未被观测}.
\end{cases}
\end{equation}
常见取值 $p_{\mathrm{occ}}=0.7$、$p_{\mathrm{free}}=0.3$。

\begin{example}[三帧观测把概率推到多少]
取先验 $p(m)=0.5$，则 $l(m)=0$，命中增量 $l_{\mathrm{occ}}=\log(0.7/0.3)=0.847$，空闲增量 $l_{\mathrm{free}}=\log(0.3/0.7)=-0.847$。假设同一个格子连续被命中 3 次：
\begin{equation}
l_3=0+3\times0.847=2.541
\;\Longrightarrow\;
p(m\mid\bm{z}_{1:3})=\frac{1}{1+e^{-2.541}}=0.927 .
\end{equation}
若连续命中 10 次，$l_{10}=8.47$，$p=0.99979$。可以看到概率会越来越``自信''，最终逼近 $1$ 但永远不到 $1$——这在数值上是好事（避免除零），但也带来一个工程问题：如果某天这个位置被清空了（搬箱子），要把它``拉回''空闲需要很多帧观测，因为它的 log-odds 已经很大了。所以工程上会把 $l$ 截断在 $[-l_{\max},l_{\max}]$ 之间（常取 $l_{\max}=2\sim4$），让地图保持``可修改''。
\end{example}

\begin{pseudo}{占据栅格的 log-odds 更新（单帧）}
  \item 输入：栅格 $l$、当前帧位姿、雷达点云
  \item 对每条射线：从传感器原点出发，沿射线方向步进
  \item 沿途经过的格子：$l\gets l+l_{\mathrm{free}}$
  \item 射线命中的格子：$l\gets l+l_{\mathrm{occ}}$
  \item 对更新过的每个格子：$l\gets\max(-l_{\max},\min(l_{\max},l))$ 截断
  \item 需要查询概率时：$p=1/(1+e^{-l})$
\end{pseudo}

把式~\eqref{eq:14-logodds-update} 落到代码上只有几行——注意工程上还要做截断与概率还原：
\begin{lstlisting}[language=C++]
// occupancy grid: log-odds update along a ray
float l_hit  = std::log(p_occ  / (1.0f - p_occ));   // hit increment
float l_miss = std::log(p_free / (1.0f - p_free));  // miss increment
for (auto& cell : cells_along_ray) {
    cell.log_odds += cell.is_endpoint ? l_hit : l_miss;
    // clamp so the map stays updatable
    cell.log_odds = std::clamp(cell.log_odds, -l_max, l_max);
}
// convert back to probability for display / planning
float p = 1.0f / (1.0f + std::exp(-cell.log_odds));
\end{lstlisting}

\begin{figure}[htbp]
  \centering
  \begin{tikzpicture}[scale=0.8]
    % grid
    \foreach \i in {0,...,7} \foreach \j in {0,...,5}
      \draw[black!20] (\i,\j) rectangle ++(1,1);
    % free cells along three rays
    \foreach \x in {1,...,6} \fill[structurecolor!16] (\x,0) rectangle ++(1,1);
    \foreach \i in {1,...,4} \fill[structurecolor!16] (\i,\i) rectangle ++(1,1);
    \foreach \y in {1,...,4} \fill[structurecolor!16] (0,\y) rectangle ++(1,1);
    % hit cells
    \fill[black!70] (7,0) rectangle ++(1,1);
    \fill[black!70] (5,5) rectangle ++(1,1);
    \fill[black!70] (0,5) rectangle ++(1,1);
    % rays
    \draw[red!60!black, thick] (0.5,0.5) -- (7.0,0.5);
    \draw[red!60!black, thick] (0.5,0.5) -- (5.0,5.0);
    \draw[red!60!black, thick] (0.5,0.5) -- (0.5,5.0);
    % sensor
    \fill[red!60!black] (0.5,0.5) circle (3pt);
    \node[label, anchor=north] at (0.9,-0.12) {传感器};
    % legend
    \fill[black!70] (0.2,6.15) rectangle ++(0.5,0.5);
    \node[label, anchor=west] at (0.8,6.4) {命中：占据};
    \fill[structurecolor!16] (3.7,6.15) rectangle ++(0.5,0.5);
    \node[label, anchor=west] at (4.3,6.4) {穿过：空闲};
    \draw[black!25] (7.0,6.15) rectangle ++(0.5,0.5);
    \node[label, anchor=west] at (7.6,6.4) {未观测：未知};
  \end{tikzpicture}
  \caption{一帧二维激光扫描后的占据栅格。三条射线从传感器出发，穿过的格子被更新为``空闲''，命中的格子被更新为``占据''，其余格子保持``未知''。log-odds 让多次观测可以像记账一样累加。}\label{fig:14-occupancy}
\end{figure}

\section{八叉树地图}\label{sec:14-octree}

占据栅格的内存问题在 3D 下被放大到不可接受：$100\ \mathrm{m}\times100\ \mathrm{m}\times10\ \mathrm{m}$ 的场地、$10\ \mathrm{cm}$ 分辨率，有 $1000\times1000\times100=10^{8}$ 个体素；每个体素一个概率就是几百 MB，而且绝大多数体素是空的、根本不需要存。

八叉树（octree）的基本思想是``\textbf{只存有差异的地方，相同的兄弟节点合并成一个父节点}''。把整个空间看成一个立方体（根节点），递归地切成 8 个小立方体（8 个孩子），直到达到目标分辨率。每个节点存一个占据概率：若 8 个孩子都一样（都空、都占据或都未知），就\textbf{剪枝}——只保留父节点，删掉孩子。

\begin{definition}[八叉树]
八叉树（octree）是一种每个内部节点恰好有 8 个子节点的树结构，用于对三维空间做递归八分。叶节点对应分辨率最高的体素，内部节点代表其覆盖区域的整体状态。若一个节点的所有子节点状态相同，则用该节点替代其子树（称为剪枝），从而以远小于稠密体素的内存表示大范围空间。
\end{definition}

八叉树的深度与分辨率的关系是：若根节点边长 $L$、叶节点边长 $r$，则深度
\begin{equation}
d=\left\lceil\log_2\frac{L}{r}\right\rceil .
\end{equation}
定位一个空间点的过程就是``每层比较 $x,y,z$ 与中点的大小，走 8 个分支之一''，$O(d)$ 次比较即可。更新一个体素就是沿这条路径走到叶节点、更新其概率，再沿路返回做剪枝判断。OctoMap\cite{hornung2013octomap} 就是用 log-odds 作为节点概率、并带剪枝的经典实现。

\begin{figure}[htbp]
  \centering
  \begin{tikzpicture}[font=\footnotesize]
    \node[block] (root) at (4.5,3.6) {根节点};
    \foreach \i in {0,...,7}{
      \node[rectangle, draw=structurecolor, thick, fill=structurecolor!8,
            minimum size=4mm, inner sep=0pt] (c\i) at (1.2*\i+0.9,2.2) {};
      \draw[black!55] (root) -- (c\i);
    }
    \node[label] at (0.9,1.8) {8 个孩子};
    % expand the 2nd child into 8 leaves
    \foreach \i in {0,...,7}{
      \node[rectangle, draw=structurecolor!60, fill=structurecolor!25,
            minimum size=2.4mm, inner sep=0pt] (l\i) at (0.62*\i+0.25,0.4) {};
      \draw[black!40] (c1) -- (l\i);
    }
    \node[label] at (7.4,0.4) {继续八分，直到叶节点};
  \end{tikzpicture}
  \caption{八叉树的节点结构示意（示意图，实际每个节点对应一个空间立方体）。根节点代表整个空间，每层切成 8 份；当 8 个孩子状态相同时合并回父节点，这就是八叉树省内存的关键。}\label{fig:14-octree}
\end{figure}

内存对比给出直观量级。同样是 $100\ \mathrm{m}\times100\ \mathrm{m}\times10\ \mathrm{m}$、$10\ \mathrm{cm}$ 分辨率：
\begin{itemize}
  \item 稠密体素：$10^{8}$ 个体素，若每个 1 字节也要 $100\ \mathrm{MB}$，若存 float 则 $400\ \mathrm{MB}$，实际还会因为频繁查询而更慢；
  \item 八叉树：真正需要表示的只有``表面附近''的体素。室内场景的表面积约 $2000\ \mathrm{m}^{2}$，每个 $0.1\ \mathrm{m}$ 体素覆盖 $0.01\ \mathrm{m}^{2}$，约 $2\times10^{5}$ 个叶节点；加上内部节点，总节点数在 $10^{5}\sim10^{6}$ 量级，内存只需几 MB 到几十 MB\cite{hornung2013octomap}。
\end{itemize}
两个数量级的差距，就是八叉树成为 3D 占据地图标准做法的原因。

\section{TSDF：截断符号距离场}\label{sec:14-tsdf}

占据栅格只回答``这个体素有没有东西''，不回答``表面在哪、有多远''。要做高质量表面重建（比如 Marching Cubes 出网格、给规划做精细避障），需要一种连续的、带距离信息的地图——截断符号距离场（Truncated Signed Distance Field，TSDF）。

\begin{definition}[TSDF]
TSDF 是把空间中每个体素 $x$ 映射到一个带符号的距离值 $D(x)$ 的场：$D(x)>0$ 表示 $x$ 在表面前方（自由空间），$D(x)<0$ 表示在表面后方（被物体占据），$|D(x)|$ 是到最近表面的距离。为了控制内存与计算量，只保留 $|D(x)|\le\mu$ 的部分（$\mu$ 称为截断距离），超出部分截断为 $\pm\mu$。
\end{definition}

为什么必须截断？因为远离表面的距离对建图没有用（我们只关心表面附近那一层薄壳），而``到处都要算到最近表面的精确距离''的计算量极大。截断 $\mu$ 通常取几倍体素分辨率（如 $\mu=4r$）。

\subsection{单帧观测给出的 SDF 值}

设传感器在 $\bm{o}$，某条射线上的深度观测为 $z$（沿射线到表面的距离）。对射线上的一个体素 $x$，它到传感器的距离是 $\|\bm{x}-\bm{o}\|$。若这个体素位于表面之前，则它到表面的距离（沿射线方向）约为
\begin{equation}
d(\bm{x})=z-\|\bm{x}-\bm{o}\| ,
\end{equation}
即``表面比体素远多少''。$d>0$ 表示体素在表面前方，$d<0$ 表示在后方，正是带符号距离。截断后
\begin{equation}
\hat d(\bm{x})=\mathrm{clamp}\big(d(\bm{x}),\,-\mu,\,+\mu\big)
=\begin{cases}
+\mu, & d>\mu,\\
d, & -\mu\le d\le\mu,\\
-\mu, & d<-\mu .
\end{cases}
\label{eq:14-tsdf-clamp}
\end{equation}

\subsection{多帧融合：加权平均就是最小二乘}

一帧观测不够准（噪声、视角）。多帧观测同一个体素时，最自然的做法是加权平均。设第 $t$ 帧给出的截断距离是 $\hat d_t(\bm{x})$，权重 $w_t(\bm{x})\ge0$，则
\begin{equation}
D_t(\bm{x})=\frac{W_{t-1}(\bm{x})\,D_{t-1}(\bm{x})+w_t(\bm{x})\,\hat d_t(\bm{x})}
{W_{t-1}(\bm{x})+w_t(\bm{x})},
\qquad
W_t(\bm{x})=W_{t-1}(\bm{x})+w_t(\bm{x}).
\label{eq:14-tsdf-fusion}
\end{equation}
这个式子是``加权平均''的递归写法。它为什么是最优的？因为它是下面最小二乘问题的解：
\begin{equation}
D_t(\bm{x})=\argmin_{D}\;\sum_{k=1}^{t}w_k(\bm{x})\big(D-\hat d_k(\bm{x})\big)^{2}.
\label{eq:14-tsdf-ls}
\end{equation}
对 $D$ 求导并令其为零：
\begin{equation}
\frac{\partial}{\partial D}\sum_{k=1}^{t}w_k(D-\hat d_k)^{2}
=2\sum_{k=1}^{t}w_k(D-\hat d_k)=0
\;\Longrightarrow\;
D\sum_{k=1}^{t}w_k=\sum_{k=1}^{t}w_k\hat d_k .
\end{equation}
所以
\begin{equation}
D_t=\frac{\sum_{k=1}^{t}w_k\hat d_k}{\sum_{k=1}^{t}w_k}.
\label{eq:14-tsdf-closed}
\end{equation}
把式~\eqref{eq:14-tsdf-closed} 的分子分母同时拆成``前 $t-1$ 项 $+$ 第 $t$ 项''，并记 $W_{t-1}=\sum_{k=1}^{t-1}w_k$、$W_{t-1}D_{t-1}=\sum_{k=1}^{t-1}w_k\hat d_k$，立刻得到式~\eqref{eq:14-tsdf-fusion}。也就是说，\textbf{TSDF 的递归更新不是``工程上的近似''，而是最小二乘解的等价递推}——每来一帧只需 $O(1)$ 的更新，不需要保存全部历史。

\subsection{权重设计与表面提取}

权重 $w_t$ 怎么选？物理上有两条直觉：\textbf{距离越远越不准}（深度误差随距离平方增长），\textbf{入射角越斜越不准}（斜射时同样的深度误差对应更大的表面位置误差）。所以常用
\begin{equation}
w_t(\bm{x})=\frac{\cos\theta}{\|\bm{x}-\bm{o}\|^{2}},
\end{equation}
其中 $\theta$ 是射线方向与表面法向的夹角。工程上还会给累计权重设上限 $W\le W_{\max}$，使得``很久以前的观测''权重不再增长，地图才能慢慢忘掉旧东西、适应环境变化。

有了 $D(x)$ 的体素网格，用 Marching Cubes 在每个小立方体的 8 个角上找 $D$ 的零交叉点，插值出表面三角形，就得到网格模型。这就是``重建''。

\begin{figure}[htbp]
  \centering
  \begin{tikzpicture}[x=1.1cm,y=1.0cm,font=\footnotesize]
    % truncation band (drawn first, so the profile stays visible)
    \fill[structurecolor!12] (1.6,-1.4) rectangle (4.4,1.4);
    % axes
    \draw[->, thick] (-0.2,0) -- (6.4,0) node[label, right] {沿射线的距离};
    \draw[->, thick] (0,-2.0) -- (0,2.0) node[label, above] {$D$};
    % truncation bounds
    \draw[dashed, black!45] (0,1.4) -- (6.2,1.4);
    \draw[dashed, black!45] (0,-1.4) -- (6.2,-1.4);
    \node[label, anchor=west] at (6.25,1.4) {$+\mu$};
    \node[label, anchor=west] at (6.25,-1.4) {$-\mu$};
    % surface
    \draw[dashed, structurecolor, thick] (3.0,-2.0) -- (3.0,2.0);
    \node[label, structurecolor, above] at (3.0,2.0) {表面};
    % tsdf profile (drawn last)
    \draw[very thick, red!60!black] (0.1,1.4) -- (1.6,1.4) -- (3.0,0) -- (4.4,-1.4) -- (6.2,-1.4);
    \node[label, align=center] at (2.75,-1.0) {截断区\\$|D|\le\mu$};
    \node[label, align=center] at (5.35,0.65) {远景\\截断为常数};
  \end{tikzpicture}
  \caption{TSDF 沿一条射线的剖面。表面处 $D=0$，表面前方 $D>0$、后方 $D<0$；只有 $|D|\le\mu$ 的薄壳被保留（有色区域），远处一律截断为 $\pm\mu$。多帧观测在薄壳内做加权平均。}\label{fig:14-tsdf}
\end{figure}

\section{ESDF：欧氏符号距离场}\label{sec:14-esdf}

TSDF 的距离是``沿射线的近似距离''，够用来重建表面，但不够用来做规划——轨迹优化需要的是``空间中每一点到最近障碍的\textbf{欧氏}距离''。这就是 ESDF（Euclidean Signed Distance Field）。

\begin{definition}[ESDF]
记障碍物点集为 $\mathcal{O}$（通常取 TSDF 的零交叉面或占据体素），ESDF 在每个点 $\bm{x}$ 处的值定义为
\begin{equation}
d(\bm{x})=\min_{\bm{y}\in\mathcal{O}}\norm{\bm{x}-\bm{y}}_2 ,
\end{equation}
即 $\bm{x}$ 到最近障碍点的欧氏距离；符号约定为障碍外部为正、内部为负（只取外部时称为距离场）。
\end{definition}

ESDF 最关键的性质是它的\textbf{梯度}。设 $\bm{x}$ 的最近障碍点是 $\bm{y}^{\star}$，则 $d(\bm{x})=\norm{\bm{x}-\bm{y}^{\star}}_2$。直接对 $\bm{x}$ 求导（$\bm{y}^{\star}$ 是使距离最小的点，由包络定理可以当作常数处理）：
\begin{equation}
\nabla d(\bm{x})
=\frac{\bm{x}-\bm{y}^{\star}}{\norm{\bm{x}-\bm{y}^{\star}}_2},
\label{eq:14-esdf-grad}
\end{equation}
这是一个\textbf{单位向量}，方向是``从最近障碍指向 $\bm{x}$''。它有两个直接推论：第一，$\norm{\nabla d}=1$（处处成立，除了障碍边界上的奇点）；第二，沿着 $-\nabla d$ 走就是``最快远离障碍''的方向。

这正是规划器想要的：在轨迹优化（第~\ref{ch:20-local-planning}~章）里，把``离障碍太近''写成惩罚项 $c(\bm{x})=\max(0,\,d_{\mathrm{safe}}-d(\bm{x}))^{2}$，它的梯度可以直接由 ESDF 给出，从而把轨迹``推离''障碍。没有 ESDF，就需要对每个轨迹点做最近邻查询，代价高且不连续。

怎么高效算 ESDF？最直接的做法是对每个体素暴力搜索最近障碍点，复杂度 $O(N\cdot|\mathcal{O}|)$，不可接受。工程上有两类做法：一是\textbf{波前传播}（从障碍体素出发做 BFS，把距离值一层层推出去，即距离变换）；二是 Voxblox\cite{oleynikova2017voxblox} 的增量做法——先建 TSDF，再只在``新观测改变了的区域''重新计算 ESDF，未变化区域沿用旧值。后者让 ESDF 可以在机器人运动时实时维护，是目前 3D 导航系统的常用方案。

\section{点云地图与地图管理}\label{sec:14-pointcloud}

点云地图是最``原始''的地图：把每帧点云按位姿变换到世界系，拼起来即可。它信息最全（保留了全部几何细节），但也最重（一帧 10 万点、1000 帧就是 1 亿点，几十 GB）。所以实际系统必须做管理：

\begin{itemize}
  \item \textbf{体素降采样}：用第~\ref{ch:12-frontend}~章的体素栅格把点云稀疏化，通常按 $0.1\sim0.2\ \mathrm{m}$ 的体素保留质心，点云规模能降一到两个数量级。
  \item \textbf{动态点清理}：环境里的行人、车辆会留下``鬼影''。常用做法是记录每个点的``被再次观测次数''，只保留被多帧一致观测到的点；或者用可见性检查（从其他位姿看过去应该看得见，却看不见的点判为动态）。
  \item \textbf{分块存储}：把地图切成块（如 $10\ \mathrm{m}\times10\ \mathrm{m}$），每块独立保存与加载。这样定位时只加载机器人附近的块，内存可控。
  \item \textbf{地图格式}：PCD 保存点云（含强度、法向），PLY 保存网格。工程上还要保存位姿图、关键帧索引、地图坐标系原点等元信息——这就是``地图包''（map bundle）的概念，第~\ref{ch:23-zlocalization}~章会讲 Z-Localization 如何用先验地图包做定位。
\end{itemize}

\section{地图精度评价}\label{sec:14-eval}

建完图，怎么判断``好不好''？单看``好看不好看''不行，要量化。常用的四个指标如下。

\textbf{（1）覆盖率（coverage）。}重建出来的表面占真实表面的比例。做实验时用激光测距仪或全站仪测量若干条已知边长，与地图上量出的长度比较。

\textbf{（2）点到面误差（point-to-plane RMSE）。}把重建点云与真值点云（或 CAD 模型）配准后，计算每个重建点到真值表面的距离，取均方根：
\begin{equation}
\mathrm{RMSE}=\sqrt{\frac{1}{N}\sum_{i=1}^{N}\big(\bm{n}_i^{\top}(\bm{p}_i-\bm{q}_i)\big)^{2}},
\end{equation}
其中 $\bm{p}_i$ 是重建点，$\bm{q}_i$ 是真值表面上离它最近的点，$\bm{n}_i$ 是该处法向。它衡量``表面准不准''，通常要求与传感器噪声同量级（厘米级）。

\textbf{（3）地图一致性。}机器人绕一圈回到起点时，同一面墙会被观测两次。若两次重建的墙位置重合，地图是一致的；若错开，就会出现``双墙''（double wall）。一个实用的量化方法是测量同一面墙在重复观测处的厚度：理想情况下应等于单帧点云噪声，若明显更厚，说明回环后仍有残差。

\textbf{（4）地图可用性。}最终标准还是``下游任务能不能用''：全局规划能否找到路、局部避障会不会撞、定位重定位能否成功。工程上要拿真实场景跑一遍导航全流程，而不是只看点云好看不好看。

\begin{summary}
\begin{enumerate}
  \item 地图的形式由用途决定（表~\ref{tab:14-map-usage}）：定位要似然（点云、栅格），规划要代价（栅格、ESDF），重建要距离（TSDF）。
  \item 占据栅格存``被占据的概率''；由贝叶斯公式与观测独立性假设，可推出 log-odds 的加性更新 $l_t=l_{t-1}+l(m\mid\bm{z}_t)-l(m)$（式~\eqref{eq:14-logodds-update}），把连乘变成连加，并可用截断保持地图可修改。
  \item 命中/空闲增量由逆观测模型给出（常用 $p_{\mathrm{occ}}=0.7$、$p_{\mathrm{free}}=0.3$，增量 $\pm0.847$）；3 次命中后概率约 $0.93$，10 次后约 $0.9998$。
  \item 八叉树用``递归八分 + 同质剪枝''把 3D 占据地图的内存从百 MB 量级降到几 MB 量级，深度 $d=\lceil\log_2(L/r)\rceil$。
  \item TSDF 存带符号、截断的距离；多帧融合的递归式~\eqref{eq:14-tsdf-fusion} 正是加权最小二乘解~\eqref{eq:14-tsdf-closed} 的等价递推，权重常取 $\cos\theta/\|\bm{x}-\bm{o}\|^{2}$。
  \item ESDF 是到最近障碍的欧氏距离，梯度 $\nabla d$ 是单位向量、指向远离障碍的方向（式~\eqref{eq:14-esdf-grad}），因此可直接为规划提供``推离障碍''的方向；Voxblox 用增量更新维护它。
  \item 点云地图要配合降采样、动态点清理、分块存储与地图包管理；地图质量要用覆盖率、点到面 RMSE、一致性与下游可用性来评价。
\end{enumerate}


\end{summary}

\begin{problemset}
  \item 若把 $p_{\mathrm{occ}}$ 从 $0.7$ 改成 $0.9$，命中一次的 log-odds 增量是多少？同样的 ``3 次命中'' 后概率变成多少？这说明逆观测模型的取值如何影响地图``反应速度''。
  \item 式~\eqref{eq:14-logodds-update} 中假设了观测独立性。现实中同一面墙在相近位置被连续观测，这个假设成立吗？如果不成立，地图会出现什么偏差？
  \item 为什么八叉树的剪枝要以``8 个子节点状态相同''为条件？如果只比较其中 7 个会怎样？
  \item 在式~\eqref{eq:14-tsdf-fusion} 中，若某天开始权重不再增加（$W$ 达到上限），新观测的影响会怎样体现？这与``地图要能变旧''的准则有什么关系？
  \item 证明 $\nabla d(\bm{x})$ 是单位向量。提示：对 $d^{2}(\bm{x})=\norm{\bm{x}-\bm{y}^{\star}}^{2}$ 求导，并注意 $\bm{y}^{\star}$ 是最小值点。
  \item 你的机器人在一个 $200\ \mathrm{m}^{2}$ 的场地里建 2D 栅格图，分辨率 $5\ \mathrm{cm}$。估算纯占据概率存储需要多少内存；若改用 log-odds 并限制最大范围为 $120\ \mathrm{m}\times120\ \mathrm{m}$，又是多少？
\end{problemset}

\begin{reading}
\begin{enumerate}
  \item Thrun, Burgard, Fox，《Probabilistic Robotics》第 9 章\cite{thrun2005probabilistic}：占据栅格与 log-odds 更新的标准参考，本章的推导与它一致。
  \item Hornung 等，OctoMap\cite{hornung2013octomap}：八叉树的工程实现，重点看节点概率模型与剪枝策略。
  \item Curless \& Levoy\cite{curless1996tsdf}：TSDF 的原始论文，加权融合与表面提取的经典做法。
  \item Oleynikova 等，Voxblox\cite{oleynikova2017voxblox}：从 TSDF 增量构建 ESDF，是第~\ref{ch:20-local-planning}~章轨迹优化常用的地图。
\end{enumerate}

\end{reading}

% 本章使用 TikZ 的 arrows.meta / positioning 库（main.tex 与 tests 均已加载，此行幂等无害）
\usetikzlibrary{arrows.meta, positioning}

\chapter{2D 激光 SLAM：Cartographer}\label{ch:15-cartographer}

\chapterintro{Cartographer 是 Google 在 2016 年开源的激光 SLAM 系统\cite{hess2016cartographer}，也是 RoboMaster 赛场上最常用的 2D 建图方案之一。它把前面几章的知识组装成了一个完整可运行的工程：用``子图''组织地图，用``帧到子图匹配''做前端，用``分支定界''在离散位姿空间里快速找到最优位姿，用``位姿图优化''做全局闭环。本章会按这个顺序逐层拆开：先看系统架构与子图概念，再完整推导扫描匹配的得分函数、双线性插值及其雅可比，然后重点推导分支定界的``上界''为什么成立、为什么能剪枝，最后讲全局 SLAM、IMU/里程计的作用与关键参数。读完本章，你应该能对着 Cartographer 的配置文件说清每个参数在调什么。}

\section{Cartographer 的定位与系统架构}\label{sec:15-arch}

在 Cartographer 之前，2D 激光 SLAM 的主流有两种：Gmapping\cite{grisetti2007gmapping} 用粒子滤波（第~\ref{ch:10-global-loc}~章的粒子滤波在这里被用来``同时定位与建图''），Hector 用无回环的扫描匹配。三者的对比如表~\ref{tab:15-compare}。

\begin{table}[htbp]
  \centering
  \caption{三种 2D 激光 SLAM 方案的对比}\label{tab:15-compare}
  \footnotesize
  \begin{tabular}{p{0.14\linewidth}p{0.18\linewidth}p{0.20\linewidth}p{0.34\linewidth}}
    \toprule
    方案 & 核心方法 & 回环 & 特点与局限 \\
    \midrule
    Gmapping & 粒子滤波 + 扫描匹配 & 有（粒子重采样） & 小场景效果好；粒子数随地图增大而增多，大场景计算与内存吃紧 \\
    Hector & 高斯-牛顿扫描匹配 & 无 & 不需要里程计；但无回环，长距离必漂移 \\
    Cartographer & 子图 + 分支定界 + 位姿图 & 有（子图间约束） & 大场景实时、回环可靠；配置参数多，需调参 \\
    \bottomrule
  \end{tabular}
\end{table}

Cartographer 的整体数据流如图~\ref{fig:15-arch}，可以按第~\ref{ch:11-slam-framework}~章的语境一分成两半：
\begin{itemize}
  \item \textbf{Local SLAM（前端）}：把每一帧 scan 与当前子图做扫描匹配，求出相对位姿，然后把 scan 插入子图。它连续运行、频率高。
  \item \textbf{Global SLAM（后端）}：每当一个子图``完成''，就把它与其他历史子图做约束搜索（回环检测），并把所有约束丢进位姿图优化。它低频运行，负责全局一致。
\end{itemize}

\begin{figure}[htbp]
  \centering
  \begin{tikzpicture}[font=\small]
    \node[block] (s)     at (0,0)    {传感器\\{\footnotesize scan / IMU / odom}};
    \node[block2] (local) at (3.6,0)  {Local SLAM\\{\footnotesize 扫描匹配 + 插入}};
    \node[block] (sub)   at (7.2,0)  {子图\\{\footnotesize 概率栅格}};
    \node[block2] (gs)    at (7.2,-1.9) {Global SLAM\\{\footnotesize 约束 + 位姿图}};
    \node[block] (out)   at (3.6,-1.9) {优化后位姿\\{\footnotesize 地图 / TF}};
    \draw[arrow] (s) -- node[label,above]{scan} (local);
    \draw[arrow] (local) -- node[label,above]{插入} (sub);
    \draw[arrow] (sub) -- node[label,right]{完成的子图} (gs);
    \draw[arrow] (gs) -- node[label,below]{全局修正} (out);
    \draw[arrow, dashed] (out.north) -- node[label,left,align=center]{初值} (local.south);
    \node[label, anchor=west, align=left] at (9.3,0) {前端：\\局部一致};
    \node[label, anchor=west, align=left] at (9.3,-1.9) {后端：\\全局一致};
  \end{tikzpicture}
  \caption{Cartographer 的系统架构。传感器数据进入 Local SLAM 做帧到子图匹配并插入子图；子图完成后进入 Global SLAM 做约束构建与位姿图优化，优化结果反过来给前端提供位姿初值，并发布地图与 TF。}\label{fig:15-arch}
\end{figure}

\section{子图：地图的基本单位}\label{sec:15-submap}

\begin{definition}[子图]
子图（submap）是 Cartographer 中地图的基本单位：它是一张以概率栅格形式存储的局部地图，由连续插入的若干帧 scan 融合而成。每个子图有自己的局部坐标系，子图内的栅格存的是``被占据的 log-odds 值''（与第~\ref{ch:14-mapping}~章完全一致）。当插入的 scan 达到设定数量（参数 \lstinline{num_range_data}，默认 90）后，该子图``完成''，不再更新。
\end{definition}

用子图而不是一张全局大地图，有三个好处：第一，插入 scan 时只需查局部栅格，内存可控；第二，每个子图有独立的``原点''，避免全局地图越来越大导致的数值精度问题；第三，全局优化时只需要调整``子图位姿''（每个子图一个位姿），变量数大大减少——一张由 20 个子图拼成的地图，位姿图只有 20 多个节点，而不是几千帧。

子图的概率更新沿用第~\ref{ch:14-mapping}~章的 log-odds 公式：射线穿过的格子加``空闲''增量，命中格子加``占据''增量。查询时用
\begin{equation}
p(m)=\frac{1}{1+e^{-l(m)}},\qquad \mathrm{odds}(m)=\frac{p(m)}{1-p(m)}=e^{l(m)} .
\end{equation}

\section{扫描匹配：把 scan 对齐到子图}\label{sec:15-match}

\subsection{得分函数}

设当前 scan 有 $K$ 个点，在 scan 系下的坐标为 $\bm{h}_k=(h_{k,x},h_{k,y})$。要估计的是 scan 系到子图系的位姿 $\xi=(t_x,t_y,\theta)$，写成二维刚体变换
\begin{equation}
T_\xi\bm{h}_k=\bm{R}(\theta)\bm{h}_k+\bm{t}
=\begin{bmatrix}\cos\theta&-\sin\theta\\ \sin\theta&\cos\theta\end{bmatrix}\bm{h}_k+\begin{bmatrix}t_x\\ t_y\end{bmatrix}.
\end{equation}
把变换后的点拿到子图上查概率，得分定义为所有点概率之和：
\begin{equation}
S(\xi)=\sum_{k=1}^{K}M_{\mathrm{smooth}}\big(T_\xi\bm{h}_k\big),
\label{eq:15-score}
\end{equation}
其中 $M_{\mathrm{smooth}}$ 是子图栅格的双线性插值版本。$S$ 越大，说明``按这个位姿把 scan 放上去，点落在已占据格子里的越多''，位姿越可能正确。Cartographer 要找的就是
\begin{equation}
\xi^{\star}=\argmax_{\xi} S(\xi).
\label{eq:15-argmax}
\end{equation}

\subsection{双线性插值及其导数}

栅格是离散的，点 $\bm{p}=(x,y)$ 一般落在四个格子中间。最近邻插值会让得分函数变成阶梯状、处处不可导，梯度法就没法用；双线性插值给出连续可导的得分面，所以是标配。

设栅格分辨率为 $r$，把 $\bm{p}$ 换算成栅格坐标：
\begin{equation}
\frac{x}{r}=x_0+u,\qquad \frac{y}{r}=y_0+v,\qquad x_0=\left\lfloor\frac{x}{r}\right\rfloor,\ y_0=\left\lfloor\frac{y}{r}\right\rfloor,\ u,v\in[0,1).
\end{equation}
记四个相邻格子的概率为 $M_{00}=M(x_0,y_0)$、$M_{10}=M(x_0+1,y_0)$、$M_{01}=M(x_0,y_0+1)$、$M_{11}=M(x_0+1,y_0+1)$。先沿 $x$ 方向插值两次，再沿 $y$ 方向插值一次：
\begin{align}
M_{\mathrm{smooth}}(\bm{p})
&=(1-v)\Big[(1-u)M_{00}+uM_{10}\Big]+v\Big[(1-u)M_{01}+uM_{11}\Big]\notag\\
&=(1-u)(1-v)M_{00}+u(1-v)M_{10}+(1-u)vM_{01}+uvM_{11}.
\label{eq:15-bilinear}
\end{align}
式~\eqref{eq:15-bilinear} 就是双线性插值：四个角的权重是``到对角的距离''之积，离哪个角越近权重越大。

对 $\bm{p}$ 求偏导时要注意链式法则：$\partial u/\partial x=1/r$、$\partial v/\partial y=1/r$，而 $x_0,y_0$ 在格子内部是常数。于是
\begin{align}
\frac{\partial M_{\mathrm{smooth}}}{\partial x}
&=\frac{1}{r}\Big[(1-v)\big(M_{10}-M_{00}\big)+v\big(M_{11}-M_{01}\big)\Big],
\label{eq:15-dMdx}\\
\frac{\partial M_{\mathrm{smooth}}}{\partial y}
&=\frac{1}{r}\Big[(1-u)\big(M_{01}-M_{00}\big)+u\big(M_{11}-M_{10}\big)\Big].
\label{eq:15-dMdy}
\end{align}
这两个式子是后面雅可比的``上游''：它们告诉你``点移动一点，得分变多少''。

\subsection{非线性最小二乘与雅可比}

把式~\eqref{eq:15-argmax} 的求最大值改写为求最小值：定义残差
\begin{equation}
r_k(\xi)=1-M_{\mathrm{smooth}}\big(T_\xi\bm{h}_k\big)\in\R,
\end{equation}
则 $\min_\xi\sum_k r_k^2$ 与 $\max_\xi S$ 等价（常数 $K$ 不影响极值点）。对残差做高斯-牛顿迭代，需要残差对位姿的雅可比。链式法则分成两段：\textbf{点对位姿}、\textbf{得分对点}。

第一段：变换后的点对位姿求导。记 $\bm{R}'\bm{h}_k$ 为旋转矩阵对 $\theta$ 求导再作用在点上的结果：
\begin{equation}
\frac{\partial (T_\xi\bm{h}_k)}{\partial t_x}=\begin{bmatrix}1\\0\end{bmatrix},\qquad
\frac{\partial (T_\xi\bm{h}_k)}{\partial t_y}=\begin{bmatrix}0\\1\end{bmatrix},\qquad
\frac{\partial (T_\xi\bm{h}_k)}{\partial\theta}=\begin{bmatrix}-(h_x\sin\theta+h_y\cos\theta)\\ h_x\cos\theta-h_y\sin\theta\end{bmatrix}.
\end{equation}
注意最后一项可以写成 $\bm{J}_2\bm{R}(\theta)\bm{h}_k$，其中
\begin{equation}
\bm{J}_2=\begin{bmatrix}0&-1\\1&0\end{bmatrix}
\end{equation}
是``逆时针转 $90^{\circ}$''的矩阵。合并成 $2\times3$ 的矩阵：
\begin{equation}
\frac{\partial (T_\xi\bm{h}_k)}{\partial\xi}
=\begin{bmatrix}\bm{I}_2 & \bm{J}_2\bm{R}(\theta)\bm{h}_k\end{bmatrix}\in\R^{2\times3}.
\label{eq:15-dT-dxi}
\end{equation}

第二段：把得分梯度式~\eqref{eq:15-dMdx}、式~\eqref{eq:15-dMdy} 与式~\eqref{eq:15-dT-dxi} 相乘，得到残差对位姿的 $1\times3$ 雅可比：
\begin{equation}
\bm{J}_k=\frac{\partial r_k}{\partial\xi}
=-\Big[\frac{\partial M}{\partial x},\ \frac{\partial M}{\partial y}\Big]
\begin{bmatrix}\bm{I}_2 & \bm{J}_2\bm{R}\bm{h}_k\end{bmatrix}\in\R^{1\times3}.
\label{eq:15-jac}
\end{equation}
它的每一列分别对应``平移 $x$、平移 $y$、旋转 $\theta$ 一点点，得分变多少''。有了 $\bm{J}_k$，高斯-牛顿方程就是（残差是一维，信息矩阵取 1）：
\begin{equation}
\Big(\sum_{k}\bm{J}_k^{\top}\bm{J}_k\Big)\delta\bm{\xi}
=-\sum_{k}\bm{J}_k^{\top}r_k .
\label{eq:15-gn}
\end{equation}
解出 $\delta\bm{\xi}$ 后按 $\xi\leftarrow\xi\oplus\delta\bm{\xi}$ 更新（二维时 $\oplus$ 就是直接相加），重复几轮即可收敛。

\begin{keybox}{扫描匹配的三个``为什么''}
\textbf{为什么用双线性插值？}为了让得分函数连续可导；最近邻会让 GN 拿不到梯度。\\
\textbf{为什么只迭代几轮？}匹配只是在给定位姿附近做局部精化，初值由里程计/IMU/上一次位姿给出；跑太多轮既慢又可能被局部极值带偏。\\
\textbf{为什么还要分支定界？}GN 只能做局部精化，初值差得多时（比如刚开机或回环搜索）会掉进错误的局部极小；此时需要全局搜索，而全局搜索靠分支定界加速。
\end{keybox}

\section{分支定界：在离散位姿空间里找全局最优}\label{sec:15-bnb}

\subsection{问题与离散化}

回环检测或重定位时，我们不知道大致位姿，只有``可能在这一带''的先验。此时要解的是全局的式~\eqref{eq:15-argmax}。把位姿空间离散化：平移按子图分辨率 $r$ 取整数格，旋转按固定角步长 $\delta_\theta$ 取整。于是在一个窗口内，候选位姿数量是
\begin{equation}
N=N_x\cdot N_y\cdot N_\theta ,
\end{equation}
典型的回环搜索窗口 $N_x=N_y=100$、$N_\theta=180$（约 $2^{\circ}$ 一档），总数约 $1.8\times10^{6}$。每个候选都要对 $K\sim10^{3}$ 个点查一次子图，直接算就是 $10^{9}$ 次查表，实时不可能。

分支定界（Branch and Bound）是解决``在离散空间里找最大值''的经典方法。它把搜索空间组织成一棵树：根节点代表整个搜索窗口，每个节点代表一个``位姿子集''，节点一分为四（平移二分 $\times$ 旋转二分）。对每个节点，我们不计算精确得分，而是计算一个\textbf{上界}——该子集里任何位姿的得分都不会超过这个上界。于是：若某节点的上界低于``目前已经找到的最好得分''，整棵子树直接剪掉，一个叶子都不用算；否则继续往下分。

\begin{definition}[分支定界]
分支定界（Branch and Bound）是一类求解离散最优化问题的方法：把候选解集合递归地划分成子集（分支），并给每个子集计算一个``其上所有解的目标值都不会超过''的上界（定界）。搜索时优先处理上界高的子集，一旦某子集的上界低于当前最优解，就剪掉该子集。若上界是紧的（接近真实最大值），可以在远少于穷举的节点数内找到最优解。
\end{definition}

\subsection{上界的推导}

关键问题：怎么给一个位姿子集算上界，而且算得便宜？答案是用\textbf{预计算的多分辨率栅格}。

设一个节点代表的位姿子集是：平移在以 $\xi_c$ 为中心的半宽 $c_h$ 个格子内，旋转在半宽 $\delta_h$ 内。对子集里的任意位姿 $\xi$ 与任意 scan 点 $\bm{h}_k$，有
\begin{align}
\big\|T_\xi\bm{h}_k-T_{\xi_c}\bm{h}_k\big\|
&=\big\|\bm{R}(\theta)\bm{h}_k+\bm{t}-\bm{R}(\theta_c)\bm{h}_k-\bm{t}_c\big\|\notag\\
&\le\underbrace{\big\|\bm{R}(\theta)-\bm{R}(\theta_c)\big\|\cdot\|\bm{h}_k\|}_{\text{旋转误差}}
+\underbrace{\|\bm{t}-\bm{t}_c\|}_{\text{平移误差}} .
\label{eq:15-bound-pos}
\end{align}
二维旋转有个好性质：两旋转矩阵之差的作用可以界住角度差，
\begin{equation}
\big\|\bm{R}(\theta)-\bm{R}(\theta_c)\big\|_2
=2\left|\sin\frac{\theta-\theta_c}{2}\right|\le|\theta-\theta_c| .
\label{eq:15-rot-bound}
\end{equation}
（用 $\sin\alpha\le\alpha$ 即得。）由节点定义，$|\theta-\theta_c|\le\delta_h$，$\|\bm{t}-\bm{t}_c\|\le\sqrt{2}\,c_h r$（二维平移偏差最大值是对角线），代回式~\eqref{eq:15-bound-pos}：
\begin{equation}
\big\|T_\xi\bm{h}_k-T_{\xi_c}\bm{h}_k\big\|
\le\delta_h\|\bm{h}_k\|+\sqrt{2}\,c_h r
\;\triangleq\;\rho_h(\bm{h}_k).
\label{eq:15-radius}
\end{equation}
也就是说：\textbf{子集内任意位姿下，第 $k$ 个点落到的位置，一定在以 $T_{\xi_c}\bm{h}_k$ 为圆心、$\rho_h(\bm{h}_k)$ 为半径的圆内}。

于是对每个深度 $h$，预先算一张``膨胀栅格''：把子图里每个格子的值替换成它周围半径 $\rho_h$ 范围内的最大值，
\begin{equation}
M^{h}_{\mathrm{precomp}}(\bm{p})=\max_{\|\bm{q}-\bm{p}\|\le\rho_h}M_{\mathrm{smooth}}(\bm{q}).
\label{eq:15-precomp}
\end{equation}
由式~\eqref{eq:15-radius}，子集内任意位姿下 $M_{\mathrm{smooth}}(T_\xi\bm{h}_k)\le M^{h}_{\mathrm{precomp}}(T_{\xi_c}\bm{h}_k)$。对所有点求和：
\begin{equation}
S(\xi)=\sum_k M_{\mathrm{smooth}}(T_\xi\bm{h}_k)
\le\sum_k M^{h}_{\mathrm{precomp}}\big(T_{\xi_c}\bm{h}_k\big)
\;\triangleq\;\overline{S}_h(\xi_c).
\label{eq:15-upperbound}
\end{equation}
右边只依赖节点中心位姿 $\xi_c$ 与预计算栅格，查表即可算出——这就是\textbf{上界}。

两个必须讲清的性质：
\begin{enumerate}
  \item \textbf{上界是``真上界''}：式~\eqref{eq:15-upperbound} 是严格推导出来的不等式，不是启发式，所以剪枝不会漏掉最优解。
  \item \textbf{上界随深度单调收紧}：深度越大，$c_h$ 与 $\delta_h$ 越小，由式~\eqref{eq:15-radius} 得到的 $\rho_h$ 越小；$\rho_h$ 越小，式~\eqref{eq:15-precomp} 里取 $\max$ 的范围越小，$M^h_{\mathrm{precomp}}$ 逐点越小。所以\textbf{父节点的上界一定不小于子节点的上界}，这保证搜索树自顶向下``越来越紧''，剪枝才有意义。
\end{enumerate}

\subsection{搜索策略与伪代码}

搜索用\textbf{优先队列}（每次取出上界最大的节点）最直观：由于父节点的上界不小于子孙，第一个被取出并展开到叶子的节点，其真实得分一定不低于队列里其他所有节点的上界，因此\textbf{第一个到达的叶子就是全局最优解}。这就是``最优优先''的分支定界。

\begin{figure}[htbp]
  \centering
  \begin{tikzpicture}[font=\footnotesize]
    \tikzset{sn/.style={circle, draw=structurecolor, thick, fill=structurecolor!10,
                        minimum size=8mm, inner sep=1pt}}
    \node[sn] (r) at (0,0) {9.1};
    \node[sn] (a) at (-3.6,-1.5) {8.6};
    \node[sn] (b) at (-1.2,-1.5) {8.9};
    \node[sn] (c) at (1.2,-1.5) {7.8};
    \node[sn] (d) at (3.6,-1.5) {8.4};
    \draw (r) -- (a); \draw (r) -- (b); \draw (r) -- (c); \draw (r) -- (d);
    \node[sn] (b1) at (-2.4,-3.0) {8.5};
    \node[sn, fill=structurecolor!35] (b2) at (-1.2,-3.0) {8.8};
    \node[sn] (b3) at (0.0,-3.0) {8.7};
    \node[sn] (c1) at (1.2,-3.0) {7.7};
    \draw (b) -- (b1); \draw (b) -- (b2); \draw (b) -- (b3);
    \draw (c) -- (c1);
    % pruned subtree of d
    \draw[red!60!black, dashed] (d) -- (3.6,-3.0);
    \node[sn, draw=red!60!black, fill=red!10] at (3.6,-3.0) {$\times$};
    \node[label, align=center] at (-1.2,-3.9) {最优叶子：得分 $8.8$};
    \node[label, align=center] at (3.6,-3.9) {上界 $7.8$ 低于当前最优\\$8.5$：整棵剪掉};
  \end{tikzpicture}
  \caption{分支定界的示意。圆圈内数字是该节点的上界。优先队列先展开上界 9.1 的根，再展开 8.9 的分支；当 8.8 的叶子被找到后，所有上界小于 8.8 的节点（如 7.8、8.4、8.5）整棵剪掉。上界必须``真''，否则会剪掉最优解。}\label{fig:15-bnb}
\end{figure}

\begin{pseudo}{分支定界搜索最优位姿}
  \item 输入：先验位姿 $\xi_0$、搜索窗口、预计算的多层膨胀栅格
  \item 根节点 $\leftarrow$ 整个窗口，用式~\eqref{eq:15-upperbound} 算上界
  \item 优先队列 $Q\leftarrow\{\text{根节点}\}$；最优得分 $\mathrm{best}\leftarrow-\infty$
  \item \textbf{while} $Q$ 非空 \textbf{do}
  \item \quad 从 $Q$ 取出上界最大的节点 $n$
  \item \quad \textbf{if} $n.\text{上界}<\mathrm{best}$ \textbf{then} 剪枝，\textbf{continue}
  \item \quad \textbf{if} $n$ 是叶子 \textbf{then} $\mathrm{best}\leftarrow S(n)$，$\xi^{\star}\leftarrow n.\text{位姿}$，\textbf{continue}
  \item \quad 把 $n$ 划分成子节点；对每个子节点算上界并放入 $Q$
  \item \textbf{end while}
  \item \textbf{return} $\xi^{\star}$
\end{pseudo}

\subsection{为什么能实时}

遍历所有候选是 $O(N)$，其中 $N=N_xN_yN_\theta$ 可达 $10^{6}$；每个候选算得分又要 $O(K)$，总代价 $O(NK)$。分支定界不会遍历所有叶子：由于上界随深度收紧得很快，大部分分支在高层就被剪掉。实际运行中，一次回环搜索通常只访问 $10^{3}\sim10^{4}$ 个节点，相比 $10^{6}$ 个叶子少了两三个数量级——这就是 Cartographer 能``在后台实时做回环''的原因。当然，最坏情况下（得分面非常平坦、上界很松）它仍会退化为穷举，所以工程上还要限制搜索窗口大小与层数。

\section{全局 SLAM：把子图连成一张图}\label{sec:15-global}

Local SLAM 只保证``子图内部一致''，子图与子图之间会漂移。Global SLAM 要把它们缝起来：

\begin{enumerate}
  \item \textbf{约束构建（constraint building）}：当有新的 scan 或新子图完成时，在``历史子图''里找可能重叠的那些，用分支定界在粗分辨率上搜索候选相对位姿。若该位姿下的匹配得分超过阈值（参数 \lstinline{min_score}），就把它作为一条``子图—scan 约束''加入位姿图。
  \item \textbf{回环检测}：本质上就是``在子图集合上做约束搜索''。Cartographer 不是拿两帧点云直接比，而是拿``当前 scan 对历史子图''匹配，这比 Scan Context 那样的描述子更准（直接优化得分），代价是搜索更重，所以要靠分支定界加速。
  \item \textbf{位姿图优化}：节点是各子图的位姿，边是约束（含协方差），用 Ceres\cite{hess2016cartographer} 做稀疏最小二乘求解（第~\ref{ch:13-backend}~章）。优化后，所有子图位姿被修正，地图重新拼合。
\end{enumerate}

Cartographer 的位姿图规模很小：一张地图可能只有几十到几百个子图，每条边只连两个节点，海森矩阵极其稀疏，一次优化通常几十毫秒。这与前端每帧的扫描匹配（几毫秒）配合起来，整体可以实时运行。

\section{IMU 与里程计的作用}\label{sec:15-imu}

2D 情形下，IMU 的价值不像 3D LIO 那么大，但仍然有用：
\begin{itemize}
  \item \textbf{提供旋转初值}：陀螺仪积分给出 yaw 的短时变化，可以作为扫描匹配的旋转初值，把三维搜索（$x,y,\theta$）里 $\theta$ 的范围缩小。
  \item \textbf{提供重力方向}：如果雷达不是水平安装，重力方向可以帮助把 scan 投影到水平面；纯 2D 系统里用得少，3D 版本则必须用。
  \item \textbf{给位姿图加约束}：IMU 的航向角可以作为一条``绝对朝向''约束加进优化，抑制长走廊里的旋转漂移。
\end{itemize}
里程计（轮速）的作用更直接：它给出相邻帧之间的平移与旋转估计，作为扫描匹配的初值。没有里程计时，Cartographer 用``上一帧位姿''或匀速模型预测作初值；有里程计时初值更准，匹配更稳。三者都必须做\textbf{时间同步}（第~\ref{ch:03-sensors}~章）：scan、IMU、odom 的时间戳若差几十毫秒，在机器人快速运动时会直接变成位姿误差。

\section{关键参数与调参}\label{sec:15-tuning}

Cartographer 的配置文件（\lstinline{.lua}）参数很多，真正影响效果的是表~\ref{tab:15-params} 这几个。

\begin{table}[htbp]
  \centering
  \caption{Cartographer 关键参数及其影响}\label{tab:15-params}
  \footnotesize
  \begin{tabular}{p{0.36\linewidth}p{0.12\linewidth}p{0.44\linewidth}}
    \toprule
    参数 & 默认值 & 调大 / 调小的后果 \\
    \midrule
    \lstinline{submaps.resolution} & 0.05 m & 调小：更精细但内存与计算增大；调大：省资源但细节丢失 \\
    \lstinline{submaps.num_range_data} & 90 & 调大：子图更大、单图更稳、回环机会少；调小：子图碎、约束多、优化慢 \\
    \lstinline{constraint_builder.min_score} & 0.55 & 调大：只接受高置信回环，漏检多；调小：回环多但误检风险高 \\
    \lstinline{optimization_problem.huber_scale} & 1.0 & 调小：对错边更鲁棒；调大：更信任所有约束 \\
    \lstinline{TRAJECTORY_BUILDER_2D.voxel_filter_size} & 0.025 m & 调大：降采样更狠、更快；调小：保留细节但更慢 \\
    \lstinline{POSE_GRAPH.optimize_every_n_nodes} & 90 & 调小：优化频繁、更跟手但更耗 CPU；调大：省 CPU 但闭环延迟 \\
    \bottomrule
  \end{tabular}
\end{table}

调参的顺序建议是：先保证\textbf{数据质量}（时间同步、外参、雷达去畸变），再调\textbf{分辨率与降采样}（决定计算量），最后调\textbf{回环相关参数}（\lstinline{min_score}、优化频率）。顺序反了会陷入``调一个坏一个''的循环。

\section{优缺点与适用场景}\label{sec:15-pros}

\textbf{优点：}回环可靠（子图 + 分支定界 + 位姿图，长距离建图能闭环）；实时性好（分支定界把全局搜索压到可接受的时间）；工程成熟（ROS/ROS 2 集成好，配置项清晰）；只依赖 2D 激光也能跑（可选 IMU/odom），成本低。

\textbf{局限：}第一，2D 假设机器人在地面单层运动，多层场地、坡道、空中场景都不适用，这时要换 3D SLAM（第~\ref{ch:16-lio-evolution}~章与第~\ref{ch:17-advanced-lio}~章）；第二，长走廊、空旷场地存在几何退化（第~\ref{ch:12-frontend}~章），单靠激光会漂移，必须靠 IMU/odom 兜底；第三，参数较多，场景变化大时需要重新调参；第四，子图是``局部''的，超大场景（几万平方米）需要分块建图再拼接。

对 RoboMaster 场景：赛场是单层、结构相对规整、有围挡与固定道具，2D Cartographer 是非常合适的选择；若场地有坡道或需要 3D 避障，则应采用第~\ref{ch:16-lio-evolution}~章起的 LIO 方案。

\begin{summary}
\begin{enumerate}
  \item Cartographer 由 Local SLAM（前端：帧到子图匹配 + 插入）与 Global SLAM（后端：约束构建 + 位姿图优化）组成（图~\ref{fig:15-arch}），是第~\ref{ch:11-slam-framework}~章语境一的典型实现。
  \item 子图是地图的基本单位，本质是一张概率栅格（log-odds），插入 \lstinline{num_range_data} 帧后完成；用子图位姿做全局优化，变量数从``帧数''降到``子图数''。
  \item 扫描匹配的得分是 $S(\xi)=\sum_k M_{\mathrm{smooth}}(T_\xi\bm{h}_k)$（式~\eqref{eq:15-score}），双线性插值式~\eqref{eq:15-bilinear} 让得分连续可导；其梯度由式~\eqref{eq:15-dMdx}、式~\eqref{eq:15-dMdy} 给出，残差雅可比为式~\eqref{eq:15-jac}，用高斯-牛顿式~\eqref{eq:15-gn} 迭代求解。
  \item 分支定界把位姿空间组织成树，用预计算膨胀栅格式~\eqref{eq:15-precomp} 给出上界式~\eqref{eq:15-upperbound}；上界严格成立、且随深度单调收紧，因此剪枝不会漏掉最优解，实际访问节点数比穷举少两三个数量级。
  \item 全局 SLAM 用约束搜索（分支定界）找子图间回环，用 Ceres 做位姿图优化；IMU/odom 提供初值与旋转约束，时间同步是前提。
  \item 参数调优顺序：数据质量 $\rightarrow$ 分辨率/降采样 $\rightarrow$ 回环与优化频率；2D 单层场景适用，多层与坡道需换 LIO。
\end{enumerate}


\end{summary}

\begin{problemset}
  \item 式~\eqref{eq:15-bilinear} 中，若点恰好落在格子中心（$u=v=0.5$），四个格子的权重各是多少？这个结果符合直觉吗？
  \item 为什么扫描匹配的残差定义为 $1-M$ 而不是 $M$？两者对高斯-牛顿的迭代有什么影响？
  \item 若把式~\eqref{eq:15-radius} 中的平移偏差从 $\sqrt{2}c_h r$ 改成 $2c_h r$，上界会变松还是变紧？变松之后剪枝效率如何变化？
  \item 分支定界的上界如果取成``节点中心位姿的真实得分''，会有什么后果？为什么说``上界必须真''？
  \item Cartographer 的位姿图节点是子图而不是每帧 scan，这样做省下了什么？又牺牲了什么？
  \item 在长走廊里建图，Cartographer 的哪些环节最容易出问题？你会调整表~\ref{tab:15-params} 中的哪些参数、或者引入什么传感器？
\end{problemset}

\begin{reading}
\begin{enumerate}
  \item Hess 等，Cartographer 论文\cite{hess2016cartographer}：本章的原始文献，重点读第 III 节的 ``Local SLAM''、第 IV 节的 ``Branch-and-Bound'' 与第 V 节的 ``Global SLAM''。
  \item Grisetti 等，Gmapping\cite{grisetti2007gmapping}：粒子滤波 2D SLAM 的代表，用来对照表~\ref{tab:15-compare}。
  \item 第~\ref{ch:14-mapping}~章：子图的概率栅格更新与 log-odds，本章直接沿用。
  \item Cartographer 官方文档与 \lstinline{cartographer_ros} 配置示例：调参时对照表~\ref{tab:15-params} 阅读。
\end{enumerate}

\end{reading}

\chapter{3D 激光惯性里程计：从 LOAM 到 LIO}\label{ch:16-lio-evolution}

\usetikzlibrary{arrows.meta, positioning, calc, shapes.geometric}

\chapterintro{第~\ref{ch:15-cartographer}~章的 Cartographer 只处理机器人在平面内运动的情形（状态只有 $x,y,\theta$ 三个自由度）。可一旦机器人加上云台、上下坡、被撞飞，或者换成无人机、机械臂，运动就是 6 自由度的。本章把问题升到三维，回答三个问题：为什么 3D 激光必须要和 IMU 一起用？一帧点云内部的运动畸变到底有多严重、怎么消掉？从 2014 年的 LOAM 到 2023 年的 Point-LIO，这十年里每个算法究竟解决了前一个算法的什么毛病？读完本章，你应该能看懂任意一个 LIO 开源包的数据流，并能判断"我这个场景该选哪一个"。本章只讲演化脉络与共性原理，FAST-LIO2 和 Point-LIO 的逐行数学推导放在第~\ref{ch:17-advanced-lio}~章。}

\section{为什么需要激光惯性里程计}\label{sec:16-why}

\subsection{2D 激光 SLAM 的三条隐含假设}

Cartographer~\cite{hess2016cartographer} 之所以能把计算量压到嵌入式平台可接受的水平，靠的是三条假设：机器人只在水平面内运动、激光只有一个扫描平面、地图用一个二维栅格就能表达。这三条在 RoboMaster 的场地上其实已经摇摇欲坠：

\begin{enumerate}
  \item 哨兵和步兵会过坡道、飞坡、被撞得整个车身翻滚，云台的 pitch 角也可能连续变化，此时"$z$ 坐标恒定、roll/pitch 恒为零"直接失效；
  \item 单线激光只切出场景的一个水平薄片，两块高度不同的挡板在 2D 地图上完全一样，机器人无法区分"前面是 20\,cm 高的坎"还是"前面是三米高的墙"；
  \item 换成 3D 激光后，一帧点云动辄 $10^4\sim10^5$ 个点，比 2D 单线激光多两个数量级，Cartographer 那套"子图 + 分支定界"的暴力搜索代价立刻变得不可接受。
\end{enumerate}

所以 3D 场景需要 6 自由度的位姿 $(R,\bm{p})\in\SE(3)$（李群的定义见第~\ref{ch:02-math}~章），计算量又要求我们不能对全部点做穷举搜索。这两条约束一起，把 3D 激光里程计推向了"用 IMU 提供运动先验、用点云提供几何约束"的路子。

\subsection{纯激光里程计的三种失效模式}

回忆第~\ref{ch:07-lidar-odom}~章的扫描匹配：给定两帧点云，通过最小化点到点或点到面残差求相对位姿。它的信息矩阵是 $\bm{H}^{\top}\bm{H}$（$\bm{H}$ 是残差对位姿的雅可比堆叠），当这个矩阵在某个方向上接近奇异时，那个方向上的位姿就解不准，这叫\textbf{几何退化}（degeneracy）。具体有三种典型情形：

\begin{itemize}
  \item \textbf{长走廊、隧道、平行墙}：沿走廊轴线方向平移，点云外观几乎不变，$\bm{H}^{\top}\bm{H}$ 沿该方向的特征值接近零，机器人会以为"自己没动"，位置估计沿轴线无界漂移；
  \item \textbf{开阔场地或单一平面}：所有点几乎落在一个平面上，只能约束垂直于平面的平移和两个转角，沿平面内的两个平移和绕平面法向的转角都约束不足；
  \item \textbf{剧烈运动与稀疏场景}：相邻两帧重叠区域太小，或者机器人转得太快导致两帧几乎没有共视点，匹配的初值差，ICP 类算法会收敛到错误的局部极小。
\end{itemize}

第二种失效其实不是几何问题，而是时间问题：3D 激光的"一帧"并不是一个瞬间，而是几十到一百毫秒的连续采集，帧内机器人在动，点云自己就被"拉花"了。这就是下一节要讲的去畸变。

\subsection{IMU 补上了什么}

IMU 的特点是频率高（$100\sim1000$\,Hz）、短时间内的角速度和比力测量准，但积分会漂移。这两个特点和激光恰好互补，见表~\ref{tab:16-complement}。

\begin{table}[htbp]
  \centering
  \caption{激光与 IMU 的互补性}\label{tab:16-complement}
  \small
  \begin{tabular}{lll}
    \toprule
    特性 & 3D 激光 & IMU \\
    \midrule
    输出频率 & $10$\,Hz（一帧 $100$\,ms） & $100\sim1000$\,Hz \\
    擅长的时间尺度 & $>0.1$\,s 的绝对几何 & $<0.1$\,s 的高动态运动 \\
    提供的量 & 场景几何（绝对约束） & 角速度、比力（增量约束） \\
    随时间的误差 & 缓慢漂移、可回环消除 & 二重积分，秒级发散 \\
    典型失效 & 几何退化、帧内畸变、快速旋转 & 饱和、冲击、温度漂移 \\
    \bottomrule
  \end{tabular}
\end{table}

用两个数字感受一下 IMU 靠不住的程度：一块常见 MEMS 陀螺的零偏若为 $0.01$\,rad/s（约 $0.57^{\circ}$/s），纯积分 $10$\,s 就会产生 $0.01\times10\times180/\pi\approx5.7^{\circ}$ 的姿态误差；加速度计零偏若为 $0.02\,\mathrm{m/s^2}$，同样的 $10$\,s 内位置误差为 $\frac{1}{2}\times0.02\times10^2=1$\,m。反过来，激光在 $100$\,ms 内几乎不可能漂移，所以把两者放进同一个估计器，正好让 IMU 负责"帧内"、激光负责"帧间"。

\begin{definition}[激光惯性里程计 LIO]
以 3D 激光点云和 IMU 测量为输入，在一个估计器中联合估计载体位姿 $T_{WB}$（以及速度、IMU 零偏等辅助状态）并增量维护地图的算法，称为\textbf{激光惯性里程计}（LiDAR-Inertial Odometry，LIO）。若激光与 IMU 的原始测量被送进同一个估计器的同一次状态更新中，称为\textbf{紧耦合}；若先由激光独立的扫描匹配算出位姿，再把这个位姿当作观测去修正 IMU 的递推结果，则称为\textbf{松耦合}。两者的数学差别在 \S\ref{sec:16-tight} 展开。
\end{definition}

\begin{keybox}{一句话记住 LIO 的价值}
纯激光里程计在几何退化处"不知道自己错了"，纯 IMU 在几秒后"必然错"。LIO 把 IMU 的高频运动先验和激光的绝对几何约束放进同一个概率模型，使二者互相纠错：退化方向上靠 IMU 撑住，IMU 漂移靠激光拉回。
\end{keybox}

\section{去畸变：把一帧点云对齐到同一时刻}\label{sec:16-deskew}

\subsection{问题有多严重}

以常见的机械式 16 线激光为例：转速 $10$\,Hz，即一帧 $100$\,ms；一帧内每线约 $1800$ 个点，相邻点间隔约 $55\,\mu$s。也就是说，一帧点云里第一个点和最后一个点的采集时刻相差 $100$\,ms。若机器人此时角速度为 $90^{\circ}$/s，这 $100$\,ms 内姿态变化 $9^{\circ}$；若点距离 $r=5$\,m，则帧首帧尾的点在空间上会被错开约 $r\sin9^{\circ}\approx0.78$\,m。如果机器人同时以 $2$\,m/s 前进，平移误差再叠加 $0.2$\,m。而地图体素分辨率通常只有 $0.1\sim0.2$\,m，这样的畸变会让整个配准直接崩溃。

\begin{definition}[运动畸变与去畸变]
激光雷达一帧点云中的每个点都带有自己的采集时刻 $t_i$（相对扫描起始时刻的偏移）。由于载体在帧内运动，若把整帧点云当作同一时刻 $t_{\mathrm{ref}}$ 的刚体扫描来处理，就会引入系统性形变，称为\textbf{运动畸变}。按每个点自己的时刻把它的坐标变换到同一参考时刻 $t_{\mathrm{ref}}$ 的过程，称为\textbf{去畸变}（deskewing / motion compensation）。
\end{definition}

\subsection{去畸变公式}

设 $t_i$ 时刻点 $i$ 在雷达系中的坐标为 $\bm{p}_i^{L}$，该时刻雷达到世界系的位姿为 $T_{WL}(t_i)=(R(t_i),\bm{p}(t_i))$。那么该点在世界上真正的位置是
\begin{equation}\label{eq:16-world-point}
  \bm{p}_i^{W}=R(t_i)\,\bm{p}_i^{L}+\bm{p}(t_i).
\end{equation}
去畸变就是把同一批世界点重新表达在参考时刻的雷达系中。把式~\eqref{eq:16-world-point} 反解，得到
\begin{equation}\label{eq:16-deskew}
  \bm{p}_i^{L_{\mathrm{ref}}}=T_{WL}(t_{\mathrm{ref}})^{-1}\bm{p}_i^{W}
  =R(t_{\mathrm{ref}})^{\top}\Bigl(R(t_i)\bm{p}_i^{L}+\bm{p}(t_i)-\bm{p}(t_{\mathrm{ref}})\Bigr).
\end{equation}
式~\eqref{eq:16-deskew} 的几何含义很直白：先把点搬回世界系，再用参考时刻的位姿搬进参考帧。若帧内没有运动（$R(t_i)=R(t_{\mathrm{ref}})$，$\bm{p}(t_i)=\bm{p}(t_{\mathrm{ref}})$），它就退化成恒等变换，说明去畸变只在"帧内动"的情况下才需要。

\subsection{匀速模型与 IMU 辅助模型}

式~\eqref{eq:16-deskew} 需要帧内任意时刻的位姿 $T_{WL}(t)$，而我们只在一个帧边界上知道它。最简单的近似是\textbf{匀速模型}：假设帧内线速度和角速度恒定，用前后两帧的位姿插值：
\begin{equation}\label{eq:16-constant-velocity}
  R(t)=R(t_0)\exp\!\Bigl(s\,\bigl[\mathrm{Log}(R(t_0)^{\top}R(t_1))\bigr]_\times\Bigr),\qquad
  \bm{p}(t)=(1-s)\bm{p}(t_0)+s\,\bm{p}(t_1),\quad s=\frac{t-t_0}{t_1-t_0},
\end{equation}
其中 $\mathrm{Log}(\cdot)$ 是 $\SO(3)$ 到旋转向量的对数映射（见第~\ref{ch:02-math}~章），$s$ 是归一化时间。式~\eqref{eq:16-constant-velocity} 就是 $\SE(3)$ 上的线性插值：平移线性、旋转走测地线（SLERP）。

匀速模型在加速度不大时够用。它的误差量级可以估出来：帧内加速度 $a=5\,\mathrm{m/s^2}$、时间 $0.1$\,s 时，位置上的二阶误差约 $\frac12 a \Delta t^2=2.5$\,cm，可以接受；但角加速度大或者有撞击时就不行了。

更可靠的做法是用 IMU 把 $T_{WL}(t)$ 真正积分出来。对 IMU 相邻两个采样时刻 $t_\alpha,t_\beta$ 之间的测量做积分（积分公式在第~\ref{ch:06-imu}~章与第~\ref{ch:17-advanced-lio}~章给出），得到一串离散位姿 $\{T_{WL}(t_\alpha)\}$；对落在 $[t_\alpha,t_\beta]$ 内的点，再按式~\eqref{eq:16-constant-velocity} 在相邻两个 IMU 位姿之间插值即可。至于更精细的连续时间表示（例如把整条轨迹参数化为 B 样条再做三次插值），DLIO 等工作专门讨论了这一点~\cite{chen2023dlio}。实际工程中按"分段线性"处理已经能覆盖绝大多数场景。

\begin{note}
去畸变只需要帧内\textbf{相对}运动，不需要绝对位姿。IMU 零偏在 $100$\,ms 内造成的姿态漂移通常远小于 $0.1^{\circ}$，所以去畸变对零偏不敏感；这也是为什么即使零偏还没收敛，FAST-LIO2 一类算法也能先把点云去畸变。
\end{note}

\begin{figure}[htbp]
  \centering
  \begin{tikzpicture}[font=\footnotesize]
    % 上半部分：真实场景，雷达从 A 移动到 B
    \draw[thick, black!70] (0,0) -- (8,0) node[right]{实际墙面};
    \fill[structurecolor] (2.0,2.2) circle (2pt) node[above right] {$A$（$t_0$ 时刻）};
    \fill[structurecolor!55] (3.2,2.2) circle (2pt) node[above right] {$B$（$t_1$ 时刻）};
    \draw[arrow] (2.0,2.2) -- (3.2,2.2);
    \foreach \x in {0.6,1.5,2.4,3.3,4.2,5.1} {
      \draw[black!35, thin] (2.0,2.2) -- (\x,0);
    }
    \node[label] at (0.4,-0.6) {(a) 雷达在帧内从 $A$ 平移到 $B$，同一堵墙被不同时刻“看到”};
    % 下半部分：点云
    \begin{scope}[yshift=-4.6cm]
      \draw[black!25, dashed] (0.4,0.5) -- (7.6,0.5);
      \node[label] at (7.9,0.5) {理想};
      \foreach \x in {1,2,3,4,5,6,7} {
        \fill[red!60!black] (\x,{0.5+0.10*(\x-4)*(\x-4)}) circle (1.2pt);
      }
      \foreach \x in {1,2,3,4,5,6,7} {
        \fill[structurecolor] (\x,{0.5-0.22}) circle (1.2pt);
      }
      \node[label] at (0.4,-0.1) {(b) 红点：整帧被当作同一时刻，沿墙“拉弯”；蓝点：去畸变后回到一条直线};
    \end{scope}
  \end{tikzpicture}
  \caption{帧内运动造成的点云畸变。雷达在采集一帧的过程中从 $A$ 走到 $B$，若不做补偿，同一堵平面墙会在点云中被弯成一条弧线；按每个点的时间戳去畸变后，点重新落到同一平面上。}\label{fig:16-deskew}
\end{figure}

\section{演化时间线：从 LOAM 到 Point-LIO}\label{sec:16-history}

十年间的路线可以概括成一句话：\textbf{特征越来越少，点用得越来越多，更新越来越频繁}。LOAM 靠人工设计的线/面特征做两阶段匹配；LeGO-LOAM 用地面分割把特征做得更可靠、更省算力；LIO-SAM 把估计框架换成因子图并首次把 IMU 与激光紧耦合地写进因子图；FAST-LIO 把后端换成迭代 EKF，用紧耦合把点面残差直接作用于状态；FAST-LIO2 干脆去掉特征提取，直接对原始点做配准，并配增量式 iKD-Tree 让"每帧都对新点做最近邻"变得便宜；Point-LIO 再进一步，把一个点当成一次更新，彻底摆脱"帧"的束缚。

\begin{figure}[htbp]
  \centering
  \begin{tikzpicture}[font=\footnotesize]
    \draw[-{Stealth[length=2.5mm]}, thick, black!70] (0,0) -- (13.8,0);
    \foreach \x/\yr/\nm in {
      0.7/2014/{LOAM}, 2.9/2018/{LeGO-LOAM}, 5.1/2020/{LIO-SAM},
      7.1/2021/{FAST-LIO}, 9.1/2022/{FAST-LIO2},
      10.9/2023/{Point-LIO}, 12.7/2024/{LIVO2}}{
      \draw[thick, structurecolor] (\x,-0.10) -- (\x,0.10);
      \fill[structurecolor] (\x,0) circle (1.6pt);
      \node[above=3pt, font=\scriptsize, align=center] at (\x,0) {\nm};
      \node[below=3pt, font=\scriptsize, text=black!60] at (\x,0) {\yr};
    }
    \node[font=\scriptsize, text=black!60, align=left] at (6.8,1.15) {特征法 $\longrightarrow$ 直接法};
    \node[font=\scriptsize, text=black!60, align=left] at (6.8,-0.95) {滤波 $\longleftrightarrow$ 优化};
  \end{tikzpicture}
  \caption{3D 激光惯性里程计的演化主线。VoxelMap（2022）与 Point-LIO（2023）几乎同时探索"地图表示"与"更新粒度"两个方向；2022 年之后视觉信息也被引入，形成 FAST-LIVO、R3LIVE 等 LVI 系统。}\label{fig:16-timeline}
\end{figure}

\begin{table}[htbp]
  \centering
  \caption{代表性 LIO 系统一览}\label{tab:16-milestones}
  \footnotesize
  \begin{tabular}{lllll}
    \toprule
    方法 & 年份 & 核心贡献 & 后端 & 典型场景 \\
    \midrule
    LOAM~\cite{zhang2014loam} & 2014 & 线/面特征 + 帧间帧到图两阶段 & 优化 & 无 IMU 的室外测绘 \\
    LeGO-LOAM~\cite{shan2018lego} & 2018 & 地面分割 + 聚类 + 两步优化 & 优化 & 嵌入式、起伏地面 \\
    LIO-SAM~\cite{shan2020liosam} & 2020 & IMU 紧耦合 + 因子图 + 回环/GPS & 因子图 & 有回环需求的室外 \\
    FAST-LIO~\cite{xu2021fastlio} & 2021 & 紧耦合 IEKF，点面残差 & IEKF & 高动态、室内外通用 \\
    FAST-LIO2~\cite{xu2022fastlio2} & 2022 & 去特征 + 原始点配准 + iKD-Tree & IEKF & 通用主力方案 \\
    VoxelMap~\cite{yuan2022voxelmap} & 2022 & 自适应体素 + 概率平面地图 & IEKF & 内存敏感、结构化场景 \\
    Point-LIO~\cite{he2023pointlio} & 2023 & 逐点更新、抗 IMU 饱和 & 逐点 KF & 超高速旋转、IMU 量程小 \\
    FAST-LIVO2~\cite{yuan2024fastlivo2} & 2024 & 激光-惯性-视觉直接法紧耦合 & IEKF & 退化场景 + 纹理丰富 \\
    \bottomrule
  \end{tabular}
\end{table}

\subsection{LOAM（2014）：两阶段、特征法、无 IMU}

LOAM~\cite{zhang2014loam} 是这一支的鼻祖。它先按曲率把每帧点云分成"线特征"（曲率大，如墙角、杆子）和"面特征"（曲率小，如墙面、地面），然后用点到线、点到面两种残差做配准。它的关键设计是\textbf{两阶段分工}：高频（约 $10$\,Hz）的里程计只拿当前帧和上一帧匹配，保证实时；低频（约 $1$\,Hz）的建图模块把当前帧和已经建好的地图匹配，消除累积漂移。两阶段之间用位姿变换把新帧"投"进地图坐标系，避免了地图不断被重写。

LOAM 的局限也很清楚：没有 IMU，帧内运动畸变只能用匀速模型补偿，机器人急转时点云会被拉花；特征提取依赖人工阈值，场景纹理弱或点数稀疏时特征不稳定；两阶段之间的耦合只是"传位姿"，本质仍是松耦合。

\subsection{LeGO-LOAM（2018）：先分割再匹配}

LeGO-LOAM~\cite{shan2018lego} 的出发点是"嵌入式平台算力不够"。它做了三件事：（1）\textbf{地面分割}：用相邻点的俯仰角把地面点分离出来；（2）\textbf{点云聚类}：对非地面点做聚类，去掉小于 30 点的杂簇，抑制噪声；（3）\textbf{两步优化}：先用地面点估计 $(z,\mathrm{roll},\mathrm{pitch})$，再用非地面点估计 $(x,y,\mathrm{yaw})$，把 6 自由度问题拆成两个 3 自由度问题，计算量大幅下降。代价是：强依赖地面假设，机器人在空中或者地面点不足时失效；两步拆分也丢掉了部分耦合信息。

\subsection{LIO-SAM（2020）：紧耦合 + 因子图}

LIO-SAM~\cite{shan2020liosam} 是"滤波 vs 优化"这条分水岭上的代表作。它把 IMU 与激光首次真正紧耦合地放进一个因子图，用 GTSAM 求解。图中共有四类因子（见图~\ref{fig:16-factorgraph}）：IMU 预积分因子、激光里程计因子（点到面配准得到的相对位姿）、GPS 因子、回环因子。整个问题写成最大后验估计：
\begin{equation}\label{eq:16-map}
  \mathcal{X}^{\star}=\argmin_{\mathcal{X}}\;\frac12\sum_{k}\norm{\bm{r}_k(\mathcal{X})}^{2}_{\bm{\Sigma}_k},
\end{equation}
其中 $\mathcal{X}=\{T_1,\dots,T_N\}$ 是所有关键帧位姿，$\bm{r}_k$ 是第 $k$ 个因子的残差（例如预积分因子预测的相对位姿与状态中相对位姿之差），$\norm{\bm{r}}^{2}_{\bm{\Sigma}}=\bm{r}^{\top}\bm{\Sigma}^{-1}\bm{r}$ 是马氏距离的平方（定义见第~\ref{ch:02-math}~章）。式~\eqref{eq:16-map} 用高斯-牛顿迭代求解，每次迭代解一个稀疏线性方程组（细节见第~\ref{ch:13-backend}~章）。

因子图相对滤波器的根本优势是：\textbf{线性化点可以反复修正}。滤波器一旦把某一帧线性化掉、边缘化掉，就再也改不了；而因子图保留历史状态，回环时可以把整条轨迹一起拉直。代价是计算量随关键帧数量增长，必须靠滑窗边缘化控制规模。

\begin{figure}[htbp]
  \centering
  \begin{tikzpicture}[font=\footnotesize]
    % 状态节点
    \foreach \i in {0,...,4} {
      \node[circle, draw=structurecolor, thick, fill=white, inner sep=1.4pt] (x\i) at (2*\i,0) {$x_{\i}$};
    }
    % IMU 预积分因子（上方方块）
    \foreach \i [evaluate=\i as \j using int(\i+1)] in {0,...,3} {
      \node[draw=structurecolor!70, fill=structurecolor!10, minimum size=3.2mm, inner sep=0pt] (f\i) at ({2*\i+1},1.2) {};
      \draw[black!60] (x\i) -- (f\i);
      \draw[black!60] (x\j) -- (f\i);
    }
    \node[label] at (4,1.85) {IMU 预积分因子 $f^{\mathrm{imu}}$};
    % 激光里程计因子（下方方块）
    \foreach \i in {1,2,3} {
      \node[draw=structurecolor!70, fill=structurecolor!10, minimum size=3.2mm, inner sep=0pt] (l\i) at ({2*\i},-1.2) {};
      \draw[black!60] (x\i) -- (l\i);
    }
    \node[label] at (4,-1.85) {激光里程计因子 $f^{\mathrm{lidar}}$};
    % GPS 因子
    \node[draw=orange!80!black, fill=orange!80!black!12, minimum size=3.2mm, inner sep=0pt] (g) at (4,2.5) {};
    \draw[black!60] (x2) -- (g);
    \node[label] at (5.4,2.5) {GPS 因子（可选）};
    % 回环因子
    \node[draw=red!60!black, fill=red!60!black!10, minimum size=3.2mm, inner sep=0pt] (c) at (4,-2.6) {};
    \draw[black!60, dashed] (x0) to[bend right=28] (c);
    \draw[black!60, dashed] (x4) to[bend left=28] (c);
    \node[label] at (5.9,-2.6) {回环因子};
  \end{tikzpicture}
  \caption{LIO-SAM 的因子图结构。圆形是待估状态（关键帧位姿），方块是因子（约束）。IMU 预积分因子约束相邻帧之间的相对运动，激光里程计因子给出相邻帧间的相对位姿观测，回环因子把相距很远的帧连起来，GPS 因子提供绝对位置。}\label{fig:16-factorgraph}
\end{figure}

\subsection{FAST-LIO（2021）与 FAST-LIO2（2022）：把后端换成迭代 EKF}

FAST-LIO~\cite{xu2021fastlio} 的贡献是把"紧耦合"做到极致：状态里不仅有姿态、位置、速度、零偏，还包含重力方向；观测不再是"位姿"，而是每个点与地图平面的点到面残差；求解器是\textbf{迭代扩展卡尔曼滤波}（IEKF），在一次更新里反复线性化、反复重新建立数据关联。为了保证实时，它引入了流形上的迭代公式，让每次迭代的更新量保持在切空间里，避免了四元数归一化之类的数值问题。

FAST-LIO2~\cite{xu2022fastlio2} 做了一个看起来很小、实际上影响深远的决定：\textbf{去掉特征提取}。既然点到面残差已经足够，为什么还要花时间算曲率、挑特征？它直接拿降采样后的原始点云去和地图配准，配合\textbf{增量式 KD 树}（iKD-Tree）——地图点可以动态插入、惰性删除、按包围盒批量操作，不必每帧重建整棵树。这一改动同时带来了三个好处：一是省掉了特征提取的时间，二是弱纹理、非结构化场景也能用，三是原始点携带的信息比几百个特征点丰富得多。

\begin{keybox}{FAST-LIO 系列的设计哲学}
"不要为了省计算量而丢信息。"在 CPU 上，一个点到面残差的计算只要几十次浮点运算，而丢掉的几何信息可能让配准在退化场景下失败。FAST-LIO2 用更便宜的数据结构（iKD-Tree）换来了"用全部点"的奢侈。
\end{keybox}

\subsection{Point-LIO（2023）与 VoxelMap（2022）：两个方向的突破}

Point-LIO~\cite{he2023pointlio} 质疑了"帧"这个概念本身：为什么一定要等一帧点云攒齐再更新？它让\textbf{每个点一到就做一次更新}，点的时间戳直接用于状态传播，因此状态更新的频率等于点的频率（可达 $10^5$\,Hz 量级），并且天然免疫 IMU 饱和——当 IMU 读数不可信时，状态靠逐点的激光观测一步一步拉回来。VoxelMap~\cite{yuan2022voxelmap} 走的是地图表示路线：不用点云地图，而用\textbf{自适应体素}——每个体素里的点如果近似共面，就只存一个平面参数和它的协方差；不够平就继续八叉树细分裂开。这样地图又小又快，还能给出平面的不确定性。两者都会在第~\ref{ch:17-advanced-lio}~章展开。

\subsection{视觉-激光-惯性：FAST-LIVO、R3LIVE、LVI-SAM}

激光在几何退化（长走廊、纯平面墙）时无解，而相机恰好能提供纹理约束；相机在弱光、无纹理、快速运动时失效，激光恰好不受影响。于是 2021 年之后出现了一批 LVI 系统：LVI-SAM~\cite{shan2021lvisam} 用两个因子图（视觉-惯性和激光-惯性）互相提供初值；R3LIVE~\cite{lin2022r3live} 同时建几何地图和视觉纹理地图；FAST-LIVO~\cite{yuan2022fastlivo} 与 FAST-LIVO2~\cite{yuan2024fastlivo2} 把点到面残差与光度误差放进同一个 IEKF 中做直接法紧耦合。它们的共性是：状态估计的骨架仍然是 LIO，视觉只是"多一类残差"。本书只在第~\ref{ch:17-advanced-lio}~章用一小节说明其定位，不展开推导。

\section{两条技术路线：特征法与直接法、滤波与优化}\label{sec:16-routes}

把上面这些方法按两个维度切开，就得到表~\ref{tab:16-routes}。这两个维度经常被混为一谈，其实它们回答的是完全不同的两个问题：\textbf{用什么残差}（特征法 vs 直接法）和\textbf{用什么估计框架}（滤波 vs 优化）。

\begin{table}[htbp]
  \centering
  \caption{LIO 方法的二维分类}\label{tab:16-routes}
  \small
  \begin{tabular}{lll}
    \toprule
    & 滤波（递归） & 优化（因子图/批量） \\
    \midrule
    特征法 & FAST-LIO（线/面点 + IEKF） & LOAM、LeGO-LOAM、LIO-SAM \\
    直接法 & FAST-LIO2、Point-LIO、VoxelMap & 较少见（需把配准残差写成因子） \\
    \bottomrule
  \end{tabular}
\end{table}

\textbf{特征法 vs 直接法}。特征法先挑出少量"好"的点（线点、面点），只对这些点建残差；直接法对全部（降采样后的）点建残差。前者的好处是残差数量少、每次迭代便宜，坏处是特征提取的阈值对场景敏感，且在结构化程度低的场景里挑不出特征；后者好处是信息完整、对场景没有假设，坏处是点数多，必须靠高效的数据结构（iKD-Tree、体素哈希）和降采样来控制计算量。值得注意的是，"直接法"在 LIO 语境里指的是直接使用原始点，而不是视觉里的"直接法"（不提取特征、直接最小化光度误差），两者思想相通但形式不同。

\textbf{滤波 vs 优化}。滤波（EKF/IEKF）只维护当前时刻的状态和协方差，来一帧就"吃掉"一帧，计算量与历史长度无关；优化（因子图）保留全部关键帧，一次性求解所有状态，计算量随规模增长，但可以反复线性化、天然支持回环和全局一致性。第~\ref{ch:09-fusion}~章和第~\ref{ch:13-backend}~章分别讲过两种框架的概率基础，这里只强调一点：\textbf{在里程计这个局部问题上，二者的精度差距远小于在 SLAM 全局问题上的差距}。里程计只需要局部一致，滤波的"不能回头改"并不致命；而一旦要做回环，滤波器就必须额外加一个位姿图模块，否则误差永远补不回来。

\section{紧耦合到底“紧”在哪}\label{sec:16-tight}

\subsection{两种耦合的数学形式}

\paragraph{松耦合（级联）} 的做法是两级估计。第一步，独立的扫描匹配给出位姿观测 $\bm{z}=T_{k}^{L}$ 及其协方差 $\bm{\Sigma}_L$；第二步，把 IMU 递推结果当作先验，把 $\bm{z}$ 当作观测做一次 EKF：
\begin{equation}\label{eq:16-loose}
  \begin{aligned}
    \bar{\bm{x}}_k&=\bm{f}(\bm{x}_{k-1},\bm{u}_k), &\bar{\bm{P}}_k&=\bm{F}\bm{P}_{k-1}\bm{F}^{\top}+\bm{Q},\\
    \bm{K}&=\bar{\bm{P}}_k\bm{H}^{\top}\bigl(\bm{H}\bar{\bm{P}}_k\bm{H}^{\top}+\bm{\Sigma}_L\bigr)^{-1}, & \bm{H}&=\begin{bmatrix}\bm{I}_6&\bm{0}_{6\times12}\end{bmatrix}.
  \end{aligned}
\end{equation}
式~\eqref{eq:16-loose} 中 $\bm{H}$ 只把状态里的位姿部分"取出来"和观测比较，观测维度是 6（位姿）。这确实是紧耦合公式的简化版，但信息流是单向的：扫描匹配不知道 IMU 的存在，它算出来的 $\bm{\Sigma}_L$ 也没有考虑 IMU 先验；而 IMU 只知道"激光说我在这个位姿"，不知道这个位姿是靠哪几个点、在哪个方向上算出来的。

\paragraph{紧耦合（联合）} 的做法是只保留一个估计器，观测不是 6 维位姿，而是 $m$ 个点的点面残差：
\begin{equation}\label{eq:16-tight}
  \bm{K}=\bar{\bm{P}}_k\bm{H}^{\top}\bigl(\bm{H}\bar{\bm{P}}_k\bm{H}^{\top}+\bm{R}_m\bigr)^{-1},\qquad
  \bm{H}=\frac{\partial\bm{r}}{\partial\delta\bm{x}}\in\R^{m\times18},
\end{equation}
其中 $\bm{r}=[r_1,\dots,r_m]^{\top}$，$r_i=\bm{n}_i^{\top}(\bm{R}\bm{p}_i+\bm{p}-\bm{q}_i)$ 是第 $i$ 个点的点到面残差（$\bm{n}_i$ 是从地图拟合出的平面法向，$\bm{q}_i$ 是平面上一点），$\bm{R}_m$ 是各点观测噪声拼成的对角阵。式~\eqref{eq:16-tight} 和式~\eqref{eq:16-loose} 的代数形式完全一样，差别全在 $\bm{H}$ 和观测的含义上：前者是 6 行的"位姿投影"，后者是 $m$ 行的"几何残差"，$m$ 通常是几百到几千。

\begin{figure}[htbp]
  \centering
  \begin{tikzpicture}[font=\small]
    \node[label, text=structurecolor] at (0,1.15) {\textbf{松耦合：级联估计}};
    \node[block] (l0) at (0,0) {3D 激光};
    \node[block] (l1) at (2.4,0) {扫描匹配};
    \node[block] (i0) at (0,-1.7) {IMU};
    \node[block] (i1) at (2.4,-1.7) {积分/预积分};
    \node[block2] (f0) at (5.0,-0.85) {EKF 融合};
    \node[block] (o0) at (7.6,-0.85) {位姿};
    \draw[arrow] (l0) -- (l1);
    \draw[arrow] (i0) -- (i1);
    \draw[arrow] (l1) -- node[label, above right] {位姿观测 $\bm{z}$} (f0);
    \draw[arrow] (i1) -- node[label, below right] {先验 $\bar{\bm{x}},\bar{\bm{P}}$} (f0);
    \draw[arrow] (f0) -- (o0);
    % 紧耦合
    \begin{scope}[yshift=-5.6cm]
      \node[label, text=structurecolor] at (0,1.15) {\textbf{紧耦合：联合估计}};
      \node[block] (a0) at (0,0) {3D 激光};
      \node[block] (a1) at (2.4,0) {点面残差 $\bm{r}$};
      \node[block] (b0) at (0,-1.7) {IMU};
      \node[block] (b1) at (2.4,-1.7) {状态传播};
      \node[block2] (c0) at (5.0,-0.85) {IEKF 联合更新};
      \node[block] (d0) at (7.6,-0.85) {位姿};
      \node[block] (m0) at (2.4,-3.0) {iKD-Tree 地图};
      \draw[arrow] (a0) -- (a1);
      \draw[arrow] (b0) -- (b1);
      \draw[arrow] (a1) -- node[label, above right] {$m$ 维残差} (c0);
      \draw[arrow] (b1) -- node[label, below right] {同一个状态 $\bm{x}$} (c0);
      \draw[arrow] (c0) -- (d0);
      \draw[arrow, dashed] (m0) -- node[label, right] {最近邻/平面} (a1);
      \draw[arrow, dashed] (c0.south) to[bend left=12] (m0.east);
    \end{scope}
  \end{tikzpicture}
  \caption{松耦合与紧耦合的信息流对比。松耦合中扫描匹配是独立模块，只把 6 自由度位姿交给滤波器；紧耦合中点云直接产生残差，与 IMU 传播结果在同一个状态向量、同一个协方差里融合，地图也参与其中。}\label{fig:16-coupling}
\end{figure}

\subsection{紧耦合为什么更鲁棒}

把两者的差别想成一个很具体的场景：机器人正沿一条长走廊前进，走廊方向几何退化，扫描匹配在轴线方向上算不出位移。松耦合下，扫描匹配只好输出一个"轴线方向几乎没动"的位姿，并且给出一个看起来不大、但实际上错误的 $\bm{\Sigma}_L$；滤波器看到这个观测后会相信它，把 IMU 预积分算出来的正确位移拽回去。紧耦合下，$\bm{H}^{\top}\bm{H}$ 在轴线方向上同样接近奇异，但这个奇异方向上的不确定度由先验协方差 $\bar{\bm{P}}_k$ 提供（IMU 已经算出了位移），贝叶斯融合会自动在"激光没约束"的方向上保留 IMU 的推断，只在激光有约束的方向上修正。换句话说，\textbf{松耦合把退化藏进了 $\bm{\Sigma}_L$，紧耦合让退化在 $\bm{H}$ 里显式暴露出来}，滤波器就能按方向分别处理。

\begin{remark}[一个容易误解的点]
在 FAST-LIO2 中，点到面残差对零偏的雅可比恰好为零（点的坐标变换不依赖零偏），那零偏是怎么被激光修正的？答案是：通过传播形成的先验互协方差 $\bm{P}_{\theta b_g}$。姿态误差与陀螺零偏在预测步就被耦合在一起，激光残差修正姿态时，$\bm{K}$ 中对应的行会把零偏一起拉动。这仍然是紧耦合——两者共享同一个状态和同一个协方差矩阵；只是耦合路径是"先验相关性"而不是"残差显式依赖"。
\end{remark}

\section{选型建议}\label{sec:16-select}

\begin{table}[htbp]
  \centering
  \caption{LIO 方案选型决策表}\label{tab:16-select}
  \footnotesize
  \begin{tabular}{lll}
    \toprule
    需求/场景 & 推荐方案 & 理由 \\
    \midrule
    纯 2D 平面、算力极低 & Cartographer（第~\ref{ch:15-cartographer}~章） & 3 自由度，代价最小 \\
    地面机器人、需要回环与 GPS & LIO-SAM~\cite{shan2020liosam} & 因子图可回环、可融合绝对观测 \\
    通用 3D、稳定优先 & FAST-LIO2~\cite{xu2022fastlio2} & 无特征、直接法、社区成熟 \\
    剧烈旋转、IMU 量程不足 & Point-LIO~\cite{he2023pointlio} & 逐点更新，帧内不依赖 IMU 可信 \\
    长走廊/隧道等退化场景 & 视觉 + 激光紧耦合（LVI） & 纹理补几何 \\
    内存/算力受限、结构化场景 & VoxelMap~\cite{yuan2022voxelmap} & 平面地图，存储与查询都省 \\
    无 IMU 可用 & LOAM 类纯激光 & 不依赖惯性观测 \\
    \bottomrule
  \end{tabular}
\end{table}

选型时还要考虑三个工程约束：（1）\textbf{时间同步}：激光与 IMU 的时间戳必须同步到毫秒级，否则去畸变和传播都会错，这是 LIO 现场翻车的第一大原因；（2）\textbf{外参标定}：雷达到 IMU 的旋转和平移要足够准，$1^{\circ}$ 的外参误差在 $10$\,m 处就是 $17$\,cm 的偏差；（3）\textbf{IMU 量程}：RoboMaster 的撞击瞬间角速度很容易超过 $2000^{\circ}$/s，选型时要看陀螺量程和算法对饱和的鲁棒性。

\begin{summary}

\begin{enumerate}
  \item 3D 场景需要 6 自由度位姿，纯激光里程计在几何退化、帧内畸变、快速运动三种情况下失效；IMU 频率高、短时准，恰好补上帧内运动，激光则把 IMU 的漂移拉回来。
  \item LIO 的定义是"激光与 IMU 在同一个估计器中联合估计位姿及辅助状态"。紧耦合 = 原始测量进同一次更新；松耦合 = 一个传感器的估计结果当另一个的观测。
  \item 运动畸变来自"一帧点云是一个 $100$\,ms 的时间窗"，去畸变的本质是把每个点按自己的时间戳变换到参考时刻：$\bm{p}_i^{L_{\mathrm{ref}}}=R(t_{\mathrm{ref}})^{\top}(R(t_i)\bm{p}_i^{L}+\bm{p}(t_i)-\bm{p}(t_{\mathrm{ref}}))$。
  \item 去畸变只需要帧内相对运动；匀速模型误差约 $\frac12 a\Delta t^2$，IMU 辅助模型用积分轨迹插值，连续时间表示（B 样条）精度更高但更贵。
  \item 演化主线：LOAM（特征法两阶段）$\to$ LeGO-LOAM（地面分割、轻量化）$\to$ LIO-SAM（紧耦合因子图、可回环）$\to$ FAST-LIO（IEKF + 点面残差）$\to$ FAST-LIO2（去特征 + iKD-Tree）$\to$ Point-LIO（逐点更新、抗饱和）。
  \item 两个正交的维度：特征法/直接法回答"用什么残差"，滤波/优化回答"用什么估计框架"，不要混为一谈。
  \item 紧耦合更鲁棒的根本原因：退化方向上的不确定度由先验协方差承担，滤波器可以"按方向"融合，而不是把一个退化的 6 维位姿当成可靠观测。
\end{enumerate}


\end{summary}

\begin{problemset}

  \item 设激光雷达转速 $10$\,Hz，机器人以 $1$\,m/s 平移、同时以 $60^{\circ}$/s 绕竖直轴旋转，点距离 $4$\,m。估算帧首帧尾的点在空间上错开多少（分别给出平移项与旋转项）。
  \item 用式~\eqref{eq:16-constant-velocity} 推导匀速模型下帧中点（$s=1/2$）的位姿，并说明若真实轨迹在帧内有恒定角加速度 $\alpha$，插值误差的一阶项是多少。
  \item 证明式~\eqref{eq:16-loose} 与式~\eqref{eq:16-tight} 在 $\bm{H}$ 取 $\begin{bmatrix}\bm{I}_6&\bm{0}\end{bmatrix}$、且扫描匹配的 $\bm{\Sigma}_L$ 恰好等于"把全部点残差边缘化掉"的结果时等价。这个等价条件在实践中为什么几乎不成立？
  \item 长走廊场景下，$\bm{H}^{\top}\bm{H}$ 为什么会奇异？请写出沿走廊方向的零空间向量，并解释为什么先验协方差可以补上这个方向。
  \item LIO-SAM 用因子图、FAST-LIO2 用 IEKF，两者都能跑出厘米级轨迹。请各举一个"其中一个明显更优"的场景，并说明原因。
\end{problemset}

\begin{reading}

\begin{itemize}
  \item FAST-LIO2 论文~\cite{xu2022fastlio2}：读它之前先读第~\ref{ch:17-advanced-lio}~章，重点看状态定义、IEKF 更新式与 iKD-Tree 的算法描述。
  \item LIO-SAM 论文~\cite{shan2020liosam} 与 GTSAM 文档：理解因子图如何把 IMU 预积分、激光位姿、回环写进同一个优化问题。
  \item Point-LIO 论文~\cite{he2023pointlio}：重点看它如何用逐点更新绕开"帧内 IMU 不可信"的问题，以及饱和检测的处理方式。
  \item VoxelMap 论文~\cite{yuan2022voxelmap}：理解自适应体素与概率平面地图，对比 iKD-Tree 的存储与查询代价。
\end{itemize}


\end{reading}

\chapter{FAST-LIO2、Point-LIO 与先进 LIO 深度解析}\label{ch:17-advanced-lio}

\usetikzlibrary{arrows.meta, positioning, calc, shapes.geometric}

\chapterintro{这是全书数学密度最高的一章。前 16 章里我们反复听到两个名字：FAST-LIO2 和 Point-LIO。它们为什么能在普通 CPU 上跑到 $10$\,kHz 的状态输出、能在剧烈撞击下不发散？本章把它们的推导从头拆到底：先在 $\SO(3)\times\R^{15}$ 上定义状态和误差状态，再把 IMU 连续运动方程积分成离散传播，然后一块一块推出 $18\times18$ 的误差状态转移矩阵 $F$，接着用最大后验估计重新导出迭代 EKF（IEKF）的更新式与收敛判据，逐项推导点到面残差对状态的雅可比，最后讲清 iKD-Tree 的增量插入、box-wise 删除与 Point-LIO 的逐点更新。读完本章，你应该能对着论文公式直接找到源码里对应的函数，也能自己改一版滤波器。}

\section{FAST-LIO2 的总体架构}\label{sec:17-arch}

FAST-LIO2~\cite{xu2022fastlio2} 的程序结构可以概括成"一条主线、两个频率"：

\begin{itemize}
  \item \textbf{主线}：状态 $\bm{x}$ 与协方差 $\bm{P}$ 是唯一被维护的东西。IMU 来了就\textbf{预测}（把 $\bm{x},\bm{P}$ 沿运动方程往前推），激光来了就\textbf{更新}（用点面残差把 $\bm{x},\bm{P}$ 修正）。
  \item \textbf{两个频率}：IMU 一般 $100\sim200$\,Hz（$\Delta t\approx5\sim10$\,ms），激光 $10$\,Hz（一帧 $100$\,ms）。预测在 IMU 频率上进行，状态输出（odom）也在 IMU 频率上；更新只在每帧点云到达时进行。
  \item \textbf{地图}：所有去畸变后的点按世界坐标插入\textbf{增量式 KD 树}（iKD-Tree），既作为下一帧最近邻查询的地图，也作为建图的输出。
\end{itemize}

\begin{figure}[htbp]
  \centering
  \begin{tikzpicture}[font=\small]
    \node[block] (imu) at (0,0) {IMU\\$200$\,Hz};
    \node[block] (lid) at (0,-2.2) {3D 激光\\$10$\,Hz};
    \node[block] (pre) at (2.8,0) {IMU 预积分\\（传播）};
    \node[block] (des) at (2.8,-2.2) {去畸变\\+ 降采样};
    \node[block] (map) at (2.8,-4.2) {iKD-Tree\\增量地图};
    \node[block] (res) at (5.6,-2.2) {点面残差\\$\bm{r},\bm{H}$};
    \node[block2] (ekf) at (8.2,-1.0) {迭代 EKF\\$\bm{x}\leftarrow\bm{x}\boxplus\delta\bm{x}$};
    \node[block] (out) at (11.0,-1.0) {odom\\去畸变点云};
    \draw[arrow] (imu) -- (pre);
    \draw[arrow] (lid) -- (des);
    \draw[arrow] (pre) -- node[label, above] {先验 $\bar{\bm{x}},\bar{\bm{P}}$} (ekf);
    \draw[arrow] (des) -- (res);
    \draw[arrow] (res) -- node[label, above] {$\bm{H},\bm{R}_m$} (ekf);
    \draw[arrow] (ekf) -- (out);
    \draw[arrow, dashed] (map) -- node[label, right] {最近邻/平面} (res);
    \draw[arrow, dashed] (ekf.south east) to[bend left=18] node[label, right] {插入新点} (map.east);
    \draw[arrow, dashed] (ekf.north east) to[bend right=18] node[label, right] {位姿预测} (des.east);
  \end{tikzpicture}
  \caption{FAST-LIO2 的数据流。IMU 以 $200$\,Hz 驱动状态传播并持续输出里程计；激光帧到达后，先用传播位姿把点云去畸变、降采样，再为每个点在 iKD-Tree 中找平面、算点面残差，最后用迭代 EKF 一次性融合全部残差，并把新点插回地图。}\label{fig:17-arch}
\end{figure}

一句话总结 FAST-LIO2 的三个关键设计：状态里放重力并跟随估计；误差状态定义在流形切空间上，更新量小、线性化准；\textbf{一次更新用上全部点}，不做特征提取，靠 iKD-Tree 把"对每个点找最近邻"的成本压到可接受。

\section{状态定义与流形更新}\label{sec:17-state}

\subsection{主状态}

FAST-LIO2 的主状态取为
\begin{equation}\label{eq:17-state}
  \bm{x}=\bigl[\bm{R},\;\bm{p},\;\bm{v},\;\bm{b}_g,\;\bm{b}_a,\;\bm{g}\bigr],
\end{equation}
其中 $\bm{R}\in\SO(3)$ 是 IMU 系到世界系的旋转（$T_{WB}$ 的旋转部分），$\bm{p}\in\R^3$ 是 IMU 系原点在世界系中的位置，$\bm{v}\in\R^3$ 是世界系速度，$\bm{b}_g,\bm{b}_a\in\R^3$ 是陀螺与加速度计零偏（定义见第~\ref{ch:03-sensors}~章），$\bm{g}\in\R^3$ 是重力向量，默认 $\bm{g}=[0,0,-9.81]^{\top}$。注意 $\bm{R}$ 不是 $\R^3$ 里的普通向量，而是李群 $\SO(3)$ 上的元素，所以状态空间是流形
\begin{equation}
  \mathcal{M}=\SO(3)\times\R^{15},
\end{equation}
$\R^{15}$ 来自 $\bm{p},\bm{v},\bm{b}_g,\bm{b}_a$（各 3 维）与 $\bm{g}$（按 3 维写）。严格说重力是单位向量，只有 2 个自由度，实际实现把它参数化在球面 $S^2$ 上，于是误差状态是 17 维；本节为了式子整齐按 3 维 $\delta\bm{g}$ 推导，并对两者的关系给出说明。

\subsection{为什么把重力放进状态}

初学者常问：重力不是常数吗，为什么要估计？原因有三。其一，\textbf{初始姿态未知}：开机时并不知道机器人的 roll/pitch，而静止的加速度计测到的正是 $-R^{\top}\bm{g}$，所以"重力方向"与"初始姿态"是同一件事的两面，把 $\bm{g}$ 放进状态等于让滤波器自己去对齐重力。其二，\textbf{标定残差}：加速度计的零偏与标度因数误差在长时间上表现为一个"等效的重力方向偏差"，在线估 $\bm{g}$ 能吸收一部分。其三，\textbf{一致性}：若把 $\bm{g}$ 写死而初始姿态差 $2^{\circ}$，这个误差会被永久固化在重力方向上，速度持续被一个错误的常值加速度污染，位置随时间二次发散；让 $\bm{g}$ 可估，滤波器就能把"一开始估错了"改过来。

\subsection{误差状态与两个流形运算}

我们不直接在 $\bm{x}$ 上做加法（$\bm{R}+\delta\bm{R}$ 会破坏正交性），而是定义\textbf{误差状态}
\begin{equation}\label{eq:17-error-state}
  \delta\bm{x}=\bigl[\delta\bm{\theta},\;\delta\bm{p},\;\delta\bm{v},\;\delta\bm{b}_g,\;\delta\bm{b}_a,\;\delta\bm{g}\bigr]\in\R^{18},
\end{equation}
并用两个流形运算把"状态"和"误差"联系起来：
\begin{equation}\label{eq:17-plus-minus}
  \bm{x}\boxplus\delta\bm{x}\;\triangleq\;
  \begin{cases}
    \bm{R}\leftarrow\bm{R}\exp\bigl([\delta\bm{\theta}]_\times\bigr),\\
    \bm{p}\leftarrow\bm{p}+\delta\bm{p},\quad \bm{v}\leftarrow\bm{v}+\delta\bm{v},\\
    \bm{b}_g\leftarrow\bm{b}_g+\delta\bm{b}_g,\quad \bm{b}_a\leftarrow\bm{b}_a+\delta\bm{b}_a,\quad \bm{g}\leftarrow\bm{g}+\delta\bm{g},
  \end{cases}
\end{equation}
\begin{equation}
  \bm{x}\boxminus\bar{\bm{x}}\;\triangleq\;
  \bigl[\mathrm{Log}(\bar{\bm{R}}^{\top}\bm{R}),\;\bm{p}-\bar{\bm{p}},\;\bm{v}-\bar{\bm{v}},\;\bm{b}_g-\bar{\bm{b}}_g,\;\bm{b}_a-\bar{\bm{b}}_a,\;\bm{g}-\bar{\bm{g}}\bigr].
\end{equation}
式~\eqref{eq:17-plus-minus} 中的旋转部分叫\textbf{右扰动}：新姿态等于旧姿态右乘一个由 $\delta\bm{\theta}$ 生成的小旋转。$\boxminus$ 是它的逆运算，$\bar{\bm{R}}^{\top}\bm{R}$ 表示"从参考姿态到当前姿态还需要转多少"，再取对数变回旋转向量。这两个运算合起来构成切空间上的"加法"，是 ESKF 与流形滤波的标准做法~\cite{sola2017eskf}。

为什么必须用误差状态？因为 $\bm{R}$ 只有 3 个自由度却用 9 个数表示，直接对 9 个数加噪声会破坏 $\bm{R}^{\top}\bm{R}=\bm{I}$；而 $\delta\bm{\theta}$ 是 3 维的小量，可以放心地做线性化、做高斯假设。另一个好处是：误差状态在机器人正常运动时始终接近零，一阶线性化的误差极小，这正是一阶 EKF 能这么准的原因。

这里补充两点代码层面的约定。第一，官方实现把各状态块命名为 \texttt{pos}、\texttt{rot}、\texttt{vel}、\texttt{bg}、\texttt{ba}、\texttt{grav}，并默认把 LiDAR 到 IMU 的外参 $\bm{R}_L,\bm{t}_L$（各 3 维）也放进状态，所以误差状态是 $23$ 维而不是 $17$ 维，多出来的 6 维是外参；外参块的推导与位姿块完全同构，本书为突出主线只推导 $18$ 维。第二，重力按 $S^2$ 的 2 维切空间处理，最后一块是 $3\times2$ 的窄矩阵，读者在代码里看到 \texttt{dim} 为 23 时不要惊讶。

\section{IMU 运动学：从连续方程到离散传播}\label{sec:17-imu}

\subsection{测量模型}

IMU 输出的是角速度 $\tilde{\bm{\omega}}$ 与比力 $\tilde{\bm{a}}$（加速度计测的是"除重力以外的比力"，不是加速度本身）。设真实角速度 $\bm{\omega}$ 与真实加速度 $\bm{a}$，则测量模型为
\begin{equation}\label{eq:17-meas-model}
  \tilde{\bm{\omega}}=\bm{\omega}+\bm{b}_g+\bm{n}_g,\qquad
  \tilde{\bm{a}}=\bm{R}^{\top}\bigl(\bm{a}-\bm{g}\bigr)+\bm{b}_a+\bm{n}_a,
\end{equation}
其中 $\bm{n}_g,\bm{n}_a$ 是零均值高斯白噪声，$\bm{b}_g,\bm{b}_a$ 是缓慢变化的零偏（建模为随机游走）。式~\eqref{eq:17-meas-model} 的第二个式子说明：加速度计在机体系里测到的是"世界系加速度 $\bm{a}$ 减去重力 $\bm{g}$"，再转到机体系并加上零偏噪声。静止时 $\bm{a}=0$，读数就是 $\bm{R}^{\top}(-\bm{g})$，这就是为什么静止时可以用加速度计测 roll/pitch。

\subsection{连续运动方程}

对式~\eqref{eq:17-meas-model} 反解出真实量，代入运动学，得到
\begin{equation}\label{eq:17-continuous}
  \begin{cases}
    \dot{\bm{R}}=\bm{R}\,[\bm{\omega}]_\times,\\
    \dot{\bm{p}}=\bm{v},\\
    \dot{\bm{v}}=\bm{R}\,\bigl(\tilde{\bm{a}}-\bm{b}_a-\bm{n}_a\bigr)+\bm{g},\\
    \dot{\bm{b}}_g=\bm{n}_{bg},\qquad \dot{\bm{b}}_a=\bm{n}_{ba},\qquad \dot{\bm{g}}=\bm{0},
  \end{cases}
\end{equation}
其中 $[\bm{\omega}]_\times$ 是由 $\bm{\omega}=[\omega_x,\omega_y,\omega_z]^{\top}$ 生成的反对称矩阵：
\begin{equation}
  [\bm{\omega}]_\times=\begin{bmatrix}0&-\omega_z&\omega_y\\ \omega_z&0&-\omega_x\\ -\omega_y&\omega_x&0\end{bmatrix},\qquad
  [\bm{\omega}]_\times\bm{u}=\bm{\omega}\times\bm{u}.
\end{equation}

第一条 $\dot{\bm{R}}=\bm{R}[\bm{\omega}]_\times$ 值得单独推一遍。设 $t$ 到 $t+\Delta t$ 内机体系下的角速度恒为 $\bm{\omega}$，这段时间机体系相对世界系转过一个小角度 $\bm{\omega}\Delta t$，于是
\begin{equation}
  \bm{R}(t+\Delta t)=\bm{R}(t)\exp\bigl([\bm{\omega}\Delta t]_\times\bigr)
  \;\approx\;\bm{R}(t)\bigl(\bm{I}+[\bm{\omega}]_\times\Delta t\bigr),
\end{equation}
移项、两边除以 $\Delta t$ 并令 $\Delta t\to0$，即得 $\dot{\bm{R}}=\bm{R}[\bm{\omega}]_\times$。注意 $\bm{\omega}$ 是\textbf{机体系}下的角速度，所以要右乘；若用世界系角速度 $\bm{\omega}^{W}$，则应是 $\dot{\bm{R}}=[\bm{\omega}^{W}]_\times\bm{R}$。这两种写法等价，但混用会导致符号错误，本书统一用机体系角速度、右乘。

式~\eqref{eq:17-continuous} 前三条的物理含义：姿态由角速度积分，位置由速度积分，速度由"比力转到世界系再加上重力"积分。后三条说明零偏和重力在\textbf{两次 IMU 采样之间}被视为常量，它们的随机游走只在长时间尺度上起作用。

\subsection{离散传播}

IMU 采样间隔 $\Delta t$ 通常为 $5\sim10$\,ms，远小于激光帧长 $100$\,ms。设第 $k$ 个 IMU 采样时刻为 $t_k$，把式~\eqref{eq:17-continuous} 在 $[t_k,t_k+\Delta t]$ 上积分。假设区间内 $\bm{\omega}$、$\bm{a}$ 恒定（零阶保持），且 $t_k$ 时刻的状态已知，则对任意 $\tau\in[0,\Delta t]$：
\begin{equation}\label{eq:17-exp-rot}
  \bm{R}(t_k+\tau)=\bm{R}_k\exp\bigl([\bm{\omega}_k\tau]_\times\bigr),\qquad
  \bm{\omega}_k\triangleq\tilde{\bm{\omega}}_k-\bm{b}_{g,k}.
\end{equation}
速度通过对加速度积分得到：
\begin{equation}
  \bm{v}(t_k+\tau)=\bm{v}_k+\int_0^{\tau}\Bigl[\bm{R}_k\exp\bigl([\bm{\omega}_k s]_\times\bigr)\bm{a}_k+\bm{g}\Bigr]\diff s,
  \qquad \bm{a}_k\triangleq\tilde{\bm{a}}_k-\bm{b}_{a,k}.
\end{equation}
把 $\exp([\bm{\omega}_k s]_\times)=\bm{I}+[\bm{\omega}_k]_\times s+O(s^2)$ 代入，只保留到一阶：
\begin{equation}
  \bm{v}(t_k+\tau)=\bm{v}_k+\bm{R}_k\bm{a}_k\,\tau+\bm{g}\,\tau+O(\tau^2).
\end{equation}
位置再由 $\bm{p}(t_k+\tau)=\bm{p}_k+\int_0^\tau\bm{v}(t_k+s)\diff s$ 积分一次，保留到 $\tau^2$：
\begin{equation}
  \bm{p}(t_k+\tau)=\bm{p}_k+\bm{v}_k\tau+\frac12\bigl(\bm{R}_k\bm{a}_k+\bm{g}\bigr)\tau^2+O(\tau^3).
\end{equation}
取 $\tau=\Delta t$，再补上零偏的零阶保持，得到工程上最常用的离散传播公式：
\begin{equation}\label{eq:17-discrete}
  \boxed{
  \begin{aligned}
    \bm{R}_{k+1}&=\bm{R}_k\exp\bigl([\bm{\omega}_k\Delta t]_\times\bigr),\\
    \bm{v}_{k+1}&=\bm{v}_k+\bigl(\bm{R}_k\bm{a}_k+\bm{g}_k\bigr)\Delta t,\\
    \bm{p}_{k+1}&=\bm{p}_k+\bm{v}_k\Delta t+\frac12\bigl(\bm{R}_k\bm{a}_k+\bm{g}_k\bigr)\Delta t^2,\\
    \bm{b}_{g,k+1}&=\bm{b}_{g,k},\qquad \bm{b}_{a,k+1}=\bm{b}_{a,k},\qquad \bm{g}_{k+1}=\bm{g}_k.
  \end{aligned}}
\end{equation}
式~\eqref{eq:17-discrete} 是"用第 $k$ 个采样点的测量把状态推到第 $k+1$ 个采样点"，$100$\,ms 的一帧里要做约 20 次这样的传播，这就是预测步。注意离散化丢掉了 $O(\Delta t^2)$ 的项（旋转对加速度方向的影响、角速度变化的影响等）；由于 $\Delta t$ 只有几毫秒，这些项通常可以忽略。实际实现中，为了进一步减小误差，常见做法是用\textbf{中点积分}（用 $t_k$ 与 $t_{k+1}$ 的测量平均值）或者把 $\bm{R}$ 的旋转放在半步位置上计算，代价是代码复杂一些。

指数映射 $\exp([\bm{\phi}]_\times)$ 用 Rodrigues 公式计算：
\begin{equation}\label{eq:17-rodrigues}
  \exp\bigl([\bm{\phi}]_\times\bigr)=\bm{I}+\frac{\sin\theta}{\theta}[\bm{\phi}]_\times+\frac{1-\cos\theta}{\theta^2}[\bm{\phi}]_\times^2,\qquad \theta=\norm{\bm{\phi}}.
\end{equation}
当 $\theta$ 很小时用一阶近似 $\exp([\bm{\phi}]_\times)\approx\bm{I}+[\bm{\phi}]_\times$。例如以 $1$\,rad/s 转、$\Delta t=5$\,ms 时 $\bm{\phi}=(0,0,0.005)$、$\theta=0.005$，Rodrigues 公式给出的各项系数为 $0.999996$ 与 $0.4999997$，与一阶近似的差别在 $10^{-6}$ 量级——这就是为什么小角度下两者可以互换。

\section{误差状态传播：$F$ 矩阵逐块推导}\label{sec:17-f}

这一节是全章的核心。我们要回答的问题是：如果 $t_k$ 时刻的状态有一个小误差 $\delta\bm{x}_k$，那么经过一次 IMU 传播后，$t_{k+1}$ 时刻的误差 $\delta\bm{x}_{k+1}$ 是多少？答案写成一阶线性形式
\begin{equation}\label{eq:17-f-linear}
  \delta\bm{x}_{k+1}=\bm{F}_k\,\delta\bm{x}_k+\bm{G}_k\,\bm{n}_k,
\end{equation}
其中 $\bm{F}_k\in\R^{18\times18}$ 是误差状态转移矩阵，$\bm{G}_k\in\R^{18\times12}$ 把过程噪声 $\bm{n}=[\bm{n}_g^{\top},\bm{n}_a^{\top},\bm{n}_{bg}^{\top},\bm{n}_{ba}^{\top}]^{\top}$ 映射进误差状态。连续时间下，$\dot{\delta\bm{x}}=\bm{F}\delta\bm{x}+\bm{G}\bm{n}$，离散化取 $\bm{F}_k\approx\bm{I}+\bm{F}\Delta t$。

下面逐块推导。\textbf{约定}：真实状态 = 名义状态 $\boxplus$ 误差，即
\begin{equation}
  \bm{R}=\hat{\bm{R}}\exp([\delta\bm{\theta}]_\times),\quad
  \bm{p}=\hat{\bm{p}}+\delta\bm{p},\quad
  \bm{v}=\hat{\bm{v}}+\delta\bm{v},\quad
  \bm{b}_g=\hat{\bm{b}}_g+\delta\bm{b}_g,\quad
  \bm{b}_a=\hat{\bm{b}}_a+\delta\bm{b}_a,\quad
  \bm{g}=\hat{\bm{g}}+\delta\bm{g},
\end{equation}
带帽的量是滤波器当前的名义估计，且所有推导只保留 $\delta$ 的一阶项，忽略 $O(\delta^2)$。

\subsection{姿态误差 $\delta\bm{\theta}$}

对 $\bm{R}=\hat{\bm{R}}\exp([\delta\bm{\theta}]_\times)$ 求导。先算左边：
\begin{equation}
  \frac{\diff}{\diff t}\Bigl(\hat{\bm{R}}\exp([\delta\bm{\theta}]_\times)\Bigr)
  =\dot{\hat{\bm{R}}}\exp([\delta\bm{\theta}]_\times)+\hat{\bm{R}}\frac{\diff}{\diff t}\exp([\delta\bm{\theta}]_\times).
\end{equation}
其中 $\dot{\hat{\bm{R}}}=\hat{\bm{R}}[\hat{\bm{\omega}}]_\times$，且由 $\exp([\delta\bm{\theta}(t+\diff t)]_\times)\approx\exp([\delta\bm{\theta}]_\times)\exp([\dot{\delta\bm{\theta}}\diff t]_\times)$ 可得 $\frac{\diff}{\diff t}\exp([\delta\bm{\theta}]_\times)=\exp([\delta\bm{\theta}]_\times)[\dot{\delta\bm{\theta}}]_\times$。于是左边为
\begin{equation}
  \hat{\bm{R}}[\hat{\bm{\omega}}]_\times\exp([\delta\bm{\theta}]_\times)+\hat{\bm{R}}\exp([\delta\bm{\theta}]_\times)[\dot{\delta\bm{\theta}}]_\times.
\end{equation}
右边由真实运动方程 $\dot{\bm{R}}=\bm{R}[\bm{\omega}]_\times$ 给出，其中 $\bm{\omega}=\tilde{\bm{\omega}}-\bm{b}_g-\bm{n}_g=\hat{\bm{\omega}}-\delta\bm{b}_g-\bm{n}_g$，所以
\begin{equation}
  \dot{\bm{R}}=\hat{\bm{R}}\exp([\delta\bm{\theta}]_\times)[\hat{\bm{\omega}}-\delta\bm{b}_g-\bm{n}_g]_\times.
\end{equation}
令左右相等，并在两边左乘 $\hat{\bm{R}}^{\top}$：
\begin{equation}
  [\hat{\bm{\omega}}]_\times\exp([\delta\bm{\theta}]_\times)+\exp([\delta\bm{\theta}]_\times)[\dot{\delta\bm{\theta}}]_\times
  =\exp([\delta\bm{\theta}]_\times)[\hat{\bm{\omega}}-\delta\bm{b}_g-\bm{n}_g]_\times.
\end{equation}
用一阶近似 $\exp([\delta\bm{\theta}]_\times)\approx\bm{I}+[\delta\bm{\theta}]_\times$ 展开两边，保留一阶项：
\begin{align}
  \text{左}&=[\hat{\bm{\omega}}]_\times+[\hat{\bm{\omega}}]_\times[\delta\bm{\theta}]_\times+[\dot{\delta\bm{\theta}}]_\times,\\
  \text{右}&=[\hat{\bm{\omega}}]_\times-[\delta\bm{b}_g]_\times-[\bm{n}_g]_\times+[\delta\bm{\theta}]_\times[\hat{\bm{\omega}}]_\times.
\end{align}
移项，利用反对称矩阵的交换子恒等式 $[\bm{a}]_\times[\bm{b}]_\times-[\bm{b}]_\times[\bm{a}]_\times=[\bm{a}\times\bm{b}]_\times$：
\begin{equation}
  [\dot{\delta\bm{\theta}}]_\times=[\delta\bm{\theta}]_\times[\hat{\bm{\omega}}]_\times-[\hat{\bm{\omega}}]_\times[\delta\bm{\theta}]_\times-[\delta\bm{b}_g]_\times-[\bm{n}_g]_\times
  =[\delta\bm{\theta}\times\hat{\bm{\omega}}-\delta\bm{b}_g-\bm{n}_g]_\times.
\end{equation}
去掉两边的 $[\cdot]_\times$，并注意 $\delta\bm{\theta}\times\hat{\bm{\omega}}=-\hat{\bm{\omega}}\times\delta\bm{\theta}=-[\hat{\bm{\omega}}]_\times\delta\bm{\theta}$，得到
\begin{equation}\label{eq:17-dtheta}
  \dot{\delta\bm{\theta}}=-[\hat{\bm{\omega}}]_\times\,\delta\bm{\theta}-\delta\bm{b}_g-\bm{n}_g.
\end{equation}
式~\eqref{eq:17-dtheta} 三项各有物理含义：$-\delta\bm{b}_g$ 是"零偏估错多少，姿态就以多快速度偏多少"；$-\bm{n}_g$ 是陀螺噪声直接积分成角度误差；$-[\hat{\bm{\omega}}]_\times\delta\bm{\theta}$ 是因为误差 $\delta\bm{\theta}$ 表达在机体系里，而机体系自己在转，所以这个误差向量要"跟着坐标系转"，其变化率为 $-\bm{\omega}\times\delta\bm{\theta}$。第三项是初学者最容易漏掉的耦合项。

\subsection{位置误差 $\delta\bm{p}$}

由 $\dot{\bm{p}}=\bm{v}$ 直接得
\begin{equation}\label{eq:17-dp}
  \dot{\delta\bm{p}}=\delta\bm{v}.
\end{equation}
物理含义：位置误差的变化率就是速度误差，这是纯运动学关系，没有任何近似。

\subsection{速度误差 $\delta\bm{v}$}

真实加速度为 $\bm{a}=\bm{R}(\tilde{\bm{a}}-\bm{b}_a-\bm{n}_a)+\bm{g}$，名义加速度为 $\hat{\bm{a}}=\hat{\bm{R}}(\tilde{\bm{a}}-\hat{\bm{b}}_a)+\hat{\bm{g}}$。相减：
\begin{align}
  \delta\dot{\bm{v}}&=\bm{R}\bigl(\tilde{\bm{a}}-\bm{b}_a-\bm{n}_a\bigr)+\bm{g}-\Bigl[\hat{\bm{R}}\bigl(\tilde{\bm{a}}-\hat{\bm{b}}_a\bigr)+\hat{\bm{g}}\Bigr]\nonumber\\
  &\approx\hat{\bm{R}}\bigl(\bm{I}+[\delta\bm{\theta}]_\times\bigr)\bigl(\tilde{\bm{a}}-\hat{\bm{b}}_a-\delta\bm{b}_a-\bm{n}_a\bigr)-\hat{\bm{R}}\bigl(\tilde{\bm{a}}-\hat{\bm{b}}_a\bigr)+\delta\bm{g}.
\end{align}
记 $\bm{a}\triangleq\tilde{\bm{a}}-\hat{\bm{b}}_a$，展开并只保留一阶项：
\begin{align}
  \delta\dot{\bm{v}}&=\hat{\bm{R}}\bm{a}+\hat{\bm{R}}[\delta\bm{\theta}]_\times\bm{a}-\hat{\bm{R}}\delta\bm{b}_a-\hat{\bm{R}}\bm{n}_a-\hat{\bm{R}}\bm{a}+\delta\bm{g}+O(\delta^2)\nonumber\\
  &=\hat{\bm{R}}[\delta\bm{\theta}]_\times\bm{a}-\hat{\bm{R}}\delta\bm{b}_a-\hat{\bm{R}}\bm{n}_a+\delta\bm{g}.
\end{align}
再用 $[\delta\bm{\theta}]_\times\bm{a}=\delta\bm{\theta}\times\bm{a}=-[\bm{a}]_\times\delta\bm{\theta}$，得
\begin{equation}\label{eq:17-dv}
  \dot{\delta\bm{v}}=-\hat{\bm{R}}[\tilde{\bm{a}}-\hat{\bm{b}}_a]_\times\,\delta\bm{\theta}-\hat{\bm{R}}\,\delta\bm{b}_a-\hat{\bm{R}}\,\bm{n}_a+\delta\bm{g}.
\end{equation}
式~\eqref{eq:17-dv} 的四项：姿态误差让比力的方向偏了，于是产生加速度误差 $-\hat{\bm{R}}[\bm{a}]_\times\delta\bm{\theta}$；加速度计零偏误差直接产生 $\hat{\bm{R}}\delta\bm{b}_a$ 的加速度偏差；噪声 $\bm{n}_a$ 同理；重力估计误差 $\delta\bm{g}$ 原封不动地进入加速度误差。做个数量级估计：姿态差 $1^{\circ}=0.017$\,rad，比力 $10\,\mathrm{m/s^2}$，则加速度误差约 $0.17\,\mathrm{m/s^2}$，$1$\,s 就积累出 $0.17$\,m/s 的速度误差、$8.5$\,cm 的位置误差——所以姿态和速度是强耦合的，不能分开估。

\subsection{零偏与重力}

零偏建模为随机游走，其连续时间导数是白噪声：
\begin{equation}\label{eq:17-dbias}
  \dot{\delta\bm{b}}_g=\bm{n}_{bg},\qquad \dot{\delta\bm{b}}_a=\bm{n}_{ba}.
\end{equation}
重力在世界系中是常量（忽略地球曲率与自转），名义估计的导数为零，所以
\begin{equation}\label{eq:17-dg}
  \dot{\delta\bm{g}}=\bm{0}.
\end{equation}
式~\eqref{eq:17-dg} 看起来"重力误差不会变"，但它会通过式~\eqref{eq:17-dv} 影响速度；而且在流形参数化下，$\bm{g}$ 的误差被约束在切空间（球面的切平面）里，模长保持不变，这才是"重力是单位向量"的正确表达。

\subsection{组装 $F$ 与 $Q$}

把式~\eqref{eq:17-dtheta}、\eqref{eq:17-dp}、\eqref{eq:17-dv}、\eqref{eq:17-dbias}、\eqref{eq:17-dg} 按 $\delta\bm{x}=[\delta\bm{\theta},\delta\bm{p},\delta\bm{v},\delta\bm{b}_g,\delta\bm{b}_a,\delta\bm{g}]$ 的顺序排成矩阵 $\dot{\delta\bm{x}}=\bm{F}\delta\bm{x}+\bm{G}\bm{n}$，得到
\begin{equation}\label{eq:17-F}
  \bm{F}=
  \begin{bmatrix}
    -[\hat{\bm{\omega}}]_\times & \bm{0} & \bm{0} & -\bm{I} & \bm{0} & \bm{0}\\
    \bm{0} & \bm{0} & \bm{I} & \bm{0} & \bm{0} & \bm{0}\\
    -\hat{\bm{R}}[\tilde{\bm{a}}-\hat{\bm{b}}_a]_\times & \bm{0} & \bm{0} & \bm{0} & -\hat{\bm{R}} & \bm{I}\\
    \bm{0} & \bm{0} & \bm{0} & \bm{0} & \bm{0} & \bm{0}\\
    \bm{0} & \bm{0} & \bm{0} & \bm{0} & \bm{0} & \bm{0}\\
    \bm{0} & \bm{0} & \bm{0} & \bm{0} & \bm{0} & \bm{0}
  \end{bmatrix},
  \qquad
  \bm{G}=
  \begin{bmatrix}
    -\bm{I} & \bm{0} & \bm{0} & \bm{0}\\
    \bm{0} & \bm{0} & \bm{0} & \bm{0}\\
    \bm{0} & -\hat{\bm{R}} & \bm{0} & \bm{0}\\
    \bm{0} & \bm{0} & \bm{I} & \bm{0}\\
    \bm{0} & \bm{0} & \bm{0} & \bm{I}\\
    \bm{0} & \bm{0} & \bm{0} & \bm{0}
  \end{bmatrix}.
\end{equation}
$\bm{G}$ 的列按 $\bm{n}=[\bm{n}_g,\bm{n}_a,\bm{n}_{bg},\bm{n}_{ba}]$ 排列。$\bm{Q}$ 是过程噪声的功率谱矩阵：
\begin{equation}\label{eq:17-Q}
  \bm{Q}=\diag\bigl(\sigma_g^2\bm{I}_3,\;\sigma_a^2\bm{I}_3,\;\sigma_{bg}^2\bm{I}_3,\;\sigma_{ba}^2\bm{I}_3\bigr),
\end{equation}
其中 $\sigma_g$ 是陀螺的角度随机游走密度（单位 $\mathrm{rad/s/\sqrt{Hz}}$），$\sigma_a$ 是加速度计的速度随机游走密度（$\mathrm{m/s^2/\sqrt{Hz}}$），$\sigma_{bg},\sigma_{ba}$ 是零偏随机游走密度。这些参数一般从 IMU 数据手册的 Allan 方差曲线读出（见第~\ref{ch:03-sensors}~章）。

\begin{table}[htbp]
  \centering
  \caption{$\bm{F}$ 矩阵分块逐块解释（行/列顺序：$\delta\bm{\theta},\delta\bm{p},\delta\bm{v},\delta\bm{b}_g,\delta\bm{b}_a,\delta\bm{g}$）}\label{tab:17-F}
  \footnotesize
  \begin{tabular}{llll}
    \toprule
    块 & 表达式 & 物理来源 & 是否为零 \\
    \midrule
    $(1,1)$ $\partial\dot{\delta\bm{\theta}}/\partial\delta\bm{\theta}$ & $-[\hat{\bm{\omega}}]_\times$ & 误差向量随机体系旋转 & 非线性项 \\
    $(1,4)$ $\partial\dot{\delta\bm{\theta}}/\partial\delta\bm{b}_g$ & $-\bm{I}$ & 陀螺零偏误差直接积分 & 常数 \\
    $(2,3)$ $\partial\dot{\delta\bm{p}}/\partial\delta\bm{v}$ & $\bm{I}$ & $\dot{\bm{p}}=\bm{v}$ & 常数 \\
    $(3,1)$ $\partial\dot{\delta\bm{v}}/\partial\delta\bm{\theta}$ & $-\hat{\bm{R}}[\tilde{\bm{a}}-\hat{\bm{b}}_a]_\times$ & 姿态误差使比力方向偏转 & 关键耦合 \\
    $(3,5)$ $\partial\dot{\delta\bm{v}}/\partial\delta\bm{b}_a$ & $-\hat{\bm{R}}$ & 加速度计零偏误差 & 常数 \\
    $(3,6)$ $\partial\dot{\delta\bm{v}}/\partial\delta\bm{g}$ & $\bm{I}$ & 重力误差直接进入加速度 & 常数 \\
    $(4,4)$,$(5,5)$,$(6,6)$ & $\bm{0}$ & 零偏/重力自身不传播（仅由噪声驱动） & --- \\
    \bottomrule
  \end{tabular}
\end{table}

\subsection{协方差传播}

连续时间下协方差满足 Lyapunov 微分方程 $\dot{\bm{P}}=\bm{F}\bm{P}+\bm{P}\bm{F}^{\top}+\bm{G}\bm{Q}\bm{G}^{\top}$。工程上做一阶离散化，令 $\bm{F}_k=\bm{I}+\bm{F}\Delta t$，得到
\begin{equation}\label{eq:17-cov-prop}
  \bar{\bm{P}}_{k+1}=\bm{F}_k\bar{\bm{P}}_k\bm{F}_k^{\top}+\bm{Q}_k,\qquad
  \bm{Q}_k\approx\bm{G}\bm{Q}\bm{G}^{\top}\Delta t.
\end{equation}
把 $\bm{G}\bm{Q}\bm{G}^{\top}\Delta t$ 展开，得到每个 IMU 步注入的噪声协方差：
\begin{equation}\label{eq:17-Qd}
  \bm{Q}_k=\diag\bigl(\sigma_g^2\Delta t\,\bm{I}_3,\;\bm{0},\;\sigma_a^2\Delta t\,\bm{I}_3,\;\sigma_{bg}^2\Delta t\,\bm{I}_3,\;\sigma_{ba}^2\Delta t\,\bm{I}_3\bigr).
\end{equation}
注意 $\delta\bm{p}$ 那一块是零：位置没有直接的噪声注入，它的不确定性完全来自速度的积分。这也解释了为什么位置误差随时间按 $\sqrt{t^3}$ 增长（速度误差按 $\sqrt{t}$，积分一次变 $\sqrt{t^3}$），是惯性导航最经典的误差增长规律。

推导过程本身只有矩阵乘除，但含义很重要：$\bm{F}$ 描述"误差如何生长和互相传染"，$\bm{Q}$ 描述"每一步注入多少新误差"。$\bm{P}$ 就是这两者的累积结果。$\bm{P}$ 中非对角块（如 $\bm{P}_{\theta b_g}$）记录了"姿态误差和零偏误差有多相关"，第~\ref{sec:17-tight-remark}~节会看到它正是零偏被激光间接修正的通道。

一阶形式与 FAST-LIO2 论文附录中"对旋转块做精确指数处理"的版本在 $\Delta t\to0$ 时一致；$\Delta t=5$\,ms、转速 $1000^{\circ}$/s 时两者相差约 $10^{-4}$ 量级。本书按一阶形式推导是为了让"每一项从哪来"看得清楚，读代码时若看到形如 $\bm{A}(\bm{\omega}\Delta t)$ 的矩阵，那就是论文里精确处理的版本。

\section{迭代扩展卡尔曼滤波（IEKF）}\label{sec:17-iekf}

\subsection{标准 EKF 更新回顾}

高斯假设下，预测步给出先验 $\bm{x}\sim\mathcal{N}(\bar{\bm{x}},\bar{\bm{P}})$，观测为 $\bm{z}=\bm{h}(\bm{x})+\bm{v}$，$\bm{v}\sim\mathcal{N}(\bm{0},\bm{R})$。EKF 在 $\bar{\bm{x}}$ 处线性化 $\bm{h}$，得到
\begin{equation}\label{eq:17-ekf}
  \bm{K}=\bar{\bm{P}}\bm{H}^{\top}\bigl(\bm{H}\bar{\bm{P}}\bm{H}^{\top}+\bm{R}\bigr)^{-1},\qquad
  \bm{x}=\bar{\bm{x}}\boxplus\bm{K}\bigl(\bm{z}-\bm{h}(\bar{\bm{x}})\bigr),\qquad
  \bm{P}=(\bm{I}-\bm{K}\bm{H})\bar{\bm{P}}.
\end{equation}
EKF 只线性化一次。问题在于：在 LIO 里 $\bm{h}$ 是通过"找最近平面"定义的，而"最近"这件事本身依赖于当前位姿估计——位姿差一点，找到的平面就可能不同。只线性化一次，等于把"找平面"和"估位姿"割裂开，初值差时容易收敛到错误的对应关系。

\subsection{从最大后验估计推导迭代更新}

IEKF 的想法是：把更新步看成一个最大后验（MAP）问题，用高斯-牛顿反复求解，每次迭代都在新的估计点上重新线性化、重新找平面。要求解的是
\begin{equation}\label{eq:17-map}
  \bm{x}^{\star}=\argmin_{\bm{x}}\;\frac12\norm{\bm{x}\boxminus\bar{\bm{x}}}^{2}_{\bar{\bm{P}}}
  +\frac12\norm{\bm{z}-\bm{h}(\bm{x})}^{2}_{\bm{R}},
\end{equation}
其中第一项是"别离先验太远"，第二项是"要拟合观测"，$\norm{\bm{a}}^{2}_{\bm{\Sigma}}=\bm{a}^{\top}\bm{\Sigma}^{-1}\bm{a}$。式~\eqref{eq:17-map} 就是把贝叶斯公式 $p(\bm{x}|\bm{z})\propto p(\bm{z}|\bm{x})p(\bm{x})$ 取负对数后忽略常数项的结果——这也说明了滤波器与优化本质上是同一个问题，只是滤波器每来一帧就消掉历史。

设当前迭代点为 $\bm{x}_j$，令 $\bm{x}=\bm{x}_j\boxplus\delta\bm{x}$。在流形上做一阶展开：
\begin{equation}\label{eq:17-linearize}
  \bm{x}\boxminus\bar{\bm{x}}\approx\bigl(\bm{x}_j\boxminus\bar{\bm{x}}\bigr)+\delta\bm{x},\qquad
  \bm{h}(\bm{x})\approx\bm{h}(\bm{x}_j)+\bm{H}_j\delta\bm{x},\qquad
  \bm{H}_j=\left.\frac{\partial\bm{h}}{\partial\delta\bm{x}}\right|_{\bm{x}_j}.
\end{equation}
第一个近似能用，是因为 $\boxplus$ 在切空间里就是加法：从 $\bar{\bm{x}}$ 出发先走 $\bm{x}_j\boxminus\bar{\bm{x}}$ 再走 $\delta\bm{x}$，一阶上等于直接走两者之和。代入式~\eqref{eq:17-map}，代价函数变成 $\delta\bm{x}$ 的二次型：
\begin{equation}
  J(\delta\bm{x})=\frac12\norm{\bigl(\bm{x}_j\boxminus\bar{\bm{x}}\bigr)+\delta\bm{x}}^{2}_{\bar{\bm{P}}}
  +\frac12\norm{\bm{z}-\bm{h}(\bm{x}_j)-\bm{H}_j\delta\bm{x}}^{2}_{\bm{R}}.
\end{equation}
对 $\delta\bm{x}$ 求导并令其为零（这正是"正规方程"的流形版本，见第~\ref{ch:02-math}~章）：
\begin{equation}
  \bar{\bm{P}}^{-1}\bigl[(\bm{x}_j\boxminus\bar{\bm{x}})+\delta\bm{x}\bigr]
  -\bm{H}_j^{\top}\bm{R}^{-1}\bigl[\bm{z}-\bm{h}(\bm{x}_j)-\bm{H}_j\delta\bm{x}\bigr]=\bm{0},
\end{equation}
整理成
\begin{equation}\label{eq:17-normal}
  \bigl(\bar{\bm{P}}^{-1}+\bm{H}_j^{\top}\bm{R}^{-1}\bm{H}_j\bigr)\delta\bm{x}
  =-\bar{\bm{P}}^{-1}\bigl(\bm{x}_j\boxminus\bar{\bm{x}}\bigr)+\bm{H}_j^{\top}\bm{R}^{-1}\bigl(\bm{z}-\bm{h}(\bm{x}_j)\bigr).
\end{equation}
现在要解出 $\delta\bm{x}$，需要对左边那个 $18\times18$ 矩阵求逆。用两次矩阵求逆引理（Woodbury 恒等式）把它化成和 EKF 一样的形式。记 $\bm{M}=\bar{\bm{P}}^{-1}+\bm{H}_j^{\top}\bm{R}^{-1}\bm{H}_j$，先用 $(\bm{A}+\bm{BCD})^{-1}=\bm{A}^{-1}-\bm{A}^{-1}\bm{B}(\bm{C}^{-1}+\bm{DA}^{-1}\bm{B})^{-1}\bm{DA}^{-1}$（取 $\bm{A}=\bar{\bm{P}}^{-1}$、$\bm{B}=\bm{H}_j^{\top}$、$\bm{C}=\bm{R}^{-1}$、$\bm{D}=\bm{H}_j$）：
\begin{equation}
  \bm{M}^{-1}=\bar{\bm{P}}-\bar{\bm{P}}\bm{H}_j^{\top}\bigl(\bm{R}+\bm{H}_j\bar{\bm{P}}\bm{H}_j^{\top}\bigr)^{-1}\bm{H}_j\bar{\bm{P}}.
\end{equation}
右乘 $\bm{H}_j^{\top}\bm{R}^{-1}$：
\begin{align}
  \bm{M}^{-1}\bm{H}_j^{\top}\bm{R}^{-1}
  &=\bar{\bm{P}}\bm{H}_j^{\top}\bm{R}^{-1}-\bar{\bm{P}}\bm{H}_j^{\top}\bigl(\bm{R}+\bm{H}_j\bar{\bm{P}}\bm{H}_j^{\top}\bigr)^{-1}\bm{H}_j\bar{\bm{P}}\bm{H}_j^{\top}\bm{R}^{-1}\nonumber\\
  &=\bar{\bm{P}}\bm{H}_j^{\top}\Bigl[\bm{I}-\bigl(\bm{R}+\bm{H}_j\bar{\bm{P}}\bm{H}_j^{\top}\bigr)^{-1}\bm{H}_j\bar{\bm{P}}\bm{H}_j^{\top}\Bigr]\bm{R}^{-1}\nonumber\\
  &=\bar{\bm{P}}\bm{H}_j^{\top}\bigl(\bm{R}+\bm{H}_j\bar{\bm{P}}\bm{H}_j^{\top}\bigr)^{-1}\bm{R}\,\bm{R}^{-1}
  =\bm{K}_j,
\end{align}
倒数第二步用了 $\bm{I}-(\bm{R}+\bm{S})^{-1}\bm{S}=(\bm{R}+\bm{S})^{-1}\bm{R}$（记 $\bm{S}=\bm{H}_j\bar{\bm{P}}\bm{H}_j^{\top}$）。同理右乘 $\bar{\bm{P}}^{-1}$：
\begin{equation}
  \bm{M}^{-1}\bar{\bm{P}}^{-1}=\bm{I}-\bar{\bm{P}}\bm{H}_j^{\top}\bigl(\bm{R}+\bm{H}_j\bar{\bm{P}}\bm{H}_j^{\top}\bigr)^{-1}\bm{H}_j=\bm{I}-\bm{K}_j\bm{H}_j.
\end{equation}
两式代回式~\eqref{eq:17-normal}，得到迭代更新公式：
\begin{equation}\label{eq:17-iekf-update}
  \boxed{
  \delta\bm{x}_j=\bm{K}_j\bigl(\bm{z}-\bm{h}(\bm{x}_j)\bigr)-\bigl(\bm{I}-\bm{K}_j\bm{H}_j\bigr)\bigl(\bm{x}_j\boxminus\bar{\bm{x}}\bigr),\qquad
  \bm{x}_{j+1}=\bm{x}_j\boxplus\delta\bm{x}_j,}
\end{equation}
其中 $\bm{K}_j=\bar{\bm{P}}\bm{H}_j^{\top}(\bm{H}_j\bar{\bm{P}}\bm{H}_j^{\top}+\bm{R})^{-1}$ 每次都重新计算（因为 $\bm{H}_j$ 变了），而 $\bar{\bm{P}}$ 在整个迭代过程中保持不变。

式~\eqref{eq:17-iekf-update} 有两项，物理含义清楚：第一项 $\bm{K}_j(\bm{z}-\bm{h}(\bm{x}_j))$ 是"按当前残差做一次标准的 EKF 修正"；第二项 $(\bm{I}-\bm{K}_j\bm{H}_j)(\bm{x}_j\boxminus\bar{\bm{x}})$ 是"把当前迭代点相对先验的偏移扣掉"——因为 $\bm{K}_j$ 的设计是让更新"从 $\bar{\bm{x}}$ 出发"，而我们现在站在 $\bm{x}_j$ 上，必须先把这段位移还回去，否则会重复计算先验信息。初值取 $\bm{x}_0=\bar{\bm{x}}$ 时第二项为零，第一次迭代就退化成标准 EKF 更新；随着迭代进行，$\bm{x}_j\boxminus\bar{\bm{x}}$ 不再为零，第二项开始起作用，保证每次迭代都在解同一个 MAP 问题而不是越走越远。

\subsection{收敛判据与实现细节}

迭代终止条件通常取下面三条之一（或组合）：
\begin{enumerate}
  \item 更新量足够小：$\norm{\delta\bm{x}_j}<\varepsilon_x$（典型 $\varepsilon_x=10^{-3}$，各分量单位不同，实现中常按块分别设阈值）；
  \item 达到最大迭代次数 $j_{\max}$（典型取 $4$，也有取 $3\sim5$ 的实现）；
  \item 残差平方和不再明显下降。
\end{enumerate}
FAST-LIO2 的工程实现主要用第 2 条：固定迭代若干次，因为残差与协方差都在实时变化，动态判据的收益不大。理论上，当 $\bm{H}_j$ 变化很小时，迭代式~\eqref{eq:17-iekf-update} 会以二次速率收敛到 MAP 解——这是高斯-牛顿法的标准结论；实际中限制收敛速度的是"最近邻对应关系"的离散跳变：某次迭代中某个点换了平面，$\bm{h}$ 就会跳一下。

更新完成后协方差按 $\bm{P}=(\bm{I}-\bm{K}\bm{H})\bar{\bm{P}}$ 计算（$\bm{K},\bm{H}$ 用最后一次迭代的值），或者用数值更稳的 Joseph 形式
\begin{equation}\label{eq:17-joseph}
  \bm{P}=(\bm{I}-\bm{K}\bm{H})\bar{\bm{P}}(\bm{I}-\bm{K}\bm{H})^{\top}+\bm{K}\bm{R}\bm{K}^{\top},
\end{equation}
它在有限精度下能保证 $\bm{P}$ 对称正定。Point-LIO 因为每次只做一维观测更新，特别依赖 Joseph 形式来避免协方差失去正定性。

\begin{pseudo}{FAST-LIO2 主循环（简化）}
  \item 初始化：读第一帧点云建立初始地图；用静止段估计初始重力方向与零偏
  \item 收到 IMU 测量 $(\tilde{\bm{\omega}}_k,\tilde{\bm{a}}_k)$：按式~\eqref{eq:17-discrete} 传播状态，按式~\eqref{eq:17-cov-prop} 传播协方差；发布 odom
  \item 收到一帧点云：用传播轨迹把每个点去畸变到扫描末时刻，再按距离/体素降采样
  \item 对每个点：在 iKD-Tree 中找最近 $5$ 点，PCA 拟合平面得到 $(\bm{n}_i,\bm{q}_i)$，算残差 $r_i$ 与雅可比 $\bm{H}_i$
  \item 令 $\bar{\bm{x}}=\bm{x}$（传播结果），$\bm{x}_0=\bar{\bm{x}}$，迭代 $j=0,1,\dots,j_{\max}$：用式~\eqref{eq:17-iekf-update} 计算 $\delta\bm{x}_j$ 并更新 $\bm{x}_{j+1}$
  \item 用最后一次的 $\bm{H}$ 按式~\eqref{eq:17-joseph} 更新协方差 $\bm{P}$
  \item 把去畸变后的点用最终位姿变换到世界系，插入 iKD-Tree；用 box-wise 删除清掉远离当前位置的老点
\end{pseudo}

\section{点到面观测模型与雅可比}\label{sec:17-jacobian}

\subsection{数据关联：从最近邻到平面}

对去畸变后的第 $i$ 个点，令它在雷达系中的坐标为 $\bm{p}_i^{L}$。先把它转到 IMU 系：
\begin{equation}\label{eq:17-extrinsic}
  \bm{p}_i^{I}=\bm{R}_L\,\bm{p}_i^{L}+\bm{t}_L,
\end{equation}
其中 $(\bm{R}_L,\bm{t}_L)$ 是 LiDAR 到 IMU 的外参。再转到世界系：
\begin{equation}\label{eq:17-world}
  \bm{p}_i^{W}=\bm{R}\,\bm{p}_i^{I}+\bm{p}.
\end{equation}
在 iKD-Tree 中取 $\bm{p}_i^{W}$ 的最近 $5$ 个地图点 $\{\bm{q}_{i,1},\dots,\bm{q}_{i,5}\}$，计算它们的均值与协方差：
\begin{equation}\label{eq:17-planefit}
  \bar{\bm{q}}_i=\frac15\sum_{s=1}^{5}\bm{q}_{i,s},\qquad
  \bm{\Sigma}_i=\frac15\sum_{s=1}^{5}\bigl(\bm{q}_{i,s}-\bar{\bm{q}}_i\bigr)\bigl(\bm{q}_{i,s}-\bar{\bm{q}}_i\bigr)^{\top}.
\end{equation}
对 $\bm{\Sigma}_i$ 做特征分解，最小特征值对应的特征向量取作平面法向 $\bm{n}_i$（$\norm{\bm{n}_i}=1$），平面过点 $\bar{\bm{q}}_i$。如果最小特征值不够小（点不共面，如树叶、边缘处），该点应被剔除——这就是最朴素的平面性检验。取 5 个点而不是 1 个，是因为 5 个点刚好够定出一个平面（3 个自由度）并附带一个可判断"平不平"的协方差。

\subsection{残差与雅可比}

点到面残差定义为世界系中点到平面的有符号距离：
\begin{equation}\label{eq:17-residual}
  r_i=\bm{n}_i^{\top}\bigl(\bm{R}\,\bm{p}_i^{I}+\bm{p}-\bar{\bm{q}}_i\bigr)
  =\bm{n}_i^{\top}\bigl(\bm{p}_i^{W}-\bar{\bm{q}}_i\bigr).
\end{equation}
式~\eqref{eq:17-residual} 的几何含义：$\bm{n}_i^{\top}(\cdot)$ 把三维误差投影到法向上，只保留"垂直于平面"的那一个分量；平面内的偏差被投影掉了，因为平面内往哪儿挪都还在同一个平面上，这个残差无法也不应该约束它。这也正是点到面残差只提供 1 维约束、$m$ 个点最多提供 $m$ 维约束的原因。

下面求偏导。用右扰动约定 $\bm{R}=\hat{\bm{R}}\exp([\delta\bm{\theta}]_\times)$（与 \S\ref{sec:17-f} 一致），$\bm{p}=\hat{\bm{p}}+\delta\bm{p}$。把扰动形式代入式~\eqref{eq:17-residual}：
\begin{align}
  r_i(\delta\bm{x})&=\bm{n}_i^{\top}\Bigl(\hat{\bm{R}}\exp([\delta\bm{\theta}]_\times)\bm{p}_i^{I}+\hat{\bm{p}}+\delta\bm{p}-\bar{\bm{q}}_i\Bigr)\nonumber\\
  &\approx\bm{n}_i^{\top}\Bigl(\hat{\bm{R}}\bigl(\bm{I}+[\delta\bm{\theta}]_\times\bigr)\bm{p}_i^{I}+\hat{\bm{p}}+\delta\bm{p}-\bar{\bm{q}}_i\Bigr)\nonumber\\
  &=r_i(\bm{0})+\bm{n}_i^{\top}\hat{\bm{R}}[\delta\bm{\theta}]_\times\bm{p}_i^{I}+\bm{n}_i^{\top}\delta\bm{p}.
\end{align}
再用 $[\delta\bm{\theta}]_\times\bm{p}_i^{I}=\delta\bm{\theta}\times\bm{p}_i^{I}=-[\bm{p}_i^{I}]_\times\delta\bm{\theta}$（这是"叉乘的反对称矩阵表示"，见第~\ref{ch:02-math}~章）：
\begin{equation}\label{eq:17-jacobian}
  r_i(\delta\bm{x})\approx r_i(\bm{0})
  \underbrace{-\bm{n}_i^{\top}\hat{\bm{R}}[\bm{p}_i^{I}]_\times}_{\partial r_i/\partial\delta\bm{\theta}}\delta\bm{\theta}
  +\underbrace{\bm{n}_i^{\top}}_{\partial r_i/\partial\delta\bm{p}}\delta\bm{p}.
\end{equation}
于是第 $i$ 个点的雅可比是一个 $1\times18$ 的行向量：
\begin{equation}\label{eq:17-Hi}
  \boxed{
  \bm{H}_i=\frac{\partial r_i}{\partial\delta\bm{x}}
  =\begin{bmatrix}
    -\bm{n}_i^{\top}\hat{\bm{R}}[\bm{p}_i^{I}]_\times, & \bm{n}_i^{\top}, & \bm{0}_{3}^{\top}, & \bm{0}_{3}^{\top}, & \bm{0}_{3}^{\top}, & \bm{0}_{3}^{\top}
  \end{bmatrix},}
\end{equation}
按 $[\delta\bm{\theta},\delta\bm{p},\delta\bm{v},\delta\bm{b}_g,\delta\bm{b}_a,\delta\bm{g}]$ 分块，非零的只有前两块。

逐块解释：$\partial r_i/\partial\delta\bm{\theta}=-\bm{n}_i^{\top}\hat{\bm{R}}[\bm{p}_i^{I}]_\times$ 表"姿态转一点，点在世界上移动 $\bm{R}\delta\bm{\theta}\times\bm{p}_i^{I}$，再投影到法向"；$\partial r_i/\partial\delta\bm{p}=\bm{n}_i^{\top}$ 表"位置挪一点，点在世界上跟着挪，投影到法向"；$\delta\bm{v},\delta\bm{b}_g,\delta\bm{b}_a,\delta\bm{g}$ 的块为零，因为点的坐标变换直接看不出速度、零偏、重力——它们只能通过先验协方差间接被修正（见 \S\ref{sec:17-tight-remark}）。

\begin{table}[htbp]
  \centering
  \caption{点到面观测的雅可比分块与含义}\label{tab:17-jacobian}
  \small
  \begin{tabular}{llll}
    \toprule
    分块 & 维度 & 表达式 & 含义 \\
    \midrule
    $\partial r_i/\partial\delta\bm{\theta}$ & $1\times3$ & $-\bm{n}_i^{\top}\hat{\bm{R}}[\bm{p}_i^{I}]_\times$ & 姿态扰动使点位置绕机体系转动 \\
    $\partial r_i/\partial\delta\bm{p}$ & $1\times3$ & $\bm{n}_i^{\top}$ & 平移沿法向的分量 \\
    $\partial r_i/\partial\delta\bm{v}$ & $1\times3$ & $\bm{0}$ & 点坐标不依赖速度 \\
    $\partial r_i/\partial\delta\bm{b}_g$ & $1\times3$ & $\bm{0}$ & 通过先验协方差间接作用 \\
    $\partial r_i/\partial\delta\bm{b}_a$ & $1\times3$ & $\bm{0}$ & 通过先验协方差间接作用 \\
    $\partial r_i/\partial\delta\bm{g}$ & $1\times3$ & $\bm{0}$ & 通过先验协方差间接作用 \\
    \midrule
    $\partial r_i/\partial\bm{R}_L$ & $1\times3$ & $-\bm{n}_i^{\top}\hat{\bm{R}}\hat{\bm{R}}_L[\bm{p}_i^{L}]_\times$ & 仅在外参可估时存在 \\
    $\partial r_i/\partial\bm{t}_L$ & $1\times3$ & $\bm{n}_i^{\top}\hat{\bm{R}}$ & 仅在外参可估时存在 \\
    \bottomrule
  \end{tabular}
\end{table}

\begin{remark}[左扰动与右扰动]
如果改用左扰动 $\bm{R}=\exp([\delta\bm{\theta}]_\times)\hat{\bm{R}}$，同样的一阶展开给出
$r_i\approx r_i(\bm{0})-\bm{n}_i^{\top}[\hat{\bm{R}}\bm{p}_i^{I}]_\times\delta\bm{\theta}$，即雅可比是 $-\bm{n}_i^{\top}[\bm{p}_i^{W}-\hat{\bm{p}}]_\times$。两种约定的差别只是把旋转误差表达在世界系还是机体系，数值结果完全等价（相差一个 $\bm{R}$ 的相似变换）。\textbf{关键是不能混用}：如果雅可比按右扰动推导，更新步就必须按右扰动施加（$\bm{R}\leftarrow\bm{R}\exp(\delta\bm{\theta})$），否则会引入一个随姿态变化的系统性错误。FAST-LIO2 的代码用右扰动，本书全程保持一致。
\end{remark}

\subsection{堆叠与退化分析}

一帧里有 $m$ 个有效点，把它们堆起来：
\begin{equation}\label{eq:17-Hstack}
  \bm{r}=\begin{bmatrix}r_1\\\vdots\\r_m\end{bmatrix}\in\R^m,\qquad
  \bm{H}=\begin{bmatrix}\bm{H}_1\\\vdots\\\bm{H}_m\end{bmatrix}\in\R^{m\times18},\qquad
  \bm{R}_m=\diag(\sigma_1^2,\dots,\sigma_m^2).
\end{equation}
观测方程是 $\bm{0}=\bm{r}(\bm{x})+\bm{v}$，即"理想情况下残差应为零"，所以式~\eqref{eq:17-iekf-update} 里的 $\bm{z}$ 取 $\bm{0}$、$\bm{h}(\bm{x}_j)$ 就是 $\bm{r}(\bm{x}_j)$——这是把"配准"写进滤波器的标准技巧：残差的期望是零，方差是点的测量噪声。

只看平移部分，$\bm{H}^{\top}\bm{H}$ 的平移块是 $\sum_i\bm{n}_i\bm{n}_i^{\top}$。如果所有点都在同一面墙上，所有 $\bm{n}_i$ 相同，这个矩阵的秩是 1：\textbf{沿墙面方向的两个平移没有被任何残差约束}，这就是几何退化。如果场景有两面互相垂直的墙，秩变成 2；三面互相垂直的墙才满秩。这个观察直接给出退化检测方法：对 $\sum_i\bm{n}_i\bm{n}_i^{\top}$ 做特征分解，若最小特征值小于阈值（或小于最大特征值的某个比例），说明机器人在该特征向量方向上"看不见"约束，此时应当依赖 IMU 先验，或者降低该方向的更新增益。

为什么用点到面而不是点到点？点到点残差是 $\bm{p}_i^{W}-\bm{q}_i$ 的三维向量，表面上看提供了 3 维约束，但地图点本身有采样噪声，点到点残差会把"地图点采在哪"的误差当成位姿误差，导致收敛慢、精度低（经典 ICP 的问题，见第~\ref{ch:07-lidar-odom}~章）。点到面残差把沿平面方向的噪声投影掉，只保留真正有信息的法向分量，所以同样的点能给出更准的位姿。代价是：法向估计本身有误差，平面拟合不可靠的点必须剔除，否则会引入系统偏差。

\subsection{零偏与重力如何被间接修正}\label{sec:17-tight-remark}

表~\ref{tab:17-jacobian} 里 $\delta\bm{v},\delta\bm{b}_g,\delta\bm{b}_a,\delta\bm{g}$ 的分块全为零，很多读者由此误以为"激光修正不了零偏"。修正其实走的是先验协方差这条路：把 $\bm{H}$ 分成位姿块 $\bm{H}_p$ 与零块，则 $\bm{K}$ 的零偏部分为
\begin{equation}\label{eq:17-K-bias}
  \bm{K}_{b}=\bar{\bm{P}}_{bp}\bm{H}_p^{\top}\bigl(\bm{H}_p\bar{\bm{P}}_{p}\bm{H}_p^{\top}+\bm{R}\bigr)^{-1},
\end{equation}
其中 $\bar{\bm{P}}_{bp}$ 是零偏与位姿的互协方差。只要 $\bar{\bm{P}}_{bp}\neq\bm{0}$，零偏就会被激光残差拉动。用一维例子验证：设 $\theta_{k+1}=\theta_k-b_k\Delta t$、$b_{k+1}=b_k$，则 $\bm{F}_k=\begin{bmatrix}1&-\Delta t\\0&1\end{bmatrix}$，传播一步后 $\bar{P}_{\theta b}=P_{\theta b}-\Delta tP_{bb}$——只要 $P_{bb}>0$ 就非零。此时观测到角度（$\bm{h}=[1,0]$）给出 $K_b=\bar{P}_{\theta b}/(\bar{P}_{\theta\theta}+\sigma^2)\neq0$，零偏被修正，尽管观测对零偏的雅可比是零。

这就是"紧耦合"在数学上的真正含义：\textbf{两个传感器共享同一个状态和同一个协方差矩阵，一个传感器的观测可以通过相关性修正另一个传感器的状态}。$\bm{F}$ 的非对角块在传播中建立起来的 $P_{\theta b_g},P_{v b_a},P_{v g}$ 就是这条通道；没有传播，就没有相关性，也就没有紧耦合。

\section{iKD-Tree：增量式地图}\label{sec:17-ikdtree}

\subsection{为什么不能每帧重建}

FAST-LIO2 每帧要做的事是：对每个去畸变后的点，在地图里找最近的 $5$ 个点。地图点数 $n$ 随运行时间增长，典型值 $10^5\sim10^6$。若每帧都把地图点重新建一棵 KD 树，建树复杂度 $O(n\log n)$；$n=10^6$、$10$\,Hz 时每帧约需 $10^8$ 次比较，不可能实时。而实际上每帧只新增几千到几万个点、只删掉远处的一小部分——为一小部分变化付出重建整棵树的代价显然浪费。iKD-Tree（incremental k-d tree）的思路就是：\textbf{把"建树"拆成"插入、删除、局部重建"三种便宜的操作}，让树一帧一帧地长出来。

\subsection{节点里存什么}

iKD-Tree 的每个节点存五类信息：\textbf{一个地图点}；\textbf{分割轴与分割值}（用哪一维 $x/y/z$ 把子树分成两半）；\textbf{轴对齐包围盒}（AABB，即子树所有点的最小/最大坐标）；\textbf{子树规模}（有效点数与总点数，供平衡与重建判断）；\textbf{删除标记}。AABB 是最近邻搜索能剪枝的关键：若"查询点到某子树 AABB 的最近距离"已大于当前找到的第 $k$ 近邻距离，整棵子树都不用看；没有 AABB，KD 树就退化成链表。

\subsection{增量插入}

插入一个点 $\bm{q}$ 的过程就是从根往下走：遇到包围盒为空的节点就新建节点存放 $\bm{q}$；否则按该节点的分割轴取出 $\bm{q}$ 的对应坐标，与分割值比较决定去左还是右子树，递归直到插入完成，再沿路径回溯更新每个祖先的 AABB 与计数。插入后做平衡检查：若某节点左右子树规模之比超过 $\alpha$（典型 $\alpha=0.7$），说明树开始"歪"，把该子树的所有有效点收集起来重建为一棵平衡树。这个条件与"替罪羊树"（scapegoat tree）的思想一致：大部分插入不破坏平衡，偶尔一次重建把子树摊平，均摊到每次插入仍是 $O(\log n)$。若被删除点占该子树的比例超过阈值（典型 $30\%$），同样触发重建，把删除标记真正"兑现"掉——否则树里全是"幽灵节点"，搜索时白白遍历。

\subsection{惰性删除与 box-wise 操作}

删除一个点在 KD 树里是很麻烦的：不能直接摘掉节点（会破坏树的结构），重建整棵树又太贵。iKD-Tree 用\textbf{惰性删除}（lazy delete）：只把节点标记为"已删除"，搜索时跳过它，树的形状不变。这样做的问题是被删节点还占着内存、还参与遍历；解决办法是在满足条件时重建（见上一段）。

更妙的是\textbf{box-wise 操作}。LIO 运行时经常需要"删掉离当前位置超过 $R$ 的所有点"或"把某个立方体里的点全部插入"。这类操作对应一个轴对齐盒子，iKD-Tree 可以直接利用 AABB 判断：
\begin{itemize}
  \item \textbf{box-wise 删除}：从根开始，若某节点的 AABB 完全落在删除盒内，直接把这棵子树整体标记删除（$O(1)$），不再往下递归；若不相交，直接返回；只有部分相交时才递归到子节点。
  \item \textbf{box-wise 插入}：类似地，若插入盒完全包含某子树，可以把新点集整体挂成一个子树，而不是逐个插入。
\end{itemize}
一次 box-wise 删除的实际代价与"落在盒边界上的节点数"成正比，而不是与盒内点数成正比——这正是 LIO 里"滑窗式局部地图"能做到实时的原因。

\begin{figure}[htbp]
  \centering
  \begin{tikzpicture}[font=\footnotesize]
    % 左：二维空间划分
    \draw[thick, black!60] (0,0) rectangle (6,4);
    \draw[thick, structurecolor] (3,0) -- (3,4);
    \draw[thick, structurecolor] (0,2.3) -- (3,2.3);
    \draw[thick, structurecolor] (3,1.7) -- (6,1.7);
    \draw[thick, structurecolor] (1.5,2.3) -- (1.5,4);
    \draw[thick, structurecolor] (4.5,1.7) -- (4.5,0);
    \foreach \x/\y in {0.6/0.7,1.1/1.3,2.2/0.9,2.6/1.4,3.8/0.5,4.2/1.1,5.3/0.6,5.6/1.3,
                       1.0/3.1,1.9/3.5,2.5/2.8,3.6/3.2,4.2/2.6,5.2/3.3}{
      \fill[black!70] (\x,\y) circle (1.3pt);
    }
    \draw[red!60!black, thick, dashed] (3.1,0.1) rectangle (5.9,1.6);
    \node[label, text=red!60!black] at (4.5,0.35) {box-wise 删除};
    \node[label] at (3,-0.55) {(a) 平面划分与删除盒};
    % 右：树结构
    \begin{scope}[xshift=7.9cm, yshift=1.1cm]
      \node[draw=structurecolor, fill=structurecolor!10, rounded corners=1pt, inner sep=1.5pt] (r) at (2.5,2.6) {轴 $x$};
      \node[draw=structurecolor, fill=structurecolor!10, rounded corners=1pt, inner sep=1.5pt] (l) at (1.1,1.3) {轴 $y$};
      \node[draw=structurecolor, fill=structurecolor!10, rounded corners=1pt, inner sep=1.5pt] (rr) at (3.9,1.3) {轴 $y$};
      \node[draw=black!40, fill=white, rounded corners=1pt, inner sep=1.5pt] (ll) at (0.2,0) {叶};
      \node[draw=black!40, fill=white, rounded corners=1pt, inner sep=1.5pt] (lr) at (1.9,0) {叶};
      \node[draw=black!40, fill=white, rounded corners=1pt, inner sep=1.5pt] (rl) at (3.0,0) {叶};
      \node[draw=red!60!black, fill=red!60!black!12, rounded corners=1pt, inner sep=1.5pt] (rrt) at (4.9,0) {子树};
      \draw[black!60] (r) -- (l);
      \draw[black!60] (r) -- (rr);
      \draw[black!60] (l) -- (ll);
      \draw[black!60] (l) -- (lr);
      \draw[black!60] (rr) -- (rl);
      \draw[black!60] (rr) -- (rrt);
      \node[label] at (2.5,-0.7) {(b) 对应的二叉树：红框子树整体标记删除};
    \end{scope}
  \end{tikzpicture}
  \caption{iKD-Tree 的空间划分与树结构。左图是一条条按轴切分的划分线，红框表示一次"删除这个盒子里的点"的请求；右图中红框对应的整棵子树只需在根上打一个删除标记，搜索时遇到它也只需检查一次标记，不用逐点遍历。}\label{fig:17-ikdtree}
\end{figure}

\subsection{复杂度与内存}

把每次操作摊开算：普通 KD 树每帧重建要 $O(n\log n)$，而 iKD-Tree 只在初始化时建一次树，之后每次插入均摊 $O(\log n)$、单点删除 $O(1)$（打标记）、box-wise 删除约 $O(\text{边界节点数})$、$k$ 近邻查询 $O(\log n)$（含剪枝）；最坏情况由 $\alpha$ 平衡保证不退化。这正是把"每帧 $O(m\log n)$ 的查询 + $O(n\log n)$ 的重建"变成"只有查询和增量插入"的关键。

"地图只增不减"会带来内存问题：$10^5$ 点/帧、$10$\,Hz，一分钟就是 $6\times10^7$ 个点。FAST-LIO2 的对策是\textbf{局部地图}：只维护以机器人当前位置为中心、边长 \texttt{cube\_side\_length}（典型 $100\sim200$\,m）的立方体内的点，超出范围的用 box-wise 删除清掉；此外当前帧按 \texttt{filter\_size\_surf} 体素降采样（典型 $0.5$\,m），插入地图前再按 \texttt{filter\_size\_map} 降采样（典型 $0.5$\,m）。这样地图点数被限制在"场景表面积 / 体素面积"的量级，跑几小时也不会爆内存。并行方面，建树与查询常用 OpenMP 在左右子树上分任务，实际加速比受内存带宽限制。

\section{Point-LIO：逐点更新}\label{sec:17-pointlio}

\subsection{FAST-LIO2 留下的两个"帧"假设}

FAST-LIO2 虽然把状态更新写得很漂亮，但它有两个隐含假设：其一，\textbf{帧内 IMU 可信}——$100$\,ms 内的状态全靠 IMU 积分外推，直到帧末才被激光修正一次，若这 $100$\,ms 里陀螺饱和（读数被截断在量程上限）或遭遇撞击（加速度计超量程），外推就会跑偏，帧末的数据关联（找最近平面）也会找错，滤波器可能被错误观测带得更偏；其二，\textbf{更新是帧同步的}——两次更新之间的 $100$\,ms 内位姿本质上在"盲推"，高速旋转时 $100$\,ms 就是几十度。直观地说：机器人以 $2000^{\circ}$/s 旋转（很多 MEMS 陀螺的量程上限）时，$100$\,ms 就是 $200^{\circ}$——半圈，这段时间里 IMU 一旦饱和，谁也不知道它转了多少。

\subsection{核心思想：点一到就更新}

Point-LIO~\cite{he2023pointlio} 的做法非常直接：\textbf{每收到一个激光点，就做一次状态更新}。激光点的到达间隔是微秒量级（$10$\,Hz、每帧 $10^4$ 点时约 $10\,\mu$s），于是状态更新的频率从 $10$\,Hz 变成 $10^4\sim10^5$\,Hz。状态不再是"每帧修正一次"，而是"被每个点持续地拽着走"。

\begin{figure}[htbp]
  \centering
  \begin{tikzpicture}[font=\footnotesize]
    % 时间轴
    \draw[->, black!60] (0,-0.2) -- (12.6,-0.2) node[right]{$t$};
    % 上：FAST-LIO2
    \node[label, anchor=west] at (0,1.9) {\textbf{FAST-LIO2}};
    \draw[thick, structurecolor] (0.2,1.2) -- (12.2,1.2);
    \foreach \x in {0.2,0.6,...,12.2} { \draw[structurecolor!70] (\x,1.12) -- (\x,1.28); }
    \foreach \x in {3.2,6.8,10.4} {
      \fill[red!60!black] (\x,1.2) -- (\x-0.12,1.48) -- (\x+0.12,1.48) -- cycle;
      \draw[red!60!black, dashed] (\x,1.2) -- (\x,0.05);
    }
    \node[label, text=red!60!black, align=left] at (1.6,1.55) {只在帧末更新（每 $100$\,ms 一次）};
    % 下：Point-LIO
    \node[label, anchor=west] at (0,-0.85) {\textbf{Point-LIO}};
    \draw[thick, structurecolor] (0.2,-1.5) -- (12.2,-1.5);
    \foreach \x in {0.2,0.6,...,12.2} {
      \fill[green!45!black] (\x,-1.5) circle (1.4pt);
    }
    \node[label, text=green!45!black, align=left] at (1.6,-1.95) {每个点一次标量更新（$10\,\mu$s 量级一次）};
    % 帧边界
    \foreach \x in {0.2,3.2,6.8,10.4} {
      \draw[black!35, dotted] (\x,0.05) -- (\x,-0.2);
    }
    \draw[decorate, decoration={brace, amplitude=4pt}, black!60] (3.2,2.05) -- (6.8,2.05)
      node[midway, above=4pt, label] {一帧（$100$\,ms）};
  \end{tikzpicture}
  \caption{两种更新粒度的时序对比。上：FAST-LIO2 在每帧末尾做一次 $m$ 维（多残差）更新，两次更新之间状态只靠 IMU 外推；下：Point-LIO 对每个点做一次一维（单残差）更新，状态几乎连续地被观测修正。}\label{fig:17-pointlio}
\end{figure}

\subsection{逐点传播与标量更新}

设在 $t_{i-1}$ 时刻状态为 $\bm{x}_{i-1}$、协方差为 $\bm{P}_{i-1}$，第 $i$ 个点的时间戳为 $t_i$。一次逐点处理的完整流程是：

\textbf{（1）传播。} 用 $[t_{i-1},t_i]$ 之间的 IMU 数据把状态推到 $t_i$：
\begin{equation}\label{eq:17-point-prop}
  \bar{\bm{x}}_i=\bm{f}\bigl(\bm{x}_{i-1},\,\mathcal{U}_{[t_{i-1},t_i]}\bigr),\qquad
  \bar{\bm{P}}_i=\bm{F}_i\bm{P}_{i-1}\bm{F}_i^{\top}+\bm{Q}_i,
\end{equation}
当 $t_i-t_{i-1}$ 只有几十微秒时，通常这段区间内没有 IMU 采样，直接用最近一次的 IMU 测量外推即可（$\bm{F}_i\approx\bm{I}+\bm{F}\Delta t_i$）。

\textbf{（2）关联。} 用 $\bar{\bm{x}}_i$ 把点变换到世界系，在 iKD-Tree 中取最近 $5$ 点、拟合平面，得到 $(\bm{n}_i,\bar{\bm{q}}_i)$ 与残差
\begin{equation}\label{eq:17-point-residual}
  r_i=\bm{n}_i^{\top}\bigl(\bm{R}_i\bm{p}_i^{I}+\bm{p}_i-\bar{\bm{q}}_i\bigr),
\end{equation}
以及雅可比 $\bm{h}_i=\partial r_i/\partial\delta\bm{x}\in\R^{1\times18}$（形式与式~\eqref{eq:17-Hi} 完全相同）。

\textbf{（3）更新。} 这是一次\textbf{一维观测}的卡尔曼更新：
\begin{equation}\label{eq:17-scalar-k}
  \bm{K}_i=\frac{\bar{\bm{P}}_i\bm{h}_i^{\top}}{\bm{h}_i\bar{\bm{P}}_i\bm{h}_i^{\top}+\sigma_i^2},\qquad
  \delta\bm{x}_i=\bm{K}_i\,r_i,\qquad
  \bm{x}_i=\bar{\bm{x}}_i\boxplus\delta\bm{x}_i,
\end{equation}
\begin{equation}\label{eq:17-scalar-p}
  \bm{P}_i=(\bm{I}-\bm{K}_i\bm{h}_i)\bar{\bm{P}}_i(\bm{I}-\bm{K}_i\bm{h}_i)^{\top}+\bm{K}_i\sigma_i^2\bm{K}_i^{\top}.
\end{equation}
式~\eqref{eq:17-scalar-k} 的推导是标准高斯条件化：$\bm{K}_i$ 就是"状态与观测的互协方差除以观测的总方差"，分母是个\textbf{标量}，所以这里没有矩阵求逆，只有一次除法。式~\eqref{eq:17-scalar-p} 用 Joseph 形式，因为每次更新的步长很小、$\bm{P}$ 的更新次数极多（每秒上万次），必须保证数值上始终对称正定。每更新一次就立刻发布当前状态，于是里程计输出频率可以做到与点频同量级（论文报告可达 $1$\,kHz 以上）。

\subsection{顺序更新与批量更新等价吗}

一个自然的问题：把一帧的 $m$ 个点拆成 $m$ 次一维更新，和 FAST-LIO2 一次 $m$ 维更新，结果一样吗？在\textbf{线性高斯}情形下答案是一样：批量更新的后验信息矩阵 $\bm{\Lambda}=\bar{\bm{P}}^{-1}+\sum_i\bm{h}_i^{\top}\bm{h}_i/\sigma_i^2$ 只是把各点贡献相加，与顺序无关；顺序更新每做一次相当于用 Sherman--Morrison 公式在信息矩阵上加同一项，做完 $m$ 次得到同一个和。所以在\textbf{线性}情形下二者的后验均值与协方差完全相同，区别只在计算路径。

在\textbf{非线性}情形下两者不再严格相等：Point-LIO 在每个点处重新线性化、重新关联，相当于把 IEKF 的迭代分散到时间上，对剧烈运动更友好；但单点的噪声、外点会直接污染状态，必须配合门限。计算上批量更新每帧 $O(m\cdot18+18^3)$，顺序更新是 $m$ 次 $O(18^2)$，后者其实更贵——逐点更新换来的不是算力节省，而是\textbf{更细的时间分辨率和更强的鲁棒性}。

\subsection{IMU 饱和时：把 IMU 从"传播输入"降级为"大噪声"}

这是 Point-LIO 最关键的设计，也是它和 FAST-LIO2 最大的区别。

正常工作时，IMU 读数是"传播的输入"：滤波器相信它，用它积分出两个点之间的状态变化。但陀螺或加速度计一旦饱和，读数被截断，继续信它就会把错误积分进去。Point-LIO 的处理思路是：\textbf{检测到饱和/不可信时，放弃"用 IMU 积分"这件事，改成随机过程模型}——也就是把状态传播写成
\begin{equation}\label{eq:17-random-walk}
  \bar{\bm{x}}_i=\bm{x}_{i-1}+\bm{w}_i,\qquad \bm{w}_i\sim\mathcal{N}(\bm{0},\bm{Q}_i^{\mathrm{sat}}),
\end{equation}
即"两个点之间的状态怎么变，我不知道，用一个很大的过程噪声 $\bm{Q}_i^{\mathrm{sat}}$ 来表示这种无知"。物理上这相当于假设姿态、速度在两个点之间做随机游走，而不是按某个确定的动力学演化。$\bm{Q}_i^{\mathrm{sat}}$ 的量级由饱和程度决定：读数越接近量程、越不可信，$\bm{Q}$ 就放得越大（论文中的做法是按饱和检测结果自适应放大；不同版本实现细节略有差异，但思想都是"越不可信，越不依赖它"）。

为什么这样做能行？因为逐点更新的频率极高：即使"完全不知道"两个点之间怎么动的，每过 $10\,\mu$s 就有一个激光点来告诉滤波器"你现在应该在这个平面上"。用一维卡尔曼滤波的稳态解可以估出误差上界。标量情形下，每步过程噪声方差 $q$、观测噪声方差 $\sigma^2$，协方差递推为 $\bar{P}=P+q$、$P'=(P+q)\sigma^2/(P+q+\sigma^2)$；令 $P'=P$ 得二次方程 $P^2+Pq-q\sigma^2=0$，正根为
\begin{equation}\label{eq:17-riccati}
  P_\infty=\frac{-q+\sqrt{q^2+4q\sigma^2}}{2}.
\end{equation}
代入 $q=10^{-4}\,\mathrm{m^2}$（每步"完全不知道"的极端情况）、$\sigma^2=10^{-4}\,\mathrm{m^2}$（单点平面观测标准差 $1$\,cm），得 $P_\infty\approx6.2\times10^{-5}\,\mathrm{m^2}$，即 $\sqrt{P_\infty}\approx7.9$\,mm。也就是说：\textbf{即使动力学模型完全失效、每步注入 $1$\,cm 量级的不确定性，只要每个点都提供一次 $1$\,cm 精度的观测，稳态误差仍被压在 $8$\,mm 左右}——这就是"高频更新"换来的鲁棒性。代价是：如果点云也退化（长走廊等），逐点更新同样失效，只能靠其他传感器或运动约束。

\subsection{代价与收益}

\begin{table}[htbp]
  \centering
  \caption{FAST-LIO2 与 Point-LIO 的对比}\label{tab:17-compare}
  \footnotesize
  \begin{tabular}{lll}
    \toprule
    维度 & FAST-LIO2 & Point-LIO \\
    \midrule
    更新粒度 & 每帧一次（$m$ 维） & 每个点一次（$1$ 维） \\
    更新频率 & $10$\,Hz & $10^4\sim10^5$\,Hz（点频） \\
    odom 输出 & IMU 频率（$100\sim200$\,Hz） & 点频（可达 $1$\,kHz 以上） \\
    单次更新代价 & $O(m\cdot18+18^3)$ & $O(18^2)$ \\
    单帧总代价 & 较低 & 较高（每点一次矩阵更新） \\
    帧内 IMU 依赖 & 强（$100$\,ms 全靠它外推） & 弱（每个点都在纠偏） \\
    IMU 饱和 & 可能发散 & 降级为随机过程模型，通常仍稳定 \\
    外点敏感性 & 一帧内平均，较鲁棒 & 单点直接进状态，需门限剔除 \\
    实现复杂度 & 中 & 高（时间戳粒度、饱和检测、门限） \\
    适用场景 & 通用、高动态但 IMU 够用 & 超高速旋转、IMU 量程不足、需要极高频输出 \\
    \bottomrule
  \end{tabular}
\end{table}

表~\ref{tab:17-compare} 再补充两点工程经验。第一，Point-LIO 的"逐点"依赖激光点自带精确时间戳；机械式多线激光的每个点时间戳由方位角推算，扫描一圈内的时间精度是微秒级的，可以用；但某些雷达只给整帧时间戳，此时要先用帧内运动模型给每个点分配时间，效果会打折。第二，逐点更新对外点极其敏感：一个点如果落在了错误的平面（比如打到移动的人身上），它会把状态直接拽偏。实践中必须做残差门限（例如把 $|r_i|/\sigma$ 超过若干倍的点丢弃）或者用鲁棒核函数降权。

\section{VoxelMap 与地图表示对比}\label{sec:17-voxelmap}

iKD-Tree 把地图存成"点"，每次查询做 $k$ 近邻搜索再拟合平面。VoxelMap~\cite{yuan2022voxelmap} 换了一个思路：\textbf{把地图直接存成平面}。

\subsection{自适应体素与平面模型}

VoxelMap 把空间划分成体素（初始边长如 $1$\,m），对每个体素内的点做 PCA：算出均值 $\bar{\bm{q}}$ 与协方差矩阵 $\bm{\Sigma}$，其特征值为 $\lambda_1\ge\lambda_2\ge\lambda_3$、对应特征向量 $\bm{u}_1,\bm{u}_2,\bm{u}_3$。若最小特征值 $\lambda_3$ 足够小（相对 $\lambda_1$，或小于一个绝对阈值），说明这些点近似共面，就把该体素建模为一个平面 $\pi=(\bm{n},d)$，其中 $\bm{n}=\bm{u}_3$（最小特征值对应的法向）、$d=\bm{n}^{\top}\bar{\bm{q}}$；否则把体素八等分，递归处理八个子体素，直到每个体素"足够平"或者小到不再细分。这就是\textbf{自适应体素}：平面区域用大块，复杂区域用小块。

每个平面不只存参数，还存它的不确定性：平面拟合的协方差可由点分布协方差近似（$\bm{n}$ 的不确定性大致与 $\lambda_3/\lambda_2$ 相关），也可用递归最小二乘在信息形式下增量更新。这样平面参数随观测增多越来越准，残差的噪声方差可直接取平面参数的方差——比 iKD-Tree 里"所有点用同一个常数 $\sigma$"更精细。

\begin{table}[htbp]
  \centering
  \caption{iKD-Tree 与 VoxelMap 的地图表示对比}\label{tab:17-map}
  \small
  \begin{tabular}{lll}
    \toprule
    维度 & iKD-Tree（点地图） & VoxelMap（平面地图） \\
    \midrule
    存储单元 & 点 & 体素/平面（含协方差） \\
    查询 & $k$ 近邻 + 现场 PCA，$O(\log n)$ & 哈希到体素，直接取平面，近似 $O(1)$ \\
    平面拟合 & 每次查询现算 & 增量维护，可复用 \\
    精度 & 高（点云原始信息全保留） & 平面区域高，复杂区域退化 \\
    内存 & 与点数成正比，需滑窗删除 & 与体素数成正比，可预测 \\
    退化场景 & 依赖法向估计 & 平面本身即法向，更稳 \\
    非结构化场景 & 好 & 大量细分，优势下降 \\
    实现复杂度 & 中（树操作） & 中高（八叉树 + 概率更新） \\
    \bottomrule
  \end{tabular}
\end{table}

选型上可以这样记：\textbf{场景以平面为主、内存紧张、想要更快}，选 VoxelMap 类；\textbf{场景复杂、点云信息宝贵、想要更通用}，选 iKD-Tree 类。二者的残差与雅可比形式完全一样（都是点到面），差别只在地图这一层。

\section{视觉-激光-惯性：FAST-LIVO2 一句话}\label{sec:17-livo}

FAST-LIVO2~\cite{yuan2024fastlivo2} 沿着同一个 IEKF 框架把相机接了进来：状态不变，观测从"点到面残差"扩充为"点到面残差 + 光度误差"。光度误差为 $r_j^{\mathrm{vis}}=\bm{I}_j(\bm{\pi}(\bm{R},\bm{p},\bm{\rho}_j))-\bar{I}_j$，即把地图点按当前位姿投影到图像上取灰度、与参考帧同名点相减，$\bm{\pi}(\cdot)$ 是针孔投影。两类残差在同一个 $\bm{H}$ 里堆叠、用同一个 $\bm{P}$ 融合，视觉便在激光退化的方向上补充约束。适用条件是环境有纹理、光照稳定；纯黑、纯白、强反光的场地仍只能靠激光和 IMU。

\section{论文公式与源码对照}\label{sec:17-code}

\subsection{对照表}

表~\ref{tab:17-code} 把本章的公式与 FAST-LIO2 官方仓库（以及同源的 Point-LIO）中常见的文件/函数对应起来。不同版本的仓库可能有重命名，但"状态结构体 $\to$ 传播函数 $\to$ 更新函数 $\to$ 残差函数"这条主线是不变的。

\begin{table}[htbp]
  \centering
  \caption{论文公式与源码对照（以官方仓库某一版本为例，函数名可能随版本微调）}\label{tab:17-code}
  \scriptsize
  \begin{tabular}{llll}
    \toprule
    概念/公式 & 本章公式 & 文件 & 典型函数/变量 \\
    \midrule
    状态定义 & 式~\eqref{eq:17-state} & \texttt{use-ikfom.hpp} & \texttt{struct state\_ikfom} \{pos, rot, vel, bg, ba, grav\} \\
    流形加减 & 式~\eqref{eq:17-plus-minus} & \texttt{use-ikfom.hpp} & \texttt{operator+}、\texttt{operator-} \\
    IMU 传播 & 式~\eqref{eq:17-discrete} & \texttt{IMU\_Processing.hpp} & \texttt{ProcessModel()} \\
    去畸变 & 见第~\ref{ch:16-lio-evolution}~章 & \texttt{IMU\_Processing.hpp} & \texttt{UndistortPcl()} \\
    协方差传播 & 式~\eqref{eq:17-cov-prop} & \texttt{esekfom.hpp}（IKFoM） & \texttt{predict()}、\texttt{predict\_dyn()} \\
    IEKF 迭代更新 & 式~\eqref{eq:17-iekf-update} & 同上 & \texttt{update\_iterated\_dyn\_share\_modified()} \\
    点面残差与雅可比 & 式~\eqref{eq:17-residual},\eqref{eq:17-Hi} & \texttt{laserMapping.cpp} & \texttt{h\_share\_model()} \\
    主循环与回调 & 伪代码骨架 & \texttt{laserMapping.cpp} & \texttt{main()}、\texttt{imu\_cbk()} \\
    地图查询/插入/删除 & \S\ref{sec:17-ikdtree} & \texttt{ikd\_Tree.cpp} & \texttt{Nearest\_Search()} 等三接口 \\
    逐点更新 & 式~\eqref{eq:17-scalar-k} & Point-LIO 仓库 & 逐点更新函数（与 FAST-LIO 同源） \\
    \bottomrule
  \end{tabular}
\end{table}

\subsection{一段关键代码的骨架}

下面这段代码是\texttt{h\_share\_model()} 的骨架（为可读性做了简化，变量名沿用官方风格；真正的实现还要处理外参、盲区、点的时间戳和并行化）：

\begin{lstlisting}[language=C++, caption={点到面残差与雅可比（示意图，非逐字源码）}, label={lst:17-hshare}]
// schematic skeleton, names follow the official FAST-LIO2 layout
for (int i = 0; i < feats_down_size; ++i) {
    Vector3d p_imu = offset_R_L_I * p_lidar + offset_T_L_I;   // to IMU frame
    Vector3d p_world = rot * p_imu + pos;                     // to world frame
    ikdtree.Nearest_Search(p_world, NUM_MATCH_POINTS,
                           nearest_points, sq_dist);          // 5-NN
    if (nearest_points.size() < NUM_MATCH_POINTS) continue;
    plane = fit_plane(nearest_points);                        // PCA normal
    if (!plane.is_planar) continue;                           // reject
    residual[i] = plane.normal.dot(p_world - plane.centroid);
    H.block<1, 3>(i, 0) = -plane.normal.transpose() * rot * SKEW(p_imu);
    H.block<1, 3>(i, 3) =  plane.normal.transpose();
}
// H is stacked as (m x 23); only pose (and extrinsics) blocks are non-zero
\end{lstlisting}

这段代码与式~\eqref{eq:17-Hi} 一一对应：\texttt{p\_imu} 是式~\eqref{eq:17-extrinsic} 的 $\bm{p}_i^{I}$，\texttt{p\_world} 是式~\eqref{eq:17-world} 的 $\bm{p}_i^{W}$，\texttt{residual[i]} 是式~\eqref{eq:17-residual} 的 $r_i$，\texttt{H.block(0)} 与 \texttt{H.block(3)} 分别是表~\ref{tab:17-jacobian} 的前两块。注意 \texttt{SKEW(p\_imu)} 就是 $[\bm{p}_i^{I}]_\times$，负号来自 $[\delta\bm{\theta}]_\times\bm{p}=-\bm{p}\times\delta\bm{\theta}$。读代码时按"状态 $\to$ 传播 $\to$ 更新 $\to$ 残差 $\to$ 地图"的顺序走，正好对应表~\ref{tab:17-code} 的行序。

\section{工程实践：参数、发散与开销}\label{sec:17-tuning}

最常调的几个参数：\texttt{filter\_size\_surf} 与 \texttt{filter\_size\_map}（当前帧/入图前的降采样体素，典型 $0.5$\,m，调大则点少更快但精度降）、\texttt{cube\_side\_length}（局部地图边长，典型 $100\sim200$\,m，调大保留远处点但内存涨）、\texttt{point\_filter\_num}（每 $N$ 点取 $1$，典型 $3$）、\texttt{maximum\_iter}（IEKF 迭代次数，典型 $4$，调大更准更慢）、\texttt{extrinsic\_T/R}（标定值，差 $1^{\circ}$ 在 $10$\,m 处偏 $17$\,cm）。

\textbf{常见发散原因}（按现场出现频率排序）：\textbf{时间同步}——激光与 IMU 时间戳差 $10$\,ms，$2$\,m/s 下就是 $2$\,cm 错位，这是"代码没问题但一跑就飘"的头号原因；\textbf{外参错误}——标定不准或把 LiDAR 系/IMU 系搞反，表现为"静止正常、一转就飘"；\textbf{IMU 饱和}——撞击瞬间角速度超量程，积分立刻跑偏；\textbf{初始化时机器人在动}——首帧需要近似静止来对齐重力、估零偏；\textbf{场景退化}——长走廊、纯平面、玻璃幕墙，激光约束不足；\textbf{点太少或太远}；\textbf{迭代次数不够}——被撞后初值差，$4$ 次迭代可能没收敛，表现为一两帧内位姿跳变。

\textbf{开销量级与调试}（x86 单核、$10$\,Hz、每帧 $10^5$ 原始点降采样到 $5\times10^3\sim2\times10^4$）：$k$ 近邻查询约占一半时间、残差与矩阵运算占另一半；整机 CPU 占用通常在单核 $20\%\sim40\%$，内存随局部地图规模在百 MB 量级。调参时先看 odom 是否连续（有无跳变），再看残差均值是否稳定在厘米级、协方差 trace 是否收敛；把 iKD-Tree 点数打出来看是否持续增长（增长说明滑窗删除没生效）；最后用 rosbag 回放对照真值。第~\ref{ch:27-tuning}~章给出更系统的排错流程。

\begin{summary}

\begin{enumerate}
  \item FAST-LIO2 的骨架是"IMU 高频预测 + 激光帧更新"：预测用式~\eqref{eq:17-discrete}，协方差用式~\eqref{eq:17-cov-prop}，更新用式~\eqref{eq:17-iekf-update}；状态 $\bm{x}=[\bm{R},\bm{p},\bm{v},\bm{b}_g,\bm{b}_a,\bm{g}]$ 中姿态在 $\SO(3)$、重力在 $S^2$ 上，误差状态用 $\boxplus/\boxminus$ 在切空间运算（官方实现含外参，共 $23$ 维）。
  \item $\bm{F}$ 的非零块只有六个：$-[\hat{\bm{\omega}}]_\times$、$-\bm{I}$（零偏进姿态）、$\bm{I}$（速度进位置）、$-\hat{\bm{R}}[\tilde{\bm{a}}-\hat{\bm{b}}_a]_\times$、$-\hat{\bm{R}}$（零偏进速度）、$\bm{I}$（重力进速度）；位置块不注入噪声，位置误差靠速度误差积分增长。
  \item IEKF 是"对 MAP 问题做高斯-牛顿"：$\delta\bm{x}_j=\bm{K}_j(\bm{z}-\bm{h}(\bm{x}_j))-(\bm{I}-\bm{K}_j\bm{H}_j)(\bm{x}_j\boxminus\bar{\bm{x}})$，每次迭代重新线性化、重新关联，而 $\bar{\bm{P}}$ 保持不变。
  \item 点到面残差 $r_i=\bm{n}_i^{\top}(\bm{R}\bm{p}_i^{I}+\bm{p}-\bar{\bm{q}}_i)$ 对右扰动的雅可比只有两块：$-\bm{n}_i^{\top}\bm{R}[\bm{p}_i^{I}]_\times$ 与 $\bm{n}_i^{\top}$；零偏、重力、速度的块为零，它们靠先验互协方差被间接修正。
  \item iKD-Tree 用"增量插入 + 惰性删除 + box-wise 操作 + 局部重建"把地图维护降到每帧 $O(m\log n)$；Point-LIO 把更新粒度从"帧"降到"点"，每点一次标量卡尔曼更新，IMU 饱和时改用随机过程模型、放大 $\bm{Q}$，稳态误差由式~\eqref{eq:17-riccati} 给出。
\end{enumerate}


\end{summary}

\begin{problemset}

  \item 从 $\dot{\bm{R}}=\bm{R}[\bm{\omega}]_\times$ 出发，推导式~\eqref{eq:17-discrete} 中 $\bm{v}_{k+1}$ 的表达式，并写出被忽略的二阶项（提示：把 $\exp([\bm{\omega}s]_\times)$ 展开到 $s^2$ 后积分）。
  \item 若把重力从状态里去掉（当作常量 $\bm{g}=[0,0,-9.81]^{\top}$），误差状态维数变成多少？初始姿态差 $2^{\circ}$ 会导致什么后果？请用式~\eqref{eq:17-dv} 解释。
  \item 用左扰动 $\bm{R}=\exp([\delta\bm{\theta}]_\times)\hat{\bm{R}}$ 重新推导式~\eqref{eq:17-Hi}，验证雅可比为 $-\bm{n}_i^{\top}[\bm{R}\bm{p}_i^{I}]_\times$，并说明它与右扰动版本的关系。
  \item 证明线性高斯下顺序标量更新与批量更新的后验相同（提示：信息矩阵可加），再推导式~\eqref{eq:17-riccati} 并计算 $\sigma=0.02$\,m、$q=10^{-3}\,\mathrm{m^2}$ 时的稳态标准差。
  \item 一帧点云只有一面墙：写出 $\bm{H}^{\top}\bm{H}$ 的平移块、秩与零空间方向，说明此时为何要降低更新增益；再解释 box-wise 删除为何是 $O(\text{边界节点数})$。
\end{problemset}

\begin{reading}

\begin{itemize}
  \item FAST-LIO2 论文~\cite{xu2022fastlio2} 与官方仓库：按表~\ref{tab:17-code} 的顺序读代码，核对论文附录 $\bm{F}$ 与式~\eqref{eq:17-F}，前置读 ESKF 教程~\cite{sola2017eskf}。
  \item Point-LIO~\cite{he2023pointlio}、VoxelMap~\cite{yuan2022voxelmap}、FAST-LIVO2~\cite{yuan2024fastlivo2}：逐点更新、体素地图、视觉残差三支。
\end{itemize}



\end{reading}


\part{导航篇：从地图到运动}

\chapter{导航系统总览与 Nav2}\label{ch:18-nav-overview}

\usetikzlibrary{arrows.meta,positioning,calc,fit,shapes.geometric,patterns}

\chapterintro{从这一章开始，我们进入全书第四部分：导航。里程计篇回答“机器人现在在哪”，SLAM 篇回答“地图长什么样”，而导航篇要回答一个更接近任务的问题——“知道自己在哪、也知道地图之后，怎么把机器人安全、平稳地送到目标点”。本章先建立导航系统的全局骨架：导航由定位、建图、规划、控制四块组成；然后正面解释一个几乎每个新队员都会问、但网上说法互相矛盾的概念——“导航的前端/后端”；最后以 ROS 2 的 Nav2 为实例，讲清行为树、代价地图、TF 树、全局/局部规划的分工。学完本章，你应该能画出 Nav2 的节点数据流图，并说出每个节点负责哪一件事、以什么频率运行、失败时由谁兜底。}

\section{导航到底在做什么}\label{sec:18-what}

用一句话概括：\textbf{导航是在已知或部分已知的地图上，依据自身的位姿估计，生成一条满足机器人运动约束、并且不撞障碍的可行轨迹，并把轨迹变成轮子/电机的指令。}把它写成优化问题时，输入是地图 $M$、当前位姿 $\bm{x}_0$、目标位姿 $\bm{x}_g$、机器人模型（尺寸、最大速度、最大加速度、最小转弯半径），输出是一段控制序列 $\bm{u}_{0:T}$ 或一段轨迹 $\bm{x}_{0:T}$，约束是\emph{不碰撞}、\emph{运动学可行}、\emph{动力学可行}，目标通常是最短、最平滑或最省时间。

\begin{figure}[htbp]
  \centering
  \begin{tikzpicture}[node distance=7mm and 8mm, font=\small]
    \node[block] (sensor) {传感器\\{\footnotesize 点云/图像/IMU/轮速}};
    \node[block, right=of sensor] (loc) {定位与建图\\{\footnotesize 第 10--17 章}};
    \node[block, right=of loc] (plan) {规划\\{\footnotesize 第 19--21 章}};
    \node[block, right=of plan] (ctrl) {控制\\{\footnotesize 第 20、22 章}};
    \node[block, below=10mm of ctrl] (phys) {真实世界\\{\footnotesize 机器人 + 环境}};
    \draw[arrow] (sensor) -- (loc);
    \draw[arrow] (loc) -- (plan);
    \draw[arrow] (plan) -- (ctrl);
    \draw[arrow] (ctrl) -- (phys);
    \draw[arrow] (phys.west) -| (sensor.south);
    \draw[arrow, dashed] (loc.south) |- ++(0,-4mm) -| (ctrl.south);
  \end{tikzpicture}
  \caption{导航系统的闭环。上半行是“信息流”：传感器数据经定位与建图变成位姿和地图，再经规划变成轨迹，经控制变成指令；下半行是“物理流”：指令驱动真实世界，真实世界的变化又被传感器观测到。虚线表示规划与控制都要使用定位结果（位姿）。}
  \label{fig:18-loop}
\end{figure}

图~\ref{fig:18-loop} 里有两点容易忽略，却决定了工程上的成败。第一，\textbf{闭环里每一环都有延迟}，位姿估计可能比真实状态晚几十毫秒，控制指令下发又要经过通信与执行器，所以“用当前位姿算出的避障动作”实际上是在处理几十毫秒前的世界；速度越快、延迟越大，越容易撞。第二，\textbf{每一环的频率差一到两个数量级}，定位 10--50 Hz，全局规划 1 Hz 量级，局部规划 10--20 Hz，电机控制 100 Hz 以上。频率不匹配就必须靠“缓存最近的位姿与地图”和“降频重规划”来衔接，这正是 Nav2 用生命周期节点、代价地图和速度平滑器要处理的事情。

\begin{table}[htbp]
  \centering
  \caption{导航系统的四个环节及其在本书中的位置}
  \label{tab:18-pipeline}
  \small
  \begin{tabular}{llll}
    \toprule
    环节 & 回答的问题 & 典型频率 & 对应章节 \\
    \midrule
    定位与建图 & 我在哪？地图什么样？ & 10--50 Hz & 第 10--17 章 \\
    全局规划 & 大致该往哪走？ & 1 Hz & 第 19、21 章 \\
    局部规划 & 接下来 1--3 秒怎么走？ & 10--20 Hz & 第 20、21 章 \\
    控制与执行 & 左右轮各转多快？ & 50--200 Hz & 第 20、22 章 \\
    \bottomrule
  \end{tabular}
\end{table}

\section{“前端/后端”到底指什么：三种语境}\label{sec:18-fb}

这是本章最需要说清楚的一节。新队员常看到三句互相打架的话：“SLAM 的后端是图优化”“导航的后端是规划与控制”“定位系统的后端负责全局修正”。三句话都不算错，但它们在说三个不同的东西。之所以混乱，是因为“前端/后端”原先是软件工程里的分层说法（\textbf{前端}负责与外界交互、做即时处理，\textbf{后端}负责汇总、优化、做全局决策），后来被三个不同领域各自借用，含义就分叉了。

\subsection{语境一：SLAM 的前端/后端（第 11--13 章）}

在 SLAM 里，分界线非常明确：\textbf{前端}是帧间/扫描间的数据关联与相对位姿估计（特征匹配、ICP、直接法），它只关心相邻几帧，结果快但有累积漂移；\textbf{后端}是把这些相对约束与回环约束放在一起做全局优化（位姿图、因子图、BA），它关心整段轨迹的一致性，慢但能把漂移“摊平”。判断标准是：\emph{是否对全局变量做联合优化}。前端 10--30 Hz，后端 1 Hz 或只在回环时触发。第~\ref{ch:11-slam-framework}~章给出的定义是本书的基准，后面两节都要与它对照。

\subsection{语境二：定位系统的前端/后端（第 9、10、23 章）}

在定位（尤其是多传感器融合定位）里，\textbf{前端}指“用高频、局部的观测做状态递推与更新”，例如 IMU 预积分 + 激光/视觉里程计的 ESKF 滤波；\textbf{后端}指“用低频、全局的观测做长期修正”，例如 AMCL 粒子滤波重定位、UWB/二维码/UWB 基站给的绝对位置、回环检测给出的位姿修正。这里的分界不再是“有没有联合优化”，而是\emph{观测是局部的还是全局的、修正的是短期还是长期漂移}。Z-Localization 那类系统把“前端滤波 + 后端全局约束检验”拆得更细，第~\ref{ch:23-zlocalization}~章会展开。

\subsection{语境三：导航系统的前端/后端（本章、第 19--22 章）}

在导航系统里，最常见的说法是：\textbf{前端 = 定位与地图}（回答“我在哪、哪里能走”），\textbf{后端 = 规划与控制}（回答“往哪走、怎么走”）。这个说法流行，是因为它对应了数据流的方向：感知与状态估计在前，决策与控制在后，和 Web 开发里“前端展示、后端计算”的直觉很像，讲课时容易理解。

但必须指出它\textbf{不严谨}的三点：

\begin{enumerate}
  \item 导航的“后端”里并没有 SLAM 后端那种\emph{全局联合优化}。规划是搜索/采样/优化一个轨迹，控制是求解一个跟踪问题，两者都不优化“整段运行历史”。把规划叫后端，只是借了“做决策”的意思。
  \item 导航自己内部还有一层“全局/局部”的层次（全局规划 1 Hz、局部规划 10--20 Hz），这层“全局/局部”与“前端/后端”\emph{不是一回事}。很多混淆就来自把“全局规划”误说成“导航后端”。
  \item 被划进“导航前端”的定位模块内部，\emph{自己又有}前端/后端（语境二）。所以“导航的前端”内部还套着一个“定位的后端”，逐层嵌套时字面意思会打架。
\end{enumerate}

\begin{table}[htbp]
  \centering
  \caption{三种“前端/后端”语境对照（建议背下这张表）}
  \label{tab:18-fb}
  \small
  \begin{tabular}{p{1.6cm}p{3.6cm}p{3.9cm}p{3.5cm}}
    \toprule
    语境 & 前端 & 后端 & 分界依据 \\
    \midrule
    SLAM & 帧间配准、数据关联 & 位姿图/因子图全局优化 & 是否联合优化全局变量 \\
    定位系统 & 高频滤波、局部观测递推 & 全局观测重定位、长期漂移修正 & 观测是局部还是全局 \\
    导航系统 & 定位与地图 & 规划与控制（口语说法） & 数据流先后（非严格定义） \\
    \bottomrule
  \end{tabular}
\end{table}

\begin{keybox}{一句话判断法}
听到“后端”时，先看它后面接什么：接“图优化/回环”$\rightarrow$ SLAM 后端；接“全局定位/绝对观测修正”$\rightarrow$ 定位后端；接“规划/控制/行为树”$\rightarrow$ 导航后端（口语）。本书此后一律使用“SLAM 前后端”“定位前后端”这种带限定词的说法，避免歧义。
\end{keybox}

\section{Nav2：ROS 2 的导航事实标准}\label{sec:18-nav2}

Nav2 是 ROS 2 的导航框架，前身是 ROS 1 的 \texttt{move\_base}，其设计在 \cite{macenski2020nav2} 中系统给出，整体软件架构与 ROS 2 的设计哲学见 \cite{macenski2022ros2}。它不是一个“程序”，而是一组\textbf{可组合、可替换、由行为树调度}的节点。

\subsection{节点分工}

表~\ref{tab:18-nav2nodes} 是必须记住的节点清单。核心思路是：\texttt{bt\_navigator} 不自己算路径，它只负责“按行为树的剧本调用别人”；真正算路的是 \texttt{planner\_server}（全局）与 \texttt{controller\_server}（局部）；\texttt{behavior\_server} 提供“原地旋转、后退、等待”等恢复动作；\texttt{lifecycle\_manager} 负责按依赖顺序把大家从“未配置”一路带到“激活”。

\begin{table}[htbp]
  \centering
  \caption{Nav2 主要节点及其职责}
  \label{tab:18-nav2nodes}
  \small
  \begin{tabular}{lll}
    \toprule
    节点 & 职责 & 关键插件/接口 \\
    \midrule
    \texttt{bt\_navigator} & 行为树调度，对外提供导航 action & BehaviorTree.CPP \\
    \texttt{planner\_server} & 全局路径规划 & NavFn、Smac、Smac Hybrid \\
    \texttt{controller\_server} & 局部轨迹跟踪与避障 & DWB、TEB、MPPI、RPP \\
    \texttt{smoother\_server} & 路径平滑（减少拐点） & Simple、Savitzky-Golay \\
    \texttt{behavior\_server} & 恢复行为：旋转、后退、等待 & spin、backup、wait \\
    \texttt{waypoint\_follower} & 多点巡航 & NavigateThroughPoses \\
    \texttt{velocity\_smoother} & 速度限幅与加速度平滑 & 一阶低通 \\
    \texttt{map\_server / amcl} & 加载地图 / 粒子滤波定位 & 第 10、14 章 \\
    \texttt{lifecycle\_manager} & 生命周期编排 & 按序 configure/activate \\
    \bottomrule
  \end{tabular}
\end{table}

\begin{figure}[htbp]
  \centering
  \begin{tikzpicture}[node distance=7mm and 10mm, font=\small]
    \node[block] (client) {任务层/RViz2\\{\footnotesize NavigateToPose}};
    \node[block2, below=of client] (bt) {\texttt{bt\_navigator}\\{\footnotesize 行为树}};
    \node[block, below left=of bt] (pl) {\texttt{planner\_server}\\{\footnotesize 全局规划}};
    \node[block, below right=of bt] (cr) {\texttt{controller\_server}\\{\footnotesize 局部控制}};
    \node[block, left=of pl] (ms) {\texttt{map\_server}\\{\footnotesize + AMCL}};
    \node[block, right=of cr] (vs) {\texttt{velocity\_smoother}};
    \node[block, below=of vs] (base) {底盘驱动\\{\footnotesize /cmd\_vel}};
    \draw[arrow] (client) -- node[label, right]{action} (bt);
    \draw[arrow] (bt) -- node[label, above left]{\footnotesize ComputePathToPose} (pl);
    \draw[arrow] (bt) -- node[label, above right]{\footnotesize FollowPath} (cr);
    \draw[arrow, dashed] (ms) -- (pl);
    \draw[arrow, dashed] (ms.south) |- (cr.west);
    \draw[arrow] (cr) -- (vs);
    \draw[arrow] (vs) -- (base);
  \end{tikzpicture}
  \caption{Nav2 的节点数据流（简化）。箭头方向是“谁调用谁/数据流向谁”。\texttt{bt\_navigator} 通过 action 调用规划与控制服务；两者都以 \texttt{map\_server} 的静态地图和 AMCL 的位姿为输入；\texttt{controller\_server} 输出的速度指令经速度平滑器限幅后下发底盘。行为树的好处是：换一套 XML 脚本就能改变“先规划、失败则旋转、再失败则后退”的策略，而不用改 C++ 代码。}
  \label{fig:18-nav2}
\end{figure}

\subsection{action 接口：为什么不是 topic}

导航是“有开始、有结束、有反馈、可能失败、可取消”的长任务，正适合 ROS 2 的 \textbf{action}（动作）接口。ROS 2 的 action 由三部分组成：goal（目标）、feedback（周期反馈）、result（最终结果）；底层复用 topic 与 service，因此天然支持“发起后先返回、执行中持续反馈、随时取消”。Nav2 最常用的 action 是 \texttt{nav2\_msgs/action/NavigateToPose}，可以这样从命令行发一个目标：

\begin{lstlisting}[language=bash]
# Send a navigation goal from the command line (Nav2 must be running)
ros2 action send_goal /navigate_to_pose nav2_msgs/action/NavigateToPose \
  "{pose: {header: {frame_id: map}, pose: {position: {x: 1.0, y: 2.0},
   orientation: {w: 1.0}}}}" --feedback
\end{lstlisting}

反馈里会周期性给出当前位姿、已走时间、剩余距离、已执行的恢复次数。工程上只要盯住“剩余距离长时间不下降 + 恢复次数持续增加”，基本就能判定机器人卡住了——这比看日志快得多。

\subsection{与 ROS 1 \texttt{move\_base} 的差异}

\begin{table}[htbp]
  \centering
  \caption{\texttt{move\_base} 与 Nav2 的对比}
  \label{tab:18-movebase}
  \small
  \begin{tabular}{p{2.6cm}p{4.6cm}p{5.4cm}}
    \toprule
    维度 & ROS 1 \texttt{move\_base} & Nav2 \\
    \midrule
    任务编排 & 固定流程 + 硬编码恢复行为 & 行为树 XML，可在线切换策略 \\
    生命周期 & 无统一管理，启动顺序靠脚本 & 生命周期节点 + \texttt{lifecycle\_manager} \\
    插件化 & 规划器/控制器插件，扩展有限 & 规划、控制、平滑、行为全部插件化 \\
    接口 & \texttt{MoveBaseAction} & \texttt{NavigateToPose} 等一组 action \\
    多机器人 & 需命名空间硬隔离 & 天然支持命名空间与多实例 \\
    \bottomrule
  \end{tabular}
\end{table}

\section{代价地图：导航的“地面真相”}\label{sec:18-costmap}

\begin{definition}[代价地图]
代价地图（costmap）是叠加在二维平面上的栅格，每个格子存一个 8 位整数 $c\in\{0,1,\dots,255\}$，表示“机器人中心经过该格的风险”。约定 \textbf{0 表示自由}、\textbf{254 表示致命障碍}（\texttt{LETHAL\_OBSTACLE}）、\textbf{253 表示内切膨胀障碍}（\texttt{INSCRIBED\_INFLATED\_OBSTACLE}，即“机器人中心到这里必然碰撞”）、\textbf{255 表示未知}（\texttt{NO\_INFORMATION}）。中间值表示“离障碍越近代价越高”。
\end{definition}

举个例子：$10\ \mathrm{m}\times 10\ \mathrm{m}$ 的场地、分辨率 $0.05\ \mathrm{m}$，则每边 $10/0.05=200$ 格，共 $4\times10^4$ 格；每格 1 字节，单层约 39 KB，四层也不过 160 KB 左右，完全可以每帧全量更新。这就是为什么 2D 代价地图能在 5--20 Hz 下实时刷新。

\subsection{分层结构}

Nav2 的 2D 代价地图由若干 \textbf{layer}（层）叠加而成，自下而上依次是：

\begin{itemize}
  \item \textbf{静态层}（static layer）：读 \texttt{map\_server} 的 pgm 地图，给出已知障碍与未知区域，基本不变。
  \item \textbf{障碍层/体素层}（obstacle layer / voxel layer）：把实时传感器（激光、深度相机、3D 点云）观测到的障碍投影到 2D 栅格。体素层会把 3D 点云压成带高度过滤的 2D 障碍，是 3D 传感器进入 2D 导航的桥梁。
  \item \textbf{膨胀层}（inflation layer）：把障碍向外“扩散”出一圈代价梯度，让规划器主动远离墙壁。它不新增障碍，只改变代价值。
\end{itemize}

层的叠加规则是“取最大”：$c(x,y)=\max_k c_k(x,y)$。因此静态层里的致命障碍不会被障碍层覆盖成自由，而障碍层新看到的障碍会立即压过静态层的自由空间。

\subsection{内切半径、外切半径与膨胀}

设机器人俯视轮廓的最小外接圆半径为\textbf{外切半径} $r_{\mathrm{out}}$，机器人能塞进的最窄通道半宽为\textbf{内切半径} $r_{\mathrm{in}}$。对 $0.6\ \mathrm{m}\times 0.6\ \mathrm{m}$ 的方形底盘，$r_{\mathrm{in}}=0.3\ \mathrm{m}$，$r_{\mathrm{out}}=\sqrt{0.3^2+0.3^2}=0.424\ \mathrm{m}$。

膨胀层的代价值由下式给出：
\begin{equation}
  c(d)=
  \begin{cases}
    254, & d<r_{\mathrm{in}} \text{（中心在障碍上）},\\[2pt]
    253, & r_{\mathrm{in}}\le d \le r_{\mathrm{in}}+\varepsilon,\\[2pt]
    \left\lfloor 253\,\exp\!\left[-k\,(d-r_{\mathrm{in}})\right]\right\rfloor, & r_{\mathrm{in}}<d<r_{\mathrm{infl}},\\[2pt]
    0, & d\ge r_{\mathrm{infl}},
  \end{cases}
  \label{eq:18-inflation}
\end{equation}
其中 $d$ 是该格到最近障碍的欧氏距离，$r_{\mathrm{infl}}$ 是 \texttt{inflation\_radius}（膨胀半径），$k$ 是 \texttt{cost\_scaling\_factor}（代价衰减系数）。$d$ 用距离变换在整张图上一次算出，所以膨胀本身是 $O(N)$ 的。

取 $r_{\mathrm{in}}=0.3$、$k=3.0$、$r_{\mathrm{infl}}=0.55$（单位 m），算三个数：
\begin{itemize}
  \item $d=0.35$：$c=253e^{-3\times0.05}=253\times0.861=217.8\to217$，很高，规划器会避开；
  \item $d=0.45$：$c=253e^{-3\times0.15}=253\times0.638=161.4\to161$，仍然偏高；
  \item $d=0.55$：$c=253e^{-3\times0.25}=253\times0.472=119.5\to119$，但按上界已截断为 0。
\end{itemize}
注意最后一行的截断：只要 $d\ge r_{\mathrm{infl}}$ 就强制归零，否则指数衰减永远到不了 0，整张图都会“有代价”。

\begin{keybox}{实例：机器人过窄门}
门宽 $0.9\ \mathrm{m}$，方形底盘 $0.6\ \mathrm{m}$。机器人中心必须离两侧门框各 $\ge r_{\mathrm{in}}=0.3\ \mathrm{m}$，可行带宽度只剩 $0.9-2\times0.3=0.3\ \mathrm{m}$，按 $0.05\ \mathrm{m}$ 分辨率约 6 格。若把 \texttt{inflation\_radius} 设成 $0.8$、$k=5$，门中心 $d=0.45$ 处代价为 $253e^{-5\times0.15}=253\times0.472=119$，两侧墙的膨胀区在门中央几乎重叠，规划器会认为“门口很贵”而绕路甚至报“无可行路径”。正确做法是：$r_{\mathrm{in}}$ 按真实车体设准（别虚大），\texttt{inflation\_radius} 取略大于 $r_{\mathrm{out}}$（如 $0.45$--$0.5$），$k$ 取 2--3；或给规划器设置“允许穿越高代价区”的阈值。
\end{keybox}

\subsection{更新频率与发布频率}

\texttt{update\_frequency} 是代价地图\emph{重新计算}（融合新观测、重算膨胀）的频率，\texttt{publish\_frequency} 是把它\emph{发布}给其他节点/RViz 的频率。前者决定避障反应速度，后者只影响可视化与订阅端。常见配置是 update 5 Hz、publish 2 Hz；全局代价地图可以更低（1 Hz），因为它只服务 1 Hz 的全局规划。

\section{TF 树与坐标系}\label{sec:18-tf}

\begin{definition}[TF 树]
TF 树是以坐标系为节点、以带时间戳的刚体变换为边的树。任意两个坐标系 $A$、$C$ 之间的变换，通过沿树从 $A$ 走到 $C$ 的路径逐段复合得到：若路径是 $A\to B\to C$，则 $T_{AC}=T_{AB}\,T_{BC}$。TF 库（tf2）沿路径查询并插值到指定时刻，因此每个坐标系在任意时刻都只有一个父节点——这正是“树”而不是“图”的要求。
\end{definition}

举个复合的例子。机体在世界系中位姿 $T_{WB}$：平移 $(1,2)\ \mathrm{m}$、偏航 $30^\circ$；激光雷达在机体系下 $(0.2,0)\ \mathrm{m}$、无旋转。则
\begin{equation}
  {}^{W}\bm{p}_{L}
  = {}^{W}\bm{R}_{B}\,{}^{B}\bm{p}_{L}+{}^{W}\bm{p}_{B}
  = \begin{bmatrix}\cos30^\circ & -\sin30^\circ\\ \sin30^\circ & \cos30^\circ\end{bmatrix}
    \begin{bmatrix}0.2\\0\end{bmatrix}+\begin{bmatrix}1\\2\end{bmatrix}
  = \begin{bmatrix}0.173\\0.1\end{bmatrix}+\begin{bmatrix}1\\2\end{bmatrix}
  = \begin{bmatrix}1.173\\2.1\end{bmatrix}.
\end{equation}
也就是说，雷达看到的世界坐标要先减去这个偏移才能变回机体系——外参错了，障碍物在地图上的位置就会整体偏移，表现为“代价地图里的墙和真实墙错位”。

\begin{figure}[htbp]
  \centering
  \begin{tikzpicture}[node distance=8mm and 10mm, font=\small,
      frame/.style={rectangle, rounded corners=2pt, draw=structurecolor, thick,
                    fill=structurecolor!8, minimum height=7mm, minimum width=20mm, align=center}]
    \node[frame] (map) {\texttt{map}};
    \node[frame, below=of map] (odom) {\texttt{odom}};
    \node[frame, below=of odom] (foot) {\texttt{base\_footprint}};
    \node[frame, below=of foot] (base) {\texttt{base\_link}};
    \node[frame, below left=of base] (laser) {\texttt{laser}};
    \node[frame, below right=of base] (imu) {\texttt{imu}};
    \draw[arrow] (map) -- node[label, right]{\footnotesize AMCL/SLAM：可跳变} (odom);
    \draw[arrow] (odom) -- node[label, right]{\footnotesize 轮速+IMU：连续} (foot);
    \draw[arrow] (foot) -- node[label, right]{\footnotesize 静态} (base);
    \draw[arrow] (base) -- node[label, left]{\footnotesize 静态外参} (laser);
    \draw[arrow] (base) -- node[label, right]{\footnotesize 静态外参} (imu);
  \end{tikzpicture}
  \caption{2D 导航的标准 TF 树。\texttt{map}$\to$\texttt{odom} 由定位模块发布，可以跳变（重定位时瞬间修正）；\texttt{odom}$\to$\texttt{base\_footprint} 由里程计发布，必须连续无跳变；\texttt{base\_link} 以下的传感器外参是静态的，由 URDF 经 \texttt{robot\_state\_publisher} 广播。}
  \label{fig:18-tf}
\end{figure}

图~\ref{fig:18-tf} 的树结构决定了三件事：

\begin{enumerate}
  \item \textbf{谁发布}：\texttt{map}$\to$\texttt{odom} 只能由定位模块（AMCL 或 SLAM）发布；\texttt{odom}$\to$\texttt{base\_footprint} 只能由里程计发布；传感器外参由 URDF 发布。重复发布会让 tf2 报“TF 树出现多个父节点”，整棵树失效。
  \item \textbf{为什么不直接 \texttt{map}$\to$\texttt{base\_link}}：因为定位会跳变，而局部规划、控制需要平滑连续的位姿。把跳变隔离在 \texttt{map}$\to$\texttt{odom} 这一段，局部控制始终用连续的 \texttt{odom} 系，就不会因为一次重定位而急打方向。
  \item \textbf{时间戳}：tf2 默认只缓存 10 s，查询时刻超出缓存就抛 \texttt{ExtrapolationException}。常见原因是某个节点的时间戳用了系统时间而其他节点用 ROS 时间（仿真），或者传感器驱动停了。
\end{enumerate}

\begin{lstlisting}[language=bash]
# Check whether map->odom is actually published and how it evolves
ros2 run tf2_ros tf2_echo map odom
# Visualize the whole TF tree; look for disconnected branches
ros2 run tf2_tools view_frames
\end{lstlisting}

\section{全局规划与局部规划的分工}\label{sec:18-global-local}

\begin{definition}[全局规划与局部规划]
\textbf{全局规划}（global planning）在静态或准静态的全局地图上，从起点到目标点求一条（近似）最优的几何路径，只考虑“能不能走”，不考虑“怎么拐”；\textbf{局部规划}（local planning，也叫局部轨迹规划/控制）在机器人附近几米的滚动窗口内，结合实时障碍与自身运动学、动力学约束，生成未来 1--3 秒的可执行轨迹并跟踪全局路径。
\end{definition}

两者的输入输出完全不同：全局规划输入是起点、目标、全局代价地图，输出是一串路径点（path，只有几何没有时间）；局部规划输入是这条路径、局部代价地图、当前位姿与速度，输出是带时间信息的轨迹或直接的速度指令 $(\bm{v},\omega)$。频率上差一个数量级，全局 1 Hz 足够（路径不会每秒都变），局部必须 10--20 Hz（障碍随时出现）。

评价一条全局路径看四点：\textbf{长度}（是否绕远）、\textbf{平滑度}（曲率是否连续，曲率突变会让底盘抖动）、\textbf{安全性}（离障碍的净空）、\textbf{可行性}（是否满足最小转弯半径）。评价局部规划则看：\textbf{跟踪误差}、\textbf{实时性}、\textbf{避障成功率}、\textbf{舒适性}（加速度、 jerk）。第~\ref{ch:19-global-planning}~章讲全局规划的算法，第~\ref{ch:20-local-planning}~章讲局部规划与控制，第~\ref{ch:21-3d-nav}~章把两者搬到 3D 地形上，第~\ref{ch:22-nav2-practice}~章把它们串成完整链路。

\begin{summary}

\begin{enumerate}
  \item 导航 = 定位/建图 + 全局规划 + 局部规划 + 控制，是一个带延迟、频率相差 1--2 个数量级的闭环；理解延迟与频率差是调参的前提。
  \item “前端/后端”有三种语境：SLAM 看是否联合优化全局变量；定位系统看观测是局部还是全局；导航系统口语里指“定位与地图 / 规划与控制”，后者不严谨，本书一律加限定词。
  \item Nav2 是行为树调度的插件化节点集合：\texttt{bt\_navigator} 编排，\texttt{planner\_server} 与 \texttt{controller\_server} 干活，\texttt{behavior\_server} 兜底，\texttt{lifecycle\_manager} 负责启动顺序，对外用 action 接口。
  \item 代价地图是 0--255 的分层栅格，取最大叠加；内切半径决定“必然碰撞区”，膨胀半径与衰减系数决定“多怕墙”。式~\eqref{eq:18-inflation} 是调参的量化依据。
  \item TF 树是一条链：\texttt{map}$\to$\texttt{odom} 可跳变、\texttt{odom}$\to$\texttt{base} 必须连续、传感器外参静态；重复发布与时间戳越界是最常见的两类故障。
  \item 全局规划求几何路径（1 Hz，重长度/平滑/可行），局部规划求带时间的轨迹（10--20 Hz，重实时/避障）。
\end{enumerate}


\end{summary}

\begin{problemset}

  \item 用式~\eqref{eq:18-inflation} 计算：$r_{\mathrm{in}}=0.25$、$k=2.5$、$r_{\mathrm{infl}}=0.6$ 时，$d=0.3,0.4,0.5,0.6$ 四处的代价值，并解释为什么 $d=0.6$ 必须归零。
  \item 一辆 $0.5\ \mathrm{m}\times0.4\ \mathrm{m}$ 的底盘要通过 $0.75\ \mathrm{m}$ 宽的门。请给出内切半径 $r_{\mathrm{in}}$ 与外切半径 $r_{\mathrm{out}}$，再给出一组可用的膨胀半径与代价衰减系数，并说明理由。
  \item 如果 AMCL 和另一个建图节点同时发布 \texttt{map}$\to$\texttt{odom}，会发生什么？如何用命令行确认是谁在发布？
  \item 为什么全局规划不需要 10 Hz？如果强行提高到 10 Hz，会带来什么代价？
  \item 用一句话分别回答“SLAM 的后端是什么”“定位系统的后端是什么”“导航的后端是什么”，并说明三者的分界依据。
\end{problemset}

\begin{reading}

\begin{itemize}
  \item \cite{macenski2020nav2}：Nav2 的原始设计论文，读它解决“为什么要有行为树、为什么拆成这么多 server”的问题。
  \item \cite{macenski2022ros2}：ROS 2 的设计与架构，理解 action、生命周期节点、组件化之后再读 Nav2 源码会轻松很多。
  \item \cite{thrun2005probabilistic}：第 1 章与第 5--6 章，概率机器人视角下“定位—规划—控制”的统一框架。
  \item Nav2 官方文档 \href{https://docs.nav2.org}{docs.nav2.org}：配置项、默认参数、行为树 XML 示例，工程落地时按图索骥。
\end{itemize}


\end{reading}

\chapter{全局路径规划}\label{ch:19-global-planning}

\usetikzlibrary{arrows.meta,positioning,calc,fit,patterns,decorations.pathreplacing}

\chapterintro{第~\ref{ch:18-nav-overview}~章把导航拆成“全局规划 + 局部规划”两层，本章专门讲全局规划：在一张已知的地图（栅格、路网、可通行区域）上，为机器人找一条从起点到目标点的路径。我们按“从严谨到实用”的顺序展开：先定义问题与评价标准，再讲 Dijkstra、A* 及其可采纳性证明，然后讲工业界最常用的三种加速/扩展——JPS（跳点搜索）、Hybrid A*（满足最小转弯半径）与 RRT*（高维采样），最后落到 Nav2 的规划器插件与路径平滑。学完本章，你应该能手推 A* 为什么最优、能根据底盘类型选规划器、并知道代价地图的代价值如何变成搜索的边权。}

\section{问题定义与评价标准}\label{sec:19-problem}

全局规划可以抽象成\textbf{图搜索}问题。把机器人可能处于的每个状态抽象成一个\textbf{节点}（node），把一次允许的转移抽象成一条\textbf{边}（edge），边上带非负代价 $c(n,n')>0$；给定起点 $s$ 与目标 $t$，要求一条路径 $\pi=(s=n_0,n_1,\dots,n_k=t)$，使其总代价
\begin{equation}
  C(\pi)=\sum_{i=0}^{k-1} c(n_i,n_{i+1})
  \label{eq:19-pathcost}
\end{equation}
最小。2D 栅格地图上，节点就是格子，边是 4 邻接或 8 邻接的相邻格，$c$ 通常取 1（直走）与 $\sqrt2$（斜走），再加上该格的代价值惩罚。

先用一个具体例子把符号落地：分辨率 $0.05\ \mathrm{m}$ 的栅格里，直走一格代价 1、实际长度 $0.05\ \mathrm{m}$；斜走一格代价 $\sqrt2\approx1.414$、实际长度 $0.0707\ \mathrm{m}$。若某格代价值 $c_{\mathrm{cell}}=100$（比较靠近障碍），工程上把它折算成额外代价，例如边权取 $w=1+3\cdot(100/254)\approx2.18$，于是规划器会“宁可绕一格也不贴着障碍走”。

\begin{definition}[完备性、最优性、启发函数]
\textbf{完备性}（completeness）：如果解存在，算法一定能在有限时间内找到一条解；\textbf{最优性}（optimality）：算法找到的解的代价等于最小代价 $C^*$；\textbf{启发函数}（heuristic）$h(n)$ 是对“从节点 $n$ 到目标的剩余代价”的估计，真实剩余最优代价记作 $h^*(n)$。若对所有 $n$ 都有 $0\le h(n)\le h^*(n)$，称 $h$ \textbf{可采纳}（admissible）；若还满足 $h(n)\le c(n,n')+h(n')$（三角不等式），称 $h$ \textbf{一致}（consistent，或单调）。
\end{definition}

用数字感受一下：在 $10\times10$ 栅格里从 $(1,1)$ 到 $(8,8)$，若只许四邻接直走，真实剩余代价 $h^*=7+7=14$。曼哈顿距离 $h=|8-1|+|8-1|=14$，恰好等于 $h^*$，既紧又可采纳；欧氏距离 $h=\sqrt{7^2+7^2}=9.90$，也可采纳但偏小；若误用“直线距离 $\times2$”$=19.8>h^*$，就\textbf{不可采纳}，A* 不再保证最优。

\section{Dijkstra 算法}\label{sec:19-dijkstra}

Dijkstra 是“无启发”的最短路：维护一个按 $g(n)$（从起点到 $n$ 的当前最优代价）排序的优先队列，每次弹出 $g$ 最小的节点并松弛它的邻居。

\begin{pseudo}{Dijkstra（单源最短路）}
  \item 初始化：$g(s)=0$，其余 $g(n)=+\infty$；优先队列 $Q$ 只放 $s$；closed 集为空
  \item 当 $Q$ 非空：弹出 $g$ 最小的节点 $n$，若 $n=t$ 则回溯得到路径并返回
  \item 若 $n\in$ closed 则跳过；把 $n$ 加入 closed
  \item 对 $n$ 的每个邻居 $n'$：若 $g(n)+c(n,n')<g(n')$，则更新 $g(n')$，记录父节点为 $n$，把 $n'$ 压入 $Q$
  \item 若 $Q$ 空了仍未到 $t$，说明不存在可行路径
\end{pseudo}

\textbf{正确性要点}：当 $n$ 被弹出时，$g(n)$ 已经是最优的。因为所有边权非负，任何“更短、但尚未发现的路径”都必须经过 $Q$ 中某个 $g$ 不小于 $g(n)$ 的节点，而它之后再走非负代价只会更贵——这与 $n$ 是当前最小矛盾。复杂度：用二叉堆实现时，每个节点入队/出队各 $O(\log V)$，每条边松弛一次，总计 $O((V+E)\log V)$；栅格图上 $E\approx4V$ 或 $8V$，故约 $O(V\log V)$。

$10\times10$ 的图有 $V=100$ 个节点，计算量可忽略；换成 $1000\times1000$ 的室内地图，$V=10^6$，Dijkstra 就明显偏慢，这正是需要启发函数的原因。

\section{A*：用启发函数剪枝}\label{sec:19-astar}

A* 在 Dijkstra 上只改一处：优先队列按
\begin{equation}
  f(n)=g(n)+h(n)
  \label{eq:19-f}
\end{equation}
排序。$g(n)$ 是“已经花的钱”，$h(n)$ 是“预计还要花的钱”，$f(n)$ 是“整条路径的预算估计”。因为 $h$ 把目标方向的信息带了进来，A* 会优先朝目标扩展，扩展节点数通常远少于 Dijkstra\cite{hart1968astar}。

\subsection{可采纳性：为什么 $h\le h^*$ 就保证最优}

这是本章最需要“真的推一遍”的定理。

\begin{theorem}[A* 的最优性]
若启发函数可采纳（$0\le h(n)\le h^*(n)$ 且 $h(t)=0$），边权非负，则 A* 返回的路径代价等于最优代价 $C^*$。
\end{theorem}

\textbf{证明}（反证法）。设最优路径为 $\pi^*$，代价 $C^*$。假设 A* 在弹出目标 $t$ 时返回的路径代价 $C>C^*$，即 $f(t)=g(t)+h(t)=C>C^*$。

第一步，证明“前沿里始终存在 $\pi^*$ 上的节点，且其 $g$ 已是最优”。A* 每弹出一个节点就松弛其所有邻居，因此只要 $\pi^*$ 上某节点的前驱被扩展过，该节点就会被放进 open 表。初始时 $s=n_0$ 在 open 表且 $g(s)=0=g^*(s)$。假设弹出 $t$ 之前，$\pi^*$ 上已有 $n_0,\dots,n_{k-1}$ 被扩展且 $g(n_i)=g^*(n_i)$，则扩展 $n_{k-1}$ 时松弛边 $(n_{k-1},n_k)$，得到
\begin{equation}
  g(n_k)\le g(n_{k-1})+c(n_{k-1},n_k)=g^*(n_{k-1})+c(n_{k-1},n_k)=g^*(n_k).
\end{equation}
而 $g(n_k)$ 不可能小于最优前缀代价 $g^*(n_k)$（否则 $g^*(n_k)$ 就不是最优），故 $g(n_k)=g^*(n_k)$。归纳得：在弹出 $t$ 的时刻，$\pi^*$ 上第一个尚未被扩展的节点 $n$ 必在 open 表中，且 $g(n)=g^*(n)$。

第二步，估计这个 $n$ 的 $f$ 值。由 $\pi^*$ 的最优性，$g^*(n)+h^*(n)=C^*$；又 $h$ 可采纳，$h(n)\le h^*(n)$，于是
\begin{equation}
  f(n)=g(n)+h(n)=g^*(n)+h(n)\le g^*(n)+h^*(n)=C^*.
\end{equation}

第三步，利用 A* 的“按 $f$ 最小弹出”性质。既然 $t$ 被弹出，说明它对所有 open 节点都有 $f(t)\le f(n)$，于是
\begin{equation}
  C=f(t)\le f(n)\le C^*,
\end{equation}
这与假设 $C>C^*$ 矛盾。因此 $C=C^*$，A* 返回最优解。证毕。

这个证明用了三个前提：边权非负（保证“绕远不会更便宜”）、$h(t)=0$（目标处预算就是已花费）、$h$ 可采纳（保证估计不会高过真实），缺一不可。

\subsection{一致性：为什么可以不重复扩展}

若 $h$ 还满足一致性 $h(n)\le c(n,n')+h(n')$，则沿任意一条边 $n\to n'$，
\begin{equation}
  f(n)=g(n)+h(n)\le g(n)+c(n,n')+h(n')=f(n'),
\end{equation}
即 $f$ 沿路径单调不减。注意等式右端正是 $f(n')$，因为 $g(n')=g(n)+c(n,n')$ 是沿这条边到达 $n'$ 的代价（可能还不是最优，但上界成立）。

由此可证“每个节点只扩展一次就得到最优 $g$”。设节点 $n$ 被弹出时 $g(n)>g^*(n)$。取最优路径上第一个尚未被扩展的节点 $n'$，由第一步同样的归纳可知 $g(n')=g^*(n')$。沿最优路径从 $n'$ 到 $n$ 反复使用一致性 $h(n_i)\le c(n_i,n_{i+1})+h(n_{i+1})$ 并求和，得到
\begin{equation}
  h(n')\le \bigl(g^*(n)-g^*(n')\bigr)+h(n).
\end{equation}
两边同加 $g^*(n')$：
\begin{equation}
  f(n')=g^*(n')+h(n')\le g^*(n)+h(n)<g(n)+h(n)=f(n),
\end{equation}
最后一个不等号用到了假设 $g^*(n)<g(n)$。于是 $f(n')<f(n)$，$n'$ 会被先弹出，与“$n$ 先弹出”矛盾。所以弹出时必有 $g(n)=g^*(n)$，closed 集不会误封节点，A* 无需重新打开已关闭节点。

\subsection{启发函数的设计}

设当前节点 $(x,y)$、目标 $(x_g,y_g)$，常用三种：

\begin{itemize}
  \item \textbf{曼哈顿距离} $h=|x-x_g|+|y-y_g|$：四邻接栅格的精确下界，可采纳且一致；上例 $h=14$。
  \item \textbf{欧氏距离} $h=\sqrt{(x-x_g)^2+(y-y_g)^2}$：任意方向移动都适用；上例 $h=\sqrt{98}=9.90$，比曼哈顿松，扩展更多。
  \item \textbf{对角距离}（八邻接、直走 1 斜走 $\sqrt2$）：设 $d_x=|x-x_g|$、$d_y=|y-y_g|$，则 $h=\max(d_x,d_y)+(\sqrt2-1)\min(d_x,d_y)$；上例 $h=7+0.414\times7=9.90$，数值与欧氏相同，但在八邻接图上更精确。
\end{itemize}

\textbf{启发函数越紧（越接近 $h^*$ 且不超过），扩展越少，但计算 $h$ 本身也越贵。}这是工程上的经典取舍：曼哈顿最便宜，欧氏最通用，八邻接地图上用对角距离。

\subsection{$10\times10$ 网格手算例子}

\begin{figure}[htbp]
  \centering
  \begin{tikzpicture}[scale=0.62, font=\footnotesize]
    \draw[step=1cm, black!22, thin] (0,0) grid (7,7);
    \foreach \x/\y in {0/0,1/0,0/1,1/1,2/0,2/1,2/2,2/3,2/4,4/4,5/4,4/5,4/6}
      \fill[structurecolor!12] (\x,\y) rectangle ++(1,1);
    \foreach \x/\y in {3/1,3/2,3/3}
      \fill[black!55] (\x,\y) rectangle ++(1,1);
    \draw[arrow, structurecolor, line width=1.1pt]
      (0.5,0.5) -- (0.5,2.5) -- (2.5,2.5) -- (2.5,4.5)
      -- (4.5,4.5) -- (4.5,6.5) -- (6.5,6.5);
    \node[circle, fill=green!45!black, inner sep=1.2pt] at (0.5,0.5) {};
    \node[below left, label] at (0.5,0.5) {$s$};
    \node[circle, fill=red!70!black, inner sep=1.2pt] at (6.5,6.5) {};
    \node[above right, label] at (6.5,6.5) {$t$};
    \node[black!55, font=\scriptsize] at (3.5,2.5) {障碍};
    \node[structurecolor, font=\scriptsize] at (1.5,1.5) {已扩展};
  \end{tikzpicture}
  \caption{A* 在栅格上的扩展与最终路径示意。浅色格是 $f$ 值相同、被优先扩展的“等值带”；深色格是障碍；折线是最终路径。启发函数的作用就是把扩展限制在朝目标的等值带内，而不是像 Dijkstra 那样铺满整张图。}
  \label{fig:19-astar}
\end{figure}

设栅格 $10\times10$，坐标 $(x,y)$ 取 $0,\dots,9$，四邻接直走代价 1。起点 $s=(1,1)$，目标 $t=(8,8)$。障碍是一条横墙：第 $y=2$ 行的 $x=0,\dots,8$ 全部不可通行，只在 $(9,2)$ 留一个缺口，于是机器人必须先绕到最右侧，才能上到墙的上方。曼哈顿启发 $h=|8-x|+|8-y|$，扩展过程见表~\ref{tab:19-astar}。

\begin{table}[htbp]
  \centering
  \caption{手算 A*：关键节点的 $g$、$h$、$f$}
  \label{tab:19-astar}
  \small
  \begin{tabular}{llccc}
    \toprule
    节点 $(x,y)$ & 到达方式 & $g$ & $h$ & $f=g+h$ \\
    \midrule
    $(1,1)$ & 起点 & 0 & 14 & 14 \\
    $(2,1)$ & 向右一步 & 1 & 13 & 14 \\
    $(5,1)$ & 向右四步 & 4 & 10 & 14 \\
    $(8,1)$ & 向右七步 & 7 & 7 & 14 \\
    $(9,1)$ & 再向右一步 & 8 & 8 & 16 \\
    $(9,3)$ & 经缺口上移 & 10 & 6 & 16 \\
    $(9,7)$ & 沿右侧上行 & 14 & 2 & 16 \\
    $(8,8)$ & 向左一步到达 & 16 & 0 & 16 \\
    \bottomrule
  \end{tabular}
\end{table}

读表~\ref{tab:19-astar} 要注意两点。第一，墙下方一整排节点的 $f$ 都等于 14，A* 无法区分它们，只能全部扩展，这就是图~\ref{fig:19-astar} 里的“等值带”；启发函数在陷阱地形中会变“钝”。第二，真正的路径代价是 16，而 $h(s)=14$，说明曼哈顿距离是下界（可采纳），但低估了 2 步——多出来的 2 步正是绕墙的代价。若换用欧氏距离，$h(s)=9.90$，低估更多，扩展的节点只会更多。这解释了“为什么在长走廊、凹形房间里 A* 会退化成 Dijkstra”。

\section{JPS：跳点搜索}\label{sec:19-jps}

A* 的扩展量还可以再压：栅格上有大量节点处于同一方向、同一代价的“直道”上，逐个扩展毫无意义。JPS（Jump Point Search）\cite{harabor2011jps} 的思路是\textbf{只保留必须做决策的“跳点”}（jump point），直道中间的点直接跳过。

核心概念是\textbf{强制邻居}（forced neighbor）：假设从父节点 $p$ 沿方向 $\bm{d}$ 走到当前节点 $x$。若 $x$ 的某个邻居 $n$ 不在 $\bm{d}$ 的“自然前进方向”上，但\textbf{不经过 $x$ 就无法从 $p$ 以最优代价到达 $n$}（典型情形是 $n$ 旁边紧贴障碍，$x$ 是该障碍的拐角），则 $n$ 是 $x$ 的强制邻居，$x$ 必须作为跳点保留。JPS 的规则是：沿方向 $\bm{d}$ 一直“跳”，直到（1）遇到强制邻居，（2）到达目标，（3）撞到障碍；遇到（1）（2）就把该点加入 open 表，遇到（3）则此方向作废。

\begin{pseudo}{JPS 的跳点函数 Jump($x$, $\bm{d}$)}
  \item $x' \gets x+\bm{d}$（沿方向走一步）
  \item 若 $x'$ 不可通行，返回空
  \item 若 $x'=t$，返回 $x'$
  \item 若 $x'$ 存在强制邻居，返回 $x'$
  \item 若 $\bm{d}$ 是对角方向：对两个正交分量分别递归 Jump($x'$, $\bm{d}_x$)、Jump($x'$, $\bm{d}_y$)，任一返回非空则返回 $x'$
  \item 否则返回 Jump($x'$, $\bm{d}$)
\end{pseudo}

JPS 的最坏复杂度与 A* 同阶，但实际扩展节点数常减少一个数量级（开阔场景中甚至接近直线上的点数）。它的代价是：强制邻居的判定与对角递归让代码复杂，且依赖栅格的 8 邻接规则；在代价值不均匀的代价地图上，JPS 的“对称路径假设”不成立，效果会打折扣。

\section{Hybrid A*：给车用的 A*}\label{sec:19-hybrid}

普通 A* 的路径由栅格边组成，会出现 $45^\circ$ 折线，且隐含“可以原地转向、可以任意小半径转弯”的假设。轮式底盘（尤其阿克曼式）做不到：它有最小转弯半径 $R_{\min}$，不能横移，倒车也受限制。Hybrid A*\cite{dolgov2010hybridastar} 的做法是：\textbf{搜索状态仍是连续的 $(x,y,\theta)$，但用栅格做“已访问”去重}，扩展时不再连相邻格，而是用一串符合车辆运动学的\textbf{运动基元}（motion primitive）。

最常用的基元来自自行车模型。设轴距 $L$、前轮转角 $\delta$、速度 $v$（可正可负），则
\begin{equation}
  \dot x=v\cos\theta,\qquad \dot y=v\sin\theta,\qquad \dot\theta=\frac{v\tan\delta}{L}.
  \label{eq:19-bicycle}
\end{equation}
把 $\delta$ 固定为几个离散值（如 $-\delta_{\max},0,+\delta_{\max}$），把时间离散为步长 $\Delta t$，就得到一条短弧：
\begin{equation}
  x_{k+1}=x_k+v\Delta t\cos\theta_k,\quad
  y_{k+1}=y_k+v\Delta t\sin\theta_k,\quad
  \theta_{k+1}=\theta_k+\frac{v\Delta t\tan\delta_k}{L}.
\end{equation}
例如 $L=0.5\ \mathrm{m}$、$\delta_{\max}=30\degree$、$v=1\ \mathrm{m/s}$、$\Delta t=0.2\ \mathrm{s}$，则每步前进 $0.2\ \mathrm{m}$，转角增量 $\Delta\theta=0.2\times\tan30\degree/0.5=0.231\ \mathrm{rad}=13.2\degree$，对应的最小转弯半径 $R_{\min}=L/\tan\delta_{\max}=0.5/0.577=0.866\ \mathrm{m}$。

“Hybrid”体现在两处：状态连续（$x,y,\theta$ 是实数），去重离散（把状态映射到索引 $(\lfloor x/\rho\rfloor,\lfloor y/\rho\rfloor,\lfloor\theta/\Delta\theta\rfloor)$ 记入 closed 表，$\rho$ 是栅格分辨率）。

启发函数取两个下界的最大值：
\begin{itemize}
  \item $h_{\mathrm{kin}}$：忽略障碍、只考虑运动学的最短距离（Reeds-Shepp 曲线，允许前进 + 倒车；只许前进时用 Dubins 曲线）；
  \item $h_{\mathrm{grid}}$：忽略运动学、用 2D 栅格 Dijkstra 算出的带障碍最短距离。
\end{itemize}
两者都是真实代价的下界（各自忽略了一部分约束），所以 $h=\max(h_{\mathrm{kin}},h_{\mathrm{grid}})$ 仍可采纳；取最大还能让搜索同时被“运动学可行”和“避开障碍”两个方向引导。

代价函数除了路径长度，还要加惩罚项：倒车惩罚、换挡惩罚、转角变化惩罚、贴近障碍惩罚。工程上这些权重通常设成“1 m 倒车约等于 2--3 m 前进”的量级。

\begin{figure}[htbp]
  \centering
  \begin{tikzpicture}[scale=0.95, font=\footnotesize]
    \fill[structurecolor!8] (-0.4,-0.35) rectangle (0.4,0.35);
    \draw[structurecolor, thick] (-0.4,-0.35) rectangle (0.4,0.35);
    \fill[black!70] (-0.28,-0.3) rectangle (0.28,-0.22);
    \fill[black!70] (-0.28,0.22) rectangle (0.28,0.3);
    \draw[arrow] (0.4,0) -- (0.7,0) node[right, label] {$\theta$};
    \draw[arrow, green!45!black, thick] (0.4,0.0) .. controls (1.4,0.1) and (2.0,0.9) .. (2.2,1.9)
      node[above, label] {$\delta=+\delta_{\max}$};
    \draw[arrow, black!70, thick] (0.4,0) -- (2.2,0) node[right, label] {$\delta=0$};
    \draw[arrow, red!60!black, thick] (0.4,0.0) .. controls (1.4,-0.1) and (2.0,-0.9) .. (2.2,-1.9)
      node[below, label] {$\delta=-\delta_{\max}$};
  \end{tikzpicture}
  \caption{Hybrid A* 从一个状态出发的三条运动基元：左转、直行、右转。每条基元都满足式~\eqref{eq:19-bicycle} 的自行车模型，因此生成的路径天然满足最小转弯半径；搜索时把每条基元的终点作为新状态压入 open 表。}
  \label{fig:19-hybrid}
\end{figure}

Hybrid A* 适合车、差速底盘、以及一切“不能原地 $360^\circ$ 旋转”的平台。代价是状态空间变成三维，节点数明显增加；工程上常配“先粗后细”的两阶段（低分辨率搜索 + 高分辨率精修），以及 anytime 版本（先给一条可行解，再持续优化）。

\section{采样规划：RRT 与 RRT*}\label{sec:19-sampling}

当状态维数很高（机械臂 6--7 维、无人机 3 维位置 + 姿态）或地图没有显式栅格时，图搜索不再适用，改用\textbf{采样规划}：在状态空间随机撒点，把点连成一棵树，直到连到目标。

RRT（快速扩展随机树）每次采样 $q_{\mathrm{rand}}$，在树中找最近节点 $q_{\mathrm{near}}$，从 $q_{\mathrm{near}}$ 朝 $q_{\mathrm{rand}}$ 走固定步长 $\eta$ 得到 $q_{\mathrm{new}}$，若边不与障碍相交则加入树。RRT 是概率完备的（采样足够多时以概率 1 找到解），但找到的路径往往很绕、不最优。

RRT*\cite{karaman2011rrtstar} 做了两处关键改进。第一，\textbf{选择父节点}：不以“最近”为准，而是在半径 $r_n$ 内的所有节点中，选“把它接到 $q_{\mathrm{new}}$ 后总代价最小”的那个当父节点。第二，\textbf{重连}（rewire）：对半径 $r_n$ 内的每个节点 $q$，若“经过 $q_{\mathrm{new}}$ 到 $q$”比“$q$ 原来的路径”更短，就把 $q$ 的父节点改成 $q_{\mathrm{new}}$。选择半径随节点数增加而收缩，定义为
\begin{equation}
  r_n=\gamma\left(\frac{\log n}{n}\right)^{1/d}
  \label{eq:19-rrtstar}
\end{equation}
其中 $n$ 是当前节点数，$d$ 是状态维数，$\gamma$ 是常数。$r_n$ 随 $n$ 增大而缩小：早期重连范围大（快速优化），后期范围小（局部精修）。在这套机制下，RRT* 是渐近最优的：采样数 $n\to\infty$ 时，返回路径的代价以概率 1 收敛到最优。

\begin{figure}[htbp]
  \centering
  \begin{tikzpicture}[scale=1.35, font=\footnotesize]
    \coordinate (root) at (0,0);
    \coordinate (a) at (1.0,0.5);
    \coordinate (b) at (1.8,1.4);
    \coordinate (c) at (0.6,1.5);
    \coordinate (d) at (2.4,0.9);
    \coordinate (e) at (3.2,2.1);
    \coordinate (f) at (2.1,2.8);
    \coordinate (g) at (4.0,3.2);
    \coordinate (q) at (2.6,1.7);
    \draw[black!55, thick] (root) -- (a);
    \draw[black!55, thick] (a) -- (b);
    \draw[black!55, thick] (a) -- (c);
    \draw[black!55, thick] (a) -- (d);
    \draw[black!55, thick] (d) -- (e);
    \draw[black!55, thick] (b) -- (f);
    \draw[black!55, thick] (e) -- (g);
    \draw[green!45!black, thick, dashed] (q) circle (1.2);
    \draw[arrow, green!45!black, line width=1.2pt] (b) -- (q);
    \draw[arrow, red!60!black, thick, dashed] (q) -- (e);
    \foreach \p in {root,a,b,c,d,e,f,g}
      \fill[structurecolor] (\p) circle (1.5pt);
    \fill[green!45!black] (q) circle (1.8pt);
    \node[below, label, green!45!black] at (q) {$q_{\mathrm{new}}$};
    \node[circle, fill=red!70!black, inner sep=2pt] (goal) at (4.9,4.3) {};
    \node[right, label] at (goal) {目标};
    \node[green!45!black, label, above, sloped] at (2.2,1.55) {最优父节点};
    \node[red!60!black, label, above, sloped] at (2.95,1.95) {重连};
    \node[label, below left] at (root) {根};
  \end{tikzpicture}
  \caption{RRT* 的“选父节点 + 重连”。绿色虚线圈是选择半径 $r_n$；绿色实线是被选中的最优父边；红色虚线是重连后更短的连接。RRT* 用这种“不断改父节点”的方式，把一棵随机的树逐渐拉向最优解。}
  \label{fig:19-rrtstar}
\end{figure}

工程上 RRT* 的主要问题是收敛慢、结果依赖随机数，所以常配 \textbf{Informed RRT*}：先用一次搜索找到初始解，再把采样限制在“代价不超过当前最优解”的椭圆内，加速收敛。高维机械臂规划多用 RRT-Connect 或 OMPL 里的各种变体。

\section{其他经典方法与工程实践}\label{sec:19-others}

\textbf{可视图}（visibility graph）：把障碍物的顶点与起终点两两连线，保留不与障碍相交的边，在图上跑 Dijkstra。它在多边形障碍中给出最短路径，节点数少、结果精确，但“贴边”太危险，且只适合多边形表示的地图。Voronoi 图则相反：它给出“离障碍最远”的路径，安全性好、长度偏长，适合做保守规划。第~\ref{ch:21-3d-nav}~章会看到 FAR Planner 如何在 3D 中增量维护可视图。

\textbf{人工势场}（potential field）是另一类思路：把目标设为引力源、障碍设为斥力源，机器人沿合力方向走。它极快、极简单，但有局部极小值（机器人卡在凹形障碍前），第~\ref{ch:20-local-planning}~章会给出完整公式并讨论改进。

在 Nav2 中，全局规划由 \texttt{planner\_server} 的插件实现，常用的有：

\begin{itemize}
  \item \texttt{NavfnPlanner}：在代价地图上跑 Dijkstra/A*（可切换），支持 4/8 邻接，最经典、最稳。
  \item \texttt{SmacPlanner2D}：显式 A*，支持未知区域代价、可配置启发权重，适合 2D 差速底盘。
  \item \texttt{SmacPlannerHybrid}：Hybrid A*，支持倒车、最小转弯半径，适合阿克曼与“不能原地转”的平台。
  \item \texttt{SmacPlannerLattice}：状态格（lattice）规划，适合需要精确运动原语的场景。
\end{itemize}

\begin{lstlisting}[language=Python]
# nav2_params.yaml (excerpt): pick and tune the global planner plugin
planner_server:
  ros__parameters:
    planner_plugins: ["GridBased", "Hybrid"]
    GridBased:                  # 2D A*/Dijkstra for differential drive
      plugin: "nav2_navfn_planner/NavfnPlanner"
      use_astar: true
      tolerance: 0.5
    Hybrid:                     # Hybrid A* for car-like platforms
      plugin: "nav2_smac_planner/SmacPlannerHybrid"
      motion_model: "DUBINS"    # use REEDS_SHEPP if reverse is allowed
      minimum_turning_radius: 0.4
\end{lstlisting}

代价值如何进入搜索？NavFn 风格的实现把每个栅格的代价值 $c_{\mathrm{cell}}\in[0,254]$ 映射为附加代价
\begin{equation}
  w=c_{\mathrm{neutral}}+\kappa\,c_{\mathrm{cell}},
\end{equation}
其中中性代价 $c_{\mathrm{neutral}}$ 默认取 50（表示“平地上走一格的基准代价”），缩放系数 $\kappa$ 对应参数 \texttt{cost\_factor}，默认 3.0；致命障碍（254）直接禁止通行。这样“贴墙走”的边会明显变贵，规划器自然选择离墙远一点。

最后，规划器给出的路径往往仍有小折角，Nav2 用 \texttt{smoother\_server} 做后处理：以“平滑项 + 贴近原路径项 + 曲率/障碍惩罚”构成目标函数，梯度下降迭代优化路径点。平滑后的路径更贴合底盘运动学，局部控制器跟踪起来也更轻松。

\begin{summary}

\begin{enumerate}
  \item 全局规划是带非负边权的图搜索问题；完备性问“找不找得到”，最优性问“是否最短”，启发函数是对剩余代价的估计。
  \item Dijkstra 无启发、保证最优，复杂度 $O((V+E)\log V)$；A* 用 $f=g+h$ 引导，$h$ 可采纳时最优，一致时可“不重开节点”。
  \item 可采纳性证明的三步：前沿上存在最优路径节点且其 $g$ 已最优，于是该节点 $f\le C^*$；目标以最小 $f$ 弹出，故 $C\le C^*$。
  \item 启发函数曼哈顿/欧氏/对角距离都是下界；越紧扩展越少；陷阱地形中 A* 会退化为等值带扩张。
  \item JPS 用强制邻居跳过直道；Hybrid A* 用自行车模型基元 + 三维去重 + 双下界启发，解决最小转弯半径问题；RRT* 用“选父 + 重连 + 收缩半径”实现渐近最优。
  \item 工程上按底盘选插件：2D 差速用 NavFn 或 Smac2D，阿克曼用 Smac Hybrid；代价值经“中性代价 + 缩放系数”映射为边权，最后用 \texttt{smoother\_server} 平滑。
\end{enumerate}


\end{summary}

\begin{problemset}

  \item 用反证法复述 A* 最优性证明的三步，并指出每一步用了哪个前提。
  \item 在 $10\times10$ 栅格上，$s=(1,1)$、$t=(8,8)$，若某格被误标成致命障碍，曼哈顿启发是否仍然可采纳？启发值变不变？
  \item 为什么 $h=\max(h_{\mathrm{kin}},h_{\mathrm{grid}})$ 仍然可采纳？如果取两者的最小值会怎样？
  \item 给定 $L=0.6\ \mathrm{m}$、$\delta_{\max}=25\degree$，求最小转弯半径；若每步前进 $\Delta s=0.15\ \mathrm{m}$，求每步的转角增量。
  \item 设计一个“凹形房间”地图，说明 A* 与 RRT* 各会表现如何；如果你只能改一个参数来加速 A*，改什么？
\end{problemset}

\begin{reading}

\begin{itemize}
  \item \cite{hart1968astar}：A* 的原始论文，可采纳性与一致性的定义都出自这里，证明写得比教材更原始也更清楚。
  \item \cite{harabor2011jps}：JPS 原始论文，读它解决“直道上的节点为什么可以整段跳过”。
  \item \cite{dolgov2010hybridastar}：Hybrid A* 的工程实现细节（启发、基元、两阶段搜索），自动驾驶规划必读。
  \item \cite{karaman2011rrtstar}：RRT* 的渐近最优性证明与半径公式，想理解采样规划的理论极限就读它。
\end{itemize}


\end{reading}

\chapter{局部规划与控制}\label{ch:20-local-planning}

\usetikzlibrary{arrows.meta,positioning,calc,fit,patterns,decorations.pathreplacing}

\chapterintro{全局规划给出的是“大方向”，它只有几何、没有时间，也不考虑机器人此刻的速度、加速度和脚下 1 米内突然出现的障碍。局部规划（local planning）要在这条全局路径附近、在毫秒级的时间里，算出一条未来 1--3 秒真正能执行、不撞障碍、满足底盘约束的轨迹，并把它变成速度指令。本章按“采样—优化—预测控制—梯度优化”的演进线讲四种方法：DWA、TEB、MPC 与 EGO-Planner 式的 B 样条梯度法，再补上人工势场与底盘运动学约束，最后给出 Nav2 四个 controller 插件的选型对比。}

\section{局部规划要解决什么}\label{sec:20-what}

局部规划的输入是：全局路径 $\mathcal{P}$、局部代价地图、当前位姿与速度 $(\bm{x},\bm{u})$、机器人模型（尺寸、$v_{\max}$、$a_{\max}$、$\omega_{\max}$、最小转弯半径）。输出是：未来一段时间的控制序列，或直接是当前时刻的速度指令 $(v,\omega)$。它必须同时满足四件事：

\begin{enumerate}
  \item \textbf{跟踪}：轨迹要贴着全局路径走，不能偏离太远；
  \item \textbf{避障}：与障碍保持安全距离，包括代价地图里新出现的动态障碍；
  \item \textbf{可行}：满足速度、加速度、转角限制，不能算出底盘执行不了的动作；
  \item \textbf{实时}：10--20 Hz 内算完，因为局部规划要在障碍出现后 100 ms 量级内做出反应。
\end{enumerate}

先用一个数字把“实时”的分量说清楚：控制周期 50 ms（20 Hz），留给局部规划的计算时间通常要小于 20 ms（剩下的留给通信、TF 查询与控制平滑）。这意味着任何单次迭代超过 20 ms 的方法都必须降频、降精度，或者只在低维空间里做。

\section{DWA：把可行速度采样一遍}\label{sec:20-dwa}

DWA（Dynamic Window Approach）\cite{fox1997dwa} 的思路非常朴素：\textbf{不去设计轨迹，而是枚举“下一时刻可能的速度”，把每个速度对应的短轨迹算出来打分，选最高分的那个执行。}它之所以叫“动态窗口”，是因为待枚举的速度范围受当前速度和加速度限制，是一个随时间“滑动”的窗口。

\subsection{运动模型与动态窗口}

差速底盘（unicycle）的运动学为
\begin{equation}
  \dot x=v\cos\theta,\qquad \dot y=v\sin\theta,\qquad \dot\theta=\omega,
  \label{eq:20-unicycle}
\end{equation}
离散化后（步长 $\Delta t$）
\begin{equation}
  x_{k+1}=x_k+v\Delta t\cos\theta_k,\quad
  y_{k+1}=y_k+v\Delta t\sin\theta_k,\quad
  \theta_{k+1}=\theta_k+\omega\Delta t.
  \label{eq:20-dwa-model}
\end{equation}
对某一个候选 $(v,\omega)$，从当前位姿递推 $N$ 步就得到一条候选轨迹；若轨迹上任意一点落在致命障碍格上，该候选直接淘汰。

可枚举的速度集合是三个集合的交：
\begin{itemize}
  \item 硬约束 $V_s=\{(v,\omega)\mid v\in[v_{\min},v_{\max}],\ \omega\in[-\omega_{\max},\omega_{\max}]\}$；
  \item 动态窗口 $V_d=\{(v,\omega)\mid v\in[v_c-a_v\Delta t,\ v_c+a_v\Delta t],\ \omega\in[\omega_c-a_\omega\Delta t,\ \omega_c+a_\omega\Delta t]\}$，其中 $v_c,\omega_c$ 是当前速度，$a_v,a_\omega$ 是最大线/角加速度；
  \item 安全约束 $V_a=\{(v,\omega)\mid v\le\sqrt{2\,d_{\mathrm{stop}}(v,\omega)\,a_{\mathrm{brake}}}\}$，表示“以该速度走，还能在撞上之前刹住”。$d_{\mathrm{stop}}$ 是沿该轨迹到最近障碍的距离。
\end{itemize}
最终 $V_r=V_s\cap V_d\cap V_a$。举个数：$v_c=0.5\ \mathrm{m/s}$、$a_v=1.0\ \mathrm{m/s^2}$、$\Delta t=0.1\ \mathrm{s}$，则 $v$ 只能取 $[0.4,0.6]$；若 $v_{\max}=1.2$，真正限制搜索的是动态窗口而不是上限——这正是“动态”的意义：速度不能突变。

\subsection{评价函数}

对每个候选 $(v,\omega)$，用三个（或五个）归一化指标打分：
\begin{align}
  \mathrm{heading}(v,\omega) &= \pi-\abs{\theta_{\mathrm{goal}}-\theta(t_N)}, \label{eq:20-heading}\\
  \mathrm{clearance}(v,\omega) &= \min_{k\in[0,N]} d_{\mathrm{obs}}\bigl(x_k,y_k\bigr), \label{eq:20-clear}\\
  \mathrm{velocity}(v,\omega) &= \abs{v}, \label{eq:20-vel}
\end{align}
其中 $\theta_{\mathrm{goal}}=\mathrm{atan2}(y_g-y_0,\ x_g-x_0)$ 是当前点到目标的方位角，$\theta(t_N)$ 是候选轨迹终点的朝向，$d_{\mathrm{obs}}$ 是该点到最近障碍的距离（代价地图的距离变换直接给出）。三项量纲不同，必须先归一化。DWA 原文与 DWB 的常用做法是对所有候选求和后相除：
\begin{equation}
  \hat h=\frac{\mathrm{heading}}{\sum\mathrm{heading}},\quad
  \hat c=\frac{\mathrm{clearance}}{\sum\mathrm{clearance}},\quad
  \hat v=\frac{\mathrm{velocity}}{\sum\mathrm{velocity}},
\end{equation}
然后
\begin{equation}
  G(v,\omega)=\alpha\,\hat h+\beta\,\hat c+\gamma\,\hat v,\qquad \alpha+\beta+\gamma=1.
  \label{eq:20-G}
\end{equation}
$\alpha$ 大则“更想朝向目标”，$\beta$ 大则“更想远离障碍”，$\gamma$ 大则“更想跑得快”。三者的物理含义分别是“方向一致性”“安全裕度”“任务效率”。

\begin{table}[htbp]
  \centering
  \caption{一个 DWA 打分的数值例子（$T=1\ \mathrm{s}$，$\alpha=0.6$、$\beta=0.3$、$\gamma=0.1$）}
  \label{tab:20-dwa}
  \small
  \begin{tabular}{lccccc}
    \toprule
    候选 $(v,\omega)$ & 终点 & $\hat h$ & $\hat c$ & $\hat v$ & $G$ \\
    \midrule
    $(0.5,\ 0)$ 直行 & $(0.50,0)$ & 0.852 & 0.89 & 0.50 & 0.828 \\
    $(0.5,\ 0.5)$ 左转 & $(0.48,0.12)$ & 0.992 & 0.56 & 0.50 & 0.813 \\
    $(0.5,-0.5)$ 右转 & $(0.48,-0.12)$ & 0.639 & 1.00 & 0.50 & 0.733 \\
    \bottomrule
  \end{tabular}
\end{table}

表~\ref{tab:20-dwa} 里，左转的朝向得分最高（0.992），但离障碍更近导致安全得分低；右转最安全但方向偏得最多；直行两项都不极端，加权后胜出。这个例子说明两点：第一，\textbf{权重决定行为风格}，把 $\beta$ 调大就会变成“贴墙恐惧症”；第二，\textbf{分数接近时机器人会犹豫}（0.828 与 0.813 只差 0.015），实际 DWB 会对 $G$ 做时间上的低通平滑，避免相邻周期来回切换。

\subsection{为什么能实时，局限在哪}

\begin{figure}[htbp]
  \centering
  \begin{tikzpicture}[font=\footnotesize]
    \draw[arrow] (0,0) -- (3.3,0) node[right] {$v$};
    \draw[arrow] (0,-1.5) -- (0,1.7) node[above] {$\omega$};
    \draw[structurecolor!45, dashed] (0.4,-1.3) rectangle (3.0,1.4);
    \draw[structurecolor, thick] (1.0,-0.9) rectangle (2.4,0.85);
    \foreach \x in {1.15,1.55,1.95,2.35}
      \foreach \y in {-0.7,-0.35,0,0.35,0.7}
        \fill[black!60] (\x,\y) circle (0.6pt);
    \node[label, structurecolor!70] at (0.5,1.55) {$V_d$};
    \node[label, structurecolor!70] at (2.85,1.55) {$V_s$};
    \node[label] at (1.6,1.95) {速度空间};
    \begin{scope}[xshift=5.4cm]
      \draw[arrow] (0,0) -- (2.9,0) node[right] {$x$};
      \draw[arrow] (0,0) -- (0,2.3) node[above] {$y$};
      \fill[black!18] (2.15,1.35) circle (0.32);
      \node[black!55, label] at (2.15,0.85) {障碍};
      \fill[structurecolor!25] (-0.22,-0.18) rectangle (0.22,0.18);
      \draw[arrow, black!65] (0,0) .. controls (0.7,0.02) and (1.4,0.05) .. (2.1,0.12);
      \draw[arrow, green!45!black] (0,0) .. controls (0.45,0.5) and (0.85,1.3) .. (0.85,1.85);
      \draw[arrow, red!60!black] (0,0) .. controls (0.9,-0.35) and (1.9,-0.1) .. (2.35,0.55);
      \node[label] at (1.45,2.15) {候选轨迹};
    \end{scope}
  \end{tikzpicture}
  \caption{DWA 的两层搜索。左：在速度空间 $(v,\omega)$ 中，粗框是动态窗口 $V_d$，虚框是硬约束 $V_s$，黑点是采样点；右：每个采样点积分出一条短轨迹，撞障碍的淘汰，其余按式~\eqref{eq:20-G} 打分。}
  \label{fig:20-dwa}
\end{figure}

计算量很容易估：以 $v$ 方向 10 个、$\omega$ 方向 15 个采样计，共 150 条候选；每条积分 $N=30$ 步、每步 3 个三角函数与一次碰撞查询，总计约 $1.5\times10^4$ 次基本运算，在 20 Hz 的周期里绰绰有余。这就是 DWA 能在 2000 年代的嵌入式 CPU 上跑实时避障的原因。

但 DWA 的缺点同样明显：它只在一步速度空间里做贪心，\textbf{没有全局视野}，容易陷在局部最优（U 形障碍前反复打转）；对“必须倒车才能脱困”的场景无能为力；$T$ 只有 1--2 s，对高速或大惯量平台不够；参数（$\alpha,\beta,\gamma$、分辨率、$T$）与场地强相关，换场地要重调。

\section{TEB：把轨迹变成图优化}\label{sec:20-teb}

TEB（Timed Elastic Band）\cite{rosmann2012teb} 换了个思路：不再枚举控制，而是把轨迹表示成一串带时间戳的位姿，\textbf{把“轨迹好不好”写成一个目标函数，用图优化求解}。

轨迹变量是位姿序列 $\mathcal{Q}=\{\bm{s}_0,\bm{s}_1,\dots,\bm{s}_n\}$，其中 $\bm{s}_k=(x_k,y_k,\beta_k)$（$\beta_k$ 为朝向），相邻位姿之间还有时间间隔 $\Delta T_k>0$。目标函数由五类惩罚组成：
\begin{equation}
\begin{aligned}
J(\mathcal{Q},\Delta T)=\;&
w_h\sum_k \bigl[\max(0,r_h-d(\bm{s}_k,\mathcal{P}))\bigr]^2
+w_o\sum_k \bigl[\max(0,r_o-d(\bm{s}_k,\mathcal{O}))\bigr]^2\\
&+J_{\mathrm{kin}}+J_{\mathrm{dyn}}+J_{\mathrm{time}},
\end{aligned}
\label{eq:20-teb}
\end{equation}
其中
\begin{align}
  J_{\mathrm{dyn}} &= \textstyle\sum_k \Bigl[\max(0,\abs{v_k}-v_{\max})^2+\max(0,\abs{a_k}-a_{\max})^2\Bigr], \label{eq:20-teb-dyn}\\
  J_{\mathrm{kin}} &= \textstyle\sum_k \bigl[(y_{k+1}-y_k)\cos\beta_k-(x_{k+1}-x_k)\sin\beta_k\bigr]^2, \label{eq:20-teb-kin}\\
  J_{\mathrm{time}} &= \Bigl(\textstyle\sum_k \Delta T_k\Bigr)^2. \label{eq:20-teb-time}
\end{align}

逐项解释：$d(\bm{s}_k,\mathcal{P})$ 是位姿到全局路径的距离，超过阈值 $r_h$ 才罚，保证“不偏离大方向”；$d(\bm{s}_k,\mathcal{O})$ 是到障碍的距离，小于安全裕度 $r_o$ 才罚，保证“不贴墙”；$J_{\mathrm{kin}}$ 惩罚“斜着走”（速度方向与朝向不一致），这是差速底盘的硬约束；$J_{\mathrm{dyn}}$ 用 $v_k=\Delta s_k/\Delta T_k$ 与 $a_k=(v_{k+1}-v_k)/\Delta T_k$ 约束速度、加速度；$J_{\mathrm{time}}$ 让总时间尽量短（否则机器人会磨蹭）。

式~\eqref{eq:20-teb} 的每一项都只涉及 1--2 个位姿节点，天然是\textbf{稀疏图}：位姿是节点，惩罚项是超边（hyper-edge），时间间隔也是节点。用 g2o\cite{kummerle2011g2o} 建图后，做稀疏 Levenberg--Marquardt 迭代，通常 2--5 次迭代收敛；把上一周期的解作为初值（warm start），10--20 Hz 完全能跑。位姿数 $n$ 通常取 10--30，对应 1--3 s 的轨迹。

\begin{figure}[htbp]
  \centering
  \begin{tikzpicture}[font=\footnotesize, scale=1.0]
    \coordinate (p0) at (0,0.3);
    \coordinate (p1) at (1.1,0.7);
    \coordinate (p2) at (2.2,1.3);
    \coordinate (p3) at (3.1,2.1);
    \coordinate (p4) at (3.9,3.0);
    \draw[structurecolor, thick] (p0) -- (p1) -- (p2) -- (p3) -- (p4);
    \foreach \p in {p0,p1,p2,p3,p4}
      \fill[structurecolor!25, draw=structurecolor, thick] (\p) circle (0.09);
    \fill[black!18] (4.2,1.2) circle (0.45);
    \node[black!60, label] at (4.2,0.6) {障碍};
    \draw[red!60!black, thick, dashed] (p2) -- (3.75,1.22);
    \draw[red!60!black, thick, dashed] (p3) -- (3.85,1.49);
    \node[red!60!black, label] at (3.15,0.95) {$d(\bm{s}_k,\mathcal{O})$};
    \node[label, fill=white, inner sep=1pt] at (0.55,0.5) {$\Delta T_1$};
    \node[label, fill=white, inner sep=1pt] at (1.65,1.0) {$\Delta T_2$};
    \node[label, fill=white, inner sep=1pt] at (2.65,1.7) {$\Delta T_3$};
    \node[label, fill=white, inner sep=1pt] at (3.5,2.55) {$\Delta T_4$};
    \node[below, label] at (p0) {$\bm{s}_0$};
    \node[below, label] at (p1) {$\bm{s}_1$};
    \node[below, label] at (p2) {$\bm{s}_2$};
    \node[below, label] at (p3) {$\bm{s}_3$};
    \node[right, label] at (p4) {$\bm{s}_4$};
  \end{tikzpicture}
  \caption{TEB 的图结构：圆形是位姿节点 $\bm{s}_k$（优化变量），相邻位姿之间的连边带时间间隔 $\Delta T_k$，红色虚线是到障碍的“距离惩罚”超边。优化时全体位姿一起移动、时间间隔一起缩放，所以 TEB 能同时改变“路径形状”和“速度分配”。}
  \label{fig:20-teb}
\end{figure}

TEB 的好处是轨迹平滑、能处理车式底盘（非完整约束写进 $J_{\mathrm{kin}}$）、能把“时间”也当作优化变量，因此停车、让行、窄道通过都比较自然。代价是：目标函数非凸，初值不好会收敛到局部极小；每周期都要建图求解，CPU 占用比 DWB 高一个量级；参数（$w_h,w_o$、$r_h,r_o$）比 DWA 更难调。

\section{MPC：滚动时域的最优控制}\label{sec:20-mpc}

MPC（Model Predictive Control）\cite{rawlings2017mpc} 是控制理论的标准答案：用模型预测未来 $N$ 步，在每条预测轨迹上优化控制序列，只执行第一步，下一周期再重新预测——这就是\textbf{滚动时域}（receding horizon）。

设离散状态空间模型
\begin{equation}
  \bm{x}_{k+1}=\bm{A}\bm{x}_k+\bm{B}\bm{u}_k,
  \label{eq:20-mpc-model}
\end{equation}
其中 $\bm{x}=[x,y,\theta,v,\omega]^{\top}$，$\bm{u}=[a,\alpha]^{\top}$（线加速度、角加速度），$\bm{A},\bm{B}$ 由式~\eqref{eq:20-unicycle} 在当前工作点线性化得到。每个控制周期求解
\begin{equation}
  \begin{aligned}
    \min_{\bm{x}_{1:N},\bm{u}_{0:N-1}}\quad
    &\sum_{k=0}^{N-1}\bigl(\bm{x}_k-\bm{x}_k^{\mathrm{ref}}\bigr)^{\top}\bm{Q}\bigl(\bm{x}_k-\bm{x}_k^{\mathrm{ref}}\bigr)
     +\bm{u}_k^{\top}\bm{R}\bm{u}_k
     +\bigl(\bm{x}_N-\bm{x}_N^{\mathrm{ref}}\bigr)^{\top}\bm{Q}_f\bigl(\bm{x}_N-\bm{x}_N^{\mathrm{ref}}\bigr)\\
    \text{s.t.}\quad
    &\bm{x}_{k+1}=\bm{A}\bm{x}_k+\bm{B}\bm{u}_k,\quad \bm{x}_0=\hat{\bm{x}},\\
    &\bm{u}_{\min}\le\bm{u}_k\le\bm{u}_{\max},\quad \bm{x}_k\in\mathcal{X}_{\mathrm{free}}.
  \end{aligned}
  \label{eq:20-mpc}
\end{equation}
$\bm{Q}\succeq0$、$\bm{R}\succ0$ 是权重矩阵：$\bm{Q}$ 大则“更想贴合参考轨迹”，$\bm{R}$ 大则“更不想用大加速度”。约束里的 $\mathcal{X}_{\mathrm{free}}$ 是“不撞障碍”，工程上把每个附近障碍近似成半平面 $\bm{a}_j^{\top}\bm{x}\le b_j$ 加进去，于是整个问题变成\textbf{二次规划}（QP），可用 OSQP 之类的稀疏求解器在毫秒级解出。

规模感受一下：$N=20$、状态 5 维、控制 2 维，决策变量 $20\times5+20\times2=140$ 个，约束约几百条，20 Hz 下现代 CPU 完全够用。MPC 与 DWA 的关系可以这样理解：\textbf{DWA 是在离散控制集上做一步贪心，MPC 是在连续控制空间里做 $N$ 步带约束的优化}；DWA 不需要模型（只用运动学积分），MPC 需要模型和求解器，但能显式处理加速度、障碍、终端代价，且天然支持预测动态障碍（把障碍位置写成时间的函数即可）。

\begin{figure}[htbp]
  \centering
  \begin{tikzpicture}[font=\footnotesize, scale=1.05]
    \draw[arrow] (0,0) -- (7.6,0) node[right] {时间};
    \draw[arrow] (0,0) -- (0,3.2) node[above] {状态};
    \draw[structurecolor!60, thick]
      (0.2,0.6) .. controls (1.3,2.5) and (2.4,2.6) .. (3.2,2.35)
      .. controls (4.2,2.05) and (5.2,1.1) .. (6.2,0.8)
      .. controls (6.8,0.6) and (7.2,0.55) .. (7.5,0.55);
    \node[structurecolor!70, label] at (1.3,0.85) {参考轨迹};
    \draw[structurecolor, thick, dashed]
      (2.6,2.15) .. controls (3.3,2.3) and (4.0,2.1) .. (4.7,1.75);
    \foreach \x/\y in {2.85/2.28,3.3/2.3,3.75/2.22,4.2/2.05,4.65/1.82}
      \draw[structurecolor, thick] (\x-0.12,\y-0.12) rectangle (\x+0.12,\y+0.12);
    \fill[black!70] (2.6,2.15) circle (2pt);
    \node[label, below left] at (2.35,1.95) {当前状态 $\hat{\bm{x}}$};
    \node[structurecolor, label] at (3.7,2.75) {预测时域 $N$ 步};
    \draw[arrow, red!60!black, thick] (6.4,1.6) -- (7.1,1.6);
    \node[red!60!black, label, above] at (6.75,1.72) {下一周期整体右移};
  \end{tikzpicture}
  \caption{MPC 的滚动时域：每个控制周期都以当前状态为起点，预测未来 $N$ 步并优化；只执行第一步控制，下一周期把时域整体向前平移、用新的测量重新优化。}
  \label{fig:20-mpc}
\end{figure}

\section{梯度类方法：B 样条轨迹优化}\label{sec:20-gradient}

DWA 采样、TEB 图优化、MPC 求解 QP，都是在“低维控制/位姿”上做文章。无人机、四足等高速平台更常用\textbf{连续时间轨迹优化}：用 B 样条参数化整条轨迹，直接对控制点求梯度\cite{zhou2019fastplanner,zhou2021egoplanner}。

$p$ 次 B 样条曲线为
\begin{equation}
  \bm{p}(t)=\sum_{i=0}^{n}B_{i,p}(t)\,\bm{Q}_i,
  \label{eq:20-bspline}
\end{equation}
其中 $\bm{Q}_i\in\mathbb{R}^3$ 是控制点，$B_{i,p}(t)$ 是基函数。它有两个对规划极有用的性质：\textbf{局部支撑}（改一个控制点只影响相邻 $p+1$ 段）与\textbf{凸包性}（每一段曲线都在相邻 $p+1$ 个控制点的凸包内）。

目标函数通常写成三项之和：
\begin{equation}
  J=\lambda_s J_{\mathrm{smooth}}+\lambda_c J_{\mathrm{collision}}+\lambda_d J_{\mathrm{dynamic}}.
  \label{eq:20-J}
\end{equation}
平滑项用二阶差分（近似加速度/曲率）：
\begin{equation}
  J_{\mathrm{smooth}}=\sum_{i=1}^{n-1}\norm{\bm{Q}_{i+1}-2\bm{Q}_i+\bm{Q}_{i-1}}^2.
  \label{eq:20-smooth}
\end{equation}
碰撞项需要一个“到障碍的距离场”。经典做法是把地图做成 \textbf{ESDF}（欧氏符号距离场）\cite{oleynikova2017voxblox}：每个体素存“到最近障碍的距离” $d(\bm{p})$，其梯度 $\nabla d(\bm{p})$ 指向远离障碍的方向。碰撞代价取二次惩罚
\begin{equation}
  J_{\mathrm{collision}}=\sum_{i}\bigl[\max\bigl(0,\ d_0-d(\bm{Q}_i)\bigr)\bigr]^2,
  \label{eq:20-coll}
\end{equation}
$d_0$ 是安全距离。对控制点求偏导，用链式法则：
\begin{equation}
  \frac{\partial J_{\mathrm{collision}}}{\partial \bm{Q}_i}
  =-2\bigl(d_0-d(\bm{Q}_i)\bigr)\,\nabla d(\bm{Q}_i),
  \qquad d(\bm{Q}_i)<d_0.
  \label{eq:20-coll-grad}
\end{equation}
举个数值：$d_0=0.5\ \mathrm{m}$、某控制点处 $d=0.2\ \mathrm{m}$，则惩罚为 $(0.3)^2=0.09$，梯度大小为 $2\times0.3=0.6$，方向沿 $\nabla d$（远离障碍），优化后该控制点被推出去。

\begin{pseudo}{ESDF 梯度下降优化一段 B 样条轨迹}
  \item 用前端搜索（A*/Hybrid A*）给出一条初始轨迹，转成 B 样条控制点 $\bm{Q}_i$
  \item 重复直到收敛或达到迭代上限：
  \item \quad 对每个控制点，从 ESDF 查 $d(\bm{Q}_i)$ 与 $\nabla d(\bm{Q}_i)$，按式~\eqref{eq:20-coll-grad} 累加碰撞梯度
  \item \quad 按式~\eqref{eq:20-smooth} 累加平滑梯度，两项加权求和得 $\partial J/\partial\bm{Q}_i$
  \item \quad 沿负梯度更新 $\bm{Q}_i$，并检查凸包性质（控制点都不在障碍内则整段安全）
  \item 输出轨迹，交给控制器跟踪
\end{pseudo}

EGO-Planner\cite{zhou2021egoplanner} 的关键改进是\textbf{不显式建 ESDF}：它用“碰撞对 $\{\bm{p},\bm{v}\}$”（轨迹上的点与对应的障碍表面点）直接构造惩罚，并利用 B 样条凸包性质保证优化后的轨迹不穿障碍。这省掉了建 ESDF 的内存与计算，代价是碰撞对的维护逻辑更复杂。对 RoboMaster 这类 2D 主要、偶尔上坡的场地，完整的 ESDF 通常负担得起，可以先从 Fast-Planner 思路入手。

\section{人工势场：最直观也最容易卡住}\label{sec:20-apf}

人工势场把目标设成引力源、障碍设成斥力源。引力势与引力（梯度取负）为
\begin{equation}
  U_{\mathrm{att}}(\bm{q})=\tfrac12 k_a\norm{\bm{q}-\bm{q}_g}^2,
  \qquad
  \bm{F}_{\mathrm{att}}=-\nabla U_{\mathrm{att}}=k_a(\bm{q}_g-\bm{q}).
  \label{eq:20-att}
\end{equation}
斥力势（$\rho$ 为到障碍距离，$\rho_0$ 为影响半径）为
\begin{equation}
  U_{\mathrm{rep}}(\bm{q})=
  \begin{cases}
    \tfrac12 k_r\left(\dfrac{1}{\rho}-\dfrac{1}{\rho_0}\right)^2, & \rho<\rho_0,\\[6pt]
    0, & \rho\ge\rho_0,
  \end{cases}
  \qquad
  \bm{F}_{\mathrm{rep}}=k_r\left(\frac{1}{\rho}-\frac{1}{\rho_0}\right)\frac{1}{\rho^2}\,\nabla\rho.
  \label{eq:20-rep}
\end{equation}
数值例子：$k_a=2$，目标 $(5,0)$，当前位置 $(2,0)$，则 $\bm{F}_{\mathrm{att}}=(6,0)$；$k_r=0.5$、$\rho_0=1$、$\rho=0.5$，则 $U_{\mathrm{rep}}=0.25$，$\abs{\bm{F}_{\mathrm{rep}}}=0.5\times(2-1)\times4=2$，方向远离障碍。

势场法的问题是\textbf{局部极小值}：当引力与斥力恰好抵消且不在目标点时，机器人停住。凹形障碍、长墙尽头、目标位于墙后都会出现。改进思路有三类：加随机扰动/模拟退火跳出；用调和函数（满足拉普拉斯方程，理论上无局部极小，但计算贵）；最实用的是\textbf{用全局路径做引导}——沿全局路径取“虚拟目标”，让势场只负责局部避障，这就回到了“全局规划 + 局部规划”的分工。

\section{底盘约束如何进入局部规划}\label{sec:20-kinematics}

局部规划算出的 $(v,\omega)$ 必须能被执行器实现。三种常见底盘：

\textbf{差速}（两轮独立驱动）：设轮距 $b$、左右轮速 $v_l,v_r$，则
\begin{equation}
  v=\frac{v_r+v_l}{2},\qquad \omega=\frac{v_r-v_l}{b}.
  \label{eq:20-diff}
\end{equation}
例如 $v_r=1.0\ \mathrm{m/s}$、$v_l=0.6\ \mathrm{m/s}$、$b=0.4\ \mathrm{m}$，得 $v=0.8\ \mathrm{m/s}$、$\omega=1.0\ \mathrm{rad/s}$，转弯半径 $R=v/\omega=0.8\ \mathrm{m}$。约束是 $v_l,v_r\in[-v_{\max},v_{\max}]$，等价于 $v$ 与 $\omega$ 的联合约束（不是矩形）。

\textbf{麦轮}（全向）：四轮速度由底盘速度解算，例如前左轮
\begin{equation}
  v_{fl}=v_x-v_y-(l_x+l_y)\omega,
\end{equation}
其余三轮按符号变化。麦轮可以横移（$v_y\neq0$），所以局部规划的状态空间可以加一个横向速度，路径不必先转向再前进。

\textbf{阿克曼}（车式）：前轮转角 $\delta$、轴距 $L$，有
\begin{equation}
  \omega=\frac{v\tan\delta}{L},\qquad R=\frac{L}{\tan\delta}\ge R_{\min}=\frac{L}{\tan\delta_{\max}}.
  \label{eq:20-ackermann}
\end{equation}
$L=0.6\ \mathrm{m}$、$\delta_{\max}=25^\circ$ 时 $R_{\min}=0.6/0.466=1.29\ \mathrm{m}$。这条约束必须写进规划器：2D 差速用的 DWA 会算出“原地旋转”，阿克曼执行不了，所以要么用 Hybrid A* 做全局规划，要么在局部规划里用 TEB 的 $J_{\mathrm{kin}}$ 或 MPC 的模型约束。

\section{Nav2 的 controller 插件怎么选}\label{sec:20-plugins}

\begin{table}[htbp]
  \centering
  \caption{Nav2 常见 controller 插件对比}
  \label{tab:20-controllers}
  \small
  \begin{tabular}{p{2.7cm}p{2.5cm}p{3.1cm}p{3.4cm}}
    \toprule
    插件 & 类型 & 优点 & 适用/局限 \\
    \midrule
    DWB & 速度采样（DWA 的 Nav2 实现） & 稳、易调、CPU 低 & 差速底盘、中低速；窄道与局部最优弱 \\
    TEB & 图优化 & 轨迹平滑、支持阿克曼、可倒车 & CPU 高、初值敏感、参数多 \\
    MPPI & 采样 + 预测控制 & 处理动态障碍、平滑、GPU 友好 & 需要调采样数与温度参数 \\
    RPP & 纯跟踪几何法 & 极快、路径跟踪准 & 避障弱，靠代价地图限速 \\
    \bottomrule
  \end{tabular}
\end{table}

新手最常犯的错误是“上来就换 TEB”。正确的顺序是：先用 DWB 把 TF、代价地图、全局规划调通，确认能稳定跟踪；出现“窄道过不去、需要倒车、轨迹抖动”再考虑 TEB 或 MPPI。表~\ref{tab:20-controllers} 里没有“最好”，只有“匹配”。

\begin{lstlisting}[language=Python]
# Minimal DWB-style scoring (illustrative, not production code)
def score(candidates, goal, costmap, w=(0.6, 0.3, 0.1)):
    best, best_g = None, -1.0
    for v, w_ in candidates:
        traj = rollout(v, w_, T=1.0, dt=0.05)   # integrate unicycle
        if collides(traj, costmap):
            continue
        h = heading_score(traj[-1], goal)       # pi - |dtheta|
        c = clearance_score(traj, costmap)      # min obstacle distance
        s = speed_score(v)                      # |v| / v_max
        g = w[0] * h + w[1] * c + w[2] * s
        if g > best_g:
            best, best_g = (v, w_), g
    return best
\end{lstlisting}

\begin{summary}

\begin{enumerate}
  \item 局部规划在 1--3 s 的滚动窗口内，同时满足跟踪、避障、运动学可行与 20 Hz 实时四项要求。
  \item DWA 在动态窗口 $V_r=V_s\cap V_d\cap V_a$ 内采样 $(v,\omega)$，积分短轨迹，按式~\eqref{eq:20-G} 加权打分；快但贪心、易陷局部最优。
  \item TEB 把轨迹写成位姿 + 时间间隔的图，用式~\eqref{eq:20-teb} 的六类惩罚做 g2o 优化；平滑、支持非完整约束，但非凸、CPU 高。
  \item MPC 用模型预测 $N$ 步、解带约束的 QP、只执行第一步；比 DWA 多模型与约束，比 TEB 更适合显式预测动态障碍。
  \item 梯度法用 B 样条参数化轨迹，平滑项用二阶差分，碰撞项用 ESDF 的二次惩罚与梯度（式~\eqref{eq:20-coll-grad}），EGO-Planner 用碰撞对避免建 ESDF。
  \item 人工势场直观但有局部极小；底盘约束（差速式~\eqref{eq:20-diff}、阿克曼式~\eqref{eq:20-ackermann}）必须进入局部规划，否则算出的动作执行不了。
  \item Nav2 选型：DWB 打底，TEB/MPPI 解决窄道与动态障碍，RPP 用于快速纯跟踪。
\end{enumerate}


\end{summary}

\begin{problemset}

  \item 用式~\eqref{eq:20-dwa-model} 手算：初始位姿 $(0,0,0)$，$v=0.6\ \mathrm{m/s}$、$\omega=0.4\ \mathrm{rad/s}$、$\Delta t=0.1\ \mathrm{s}$，求 5 步后的位姿。
  \item 在表~\ref{tab:20-dwa} 的数据下，若把 $\beta$ 从 0.3 调到 0.8、$\alpha$ 降到 0.1（保持和为 1），重新排序三个候选，并解释机器人行为会怎样变化。
  \item TEB 的 $J_{\mathrm{kin}}$ 为什么能表达“不能横移”？如果机器人是麦轮，这一项应该怎么改？
  \item MPC 的 $\bm{Q}$ 与 $\bm{R}$ 分别调大，对轨迹的跟踪精度与舒适性有什么影响？
  \item 给定 $L=0.5\ \mathrm{m}$、$\delta_{\max}=30^\circ$，求阿克曼底盘的最小转弯半径；若要求 $R_{\min}\le0.8\ \mathrm{m}$，$\delta_{\max}$ 至少要多大？
  \item 有人说“DWA 就是 MPC 的特例”，你同意吗？请从目标函数、约束、求解方式三方面比较。
\end{problemset}

\begin{reading}

\begin{itemize}
  \item \cite{fox1997dwa}：DWA 原始论文，评价函数与动态窗口的定义都在这里，篇幅不长，适合精读。
  \item \cite{rosmann2012teb}：TEB 论文，重点看它如何把“时间间隔”也作为优化变量，以及 g2o 图的构建。
  \item \cite{rawlings2017mpc}：MPC 的标准教材，读第 1--2 章即可建立“滚动时域 + QP”的直觉。
  \item \cite{zhou2019fastplanner,zhou2021egoplanner}：B 样条轨迹优化的工程实现，重点看碰撞代价的梯度与“无 ESDF”版本的区别。
\end{itemize}


\end{reading}

\chapter{3D 导航与 CMU 自主探索栈}\label{ch:21-3d-nav}

\usetikzlibrary{arrows.meta,positioning,calc,fit,patterns,decorations.pathreplacing}

\chapterintro{前两章的规划器都假设世界是一张 2D 平面栅格：障碍是“墙”，可走区域是“空地”。但 RoboMaster 之外的很多场景——室外坡道、台阶、多层建筑、越野地形——机器人会上下起伏，2D 栅格会把“坡”压成“平地”，把“台阶”压成“墙”，规划出的路径要么上不去，要么明明能过却被判死。本章讲 3D 导航，并以用户点名的 \textbf{CMU（卡内基梅隆大学 Ji Zhang 组，CMU AirLab）自主探索开发环境}为主线：讲清 terrain\_analysis 怎么做 2.5D 地形可通行性分析、local\_planner 怎么在地形上避障、FAR Planner 怎么用可视图在未知环境里“试探着规划”、TARE 怎么做分层自主探索。学完本章，你应该能画出这套栈的模块数据流，能推导坡度/粗糙度/台阶高度的公式，并知道如何把 FAST-LIO2 的点云喂给它。}

\section{为什么 2D 导航不够}\label{sec:21-why3d}

2D 导航的本质是\textbf{降维假设}：把机器人当成一个在平面上运动的圆盘，把世界投影到 $z=0$ 的平面上。这个假设在平整室内非常好用，代价是丢掉三样信息：

\begin{enumerate}
  \item \textbf{高度}：坡道、台阶、门槛、减速带都有高度变化。2D 栅格一旦把“坡”投影成“可走”，机器人可能因坡度太陡而打滑；把“台阶”投影成“障碍”，则明明能上的小台阶被当成墙。
  \item \textbf{垂直净空}：桌子底下、桥洞、悬空管道下方，2D 地图只看到“有障碍”，3D 才知道“能从下面过”。
  \item \textbf{地形质量}：同样平坦的格子，草地、碎石、沙地的可通过性完全不同；2D 只区分“占据/空闲”。
\end{enumerate}

解决思路有两条：\textbf{真 3D}（体素地图、ESDF、无人机式轨迹优化）与 \textbf{2.5D}（把地图表示成“每个格子一个高度”的高程图）。真 3D 表达力强但内存与计算量大（$100\ \mathrm{m}\times100\ \mathrm{m}\times10\ \mathrm{m}$、$0.1\ \mathrm{m}$ 体素就是 $10^8$ 个体素）；2.5D 假设“每个 $(x,y)$ 只有一个高度”，内存是 2D 级别，又能表达坡度和台阶，是地面机器人最实用的折中。CMU 这套栈走的就是 2.5D 路线。

\section{CMU 自主探索栈的总体架构}\label{sec:21-arch}

CMU AirLab 开源了一套“点云进、速度出”的完整导航与探索系统，名为 Autonomous Exploration Development Environment（下称 CMU 栈），配套论文是 ICRA 2022 的 \cite{cao2022tare}，其中 FAR Planner 的工作另见 \cite{yang2022farplanner}。它的设计哲学是：\textbf{把“地形理解”与“规划”解耦}——先由地形分析把 3D 点云压成可通行性地图，再由规划器在可通行性地图上做局部与全局规划。

\begin{figure}[htbp]
  \centering
  \begin{tikzpicture}[font=\small,
      box/.style={rectangle, rounded corners=2pt, draw=structurecolor, thick,
                  fill=structurecolor!8, minimum height=7mm, minimum width=30mm, align=center}]
    \node[box] (lio) at (0,0) {LIO\\{\footnotesize FAST-LIO2 / LIO-SAM}};
    \node[box] (ta) at (0,-1.6) {terrain\_analysis\\{\footnotesize 2.5D 地形}};
    \node[box] (lp) at (0,-3.2) {local\_planner\\{\footnotesize 局部避障}};
    \node[box] (cmd) at (0,-4.8) {底盘\\{\footnotesize /cmd\_vel}};
    \node[box] (far) at (5.7,-1.6) {FAR Planner\\{\footnotesize 可视图全局规划}};
    \node[box] (tare) at (5.7,-3.2) {TARE\\{\footnotesize 分层自主探索}};
    \draw[arrow] (lio) -- node[label, right]{\scriptsize /registered\_scan} (ta);
    \draw[arrow] (ta) -- node[label, right]{\scriptsize /terrain\_map} (lp);
    \draw[arrow] (lp) -- node[label, right]{\scriptsize /cmd\_vel} (cmd);
    \draw[arrow] (ta) -- node[label, above]{\scriptsize /terrain\_map\_ext} (far);
    \draw[arrow] (far) -- node[label, right]{\scriptsize 探索目标} (tare);
    \draw[arrow, dashed] (far.west) -- node[label, below, sloped]{\scriptsize /way\_point} (lp.east);
  \end{tikzpicture}
  \caption{CMU 自主探索栈的模块数据流（简化）。LIO 输出配准后的点云，地形分析把它变成 2.5D 可通行性地图；局部规划器在局部地形图上避障并输出速度；FAR Planner 在全局可视图上选下一个目标点（\texttt{/way\_point}）；TARE 负责“没有目标时自己找目标”的探索任务。\texttt{terrain\_analysis\_ext}、\texttt{path\_follower} 与 \texttt{visualization\_tools} 是辅助模块，图中省略。}
  \label{fig:21-stack}
\end{figure}

各模块职责如下（频率为典型量级，具体以仓库配置为准）：

\begin{itemize}
  \item \texttt{terrain\_analysis}：订阅 LIO 输出的配准点云（\texttt{/registered\_scan}），维护一张以机器人为中心的 2.5D 高程栅格，计算每格的坡度、粗糙度、台阶高度与到障碍的距离，发布 \texttt{/terrain\_map}（每点携带“到最近障碍的距离”，供局部规划做净空判断）。
  \item \texttt{terrain\_analysis\_ext}：在更大范围内做“扩展地形清理”，输出 \texttt{/terrain\_map\_ext}，供 FAR Planner 构建全局可视图；范围大、分辨率可稍粗。
  \item \texttt{local\_planner}：在 \texttt{/terrain\_map} 上做局部路径规划（本质是带净空约束的图搜索 + 速度平滑），跟踪 \texttt{/way\_point} 并输出 \texttt{/cmd\_vel}，同时把碰撞检测做成“机器人尺寸 + 离地间隙”的模板。
  \item \texttt{far\_planner}：维护增量可视图（visibility graph），在已知/未知混合环境里规划到目标（或探索目标）的全局路线，输出下一个 waypoint；核心思想是 \emph{attemptable}（可尝试）：允许穿未知区域，走不通再重规划。
  \item \texttt{tare\_planner}：自主探索（exploration）规划器，没有人类给的目标，自己决定“去哪里看”，采用局部/全局两层结构：局部层在当前区域找边界（frontier），全局层决定远距离的探索顺序。
  \item \texttt{path\_follower}：跟踪一条给定路线（如示教轨迹），不负责避障。
  \item \texttt{visualization\_tools}：RViz 可视化与调试工具。
\end{itemize}

\begin{table}[htbp]
  \centering
  \caption{CMU 栈主要话题（名字以仓库为准，这里给出最常见的）}
  \label{tab:21-topics}
  \small
  \begin{tabular}{lll}
    \toprule
    话题 & 方向 & 内容 \\
    \midrule
    \texttt{/state\_estimation} & 输入 & LIO 位姿（odom/map） \\
    \texttt{/registered\_scan} & 输入 & 世界系下的配准点云 \\
    \texttt{/terrain\_map} & 中间 & 局部 2.5D 地形图 \\
    \texttt{/terrain\_map\_ext} & 中间 & 扩展地形图（给 FAR） \\
    \texttt{/way\_point} & 中间 & 局部规划的目标点 \\
    \texttt{/cmd\_vel} & 输出 & 底盘速度指令 \\
    \texttt{/exploration\_paths} & 可视化 & TARE 的探索路径 \\
    \bottomrule
  \end{tabular}
\end{table}

\section{地形分析：把点云压成 2.5D 高程图}\label{sec:21-terrain}

\begin{definition}[地形可通行性分析]
地形可通行性分析（terrain traversability analysis）是给每个空间位置打一个“机器人能否安全通过”的分数或标签的过程。常用四个几何量：\textbf{坡度}（slope，局部平面与水平面的夹角）、\textbf{粗糙度}（roughness，局部高度起伏的统计量）、\textbf{台阶高度}（step height，相邻格之间允许的高度突变）、\textbf{净空}（clearance，到最近障碍物的距离）。四者都满足阈值时判为可通行。
\end{definition}

\subsection{2.5D 栅格的组织}

地形图是一个栅格 $\{(i,j)\}$，分辨率 $r$ 通常取 $0.1$--$0.2\ \mathrm{m}$，每格记录三样东西：高度 $h_{ij}$（格内点云 $z$ 的中位数或最低值）、粗糙度 $\sigma_{ij}$（格内 $z$ 的标准差）、以及净空 $d_{ij}$（到最近“障碍格”的平面距离，用距离变换一次算出）。点云到来时，把每个点按 $(x,y)$ 投到格子，更新该格的统计量；机器人移动时用滑动窗口平移整张图（只重算新进入窗口的格子）。

\subsection{坡度：从高度差到角度}

对 4 邻接的两个格子，水平距离就是分辨率 $r$；高度差 $\Delta z=h_{i'j'}-h_{ij}$。它们连线与水平面的夹角就是坡度：
\begin{equation}
  \theta_{\mathrm{slope}}=\arctan\!\left(\frac{\abs{\Delta z}}{r}\right).
  \label{eq:21-slope}
\end{equation}
对 8 邻接的对角邻居，水平距离是 $\sqrt2\,r$，分母换成 $\sqrt{\Delta x^2+\Delta y^2}$。例：$r=0.1\ \mathrm{m}$，相邻格高差 $0.05\ \mathrm{m}$，则 $\theta=\arctan(0.5)=26.6^\circ$；若高差 $0.1\ \mathrm{m}$，则 $\theta=45^\circ$。工程上坡道阈值常取 $25^\circ$--$30^\circ$（约 $0.44$--$0.52\ \mathrm{rad}$），因为再陡四轮底盘就开始打滑。

若用格内全部点拟合平面，设平面法向为 $\bm{n}$，坡度还可以写成
\begin{equation}
  \theta_{\mathrm{slope}}=\arccos\!\left(\frac{\bm{n}^{\top}\bm{e}_z}{\norm{\bm{n}}}\right),
  \label{eq:21-slope-normal}
\end{equation}
其中 $\bm{e}_z=(0,0,1)^{\top}$。法向用 PCA 求：对格内点 $\{\bm{p}_i\}$ 先算均值 $\bar{\bm{p}}$，再算协方差
\begin{equation}
  \bm{C}=\frac{1}{N}\sum_{i=1}^{N}(\bm{p}_i-\bar{\bm{p}})(\bm{p}_i-\bar{\bm{p}})^{\top},
  \label{eq:21-cov}
\end{equation}
取最小特征值对应的特征向量即法向（这是第 7 章扫描匹配里“平面拟合”的同一套工具）。只有三个点时可以直接用叉积，例：$\bm{p}_1=(0,0,0)$、$\bm{p}_2=(1,0,0.1)$、$\bm{p}_3=(0,1,0.1)$，则
\begin{equation}
  \bm{n}=(\bm{p}_2-\bm{p}_1)\times(\bm{p}_3-\bm{p}_1)=(1,0,0.1)\times(0,1,0.1)=(-0.1,-0.1,1),
\end{equation}
归一化后 $\hat{\bm{n}}=(-0.0995,-0.0995,0.995)$，$\theta=\arccos(0.995)=5.7^\circ$，说明这三个点所在的平面几乎是平的。

\subsection{粗糙度与台阶高度}

\textbf{粗糙度}用邻域高度的标准差衡量。设某格及其 $3\times3$ 邻域共 9 个高度值，均值 $\bar h$，则
\begin{equation}
  \sigma=\sqrt{\frac{1}{N}\sum_{i=1}^{N}(h_i-\bar h)^2}.
  \label{eq:21-rough}
\end{equation}
例：5 个高度 $0.02,0.01,0.03,0.00,0.04$（单位 m），均值 $0.02$，偏差为 $0,-0.01,0.01,-0.02,0.02$，平方和 $=10^{-3}$，方差 $=10^{-3}/5=2\times10^{-4}$，标准差 $\sigma=0.0141\ \mathrm{m}$，即起伏约 1.4 cm——很平。若 $\sigma>0.03\ \mathrm{m}$，通常是碎石或草地，应降低速度或判为不可通行。

\textbf{台阶高度}定义为邻域内最大高度突变：
\begin{equation}
  h_{\mathrm{step}}=\max_{\substack{(i',j')\in\mathcal{N}(i,j)}}\abs{h_{i'j'}-h_{ij}}.
  \label{eq:21-step}
\end{equation}
路缘石 $0.15\ \mathrm{m}$、门槛 $0.2\ \mathrm{m}$ 是典型值；轮式底盘能过的台阶通常不超过 $0.1$--$0.2\ \mathrm{m}$，超过就必须绕行。注意台阶与坡度的区别：$0.15\ \mathrm{m}$ 的台阶在 $0.1\ \mathrm{m}$ 分辨率下坡度是 $\arctan(1.5)=56^\circ$，会被坡度阈值直接判死；但在 $0.2\ \mathrm{m}$ 分辨率下可能被平均成 $37^\circ$，于是需要台阶项单独把关。

\subsection{可通行性判定与伪代码}

\begin{pseudo}{2.5D 地形可通行性判定（单帧）}
  \item 对每个新点，按 $(x,y)$ 投到格子 $(i,j)$，更新 $h_{ij}$ 与 $\sigma_{ij}$
  \item 对每个格子，用式~\eqref{eq:21-slope} 或式~\eqref{eq:21-slope-normal} 算坡度 $\theta_{ij}$
  \item 用式~\eqref{eq:21-step} 算台阶高度 $h_{\mathrm{step}}$，用式~\eqref{eq:21-rough} 算粗糙度 $\sigma$
  \item 用距离变换算净空 $d_{ij}$（到最近障碍格的距离）
  \item 若 $\theta_{ij}>\theta_{\max}$ 或 $h_{\mathrm{step}}>h_{\max}$ 或 $\sigma>\sigma_{\max}$ 或 $d_{ij}<r_{\mathrm{robot}}$，标记为不可通行
  \item 发布 \texttt{/terrain\_map}：可通行点附带净空值 $d_{ij}$，供局部规划使用
\end{pseudo}

关键参数的量级（以仓库默认配置为参考，不同平台要重调）：分辨率 $r\approx0.1$--$0.2\ \mathrm{m}$；坡度阈值 $\theta_{\max}\approx25^\circ$--$30^\circ$；台阶阈值 $h_{\max}\approx0.2$--$0.3\ \mathrm{m}$（对应代码里的“相对高度阈值”类参数，如 \texttt{maxRelZ} 在 $0.2$--$0.3$ 量级）；粗糙度阈值 $\sigma_{\max}\approx0.03$--$0.05\ \mathrm{m}$；机器人半径 $r_{\mathrm{robot}}$ 按车体设。车越高、轮越大，这些阈值越宽松。

\subsection{2.5D 与真 3D 的折中}

2.5D 的优点是内存小、查询快、与 2D 规划器兼容（把“可通行性”当地图值，就能复用 A*/Dijkstra）。它的局限也很清楚：

\begin{itemize}
  \item \textbf{一个 $(x,y)$ 只能存一个高度}：桥洞、多层建筑、桌子下方会打架；工程上靠“分层高程图”或“多地图分层”补救。
  \item \textbf{悬空障碍}：树枝、管道只要在点云里，就会把下面的格子标成障碍；需要按高度带过滤（只保留 $[z_{\min},z_{\max}]$ 的点）。
  \item \textbf{动态物体}：行人、其他机器人会被压成“高地”，导致地形图抖动；需要时间衰减与一致性滤波。
\end{itemize}

\begin{figure}[htbp]
  \centering
  \begin{tikzpicture}[font=\footnotesize, scale=1.1]
    % side view
    \begin{scope}
      \draw[step=0.5cm, black!15] (0,0) grid (4.5,1.5);
      \draw[thick, black!70] (0,0.1) -- (1.0,0.1) -- (1.0,0.4) -- (2.0,0.4)
        -- (3.0,1.0) -- (4.5,1.0);
      \draw[structurecolor!70, dashed] (1.0,0.1) -- (1.0,0.4);
      \node[structurecolor!80, label] at (0.45,0.8) {台阶 $h_{\mathrm{step}}$};
      \draw[structurecolor!80, thin, ->] (0.75,0.68) -- (0.97,0.32);
      \node[green!45!black, label] at (3.55,1.42) {坡度 $\theta$};
      \draw[green!45!black, thin, ->] (3.4,1.34) -- (2.6,0.72);
      \draw[<->, black!50, thin] (2.0,1.28) -- (2.5,1.28);
      \node[black!60, label, above] at (2.25,1.28) {$r$};
      \node[label] at (2.25,-0.32) {（a）侧视：台阶、坡度与分辨率};
    \end{scope}
    % top view
    \begin{scope}[xshift=6.1cm]
      \fill[green!45!black!22] (0,0) rectangle (3.0,2.0);
      \fill[red!60!black!28] (2.0,1.0) rectangle (3.0,2.0);
      \fill[black!55] (0.5,0.5) rectangle (1.0,1.0);
      \draw[step=0.5cm, black!15] (0,0) grid (3.0,2.0);
      \node[label] at (1.5,-0.32) {（b）俯视：可通行性与障碍};
      \fill[green!45!black!22] (0,-0.95) rectangle (0.3,-0.75);
      \node[label, right] at (0.3,-0.85) {可通行};
      \fill[red!60!black!28] (1.45,-0.95) rectangle (1.75,-0.75);
      \node[label, right] at (1.75,-0.85) {不可通行};
      \fill[black!55] (2.75,-0.95) rectangle (3.05,-0.75);
      \node[label, right] at (3.05,-0.85) {障碍};
    \end{scope}
  \end{tikzpicture}
  \caption{2.5D 地形分析的两个视角。左（a）：侧视，格子分辨率 $r$、台阶高度 $h_{\mathrm{step}}$、坡度 $\theta$ 的定义；右（b）：俯视，每个格子按坡度/台阶/粗糙度/净空判为可通行（绿）、不可通行（红）或障碍（黑）。}
  \label{fig:21-terrain}
\end{figure}

\section{局部规划与 FAR Planner}\label{sec:21-local-far}

\subsection{local\_planner：在地形图上做带净空的局部规划}

CMU 栈的局部规划器与第~\ref{ch:20-local-planning}~章的 DWA/TEB 不同：它不是在速度空间采样，而是\textbf{在 2.5D 地形图上做带净空约束的图搜索 + 轨迹平滑}。核心步骤是：

\begin{enumerate}
  \item 以机器人为中心取局部窗口（如 $10\ \mathrm{m}\times10\ \mathrm{m}$），把窗口内的可通行格建成图；
  \item 给每个可通行格赋代价：离障碍越近代价越高（用 \texttt{/terrain\_map} 的净空值），坡度越大代价越高；
  \item 从当前位置向 \texttt{/way\_point} 做图搜索（Dijkstra/A*），得到一条初始路径；
  \item 沿路径做“碰撞模板检查”：把机器人的矩形/圆形足迹沿路径扫一遍，任何位置净空不足就剔除或绕行；
  \item 对路径做速度规划与平滑（限制线速度、角速度、加速度），输出 \texttt{/cmd\_vel}。
\end{enumerate}

与 2D 局部规划相比，它多了两件事：\textbf{地形代价}（坡度陡的格子即使可通行也贵）与\textbf{三维足迹}（车体高度、离地间隙决定“能不能从悬空物下方过”）。这也是为什么它需要 \texttt{/terrain\_map} 携带净空值，而不是只用 0/1 占据。

\subsection{FAR Planner：可视图 + attemptable}

\begin{definition}[可视图与 FAR Planner]
\textbf{可视图}（visibility graph）是以“关键点”为节点、以“两点之间无碰撞的直线段”为边的图，图上最短路就是多边形障碍环境中的几何最短路。\textbf{FAR Planner}（Fast, Attemptable Route Planner）\cite{yang2022farplanner} 是 CMU 提出的 3D 增量可视图规划器：节点是自由空间中的采样点，边是可通行的直线段，未知区域允许规划但标记为“可尝试”（attemptable），一旦观测到不通就删边重规划。
\end{definition}

它在三件事上与经典 A* 不同：

\begin{enumerate}
  \item \textbf{表示}：A* 在稠密栅格上扩展，FAR 在稀疏可视图上搜索。可视图节点数通常只有几百到几千，而 $100\ \mathrm{m}\times100\ \mathrm{m}\times0.1\ \mathrm{m}$ 的栅格有 $10^6$ 量级格子，所以 FAR 在大范围场景里快得多。
  \item \textbf{增量维护}：机器人边走边看，FAR 只更新机器人附近的可视图——新观测到的自由区域加点加边，新观测到的障碍删点删边。更新是局部的，不需要每次重建整张图。
  \item \textbf{attemptable}：这是 FAR 的核心思想。经典导航只走“已确认自由”的区域；但在未知环境里，如果坚持“确认后才走”，机器人只能原地等。FAR 把穿过未知区域的边标成“可尝试”，先规划、先走，走到新观测证明此路不通时，把这条边删掉、重新规划。这样既不会鲁莽撞墙（局部规划仍在实时避障），又不会因为未知而停滞。
\end{enumerate}

\begin{pseudo}{FAR Planner 的增量更新与规划}
  \item 收到新地形图 $\rightarrow$ 找出机器人附近新变为自由/障碍的格子
  \item 对自由格：若离已有节点足够远，采样为新节点，插入 KD 树
  \item 对每个新节点：与附近节点连线，做可见性检查；无碰撞则加边，穿过未知区则标为 attemptable
  \item 对障碍格：删除落在障碍里的节点，删除与之相交的边
  \item 在图上跑 Dijkstra（边权 = 长度 + 未知区域惩罚）
  \item 输出路径上下一个 waypoint 给 local\_planner；若 waypoint 不可达（走进死路），回到第 1 步重规划
\end{pseudo}

可见性检查（visibility check）是 FAR 里最耗时的操作：对候选边 $\overline{AB}$，沿线段以地形分辨率采样，检查每个采样点所在格的净空是否大于机器人半径。设边长为 $L$、采样步长 $\delta$，需要检查 $\lceil L/\delta\rceil$ 个点。优化手段是先用“包围盒 + 粗粒度占据检查”快速排除明显相交的边，只对可能可见的边做逐点检查。

\begin{figure}[htbp]
  \centering
  \begin{tikzpicture}[font=\footnotesize, scale=1.0]
    \fill[black!10] (3.0,0.2) rectangle (4.0,1.5);
    \node[black!55, label] at (3.5,0.85) {障碍};
    \fill[black!5] (5.6,0.1) rectangle (7.9,2.8);
    \node[black!45, label] at (6.7,0.55) {未知区域};
    \coordinate (s) at (0,0.3);
    \coordinate (a) at (1.3,0.6);
    \coordinate (b) at (2.3,1.6);
    \coordinate (c) at (4.3,2.4);
    \coordinate (d) at (5.0,0.95);
    \coordinate (t) at (7.2,2.35);
    \draw[structurecolor!80, thick] (s) -- (a) -- (b) -- (c);
    \draw[structurecolor!80, thick] (a) -- (d);
    \draw[red!60!black, thick, dashed] (a) -- (t);
    \draw[green!45!black, thick, dashed] (c) -- (t);
    \foreach \p in {s,a,b,c,d,t}
      \fill[structurecolor] (\p) circle (0.06);
    \node[red!60!black, label] at (5.3,1.33) {被障碍挡住};
    \node[green!45!black, label] at (5.0,3.05) {attemptable（穿未知区）};
    \node[structurecolor, label] at (2.2,2.3) {已确认自由边};
    \fill[green!45!black] (s) circle (0.08);
    \fill[red!70!black] (t) circle (0.08);
    \node[label, below left] at (s) {起点};
    \node[label, above] at (t) {目标};
  \end{tikzpicture}
  \caption{FAR Planner 的可视图。实线是已确认的可通行边，绿色虚线是穿过未知区域的 attemptable 边（先试着走，走不通再删），红色虚线是被障碍挡住、不允许的边。图是稀疏的，节点数远少于栅格数，因此能在大范围未知环境中保持实时。}
  \label{fig:21-far}
\end{figure}

\subsection{TARE：从“导航”到“探索”}

\begin{definition}[自主探索与 TARE]
\textbf{自主探索}（autonomous exploration）是指没有人工指定目标时，机器人自主选择“下一步去看哪里”，以最快速度把未知区域变成已知的过程。它的目标不是“到达某点”，而是“最大化信息增益 / 最小化剩余未知体积”。\textbf{TARE}（Tiered Autonomous Robot Exploration）\cite{cao2022tare} 是 CMU 栈里的探索规划器，用“局部 + 全局”两层结构解决大范围探索中的短视问题。
\end{definition}

探索与导航的区别可以先从“目标从哪来”说起：导航的目标是人给的，探索的目标是算法自己算的。TARE 的目标选择分两层：

\begin{itemize}
  \item \textbf{局部层}：在当前可视图上，从机器人附近的可视节点里找“能看到最多未知区域”的位置，计算探索增益（gain）——通常用“从该点射线投射能看到的未知体素数”近似；再除以路径代价（gain/cost），选比值最大的作为下一个 waypoint。增益/代价的比值让机器人优先去“近且能看很多”的地方。
  \item \textbf{全局层}：当局部区域没有可观的增益（当前房间看完了），再从一张低分辨率的全局地图上找远处的边界（frontier，已知与未知的交界），选一个全局目标，交给 FAR Planner 规划过去。这样避免局部层在“看完的房间里反复打转”，也避免一上来就冲向很远的未知区域。
\end{itemize}

\begin{pseudo}{TARE 式分层探索主循环}
  \item 更新地图与可视图；若没有剩余未知区域，结束
  \item 计算局部候选节点的探索增益 $g_i$ 与路径代价 $c_i$
  \item 若存在 $g_i/c_i$ 超过阈值：选最优节点，令 FAR Planner 规划过去，执行
  \item 否则：在全局低分辨率地图上找边界，选增益最大的边界作为全局目标
  \item 若全局也没有边界：探索完成；否则规划到全局目标后回到第 1 步
\end{pseudo}

“探索”这个词在 CMU 栈里还包含一层含义：\textbf{探索本身就是一个不断试错的过程}。TARE 依赖 FAR 的 attemptable 边去“试探”未知区域，走进去看到更多，再更新增益；这正是 FAR 与 TARE 配套使用的原因——一个负责“敢走”，一个负责“去哪”。

\section{部署、调参与和 LIO 的配合}\label{sec:21-deploy}

\subsection{与 LIO 的分工}

CMU 栈不自己估计位姿，它假设你已经在跑一个 LIO（仓库默认配 FAST-LIO2\cite{xu2022fastlio2} 或 LIO-SAM），并且 LIO 发布：世界系（或 odom 系）下的配准点云 \texttt{/registered\_scan} 与位姿 \texttt{/state\_estimation}。地形分析直接消费配准点云，所以：

\begin{itemize}
  \item \textbf{LIO 的漂移会变成地形图的“重影”}：同一条墙被扫成两层，坡度与净空都会算错，进而导致“幽灵障碍”。室外大范围场景要开回环（LIO-SAM）或用带全局优化的方案。
  \item \textbf{点云去畸变必须做好}：快速转向时点云拖尾会让坡度噪声变大，地形图抖动，局部规划误刹车。
  \item \textbf{外参要准}：雷达相对底盘的高度与俯仰角直接进入地形图；装歪 $2^\circ$，$10\ \mathrm{m}$ 外的高度就偏 $0.35\ \mathrm{m}$，足以把平地判成坡。
\end{itemize}

\subsection{关键参数与常见问题}

\begin{lstlisting}[language=bash]
# Typical workflow (ROS 1 / ROS 2 packages share the same idea)
roslaunch vehicle_simulator system_garage.launch   # start simulator + stack
roslaunch terrain_analysis terrain_analysis.launch # 2.5D terrain map
roslaunch local_planner local_planner.launch       # local planner
roslaunch far_planner far_planner.launch           # visibility graph
roslaunch tare_planner tare_planner.launch         # exploration
# Key parameters (names vary slightly across versions):
#   vehicle_height / sensor_height  -> footprint and clearance
#   terrain grid resolution         -> 0.1-0.2 m
#   slope / step / roughness limits -> traversability thresholds
#   maxRelZ                         -> relative height tolerance (m)
\end{lstlisting}

常见问题与排查方向：\textbf{地形图抖动}（点云噪声或外参抖动，先查 LIO 与雷达安装，再加大地形滤波窗口）；\textbf{局部规划卡住}（净空阈值太保守或车体尺寸填大，先看 \texttt{/terrain\_map} 里机器人周围是否被误判为窄）；\textbf{FAR 不更新}（\texttt{/terrain\_map\_ext} 没接上，或可视图更新半径小于移动速度$\times$更新周期）；\textbf{探索原地转}（增益阈值过高或全局层没触发，检查边界提取）。这些问题与第~\ref{ch:27-tuning}~章的排查表是同一套方法：\emph{先看数据流有没有，再看数值对不对，最后才怀疑算法}。

\begin{summary}

\begin{enumerate}
  \item 2D 导航把世界压成平面，丢掉了高度、垂直净空与地形质量；2.5D 高程图用“每格一个高度”在内存与表达力之间取得平衡，是地面 3D 导航的主流选择。
  \item CMU 栈的分工是“地形理解与规划解耦”：LIO 给配准点云，terrain\_analysis 出 2.5D 可通行性地图，local\_planner 做带净空的局部规划，FAR Planner 做增量可视图全局规划，TARE 做分层自主探索。
  \item 地形可通行性由坡度（式~\eqref{eq:21-slope}）、粗糙度（式~\eqref{eq:21-rough}）、台阶高度（式~\eqref{eq:21-step}）与净空共同判定；分辨率会“平均掉”小台阶，所以台阶项要单独把关。
  \item FAR 用稀疏可视图替代稠密栅格，用增量更新维持实时，用 attemptable 边“先走未知、走不通再删”，适合大范围未知环境。
  \item TARE 用局部/全局两层解决探索的短视问题：局部比 gain/cost，全局找 frontier；探索的目标是信息增益而不是到达某点。
  \item 这套栈依赖 LIO 的位姿与配准点云；漂移、去畸变、外参都会直接变成地形图误差，进而影响规划。
\end{enumerate}


\end{summary}

\begin{problemset}

  \item 用式~\eqref{eq:21-slope} 计算：$r=0.15\ \mathrm{m}$、相邻格高差 $0.06\ \mathrm{m}$ 时的坡度；若阈值是 $25^\circ$，该格可通行吗？
  \item 为什么在 $0.2\ \mathrm{m}$ 分辨率下台阶会被“平均”掉？请构造一个 $0.15\ \mathrm{m}$ 台阶的例子，说明坡度项与台阶项判定的差异。
  \item 2.5D 地图遇到“桥洞”会怎么错？给出两种工程上的补救方案。
  \item FAR 的 attemptable 边与 A* 在未知栅格上“把未知当自由”有什么区别？为什么后者更危险？
  \item TARE 的增益/代价比值中，如果只用增益、不看代价，机器人会表现成什么样？
  \item 如果 LIO 出现约 $1^\circ$ 的俯仰角误差，$10\ \mathrm{m}$ 外的高度误差是多少？这会怎样影响地形分析？
\end{problemset}

\begin{reading}

\begin{itemize}
  \item \cite{cao2022tare}：CMU 自主探索开发环境的系统论文，读它建立“地形分析—局部规划—FAR—TARE”的整体框架。
  \item \cite{yang2022farplanner}：FAR Planner 论文，重点看可视图的增量更新与 attemptable 策略的定义。
  \item \cite{oleynikova2017voxblox}：ESDF 的增量构建，理解“距离场”之后再看 2.5D 净空会更容易。
  \item \cite{xu2022fastlio2}：CMU 栈常用的 LIO 前端，理解它输出的配准点云是地形分析的输入。
  \item 代码仓库 \href{https://github.com/caochao39/autonomous_exploration_development_environment}{caochao39/autonomous\_exploration\_development\_environment}：有 ROS 1/ROS 2 分支与仿真场景，按 README 跑一遍仿真比读十篇论文更直观。
\end{itemize}


\end{reading}

\chapter{2D 建图 + 导航完整链路实战}\label{ch:22-nav2-practice}

\usetikzlibrary{arrows.meta,positioning,calc,fit,patterns,decorations.pathreplacing}

\chapterintro{前面各章把导航拆成定位、建图、全局规划、局部规划分别讲完了。但真正让机器人跑起来，难点往往不在单个算法，而在“把它们按正确的顺序、正确的坐标系、正确的频率接起来”。本章就是一次完整的工程串联：先用 Cartographer 建一张 2D 栅格地图并存成 pgm/yaml，再用 AMCL 在这张图上定位，最后用 Nav2 做规划与控制。我们会给出可用的配置骨架、仿真验证步骤，以及一张“现象 $\rightarrow$ 原因 $\rightarrow$ 解决”的排查表。学完本章，你应该能从零把一台差速底盘在仿真里跑通，并知道每个环节出错时该看哪里。}

\section{全链路总览}\label{sec:22-pipeline}

整条链路分两个阶段。\textbf{建图阶段}：跑 Cartographer\cite{hess2016cartographer}，订阅激光 \texttt{/scan} 与轮式里程计 \texttt{/odom}，边建图边发布 \texttt{map}$\to$\texttt{odom} 变换，最终导出 pgm 图片与 yaml 描述文件。\textbf{导航阶段}：用 \texttt{map\_server} 加载刚才的地图作为静态层，用 AMCL\cite{fox2003amcl} 在图上做粒子滤波定位并发布 \texttt{map}$\to$\texttt{odom}，Nav2\cite{macenski2020nav2} 的行为树调度全局规划与控制，输出 \texttt{/cmd\_vel} 给底盘。

\begin{figure}[htbp]
  \centering
  \begin{tikzpicture}[font=\footnotesize,
      box/.style={rectangle, rounded corners=2pt, draw=structurecolor, thick,
                  fill=structurecolor!8, minimum height=7mm, minimum width=26mm, align=center}]
    % mapping stage (left column)
    \node[box] (scan) at (0,0) {激光 \texttt{/scan}};
    \node[box] (carto) at (0,-1.3) {Cartographer\\{\scriptsize 建图}};
    \node[box] (mapfile) at (0,-2.6) {map.pgm\\{\scriptsize + map.yaml}};
    % navigation stage (right column)
    \node[box] (server) at (6.6,0) {\texttt{map\_server}};
    \node[box] (amcl) at (6.6,-1.3) {AMCL\\{\scriptsize 定位}};
    \node[box] (nav) at (6.6,-2.6) {Nav2\\{\scriptsize 规划 + 控制}};
    \node[box] (base) at (6.6,-3.9) {底盘\\{\scriptsize \texttt{/cmd\_vel}}};
    \draw[arrow] (scan) -- (carto);
    \draw[arrow] (carto) -- node[label, right]{保存} (mapfile);
    \draw[arrow] (server) -- (amcl);
    \draw[arrow] (amcl) -- node[label, right]{位姿} (nav);
    \draw[arrow] (nav) -- (base);
    \draw[arrow, dashed] (mapfile.east) -- node[label, above, sloped]{map.yaml} (server.west);
    \node[structurecolor, label] at (0,0.75) {\textbf{建图阶段}};
    \node[structurecolor, label] at (6.6,0.75) {\textbf{导航阶段}};
  \end{tikzpicture}
  \caption{2D 建图 + 导航全链路。左列是建图阶段：Cartographer 把激光与里程计变成地图文件；右列是导航阶段：map\_server 提供先验地图，AMCL 提供定位，Nav2 完成规划与控制。两个阶段共用同一张地图与同一套坐标系约定；导航阶段 AMCL 同样订阅 \texttt{/scan}，图中为清晰起见省略。}
  \label{fig:22-pipeline}
\end{figure}

理解这条链路要抓住一个关键点：\textbf{建图与导航是两个“换人不换地图”的进程}。建图时 \texttt{map}$\to$\texttt{odom} 由 Cartographer 发布；切到导航后，必须把 Cartographer 关掉，改由 AMCL 发布 \texttt{map}$\to$\texttt{odom}。如果两个节点同时发布这一段 TF，tf2 会认为树里有多个父节点，整棵 TF 树失效——这是新手最常踩的坑之一。

\section{Cartographer 建图实操}\label{sec:22-carto}

\subsection{配置文件}

Cartographer 用 Lua 配置，2D 建图的入口参数通常从一个 \texttt{backpack\_2d.lua} 派生。下面是一份精简的配置骨架（参数值按 $0.05\ \mathrm{m}$ 分辨率、室内场景给出，可按场地调整）：

\begin{lstlisting}[language=bash]
-- my_2d.lua : minimal Cartographer 2D configuration
include "backpack_2d.lua"

TRAJECTORY_BUILDER_2D.min_range = 0.3
TRAJECTORY_BUILDER_2D.max_range = 30.0
TRAJECTORY_BUILDER_2D.missing_data_ray_length = 5.0
TRAJECTORY_BUILDER_2D.use_imu_data = false
TRAJECTORY_BUILDER_2D.num_accumulated_range_data = 1
TRAJECTORY_BUILDER_2D.voxel_filter_size = 0.025
TRAJECTORY_BUILDER_2D.submaps.num_range_data = 90
TRAJECTORY_BUILDER_2D.submaps.grid_options_2d.resolution = 0.05

MAP_BUILDER.use_trajectory_builder_2d = true
POSE_GRAPH.optimize_every_n_nodes = 90
POSE_GRAPH.constraint_builder.min_score = 0.65
POSE_GRAPH.constraint_builder.global_localization_min_score = 0.7
\end{lstlisting}

\begin{sloppypar}
几个参数值得解释：\texttt{min\_range} 与 \texttt{max\_range} 是激光可用距离，设错会把地面或远处噪点当障碍；\texttt{voxel\_filter\_size} 是体素降采样尺寸，越小越精细但越耗 CPU；\texttt{submaps.num\_range\_data}=90 表示每 90 帧激光构成一个子图，太大则回环约束少、建图容易漂；\texttt{resolution}=0.05 是最终栅格分辨率，$10\ \mathrm{m}\times10\ \mathrm{m}$ 的场地对应 $200\times200$ 的栅格，正好与第~\ref{ch:18-nav-overview}~章的代价地图分辨率一致。若雷达有 IMU 且已标定，把 \texttt{use\_imu\_data} 打开能显著改善快速转向时的建图质量。
\end{sloppypar}

\subsection{建图、结束与导出}

启动建图后，遥控机器人把场地走一遍（尽量走小闭环、每个角落都扫到），然后结束轨迹并导出：

\begin{lstlisting}[language=bash]
# ROS 2: start mapping (ROS 1 uses roslaunch cartographer_ros ...)
ros2 launch cartographer_ros backpack_2d.launch.py

# After finishing the tour: close trajectory 0, then write state
ros2 service call /finish_trajectory cartographer_ros_msgs/srv/FinishTrajectory \
  "{trajectory_id: 0}"
ros2 service call /write_state cartographer_ros_msgs/srv/WriteState \
  "{filename: '${HOME}/maps/map.pbstream'}"

# Convert the pbstream into pgm + yaml for map_server
ros2 run cartographer_ros cartographer_pbstream_to_ros_map \
  -map_filestem=${HOME}/maps/map -pbstream_filename=${HOME}/maps/map.pbstream \
  -resolution=0.05
\end{lstlisting}

导出后得到 \texttt{map.pgm}（灰度图：白=自由、黑=占据、灰=未知）与 \texttt{map.yaml}（分辨率、原点、阈值）。\texttt{map.yaml} 里的 \texttt{origin} 是地图左下角在世界系中的位姿，\texttt{resolution} 必须与实际一致，否则导航时障碍会整体缩放或偏移——这类错误表现为“地图看起来对，但机器人总是撞同一面墙”。

\subsection{建图质量自检}

在建图现场就应该检查四件事，而不是等导航失败再回头：\textbf{墙是否直}（走一圈回到起点，墙若重影说明回环没生效）；\textbf{角落是否闭合}（房间四角应是直角，若歪说明里程计外参或时间同步有问题）；\textbf{动态物体是否被抹掉}（人走过留下的“鬼影”应被后续观测覆盖）；\textbf{未知区域是否太多}（大片灰色意味着没扫到，导航时会被当成障碍或直接判无路）。第~\ref{ch:15-cartographer}~章讲的子图与回环机制，正是这些现象背后的原因。

\section{Nav2 配置}\label{sec:22-nav2}

\subsection{参数文件骨架}

Nav2 的参数集中在 \texttt{nav2\_params.yaml}，按节点分段。下面是一份能跑起来的最小骨架（只保留关键项，实际使用请复制官方完整文件再改）：

\begin{lstlisting}[language=Python]
amcl:                       # localization on the saved map
  ros__parameters:
    use_sim_time: True
    min_particles: 500
    max_particles: 2000
    laser_model_type: likelihood_field
bt_navigator:               # behavior-tree scheduler
  ros__parameters:
    global_frame: map
    robot_base_frame: base_link
planner_server:             # global planning
  ros__parameters:
    expected_planner_frequency: 1.0
    planner_plugins: ["GridBased"]
    GridBased:
      plugin: "nav2_navfn_planner/NavfnPlanner"
      tolerance: 0.5
      use_astar: true
controller_server:          # local planning / tracking
  ros__parameters:
    controller_frequency: 20.0
    controller_plugins: ["FollowPath"]
    FollowPath:
      plugin: "dwb_core::DWBLocalPlanner"
      max_vel_x: 0.26
      max_vel_theta: 1.0
\end{lstlisting}

要点：\texttt{amcl} 的粒子数与激光模型决定定位稳定性，室内小场地 500--2000 个粒子够用；\texttt{planner\_server} 的 \texttt{tolerance} 表示“离目标多近算到达”，设太小会在终点附近反复规划；控制频率取 \texttt{controller\_frequency}=20.0，对应第~\ref{ch:20-local-planning}~章说的实时要求；\texttt{max\_vel\_x}=0.26 是 TurtleBot3 的典型值，实车按自己底盘改。

\subsection{地图、定位与生命周期}

导航启动时，\texttt{map\_server} 读 \texttt{map.yaml} 与 \texttt{map.pgm}，把它们发布成 \texttt{/map} 话题与静态层；AMCL 订阅 \texttt{/scan}、\texttt{/map}，发布 \texttt{map}$\to$\texttt{odom} 与 \texttt{/amcl\_pose}。\textbf{AMCL 需要一个初始位姿}：在 RViz 里用“2D Pose Estimate”点一下，或者在启动参数里给初始位姿；不给的话粒子会撒满整张地图，表现为“定位一直跳、规划路径乱飞”。

启动顺序由 \texttt{lifecycle\_manager} 保证：每个 Nav2 节点都是生命周期节点，状态依次是“未配置 $\rightarrow$ 未激活 $\rightarrow$ 激活”；管理器按依赖顺序调用 \texttt{configure} 与 \texttt{activate}，避免“控制器先于代价地图就绪”这类竞争条件。这就是第~\ref{ch:18-nav-overview}~章说“Nav2 不是一个程序，而是一组被编排的节点”的工程含义。

\begin{figure}[htbp]
  \centering
  \begin{tikzpicture}[font=\small,
      box/.style={rectangle, rounded corners=2pt, draw=structurecolor, thick,
                  fill=structurecolor!8, minimum height=8mm, minimum width=32mm, align=center}]
    \node[box] (map) at (0,2.4) {\texttt{map}$\to$\texttt{odom}\\{\scriptsize AMCL 发布，可跳变}};
    \node[box] (odom) at (0,1.2) {\texttt{odom}$\to$\texttt{base\_link}\\{\scriptsize 里程计发布，连续}};
    \node[box] (base) at (0,0) {\texttt{base\_link}$\to$\texttt{laser}\\{\scriptsize 静态外参，URDF}};
    \draw[arrow] (map) -- (odom);
    \draw[arrow] (odom) -- (base);
    \draw[decorate, decoration={brace, amplitude=5pt}, thick, structurecolor]
      (-2.35,2.9) -- (-2.35,-0.5);
    \node[structurecolor, label, rotate=90, above] at (-2.7,1.2) {导航阶段必须全部连通};
  \end{tikzpicture}
  \caption{导航阶段的 TF 链。三段变换分别由 AMCL、里程计、URDF 发布：\texttt{map}$\to$\texttt{odom} 可跳变（重定位时瞬间修正），\texttt{odom}$\to$\texttt{base\_link} 必须连续（否则局部控制抖动），传感器外参静态。任何一段断开或重复发布，都会让整条链路不可用。}
  \label{fig:22-tf}
\end{figure}

\subsection{仿真验证}

配置好之后，先在仿真里验证，再上实车。TurtleBot3 是官方示例中最省事的一套：

\begin{lstlisting}[language=bash]
# Terminal 1: simulation world
export TURTLEBOT3_MODEL=burger
ros2 launch turtlebot3_gazebo turtlebot3_world.launch.py

# Terminal 2: navigation (map + AMCL + Nav2 + RViz)
ros2 launch nav2_bringup tb3_simulation_launch.py \
  map:=${HOME}/maps/map.yaml use_sim_time:=true

# Then in RViz: set "2D Pose Estimate", send a "2D Goal Pose"
\end{lstlisting}

仿真验证的判据是三个“能”：能在地图上定位（$\texttt{/amcl\_pose}$ 稳定、粒子云收敛）、能规划（全局路径在 RViz 里出现且不穿墙）、能跟踪（机器人沿路径走、遇到新障碍能绕）。三个都过，再上实车；实车上还要额外验证急停、通信丢包与电池电压低时的行为。

\section{常见问题排查表}\label{sec:22-debug}

\begin{table}[htbp]
  \centering
  \caption{全链路常见故障：现象 $\rightarrow$ 原因 $\rightarrow$ 解决}
  \label{tab:22-debug}
  \small
  \begin{tabular}{p{3.2cm}p{4.6cm}p{4.8cm}}
    \toprule
    现象 & 可能原因 & 解决 \\
    \midrule
    地图漂移、墙重影 & 回环未触发、里程计标定差、时间不同步 & 建图走小闭环、查轮径/轮距与时间戳 \\
    定位跳变、粒子散开 & 未给初始位姿、激光与地图不匹配、场景变化大 & 给初始位姿、检查外参、必要时重定位 \\
    规划失败（无路径） & 目标在膨胀区/未知区、\texttt{tolerance} 太小 & 调小 \texttt{inflation\_radius}、给 1--2 格容差 \\
    机器人原地转 & 全局路径不可达或控制器无可行轨迹 & 看 \texttt{/local\_costmap} 是否被堵、检查恢复行为 \\
    机器人不动 & TF 断链、\texttt{/cmd\_vel} 未下发、生命周期未激活 & \texttt{tf2\_echo}、\texttt{ros2 topic hz}、看日志 \\
    代价地图不更新 & 传感器话题没接、\texttt{update\_frequency}=0 & 查 \texttt{/scan} 频率、查 layer 配置 \\
    \bottomrule
  \end{tabular}
\end{table}

表~\ref{tab:22-debug} 的排查顺序建议固定成一套流程：先确认话题在不在（\texttt{ros2 topic list} 与 \texttt{ros2 topic hz}），再确认 TF 通不通（\texttt{tf2\_echo}），然后确认数值对不对（在 RViz 里看地图、粒子与代价地图），最后才怀疑算法。绝大多数“导航不工作”都能在前两步定位，因为算法本身很少突然坏掉，坏掉的通常是数据流。

\section{性能、安全与恢复行为}\label{sec:22-safety}

\textbf{性能}：20 Hz 控制频率下，主要 CPU 开销在代价地图更新（5 Hz）与局部规划（每周期几十条轨迹）；全局规划 1 Hz，开销小。若 CPU 超载，先降 \texttt{costmap} 的更新频率与分辨率，再降局部规划采样数，最后才降控制频率（控制频率一降，避障反应就慢）。

\textbf{安全}：Nav2 的 \texttt{behavior\_server} 提供三类恢复行为——\texttt{spin}（原地旋转重新观测）、\texttt{backup}（后退脱困）、\texttt{wait}（等待障碍消失）。行为树里通常写成“规划失败 $\rightarrow$ 清空代价地图 $\rightarrow$ 旋转 $\rightarrow$ 再规划 $\rightarrow$ 后退”，层层兜底。但恢复行为不是万能的：如果定位已经错了，旋转只会让错误更大，所以\textbf{恢复行为之前必须确认定位健康}。

\textbf{急停}：工程上必须有独立于 Nav2 的急停链路（物理按钮或 \texttt{/cmd\_vel} 上的看门狗），不能只依赖行为树。看门狗的逻辑很简单：超过 $T$ 毫秒没收到新指令就下发零速度；$T$ 一般取 200--500 ms，取太小会被正常抖动误触发，取太大则失去保护意义。

\textbf{速度平滑}：\texttt{velocity\_smoother} 对 \texttt{/cmd\_vel} 做限幅与一阶低通，保证加速度不超限；调它的时候要注意，平滑太强会让避障“变肉”，反应慢半拍。

\begin{summary}

\begin{enumerate}
  \item 全链路是“建图阶段”与“导航阶段”两段：Cartographer 出 pgm/yaml，map\_server 加载，AMCL 定位，Nav2 规划控制；两段共用一张地图与一套坐标系。
  \item 建图阶段结束要 \texttt{finish\_trajectory} 并 \texttt{write\_state} 导出 pbstream，再转成 pgm/yaml；\texttt{map.yaml} 的 \texttt{origin} 与 \texttt{resolution} 必须准确。
  \item 导航阶段必须把 \texttt{map}$\to$\texttt{odom} 的发布权从 Cartographer 交给 AMCL，禁止两个节点同时发布；AMCL 必须给初始位姿。
  \item Nav2 参数按节点分段：AMCL 粒子数、planner 的频率与容差、controller 的 20 Hz 与速度上限是最常调的几项；生命周期管理器保证启动顺序。
  \item 排查固定顺序：先话题、再 TF、再数值、最后算法；表~\ref{tab:22-debug} 覆盖了最常见的六类故障。
  \item 安全靠三层：行为树的恢复行为、独立急停/看门狗、速度平滑；恢复行为之前要先确认定位健康。
\end{enumerate}


\end{summary}

\begin{problemset}

  \item 如果建图时忘了 \texttt{finish\_trajectory} 直接 \texttt{write\_state}，导出的地图可能有什么问题？
  \item AMCL 不给初始位姿时，粒子会撒在哪里？为什么表现为“路径乱飞”而不是“机器人不动”？
  \item 一个 $10\ \mathrm{m}\times10\ \mathrm{m}$ 的场地，分辨率 $0.05\ \mathrm{m}$，地图 pgm 至少多大？若分辨率改成 $0.1\ \mathrm{m}$，建图与导航各有什么变化？
  \item 机器人“原地转”与“不动”分别优先检查什么？请按本章的排查顺序各写三步。
  \item 为什么急停不能只靠 Nav2 的恢复行为？给出一个恢复行为反而让情况变糟的具体场景。
  \item 把本章的链路换成“用 FAST-LIO2 建 3D 图 + 导出 2D 栅格”的方案，需要额外做哪几步？
\end{problemset}

\begin{reading}

\begin{itemize}
  \item \cite{hess2016cartographer}：Cartographer 原始论文，理解子图、回环与位姿图之后再调 \texttt{backpack\_2d.lua} 会有的放矢。
  \item \cite{fox2003amcl}：AMCL 的自适应粒子数（KLD-Sampling），理解“为什么粒子数会自适应”有助于调定位。
  \item \cite{macenski2020nav2}：Nav2 设计论文，配置项与节点关系以它和官方文档为准。
  \item Nav2 官方文档 \href{https://docs.nav2.org}{docs.nav2.org} 与 Cartographer ROS 文档 \href{https://google-cartographer-ros.readthedocs.io}{google-cartographer-ros.readthedocs.io}：参数默认值与 launch 示例。
\end{itemize}


\end{reading}

\chapter{Z-Localization 深度解析}\label{ch:23-zlocalization}

\chapterintro{这一章解剖一套真实的、开源的、为比赛而生的定位系统：东北大学 T-DT 战队 2026 赛季开源的 Z-Localization。它不是又一篇论文的复现，而是一次工程取舍的完整展示——为什么把 LIO 和先验地图定位拆成前后端？为什么状态机要设四个状态？为什么地图里要存“点矩”而不是点？我们按“架构 → 前端 → 后端 → 使用流程 → 可借鉴之处”的顺序拆解，并把它与 FAST-LIO2、Point-LIO 放在一起对比。读完你应该能判断：什么场景该用它，什么场景该自己改。}

\section{背景与定位}\label{sec:23-background}

Z-Localization 是东北大学（Northeastern University）T-DT 实验室在 2026 赛季 RoboMaster 超级对抗赛中使用的哨兵定位方案，赛后以 GPL-2.0 协议开源\cite{zloc2026}。它的官方描述是“一套面向移动机器人的 ROS 2 激光惯性定位栈，支持里程计、点云建图、先验地图定位、回环与重定位，也可接入 UWB 提供位置辅助”。

它解决的核心问题，是比赛场景下的三个现实约束：
\begin{enumerate}
  \item \textbf{场地是动态的}：有对手、有弹丸、有观众，地图里不能有动态残影；
  \item \textbf{先验地图不一定可用}：适应性训练阶段可能来不及建图，需要“无图也能跑”；
  \item \textbf{定位必须可恢复}：云台转动、撞击、遮挡都可能导致里程计飞掉，系统要能自己找回来。
\end{enumerate}

这些约束决定了它的架构与 FAST-LIO2 那种“一个紧耦合滤波器打天下”的方案不同——它把系统拆成了\textbf{前端 ZLIO} 与\textbf{后端定位}两部分。这个拆分本身就是本章最值得学习的设计。

\section{总体架构：前后端分离}\label{sec:23-arch}

\subsection{为什么拆}

作者在设计说明中给了一个很实在的理由：直接把先验地图塞进 LIO，LIO 是紧耦合的，一旦发散很难判断——“有时候位置飞了，或者速度、加速度已经比较离谱了，才能发现里程计爆炸了”。

拆开之后，职责变得清晰：
\begin{itemize}
  \item \textbf{前端（ZLIO）}：只负责在 \textbf{odom 系}里估计运动。它维护自己的局部地图，输出高频里程计、去畸变点云、以及每一帧的位姿。
  \item \textbf{后端}：只负责维护 $\text{map}\to\text{odom}$ 的变换。它从前方送来关键帧，把关键帧的全局位姿放进因子图优化，用先验地图配准、回环、UWB 提供全局约束。
\end{itemize}

这样做有三个好处：
\begin{enumerate}
  \item \textbf{可观测性分析}：先验图配准的海森矩阵可以单独做可观测性分析，判断点云匹配是否退化、定位是否可信；
  \item \textbf{减少高频大图访问}：前端不必在每次高频更新时访问巨大的先验地图，理论上更省算力；
  \item \textbf{可恢复}：前端可以独立重启并重建 odom 系，全局位姿的恢复交给后端单独处理，不会“一崩全崩”。
\end{enumerate}

\begin{figure}[htbp]
  \centering
  \begin{tikzpicture}[node distance=0.9cm]
    \node[block, minimum width=22mm] (lidar) {LiDAR};
    \node[block, right=0.7cm of lidar, minimum width=22mm] (imu) {IMU};
    \node[block2, right=1.2cm of imu, minimum width=30mm, minimum height=16mm] (front) {前端 ZLIO\\局部里程计 + 局部地图};
    \node[block2, right=1.6cm of front, minimum width=30mm, minimum height=16mm] (back) {后端\\因子图 + $\text{map}\to\text{odom}$};
    \node[block, below=1.0cm of back, minimum width=26mm] (prior) {先验地图/描述子};
    \node[block, below=1.0cm of front, minimum width=26mm] (uwb) {UWB};
    \node[block2, below=1.2cm of lidar, minimum width=30mm] (out) {输出：位姿、速度、协方差、状态};
    \draw[arrow] (lidar) -- (front);
    \draw[arrow] (imu) -- (front);
    \draw[arrow] (front) -- node[above, label]{关键帧 + 点云} (back);
    \draw[arrow] (prior) -- (back);
    \draw[arrow] (uwb) -- (back);
    \draw[arrow] (back) -- node[right, label]{TF 校正} (out);
    \draw[arrow] (front) -- (out);
  \end{tikzpicture}
  \caption{Z-Localization 的前后端架构。前端在 odom 系估计局部运动，后端维护 map 与 odom 之间的变换并输出全局位姿。}
  \label{fig:23-arch}
\end{figure}

\subsection{与 SLAM 前后端的区别}

注意这里的“前后端”与第 11 章 SLAM 的前后端不完全一样：
\begin{itemize}
  \item SLAM 的前端 = 里程计 + 数据关联；后端 = 回环 + 全局优化。二者共同维护\textbf{一张地图}。
  \item Z-Localization 的前端 = LIO；后端 = 全局定位（图 23 章\ref{fig:23-arch}）。前端维护\textbf{局部地图}，后端维护\textbf{全局一致性}，二者通过关键帧观测连接。
\end{itemize}
这是“定位系统语境下的前端/后端”——正是第 1 章表 1 中的第二种划分。

\section{前端 ZLIO}\label{sec:23-frontend}

前端要解决四件事：误差状态怎么建、怎么启动、怎么去畸变、地图怎么存。

\subsection{误差状态建模}

常见滤波式 LIO 沿用 ESKF 的“直积流形”建模：旋转用乘法更新，位置、速度、零偏用加法更新。ZLIO 采用 \textbf{$\mathrm{SE}_2(3)\ltimes\mathfrak{se}(3)$ 半直积结构}，把导航状态与 IMU 零偏一起纳入误差变换，目的是让姿态和速度部分的无噪声误差得到精确的线性形式，减少线性化误差。

其误差方程（引自项目设计说明）为
\begin{align}
\dot{\bm{\rho}}_p &\simeq \bm{\rho}_v-\hat{\bm{p}}\times\bm{\beta}_g+\hat{\bm{p}}\times(\hat{\bm{R}}\bm{n}_\omega),\\
\dot{\bm{\theta}} &\simeq -\bm{\beta}_g+\hat{\bm{R}}\bm{n}_\omega,\\
\dot{\bm{\rho}}_v &\simeq \bm{g}\times\bm{\theta}-\bm{\beta}_a+\hat{\bm{v}}\times(\hat{\bm{R}}\bm{n}_\omega)+\hat{\bm{R}}\bm{n}_a,\\
\dot{\bm{\beta}}_g &\simeq \bigl(\hat{\bm{R}}(\bm{\omega}_m-\hat{\bm{b}}_g)\bigr)\times\bm{\beta}_g-\hat{\bm{R}}\bm{n}_{bg},\\
\dot{\bm{\beta}}_a &\simeq \bigl(\hat{\bm{R}}(\bm{a}_m-\hat{\bm{b}}_a)+\bm{g}+\hat{\bm{v}}\times\hat{\bm{R}}(\bm{\omega}_m-\hat{\bm{b}}_g)\bigr)\times\bm{\beta}_g \nonumber\\
&\quad +\bigl(\hat{\bm{R}}(\bm{\omega}_m-\hat{\bm{b}}_g)\bigr)\times\bm{\beta}_a-\hat{\bm{v}}\times(\hat{\bm{R}}\bm{n}_{bg})-\hat{\bm{R}}\bm{n}_{ba}.
\label{eq:23-err}
\end{align}
这里 $\bm{\rho}_p,\bm{\rho}_v,\bm{\theta}$ 分别对应位置、速度、姿态的误差分量，$\bm{\beta}_g,\bm{\beta}_a$ 对应零偏误差。这个模型的推导细节可参考文献\cite{equi2023}。对读者而言，重要的是理解它的动机：\textbf{标准 ESKF 在姿态误差较大时线性化误差明显，而把状态放在更“等变”的群结构上，可以让误差方程在更大范围内保持线性}。

\subsection{不要求静止的初始化}

传统 IMU 初始化要求机器人静止几秒，用加速度均值估计重力方向。但比赛里“定位飞掉重启时，云台可能还在跟随自瞄转动，几乎不能保证静止”。ZLIO 的初始化流程是：
\begin{enumerate}
  \item 先用一小段 IMU 数据给出重力方向初值；
  \item 通过 IMU 传播和点云粗配准得到初步轨迹；
  \item 积累一段数据后，把\textbf{多帧点云的共面约束}与\textbf{IMU 预积分}放在同一个窗口里联合优化，一起估计位姿、速度、IMU 零偏和重力方向；
  \item 每轮优化后，用更新后的轨迹重新去畸变，再继续优化（类似迭代最近点 + 优化的交替）；
  \item 收敛后，把坐标系对齐到重力方向，用窗口末帧状态建立滤波初值与初始地图，进入正常里程计更新。
\end{enumerate}

这套流程的价值在于：\textbf{把“初始化”从“静止假设”中解放出来}。对任何需要在运动中启动的机器人，这都是可以直接借鉴的思路。

\subsection{去畸变}

ZLIO 的去畸变采用 DLIO\cite{chen2023dlio} 的思路：用 IMU 构造\textbf{分段连续轨迹}，按每个点的采样时刻做运动补偿。具体做法是：
\begin{enumerate}
  \item 把加速度转到 odom 系并加入重力，得到实际运动加速度；
  \item 根据相邻 IMU 数据估计角加速度和加速度变化率；
  \item 构造这段时间内的姿态与位置轨迹：\textbf{位置随时间按三次多项式变化，姿态由区间内变化的角速度积分得到}；
  \item 每个点按自己的时间戳找到对应轨迹段，结合外参转换到 odom 系，再统一到本帧结束时刻的 IMU 系。
\end{enumerate}

相比“假设帧内匀速”的简单做法，这能描述扫描期间角速度和加速度的变化，显著减少云台转动、机器人加减速时的点云畸变。计算上，同一 IMU 区间内的点可以复用轨迹参数，开销可控。

\subsection{地图结构：用统计量代替点}

ZLIO 的地图沿用 VoxelMap\cite{yuan2022voxelmap} 的结构：哈希表索引体素，每个体素内部用八叉树组织平面；若体素内点能用一个平面描述就保存平面，否则继续细分，直到满足平面条件或达到最大深度。平整区域用大体素，几何复杂处用小体素。

它真正改动的地方是\textbf{历史点的保存方式}：平面建立后，不再保存原始支持点，而是保存\textbf{点数、一阶矩、二阶矩}。以体素中心 $\bm{o}$ 为局部原点，记 $\bm{x}_i=\bm{p}_i-\bm{o}$，保存
\begin{equation}
N=\sum_i 1,\qquad \bm{s}=\sum_i \bm{x}_i,\qquad \bm{Q}=\sum_i \bm{x}_i\bm{x}_i^{\top}.
\label{eq:23-moments}
\end{equation}
新点加入时，只需累加 $1$、$\bm{x}$、$\bm{x}\bm{x}^{\top}$。由这些累计量恢复平面中心与散布矩阵：
\begin{equation}
\bm{c}=\bm{o}+\frac{\bm{s}}{N},\qquad
\bm{A}=\frac{\bm{Q}}{N}-\frac{\bm{s}\bm{s}^{\top}}{N^{2}},
\label{eq:23-plane}
\end{equation}
再取 $\bm{A}$ 最小特征值对应的单位特征向量作为法向：
\begin{equation}
\bm{A}\bm{n}=\lambda_{\min}(\bm{A})\,\bm{n},\qquad \norm{\bm{n}}=1.
\label{eq:23-normal}
\end{equation}
达到更新上限后，进一步释放统计量，只保留匹配所需的平面信息。

\begin{keybox}{这一改动为什么重要}
原始点云地图的内存随点数线性增长，且每来一个新点都要重新遍历支持点拟合平面。改成“点矩 + 均值/协方差”后，每个平面只存固定大小的统计量，更新是 $O(1)$ 的加法。作者给出的实测结果是：Mid360 ramp 数据 CPU 时间减少 72.8\%、峰值内存减少 86.6\%；Mid70 TN07 数据分别减少 62.6\% 和 90.7\%。
\end{keybox}

\subsection{面建模与残差}

建图点的测量误差和位姿误差会影响拟合出的平面，因此平面的中心和法向也带有不确定性。除此之外，实际墙面、地面未必完全平整，即使平均平面估得很准，局部点仍可能偏离它。ZLIO 把这两部分都计入观测方差：从支持点沿法向的散布中扣除已知噪声能解释的部分，剩余方差记为 $\tau^{2}$。

对变换到 odom 系的扫描点 $\bm{p}_i$，记测量协方差 $\bm{\Sigma}_i$、平面参数协方差 $\bm{\Sigma}_\pi$、残差对平面参数的雅可比 $\bm{h}_i$，则点面残差及其方差为
\begin{align}
r_i &= \bm{n}^{\top}(\bm{p}_i-\bm{c}),\label{eq:23-residual}\\
\sigma_i^{2} &= \bm{n}^{\top}\bm{\Sigma}_i\bm{n}+\tau^{2}+\bm{h}_i\bm{\Sigma}_\pi\bm{h}_i^{\top}.
\label{eq:23-resvar}
\end{align}
多个点匹配到同一个地图平面时，会共享这份平面误差。如果把它们当作相互独立，点数越多，滤波器就越容易“过度相信”这个平面。ZLIO 的处理是\textbf{按平面分组}，保留平面参数误差带来的相关性。同组中两个不同点的残差协方差为
\begin{equation}
\Cov(r_i,r_j)\simeq\bm{h}_i\bm{\Sigma}_\pi\bm{h}_j^{\top},\qquad i\ne j.
\label{eq:23-corr}
\end{equation}
实现时在平面的三个误差自由度内联合处理同组观测（两个法向方向偏差 + 沿法向的位置偏差），消去共享平面误差后，再把整组观测对位姿的约束交给滤波器，从而避免构造随点数增长的大协方差矩阵。记同平面残差向量 $\bm{r}_\pi$、雅可比 $\bm{H}_\pi$，最终残差项为
\begin{equation}
\bm{R}_\pi=\diag\left(\bm{n}^{\top}\bm{\Sigma}_i\bm{n}+\tau^{2}\right)+\bm{H}_\pi\bm{\Sigma}_\pi\bm{H}_\pi^{\top},\qquad
E_\pi=\bm{r}_\pi^{\top}\bm{R}_\pi^{-1}\bm{r}_\pi.
\label{eq:23-marginal}
\end{equation}
这是全书里少见的、把“观测相关性”显式处理的例子。大多数开源 LIO 直接把点面残差当独立，点数一多就会过拟合平面；Z-Localization 的做法值得借鉴。

\subsection{时空自适应}

\textbf{时间自适应}：ZLIO 希望更新频率介于 FAST-LIO 的整帧更新和 Point-LIO 的逐点更新之间，并根据运动与观测信息量自动调整。它把点云按采样时刻切成连续小段，每段去畸变后做一次滤波更新。每次更新后，根据近期线加速度、角加速度和当前点面观测的信息矩阵决定下一次更新间隔。记机体系下加速度与角加速度为 $\bm{\kappa}$，同一坐标系下按匹配点数平均的信息矩阵为 $\bar{\bm{\Lambda}}$，允许误差为 $\varepsilon$，则二阶运动项估计的偏差为
\begin{equation}
e^{2}(\Delta t)\simeq\frac{\Delta t^{4}}{4}\bm{\kappa}^{\top}\bar{\bm{\Lambda}}\bm{\kappa},
\label{eq:23-time-err}
\end{equation}
令其等于 $\varepsilon^{2}$ 并解出
\begin{equation}
\Delta t_{\mathrm{next}}=\mathrm{clip}\left(\left(\frac{4\varepsilon^{2}}{\bm{\kappa}^{\top}\bar{\bm{\Lambda}}\bm{\kappa}}\right)^{1/4},\ \Delta t_{\min},\ \Delta t_{\max}\right).
\label{eq:23-time}
\end{equation}
运动变化大、且这些变化会明显影响当前点面残差时，缩短间隔；反之延长间隔、积累更多点。本轮更新无效或分母为零时取上限。

\textbf{空间自适应}：为减少对固定采样分辨率的依赖，让同一套参数适应不同大小的场景，ZLIO 取保留时间窗口内代表点到雷达距离的中位数并平滑，作为场景尺度；场景越开阔，目标点数越多。再根据目标点数 $N_k^{*}$ 与实际点数 $N_k$ 的差，用 PD 控制调整采样体素边长：
\begin{equation}
e_k=N_k^{*}-N_k,\qquad
\ell_{k+1}=\mathrm{clip}\left(\ell_k-K_{p,k}e_k-K_{d,k}\frac{e_k-e_{k-1}}{\Delta t_k},\ \ell_{\min},\ \ell_{\max}\right).
\label{eq:23-space}
\end{equation}
点数偏少时减小体素（保留更多点），偏多时增大体素。两个增益随场景尺度、点数误差及其变化调整。

\subsection{不用在线动态点滤除}

作者也尝试过在线滤除动态点，但没有找到能同时兼顾计算开销和滤除率的方法。比赛场景中定位特征更多来自墙壁和棚顶，加入动态点滤除前后的定位效果差距很小；考虑额外开销，最终没有在前端采用在线滤除，而是提供\textbf{离线}工具，在建图后清理保存的地图。

\section{后端定位}\label{sec:23-backend}

\subsection{因子图}

后端因子图的每个节点对应一个关键帧，优化的是它在 \textbf{map 系}中的位置和姿态。图规模由固定时长的滑动窗口控制，窗口外的旧位姿通过边缘化移出（保留约束信息，见第 13 章）。开启 UWB 在线标定时，UWB 坐标系到 map 系的对齐以及时间偏置也作为共享变量持续保留在图中。

图中的因子有：
\begin{itemize}
  \item \textbf{锚点先验}：初始化/恢复成功后，把可信位姿作为图中的锚点，形成位姿先验。
  \item \textbf{里程计因子}：前端提供相邻关键帧之间的相对运动，把轨迹前后连接起来。
  \item \textbf{先验图配准因子}：有先验地图时，扫描与地图配准提供当前关键帧的全局位姿及约束强度，用来校正前端漂移。
  \item \textbf{回环因子}：没有先验地图时，回环匹配提供当前帧与历史帧之间的相对位姿，修正累积误差。
  \item \textbf{UWB 因子}：约束天线位置，完成时间与坐标对齐后与 UWB 测量比较。
  \item \textbf{重力方向因子}：无先验地图且启用 UWB 时，用前端的重力方向约束横滚和俯仰。
\end{itemize}

把这套因子图与第 13 章的位姿图对照，你会发现它就是一个“带多种观测的位姿图”：节点是位姿，边是相对位姿约束，只是观测来源更丰富。

\subsection{状态机}

后端不仅做优化，还负责判断“现在信不信得过”。它定义了五个状态：

\begin{table}[htbp]
  \centering
  \caption{Z-Localization 后端的定位状态}\label{tab:23-states}
  \small
  \begin{tabular}{p{2.4cm}p{5.0cm}p{4.6cm}}
    \toprule
    状态 & 含义 & 进入/恢复条件 \\
    \midrule
    INITIALIZING & 等待首次定位成功 & 启动时；成功后进入 NORMAL \\
    NORMAL & 定位正常 & 全局观测与预测一致，校正被接受 \\
    OBS\_DEGRADED & 点云质量差、观测弱 & 可用全局观测与预测冲突，但证据不足以否定当前位姿；拒绝冲突校正，继续用前端相对运动 \\
    SUSPECT & 前端漂移、初值已偏 & 冲突持续达阈值，或速度异常、无法局部匹配；保留前端，用地图/描述子/UWB 重新找 TF \\
    LOST & 估计器故障或前端发散 & SUSPECT 连续恢复失败；需重置前端、重建局部里程计，再用先验图或 UWB 找回全局位姿 \\
    \bottomrule
  \end{tabular}
\end{table}

判定逻辑分两步：先看点云能提供多少有效约束，再看配准/UWB 观测与当前预测是否一致。进入 SUSPECT 或 LOST 后，若已有可信位姿，输出会\textbf{冻结}在该位姿直到恢复成功——这比输出一个错误位姿要安全得多。这个“冻结而不是乱跳”的策略，在比赛里非常实用。

\subsection{描述子}

作者对描述子的理解很到位：“描述子本质上是在设计点云查询的键值对——用描述子作为键，关联地图中的位置和局部点云。新的点云到来时，先按描述子相似度找到候选位置，再通过点云配准确认位姿。”因此他们关心的是\textbf{相似度在候选位置上的分布}：正确位置附近能否形成足够尖锐、突出的峰值。

\textbf{Scan Context（SC）}\cite{kim2018scancontext}：以雷达为中心，按距离和方位把点云划分为环和扇区，每格记录最高点高度，形成二维描述；比较时尝试不同旋转角度，找到最相似的位置与朝向，为后续配准提供候选和 yaw 初值。SC 对高度分布敏感，生成和查询时都要先做重力矫正、把点云调平。比赛地图只有约 200 个描述子，因此检索直接遍历即可，无需 KD-tree。

\textbf{G-PROBE}\cite{gprobe2026}：复现了论文中的描述子与候选检索方法，面向 Avia、奥丁之眼这类\textbf{有限视场}点云。它把点云表示为带占据概率和不确定性的二维网格，并为地图准备不同方向的视场分支；检索时先筛候选，再在不同朝向下比较共同可见区域的占据与高度分布，同时降低那些在多个朝向上都相似的候选得分。作者指出这部分测试不充分、召回率未达 100\%，建议结合自己的雷达视场调参、充分验证后再用。

\subsection{配准与验收}

正常跟踪时，后端积累近期点云，结合当前预测位姿选取附近的先验子图，做点到面配准。初值好时直接细配准；对应点不足或残差大时，先粗配准再回到细层。

弱观测下的求解参考了 DCReg 的思路。典型例子是实验室走廊：沿走廊前后移动一点，点云仍然贴住两侧墙面，配准很难判断这个方向的位置偏差。ZLIO 用\textbf{舒尔补}分别分析平移和旋转，再根据弱方向构造预条件器，改善求解的数值条件。

配准后，根据残差尺度和海森矩阵对位姿误差做\textbf{白化}：强方向权重大，弱方向权重小，不可观方向不增加约束。再结合图预测的不确定性，通过 \textbf{NIS}（归一化新息平方）和卡方检验判断是否接受这次观测。正常跟踪可以使用部分方向的约束；初始化和重定位则要求\textbf{完整位姿}通过验收，再据此确定 TF 并建立锚点。第 2 章讲过，NIS 的二次型服从卡方分布，这里正是它的直接应用。

\subsection{UWB 初始化}

没有先验地图时，UWB 能提供位置，但单个位置无法确定前端 odom 系与全局坐标系之间的 yaw。ZLIO 把同一时刻的前端天线位置与 UWB 位置配对，通过一段运动求出两条轨迹之间的平移和 yaw，横滚俯仰沿用前端的重力方向。

对水平平移和 yaw 的简化对齐模型做可观测性分析：把轨迹去中心、假设位置噪声独立各向同性且方差相同，得到信息矩阵
\begin{equation}
\mathcal{I}=\frac{1}{\sigma_p^{2}}
\begin{bmatrix}
N&0&0\\
0&N&0\\
0&0&\sum_i\norm{\tilde{\bm{p}}_i}^{2}
\end{bmatrix},
\label{eq:23-uwb-info}
\end{equation}
于是 yaw 的标准差近似为
\begin{equation}
\sigma_\psi\approx\frac{\sigma_p}{\sqrt{\sum_i\norm{\tilde{\bm{p}}_i}^{2}}}.
\label{eq:23-yawstd}
\end{equation}
这个式子非常直观：停在原地时（$\tilde{\bm{p}}_i$ 都接近零）最后一项为零，yaw 无法确定；轨迹在水平面上展开后，航向信息迅速增加。所以“启动后跑开几米，yaw 就可观测了”不是玄学，而是式~\eqref{eq:23-yawstd} 的直接结论。

\subsection{Reset}

重置时，后端冻结当前输出并请求前端重置。前端完成重置、重新提供里程计和点云后，后端再接续或重新确定 $\text{map}\to\text{odom}$ 的 TF，恢复定位输出。作者坦言这套 reset 逻辑比较复杂，主要是为了尽量避免极端情况下的问题；如果场景简单，一个普通的重置服务可能就够了。

\section{使用流程：从 bag 建图到先验定位}\label{sec:23-workflow}

官方推荐的完整流程是：\textbf{回放建图 → 保存地图 → 自动清理 → CloudCompare 精修 → 网格生成地图包与描述子 → 先验定位}。关键命令如下（节选，完整参数见官方文档）：

\begin{lstlisting}[language=bash, caption={Z-Localization 建图与定位的关键命令}]
# 1. replay a bag and build the map
ros2 launch bringup bringup.launch.py \
  frontend_config:="$PWD/config/lio.yaml" \
  backend_config:="$PWD/pcd/site.yaml" \
  mode:=mapping use_sim_time:=true rviz:=true

# 2. save the map bundle after the replay finishes
ros2 service call /backend_node/save_map std_srvs/srv/Trigger "{}"

# 3. remove dynamic ghost points offline (FreeDOM-style)
ros2 run backend clean_map_bundle \
  --in "$PWD/pcd/site_raw" --out "$PWD/pcd/site_auto_clean" \
  --lidar-origin -0.011 -0.02329 0.04412

# 4. rebuild map bundle and descriptors from the refined PCD
ros2 run backend build_pcd_map_bundle \
  --input "$PWD/pcd/site_refined.pcd" --out-bundle "$PWD/pcd/site_prior" \
  --x-min 0 --x-max 28 --y-min 0 --y-max 15 \
  --sensor-z 0.0 --bundle-step 1.5 \
  --map-voxel 0.15 --scan-voxel 0.4 \
  --descriptor-type scan_context --sc-lidar-height 0.35

# 5. run localization against the prior map
ros2 launch bringup bringup.launch.py \
  frontend_config:="$PWD/config/lio.yaml" \
  backend_config:="$PWD/pcd/site.yaml" \
  mode:=localization prior_map:="$PWD/pcd/site_prior" \
  use_sim_time:=true rviz:=true
\end{lstlisting}

流程里有几个值得注意的工程细节：
\begin{itemize}
  \item \textbf{建图与定位用同一套入口}，只是 \texttt{mode} 不同，这减少了“两套代码不一致”的风险；
  \item \textbf{离线清理动态残影}参考 FreeDOM 的思路，作者实测“非常有效”，但仍不足以直接使用，需要 CloudCompare 手工收尾；
  \item \textbf{描述子按网格重新生成}，而不是建图时直接保存——尤其雷达水平安装时，这样生成的描述子质量更好；
  \item 保存地图包时保留\textbf{关键帧和局部点云}，后续清理需要它们；只存全局点云是不够的。
\end{itemize}

\section{与其他 LIO 方案的对比}\label{sec:23-compare}

\begin{table}[htbp]
  \centering
  \caption{Z-Localization 与主流 LIO 方案的定位}\label{tab:23-compare}
  \small
  \begin{tabular}{p{2.6cm}p{3.4cm}p{3.2cm}p{4.2cm}}
    \toprule
    方案 & 核心结构 & 全局定位方式 & 适合场景 \\
    \midrule
    FAST-LIO2 & 紧耦合 IEKF + iKD-Tree & 无（需外接） & 单机建图/里程计，算力中等 \\
    Point-LIO & 逐点更新 + 无帧累积 & 无（需外接） & 剧烈运动、IMU 易饱和 \\
    LIO-SAM & 因子图 + 预积分 + 回环 & 回环 & 大范围建图，有回环 \\
    Z-Localization & 前端 ZLIO + 后端因子图 & 先验图配准 + UWB + 回环 & 有场地先验、要求可恢复的比赛场景 \\
    \bottomrule
  \end{tabular}
\end{table}

选型时的判断逻辑：
\begin{enumerate}
  \item 有没有\textbf{场地先验地图}？有 → 优先考虑 Z-Localization 这种前后端分离方案；没有 → FAST-LIO2/Point-LIO 先跑起来，或像 Z-Localization 那样用 UWB 过渡。
  \item 会不会\textbf{剧烈运动/IMU 饱和}？会 → Point-LIO 或空间自适应采样；不会 → FAST-LIO2 足够。
  \item 需不需要\textbf{定位失效后自动恢复}？需要 → 必须有后端状态机和重定位，这正是 Z-Localization 的核心价值。
  \item 算力/内存是否紧张？紧张 → 注意其“用空间换时间”的后端设计（内存占用相对较高），可通过参数降低。
\end{enumerate}

\section{可借鉴的设计思想}\label{sec:23-lessons}

即使你不用这套代码，它的几个设计决策也值得抄进自己的系统：
\begin{enumerate}
  \item \textbf{把“局部估计”和“全局校正”解耦}：不要让一个紧耦合滤波器同时承担两件事，否则发散时你无法定位问题。
  \item \textbf{状态机比算法更重要}：比赛里“输出冻结”往往比“输出一个略错的位姿”更安全。
  \item \textbf{地图存统计量而不是点}：精度损失很小，内存和更新开销大幅下降。
  \item \textbf{显式处理观测相关性}：同一平面的点不独立，忽略相关性会让滤波器过度自信。
  \item \textbf{先验地图要“精修 + 重新生成描述子”}：地图质量决定定位上限。
  \item \textbf{用可观测性分析指导调参}：弱方向、yaw 可观测性、NIS 检验，都是把“玄学”变成“可测量”的手段。
\end{enumerate}

\begin{summary}

\begin{enumerate}
  \item Z-Localization 是面向比赛场景的 ROS 2 激光惯性定位栈，核心是\textbf{前端 ZLIO + 后端因子图}的分离架构。
  \item 前端负责 odom 系内的局部估计：SE$_2$(3)$\ltimes\mathfrak{se}(3)$ 误差建模、不要求静止的初始化、DLIO 式去畸变、基于点矩的 VoxelMap。
  \item 后端负责 map$\to$odom 的变换：多源因子图、五状态状态机、Scan Context/G-PROBE 描述子、带舒尔补白化的配准与 NIS 验收、UWB 对齐、Reset。
  \item 时空自适应让前端在“整帧更新”与“逐点更新”之间自适应，兼顾精度与资源。
  \item 它的工程取舍（冻结输出、离线清图、描述子重建、统计量地图）比具体公式更值得学习。
\end{enumerate}

\end{summary}

\begin{problemset}

  \item 为什么“把先验地图直接塞进 LIO”会导致发散时难以判断？拆成前后端后，哪些量变得可观测？
  \item 由式~\eqref{eq:23-yawstd} 解释：为什么 UWB 初始化时“跑几米”就能确定 yaw？如果只走直线，yaw 能确定吗？
  \item 点矩地图用 $N,\bm{s},\bm{Q}$ 三个统计量恢复平面。请写出由它们计算平面中心与协方差矩阵的完整步骤。
  \item 状态机里 OBS\_DEGRADED 和 SUSPECT 的区别是什么？各自应该采取什么动作？
  \item 如果你要把 Z-Localization 移植到只有 2D 雷达的机器人上，哪些模块需要改？哪些思想仍然适用？
\end{problemset}

\begin{reading}

\begin{itemize}
  \item \cite{zloc2026}：项目主页与 \texttt{docs/design.md}，本章的一手材料。
  \item \cite{yuan2022voxelmap}：自适应体素地图的原始论文，理解地图结构的必读。
  \item \cite{chen2023dlio}：连续时间去畸变思路的来源。
  \item \cite{kim2018scancontext}、\cite{gprobe2026}：两种描述子的原理与适用条件。
  \item \cite{equi2023}：SE$_2$(3) 等变误差建模的理论背景。
\end{itemize}


\end{reading}



\chapter{前沿方向与展望}\label{ch:24-frontier}

\chapterintro{前面二十三章讲的都是“今天能直接拿来用”的方法。这一章换个视角：看看这个领域正在往哪走。我们不会罗列所有论文，而是挑出四条主线——语义与学习、多模态融合、长期地图维护、学习式导航——说清它们各自解决什么问题、现在到了什么程度、离工程落地还有多远。读完这一章，你应该能判断：哪些“新方法”值得投入时间，哪些还只是实验室里的玩具。}

\section{从几何 SLAM 到语义与学习}\label{sec:24-semantic}

\subsection{语义 SLAM}

传统 SLAM 输出的地图是“几何的”：点、面、栅格，不知道哪里是墙、哪里是门、哪里是人。语义 SLAM 在几何之外再估计每个点的类别或每个物体的实例，好处有三：
\begin{itemize}
  \item \textbf{动态物体剔除}：知道“这是人”，就可以把人的点从建图中剔除，而不是靠几何一致性猜；
  \item \textbf{更高层的导航}：可以下“去门口”这种指令，而不只是“去坐标 $(3,2)$”；
  \item \textbf{回环与重定位}：语义描述子对光照、季节变化更鲁棒。
\end{itemize}
代价是引入了分割网络，算力开销大、标注依赖强、分割错误会污染地图。目前工程上常见的是“离线语义标注 + 在线几何定位”的折中。

\subsection{学习式特征与描述子}

第 12 章讲的 Scan Context 是手工设计的描述子。学习式方法用神经网络直接学“两帧是否同一地点”：
\begin{itemize}
  \item 视觉侧：SuperPoint、SuperGlue、LoFTR 等把特征提取与匹配都交给网络，在弱纹理、大视角变化下明显优于 ORB；
  \item 激光侧：PointNet++、OverlapPredator、 diffusion-based 描述子等直接处理无序点云。
\end{itemize}
问题是：模型需要大量数据、推理耗时、跨场景泛化不确定。比赛场景相对固定，手工描述子往往够用；跨季节、跨场地的长期建图才更能体现学习式方法的优势。

\subsection{神经隐式表示}

NeRF（神经辐射场）和 3D Gaussian Splatting（3DGS）用神经网络或高斯集合表示场景，能渲染出照片级的新视角图像。把它们用于 SLAM 是一个活跃方向：用位姿 + 图像联合优化辐射场，同时得到高保真地图和更精确的位姿。
\begin{itemize}
  \item 优点：地图紧凑（网络权重或高斯参数）、可渲染、可做光度一致性约束；
  \item 缺点：训练/优化慢，实时性差，几何精度不一定优于点云，且对大规模场景显存压力大。
\end{itemize}
目前它更适合“离线高精地图重建”或“小场景实时”，离“比赛场上 10 Hz 定位”还有距离。

\subsection{端到端里程计}

DROID-SLAM 这类方法把特征提取、匹配、位姿优化放进一个可微的迭代网络，端到端训练。它在公开数据集上精度很高，但泛化到新传感器、新场景仍需微调，且算力需求远高于传统方法。对算法组而言，更现实的用法是“用学习式方法做前端的一部分（如特征、描述子、深度补全），后端仍用经典优化”。

\section{多模态融合}\label{sec:24-multimodal}

第 9 章讲过 LiDAR + IMU 的紧耦合。前沿方向是把更多模态塞进同一个估计器：
\begin{itemize}
  \item \textbf{LiDAR-视觉-惯性}：FAST-LIVO2\cite{yuan2024fastlivo2}、R3LIVE\cite{lin2022r3live} 把点到面残差与光度误差联合优化，在几何退化（长走廊、隧道）或视觉退化（黑暗、强光）时互相补位。
  \item \textbf{事件相机}：输出异步“亮度变化”事件，延迟极低、动态范围极大，适合高速运动与高动态场景；缺点是数据稀疏、传统视觉算法要重写。
  \item \textbf{4D 毫米波雷达}：能测距离、方位、俯仰、速度，雨雾烟尘下仍工作；缺点是点云稀疏、噪声大。
  \item \textbf{UWB 与 GNSS}：作为全局约束，与 LIO 融合（第 23 章 Z-Localization 已有实践）。
\end{itemize}

多模态融合的核心难点不是“把数据放一起”，而是\textbf{时间同步、外参标定、噪声建模、失效检测}。第 26 章讲的标定，在模态越多时越重要。

\section{长期建图与地图维护}\label{sec:24-lifelong}

比赛只需要一张场地地图，但真实应用（仓储、园区、家庭）需要机器人在几个月甚至几年里持续使用同一张地图。这就带来“长期建图”问题：
\begin{itemize}
  \item \textbf{多会话 SLAM}：多次建图的结果要能拼接、对齐到同一坐标系；
  \item \textbf{变化检测}：家具挪了、货架变了、墙刷了新漆，地图要能更新而不是被旧地图“带偏”；
  \item \textbf{地图压缩与版本管理}：点云地图动辄几个 GB，需要压缩、分层、增量加载；
  \item \textbf{定位鲁棒性}：场景变化后，描述子检索的召回率会下降，需要“变化感知”的重定位。
\end{itemize}
Z-Localization 的“先验地图包 + 描述子 + 可选 UWB”已经是一个简化版的多会话方案；工业界更完整的方案通常包含地图服务、版本管理、云端更新。

\section{学习式导航}\label{sec:24-learnnav}

第 18--21 章的规划与控制都是“经典方法”：图搜索、采样、优化、MPC。前沿方向是用学习来做决策：
\begin{itemize}
  \item \textbf{学习式局部规划}：用强化学习或模仿学习直接输出速度指令，在复杂动态场景中可能比手工规则更灵活；缺点是可解释性差、安全性难保证、训练成本高。
  \item \textbf{端到端导航}：从传感器输入直接到动作输出（如“视觉导航”），在特定场景（走廊跟随、人群穿行）表现不错，但泛化与安全仍是问题。
  \item \textbf{语义/任务导航}：输入“去厨房拿杯子”，系统自己拆解成探索、定位、抓取子任务，这是具身智能的核心问题，目前仍在实验室阶段。
  \item \textbf{多机器人协同}：多机共享地图、协同探索、任务分配，难点在通信受限下的一致性。
\end{itemize}

一个务实的判断：\textbf{短期内，学习式方法更可能作为经典规划器的“插件”（如学习式代价函数、学习式启发、学习式恢复行为），而不是整体替换。}经典方法提供安全保证与可解释性，学习提供灵活性与泛化。

\section{对 RoboMaster 的启示}\label{sec:24-rm}

结合比赛场景，给算法组三条建议：
\begin{enumerate}
  \item \textbf{先把经典方法吃透。} FAST-LIO2/Point-LIO、Cartographer、Nav2、Z-Localization 这些开源项目已经覆盖了比赛 90\% 的需求。把它们的参数、失效模式、源码吃透，比追新论文的收益高得多。
  \item \textbf{关注“可恢复性”与“可观测性”。} 比赛里最致命的不是精度差 1 cm，而是定位飞了之后起不来。Z-Localization 的状态机与前后端分离思路，比任何新算法都值得抄。
  \item \textbf{谨慎引入学习式方法。} 如果要用，优先用在“离线”或“非安全关键”环节：离线建图清理、描述子训练、仿真环境生成。在线定位与规划仍应以经典方法为主。
\end{enumerate}

\begin{summary}

\begin{enumerate}
  \item 语义 SLAM 在几何之外增加类别/实例信息，利于动态剔除、高层导航与重定位，但依赖分割网络。
  \item 学习式特征、描述子、神经隐式表示在特定场景优于手工方法，但实时性、泛化与算力仍是瓶颈。
  \item 多模态融合（LiDAR-视觉-惯性、事件相机、4D 毫米波）是提升鲁棒性的主要方向，核心难点是同步与标定。
  \item 长期建图需要多会话对齐、变化检测、地图压缩与版本管理。
  \item 学习式导航短期内更可能作为经典规划器的插件，而非整体替代。
  \item 对比赛而言，吃透经典开源方案、重视可恢复性，比追新更重要。
\end{enumerate}

\end{summary}

\begin{reading}

\begin{itemize}
  \item \cite{yuan2024fastlivo2}、\cite{lin2022r3live}：多模态紧耦合的两个代表，读完能理解“视觉如何补几何”。
  \item \cite{barfoot2017state}：第 11 章之后讨论了学习与估计的结合，适合理性看待“端到端”。
  \item 3D Gaussian Splatting 原始论文与开源实现：理解神经隐式表示的能力与局限。
  \item 近三年的 ICRA/IROS/RSS 综述：跟踪语义 SLAM、长期建图、学习式规划的最新进展。
\end{itemize}


\end{reading}




\part{工程篇：把算法跑起来}

\chapter{ROS 2 工程实践}\label{ch:25-ros2}
\usetikzlibrary{arrows.meta}

\chapterintro{本章解决一个非常具体的问题：把前面各章的算法，从论文里的公式变成能复现、别人能接手、能装到车上跑起来的工程。ROS 2 是 RM 算法组事实上的标准框架，但新手卡住的地方往往不是算法本身，而是编译不过、话题收不到、TF 查不到、参数改了没生效、时间戳对不上。本章按"先跑通最小系统、再理解通信机制、最后建立工程习惯"的顺序，把这些坑逐一填平。读完你应当能独立建一个工作空间、读懂 launch 与参数文件、看懂 TF 树，并用命令行工具定位大部分"数据没来 / 频率不够 / 时间不对"类的问题。本章是第~\ref{ch:26-calibration}~章标定与第~\ref{ch:27-tuning}~章调参的脚手架。}

\section{ROS 2 是什么：从节点到计算图}\label{sec:25-what}

先回答"为什么需要中间件"。一台 RM 机器人上，雷达驱动、IMU 驱动、LIO、代价地图、规划器、串口通信程序往往是\textbf{不同的进程}，甚至由不同人用不同语言写成。如果它们之间用共享内存或自定义 socket 直连，那么每加一个传感器、每换一个算法，就要改一遍所有程序的接口，这是灾难。中间件的作用就是把这些连接抽象成"数据管道"：每个程序只声明"我发布什么"和"我订阅什么"，连接由框架自动建立。

ROS 2 的通信层是 DDS（Data Distribution Service，数据分发服务）\cite{macenski2022ros2}。DDS 有两个关键特性：

\begin{itemize}
  \item \textbf{以数据为中心}：程序发布的是"名为 \lstinline{/scan} 的数据流"，而不是"发给某个地址的消息"。
  \item \textbf{去中心化}：没有 ROS 1 里那个必须最先启动的 \lstinline{roscore}（master）。每个节点通过组播/单播自动发现对方，任何节点先启动、后启动都能连上。
\end{itemize}

一次运行中所有节点与它们之间的连接，构成一张\textbf{计算图}（computation graph）。图上的元素只有五种，先记住这张表，后面所有工程问题都能归到其中一类：

\begin{table}[htbp]
  \centering
  \caption{ROS 2 计算图的五种元素}\label{tab:25-graph}
  \small
  \begin{tabular}{lp{5.2cm}p{5.2cm}}
    \toprule
    元素 & 语义 & 典型例子 \\
    \midrule
    节点 node & 一个进程内的计算单元（\lstinline{rclcpp::Node}） & \lstinline{fast_lio}、\lstinline{controller_server} \\
    话题 topic & 单向、多对多、连续的数据流 & \lstinline{/lidar/points}、\lstinline{/cmd_vel} \\
    服务 service & 一请求一响应的同步调用 & 清空代价地图、保存地图 \\
    动作 action & 长任务，带反馈与取消 & "导航到某点"、\lstinline{NavigateToPose} \\
    参数 parameter & 节点运行时的配置项 & \lstinline{max_vel_x}、\lstinline{filter_size_surf} \\
    \bottomrule
  \end{tabular}
\end{table}

ROS 2 与 ROS 1 的差异不是"版本升级"，而是通信层换了一整套，这决定了工程写法的不同。表~\ref{tab:25-vs1} 列出最影响写代码的几处。

\begin{table}[htbp]
  \centering
  \caption{ROS 1 与 ROS 2 的关键差异}\label{tab:25-vs1}
  \small
  \begin{tabular}{lll}
    \toprule
    维度 & ROS 1 & ROS 2 \\
    \midrule
    通信层 & 自定义 TCPROS/UDPROS & DDS（可换实现：Fast DDS、Cyclone DDS） \\
    节点发现 & 中心 master & 分布式自发现 \\
    服务质量 & 无统一 QoS & 每个话题/订阅可配 QoS \\
    构建系统 & catkin + \lstinline{catkin_make} & ament + \lstinline{colcon} \\
    进程模型 & 单 master，节点可合进程 & 节点即进程，鼓励多进程与组合 \\
    实时与安全 & 基本没有 & 支持实时执行器、SROS2 加密 \\
    \bottomrule
  \end{tabular}
\end{table}

\begin{keybox}{一句话理解}
ROS 2 把"程序之间怎么连"从代码里搬到了配置里（话题名、QoS、参数、launch）。所以调试 ROS 2 问题时，先怀疑\textbf{连接配置}，再怀疑算法代码。
\end{keybox}

\section{最小工作空间：colcon 编译流程}\label{sec:25-ws}

ROS 2 的代码组织单位是\textbf{功能包}（package），包的集合叫\textbf{工作空间}（workspace，通常命名为 \lstinline{*_ws}）。标准目录结构是 \lstinline{src/}（放源码）、\lstinline{build/}（中间文件）、\lstinline{install/}（最终产物）、\lstinline{log/}（编译日志）。新手最容易犯的错误是把源码直接丢在 \lstinline{ws} 根目录——\lstinline{colcon} 只在 \lstinline{src/} 下递归找包。

一个可用的最小流程如下（在 Ubuntu 22.04 + ROS 2 Humble 下测试）：

\begin{lstlisting}[language=bash,caption={创建并编译一个最小工作空间},label={lst:25-ws}]
# 1. create workspace
mkdir -p ~/rm_ws/src && cd ~/rm_ws/src
ros2 pkg create --build-type ament_cmake rm_lio \
    --dependencies rclcpp sensor_msgs nav_msgs

# 2. put your source into rm_ws/src/rm_lio/

# 3. build (symlink-install keeps config/launch editable without rebuild)
cd ~/rm_ws
colcon build --symlink-install --cmake-args -DCMAKE_BUILD_TYPE=Release
# build only one package and show full log when debugging
colcon build --packages-select rm_lio --event-handlers console_direct+

# 4. every new terminal must source BOTH layers
source /opt/ros/humble/setup.bash
source ~/rm_ws/install/setup.bash
ros2 pkg list | grep rm_lio
\end{lstlisting}

三点解释：

\begin{enumerate}
  \item \textbf{source 顺序不能反}。\lstinline{/opt/ros/humble} 是 underlay（系统安装的 ROS），\lstinline{install/} 是 overlay（你自己编译的）。先 underlay 后 overlay，overlay 里的同名包才会覆盖系统包。反过来会出现"改了代码但运行的是旧库"的灵异现象。
  \item \lstinline{--symlink-install} 让 \lstinline{install/} 里的 launch、YAML、URDF 变成指向 \lstinline{src/} 的符号链接，改配置文件不用重新编译。C++ 源码仍然需要重编，但省掉了大量"改个参数要编译五分钟"的痛苦。
  \item \lstinline{--cmake-args -DCMAKE_BUILD_TYPE=Release} 非常重要。默认可能是无优化编译，LIO 这类点云算法会慢 2--5 倍，直接表现为"雷达 10 Hz，里程计出不来"。
\end{enumerate}

\lstinline{package.xml} 与 \lstinline{CMakeLists.txt} 里真正需要改的字段不多，见表~\ref{tab:25-pkg}。

\begin{table}[htbp]
  \centering
  \caption{功能包两个清单文件的关键字段}\label{tab:25-pkg}
  \small
  \begin{tabular}{llp{6.4cm}}
    \toprule
    文件 & 字段 & 作用 \\
    \midrule
    \lstinline{package.xml} & \lstinline{<depend>} & 编译+运行都需要的依赖，新手统一写这个即可 \\
    \lstinline{package.xml} & \lstinline{<exec_depend>} & 只有运行期需要（如 Python 依赖、launch 包） \\
    \lstinline{CMakeLists.txt} & \lstinline{find_package()} & 找依赖；少一个就会在链接期报 undefined reference \\
    \lstinline{CMakeLists.txt} & \lstinline{ament_target_dependencies()} & 把依赖的头文件路径与库挂到目标上 \\
    \lstinline{CMakeLists.txt} & \lstinline{install(TARGETS ...)} & 不写这行，\lstinline{ros2 run} 找不到可执行文件 \\
    \lstinline{CMakeLists.txt} & \lstinline{install(DIRECTORY ...)} & 不写这行，装好的包缺少 launch 与 yaml \\
    \bottomrule
  \end{tabular}
\end{table}

\section{话题、服务与动作：什么时候用哪个}\label{sec:25-comm}

三者都是"节点之间的通信"，但语义完全不同。选错的典型后果是：用服务传 10 Hz 点云（每次请求都要等响应，吞吐崩掉），或者用话题做"开始导航"命令（没有反馈、不能取消）。

\begin{table}[htbp]
  \centering
  \caption{三种通信机制的选型表}\label{tab:25-comm}
  \small
  \begin{tabular}{lp{3.4cm}p{2.6cm}p{3.2cm}}
    \toprule
    机制 & 数据流形态 & 阻塞性 & 典型用途 \\
    \midrule
    话题 & 单向、多对多、持续 & 不阻塞（异步） & 传感器流、\lstinline{/cmd_vel}、\lstinline{/tf} \\
    服务 & 一请求一响应 & 调用方等待 & 保存地图、清空代价地图、切换参数 \\
    动作 & 长任务 + 周期反馈 + 可取消 & 可中断 & \lstinline{NavigateToPose}、机械臂抓取 \\
    \bottomrule
  \end{tabular}
\end{table}

判断规则可以压缩成三条：

\begin{enumerate}
  \item 数据是\textbf{连续产生、允许丢帧}的（点云、图像、IMU）$\Rightarrow$ 话题。
  \item 调用是\textbf{瞬时完成、需要知道成功与否}的（存图、重置）$\Rightarrow$ 服务。
  \item 任务\textbf{持续若干秒、过程中要看进度、可能要中途放弃}的（导航、探索）$\Rightarrow$ 动作。
\end{enumerate}

写代码时最容易忽略的是回调线程模型：话题回调运行在执行器（executor）的线程里。如果回调里做了 100 ms 的耗时计算（比如把点云转成另一套坐标系），那么同一执行器里的所有其他回调都会被拖慢，表现为"某个话题频率突然掉了"。第~\ref{sec:25-habit} 节会专门讲这一点。

\section{QoS：ROS 2 最容易踩的坑}\label{sec:25-qos}

QoS（Quality of Service）是 ROS 2 相对 ROS 1 最大的行为变化。它回答三个问题：

\begin{itemize}
  \item \textbf{可靠性 reliability}：\lstinline{reliable} 表示传输层保证送达（丢包重传）；\lstinline{best_effort} 表示尽力而为，丢了就丢了。
  \item \textbf{历史 history + 深度 depth}：保存最近多少条消息。\lstinline{keep_last} + \lstinline{depth=1} 表示只留最新一条，\lstinline{depth=10} 表示留 10 条。
  \item \textbf{持久性 durability}：\lstinline{volatile} 表示只给"现在在线"的订阅者；\lstinline{transient_local} 表示晚加入的订阅者也能拿到最后一条（类似"锁存"）。
\end{itemize}

为什么传感器数据要用 \lstinline{best_effort}？因为雷达 10 Hz 出点云，单帧几十 MB；如果链路拥塞，\lstinline{reliable} 会重传旧帧，越积越多，最后延迟爆炸。而定位算法只关心\textbf{最新}那帧，旧帧补过来也没用。反过来，控制指令 \lstinline{/cmd_vel} 必须 \lstinline{reliable}——丢一条可能就是撞墙。地图则要用 \lstinline{transient_local}，因为 RViz 或导航节点可能在建图节点之后启动，它仍然需要那张已经发布过的地图。

\begin{table}[htbp]
  \centering
  \caption{常见话题的推荐 QoS 配置}\label{tab:25-qos}
  \small
  \begin{tabular}{llll}
    \toprule
    话题 & reliability & history/depth & durability \\
    \midrule
    \lstinline{/lidar/points} & best\_effort & keep\_last / 5 & volatile \\
    \lstinline{/imu/data} & best\_effort & keep\_last / 100 & volatile \\
    \lstinline{/odom}、\lstinline{/tf} & best\_effort & keep\_last / 10 & volatile \\
    \lstinline{/map} & reliable & keep\_last / 1 & transient\_local \\
    \lstinline{/cmd_vel} & reliable & keep\_last / 1 & volatile \\
    \lstinline{/initialpose} & reliable & keep\_last / 1 & volatile \\
    \bottomrule
  \end{tabular}
\end{table}

\textbf{QoS 不匹配是新手最常见、也最容易被忽略的故障}。DDS 的兼容规则是：发布者提供的质量必须\textbf{不弱于}订阅者要求的质量。所以：

\begin{itemize}
  \item 发布 \lstinline{reliable}、订阅 \lstinline{best_effort} $\Rightarrow$ 兼容，能收到。
  \item 发布 \lstinline{best_effort}、订阅 \lstinline{reliable} $\Rightarrow$ \textbf{不兼容}，一条也收不到。
\end{itemize}

此时节点不会崩，只是安静地没有数据。排查的命令是：

\begin{lstlisting}[language=bash,caption={QoS 不匹配的诊断命令},label={lst:25-qosdiag}]
# check publisher/subscriber QoS of a topic
ros2 topic info /lidar/points --verbose
# expected output shows "Reliability: BEST_EFFORT" for both sides

# is data actually flowing?
ros2 topic hz /lidar/points
ros2 topic echo /lidar/points --once
\end{lstlisting}

如果 \lstinline{ros2 topic info --verbose} 显示"发布者 1 个、订阅者 1 个"但 \lstinline{ros2 topic hz} 无输出，绝大多数情况就是 QoS 不匹配或话题名不一致（注意前导斜杠与命名空间）。

\section{参数与 launch：把配置从代码里赶出去}\label{sec:25-param}

\textbf{参数}是节点的运行时配置。C++ 里必须先声明再使用，声明时给出默认值：

\begin{lstlisting}[language=C++,caption={声明与读取参数},label={lst:25-declare}]
// in constructor or on_configure()
this->declare_parameter<double>("filter_size_surf", 0.5);
this->declare_parameter<int>("max_iteration", 4);
this->declare_parameter<std::string>("lid_topic", "/lidar/points");

this->get_parameter("filter_size_surf", filter_size_surf_);
this->get_parameter("max_iteration", max_iteration_);
// dynamic update: register a callback instead of polling
param_cb_ = this->add_on_set_parameters_callback(
    [this](const std::vector<rclcpp::Parameter> & params) {
        // validate before applying; reject invalid values
        return rcl_interfaces::msg::SetParametersResult{};
    });
\end{lstlisting}

参数覆盖的优先级从低到高是：代码默认值 $<$ YAML 参数文件 $<$ launch 文件中的参数 $<$ 命令行 \lstinline{-p name:=value}。所以"改了 YAML 没生效"的第一反应应该是：检查 launch 里是不是又写了一遍参数，或者命令行的 \lstinline{-p} 把它压掉了。

\textbf{launch} 文件统一用 Python 写。一个典型的"一个 launch 起整套定位"如下：

\begin{lstlisting}[language=Python,caption={定位系统的 launch 文件片段},label={lst:25-launch}]
import os
from ament_index_python.packages import get_package_share_directory
from launch import LaunchDescription
from launch.actions import DeclareLaunchArgument
from launch.conditions import IfCondition
from launch.substitutions import LaunchConfiguration
from launch_ros.actions import Node


def generate_launch_description():
    pkg = get_package_share_directory("rm_lio")
    use_sim_time = LaunchConfiguration("use_sim_time")
    rviz = LaunchConfiguration("rviz")

    lio = Node(
        package="rm_lio", executable="lio_node", name="lio_node",
        output="screen",
        parameters=[os.path.join(pkg, "config", "lio.yaml"),
                    {"use_sim_time": use_sim_time}],
        remappings=[("points", "/lidar/points")],
    )
    rviz_node = Node(
        package="rviz2", executable="rviz2", condition=IfCondition(rviz),
        arguments=["-d", os.path.join(pkg, "rviz", "lio.rviz")],
    )
    return LaunchDescription([
        DeclareLaunchArgument("use_sim_time", default_value="false"),
        DeclareLaunchArgument("rviz", default_value="true"),
        lio, rviz_node,
    ])
\end{lstlisting}

\lstinline{output="screen"} 让节点日志打到终端而不只是写进 \lstinline{~/.ros/log}；\lstinline{remappings} 是"不改源码就换话题名"的利器，跨团队联调时几乎必用。启动后检查参数是否真的生效：

\begin{lstlisting}[language=bash,caption={检查参数实际取值},label={lst:25-paramcheck}]
ros2 param list /lio_node
ros2 param get /lio_node filter_size_surf
ros2 param dump /lio_node > /tmp/effective.yaml   # dump what is really used
\end{lstlisting}

\lstinline{ros2 param dump} 是排错神器：它输出的是\textbf{生效后}的完整参数，可以直接和你的 YAML 对比。

\section{TF2：坐标系的"宪法"}\label{sec:25-tf}

TF2 维护一棵随时间变化的\textbf{坐标系树}。定位系统的所有输出最终都要落到这棵树上，否则 RViz 里的点云会乱飞、Nav2 会报"代价地图收不到机器人位姿"。

TF 的关系是\textbf{父子}关系，每个变换有两个时间戳（发布时刻与有效时刻），查询时按时间插值。一台 RM 机器人的典型结构如图~\ref{fig:25-tf}。

\begin{figure}[htbp]
  \centering
  \begin{tikzpicture}[
      frame/.style={rectangle, rounded corners=2pt, draw=structurecolor, thick,
        fill=structurecolor!8, minimum height=7mm, minimum width=20mm, font=\small},
      lab/.style={font=\footnotesize, inner sep=1.5pt, draw=black!60, rounded corners=1pt,
        fill=white},
      arr/.style={-{Stealth[length=2mm]}, thick, draw=black!70}
    ]
    \node[frame] (map) at (0,0) {map};
    \node[frame] (odom) at (3.4,0) {odom};
    \node[frame] (base) at (6.8,0) {base\_link};
    \node[frame] (lidar) at (10.4,2.3) {lidar\_link};
    \node[frame] (imu) at (10.4,0) {imu\_link};
    \node[frame] (cam) at (10.4,-2.3) {camera\_link};
    \node[frame] (gimb) at (6.8,-3.2) {gimbal\_link};

    \draw[arr] (map) -- node[above=2pt,lab]{全局定位} (odom);
    \draw[arr] (odom) -- node[above=2pt,lab]{LIO 里程计} (base);
    \draw[arr] (base) -- node[lab,pos=0.5,above right=1pt]{静态外参} (lidar);
    \draw[arr] (base) -- node[above=2pt,lab]{静态外参} (imu);
    \draw[arr] (base) -- node[lab,pos=0.5,below right=1pt]{静态外参} (cam);
    \draw[arr] (base) -- node[right=2pt,lab]{云台反馈} (gimb);
  \end{tikzpicture}
  \caption{RM 机器人典型 TF 树。\lstinline{map}$\to$\lstinline{odom} 由全局定位修正，\lstinline{odom}$\to$\lstinline{base\_link} 由里程计给出，传感器到机体的变换是常数外参。}\label{fig:25-tf}
\end{figure}

几个必须记住的约定：

\begin{itemize}
  \item \lstinline{map} 是全局连续但会跳变的坐标系（重定位时整体平移）；\lstinline{odom} 是局部连续、绝不跳变的坐标系。所以\textbf{只有定位模块能发布 \lstinline{map}$\to$\lstinline{odom}}，且每次只发一个小修正，不能一次修正几米，否则规划器会看到机器人瞬移。
  \item \lstinline{lookupTransform(target, source, time)} 返回的是"把 \lstinline{source} 系下的向量变到 \lstinline{target} 系"的变换，即 $T_{target,source}$。方向写反是外参错误最常见的来源。
  \item 传感器到机体的外参用 \lstinline{static_transform_publisher} 发布，它们在标定后就是常数，见第~\ref{ch:26-calibration}~章。
\end{itemize}

\begin{lstlisting}[language=bash,caption={TF 的发布、查询与诊断},label={lst:25-tfcmd}]
# static extrinsics: lidar -> base_link, args: x y z yaw pitch roll
ros2 run tf2_ros static_transform_publisher 0.15 0 0.30 0 0 0 \
    base_link lidar_link

# single query at a given time (prints translation and quaternion)
ros2 run tf2_ros tf2_echo map base_link

# draw the whole tree as a pdf and show per-link publish rates
ros2 run tf2_tools view_frames
\end{lstlisting}

在 C++ 中查询 TF 时，三种异常必须分别处理，否则要么崩溃，要么拿到错误数据：

\begin{lstlisting}[language=C++,caption={lookupTransform 的三类异常},label={lst:25-tfexc}]
try {
    // TimePointZero means "latest available"; do NOT use it on recorded bags
    auto T = tf_buffer_->lookupTransform("map", "base_link",
                                         tf2::TimePointZero, tf2::durationFromSec(0.1));
} catch (const tf2::LookupException & e) {
    // frame name typo, or the publisher never existed
} catch (const tf2::ConnectivityException & e) {
    // both frames exist but no path connects them (missing static extrinsic)
} catch (const tf2::ExtrapolationException & e) {
    // requested time outside the buffer, or timestamps are out of order
}
\end{lstlisting}

\begin{note}
\textbf{时间戳乱序}是 LIO 装车后非常常见的故障：点云和 IMU 来自不同驱动，若某一路的时间戳比另一路早了 200 ms，TF 缓冲会不断收到"过去"的数据，表现为 \lstinline{ExtrapolationException} 间歇性爆发、RViz 中位姿抖动。处理办法不是调大缓冲区超时，而是回到第~\ref{ch:26-calibration}~章做时间同步。
\end{note}

\section{调试工具与工程习惯}\label{sec:25-tool}
\label{sec:25-habit}

工具用对了，排错效率差一个数量级。表~\ref{tab:25-tools} 是我自己常用的清单。

\begin{table}[htbp]
  \centering
  \caption{ROS 2 常用命令行工具}\label{tab:25-tools}
  \small
  \begin{tabular}{llp{6.2cm}}
    \toprule
    命令 & 回答什么问题 & 备注 \\
    \midrule
    \lstinline{ros2 node list} & 有哪些节点活着 & 先确认进程起来了 \\
    \lstinline{ros2 topic list -t} & 有哪些话题、类型是什么 & 类型不匹配会无法连接 \\
    \lstinline{ros2 topic hz} & 频率是多少 & 判断驱动/算法是否跟得上 \\
    \lstinline{ros2 topic echo --once} & 消息内容对不对 & 看 \lstinline{header.stamp} 与 \lstinline{frame_id} \\
    \lstinline{ros2 topic info -v} & QoS 与连接数 & QoS 不匹配的唯一快速判据 \\
    \lstinline{ros2 param get/dump} & 参数实际值是多少 & 确认 YAML 有没有被覆盖 \\
    \lstinline{ros2 bag record -a} & 把现场录下来 & 排错第一原则：先录包 \\
    \lstinline{ros2 bag play --clock} & 离线复现 + 仿真时间 & 配 \lstinline{use_sim_time:=true} \\
    \lstinline{tf2_tools view_frames} & TF 树长什么样 & 有没有断链、哪条链路频率低 \\
    \lstinline{ros2 doctor} & 系统层面有没有异常 & 网络、RMW 实现、环境变量 \\
    \lstinline{rqt_graph} & 计算图里谁连谁 & 一眼看出没连上的节点 \\
    \lstinline{rqt_plot} & 数值曲线 & 调 PID、看速度指令 \\
    \bottomrule
  \end{tabular}
\end{table}

工程习惯比工具更重要，以下五条建议按重要性排序：

\begin{enumerate}
  \item \textbf{用命名空间隔离多机/多模块}。launch 里加 \lstinline{namespace="sentry"}，所有话题变成 \lstinline{/sentry/odom}，避免与另一台车或另一个模块撞名。多机器人系统尤其要注意 \lstinline{/tf} 与 \lstinline{/tf_static} 必须显式 remap，否则全车 TF 混进一棵树。
  \item \textbf{仿真时间要一致}。跑 bag 时必须 \lstinline{use_sim_time:=true} 且播放时带 \lstinline{--clock}，所有节点的时间源才会统一为 bag 里的 \lstinline{/clock}。混用会导致 TF 全部超时。
  \item \textbf{日志分级}。\lstinline{RCLCPP_INFO} 用于状态，\lstinline{RCLCPP_WARN} 用于可恢复异常，\lstinline{RCLCPP_ERROR} 用于必须处理的问题。默认级别 INFO，现场排查时用 \lstinline{--ros-args --log-level debug} 临时打开，不要长期开 DEBUG（I/O 本身就拖慢实时性）。
  \item \textbf{执行器与回调组}。默认单线程执行器下所有回调串行执行。对于"点云处理 20 ms + 控制 1 ms"这类混合负载，应当用 \lstinline{MultiThreadedExecutor} 配合 \lstinline{MutuallyExclusive} 回调组把控制与感知分开，否则控制频率会被点云拖垮。
  \item \textbf{实时性}。对硬实时环节用 \lstinline{taskset -c 2,3} 绑定 CPU 核，避免与编译、录包抢核；不要在回调里做内存分配与文件 I/O。
\end{enumerate}

\begin{keybox}{案例：定位节点频率从 20 Hz 掉到 5 Hz}
\textbf{现象}：\lstinline{ros2 topic hz /lio/odom} 只有 5 Hz，但雷达 10 Hz 正常。\textbf{排查顺序}：(1) \lstinline{ros2 topic hz /lidar/points} 确认输入频率；(2) \lstinline{top -H} 看进程是否 CPU 打满，若单核 100\% 则多半是编译没开 Release 或多线程配置错误；(3) 若 CPU 不满，看 \lstinline{ros2 topic info -v} 是否有多个订阅者拖慢——\lstinline{reliable} 的慢订阅者会反压发布者；(4) 最后检查是否有人在同一执行器里做 \lstinline{ros2 bag record -a}，录包会把所有大消息序列化一遍。\textbf{结论}：先分清"算不完"还是"收不到"，不要一上来就改算法参数。
\end{keybox}

\begin{summary}

\begin{enumerate}
  \item ROS 2 用 DDS 取代 ROS 1 的 master，节点自发现；计算图只有节点、话题、服务、动作、参数五种元素，所有工程问题都能归到这五类。
  \item 标准编译命令是 \lstinline{colcon build --symlink-install --cmake-args -DCMAKE_BUILD_TYPE=Release}；每个终端必须按 underlay$\to$overlay 顺序 source，否则改的代码不会生效。
  \item 选通信机制的三条规则：连续数据用话题，瞬时调用用服务，长任务用动作。
  \item QoS 中最关键的是 reliability 与 durability：传感器用 \lstinline{best_effort}，控制用 \lstinline{reliable}，地图用 \lstinline{transient_local}；发布弱、订阅强会一条都收不到。
  \item 参数优先级是 代码默认 $<$ YAML $<$ launch $<$ 命令行；用 \lstinline{ros2 param dump} 看实际生效值。
  \item TF 树中 \lstinline{map}$\to$\lstinline{odom} 由全局定位发布且只能小步修正，\lstinline{odom}$\to$\lstinline{base\_link} 由里程计发布，传感器外参是静态变换。
  \item 排查问题先录包（\lstinline{ros2 bag record}），再离线用 \lstinline{topic hz}、\lstinline{topic info -v}、\lstinline{tf2_tools}、\lstinline{rqt_graph} 定位。
\end{enumerate}


\end{summary}

\begin{problemset}

  \item 你的雷达节点发布 \lstinline{/lidar/points} 用的是默认 QoS（reliable, depth 10），RViz2 订阅时选了 best\_effort。能收到数据吗？反过来呢？请说明 DDS 的兼容规则。
  \item 一个节点先于雷达驱动启动，它能在启动瞬间收到最后一帧点云吗？需要把哪个 QoS 项改成什么值？地图话题也存在同样问题，请解释为什么地图用 \lstinline{transient_local}。
  \item 写出 \lstinline{map}$\to$\lstinline{base\_link} 的完整链式变换，并说明如果 \lstinline{map}$\to$\lstinline{odom} 在某一帧突然跳了 1 m，规划器会出现什么现象、应该怎样限制单次修正量。
  \item 你在 launch 里给 \lstinline{lio_node} 传了 YAML，但 \lstinline{ros2 param get} 显示还是默认值。列出至少三种可能原因。
  \item 设计一个实验，证明"点云回调耗时 20 ms"是控制频率下降的原因，并给出两种不修改算法的缓解方案。
\end{problemset}

\begin{reading}

\begin{itemize}
  \item Macenski 等人的 ROS 2 综述 \cite{macenski2022ros2}：讲清楚 DDS 选型、实时执行器与 ROS 1 的取舍，读完能理解"为什么 ROS 2 的 QoS 长这样"。
  \item ROS 2 官方文档的 \href{https://docs.ros.org/en/humble/Concepts/About-Quality-of-Service-Settings.html}{About QoS Settings} 与 \href{https://docs.ros.org/en/humble/Tutorials/Intermediate/Tf2/Tf2-Main.html}{TF2 教程}：QoS 兼容表与 TF 广播/监听的权威说明。
  \item 第~\ref{ch:22-nav2-practice}~章的 Nav2 工程链路：把本章的 launch、参数、TF 用法放到一个真实导航系统里再走一遍。
\end{itemize}


\end{reading}

\chapter{标定与时间同步}\label{ch:26-calibration}
\usetikzlibrary{arrows.meta}

\chapterintro{标定是定位系统里"最不值得怀疑、却最常出错"的一环。算法调不好时，新手总怀疑滤波器，老手先怀疑标定：外参差 $1\degree$，10\,m 处就是 17\,cm；时间差 10\,ms，2\,m/s 时就是 2\,cm，而且这类误差会随运动放大，表现为"直着走没事，一转弯就飘"。本章把标定拆成四件事：相机内参（张正友法）、IMU 噪声（Allan 方差）、多传感器外参（手眼标定 $AX=XB$）、时间同步。每件事都给到从假设、公式到工程判据的完整链条，读完你应该能自己标定一台车，并判断结果到底"能用"还是"不能用"。}

\section{为什么标定：误差是怎么被放大的}\label{sec:26-why}

\textbf{外参误差与时间误差都是与运动耦合的误差}：静止时完全看不出来，一动就暴露。先算两个数建立量级感。

\subsection{旋转外参误差：$d\sin\theta$}

设雷达与机体之间存在未被补偿的旋转误差 $\theta$（俯仰、横滚或偏航任一分量），则距离 $d$ 处的点被放错的位置偏差约为
\begin{equation}
\delta=\norm{\bm{R}(\theta)-\bm{I}}\,d=2\sin\frac{\theta}{2}\,d\approx\theta d,
\label{eq:26-leverarm}
\end{equation}
其中 $\bm{R}(\theta)$ 是绕某轴转 $\theta$ 的旋转矩阵；$\theta$ 很小时 $2\sin(\theta/2)\approx\theta$（弧度）。

\begin{table}[htbp]
  \centering
  \caption{旋转外参误差造成的点位置偏差（cm，$\delta\approx\theta d$）}\label{tab:26-rot}
  \small
  \begin{tabular}{ccccc}
    \toprule
    $\theta$ $\backslash$ $d$ & 1\,m & 3\,m & 5\,m & 10\,m \\
    \midrule
    $0.5\degree$ & 0.9 & 2.6 & 4.4 & 8.7 \\
    $1\degree$   & 1.7 & 5.2 & 8.7 & 17.5 \\
    $2\degree$   & 3.5 & 10.5 & 17.5 & 34.9 \\
    $5\degree$   & 8.7 & 26.2 & 43.6 & 87.2 \\
    \bottomrule
  \end{tabular}
\end{table}

工程含义有两条：近处点（$1$--$3$\,m）对 $1\degree$ 误差不敏感，\textbf{所以"看起来建图还行"不能证明外参对}；定位时匹配的是 $8$--$15$\,m 外的墙与门，同样的 $1\degree$ 就是 15--25\,cm，表现为缓慢漂移或重定位错误。判据不是"图好看"，而是远处点云是否重影。

\subsection{时间误差：$v\Delta t$}

时间戳错位 $\Delta t$ 等价于把测量放错位置，偏差近似为
\begin{equation}
\delta=v\,\Delta t,
\label{eq:26-timedelay}
\end{equation}
其中 $v$ 是该点相对机体的运动速度。

\begin{table}[htbp]
  \centering
  \caption{时间不同步造成的位移偏差（mm，$\delta=v\Delta t$）}\label{tab:26-time}
  \small
  \begin{tabular}{ccccc}
    \toprule
    $\Delta t$ $\backslash$ $v$ & 0.5\,m/s & 1\,m/s & 2\,m/s & 4\,m/s \\
    \midrule
    1\,ms  & 0.5 & 1 & 2 & 4 \\
    5\,ms  & 2.5 & 5 & 10 & 20 \\
    10\,ms & 5 & 10 & 20 & 40 \\
    20\,ms & 10 & 20 & 40 & 80 \\
    50\,ms & 25 & 50 & 100 & 200 \\
    \bottomrule
  \end{tabular}
\end{table}

最容易被忽略的是\textbf{半帧偏移}：有的雷达把一帧的时间戳打在帧\textbf{起始}时刻，有的打在\textbf{结束}时刻，10\,Hz 下二者差 100\,ms，比"时钟不同步"严重一个数量级。拿到任何一款雷达，第一件事是查手册确认时间戳定义。

\section{外参的数学定义与标定链}\label{sec:26-extrinsic}

全书约定 $T_{AB}=(\bm{R}_{AB},{}^{A}\bm{p}_{B})$ 表示"把 $B$ 系向量变到 $A$ 系"：${}^{A}\bm{p}=\bm{R}_{AB}{}^{B}\bm{p}+{}^{A}\bm{p}_{B}$。于是"雷达在机体系中的外参"记作 $T_{BL}$，反之 $T_{LB}=T_{BL}^{-1}$。

\textbf{方向弄反是外参错误的第一大来源}，症状很有迷惑性：点云整体旋转/镜像，但建图"自洽"（所有点一致地错），只有和 IMU、视觉对比时才暴露。两个防错手段：一是\textbf{物理检验}——雷达水平放置、正前方 5\,m 放箱子，外参正确时应满足 $x\approx5$、$y\approx0$、$z\approx0$（机体 $x$ 前、$y$ 左、$z$ 上），任一符号反了立刻看出来；二是\textbf{只记一个方向关系}——$T_{BL}$ 的下标"先 $B$ 后 $L$"读作"$L$ 到 $B$"，代码里 \lstinline{extrinsic_T} 填 ${}^{B}\bm{p}_{L}$，\lstinline{extrinsic_R} 填 $\bm{R}_{BL}$，点云先乘 $\bm{R}_{BL}$ 再加 ${}^{B}\bm{p}_{L}$。

标定量之间有依赖：手眼标定需要内参先标好，时间偏移又需要外参有初值。顺序与判据见表~\ref{tab:26-chain}。

\begin{table}[htbp]
  \centering
  \caption{标定链：顺序、工具与验收判据}\label{tab:26-chain}
  \small
  \begin{tabular}{lllp{4.6cm}}
    \toprule
    顺序 & 标定量 & 常用工具 & 验收判据 \\
    \midrule
    1 & 相机内参、畸变 & OpenCV / Kalibr & 重投影 RMS $<0.5$\,px \\
    2 & IMU 噪声与零偏 & imu\_utils & 重力模长 $9.81\pm0.1$ \\
    3 & 相机-IMU 外参 + 时间偏移 & Kalibr & 重投影 $<1$\,px \\
    4 & 雷达-IMU 外参 + 时间偏移 & 手眼标定 & 点云墙厚 $<2$\,cm \\
    5 & 雷达-相机外参 & 标定板 + PnP & 投影边缘偏差 $<2$\,px \\
    \bottomrule
  \end{tabular}
\end{table}

\section{相机内参：针孔模型与畸变}\label{sec:26-pinhole}

相机系 $C$ 下的点 ${}^{C}\bm{P}=[X,Y,Z]^{\top}$（$Z$ 为深度）投到像素 $\bm{m}=[u,v]^{\top}$ 的过程是"先除以深度做透视，再乘内参矩阵"：
\begin{equation}
s\begin{bmatrix}u\\v\\1\end{bmatrix}=\bm{K}\begin{bmatrix}X\\Y\\Z\end{bmatrix},
\qquad
\bm{K}=\begin{bmatrix}f_x & s_0 & c_x\\ 0 & f_y & c_y\\ 0&0&1\end{bmatrix},
\label{eq:26-project}
\end{equation}
其中 $s=Z$ 是齐次尺度，$f_x,f_y$ 是以像素为单位的焦距（$f_x=f/d_x$，$d_x$ 为像元尺寸），$(c_x,c_y)$ 是\textbf{主点}，$s_0$ 是像元倾斜项（多数相机可置零）。

\textbf{数字例子}：$f_x=600$、$f_y=540$、$c_x=320$、$c_y=240$，${}^{C}\bm{P}=[1.0,0.5,2.0]^{\top}$\,m，则
\begin{equation}
u=600\times\frac{1.0}{2.0}+320=620,\qquad v=540\times\frac{0.5}{2.0}+240=375.
\end{equation}
若深度估计错了 10\%（$Z$ 从 2.0 变 2.2），$u$ 变成 592.7，横向误差 27\,px——这就是单目深度误差致命的来源。

真实镜头不满足理想针孔。把像素反投影到归一化平面 $x=X/Z,\;y=Y/Z$，记 $r^2=x^2+y^2$，Brown--Conrady 模型把理想坐标映射成\textbf{含畸变}的坐标：
\begin{equation}
\begin{aligned}
x_d &= x\bigl(1+k_1r^2+k_2r^4+k_3r^6\bigr)+2p_1xy+p_2\bigl(r^2+2x^2\bigr),\\
y_d &= y\bigl(1+k_1r^2+k_2r^4+k_3r^6\bigr)+p_1\bigl(r^2+2y^2\bigr)+2p_2xy,
\end{aligned}
\label{eq:26-distort}
\end{equation}
再乘内参得到像素。$k_1,k_2,k_3$ 是\textbf{径向}畸变（$k_1<0$ 桶形，$k_1>0$ 枕形），$p_1,p_2$ 是\textbf{切向}畸变（镜头与成像面不平行）。以 $k_1=-0.28$、$x=y=0.35$（视场边缘）为例，$r^2=0.245$，$x_d=0.35(1-0.28\times0.245)=0.326$，边缘向内拉 0.024 个归一化单位，按 $f_x=600$ 算是 14.4\,px。

\section{张正友标定法：从单应到内参}\label{sec:26-zhang}

张正友 1999 年的平面标定法是今天所有相机内参标定的基础：用同一块平面棋盘在\textbf{不同姿态}下拍若干张图，无需知道棋盘位姿即可解出内参。下面完整推导。

\subsection{第一步：平面到图像的单应}

把世界系定在棋盘平面上，$Z=0$，角点 ${}^{W}\bm{M}=[X,Y,0]^{\top}$，齐次形式 $\tilde{\bm{M}}=[X,Y,0,1]^{\top}$。投影为 $s\,\tilde{\bm{m}}=\bm{K}[\bm{R}\,|\,\bm{t}]\tilde{\bm{M}}$。把 $\bm{R}$ 按列拆开，利用 $Z=0$ 把第三列乘掉：
\begin{equation}
s\begin{bmatrix}u\\v\\1\end{bmatrix}
=\bm{K}\bigl[\bm{r}_1\;\bm{r}_2\;\bm{r}_3\;\bm{t}\bigr]
\begin{bmatrix}X\\Y\\0\\1\end{bmatrix}
=\bm{K}\bigl[\bm{r}_1\;\bm{r}_2\;\bm{t}\bigr]
\begin{bmatrix}X\\Y\\1\end{bmatrix}.
\label{eq:26-planehom}
\end{equation}
于是每张图定义一个 $3\times3$ 的\textbf{单应矩阵}：
\begin{equation}
\bm{H}=\bigl[\bm{h}_1\;\bm{h}_2\;\bm{h}_3\bigr]=\lambda\bm{K}\bigl[\bm{r}_1\;\bm{r}_2\;\bm{t}\bigr],
\label{eq:26-homography}
\end{equation}
其中 $\lambda$ 是齐次尺度。给 4 组以上角点对应（世界坐标已知、像素可检测），$\bm{H}$ 能用 DLT 线性解出（每对点 2 个方程，8 个未知量）。

\subsection{第二步：用旋转矩阵的正交性构造约束}

关键观察：$\bm{r}_1,\bm{r}_2$ 是旋转矩阵的列，满足\textbf{单位长度}与\textbf{相互正交}：$\bm{r}_1^{\top}\bm{r}_2=0$，$\norm{\bm{r}_1}=\norm{\bm{r}_2}=1$。由 $\bm{h}_i=\lambda\bm{K}\bm{r}_i$ 得 $\bm{r}_i=\lambda^{-1}\bm{K}^{-1}\bm{h}_i$，代入正交条件并约掉 $\lambda^{-2}$：
\begin{equation}
\bm{h}_1^{\top}\bm{K}^{-\top}\bm{K}^{-1}\bm{h}_2=0,
\label{eq:26-orth}
\end{equation}
代入等长条件：
\begin{equation}
\bm{h}_1^{\top}\bm{K}^{-\top}\bm{K}^{-1}\bm{h}_1=\bm{h}_2^{\top}\bm{K}^{-\top}\bm{K}^{-1}\bm{h}_2.
\label{eq:26-equal}
\end{equation}
定义对称矩阵 $\bm{B}=\bm{K}^{-\top}\bm{K}^{-1}\in\mathbb{R}^{3\times3}$，上面两式就是关于 $\bm{B}$ 的两个线性约束。\textbf{每张图给出两个约束}，而 $\bm{B}$ 只有 6 个独立元素，所以 3 张不同姿态的图即可确定 $\bm{B}$。

\subsection{第三步：把约束写成线性方程组}

把 $\bm{B}$ 的 6 个独立元素排成 $\bm{b}=[B_{11},B_{12},B_{22},B_{13},B_{23},B_{33}]^{\top}$。对任意两列 $\bm{h}_i,\bm{h}_j$，二次型可写成内积 $\bm{h}_i^{\top}\bm{B}\bm{h}_j=\bm{v}_{ij}^{\top}\bm{b}$，其中
\begin{equation}
\bm{v}_{ij}=\begin{bmatrix}
h_{i1}h_{j1}\\
h_{i1}h_{j2}+h_{i2}h_{j1}\\
h_{i2}h_{j2}\\
h_{i3}h_{j1}+h_{i1}h_{j3}\\
h_{i3}h_{j2}+h_{i2}h_{j3}\\
h_{i3}h_{j3}
\end{bmatrix}\in\mathbb{R}^{6}.
\label{eq:26-vij}
\end{equation}
式~\eqref{eq:26-orth}、\eqref{eq:26-equal} 于是成为 $\bm{v}_{12}^{\top}\bm{b}=0$ 与 $(\bm{v}_{11}-\bm{v}_{22})^{\top}\bm{b}=0$。$n$ 张图堆成 $\bm{V}\bm{b}=\bm{0}$，$\bm{V}\in\mathbb{R}^{2n\times6}$。

\textbf{尺度不确定性是本质的}：$\bm{K}$ 放大 $c$ 倍，$\bm{B}$ 缩小 $c^2$ 倍，约束完全不变。所以加 $\norm{\bm{b}}=1$，取 $\bm{V}^{\top}\bm{V}$ 最小特征值对应的特征向量（即 $\bm{V}$ 最小奇异值对应的右奇异向量）。实践中 $n\ge5$ 更抗噪，标定软件默认 20--30 张。

\subsection{第四步：由 $\bm{B}$ 闭式解出内参}

记 $f_x=\alpha$、$f_y=\beta$、$s_0=\gamma$、$c_x=u_0$、$c_y=v_0$，Zhang 给出的闭式解为
\begin{equation}
\begin{aligned}
v_0 &= \frac{B_{12}B_{13}-B_{11}B_{23}}{B_{11}B_{22}-B_{12}^{2}},\qquad
\lambda_0 = B_{33}-\frac{B_{13}^{2}+v_0\bigl(B_{12}B_{13}-B_{11}B_{23}\bigr)}{B_{11}},\\
\alpha &= \sqrt{\frac{\lambda_0}{B_{11}}},\qquad
\beta = \sqrt{\frac{\lambda_0 B_{11}}{B_{11}B_{22}-B_{12}^{2}}},\qquad
\gamma = -\frac{B_{12}\alpha^{2}\beta}{\lambda_0},\qquad
u_0 = \frac{\gamma v_0}{\beta}-\frac{B_{13}\alpha^{2}}{\lambda_0}.
\end{aligned}
\label{eq:26-intrinsic}
\end{equation}
$\lambda_0$ 正是那个被归一化掉的尺度；所有内参只依赖 $\bm{B}$ 元素之间的\textbf{比值}，因此尺度自动约掉。当像元近似正方形（$\gamma\approx0$，即 $B_{12}\approx0$）时退化为好记的形式：$u_0=-B_{13}/B_{11}$，$v_0=-B_{23}/B_{22}$，$\alpha=1/\sqrt{B_{11}}$，$\beta=1/\sqrt{B_{22}}$。

\subsection{第五步：每张图的外参}

由 $\bm{K}^{-1}\bm{h}_1=\lambda\bm{r}_1$ 取模长得 $\lambda=1/\norm{\bm{K}^{-1}\bm{h}_1}$，于是
\begin{equation}
\bm{r}_1=\frac{\bm{K}^{-1}\bm{h}_1}{\norm{\bm{K}^{-1}\bm{h}_1}},\qquad
\bm{r}_2=\frac{\bm{K}^{-1}\bm{h}_2}{\norm{\bm{K}^{-1}\bm{h}_2}},\qquad
\bm{r}_3=\bm{r}_1\times\bm{r}_2,\qquad
\bm{t}=\frac{\bm{K}^{-1}\bm{h}_3}{\norm{\bm{K}^{-1}\bm{h}_1}}.
\label{eq:26-extrinsic}
\end{equation}
两个细节：$\bm{r}_1,\bm{r}_2$ 因噪声不一定严格正交，$\bm{r}_3$ 用叉积保证右手系；$\lambda$ 的符号由"棋盘在相机前方"确定。最后对 $[\bm{r}_1\;\bm{r}_2\;\bm{r}_3]$ 做 SVD 投影到最近的旋转矩阵：$\bm{R}=\bm{U}\bm{V}^{\top}$，其中 $\bm{U}\bm{\Sigma}\bm{V}^{\top}$ 是那个近似矩阵的 SVD。

\subsection{第六步：最大似然细化}

前五步是最小化\textbf{代数误差}，不是真正的重投影误差，且还没考虑畸变。真正的目标是让所有角点的重投影误差平方和最小：
\begin{equation}
J\bigl(\{\bm{K},k_1,k_2,p_1,p_2\},\{\bm{R}_i,\bm{t}_i\}\bigr)
=\sum_{i=1}^{n}\sum_{j=1}^{m}\norm{\bm{m}_{ij}-\hat{\bm{m}}\bigl(\bm{K},k_{1:3},p_{1:2},\bm{R}_i,\bm{t}_i,{}^{W}\bm{M}_j\bigr)}^{2},
\label{eq:26-reproj}
\end{equation}
其中 $\bm{m}_{ij}$ 是第 $i$ 张图第 $j$ 个角点的检测像素，$\hat{\bm{m}}$ 由式~\eqref{eq:26-project}、\eqref{eq:26-distort} 算出。用 Levenberg--Marquardt 求解，初值取前五步结果（畸变初值取 0）。未知量 $4+5+6n$ 个，单目 20 张图约 129 个，OpenCV 利用稀疏结构加速。

\textbf{判据}：$\sqrt{J/(nm)}$ 即"重投影 RMS"，单位像素。$<0.3$\,px 很好，$0.3$--$0.5$\,px 正常，$>1$\,px 说明角点检测或模型有问题。一组 $640\times480$ 相机的正常结果是 $f_x=615.3$、$f_y=616.1$、$c_x=322.7$、$c_y=241.2$、$k_1=-0.2887$、$k_2=0.0721$、$p_1=0.0003$、$p_2=-0.0002$、RMS $0.31$\,px——注意 $f_x\approx f_y$、主点接近图像中心；若主点偏离中心超过图像宽度的 10\%，先查是不是角点顺序错了。

\textbf{实践要点}：(1) 图要覆盖画面四角与中心，畸变系数靠边缘样本；(2) 平面与像平面夹角取 $30\degree$--$60\degree$，全正对着拍会让单应退化；(3) 15--30 张，包含平移、旋转、远近变化；(4) 板子必须平且够大，纸贴不平会让 RMS 永远卡在 1\,px 以上；(5) 标定前后不能动对焦环。

\section{IMU 标定与 Allan 方差}\label{sec:26-imu}

陀螺与加速度计的测量可写成
\begin{equation}
\tilde{\bm{\omega}}=\bigl(\bm{I}+\bm{S}_g\bigr)\bm{M}_g\bm{\omega}+\bm{b}_g+\bm{n}_g,\qquad
\tilde{\bm{a}}=\bigl(\bm{I}+\bm{S}_a\bigr)\bm{M}_a\bm{a}+\bm{b}_a+\bm{n}_a,
\label{eq:26-imuerr}
\end{equation}
其中 $\bm{S}$ 是尺度因子误差，$\bm{M}$ 是轴间非正交，$\bm{b}$ 是零偏，$\bm{n}$ 是白噪声。零偏造成随时间线性累积的漂移，可标定也可在线估计；尺度/轴间是常值矩阵误差，绕一圈会表现为"转 $90\degree$ 实际积分出 $91\degree$"；白噪声不可标定，只能用统计量描述。

\textbf{零偏与尺度的六面法}：把 IMU 静置在六个正交朝向各 30--60\,s。加速度计每面测得模长应为 9.81，用六面数据可同时解出 $\bm{b}_a$ 与 $\bm{S}_a$（$\bm{M}_a$ 一般取单位阵）。陀螺零偏取长时间静置的均值，但要注意\textbf{零偏随温度漂移}，规范做法是上电预热 5--10 分钟再取。

\subsection{Allan 方差：定义}

白噪声与零偏这类随机过程不能用一个数描述，要看它在不同时间尺度上的表现。设采样间隔为 $\tau_0$，得到序列 $y[1],\dots,y[N]$（陀螺为角速度，加速度计为比力）。取簇长度 $m$，把序列分成 $K=\lfloor N/m\rfloor$ 个不重叠的簇，簇内平均为
\begin{equation}
\bar{y}_k=\frac{1}{m}\sum_{i=1}^{m}y\bigl[(k-1)m+i\bigr],\qquad k=1,\dots,K,
\label{eq:26-cluster}
\end{equation}
对应平均时间 $\tau=m\tau_0$。\textbf{Allan 方差}定义为相邻簇平均之差的方差的一半：
\begin{equation}
\sigma_A^{2}(\tau)=\frac{1}{2(K-1)}\sum_{k=1}^{K-1}\bigl(\bar{y}_{k+1}-\bar{y}_k\bigr)^{2},
\qquad \sigma_A(\tau)=\sqrt{\sigma_A^{2}(\tau)}.
\label{eq:26-allan}
\end{equation}
为什么除以 2？若相邻簇平均独立、方差各为 $\sigma^2$，差值方差是 $2\sigma^2$，除以 2 正好还原单点方差。为什么用"相邻差分"而不用"减去总均值"？因为后者会把零偏漂移当成正常波动滤掉，而 Allan 方差要保留慢变信息。

它与功率谱密度的关系是
\begin{equation}
\sigma_A^{2}(\tau)=4\int_{0}^{\infty}S_y(f)\,\frac{\sin^{4}(\pi f\tau)}{(\pi f\tau)^{2}}\,\diff f,
\label{eq:26-allanpsd}
\end{equation}
传递函数相当于一个随 $\tau$ 变形的带通滤波器：$\tau$ 小只看高频，$\tau$ 大只看低频。\textbf{所以一条 $\sigma_A(\tau)$ 曲线能同时分离出好几类噪声}，这正是它的价值。

\begin{figure}[htbp]
  \centering
  \begin{tikzpicture}[x=1.35cm,y=1.05cm,font=\small]
    \draw[-{Stealth[length=2mm]}, thick] (0,0) -- (7.2,0) node[right] {$\log\tau$};
    \draw[-{Stealth[length=2mm]}, thick] (0,0) -- (0,4.9) node[above] {$\log\sigma_A(\tau)$};
    \node[below left,font=\footnotesize] at (0,0) {$\tau_0$};
    \draw[thick, structurecolor] (0.35,4.35) -- (2.55,3.25);
    \draw[thick, structurecolor] (2.55,3.25) -- (3.75,3.15);
    \draw[thick, structurecolor] (3.75,3.15) -- (6.55,4.55);
    \node[font=\footnotesize, anchor=south west] at (0.35,4.4) {$-\tfrac12$：角度随机游走};
    \node[font=\footnotesize, anchor=south] at (3.15,3.3) {$0$：零偏不稳定性};
    \node[font=\footnotesize, anchor=north west] at (4.45,4.6) {$+\tfrac12$：速率随机游走};
    \draw[dashed, black!50] (2.55,0) -- (2.55,3.25);
    \node[below,font=\footnotesize] at (2.55,0) {$\tau=1$\,s};
    \draw[dashed, black!50] (0,3.25) -- (2.55,3.25) node[above left,font=\footnotesize,pos=0.12] {$\sigma_A(1\,\mathrm{s})=N$};
    \draw[dashed, black!50] (0,3.15) -- (3.75,3.15) node[below right,font=\footnotesize,pos=0.08] {$0.664\,B$};
  \end{tikzpicture}
  \caption{Allan 偏差曲线的双对数示意。左段斜率 $-1/2$ 是白噪声（角度/速度随机游走），中间平台是零偏不稳定性，右段斜率 $+1/2$ 是速率随机游走。}\label{fig:26-allan}
\end{figure}

\subsection{怎么从曲线上读参数}

各类噪声的双对数斜率不同，可以逐段拟合（表~\ref{tab:26-allan}）。

\begin{table}[htbp]
  \centering
  \caption{Allan 偏差的斜率与噪声类型（$y$ 为角速度或比力）}\label{tab:26-allan}
  \small
  \begin{tabular}{llcl}
    \toprule
    噪声类型 & $\sigma_A(\tau)$ 的形式 & 斜率 & 参数读法 \\
    \midrule
    量化噪声 & $\sqrt{3}\,Q/\tau$ & $-1$ & 左端最陡段 \\
    角度（速度）随机游走 & $N/\sqrt{\tau}$ & $-1/2$ & $\tau=1$\,s 处数值 $=N$ \\
    零偏不稳定性 & $0.664\,B$ & $0$ & 平台高度 $/0.664$ \\
    速率随机游走 & $K\sqrt{\tau/3}$ & $+1/2$ & 右段斜率拟合 \\
    漂移率斜坡 & $R\tau/\sqrt{2}$ & $+1$ & 最右端上扬段 \\
    \bottomrule
  \end{tabular}
\end{table}

\begin{keybox}{读数口诀（只记这一条就够用）}
\textbf{$\tau=1$\,s 处读到的 $\sigma_A$，数值上就是这个传感器的噪声密度}（陀螺单位 $^{\circ}$/s/$\sqrt{\mathrm{Hz}}$，加速度计单位 m/s$^2$/$\sqrt{\mathrm{Hz}}$）。陀螺的噪声密度乘 60 即得常用的 $^{\circ}$/$\sqrt{\mathrm{h}}$。
\end{keybox}

\textbf{数字例子}：某 MEMS 陀螺静置 2\,h 的 Allan 曲线在 $\tau=1$\,s 处读得 $\sigma_A=0.014\,^{\circ}$/s，换成 SI 单位并换算成角度随机游走：
\begin{equation}
N_g=0.014\times\frac{\pi}{180}=2.44\times10^{-4}\ \mathrm{rad/s/\sqrt{Hz}},
\qquad N_g\times60=0.84\ ^{\circ}\mathrm{/\sqrt{h}},
\end{equation}
与 datasheet 的 0.014\,$^{\circ}$/s/$\sqrt{\mathrm{Hz}}$ 一致。若 IMU 以 $f_s=200$\,Hz 输出，带宽取 $f_s/2$，则单样本标准差 $\sigma=N_g\sqrt{f_s/2}=2.4\times10^{-3}$\,rad/s$\approx0.14\,^{\circ}$/s。曲线平台在 $\tau\approx50$\,s 处约 $1.5\,^{\circ}$/h，得零偏不稳定性 $B=1.5/0.664\approx2.3\,^{\circ}$/h，即静止时零偏在几十秒尺度上有约 $1.1\times10^{-5}$\,rad/s 的随机浮动。

\subsection{从 Allan 读数到滤波器参数}

LIO/VIO 里要填的是协方差，不是曲线本身，转换只有两步：测量噪声取 $N_g^2$、$N_a^2$（FAST-LIO 的 \lstinline{gyr_cov}、\lstinline{acc_cov}），零偏随机游走取速率随机游走系数的平方 $K^2$（\lstinline{b_gyr_cov}、\lstinline{b_acc_cov}）。不同开源仓库对"连续时间噪声密度"的约定差一个 $\Delta t$ 的幂次，\textbf{同一个 IMU 在不同仓库的最优值可能相差一个常数因子，这是正常的}。工程做法：用 Allan 读数给量级初值，再在同一段 bag 上微调——噪声项调大更信 IMU，调小更信点云（见第~\ref{ch:27-tuning}~章）。

\section{雷达-IMU 外参：手眼标定 $AX=XB$}\label{sec:26-handeye}

待求量是 $X=T_{LI}$（把 IMU 系点变到雷达系），6 个自由度。已知两路轨迹：雷达里程计给出的 $T_{WL}(t_k)$、IMU 积分给出的 $T_{WI}(t_k)$。真实关系为
\begin{equation}
T_{WI}(t)=T_{WL}(t)\,X.
\label{eq:26-handerel}
\end{equation}

\textbf{推导 $AX=XB$。}取两个时刻 $t_1,t_2$。先在 $t_1$ 左乘逆：
\begin{equation}
T_{WI}(t_1)^{-1}=\bigl(T_{WL}(t_1)X\bigr)^{-1}=X^{-1}T_{WL}(t_1)^{-1}.
\end{equation}
再右乘 $t_2$ 的式子 $T_{WI}(t_2)=T_{WL}(t_2)X$：
\begin{equation}
T_{WI}(t_1)^{-1}T_{WI}(t_2)=X^{-1}T_{WL}(t_1)^{-1}T_{WL}(t_2)X.
\end{equation}
左边是 IMU 的相对运动 $B=T_{I_1I_2}$，右边括号内是雷达的相对运动 $A=T_{L_1L_2}$，于是 $B=X^{-1}AX$，等价于
\begin{equation}
A\,X=X\,B.
\label{eq:26-axxb}
\end{equation}
把 $X=(\bm{R}_X,\bm{t}_X)$ 展开，它同时给出旋转与平移两组方程：
\begin{align}
\bm{R}_A\bm{R}_X &= \bm{R}_X\bm{R}_B, \label{eq:26-axxbrot}\\
\bm{R}_A\bm{t}_X+\bm{t}_A &= \bm{R}_X\bm{t}_B+\bm{t}_X
\;\Longrightarrow\;
\bigl(\bm{R}_A-\bm{I}\bigr)\bm{t}_X=\bm{R}_X\bm{t}_B-\bm{t}_A.
\label{eq:26-axxbt}
\end{align}
\textbf{关键结构}：式~\eqref{eq:26-axxbrot} 只含 $\bm{R}_X$，所以\textbf{先解旋转、再解平移}，两步都是线性最小二乘。

\textbf{第一步：解旋转。}把旋转写成四元数，$AX=XB$ 变成
\begin{equation}
\bm{q}_A\otimes\bm{q}_X=\bm{q}_X\otimes\bm{q}_B.
\label{eq:26-quat}
\end{equation}
定义左乘与右乘矩阵（$\bm{q}=[q_0,\bm{q}_v^{\top}]^{\top}$）：
\begin{equation}
[\bm{q}]_L=\begin{bmatrix}q_0&-\bm{q}_v^{\top}\\ \bm{q}_v&q_0\bm{I}+[\bm{q}_v]_{\times}\end{bmatrix},\qquad
[\bm{q}]_R=\begin{bmatrix}q_0&-\bm{q}_v^{\top}\\ \bm{q}_v&q_0\bm{I}-[\bm{q}_v]_{\times}\end{bmatrix},
\end{equation}
它们满足 $[\bm{q}]_L\bm{q}'=\bm{q}\otimes\bm{q}'$、$[\bm{q}]_R\bm{q}'=\bm{q}'\otimes\bm{q}$。于是
\begin{equation}
\bigl([\bm{q}_A]_L-[\bm{q}_B]_R\bigr)\bm{q}_X=\bm{0}.
\label{eq:26-quatlinear}
\end{equation}
每对运动给 4 个方程，堆成 $\bm{M}\bm{q}_X=\bm{0}$（$\bm{M}\in\mathbb{R}^{4K\times4}$），取最小奇异值对应的右奇异向量再归一化即为 $\bm{q}_X$。

\textbf{可解性条件}：所有旋转轴必须"张开"。若机器人只绕竖直轴原地转，$A,B$ 的旋转轴都平行，$\bm{M}$ 秩不足，$X$ 中绕该轴的分量不可观。所以标定轨迹必须包含绕至少两个不平行轴的旋转：前后倾斜、左右横滚、加上偏航，各转 $30\degree$ 以上。

\textbf{第二步：解平移。}$\bm{R}_X$ 已知后，式~\eqref{eq:26-axxbt} 是线性方程组，堆 $K$ 对运动：
\begin{equation}
\begin{bmatrix}\bm{R}_A^{(1)}-\bm{I}\\ \vdots\\ \bm{R}_A^{(K)}-\bm{I}\end{bmatrix}\bm{t}_X
=\begin{bmatrix}\bm{R}_X\bm{t}_B^{(1)}-\bm{t}_A^{(1)}\\ \vdots\\ \bm{R}_X\bm{t}_B^{(K)}-\bm{t}_A^{(K)}\end{bmatrix},
\label{eq:26-tls}
\end{equation}
即 $\bm{D}\bm{t}_X=\bm{c}$ 的 $\mathbb{R}^{3K\times3}$ 最小二乘，解 $\hat{\bm{t}}_X=(\bm{D}^{\top}\bm{D})^{-1}\bm{D}^{\top}\bm{c}$。可解性同样要求旋转轴张开：若所有轴相同，$\bm{R}_A-\bm{I}$ 的零空间方向一致，沿该方向平移不可观。

\textbf{第三步：非线性细化。}代数解在噪声下不是最优，最后在 $\SE(3)$ 上用残差
\begin{equation}
\bm{e}_R^{(k)}=\log\!\Bigl(\bigl(\bm{R}_A^{(k)}\bm{R}_X\bigr)^{\top}\bm{R}_X\bm{R}_B^{(k)}\Bigr),\qquad
\bm{e}_t^{(k)}=\bigl(\bm{R}_A^{(k)}-\bm{I}\bigr)\bm{t}_X-\bm{R}_X\bm{t}_B^{(k)}+\bm{t}_A^{(k)}
\end{equation}
做最小二乘（Ceres/GTSAM 均可），初值取前两步结果。$\log$ 是 $\SO(3)\to\so(3)$ 的对数映射（第~\ref{ch:02-math} 章）。当雷达帧率低而运动剧烈时，两帧之间不能当成瞬时刚体变换，此时应改用连续时间轨迹模型，同时估计外参与时间偏移。

\section{时间同步与时间偏移估计}\label{sec:26-time}

一台机器人上至少四类时间源，精度差几个数量级（表~\ref{tab:26-clocksync}）。最省事又最有效的做法是\textbf{硬件触发}：让主时钟同时触发相机曝光与雷达扫描起始，驱动把触发时刻写进 \lstinline{header.stamp}，采样时刻在物理上就对齐了。

\begin{table}[htbp]
  \centering
  \caption{常见时间同步方案}\label{tab:26-clocksync}
  \small
  \begin{tabular}{llll}
    \toprule
    方案 & 典型精度 & 成本 & 适用 \\
    \midrule
    系统时钟（无同步） & 秒级漂移 & 0 & 不可用于融合 \\
    NTP（\lstinline{chrony}） & 0.1--10\,ms & 0 & 一般用途、跨机器 \\
    PTP（IEEE 1588，硬件时间戳） & $<1\,\mu$s & 需支持 PTP 的网卡 & 多相机、多雷达 \\
    GPS/PPS 秒脉冲 & $<100$\,ns & 需 GPS 模块 & 室外、多机对齐 \\
    硬件触发线 & 曝光级（$\mu$s） & 需相机/雷达支持 & 视觉-惯性标配 \\
    \bottomrule
  \end{tabular}
\end{table}

\textbf{ROS 2 里的两个开关必须分清}：\lstinline{use_sim_time:=true} 时节点的 \lstinline{now()} 来自 \lstinline{/clock} 而非系统时钟，回放 bag 必须打开它并加 \lstinline{--clock}；消息里的 \lstinline{header.stamp} 是\textbf{传感器采集时刻}，与 \lstinline{/clock} 无关，回放不会重新打时间戳，所以录制时同步必须已经做好。检查手段是 \lstinline{ros2 bag info} 看时长与消息数、\lstinline{ros2 topic echo --field header.stamp} 对比 \lstinline{date +%s\%N}，以及回放时看 TF 是否频繁报 \lstinline{ExtrapolationException}。

\textbf{偏移 $\delta t$ 的估计。}即使时钟同步，传感器内部仍有固定延迟：相机从曝光到驱动拿到图像通常 10--50\,ms，雷达时间戳可能打在帧首或帧尾。把这类偏移统一记作 $\delta t$，含义是"测量实际对应 $t-\delta t$ 时刻的状态"。它在 $\delta t$ 内产生的位姿误差是
\begin{equation}
\delta\bm{p}\approx\bm{v}\,\delta t,\qquad \delta\bm{\theta}\approx\bm{\omega}\,\delta t.
\label{eq:26-dt}
\end{equation}
三种估计方法：(1) \textbf{一维搜索}——把点云到地图的平均残差看成 $\delta t$ 的函数，在 $[-0.1,0.1]$\,s 内粗扫，再用抛物线拟合 $J(\delta t_i)\approx J_{\min}+\tfrac12J''(\delta t_i-\hat{\delta t})^2$ 求极值；(2) \textbf{联合优化}——把 $(X,\delta t)$ 一起做非线性优化，按式~\eqref{eq:26-dt} 补偿后代入式~\eqref{eq:26-axxb} 求残差，交替优化 3--5 轮；(3) \textbf{互相关}——对雷达里程计与 IMU 积分的速度曲线做互相关，峰值位置即 $\delta t$，快但只在有明显加减速时有效。

\textbf{验证}：$\delta t$ 标好后，静止时墙厚不变（时间偏移不影响静止），快速旋转时点云不再"扇开"；把 $\delta t$ 人为置零再跑一遍，若轨迹明显变差，说明这个参数确实重要。

\section{标定验证与常见错误}\label{sec:26-verify}

标定做完不等于做完，必须验证。四条判据按顺序查，任何一条不过都不要上车：

\begin{enumerate}
  \item \textbf{重力模长}：静止时 $\norm{\tilde{\bm{a}}}=9.81\pm0.1$\,m/s$^2$；否则是加速度计尺度或零偏错。
  \item \textbf{点云墙厚}：静态场景叠加 20 帧，墙与地面厚度应 $<2$\,cm。厚度随距离增大是旋转外参误差，厚度与距离无关是平移误差或时间偏移。
  \item \textbf{回环闭合}：走矩形回到起点，位置差 $<5$\,cm、角度差 $<1\degree$。
  \item \textbf{图像叠加}：点云投影到图像，边缘对齐、无系统偏移、无随位置的扭曲。
\end{enumerate}

\begin{table}[htbp]
  \centering
  \caption{标定类故障的现象、原因与处理}\label{tab:26-errors}
  \footnotesize
  \begin{tabular}{p{2.9cm}p{3.5cm}p{3.2cm}p{4.5cm}}
    \toprule
    现象 & 可能原因 & 快速检查 & 处理 \\
    \midrule
    点云整体旋转固定角度 & 外参旋转符号/方向反 & 看正前方箱子的 $x,y,z$ & 重标或按定义改正号 \\
    点云镜像 & 坐标系手性错 & 检查是否右手系 & 改约定或加镜像修正 \\
    轨迹一直向一侧漂 & 陀螺零偏未标/符号错 & 静止 60\,s 看积分角度 & 重取零偏，查外参方向 \\
    快速旋转时点云扇开 & 时间偏移大 & 静/动对比墙厚 & 估 $\delta t$ 并补偿 \\
    静止时点云就厚 & 外参平移或雷达内参错 & 近距离叠加 20 帧 & 检查平移与外参 \\
    相机图像边缘扭曲 & 畸变系数错 & 直线在图上是否直 & 重标内参，查 $k_1$ 符号 \\
    融合后重力方向偏 & 加速度计未标 & 静止模长是否 9.81 & 六面法重标 \\
    漂移快、协方差收缩过快 & 零偏随机游走参数过小 & 看协方差曲线 & 调大 \lstinline{b_gyr_cov} \\
    \bottomrule
  \end{tabular}
\end{table}

\begin{summary}

\begin{enumerate}
  \item 旋转外参误差按 $d\theta$ 放大、时间误差按 $v\Delta t$ 放大：$1\degree$ 在 10\,m 处 17.5\,cm，10\,ms 在 2\,m/s 时 2\,cm；近处建图看不出，远处定位才暴露。
  \item 标定链有依赖顺序：内参 $\to$ IMU 噪声/零偏 $\to$ 相机-IMU $\to$ 雷达-IMU $\to$ 雷达-相机，每步都有量化判据。
  \item 张正友法：由 $\bm{H}=\lambda\bm{K}[\bm{r}_1\;\bm{r}_2\;\bm{t}]$ 与旋转列的正交单位性得 $\bm{h}_1^{\top}\bm{B}\bm{h}_2=0$、$\bm{h}_1^{\top}\bm{B}\bm{h}_1=\bm{h}_2^{\top}\bm{B}\bm{h}_2$，解 $\bm{V}\bm{b}=\bm{0}$ 得 $\bm{B}$（差一尺度），再闭式解内参，最后用重投影误差做 LM 细化。
  \item 内参判据：重投影 RMS $<0.5$\,px，$f_x\approx f_y$，主点接近图像中心。
  \item Allan 方差用相邻簇平均之差的方差的一半定义；$\tau=1$\,s 处读数即噪声密度，斜率 $-1/2$ 是随机游走、$0$ 是零偏不稳定性、$+1/2$ 是速率随机游走。
  \item 手眼标定：由 $T_{WI}=T_{WL}X$ 消去世界系得 $AX=XB$，先解四元数方程得到 $\bm{R}_X$，再解 $(\bm{R}_A-\bm{I})\bm{t}_X=\bm{R}_X\bm{t}_B-\bm{t}_A$ 得 $\bm{t}_X$；要求旋转轴张开。
  \item 时间偏移影响为 $\delta\bm{p}\approx\bm{v}\delta t$、$\delta\bm{\theta}\approx\bm{\omega}\delta t$；根本解决办法是硬件触发 + PTP 或 GPS 秒脉冲。
\end{enumerate}


\end{summary}

\begin{problemset}

  \item 某雷达与 IMU 的外参偏航角差 $1.5\degree$。估算 10\,m 处点的偏差；若只在 2\,m 内建图，这个误差会被发现吗？
  \item 从 $\bm{H}=\lambda\bm{K}[\bm{r}_1\;\bm{r}_2\;\bm{t}]$ 出发，说明 3 张不同姿态的图为何恰好能定出内参；若 3 张姿态相同，方程会怎么退化？
  \item 证明 $\bm{B}$ 只能确定到相差一个尺度，并说明式~\eqref{eq:26-intrinsic} 为什么不受尺度影响。
  \item 某陀螺 Allan 曲线 $\tau=1$\,s 处 $\sigma_A=0.02\,^{\circ}$/s，$\tau=100$\,s 处平台 $\sigma_A=2\,^{\circ}$/h。求 $N_g$（两种单位）与 $B$。
  \item 手眼标定中若机器人只绕竖直轴旋转，$X$ 的哪几个分量不可观？分别从 $\bm{M}\bm{q}_X=\bm{0}$ 的秩和 $(\bm{R}_A-\bm{I})$ 的零空间解释。
  \item 雷达时间戳打在帧尾、驱动又加了 30\,ms 传输延迟，而代码把它当作帧首采集时刻。等效 $\delta t$ 是多少？10\,Hz、2\,m/s 下误差多大？
\end{problemset}

\begin{reading}

\begin{itemize}
  \item Zhang 的原始论文 \href{https://ieeexplore.ieee.org/document/888718}{A Flexible New Technique for Camera Calibration}（PAMI 2000）：本章第~\ref{sec:26-zhang} 节的全部推导都来自这里，附录 B 给出式~\eqref{eq:26-intrinsic} 的完整形式。
  \item Sola 的 IMU 误差状态文档 \cite{sola2017eskf}：噪声建模与离散化约定讲得最清楚，适合搞不清"噪声密度"与"离散协方差"换算的人。
  \item Kalibr wiki 与 OpenCV 标定文档：前者给出相机-IMU 连续时间标定的完整流程与激励要求，后者给出畸变模型与工程细节。
  \item \href{https://github.com/rpng/imu_utils}{imu\_utils}：输入一段静止数据即输出 Allan 曲线与噪声参数，可直接用表~\ref{tab:26-allan} 的读法练手。
\end{itemize}


\end{reading}

\chapter{调参与排错手册}\label{ch:27-tuning}
\usetikzlibrary{arrows.meta}

\chapterintro{同一套 FAST-LIO，别人的车跑得又稳又准，你抄了同一份配置却在原地炸开——这是新手最挫败的场景。原因很少是"参数玄学"，而是\textbf{数据质量问题被误当成参数问题}：时间戳有 100\,ms 偏移、外参差 $2\degree$、点云频率不稳、CPU 没开优化编译。本章先给一套可复现的调参方法论（一次只改一个、固定数据集、记录指标），再按 FAST-LIO/Point-LIO/Cartographer/Nav2 四套系统列出关键参数"调大/调小"的后果，最后给一份 12 行以上的排错速查表，把"症状 $\to$ 原因 $\to$ 排查命令 $\to$ 处理"直接对应起来。本章是全书最应该放在工位上的一章。}

\section{调参方法论}\label{sec:27-method}

调参最忌讳两件事：一是同时改三个参数，然后不知道是谁起的作用；二是每换一次参数就换一段数据，结果不可比。正确的做法是一条死规矩和三个习惯。

\textbf{死规矩：一次只改一个参数，其余全部冻结。}

三个习惯：

\begin{enumerate}
  \item \textbf{固定数据集}。选定一段有代表性的 bag（包含直行、转弯、急停、回环），所有对比都在它上面做。换 bag 等于换了考卷，分数没有意义。
  \item \textbf{先定量指标，再看图}。每次实验至少记录：轨迹与真值的 ATE、终点漂移、CPU 占用、平均耗时、是否丢失跟踪。只看 RViz"感觉还行"是最容易骗自己的。
  \item \textbf{先数据质量，后算法参数}。按第~\ref{ch:26-calibration} 章的四条判据检查完标定，再动 \lstinline{filter_size_surf} 这类参数。数据错了，任何参数都调不出来。
\end{enumerate}

\begin{table}[htbp]
  \centering
  \caption{调参记录表模板（强烈建议用表格软件维护）}\label{tab:27-log}
  \small
  \begin{tabular}{lllll}
    \toprule
    编号 & 改动 & ATE (m) & 平均耗时 (ms) & 现象备注 \\
    \midrule
    0  & 基线配置 & 0.42 & 18 & 转弯处有轻微拖影 \\
    1  & \lstinline{filter_size_surf} 0.5$\to$0.3 & 0.31 & 26 & 更准但更慢 \\
    2  & 在 1 基础上 \lstinline{max_iteration} 4$\to$3 & 0.39 & 17 & 精度回落，速度变好 \\
    3  & 回到 0，改外参 $z$ $+2$\,cm & 0.44 & 18 & 无明显变化，排除平移误差 \\
    \bottomrule
  \end{tabular}
\end{table}

\begin{figure}[htbp]
  \centering
  \begin{tikzpicture}[
      box/.style={rectangle, rounded corners=2pt, draw=structurecolor, thick,
        fill=structurecolor!8, align=center, font=\footnotesize, minimum height=9mm,
        text width=36mm},
      dec/.style={rectangle, rounded corners=2pt, draw=orange!80!black, thick,
        fill=orange!12, align=center, font=\footnotesize, minimum height=9mm,
        text width=36mm},
      arr/.style={-{Stealth[length=2mm]}, thick, draw=black!70}
    ]
    \node[box] (a) at (0,0) {1. 录包，离线稳定复现};
    \node[box] (b) at (0,-1.6) {2. 查数据：频率/QoS/TF/时间戳};
    \node[dec] (c) at (0,-3.2) {3. 数据与标定是否合格？};
    \node[box] (d) at (5.4,-3.2) {4. 修数据：同步、标定、编译};
    \node[box] (e) at (0,-4.8) {5. 查参数表/排错表，只改一个参数};
    \node[box] (f) at (0,-6.4) {6. 记录指标，回归验证};
    \draw[arr] (a) -- (b);
    \draw[arr] (b) -- (c);
    \draw[arr] (c) -- node[right,font=\footnotesize]{是} (e);
    \draw[arr] (c) -- node[above,font=\footnotesize]{否} (d);
    \draw[arr] (e) -- (f);
    \draw[arr] (d.north) |- (b.east);
    \draw[arr, dashed] (f.east) -- ++(1.0,0) |- (a.east);
  \end{tikzpicture}
  \caption{调参决策流程。核心是第 3 步之前不要碰算法参数；第 6 步之后必须回到第 1 步重新验证，而不是"现场再跑一次看看"。}\label{fig:27-flow}
\end{figure}

\section{FAST-LIO / FAST-LIO2：关键参数}\label{sec:27-fastlio}

FAST-LIO2 的配置分成 \lstinline{common}、\lstinline{preprocess}、\lstinline{mapping}、\lstinline{publish} 四段\cite{xu2022fastlio2}。下面是能直接跑的骨架，之后逐项解释。

\begin{lstlisting}[language=bash,caption={FAST-LIO2 配置骨架（YAML）},label={lst:27-fastlio}]
common:
  lid_topic: "/livox/lidar"
  imu_topic: "/livox/imu"
  time_sync_en: false           # true only if driver already syncs stamps
  time_offset_lidar_to_imu: 0.0 # seconds, positive: lidar lags imu

preprocess:
  lidar_type: 1                 # 1=Livox, 2=Velodyne, 3=Ouster
  scan_line: 6
  blind: 0.5                    # ignore points closer than this (m)

mapping:
  acc_cov: 0.1                  # continuous-time accel noise
  gyr_cov: 0.1                  # continuous-time gyro noise
  b_acc_cov: 0.0001             # accel bias random walk
  b_gyr_cov: 0.0001             # gyro bias random walk
  fov_degree: 180
  det_range: 100.0
  extrinsic_est_en: false       # online extrinsic estimation
  extrinsic_T: [0.0, 0.0, 0.0]
  extrinsic_R: [1.0, 0.0, 0.0,
                0.0, 1.0, 0.0,
                0.0, 0.0, 1.0]
publish:
  path_en: true
  scan_publish_en: true
  dense_publish_en: true

filter_size_surf: 0.5           # surface downsample voxel (m)
filter_size_map: 0.5            # map voxel size (m)
cube_side_length: 1000
point_filter_num: 3
max_iteration: 4
\end{lstlisting}

\begin{table}[htbp]
  \centering
  \caption{FAST-LIO2 参数调节方向与后果}\label{tab:27-fastlio}
  \footnotesize
  \begin{tabular}{p{3.3cm}p{2.6cm}p{3.9cm}p{5.2cm}}
    \toprule
    参数 & 建议范围 & 调大的后果 & 什么时候调 \\
    \midrule
    \lstinline{extrinsic_T} & 实测 & 整体偏移、点云重影 & 用卷尺+物理检验确定，别猜 \\
    \lstinline{extrinsic_R} & 实测 & 远处点云旋转错位 & 外参标定，见第~\ref{ch:26-calibration}~章 \\
    \lstinline{filter_size_surf} & 0.2--0.5\,m & 更快但细节丢失 & 点云太稠导致算不过来时调大 \\
    \lstinline{filter_size_map} & 0.2--1.0\,m & 内存小、精度低 & 内存吃紧（长走廊）时调大 \\
    \lstinline{point_filter_num} & 1--5 & 更快、更稀疏 & CPU 打满时调到 3--4 \\
    \lstinline{max_iteration} & 3--8 & 更准、更慢，易过拟合 & 转弯重影时加到 5--6 \\
    \lstinline{acc_cov} & $10^{-2}$--$1$ & 更信 IMU & 点云退化（长走廊、隧道）时调大 \\
    \lstinline{gyr_cov} & $10^{-2}$--$1$ & 更信 IMU 姿态 & 快速旋转时姿态滞后，调大 \\
    \lstinline{b_acc_cov} & $10^{-5}$--$10^{-3}$ & 零偏更快跟随 & 温漂明显时调大 \\
    \lstinline{b_gyr_cov} & $10^{-5}$--$10^{-3}$ & 零偏更快跟随 & 长时间漂移时调大 \\
    \lstinline{cube_side_length} & 500--2000 & 地图更大、内存更高 & 场地大（>100\,m）时调大 \\
    \lstinline{det_range} & 50--200\,m & 更远处点参与匹配 & 室内场地调小可提速 \\
    \bottomrule
  \end{tabular}
\end{table}

\begin{keybox}{调 FAST-LIO 的顺序}
先外参（错了怎么调都没用）$\to$ 再 \lstinline{point_filter_num}/\lstinline{filter_size_surf}（把实时性调够）$\to$ 再 \lstinline{max_iteration}（补精度）$\to$ 最后才动 \lstinline{*_cov}（处理退化场景）。顺序反了会陷入"参数互相打架"。
\end{keybox}

\section{Point-LIO：差异参数}\label{sec:27-pointlio}

Point-LIO 的卖点是\textbf{逐点更新}：每个点都做一次状态更新，而非等一帧点云凑齐，因此在剧烈运动、低帧率雷达上更稳\cite{he2023pointlio}。它的配置文件结构与 FAST-LIO 相同，差异集中在表~\ref{tab:27-pointlio}。

\begin{table}[htbp]
  \centering
  \caption{Point-LIO 相对 FAST-LIO 的关键差异}\label{tab:27-pointlio}
  \small
  \begin{tabular}{p{3.4cm}p{2.4cm}p{6.4cm}}
    \toprule
    参数 & 典型值 & 说明 \\
    \midrule
    \lstinline{use_imu_as_input} & 0 或 1 & 两种更新模式的开关，影响鲁棒性与算力；不同版本默认值不同，以仓库示例配置为准 \\
    \lstinline{imu_satu_en} & 1 & 打开 IMU 饱和检测。撞击、弹丸命中导致的瞬时超量程不会把状态带飞 \\
    \lstinline{ivox_grid_resolution} & 0.2--0.5\,m & iVox 体素边长，越小越精细也越慢 \\
    \lstinline{ivox_nearby_type} & 18 或 26 & 近邻搜索的邻域定义，26 更稳，18 更快 \\
    \lstinline{point_filter_num} & 1--2 & 逐点更新下通常不用激进降采样 \\
    \lstinline{max_iteration} & 3--5 & 单帧内部迭代次数，与逐点更新互补 \\
    \bottomrule
  \end{tabular}
\end{table}

\begin{note}
\textbf{不要把两套配置混用}。Point-LIO 与 FAST-LIO 虽然参数名大量重叠，但协方差的量纲约定并不完全一致，把 FAST-LIO 的 YAML 直接改名喂给 Point-LIO，常见结果是"能跑但一直缓慢漂移"。正确做法是从各自仓库的示例配置出发，只改外参与话题名。
\end{note}

\begin{lstlisting}[language=bash,caption={Point-LIO 关键配置（只列与 FAST-LIO 的差异项）},label={lst:27-pointlio}]
common:
  lid_topic: "/livox/lidar"
  imu_topic: "/livox/imu"
  time_sync_en: false

mapping:
  use_imu_as_input: 0        # update mode switch, see repo README
  imu_satu_en: 1             # reject saturated imu samples (hits, impacts)
  ivox_grid_resolution: 0.3  # voxel edge length of the ivox map (m)
  ivox_nearby_type: 18       # 18 or 26 neighbours
  point_filter_num: 1
  max_iteration: 4
  extrinsic_T: [0.0, 0.0, 0.0]
  extrinsic_R: [1.0, 0.0, 0.0,
                0.0, 1.0, 0.0,
                0.0, 0.0, 1.0]
\end{lstlisting}

\section{Cartographer：2D 建图参数}\label{sec:27-carto}

Cartographer 的参数分三层：前端（\lstinline{TRAJECTORY_BUILDER_2D}）、回环（\lstinline{constraint_builder}）、后端优化（\lstinline{optimization_problem}）\cite{hess2016cartographer}。它最常出的问题是\textbf{建图糊、重影、回环失败}，三者对应的参数完全不同（表~\ref{tab:27-carto}）。

\begin{table}[htbp]
  \centering
  \caption{Cartographer 2D 关键参数与症状对应}\label{tab:27-carto}
  \footnotesize
  \begin{tabular}{p{4.2cm}p{1.9cm}p{7.4cm}}
    \toprule
    参数 & 典型值 & 作用与调节方向 \\
    \midrule
    \lstinline{submaps.resolution} & 0.05\,m & 子图分辨率。调小更精细但内存涨；建图糊可先查这个 \\
    \lstinline{submaps.num_range_data} & 60--90 & 每个子图插入的帧数。调大子图更"长"，回环更少但更快 \\
    \lstinline{TRAJECTORY_BUILDER_2D.voxel_filter_size} & 0.025--0.05\,m & 前端降采样。调大提速，过大丢细节 \\
    \lstinline{adaptive_voxel_filter.max_length} & 0.3--0.5\,m & 自适应降采样上限。走廊场景调大提速 \\
    \lstinline{min_range}/\lstinline{max_range} & 0.2/30\,m & 只使用该距离内的点。室内把 \lstinline{max_range} 调小可减少动态干扰 \\
    \lstinline{missing_data_ray_length} & 5--8\,m & 无回波方向上的自由空间长度，影响建图"通透感" \\
    \lstinline{constraint_builder.min_score} & 0.55--0.65 & 回环最小匹配分数。调高少回环但更准，调低易误回环 \\
    \lstinline{constraint_builder.sampling_ratio} & 0.3 & 回环候选采样比例。调大更准更慢 \\
    \lstinline{optimization_problem.huber_scale} & 1e-2--1e1 & 鲁棒核尺度。误回环导致地图撕裂时调小 \\
    \lstinline{optimize_every_n_nodes} & 20--90 & 每隔多少节点优化一次。调小更实时，调大更稳 \\
    \lstinline{use_imu_data} & true & 有 IMU 就打开，否则纯激光在快速旋转时会糊 \\
    \lstinline{num_accumulated_range_data} & 1--2 & 旧版本参数名，等价于每帧累积的雷达数据数 \\
    \bottomrule
  \end{tabular}
\end{table}

对应的 Lua 配置片段（Cartographer 用 \lstinline{.lua} 而不是 YAML，参数按表名组织）：

\begin{lstlisting}[language=bash,caption={Cartographer 2D 关键配置（Lua）},label={lst:27-carto}]
TRAJECTORY_BUILDER_2D = {
  use_imu_data = true,
  min_range = 0.2,
  max_range = 30.,
  missing_data_ray_length = 5.,
  voxel_filter_size = 0.025,
  adaptive_voxel_filter = { max_length = 0.5 },
  submaps = {
    num_range_data = 90,
    grid_options_2d = {
      grid_type = "PROBABILITY_GRID",
      resolution = 0.05,
    },
  },
}

POSE_GRAPH.constraint_builder.min_score = 0.6
POSE_GRAPH.constraint_builder.sampling_ratio = 0.3
POSE_GRAPH.optimization_problem.huber_scale = 1e1
POSE_GRAPH.optimize_every_n_nodes = 35
\end{lstlisting}

三个高频症状的"参数线索"：

\begin{itemize}
  \item \textbf{建图糊/墙变厚}：先看 \lstinline{use_imu_data} 是否打开、IMU 是否标定；再看 \lstinline{num_range_data} 是否过大（子图太长，前端漂移积累）。
  \item \textbf{重影（同一堵墙两条）}：典型是回环没闭上。调低 \lstinline{min_score} 到 0.5 试，若出现\textbf{错误}回环（地图撕裂）说明场景太像，应该调高 \lstinline{min_score} 并调小 \lstinline{huber_scale}。
  \item \textbf{回环失败}：确认机器人回到了历史轨迹附近（雷达看不到历史区域的点就无从回环），并检查 \lstinline{sampling_ratio} 与 \lstinline{num_range_data}。
\end{itemize}

\section{Nav2：代价地图与控制器参数}\label{sec:27-nav2}

Nav2 的参数按节点分文件：\lstinline{controller_server}、\lstinline{planner_server}、\lstinline{costmap} 等\cite{macenski2020nav2}。最容易调坏的是\textbf{代价地图}与\textbf{控制器限速}，多数"原地打转/贴墙/走不动"的现象都源于此。

\begin{lstlisting}[language=bash,caption={Nav2 参数片段（局部代价地图 + 控制器）},label={lst:27-nav2}]
local_costmap:
  local_costmap:
    ros__parameters:
      update_frequency: 10.0
      rolling_window: true
      width: 6
      height: 6
      resolution: 0.05
      robot_radius: 0.28          # or footprint polygon for non-circular
      plugins: ["voxel_layer", "inflation_layer"]
      voxel_layer:
        plugin: "nav2_costmap_2d::VoxelLayer"
        max_obstacle_height: 1.5  # ignore ceiling; keep 0.3-1.5 for RM
        mark_threshold: 0
        observation_sources: scan3d
        scan3d:
          topic: /lidar/points
          data_type: "PointCloud2"
          clearing: true
          marking: true
          max_obstacle_height: 1.2
      inflation_layer:
        plugin: "nav2_costmap_2d::InflationLayer"
        inflation_radius: 0.55    # >= robot_radius + safety margin
        cost_scaling_factor: 3.0  # smaller -> wider high-cost band

controller_server:
  ros__parameters:
    controller_frequency: 20.0
    FollowPath:
      plugin: "dwb_core::DWBLocalPlanner"
      max_vel_x: 2.5
      max_vel_theta: 3.0
      acc_lim_x: 2.0
      acc_lim_theta: 4.0
      decel_lim_x: -2.0
      xy_goal_tolerance: 0.10
      yaw_goal_tolerance: 0.10
      transform_tolerance: 0.2
      vx_samples: 20
      vtheta_samples: 40
\end{lstlisting}

\begin{table}[htbp]
  \centering
  \caption{Nav2 关键参数与故障方向}\label{tab:27-nav2}
  \footnotesize
  \begin{tabular}{p{4.2cm}p{2.4cm}p{6.8cm}}
    \toprule
    参数 & 典型值 & 偏差带来的现象 \\
    \midrule
    \lstinline{robot_radius}/\lstinline{footprint} & 实测 & 偏小会贴墙蹭、偏大到不了窄通道 \\
    \lstinline{inflation_radius} & 半径$+0.2$--$0.4$\,m & 偏小常撞、偏大"到处不可通行" \\
    \lstinline{cost_scaling_factor} & 2--5 & 偏小代价带更宽、更保守；偏大贴着障碍走 \\
    \lstinline{max_vel_x} & 1.5--4 & 偏大来不及转、偏小效率低 \\
    \lstinline{acc_lim_x}/\lstinline{acc_lim_theta} & 实测 & 偏大轨迹抖动、偏小走不动 \\
    \lstinline{xy_goal_tolerance} & 0.05--0.2\,m & 偏小原地抖、偏大停不准 \\
    \lstinline{transform_tolerance} & 0.1--0.3\,s & 偏小频繁报 TF 超时、偏大位姿滞后 \\
    \lstinline{controller_frequency} & 10--30\,Hz & 偏低调速慢、偏高吃 CPU \\
    \lstinline{max_obstacle_height} & 0.3--1.5\,m & 偏大会把天花板/悬挂物当障碍 \\
    \bottomrule
  \end{tabular}
\end{table}

\section{排错速查表}\label{sec:27-lookup}

下面这张表是本章的核心，遇到问题先按"现象"找到对应行。所有排查命令都在 ROS 2 下验证过（见第~\ref{ch:25-ros2} 章的工具使用）。

\begin{table}[htbp]
  \centering
  \caption{排错速查表：现象 $\to$ 原因 $\to$ 排查命令 $\to$ 处理}\label{tab:27-trouble}
  \footnotesize
  \begin{tabular}{p{2.5cm}p{3.0cm}p{4.3cm}p{4.2cm}}
    \toprule
    现象 & 可能原因 & 排查命令 & 处理 \\
    \midrule
    定位起飞（位姿乱跳） & 外参错、点云与 IMU 时间错、初始姿态错 & \lstinline{ros2 topic hz /lio/odom} & 重标外参；估 $\delta t$；确认初值 \\
    点云拖影（一帧里点被"拉丝"） & 时间偏移大、去畸变失效 & 静止/旋转对比墙厚 & 按第~\ref{ch:26-calibration}~章估 $\delta t$ \\
    地图重影（墙变两条） & 回环没闭上、外参差、里程计漂 & 看历史轨迹是否重叠 & 调 \lstinline{min_score}；检查外参 \\
    回环错误（地图撕裂） & \lstinline{min_score} 太低、场景太像 & 回看约束残差 & 调高 \lstinline{min_score}，调小 \lstinline{huber_scale} \\
    TF 报 LookupException & 坐标系名拼错/未发布 & \lstinline{tf2_echo map base_link} & 补 static 外参或改 frame\_id \\
    TF 报 ExtrapolationException & 时间戳乱序、缓冲区不足 & \lstinline{tf2_tools view_frames} & 修同步；\lstinline{use_sim_time} 一致 \\
    时间戳乱序/回退 & 多驱动时钟不同源 & \lstinline{ros2 topic echo --field header.stamp} & 硬件触发或统一到一台主机 \\
    CPU 打满、频率掉 & 没开 Release、降采样不足 & \lstinline{top -H}、\lstinline{ros2 topic hz} & 加 \lstinline{-DCMAKE_BUILD_TYPE=Release}；调大 \lstinline{point_filter_num} \\
    内存持续增长 & 地图/容器无限增长 & \lstinline{top} 看 RSS 曲线 & 限制 cube 数量或体素边长 \\
    IMU 数据异常（跳变/饱和） & 撞击超量程、驱动错位 & 画 \lstinline{/imu/data} 曲线 & 开饱和检测；换安装位置 \\
    静止时也在缓慢漂 & 零偏未标、温度漂移 & 静置 60\,s 看积分值 & 重取零偏；调大 \lstinline{b_gyr_cov} \\
    外参符号错（点云镜像） & 坐标系手性/方向反 & 正前方放箱子看 $x,y,z$ & 按定义改正号 \\
    UWB 定位跳变 & 多径、遮挡、基站几何差 & 看 UWB 与 LIO 差值序列 & 加抗差（Huber）+ 限幅融合 \\
    导航原地转圈 & 目标不可达、代价地图全占 & \lstinline{ros2 topic echo /plan} & 查 \lstinline{inflation_radius}；确认 \lstinline{footprint} \\
    机器人贴墙蹭 & 半径过小、膨胀层不足 & 看局部代价地图热力图 & 调大 \lstinline{inflation_radius} \\
    \bottomrule
  \end{tabular}
\end{table}

\begin{note}
\textbf{所有排查的共同前提是"能复现"}。所以第一条纪律永远是先 \lstinline{ros2 bag record -a} 录一段，然后离线用 \lstinline{ros2 bag play --clock} 复现。现场改参数再跑一次，往往只是碰巧好了，问题还在。
\end{note}

\section{日志与可视化}\label{sec:27-vis}

\begin{itemize}
  \item \textbf{RViz2 看四样东西}：点云（是否拖影、墙是否重影）、轨迹（是否闭环）、TF（是否断链）、协方差（\lstinline{Odometry} 的 covariance 是否收敛）。这四个显示项能覆盖八成问题。
  \item \textbf{\lstinline{rqt_plot}} 画 \lstinline{/cmd_vel} 与 \lstinline{/odom} 的速度曲线，看限速/加速度是否被异常裁剪。
  \item \textbf{evo} 看轨迹：\lstinline{evo_traj bag rm.bag /lio/odom --ref /gt/odom}，一眼看出漂移是"慢慢偏"还是"某处突变"。
  \item \textbf{Foxglove} 适合看长时间录包的数据：支持 3D 面板与自定义布局，比 RViz2 更适合离线分析。
  \item \textbf{日志级别}：把"每帧打印"这类调试语句降到 DEBUG，否则 I/O 本身就会让 LIO 掉帧。
\end{itemize}

\begin{keybox}{一个真实的排错顺序（值得背下来）}
\textbf{现象}：建图时墙变厚、转弯处炸开。\textbf{顺序}：录包 $\to$ 离线复现 $\to$ 查 \lstinline{ros2 topic hz}（输入频率是否正常）$\to$ 查 \lstinline{tf2_tools view_frames}（有没有断链、频率是否一致）$\to$ 查时间戳差值（点云与 IMU 的时间是否差几十 ms）$\to$ 查外参（正前方箱子）$\to$ 最后才调 \lstinline{max_iteration} 和 \lstinline{filter_size_surf}。按这个顺序，九成问题在前三步就找到了。
\end{keybox}

\begin{summary}

\begin{enumerate}
  \item 调参铁律：一次只改一个参数、固定数据集、记录定量指标、先修数据质量再动算法参数。
  \item FAST-LIO 的调参顺序：外参 $\to$ 实时性（\lstinline{point_filter_num}、\lstinline{filter_size_surf}）$\to$ 精度（\lstinline{max_iteration}）$\to$ 退化处理（\lstinline{*_cov}）。
  \item Point-LIO 的差异在逐点更新与 IMU 饱和处理（\lstinline{use_imu_as_input}、\lstinline{imu_satu_en}、iVox 参数）；不要混用两套配置。
  \item Cartographer 的"糊/重影/回环失败"分别对应前端降采样、回环阈值与子图长度，\lstinline{min_score} 与 \lstinline{huber_scale} 是回环的两个关键旋钮。
  \item Nav2 的三个高危参数是 \lstinline{robot_radius}、\lstinline{inflation_radius}、\lstinline{cost_scaling_factor}，绝大多数"走不动/贴墙/原地转"都源于此。
  \item 排错第一原则是\textbf{先录包、离线复现}；第二原则是按"频率 $\to$ TF $\to$ 时间戳 $\to$ 外参 $\to$ 参数"的顺序排查。
\end{enumerate}


\end{summary}

\begin{problemset}

  \item 你的 FAST-LIO 在长走廊里严重漂移，但转弯处正常。应该优先调 \lstinline{acc_cov}/\lstinline{gyr_cov} 还是 \lstinline{max_iteration}？说明理由。
  \item Cartographer 建图出现"墙变成两条平行线"，请给出三种可能原因，并说明每种原因下应该看什么数据来区分。
  \item Nav2 机器人在离目标 5\,cm 处原地抖动，列出至少三个相关参数并说明调整方向。
  \item 为什么"一次只改一个参数"在调 LIO 时尤其重要？请结合参数之间的耦合（如降采样与迭代次数）说明。
  \item 设计一个 30 分钟的排错流程，用于判断"定位偶尔跳 30\,cm"是时间同步问题还是重定位问题。
\end{problemset}

\begin{reading}

\begin{itemize}
  \item FAST-LIO2 \cite{xu2022fastlio2} 与 Point-LIO \cite{he2023pointlio} 的原始论文：前者解释 iKD-Tree 与迭代 EKF，后者解释逐点更新为什么能提高带宽，读完再调 \lstinline{max_iteration} 会更有方向。
  \item Cartographer 论文 \cite{hess2016cartographer} 与官方 tuning 文档：子图、分支定界回环、后端优化的关系讲得最清楚。
  \item Nav2 官方文档的 \href{https://docs.nav2.org/configuration/index.html}{Configuration Guide}：每个参数的取值范围与单位，配合第~\ref{ch:22-nav2-practice} 章的实战链路使用。
  \item 第~\ref{ch:17-advanced-lio} 章：FAST-LIO2 与 Point-LIO 的内部实现细节，是理解本章参数含义的前提。
\end{itemize}


\end{reading}

\chapter{数据集、评测与复现}\label{ch:28-evaluation}
\usetikzlibrary{arrows.meta}

\chapterintro{"我的算法比 FAST-LIO 好"——这句话在组会上说出来之前，必须先回答三个问题：在哪个数据集上、用什么指标、对齐方式是什么。否则它只是一个无法验证的感觉。本章解决的就是"怎么把感觉变成数字"：先列出定位领域最常用的公开数据集与各自适合评什么，再给出 ATE/RPE 的严格定义与 Umeyama 对齐闭式解（这是所有评测工具背后真正在算的东西），然后用 \lstinline{evo} 把命令和输出对上号，最后讲自采数据的规范与可复现实验的设计。评测能力决定了你能否在一堆参数里做出正确判断，也决定了你的实验别人能不能复现。}

\section{公开数据集}\label{sec:28-datasets}

选数据集的原则是：\textbf{它必须能暴露你想改进的那种失效模式}。在 KITTI 上刷高分的算法，遇到 RM 场地的动态障碍与对称结构可能一塌糊涂；反之只在自家场地上调参，也无法证明方法本身有效。表~\ref{tab:28-datasets} 是常用的几套。

\begin{table}[htbp]
  \centering
  \caption{常用公开数据集对比}\label{tab:28-datasets}
  \footnotesize
  \begin{tabular}{p{2.3cm}p{4.2cm}p{3.4cm}p{4.0cm}}
    \toprule
    数据集 & 传感器 & 场景 & 适合评什么 \\
    \midrule
    KITTI Odometry & 双目、HDL-64、GPS/IMU & 城郊道路、高速 & 激光/视觉里程计的 ATE，室外长距离 \\
    EuRoC MAV & 双目、ADIS16448、Vicon/Leica & 室内房间与厂房 & VIO 的 ATE/RPE，标定与时间同步 \\
    TUM VI & 双目鱼眼、IMU & 室内外混合 & 视觉惯性，含相机标定基准 \\
    NTU VIRAL & Ouster、双目、IMU、UWB & 校园、室内外 & 多传感器融合、UWB 辅助定位 \\
    Newer College & Ouster、相机、IMU & 校园、含季节变化 & 长期定位、鲁棒性 \\
    M2DGR & 鱼眼/双目/事件相机、激光、IMU、GNSS & 校园地面机器人 & 多模态融合、地面平台 \\
    RM 自采 & Livox、IMU、轮式 & 比赛场地 & 你的真实失效模式 \\
    \bottomrule
  \end{tabular}
\end{table}

几点使用建议：

\begin{enumerate}
  \item \textbf{先用小数据集跑通闭环}。EuRoC 的 MH\_01 只有 3 分钟，适合快速验证编译与配置；不要在 KITTI 上调试程序。
  \item \textbf{真值要看清来源}。Vicon/动捕是 mm 级但只覆盖小场地；RTK 是 cm 级但室内失效；KITTI 的 GPS/IMU 融合真值本身也有误差。\textbf{评测精度的下限是真值精度}。
  \item \textbf{注意传感器差异}。Livox 非重复扫描的点云分布与 Velodyne 机械旋转完全不同，在一套数据上调好的 \lstinline{filter_size_surf} 换到另一套上必须重调。
  \item \textbf{自采数据不可替代}。公开数据集用于"证明方法有效"，自采数据用于"证明在 RM 场上能用"，两者都要有。
\end{enumerate}

\section{评价指标：ATE、RPE 与漂移率}\label{sec:28-metrics}

\subsection{ATE：绝对轨迹误差}

\textbf{ATE}（Absolute Trajectory Error）衡量估计轨迹与真值轨迹的整体偏差。设有 $N$ 个对应时刻，估计位姿 $\hat{T}_i$、真值位姿 $T_i$，最简单的定义是逐点位置误差
\begin{equation}
\bm{e}_i=\mathrm{trans}\bigl(\hat{T}_i\bigr)-\mathrm{trans}\bigl(T_i\bigr),
\qquad
\mathrm{ATE}_{\mathrm{RMSE}}=\sqrt{\frac{1}{N}\sum_{i=1}^{N}\norm{\bm{e}_i}^{2}},
\label{eq:28-ate}
\end{equation}
其中 $\mathrm{trans}(\cdot)$ 取位姿的平移部分。但这样直接比是\textbf{不公平的}：两个轨迹可能只差一个初始位姿（起点朝向不同）或一个全局尺度（单目 SLAM），这不代表算法差。所以必须先做\textbf{对齐}，把估计轨迹用一个刚体变换（有时还允许缩放）配到真值上，再算误差。这就是常说的 "ATE with SE(3) alignment"。

\subsection{Umeyama 对齐的闭式解}

对齐问题：给定两组三维点 $\{\bm{x}_i\}$（估计）与 $\{\bm{y}_i\}$（真值），$i=1,\dots,N$，求尺度 $s$、旋转 $\bm{R}$、平移 $\bm{t}$，使
\begin{equation}
J(s,\bm{R},\bm{t})=\frac{1}{N}\sum_{i=1}^{N}\norm{\bm{y}_i-\bigl(s\bm{R}\bm{x}_i+\bm{t}\bigr)}^{2}
\label{eq:28-umeyama}
\end{equation}
最小。下面推导闭式解（Umeyama 1991）。

\textbf{第一步：消去平移。} 对 $\bm{t}$ 求偏导并令其为零：
\begin{equation}
\frac{\partial J}{\partial\bm{t}}=-\frac{2}{N}\sum_{i=1}^{N}\bigl(\bm{y}_i-s\bm{R}\bm{x}_i-\bm{t}\bigr)=\bm{0}
\;\Longrightarrow\;
\bm{t}=\bm{\mu}_y-s\bm{R}\bm{\mu}_x,
\label{eq:28-t}
\end{equation}
其中 $\bm{\mu}_x=\frac{1}{N}\sum_i\bm{x}_i$、$\bm{\mu}_y=\frac{1}{N}\sum_i\bm{y}_i$。记中心化坐标 $\bm{x}'_i=\bm{x}_i-\bm{\mu}_x$、$\bm{y}'_i=\bm{y}_i-\bm{\mu}_y$，代回 $J$ 得
\begin{equation}
J=\frac{1}{N}\sum_{i=1}^{N}\norm{\bm{y}'_i-s\bm{R}\bm{x}'_i}^{2}
=\frac{1}{N}\sum_{i=1}^{N}\norm{\bm{y}'_i}^{2}-\frac{2s}{N}\sum_{i=1}^{N}{\bm{y}'_i}^{\top}\bm{R}\bm{x}'_i+\frac{s^{2}}{N}\sum_{i=1}^{N}\norm{\bm{x}'_i}^{2},
\label{eq:28-expand}
\end{equation}
中间一步用了 $\norm{\bm{a}-\bm{b}}^2=\norm{\bm{a}}^2-2\bm{a}^{\top}\bm{b}+\norm{\bm{b}}^2$。记 $\sigma_x^{2}=\frac{1}{N}\sum_i\norm{\bm{x}'_i}^{2}$、$\sigma_y^{2}=\frac{1}{N}\sum_i\norm{\bm{y}'_i}^{2}$，并定义 $3\times3$ 的\textbf{互协方差矩阵}
\begin{equation}
\bm{\Sigma}=\frac{1}{N}\sum_{i=1}^{N}\bm{x}'_i{\bm{y}'_i}^{\top}.
\label{eq:28-sigma}
\end{equation}
利用迹的性质 $\tr(\bm{R}\bm{\Sigma})=\frac{1}{N}\sum_i{\bm{y}'_i}^{\top}\bm{R}\bm{x}'_i$（把 $\bm{\Sigma}$ 展开后交换求和与求迹即可验证），式~\eqref{eq:28-expand} 变成
\begin{equation}
J=\sigma_y^{2}-2s\,\tr(\bm{R}\bm{\Sigma})+s^{2}\sigma_x^{2}.
\label{eq:28-J}
\end{equation}

\textbf{第二步：先定旋转。} 式~\eqref{eq:28-J} 中只有中间一项与 $\bm{R}$ 有关，所以固定 $s$ 时最小化 $J$ 等价于最大化 $\tr(\bm{R}\bm{\Sigma})$。对 $\bm{\Sigma}$ 做奇异值分解 $\bm{\Sigma}=\bm{U}\bm{D}\bm{V}^{\top}$，$\bm{D}=\diag(d_1,d_2,d_3)$，$d_1\ge d_2\ge d_3\ge0$。代入：
\begin{equation}
\tr(\bm{R}\bm{\Sigma})=\tr\bigl(\bm{R}\bm{U}\bm{D}\bm{V}^{\top}\bigr)=\tr\bigl(\bm{V}^{\top}\bm{R}\bm{U}\bm{D}\bigr)=\sum_{k=1}^{3}d_k\bigl(\bm{V}^{\top}\bm{R}\bm{U}\bigr)_{kk}.
\end{equation}
由于 $\bm{V}^{\top}\bm{R}\bm{U}$ 是正交矩阵，其对角元不超过 1，故上界是 $\sum_k d_k$，取等当且仅当 $\bm{V}^{\top}\bm{R}\bm{U}=\bm{I}$，即 $\bm{R}=\bm{V}\bm{U}^{\top}$。但正交矩阵还可能是反射（$\det=-1$），而旋转要求 $\det=1$，于是引入修正矩阵
\begin{equation}
\bm{S}=\diag\bigl(1,\,1,\,\det(\bm{V}\bm{U}^{\top})\bigr),\qquad
\bm{R}=\bm{V}\bm{S}\bm{U}^{\top}.
\label{eq:28-R}
\end{equation}
这样 $\det(\bm{R})=1$ 且仍取到最大 $\tr(\bm{R}\bm{\Sigma})=\sum_k d_k s_k$（$s_k$ 为 $\bm{S}$ 的对角元）。

\textbf{第三步：再定尺度。} 把 $\bm{R}$ 代回式~\eqref{eq:28-J}，对 $s$ 求导：
\begin{equation}
\frac{\partial J}{\partial s}=-2\tr(\bm{R}\bm{\Sigma})+2s\sigma_x^{2}=0
\;\Longrightarrow\;
s=\frac{\tr(\bm{R}\bm{\Sigma})}{\sigma_x^{2}}=\frac{\sum_{k=1}^{3}s_kd_k}{\sigma_x^{2}}.
\label{eq:28-s}
\end{equation}
\textbf{第四步：}平移由式~\eqref{eq:28-t} 给出：$\bm{t}=\bm{\mu}_y-s\bm{R}\bm{\mu}_x$。

\begin{keybox}{评测时的三个约定}
\begin{enumerate}
  \item \textbf{是否估计尺度}：单目 SLAM 必须估计 $s$（否则轨迹差一个倍数）；激光/激光惯性系统通常固定 $s=1$，只在刚体变换下对齐。
  \item \textbf{对齐范围}：是"整条轨迹对齐一次"（全局对齐）还是"逐段对齐"，结果差很多。\lstinline{evo} 默认全局对齐。
  \item \textbf{时间对齐}：估计与真值的时间戳必须先按时间关联（\lstinline{evo} 会自动关联，但存在偏移时要用 \lstinline{--t_max_diff} 或先手动补偿）。
\end{enumerate}
\end{keybox}

\subsection{RPE：相对位姿误差}

\textbf{RPE}（Relative Pose Error）衡量"固定间隔内的相对运动"的误差，定义在成对的位姿上。取间隔 $\Delta$（按帧数、时间 1\,s 或距离 1\,m），第 $i$ 对的误差为
\begin{equation}
\bm{E}_i=\bigl(T_i^{-1}T_{i+\Delta}\bigr)^{-1}\bigl(\hat{T}_i^{-1}\hat{T}_{i+\Delta}\bigr),
\label{eq:28-rpe}
\end{equation}
即"真值的相对运动"与"估计的相对运动"之差的变换。分解成平移与旋转两部分：
\begin{equation}
\mathrm{RPE}_{\mathrm{trans},i}=\norm{\mathrm{trans}(\bm{E}_i)},\qquad
\mathrm{RPE}_{\mathrm{rot},i}=\arccos\!\left(\frac{\tr\bigl(\mathrm{rot}(\bm{E}_i)\bigr)-1}{2}\right),
\end{equation}
再对 $i$ 取 RMSE。$\arccos$ 那一项是旋转矩阵的\textbf{迹公式}：绕单位轴转 $\theta$ 的旋转矩阵迹为 $1+2\cos\theta$，反解即得旋转角。

\textbf{为什么里程计看 RPE、SLAM 看 ATE？}因为两者关心的东西不同：里程计的核心是"每走一步的增量准不准"，增量误差会累积成漂移，RPE 正好测增量；SLAM 有回环与后端优化，单步可能不准但整体轨迹闭合得很好，此时 RPE 会很难看而 ATE 很好。所以报指标时通常两个都给：ATE 说整体精度，RPE 说局部一致性。

\begin{figure}[htbp]
  \centering
  \begin{tikzpicture}[x=1.05cm,y=1.05cm,font=\footnotesize]
    % ground truth (solid) and estimate (dashed), offset by a constant
    \draw[thick, black!70] plot[smooth] coordinates
      {(0,0) (1.2,0.6) (2.4,1.0) (3.6,0.8) (4.8,1.3) (6.0,1.9) (7.2,2.2)};
    \draw[thick, structurecolor, dashed] plot[smooth] coordinates
      {(0,0.55) (1.2,1.15) (2.4,1.55) (3.6,1.35) (4.8,1.85) (6.0,2.45) (7.2,2.75)};
    % ATE: point-wise error between corresponding poses
    \foreach \x/\ya/\yb in
      {0/0/0.55, 1.2/0.6/1.15, 2.4/1.0/1.55, 3.6/0.8/1.35,
       4.8/1.3/1.85, 6.0/1.9/2.45, 7.2/2.2/2.75}
      \draw[red!60!black, thick] (\x,\ya) -- (\x,\yb);
    % RPE: relative motion between pose i and pose i+Delta
    \draw[very thick, orange!80!black] (1.2,0.6) -- (3.6,0.8);
    \draw[very thick, orange!80!black, dashed] (1.2,1.15) -- (3.6,1.35);
    \foreach \x/\y in {1.2/0.6, 3.6/0.8, 1.2/1.15, 3.6/1.35}
      \fill (\x,\y) circle (1.6pt);
    % legend
    \node[anchor=west, black!70] at (7.25,2.2) {真值};
    \node[anchor=west, structurecolor] at (7.25,2.75) {估计};
    \node[anchor=north west, font=\footnotesize] at (0.05,-0.12)
      {红：ATE（逐点位置误差）};
    \node[anchor=north west, font=\footnotesize] at (0.05,-0.65)
      {橙：RPE（$\Delta$ 内相对运动的误差）};
  \end{tikzpicture}
  \caption{ATE 与 RPE 的区别。ATE 比较同一时刻的绝对位置（红色短线段），RPE 比较固定间隔 $\Delta$ 内的相对运动（橙色线段：真值一条、估计一条，二者之差才是误差）。}\label{fig:28-ate}
\end{figure}

\subsection{漂移率与工程指标}

\textbf{漂移率}把 ATE 归一化，便于跨数据集比较：
\begin{equation}
\text{drift}=\frac{\mathrm{ATE}_{\mathrm{RMSE}}}{\text{轨迹长度}}\times100\%.
\label{eq:28-drift}
\end{equation}
$1$\,km 轨迹 ATE 1\,m 是 $0.1\%$，属于很好的水平；室内 50\,m 轨迹 0.5\,m 就是 $1\%$，需要检查。

除精度外，工程上还必须报三类指标：

\begin{itemize}
  \item \textbf{实时性}：单帧平均耗时、$95$ 分位耗时、是否达到雷达频率。只报平均耗时会掩盖偶发的卡顿。
  \item \textbf{资源}：CPU 单核占用、内存峰值（RSS）、地图内存增长速度。
  \item \textbf{鲁棒性}：连续运行 $N$ 次/小时的失败次数、丢失跟踪后的恢复时间。
\end{itemize}

\section{evo：命令与输出解读}\label{sec:28-evo}

\lstinline{evo} 是评测的事实标准工具，支持 TUM/KITTI/bag 等多种格式。四条最常用的命令：

\begin{lstlisting}[language=bash,caption={evo 常用命令},label={lst:28-evo}]
# 1) inspect two trajectories: align, then plot side by side
evo_traj tum gt.tum --ref est.tum --align --plot

# 2) absolute trajectory error (RMSE), with alignment
evo_ape tum gt.tum est.tum -va --align --plot
# add --correct_scale when evaluating monocular SLAM

# 3) relative pose error over 1 m of travel
evo_rpe kitti gt.txt est.txt --delta 1 --delta_unit m -va

# 4) compare many runs: each result is a zip saved by evo_ape/evo_rpe
evo_res results/*.zip -p --use_fig --save_table table.csv
\end{lstlisting}

输出中最该看的是 \lstinline{RMSE}、\lstinline{median}、\lstinline{std} 三个数：RMSE 对大误差敏感（平方后再平均），median 反映典型水平，std 反映抖动。若 median 好而 RMSE 差，说明轨迹中\textbf{偶有跳变}；若两者都差，说明是系统性偏移。\lstinline{--align} 做的是第~\ref{sec:28-metrics} 节的 Umeyama 对齐（不含尺度），\lstinline{--correct_scale} 才估计尺度。

\begin{note}
\textbf{报数字时的三个必须}：必须写对齐方式（有无尺度）、必须写数据段（哪几段序列）、必须写运行次数（取中位数还是最好的一次）。只写"ATE 0.12"是不可复现的。
\end{note}

\section{自采数据规范}\label{sec:28-selfdata}

RM 场地的数据是公开数据集替代不了的，但自采数据最容易烂在文件夹里。规范如下：

\begin{enumerate}
  \item \textbf{时间同步先于录制}。录之前用第~\ref{ch:26-calibration} 章的方法确认两路时间戳一致，并记录同步方式（硬件触发/PTP/NTP）。
  \item \textbf{标定文件与数据同存}。每次采数据时把当次的相机内参、外参、IMU 噪声参数、激光雷达型号配置与 bag 放在同一目录，文件名带版本号。
  \item \textbf{真值获取}：动捕精度最高但场地受限；全站仪适合静态标定点；室外用 RTK；也可以先用"高精度先验地图 + ICP"得到准真值，但要注明是"准真值"。
  \item \textbf{命名与元数据}：建议 \lstinline{日期_场地_平台_传感器_序号}，并写一个 README 记录：谁录的、什么配置、场地描述、已知问题（比如"这段有人走过"）。
  \item \textbf{不要只录成功片段}。失败的数据最有价值——炸的那 10 秒决定了你能不能复现并修好问题。
\end{enumerate}

\section{可复现实验}\label{sec:28-repro}

一个实验能被别人（或三个月后的你自己）复现，必须固定四件事：

\begin{enumerate}
  \item \textbf{代码状态}：记录 git commit hash；有未提交改动时用 \lstinline{git stash} 或打 tag。
  \item \textbf{配置快照}：把当次使用的 YAML 与 launch 文件复制到结果目录（不要只记"改了 \lstinline{filter_size_surf}"）。
  \item \textbf{随机性}：ICP/采样类算法、随机种子、线程调度都可能带来波动。固定随机种子，并\textbf{至少跑 3 次取中位数}；差异小于波动幅度的改进不能算改进。
  \item \textbf{环境}：CPU 型号、是否开 Release、是否绑核、ROS 版本、bag 回放速度（是否 \lstinline{--rate 1.0}）。
\end{enumerate}

做改进对比时用消融表（表~\ref{tab:28-ablation}），每一行只允许一个变量不同，并用中位数与波动范围一起报。

\begin{table}[htbp]
  \centering
  \caption{消融实验表模板（每个配置跑 3 次，报中位数$\pm$半极差）}\label{tab:28-ablation}
  \small
  \begin{tabular}{llccc}
    \toprule
    配置 & 改动 & ATE (m) & RPE (m) & 平均耗时 (ms) \\
    \midrule
    A & 基线 & $0.42\pm0.03$ & $0.021$ & 18 \\
    B & A + iKD-Tree 体素 0.5\,m & $0.40\pm0.02$ & $0.020$ & 14 \\
    C & B + 时间偏移补偿 & $0.29\pm0.01$ & $0.014$ & 14 \\
    D & C + \lstinline{max_iteration} 6 & $0.27\pm0.02$ & $0.013$ & 21 \\
    \bottomrule
  \end{tabular}
\end{table}

\begin{note}
\textbf{不要用"最好一次"的结果参与比较}。算法的真实表现是分布，不是单点。如果两个配置的差异落进了多次运行的波动范围，就老老实实说"无显著差异"，这比编一个漂亮的数字更专业，也更省时间。
\end{note}

\begin{summary}

\begin{enumerate}
  \item 数据集要选能暴露目标失效模式的：公开数据集证明方法有效，自采数据证明在 RM 场上能用，两者都要。
  \item ATE 是逐点绝对误差的 RMSE，通常先做 Umeyama 对齐再算；单目系统要估计尺度，激光系统固定 $s=1$。
  \item Umeyama 闭式解四步：中心化求 $\bm{\Sigma}=\frac1N\sum\bm{x}'{\bm{y}'}^{\top}$ $\to$ SVD 取 $\bm{R}=\bm{V}\bm{S}\bm{U}^{\top}$（$\bm{S}$ 修正反射）$\to$ $s=\tr(\bm{R}\bm{\Sigma})/\sigma_x^{2}$ $\to$ $\bm{t}=\bm{\mu}_y-s\bm{R}\bm{\mu}_x$。
  \item RPE 测固定间隔 $\Delta$ 内的相对运动误差，是里程计的核心指标；ATE 与 RPE 一起报才完整。
  \item 漂移率 $=\mathrm{ATE}/\text{轨迹长度}$，便于跨场景比较；工程上还要报耗时分布、内存与失败率。
  \item \lstinline{evo} 四条命令覆盖检查、ATE、RPE、多结果对比；报数字必须带对齐方式、数据段与运行次数。
  \item 可复现四要素：代码 commit、配置快照、随机性控制（3 次取中位数）、环境记录。
\end{enumerate}


\end{summary}

\begin{problemset}

  \item 估计轨迹与真值只差一个初始朝向，不做对齐的 ATE 会很大。请说明这种"不公平"在什么情况下不能通过刚体对齐消除（提示：想想单目尺度）。
  \item 从式~\eqref{eq:28-expand} 出发，补全 $\tr(\bm{R}\bm{\Sigma})$ 与逐点求和等价的推导步骤（允许用"展开后交换求和与求迹"）。
  \item 为什么 $\bm{S}$ 的第三个对角元要取 $\det(\bm{V}\bm{U}^{\top})$？如果直接取 $\bm{R}=\bm{V}\bm{U}^{\top}$ 会发生什么？
  \item 某算法 ATE 的 median 是 0.05\,m 而 RMSE 是 0.6\,m。请解释这两个数为什么会差这么多，并说明你会去查什么。
  \item 你要比较两个参数配置的优劣，两者 ATE 相差 3\%，而 5 次运行的波动范围是 $\pm5\%$。应该怎么下结论？
\end{problemset}

\begin{reading}

\begin{itemize}
  \item Umeyama 的原始论文 \href{https://ieeexplore.ieee.org/document/88573}{Least-Squares Estimation of Transformation Parameters Between Two Point Patterns}（PAMI 1991）：本章对齐闭式解的全部推导来源。
  \item \lstinline{evo} 官方文档 \href{https://michaelgrupp.github.io/evo/}{michaelgrupp.github.io/evo}：所有命令行参数与输出含义，配合本章第~\ref{sec:28-evo} 节的四条命令即可上手。
  \item KITTI Odometry 与 EuRoC 官方页面：真实数据、真值格式与评测脚本的权威说明，建议在跑通自己的系统后读它们的评测协议。
  \item 第~\ref{ch:27-tuning} 章的调参记录表：评测结果最终要落到那张表上才能形成闭环。
\end{itemize}


\end{reading}

\chapter{RoboMaster 场景实践建议}\label{ch:29-rm}
\usetikzlibrary{arrows.meta}

\chapterintro{前面各章讲的是"通用技术"，本章只回答一个问题：\textbf{在自己的车上、自己的场地里该怎么选、怎么装、怎么验}。RM 赛场有几个和实验室完全不同的约束：场地对称、动态目标多、弹丸与撞击干扰传感器、算力与功耗被严格限制、比赛期间没有时间调参。因此本章写成一份"实战清单"——方案选型表、各兵种需求对比、坐标系与标定清单、上车检查表、备赛流程与验收标准、以及几个最容易踩的坑。它不展开成比赛教程，只给你决策时需要的判断依据。}

\section{赛场环境特点}\label{sec:29-env}

以近年 $28\,\mathrm{m}\times15\,\mathrm{m}$ 的标准场地为例，定位系统要面对的环境特征有六条：

\begin{enumerate}
  \item \textbf{高度对称}。场地两侧结构几乎镜像，中央高地、环形通道、基地位置对称。对全局定位（先验地图匹配）来说这是"歧义源"：局部特征几乎一样，容易匹配到镜像位置。
  \item \textbf{动态障碍多}。对方机器人（$3$--$4$\,m/s）、己方队友、飞舞的弹丸都在点云里。弹丸对激光雷达是"瞬时噪声"，对视觉是"遮挡+光斑"。
  \item \textbf{结构重复且低矮}。围栏、装甲板、基地支架都是同一种几何，垂直方向的纹理少，容易出现"走廊式退化"。
  \item \textbf{光照变化剧烈}。顶灯、装甲板 LED、能量机关的强光、以及比赛期间的明暗切换，对相机是灾难，对激光雷达影响小。
  \item \textbf{地面材质与平整度}。场地地胶、坡道、上下坡的过渡区域会让轮式里程计打滑。
  \item \textbf{观众与人员干扰}。场边走动的人会被 VIO 当成运动目标，也曾出现过观众入场导致的定位跳变。
\end{enumerate}

这些特征直接决定了技术选型：\textbf{以 3D 激光惯性里程计为主，用先验地图做全局校正，用 UWB 做兜底}，而不是反过来。

\section{定位方案选型}\label{sec:29-scheme}

\begin{table}[htbp]
  \centering
  \caption{定位方案对比（按 RM 场景的实际表现）}\label{tab:29-scheme}
  \footnotesize
  \begin{tabular}{p{2.6cm}p{4.4cm}p{4.6cm}p{2.4cm}}
    \toprule
    方案 & 优点 & 主要风险 & 建议 \\
    \midrule
    3D LIO（FAST-LIO2 / Point-LIO） & 精度 2--5\,cm，不依赖场地改造，输出频率高 & 退化场景（长走廊、空旷区）漂移；动态目标进地图 & 主方案，必装 \\
    先验地图 + 全局定位 & 消除累积漂移，冷启动可重定位 & 地图过期、场地布置变化、对称场景误匹配 & 主方案，每场前更新地图 \\
    UWB & 全局绝对、无累积误差、可跨区 & 多径、金属反射、非视距、安装基站耗时 & 辅助与兜底 \\
    雷达站全局感知 & 上帝视角，可给全队共享位置 & 标定复杂、通信延迟、遮挡 & 有条件就上 \\
    轮式 + IMU 航迹推算 & 便宜、连续、不受视觉影响 & 打滑、累积漂移快 & 只作为短时补充 \\
    \bottomrule
  \end{tabular}
\end{table}

选型的原则是\textbf{异构冗余}：激光退化时靠 UWB 兜底，UWB 跳变时靠激光与先验地图压住，先验地图失效时至少还有连续里程计能让车停下来而不是撞墙。表~\ref{tab:29-scheme} 中任何单一方案都不足以支撑一场比赛。

\section{哨兵、步兵与雷达站的差异}\label{sec:29-platform}

三个兵种对定位系统的需求完全不同（表~\ref{tab:29-platform}），用同一套参数配置是常见错误。

\begin{table}[htbp]
  \centering
  \caption{三个兵种的定位需求对比}\label{tab:29-platform}
  \footnotesize
  \begin{tabular}{p{2.0cm}p{3.2cm}p{3.2cm}p{3.2cm}}
    \toprule
    维度 & 哨兵 & 步兵 & 雷达站 \\
    \midrule
    平台 & 悬挂/滑轨，云台可 $360\degree$ 转 & 全向底盘，可冲坡 & 静止，可升降 \\
    速度 & 平移 2--3\,m/s & 冲刺 4\,m/s 以上 & 0 \\
    算力预算 & 高（Jetson 级） & 中（受限功耗） & 最高（可上台式机） \\
    动态性 & 云台转动使雷达视角变化 & 最剧烈：撞击、跳跃、急停 & 几乎静止 \\
    定位重点 & 底盘稳定 + 云台姿态解耦 & 高速下短时精度与抗打滑 & 全局精度与多目标跟踪 \\
    关键风险 & 云台转动导致 TF 抖动 & 碰撞使 IMU 饱和 & 与车载机的通信延迟 \\
    \bottomrule
  \end{tabular}
\end{table}

具体到实现：

\begin{itemize}
  \item \textbf{哨兵}：云台与底盘必须解耦。底盘定位用 \lstinline{map}$\to$\lstinline{odom}$\to$\lstinline{base\_link}，云台角由编码器反馈单独发布，不要试图用雷达去估计云台角（云台转动时雷达点云在动，容易被误当成底盘在转）。
  \item \textbf{步兵}：撞击与急停会让 IMU 瞬时超量程。Point-LIO 的 \lstinline{imu_satu_en} 与合理的零偏随机游走参数是必备的；同时要限制单帧修正量，避免定位跳变导致规划器输出暴力动作。
  \item \textbf{雷达站}：几乎静止，重点在全局精度与多目标跟踪，而不是高频里程计。可以把算力全部投给全局匹配与多机器人状态估计，并把结果通过低延迟链路广播给队友。
\end{itemize}

\section{坐标系与标定清单}\label{sec:29-frames}

\begin{figure}[htbp]
  \centering
  \begin{tikzpicture}[
      frame/.style={rectangle, rounded corners=2pt, draw=structurecolor, thick,
        fill=structurecolor!8, minimum height=7mm, minimum width=19mm, font=\small},
      lab/.style={font=\footnotesize, inner sep=1.5pt, draw=black!60, rounded corners=1pt,
        fill=white},
      arr/.style={-{Stealth[length=2mm]}, thick, draw=black!70}
    ]
    \node[frame] (map) at (0,0) {map};
    \node[frame] (odom) at (3.4,0) {odom};
    \node[frame] (base) at (6.8,0) {base\_link};
    \node[frame] (gim) at (10.2,0) {gimbal};
    \node[frame] (lidar) at (13.4,0.9) {lidar};
    \node[frame] (cam) at (13.4,-0.9) {camera};
    \node[frame, text width=20mm] (chassis) at (6.8,-1.9) {轮式/IMU};
    \draw[arr] (map) -- node[above=1pt,lab]{全局定位} (odom);
    \draw[arr] (odom) -- node[above=1pt,lab]{LIO} (base);
    \draw[arr] (base) -- node[above=1pt,lab]{云台反馈} (gim);
    \draw[arr] (gim) -- node[above=1pt,lab]{外参} (lidar);
    \draw[arr] (gim) -- node[below=1pt,lab]{外参} (cam);
    \draw[arr] (chassis) -- node[right=1pt,lab]{里程计} (base);
  \end{tikzpicture}
  \caption{RM 机器人的坐标系与模块关系。注意云台是独立的一节：底盘到云台由编码器反馈，云台到传感器是标定出的常数外参。}\label{fig:29-frames}
\end{figure}

标定顺序必须是"从底盘往外扩"：\textbf{底盘 $\to$ 雷达 $\to$ 云台 $\to$ 相机内参 $\to$ 云台-相机外参 $\to$ UWB 坐标系}。每一环都依赖上一环，跳步会导致误差叠加且难以定位。

\begin{keybox}{上车检查表（10 条，每次上车都过一遍）}
\begin{enumerate}
  \item IMU 静止 60\,s，重力模长 $9.81\pm0.1$\,m/s$^2$，三轴零偏稳定。
  \item 正前方 5\,m 放箱子，点云坐标 $x\approx5$、$y\approx0$、$z\approx0$，误差 $<5$\,cm。
  \item \lstinline{tf2_tools view_frames} 无断链，\lstinline{map}$\to$\lstinline{odom} 与 \lstinline{odom}$\to$\lstinline{base\_link} 的发布频率正常。
  \item 点云与 IMU 时间戳差值 $<5$\,ms（回放 bag 时检查）。
  \item 静止叠加 20 帧，墙厚 $<2$\,cm。
  \item 云台零位与底盘前向一致，偏差 $<0.5\degree$；转一圈云台后底盘位姿不跳。
  \item UWB 标定点与实际坐标一致，静态抖动 $<10$\,cm。
  \item 先验地图与当前场地一致（对照场地尺寸与标志物）。
  \item 冷启动重定位：10 次至少 9 次成功，耗时 $<10$\,s。
  \item 连续运行 30\,min，无内存持续增长、无 TF 超时、无频率下降。
\end{enumerate}
\end{keybox}

\section{备赛流程与验收标准}\label{sec:29-process}

备赛是一条流水线，每一步都有可验收的输出，没通过就不要进下一步。

\begin{table}[htbp]
  \centering
  \caption{备赛流程与验收标准}\label{tab:29-process}
  \footnotesize
  \begin{tabular}{p{2.4cm}p{5.6cm}p{5.6cm}}
    \toprule
    阶段 & 要做的事 & 验收标准 \\
    \midrule
    建图 & 场地空场时多圈扫描，含高地、坡道、边界 & 回环闭合误差 $<10$\,cm，无明显重影 \\
    地图清理 & 去掉动态物体、修复空洞、统一分辨率 & 地图文件可被定位节点加载且帧率达标 \\
    先验地图包 & 建立 \lstinline{map} 与场地坐标系的对齐 & 标定点误差 $<5$\,cm \\
    定位验证 & 静置、慢走、快跑、旋转四组测试 & 位置误差 $<10$\,cm，姿态 $<1\degree$ \\
    导航联调 & 与规划、控制、行为树联跑 & 10 次导航任务全部到达，无碰撞 \\
    压力测试 & 连续 1\,h、撞击、遮挡、对手干扰 & 不崩、不持续漂移、可自动恢复 \\
    \bottomrule
  \end{tabular}
\end{table}

\section{常见坑}\label{sec:29-pitfalls}

\begin{itemize}
  \item \textbf{动态残影}：对手机器人被写进先验地图，之后每次定位都在"幻影"附近拉扯。建图时必须保证场地无人，或事后人工清理动态区域。
  \item \textbf{地图过期}：场地布置（基地、能量机关、围栏位置）每场都可能改。带着上一场的地图下场，重定位会稳定地匹配到错误位置——而且"看起来很稳定"，最难发现。
  \item \textbf{UWB 多径}：金属结构、装甲板、地面反射会产生非视距测量，表现为偶发的几十厘米跳变。要加抗差（Huber/卡方检验）并做限幅融合，不要直接当观测用。
  \item \textbf{对称导致的误重定位}：场地镜像结构使全局匹配出现双解。解决办法是用不对称特征（基地、飞镖架、观众席）做优先级加权，或把上一帧位置作为先验。
  \item \textbf{算力不足}：哨兵/步兵上同时跑 LIO、导航、视觉识别，CPU 一满，LIO 首先掉帧。要在集成阶段就做资源预算，而不是比赛前才发现。
  \item \textbf{云台转动被误当成底盘运动}：雷达装在云台上时必须用云台角做补偿，否则底盘定位会被云台转动"带偏"。
  \item \textbf{弹丸与撞击}：弹丸命中雷达或车体会造成瞬时点云爆点、IMU 饱和。要有饱和检测与单帧异常剔除。
\end{itemize}

\begin{summary}

\begin{enumerate}
  \item RM 赛场的关键约束：高度对称、动态障碍多、光照剧变、算力受限、比赛期无法调参；这些决定了必须做异构冗余而不是押注单一方案。
  \item 推荐方案：3D LIO 为主、先验地图全局定位校正、UWB 兜底、雷达站有条件就上；轮式只作短时补充。
  \item 哨兵、步兵、雷达站的算力、动态性与风险差异很大，不能共用一套参数；哨兵的核心是云台与底盘解耦。
  \item 标定顺序从底盘往外扩：底盘 $\to$ 雷达 $\to$ 云台 $\to$ 相机内参 $\to$ 云台-相机 $\to$ UWB；上车前按 10 条检查表逐项确认。
  \item 备赛分六步走（建图 $\to$ 清理 $\to$ 地图包 $\to$ 定位验证 $\to$ 导航联调 $\to$ 压力测试），每步都有量化验收标准。
  \item 最容易踩的坑：动态残影、地图过期、UWB 多径、对称误匹配、算力不足、云台转动干扰、弹丸撞击。
\end{enumerate}


\end{summary}

\begin{reading}

\begin{itemize}
  \item 第~\ref{ch:23-zlocalization} 章的 Z-Localization 深度解析与官方仓库 \cite{zloc2026}：先验地图定位在 RM 场景的完整工程实现，本章"方案选型"里"先验地图 + 全局定位"一栏就是它。
  \item 第~\ref{ch:22-nav2-practice} 章的完整链路实战：建图到导航联调的具体操作，本章的验收标准可以在那里逐项执行。
  \item 第~\ref{ch:17-advanced-lio} 章：FAST-LIO2 与 Point-LIO 的内部机制，理解它们才能判断"退化场景该加什么约束"。
\end{itemize}


\end{reading}

\chapter{学习路线与资源清单}\label{ch:30-roadmap}

\chapterintro{这本书很厚，没人能一周读完。更现实的问题是：一个刚进组的新队员，应该按什么顺序读、读到哪里算"够用"、什么时候该去碰真代码。本章给出四条时间尺度的路线（0--1 月、1--3 月、3--6 月、6--12 月），每条都配"读哪几章、跑哪份代码、做成什么事、怎么算过关"；然后给出开源代码的阅读顺序与每份代码"先看哪个文件"；最后是书单、工具链与实践建议。这一章不是知识，是路线图——它可以被反复翻，每次看都能确认自己走到了哪。}

\section{分阶段路线}\label{sec:30-stages}

\subsection{第 0--1 月：把环境与概念立起来}

目标不是"懂 SLAM"，而是\textbf{能独立编译、跑通、录包、看数据}。

\begin{table}[htbp]
  \centering
  \caption{第 0--1 月：环境与概念}\label{tab:30-s1}
  \small
  \begin{tabular}{lp{3.4cm}p{3.6cm}p{3.6cm}}
    \toprule
    项 & 内容 & 验收标准 \\
    \midrule
    读书 & 第~\ref{ch:01-overview}、\ref{ch:02-math}、\ref{ch:03-sensors} 章 & 能画出定位/建图/导航的整体框图 \\
    概念 & 里程计、SLAM、TF、QoS、外参 & 能口头解释每个词解决什么问题 \\
    环境 & Ubuntu + ROS 2 + 编译工具链 & 独立从零装好并跑通 demo \\
    代码 & 跑通 FAST-LIO2 的示例 bag & 能录自己的 bag 并回放 \\
    工具 & \lstinline{ros2 topic/tf2/bag}、RViz2 & 能查出"数据没来"是哪一环 \\
    \bottomrule
  \end{tabular}
\end{table}

这个月最容易走的弯路是"一开始就啃公式"。先让数据在屏幕上动起来，再回头看公式，效率高得多。

\subsection{第 1--3 月：里程计与 IMU}

\begin{table}[htbp]
  \centering
  \caption{第 1--3 月：里程计、IMU 与 LIO 入门}\label{tab:30-s2}
  \small
  \begin{tabular}{lp{3.4cm}p{3.6cm}p{3.6cm}}
    \toprule
    项 & 内容 & 验收标准 \\
    \midrule
    读书 & 第~\ref{ch:04-odom-intro}--\ref{ch:07-lidar-odom} 章、第~\ref{ch:16-lio-evolution} 章 & 能推导位姿递推与 ICP 残差 \\
    理论 & 四元数、ESKF、预积分、扫描匹配 & 能手推一维卡尔曼更新 \\
    代码 & KISS-ICP、FAST-LIO2 & 能改一行代码并解释影响 \\
    实验 & 同一段 bag 上比较轮式/激光/IMU 轨迹 & 能画出三条轨迹与误差 \\
    标定 & 第~\ref{ch:26-calibration} 章的 IMU 与外参部分 & 独立完成一次 IMU 零偏与 Allan 分析 \\
    \bottomrule
  \end{tabular}
\end{table}

\subsection{第 3--6 月：SLAM 与导航}

\begin{table}[htbp]
  \centering
  \caption{第 3--6 月：SLAM 后端与导航链路}\label{tab:30-s3}
  \small
  \begin{tabular}{lp{3.4cm}p{3.6cm}p{3.6cm}}
    \toprule
    项 & 内容 & 验收标准 \\
    \midrule
    读书 & 第~\ref{ch:11-slam-framework}--\ref{ch:17-advanced-lio} 章 & 能解释前端/后端/回环各管什么 \\
    读书 & 第~\ref{ch:18-nav-overview}--\ref{ch:23-zlocalization} 章 & 能画出 Nav2 的节点与数据流 \\
    代码 & Cartographer、Point-LIO & 读懂 \lstinline{submap} 与逐点更新 \\
    实验 & 建图 $\to$ 保存 $\to$ 定位 $\to$ 导航全链路 & 让车在自己建的地图上走完一条路径 \\
    工程 & 第~\ref{ch:25-ros2}、\ref{ch:27-tuning} 章 & 能独立排掉一次"定位跳变" \\
    \bottomrule
  \end{tabular}
\end{table}

\subsection{第 6--12 月：工程化与独立思考}

\begin{table}[htbp]
  \centering
  \caption{第 6--12 月：工程化、评测与改进}\label{tab:30-s4}
  \small
  \begin{tabular}{lp{3.4cm}p{3.6cm}p{3.6cm}}
    \toprule
    项 & 内容 & 验收标准 \\
    \midrule
    读书 & 第~\ref{ch:24-frontier}、\ref{ch:28-evaluation} 章与本书全部 & 能判断一篇新论文值不值得跟进 \\
    评测 & evo、ATE/RPE、消融实验 & 能设计并执行一次可复现对比 \\
    改进 & 在某个模块上做一处真实改进 & 公开数据集上指标提升且可复现 \\
    落地 & 第~\ref{ch:29-rm} 章的清单 & 在实车上稳定运行一整场比赛 \\
    输出 & 写文档、做分享 & 能把一个模块讲给新队员听懂 \\
    \bottomrule
  \end{tabular}
\end{table}

\section{代码阅读顺序}\label{sec:30-code}

读代码的顺序应该\textbf{由简到繁}，每份代码只读它的"核心那几百行"。表~\ref{tab:30-code} 给出建议。

\begin{table}[htbp]
  \centering
  \caption{开源代码阅读顺序与入口文件}\label{tab:30-code}
  \footnotesize
  \begin{tabular}{p{2.5cm}p{4.4cm}p{6.6cm}}
    \toprule
    顺序 & 仓库 & 先看什么 \\
    \midrule
    1 & KISS-ICP \cite{vizzo2023kissicp} & \lstinline{src/kiss_icp} 的配准与自适应阈值；这是最小可用的激光里程计，一天能读完 \\
    2 & FAST-LIO2 \cite{xu2022fastlio2} & \lstinline{include/ikd-Tree} 与 \lstinline{src/laserMapping.cpp}：先看"点如何进地图"，再看 IEKF 迭代 \\
    3 & Point-LIO \cite{he2023pointlio} & 逐点更新的主循环与 IMU 饱和处理，理解"为什么它能扛高动态" \\
    4 & Cartographer \cite{hess2016cartographer} & \lstinline{cartographer/mapping} 下的 \lstinline{submap} 与 \lstinline{constraint_builder}：分支定界回环 \\
    5 & Nav2 \cite{macenski2020nav2} & 从 \lstinline{bt_navigator} 的行为树 XML 入口，依次看 planner 与 controller 插件接口 \\
    6 & Z-Localization \cite{zloc2026} & 先验地图定位的完整工程：从地图加载、全局匹配到 \lstinline{map}$\to$\lstinline{odom} 修正 \\
    \bottomrule
  \end{tabular}
\end{table}

读代码的三个技巧：

\begin{enumerate}
  \item \textbf{先画数据流，再读函数}。每个仓库先写出"输入话题 $\to$ 主要处理函数 $\to$ 输出话题"，再挑中间最核心的一个函数精读。上来从 \lstinline{main} 一行行读到底，必然迷路。
  \item \textbf{带着问题读}。"它的地图怎么存？""重定位失败会怎样？""时间戳在哪里对齐？"带着三个具体问题去读，比通读效率高一个量级。
  \item \textbf{改一行，跑一遍}。把降采样体素改大一倍，看轨迹怎么变；把迭代次数改成 1，看什么时候炸。动手是最快的理解方式。
\end{enumerate}

\section{书、课程与文档}\label{sec:30-books}

\begin{itemize}
  \item \textbf{概率机器人}：Thrun 等的 \emph{Probabilistic Robotics} \cite{thrun2005probabilistic}。定位与 SLAM 的概率视角经典，重点读第 2、3、6、7、8 章，配合第~\ref{ch:02-math} 章。
  \item \textbf{状态估计}：Barfoot 的 \emph{State Estimation for Robotics} \cite{barfoot2017state}。李群、SE(3) 上的估计讲得最系统，是读懂 FAST-LIO 的钥匙。
  \item \textbf{ESKF 手册}：Sola 的 \emph{Quaternion Kinematics for the Error-State Kalman Filter} \cite{sola2017eskf}。免费且短，建议打印精读。
  \item \textbf{视觉 SLAM}：高翔等的《视觉 SLAM 十四讲》。若要做视觉/视觉惯性，这是中文最好的入门书。
  \item \textbf{课程}：深蓝学院的多传感器融合定位、激光 SLAM 课程体系完整，作业即代码；Coursera 的 Robotics 专项适合补基础。
  \item \textbf{文档}：ROS 2 官方文档、Nav2 文档、各仓库 README 与 issue 区。\textbf{issue 区是被低估的教材}——你能想到的问题，九成已经有人问过并得到了作者解释。
\end{itemize}

\section{工具链清单}\label{sec:30-tools}

\begin{table}[htbp]
  \centering
  \caption{定位与建图常用工具链}\label{tab:30-tools}
  \small
  \begin{tabular}{llp{6.4cm}}
    \toprule
    类别 & 工具 & 用在哪 \\
    \midrule
    系统 & Ubuntu 22.04、ROS 2 Humble & 全部工程 \\
    语言 & C++17、Python 3 & 算法与脚本 \\
    数学 & Eigen、Sophus & 矩阵、李群 \\
    点云 & PCL、Open3D & 点云处理与可视化 \\
    优化 & Ceres Solver、GTSAM & 后端与标定 \\
    视觉 & OpenCV、Kalibr & 内参与相机-IMU 标定 \\
    可视化 & RViz2、Foxglove、CloudCompare & 调试与地图检查 \\
    评测 & evo & ATE/RPE 与轨迹对比 \\
    录包 & \lstinline{ros2 bag}（mcap） & 数据采集与复现 \\
    \bottomrule
  \end{tabular}
\end{table}

\section{给新队员的 10 条建议}\label{sec:30-advice}

\begin{enumerate}
  \item \textbf{先跑通，再改代码}。跑不通的东西，读再多论文也没用。
  \item \textbf{先看数据，再看代码}。点云长什么样、时间戳差多少、TF 对不对，比读实现更快定位问题。
  \item \textbf{记录每一次实验}。参数、指标、现象三列，坚持三个月你就比同龄人快一倍。
  \item \textbf{学会看 TF}。定位问题里有一半最终能追溯到某条 TF 断链或方向反了。
  \item \textbf{重视标定}。外参与时间戳是"调不回来"的误差，第~\ref{ch:26-calibration} 章请至少精读一遍。
  \item \textbf{从失败案例学习}。录下"炸掉"的那段数据，比录十段成功数据更有价值。
  \item \textbf{多读 issue 与提交历史}。作者的修复往往比文档更接近真相。
  \item \textbf{一次只改一个变量}。无论调参还是做实验，这条纪律都成立。
  \item \textbf{写文档}。能写清楚一个模块，说明你真的懂了；这也是团队里最稀缺的能力。
  \item \textbf{不要怕数学，但要按需学}。遇到不懂的公式，先解决眼前问题，再回头补推导。
\end{enumerate}

\section{结语}\label{sec:30-end}

写到这里，全书的主线其实只有一句话：\textbf{定位是"用带噪声的观测去估计一个不确定的位姿，并承认它有多不确定"}。里程计篇讲如何递推，SLAM 篇讲如何用历史约束修正，导航篇讲如何用这个位姿去做决策，工程篇讲如何让它在一台真实的车上稳定运行。四条线索在最后一章汇合——你会发现，前面觉得"用不上"的数学，最后都变成了排错时判断"这到底是参数问题还是数据问题"的依据。

技术会变：FAST-LIO 之后有 Point-LIO，Cartographer 之后有新的图优化框架，Nav2 也在持续演进。但"观测—模型—优化—验证"这套方法论不会变。希望这本书不只是让你会用某几个仓库，而是让你在面对一个从没见过的传感器或算法时，知道该问哪几个问题。

\begin{summary}

\begin{enumerate}
  \item 学习路线分四段：0--1 月环境与概念、1--3 月里程计与 IMU、3--6 月 SLAM 与导航、6--12 月工程化与改进，每段都有"读什么、跑什么、做成什么、怎么过关"。
  \item 代码阅读由简到繁：KISS-ICP $\to$ FAST-LIO2 $\to$ Point-LIO $\to$ Cartographer $\to$ Nav2 $\to$ Z-Localization；先画数据流，再精读核心函数。
  \item 推荐的"少而精"资源：Probabilistic Robotics、State Estimation for Robotics、Sola 的 ESKF 文档、视觉 SLAM 十四讲，以及各仓库的 issue 区。
  \item 工具链的核心是"能复现"：统一环境、固定 bag、用 evo 量化，任何改动都必须能回到同一段数据上验证。
  \item 十条建议的核心可以压缩成三条：先跑通再改、一次只改一个变量、写文档。
  \item 全书主线：观测—模型—优化—验证；技术会变，方法论不会。
\end{enumerate}


\end{summary}

\begin{reading}

\begin{itemize}
  \item 第~\ref{ch:01-overview} 章与书末附录 A、B：全书的入口与速查工具，建议在按本章路线学习时反复回看。
  \item 各开源仓库的 README 与 issue 区：KISS-ICP、FAST-LIO2、Point-LIO、Cartographer、Nav2、Z-Localization \cite{zloc2026}。遇到具体报错时先搜 issue。
  \item \href{https://github.com/gaoxiang12/slambook2}{视觉 SLAM 十四讲}配套代码：若要做视觉惯性方向，这是最好的动手起点。
\end{itemize}


\end{reading}


\printbibliography[heading=bibintoc, title=\ebibname]

\appendix
\chapter{公式与符号速查}

\section*{A.1 全局符号表}

\begin{table}[htbp]
  \centering
  \caption{全书通用符号}\label{tab:a-symbols}
  \small
  \begin{tabular}{llp{7.2cm}}
    \toprule
    符号 & 名称 & 含义 \\
    \midrule
    $W$ & 世界系 / map 系 & 全局参考坐标系，$z$ 轴向上 \\
    $O$ & odom 系 & 连续、无跳变的局部参考系 \\
    $B$ & 机体系 & 默认与 IMU 系重合 \\
    $L$、$C$ & 雷达系、相机系 & 传感器坐标系 \\
    $T_{AB}=(R_{AB},{}^{A}p_B)$ & 位姿 & 从 $B$ 系到 $A$ 系的变换 \\
    $R$、$q$ & 旋转矩阵、四元数 & $R\in\SO(3)$ \\
    $p$、$v$、$a$ & 位置、速度、加速度 & 列向量 \\
    $\omega$、$g$ & 角速度、重力 & $g=[0,0,-9.81]^{\top}$ \\
    $\tilde{\omega}$、$\tilde{a}$ & IMU 测量值 & 含零偏与噪声 \\
    $b_g$、$b_a$ & 陀螺仪/加速度计零偏 & 缓慢漂移 \\
    $x$、$\delta x$ & 状态、误差状态 & \\
    $z$、$h(\cdot)$、$r$ & 测量、观测函数、残差 & \\
    $P$、$Q$、$R$ & 状态/过程/测量协方差 & \\
    $J$、$H$ & 雅可比、海森 & \\
    $\Lambda$ & 信息矩阵 & $\Lambda=P^{-1}$ 或 $J^{\top}R^{-1}J$ \\
    $\pi=(n_\pi,d)$ & 平面参数 & $n_\pi^{\top}x=d$ \\
    \bottomrule
  \end{tabular}
\end{table}

\section*{A.2 旋转与刚体变换}

\paragraph{旋转矩阵的性质}
\begin{equation}
R^{\top}R=I,\quad \det(R)=1,\quad (R_1R_2)^{-1}=R_2^{\top}R_1^{\top}.
\end{equation}

\paragraph{Rodrigues 公式}
设 $\phi=\theta a$（$\norm{a}=1$），则
\begin{equation}
\exp(\phi^{\wedge})=I+\sin\theta\,a^{\wedge}+(1-\cos\theta)(a^{\wedge})^{2}.
\end{equation}

\paragraph{对数映射}
\begin{equation}
\theta=\arccos\frac{\tr(R)-1}{2},\qquad
\phi=\frac{\theta}{2\sin\theta}(R-R^{\top})^{\vee}.
\end{equation}

\paragraph{$\SE(3)$ 的逆}
\begin{equation}
T^{-1}=\begin{bmatrix}R^{\top}&-R^{\top}t\\0^{\top}&1\end{bmatrix}.
\end{equation}

\paragraph{位姿复合与逆}
\begin{equation}
T_{AC}=T_{AB}T_{BC},\qquad
T_{AB}^{-1}=T_{BA}.
\end{equation}

\paragraph{右扰动}
\begin{equation}
R\leftarrow R\exp(\delta\phi^{\wedge}),\qquad t\leftarrow t+R\,\delta\rho.
\end{equation}

\section*{A.3 概率与滤波}

\paragraph{贝叶斯公式}
\begin{equation}
p(x\mid z)=\frac{p(z\mid x)p(x)}{p(z)}.
\end{equation}

\paragraph{高斯分布}
\begin{equation}
p(x)=\frac{1}{(2\pi)^{n/2}\abs{\Sigma}^{1/2}}
\exp\left(-\frac12(x-\mu)^{\top}\Sigma^{-1}(x-\mu)\right).
\end{equation}

\paragraph{方差与协方差}
\begin{equation}
\Var(X)=\E[X^{2}]-(\E[X])^{2},\qquad
\Cov(X,Y)=\E[XY]-\E[X]\E[Y].
\end{equation}

\paragraph{卡尔曼滤波}
预测：
\begin{align}
\hat{x}_{k}^{-}&=F_k\hat{x}_{k-1}+B_ku_k,\\
P_k^{-}&=F_kP_{k-1}F_k^{\top}+Q_k.
\end{align}
更新：
\begin{align}
K_k&=P_k^{-}H_k^{\top}\left(H_kP_k^{-}H_k^{\top}+R_k\right)^{-1},\\
\hat{x}_k&=\hat{x}_k^{-}+K_k\left(z_k-h(\hat{x}_k^{-})\right),\\
P_k&=(I-K_kH_k)P_k^{-}.
\end{align}

\paragraph{ESKF}
误差状态传播与更新：
\begin{equation}
\delta\dot{x}=F\delta x+Gn,\qquad
K=P^{-}H^{\top}(HP^{-}H^{\top}+R)^{-1},\qquad
\delta x\leftarrow\delta x\oplus K r.
\end{equation}

\paragraph{马氏距离与卡方检验}
\begin{equation}
d^{2}=r^{\top}(HPH^{\top}+R)^{-1}r\sim\chi^{2}(k).
\end{equation}

\section*{A.4 优化}

\paragraph{最小二乘正规方程}
\begin{equation}
A^{\top}Ax=A^{\top}b.
\end{equation}

\paragraph{高斯--牛顿}
\begin{equation}
J^{\top}J\,\delta x=-J^{\top}r,\qquad x\leftarrow x+\delta x.
\end{equation}

\paragraph{Levenberg--Marquardt}
\begin{equation}
\left(J^{\top}J+\lambda I\right)\delta x=-J^{\top}r.
\end{equation}

\paragraph{Huber 鲁棒核}
\begin{equation}
\rho(s)=
\begin{cases}
s, & s\le\delta^{2},\\
2\delta\sqrt{s}-\delta^{2}, & s>\delta^{2},
\end{cases}
\qquad
w(s)=\frac{\rho'(s)}{2\sqrt{s}}=
\begin{cases}
1, & s\le\delta^{2},\\
\delta/\sqrt{s}, & s>\delta^{2}.
\end{cases}
\end{equation}

\section*{A.5 点云配准}

\paragraph{ICP 闭式旋转解}
给定去质心点对 $(p_i,q_i)$，构造
\begin{equation}
W=\sum_i p_iq_i^{\top}=U\Sigma V^{\top},\qquad
R^{*}=VU^{\top},
\end{equation}
若 $\det(R^{*})<0$，取 $V'=[v_1,v_2,-v_3]$ 后重算。平移为 $t^{*}=\bar{q}-R^{*}\bar{p}$。

\paragraph{点到面残差}
\begin{equation}
r_i=n_i^{\top}\left(Rp_i+t-q_i\right).
\end{equation}

\section*{A.6 SLAM 后端}

\paragraph{位姿图残差}
\begin{equation}
r_{ij}=\log\left(T_{ij}^{-1}T_i^{-1}T_j\right)^{\vee}\in\R^{6}.
\end{equation}

\paragraph{舒尔补}
分块系统
\begin{equation}
\begin{bmatrix}H_{cc}&H_{cp}\\H_{pc}&H_{pp}\end{bmatrix}
\begin{bmatrix}\delta x_c\\\delta x_p\end{bmatrix}
=
\begin{bmatrix}b_c\\b_p\end{bmatrix}
\end{equation}
消去相机变量后：
\begin{equation}
\left(H_{pp}-H_{pc}H_{cc}^{-1}H_{cp}\right)\delta x_p
=b_p-H_{pc}H_{cc}^{-1}b_c.
\end{equation}

\paragraph{边缘化}
用舒尔补把旧状态消去，得到新的先验因子；注意这会引入线性化误差。

\section*{A.7 导航}

\paragraph{A* 评价函数}
\begin{equation}
f(n)=g(n)+h(n),
\end{equation}
其中 $g$ 是起点到 $n$ 的实际代价，$h$ 是到目标的启发式下界（可采纳）。

\paragraph{DWA 评价函数}
\begin{equation}
G(v,\omega)=\alpha\cdot\mathrm{heading}(v,\omega)+\beta\cdot\mathrm{clearance}(v,\omega)+\gamma\cdot\mathrm{velocity}(v,\omega).
\end{equation}

\paragraph{Cartographer 扫描匹配得分}
\begin{equation}
S(T)=\sum_k M_{\mathrm{smooth}}\left(T\,h_k\right),
\end{equation}
其中 $M_{\mathrm{smooth}}$ 是双线性插值后的概率栅格。

\paragraph{控制代价（MPC）}
\begin{equation}
\min_{u_{0:N-1}}\ \sum_{k=0}^{N-1}\left(x_k-x_k^{\mathrm{ref}}\right)^{\top}Q\left(x_k-x_k^{\mathrm{ref}}\right)+u_k^{\top}Ru_k
\quad\text{s.t.}\ x_{k+1}=f(x_k,u_k).
\end{equation}

\chapter{术语中英对照表}

\noindent 本表按主题分组，给出本书出现的主要术语中英文对照及一句话解释，便于阅读英文论文与源码时查阅。

\section*{B.1 数学与估计}

\begin{table}[htbp]
  \centering
  \caption{数学与估计基础术语}\label{tab:b-math}
  \footnotesize
  \begin{tabular}{llp{6.6cm}}
    \toprule
    中文 & 英文 & 说明 \\
    \midrule
    位姿 & pose & 位置 + 姿态，$\SE(3)$ 元素 \\
    姿态 & orientation & 旋转部分，$\SO(3)$ 元素 \\
    刚体变换 & rigid transformation & 旋转 + 平移 \\
    齐次坐标 & homogeneous coordinates & 用 4 维向量表示 3 维点，便于矩阵复合 \\
    旋转矩阵 & rotation matrix & 满足 $R^{\top}R=I$、$\det R=1$ \\
    四元数 & quaternion & 用 4 个数表示旋转，无奇异 \\
    欧拉角 & Euler angle & roll/pitch/yaw，直观但有万向锁 \\
    万向锁 & gimbal lock & pitch 为 $\pm90^\circ$ 时自由度退化 \\
    李群 & Lie group & 既是群又是流形，如 $\SO(3)$、$\SE(3)$ \\
    李代数 & Lie algebra & 李群单位元处的切空间，可做加法 \\
    指数映射 & exponential map & 李代数 $\to$ 李群 \\
    对数映射 & logarithm map & 李群 $\to$ 李代数 \\
    雅可比矩阵 & Jacobian & 一阶偏导矩阵 \\
    海森矩阵 & Hessian & 二阶偏导矩阵 \\
    泰勒展开 & Taylor expansion & 用多项式局部逼近函数 \\
    最小二乘 & least squares & 最小化残差平方和 \\
    正规方程 & normal equation & $A^{\top}Ax=A^{\top}b$ \\
    高斯--牛顿 & Gauss--Newton & 非线性最小二乘的迭代解法 \\
    LM & Levenberg--Marquardt & 带阻尼的高斯--牛顿 \\
    期望 & expectation & 概率加权平均 \\
    方差 & variance & 偏离期望的平方的期望 \\
    协方差 & covariance & 两个变量共同变化的程度 \\
    协方差矩阵 & covariance matrix & 所有两两协方差组成的矩阵 \\
    高斯分布 & Gaussian distribution & 正态分布 \\
    贝叶斯公式 & Bayes' formula & 先验 + 似然 $\to$ 后验 \\
    后验 & posterior & 观测后的状态分布 \\
    似然 & likelihood & 给定状态下观测的概率 \\
    卡方分布 & chi-square distribution & 标准正态平方和；用于一致性检验 \\
    马氏距离 & Mahalanobis distance & 按协方差加权的距离 \\
    信息矩阵 & information matrix & 协方差矩阵的逆，衡量约束强度 \\
    特征值 & eigenvalue & $Av=\lambda v$ 中的 $\lambda$ \\
    奇异值分解 & SVD & $A=U\Sigma V^{\top}$ \\
    \bottomrule
  \end{tabular}
\end{table}

\section*{B.2 定位、里程计与 SLAM}

\begin{table}[htbp]
  \centering
  \caption{定位与 SLAM 术语}\label{tab:b-slam}
  \footnotesize
  \begin{tabular}{llp{6.6cm}}
    \toprule
    中文 & 英文 & 说明 \\
    \midrule
    定位 & localization & 估计机器人在参考系中的位姿 \\
    绝对定位 & absolute localization & 直接在地图系中定位 \\
    相对定位 & relative localization & 相对起点递推，会累积误差 \\
    全局定位 & global localization & 初始位姿未知时的定位 \\
    重定位 & relocalization & 跟丢后重新找回位姿 \\
    绑架问题 & kidnapped robot problem & 机器人被瞬间移动后的定位问题 \\
    里程计 & odometry & 用传感器增量递推位姿 \\
    航迹推算 & dead reckoning & 基于速度/转向积分推位姿 \\
    漂移 & drift & 累积误差导致的缓慢偏离 \\
    前端 & front-end & 数据关联 + 局部位姿估计 \\
    后端 & back-end & 全局优化、回环、地图修正 \\
    数据关联 & data association & 确定观测之间的对应关系 \\
    扫描匹配 & scan matching & 用两帧点云估计相对位姿 \\
    点云配准 & point cloud registration & 同上，更通用的说法 \\
    迭代最近点 & ICP & 交替找最近点与求最优变换 \\
    正态分布变换 & NDT & 用体素内高斯分布建模点云 \\
    关键帧 & keyframe & 被选入地图的代表帧 \\
    局部地图 & local map & 当前帧附近的地图子集 \\
    先验地图 & prior map & 离线建好、定位时加载的地图 \\
    回环检测 & loop closure detection & 判断是否回到历史地点 \\
    位姿图 & pose graph & 节点=位姿、边=相对约束的图 \\
    因子图 & factor graph & 用因子表示观测的图模型 \\
    光束法平差 & bundle adjustment & 同时优化位姿与路标的非线性最小二乘 \\
    舒尔补 & Schur complement & 分块消元，BA 的核心技巧 \\
    边缘化 & marginalization & 消去旧状态但保留其信息 \\
    鲁棒核 & robust kernel & 降低大残差影响的损失函数 \\
    后端优化 & back-end optimization & 全局最小二乘求解 \\
    建图 & mapping & 把观测融合成环境表示 \\
    占据栅格 & occupancy grid & 每格记录占据概率的 2D 地图 \\
    八叉树 & octree & 空间递归划分的 3D 地图 \\
    点云 & point cloud & 3D 点的集合 \\
    体素 & voxel & 三维栅格单元 \\
    降采样 & downsampling & 减少点数，常用体素质心法 \\
    去畸变 & deskewing & 补偿单帧点云内的运动畸变 \\
    运动畸变 & motion distortion & 扫描期间运动导致的点云变形 \\
    可观测性 & observability & 某方向能否由观测唯一确定 \\
    退化 & degeneracy & 观测对某方向没有约束 \\
    紧耦合 & tightly coupled & 多传感器放进同一估计器 \\
    松耦合 & loosely coupled & 各传感器独立估计后融合 \\
    预积分 & preintegration & IMU 在两帧间的积分量与初值解耦 \\
    零偏 & bias & 传感器零输入时的非零输出 \\
    随机游走 & random walk & 零偏等量随时间的随机漂移 \\
    惯性测量单元 & IMU & 陀螺仪 + 加速度计 \\
    激光惯性里程计 & LIO & LiDAR + IMU 里程计 \\
    视觉惯性里程计 & VIO & 相机 + IMU 里程计 \\
    \bottomrule
  \end{tabular}
\end{table}

\section*{B.3 导航、工程与其他}

\begin{table}[htbp]
  \centering
  \caption{导航与工程术语}\label{tab:b-nav}
  \footnotesize
  \begin{tabular}{llp{6.6cm}}
    \toprule
    中文 & 英文 & 说明 \\
    \midrule
    导航 & navigation & 定位 + 地图 + 规划 + 控制 \\
    全局规划 & global planning & 在全局地图上找可行路径 \\
    局部规划 & local planning & 跟踪全局路径并实时避障 \\
    代价地图 & costmap & 栅格化的代价表示，供规划查询 \\
    膨胀层 & inflation layer & 按距离给障碍物周围加代价 \\
    行为树 & behavior tree & Nav2 的任务编排机制 \\
    生命周期 & lifecycle & ROS 2 节点的状态管理 \\
    动态窗口法 & DWA & 在速度空间采样避障 \\
    时间弹性带 & TEB & 把轨迹表示为带时间戳的图优化 \\
    模型预测控制 & MPC & 滚动时域优化控制 \\
    人工势场 & artificial potential field & 引力 + 斥力避障 \\
    可视图 & visibility graph & 连接可见节点的图，用于全局规划 \\
    自主探索 & autonomous exploration & 在未知环境中边建图边找目标 \\
    地形分析 & terrain analysis & 从点云提取坡度、粗糙度、台阶 \\
    可通行性 & traversability & 某区域能否安全通过 \\
    探索增益 & exploration gain & 探索某区域能获得的新信息量 \\
    外参 & extrinsic & 传感器之间的位姿变换 \\
    内参 & intrinsic & 相机内部成像参数 \\
    标定 & calibration & 估计内参/外参/时间偏移 \\
    时间同步 & time synchronization & 统一多传感器时间基准 \\
    时间戳 & timestamp & 数据采集时刻 \\
    点频 & point rate & 激光雷达每秒输出点数 \\
    线数 & number of channels & 激光雷达垂直方向束数 \\
    角分辨率 & angular resolution & 相邻光束的角度间隔 \\
    盲区 & blind zone & 雷达无法测距的最近距离 \\
    视场角 & field of view (FOV) & 传感器可观测的角度范围 \\
    飞行时间 & time of flight (ToF) & 通过光/声飞行时间测距 \\
    超宽带 & UWB & 室内无线定位技术 \\
    全球导航卫星系统 & GNSS & GPS/北斗等的统称 \\
    实时动态差分 & RTK & GNSS 厘米级定位技术 \\
    惯性测量单元 & IMU & 陀螺仪 + 加速度计 \\
    机器人操作系统 & ROS & 机器人软件框架 \\
    数据分发服务 & DDS & ROS 2 的通信中间件 \\
    服务质量 & QoS & ROS 2 通信的可靠性等配置 \\
    坐标变换 & TF & ROS 中的坐标系变换树 \\
    话题 & topic & 发布/订阅通信 \\
    服务 & service & 请求/响应通信 \\
    动作 & action & 带反馈的长任务通信 \\
    数据包 & rosbag & 传感器数据的录制文件 \\
    绝对轨迹误差 & ATE & 估计轨迹与真值的整体偏差 \\
    相对位姿误差 & RPE & 固定间隔内的相对误差 \\
    \bottomrule
  \end{tabular}
\end{table}


\end{document}
